\documentclass[10pt,oneside]{book}

\usepackage{fontspec}
\setmainfont{Malgun Gothic}[
  ItalicFont={Malgun Gothic},
  ItalicFeatures={FakeSlant=0.2},
  BoldItalicFont={Malgun Gothic Bold},
  BoldItalicFeatures={FakeSlant=0.2}
]
\setsansfont{Malgun Gothic}
\XeTeXlinebreaklocale "ko"
\XeTeXlinebreakskip=0pt plus 1pt
\usepackage{amsmath,amssymb,mathrsfs}
\usepackage[all]{xy}
\usepackage{tikz-cd}
\usepackage[a4paper,margin=22mm]{geometry}
\usepackage[
  colorlinks=true,
  linkcolor=blue,
  hypertexnames=false,
  unicode=true,
  pdfencoding=auto
]{hyperref}
\hypersetup{
  pdftitle={대수기하학 원론 -- 한국어 누적판},
  pdfauthor={알렉산더 그로텐디크와 장 디외도네},
  pdfsubject={NUMDAM 프랑스어 권위본에 따른 EGA 한국어 누적 번역}
}

\renewcommand{\thesection}{\arabic{section}}
\renewcommand{\thesubsection}{\thesection.\arabic{subsection}}
\renewcommand{\bibname}{참고문헌}
\renewcommand{\contentsname}{목차}
\makeatletter
\renewcommand{\@pnumwidth}{2em}
\renewcommand{\@tocrmarg}{3em}
\makeatother
\newcommand{\oldpage}[2][]{%
  \setcounter{footnote}{0}%
  \marginpar{\scriptsize\textbf{#1}~|~#2}\ignorespaces}
\newcommand{\agref}[2]{%
  \ifcsname r@#1\endcsname\hyperref[#1]{#2}\else#2\fi}
\newenvironment{env}[1][]{\par\medskip\noindent\textbf{(#1)}\ }{\par}
\newenvironment{definition}[1][]{%
  \par\medskip\noindent\itshape 정의 (#1). ---\ }{\par}
\newenvironment{proposition}[1][]{%
  \par\medskip\noindent\itshape 명제 (#1). ---\ }{\par}
\newenvironment{theorem}[1][]{%
  \par\medskip\noindent\itshape 정리 (#1). ---\ }{\par}
\newenvironment{corollary}[1][]{%
  \par\medskip\noindent\itshape 따름정리 (#1). ---\ }{\par}
\newenvironment{remark}[1][]{%
  \par\medskip\noindent\itshape 주 (#1). ---\ }{\par}
\newenvironment{lemma}[1][]{%
  \par\medskip\noindent\itshape 보조정리 (#1). ---\ }{\par}
\newenvironment{example}[1][]{%
  \par\medskip\noindent\textit{예 (#1). ---\ }}{\par}
\newenvironment{examples}[1][]{%
  \par\medskip\noindent\textit{예들 (#1). ---\ }}{\par}
\newenvironment{notation}[1][]{%
  \par\medskip\noindent\textit{표기 (#1). ---\ }}{\par}
\newenvironment{proof}{%
  \par\medskip\noindent\textit{증명. ---\ }}{\par}
\newenvironment{namedtheorem}[2][]{%
  \par\medskip\noindent\itshape #2 (#1). ---\ }{\par}
\newenvironment{namedlemma}[2][]{%
  \par\medskip\noindent\itshape #2 (#1). ---\ }{\par}

\begin{document}
% Authority: EGA I frontmatter-fr.tex, SHA-256 D506B87684E2136E3F87495190EECDF40B79DB8887ED71093C4CD56648E282A9.
\begin{titlepage}
\centering

{\Large 프랑스 고등과학연구소\par}
\vspace{0.4em}
{\large INSTITUT DES HAUTES ÉTUDES SCIENTIFIQUES\par}

% 원 인쇄본에서는 이 자리에 연구소 문장이 놓여 있다.
\vfill

{\Large\bfseries 대수기하학 원론\par}
\vspace{1.5em}
{\large\itshape A. 그로텐디크 지음\par}
\vspace{0.7em}
{\itshape J. 디외도네의 협력으로 편집\par}

\vspace{2em}
I

\vspace{1.5em}
{\Large\bfseries 스킴의 언어\par}

\vspace{1.6em}
{\large 한국어 누적판\par}
\vspace{0.5em}
{\small 현재 범위: 서문 및 EGA $0_{\mathrm I}$ 제4.1.7절까지\par}

\vfill
{\Large 1960\par}
\vspace{1.2em}
{\Large PUBLICATIONS MATHÉMATIQUES, 제4호\par}
\vspace{0.5em}
{\large 5, ROND-POINT BUGEAUD --- PARIS (XVI\textsuperscript{e})\par}
\end{titlepage}

\thispagestyle{empty}
\begin{center}
{\Large\bfseries 한국어판 정보\par}
\end{center}

\vspace{1.5em}
이 책은 알렉산더 그로텐디크와 장 디외도네의
\emph{Éléments de géométrie algébrique}를 한국어로 옮기는 누적판이다.
이 정확한 공개판은 EGA I의 앞부분, 서론, 그리고 NUMDAM 기반 프랑스어 권위본
\texttt{ega0-1-fr.tex}의 제1행부터 제2313행, 곧 제4.1.7절 끝까지를 포함한다.
EGA 전체의 한국어 번역은 같은 DOI 계열과 저장소에서 끊김 없이 계속된다.

\medskip
\noindent\textbf{한국어판 안정 DOI:}
\href{https://doi.org/10.5281/zenodo.21921513}{10.5281/zenodo.21921513}

\smallskip
\noindent\textbf{이 정확한 판의 DOI:}
\href{https://doi.org/10.5281/zenodo.21921514}{10.5281/zenodo.21921514}

\smallskip
\noindent\textbf{공개 저장소:}
\href{https://github.com/KokunoYumeto/ega-ko}{github.com/KokunoYumeto/ega-ko}

\smallskip
\noindent\textbf{전역 EGA 자료실:}
\href{https://doi.org/10.5281/zenodo.20414353}{10.5281/zenodo.20414353}

\medskip
한국어 번역과 조판은 Interlanguage 프로젝트를 위하여 OpenAI Codex가 수행하였다.
편집 가능한 TeX, 프랑스어 원문의 정확한 식별자와 해시, 모든 번역\,·\,용어\,·\,조판 결정,
난점과 미해결 항목, 빌드 기록, 추출 검사, 렌더 검사 및 SHA-256 목록은 이 판의
출처 및 증거 묶음에 들어 있다. 이 판은 원문 동기화와 모델 기반 검사를 거친
누적 작업판이며, 비평판 또는 사람의 외부 인증을 주장하지 않는다.

\vfill
\noindent\textbf{언어:} 한국어 (ko; Zenodo: kor)\par
\noindent\textbf{판 식별자:} 2026-08-13 EGA-I-through-4.1.7\par
\clearpage

\thispagestyle{empty}
\vspace*{\fill}
\begin{center}
\fbox{%
  \begin{minipage}{0.50\textwidth}
  \centering
  납본

  \smallskip
  초판 \dotfill 1960년 제3사분기

  \smallskip
  모든 권리

  모든 국가에서 유보됨

  \smallskip
  \copyright\ 1960, \emph{Institut des Hautes Études Scientifiques}
  \end{minipage}%
}
\end{center}
\vspace*{\fill}
\clearpage

\tableofcontents
\clearpage
\markboth{}{}
% Authority: EGA I intro-fr.tex, SHA-256 CE78400CD9DBC36A4D11CC933B7BE18BEAC0C69AE6A04FBA3ECFA86053572980.
\section*{서론}

\begin{center}
\emph{오스카 자리스키와 앙드레 베유에게.}
\end{center}

\oldpage[I]{5}

이 논문과 뒤이어 나올 많은 논문은 대수기하학의 기초에 관한 하나의 논고를 이루기 위한 것이다.
원칙적으로 이 분야에 관한 어떤 특별한 지식도 전제하지 않는다. 오히려 그러한 지식은 분명한 장점이 있음에도 불구하고, 그것이 수반하는 쌍유리적 관점에 지나치게 익숙해진 탓에 여기서 제시하는 관점과 기법에 익숙해지려는 이에게 때로 방해가 될 수 있음이 드러났다.
반면 독자가 다음 주제들을 잘 알고 있다고 가정한다.

\begin{enumerate}
  \item[a)] \emph{가환대수학}. 예를 들어 현재 준비 중인 N.~부르바키의 \emph{원론} 여러 권에서 전개되는 내용이다(이 책들이 출간되기 전까지는 사뮈엘-자리스키~\cite{I-13}와 사뮈엘~\cite{I-11},~\cite{I-12}을 참조할 수 있다).
  \item[b)] \emph{호몰로지 대수학}. 카르탕-아일렌베르크~\cite{I-2}(이하 (M)), 고드망~\cite{I-4}(이하 (G)), 그리고 A.~그로텐디크의 최근 논문~\cite{I-6}(이하 (T))을 참조한다.
  \item[c)] \emph{층 이론}. 주된 참고문헌은 (G)와 (T)이다. 이 이론은 가환대수학의 핵심 개념들을 ``기하학적'' 용어로 해석하고 ``전역화''하는 데 필수적인 언어를 제공한다.
  \item[d)] 끝으로 독자가 \emph{함자 언어}에 어느 정도 익숙하면 유용할 것이다. 이 논고에서는 이 언어를 계속 사용한다. (M), (G), 특히 (T)를 참조할 수 있으며, 그 원리와 일반 함자론의 주요 결과들은 이 논고의 저자들이 준비 중인 별도의 저서에서 더 자세히 설명할 예정이다.
\end{enumerate}

\bigskip
\begin{center}*\quad*\quad*\end{center}
\bigskip

이 서론은 대수기하학에서 ``스킴''이라는 관점이 무엇인지 개략적으로 설명하거나, 이 관점을 채택해야 했던 수많은 이유를 열거할 자리가 아니다. 특히 우리가 다루는 ``다양체''의 국소환에서 멱영원을 체계적으로 받아들여야 하는 이유도 여기서는 설명하지 않는다. 이 수용은 필연적으로 유리 사상의 개념을 뒤로 물리고, 정칙 사상 또는 ``사상''의 개념을 앞세운다.
이 논고의 목적은 바로 ``스킴''의 언어를 체계적으로 전개하고, 바라건대 그 필요성을 입증하는 데 있다.
그렇게 하기는 어렵지 않겠지만, 제1장에서 전개하는 개념들에 대한 ``직관적'' 입문도 여기서는 시도하지 않겠다.

\oldpage[I]{6}
이 논고의 내용을 미리 훑어보고 싶은 독자는 A.~그로텐디크가 1958년 에든버러 국제수학자대회에서 한 강연~\cite{I-7}과 같은 저자의 발표~\cite{I-8}을 참조할 수 있다.
J.-P.~세르의 논문~\cite{I-14}(이하 (FAC))도 대수기하학의 고전적 관점과 스킴의 관점 사이를 잇는 설명으로 볼 수 있다. 따라서 이 논문을 읽는 것은 우리의 \emph{원론}을 읽기 위한 훌륭한 준비가 될 수 있다.

\bigskip
\begin{center}*\quad*\quad*\end{center}
\bigskip

참고로 이 논고에서 예정한 전체 구성을 아래에 제시한다. 특히 뒤쪽 장들의 구성은 나중에 바뀔 수 있다.

\medskip
\begin{tabular}{rl}
  제1장. & --- 스킴의 언어.\\
  --- \textbf{II}. & --- 몇몇 사상 부류에 대한 초보적 전역 연구.\\
  --- \textbf{III}. & --- 연접 대수적 층의 코호몰로지. 응용.\\
  --- \textbf{IV}. & --- 사상의 국소적 연구.\\
  --- \textbf{V}. & --- 스킴의 초보적 구성법.\\
  --- \textbf{VI}. & --- 강하 기법. 스킴 구성의 일반적 방법.\\
  --- \textbf{VII}. & --- 군 스킴, 주다발 공간.\\
  --- \textbf{VIII}. & --- 다발 공간의 미분론적 연구.\\
  --- \textbf{IX}. & --- 기본군.\\
  --- \textbf{X}. & --- 유수와 쌍대성.\\
  --- \textbf{XI}. & --- 교차이론, 천 특성류, 리만-로흐 정리.\\
  --- \textbf{XII}. & --- 아벨 스킴과 피카르 스킴.\\
  --- \textbf{XIII}. & --- 베유 코호몰로지.\\
\end{tabular}
\medskip

원칙적으로 모든 장은 열려 있는 것으로 간주하며, 언제든지 절을 덧붙일 수 있다. 채택한 출판 방식의 불편을 줄이기 위해 그러한 절은 별도의 분책으로 출간한다.
한 장이 출간될 때 어떤 절이 예정되어 있거나 준비 중이면, 긴급성에 따른 순서 때문에 실제 출간이 훨씬 뒤로 미뤄지더라도 그 장의 목차에 이를 적어 둔다.
독자의 편의를 위해 이 논고의 여러 장에서 사용하는 가환대수학, 호몰로지 대수학, 층 이론의 여러 보충 결과를 ``제0장''에 모았다. 이 결과들은 어느 정도 알려져 있지만 편리한 참고문헌을 제시할 수 없었던 것들이다.
독자는 논고 본문을 읽다가 우리가 인용하는 결과가 충분히 익숙하지 않을 때에만 제0장을 참조하기를 권한다.

\oldpage[I]{7}
이렇게 읽으면 이 논고는 초심자가 가환대수학과 호몰로지 대수학에 익숙해지는 좋은 방법이 될 수 있다고 생각한다. 구체적인 응용을 동반하지 않는 이 두 분야의 공부는 적지 않은 이들에게 따분하고 심지어 침울하기까지 한 것으로 여겨지기 때문이다.

\bigskip
\begin{center}*\quad*\quad*\end{center}
\bigskip

이 서론에서 여기 제시하는 개념과 결과의 역사를, 아무리 간략하게라도 개관하는 일은 우리의 능력 밖이다.
본문에는 이해에 특히 유용하다고 판단한 참고문헌만 싣고, 가장 중요한 결과들에 대해서만 그 기원을 밝힐 것이다.
적어도 형식적으로는 이 저서가 다루는 주제들이 상당히 새롭다. 그래서 우리가 그 연구를 전해 들었을 뿐인 19세기와 20세기 초 대수기하학의 선구자들을 거의 인용하지 않는 것이다.
그러나 저자들에게 가장 직접적인 영향을 주고 스킴 관점의 발전에 기여한 저작들에 관해서는 여기서 몇 마디 할 필요가 있다.
무엇보다 먼저 J.-P.~세르의 기초적인 논문 (FAC)을 들 수 있다. 이 논문은 A.~베유의 고전적인 \emph{Foundations}~\cite{I-18}가 너무 건조하다고 느낀 여러 젊은 입문자들(그 가운데에는 이 논고의 저자 한 명도 있다)에게 대수기하학의 입문서 역할을 했다.
여기서 처음으로 ``추상'' 대수다양체의 ``자리스키 위상''이 대수적 위상수학의 특정 기법들을 적용하기에 완전히 적합하며, 특히 코호몰로지 이론을 낳는다는 사실이 증명되었다.
더 나아가 그곳에서 주어진 대수다양체의 정의는 우리가 여기서 전개하는 개념의 확장에 가장 자연스럽게 맞는다\footnote{J.-P.~세르가 지적했듯이, 환의 층을 주어 다양체의 구조를 정의한다는 발상은 H.~카르탕에게서 비롯되었으며, 카르탕은 이를 해석공간 이론의 출발점으로 삼았다.
물론 대수기하학에서와 마찬가지로 ``해석기하학''에서도 해석공간의 국소환 안에 있는 멱영원에 정당한 지위를 부여해야 한다.
H.~카르탕과 J.-P.~세르의 정의를 이렇게 확장하는 문제는 최근 H.~그라우어트~\cite{I-5}가 다루기 시작했으며, 이 일반적 틀 안에서 해석기하학을 체계적으로 설명하는 저서가 머지않아 나오기를 기대할 수 있다.
이 논고에서 전개하는 개념과 기법이 해석기하학에서도 의미를 유지한다는 점은 분명하다. 다만 후자의 이론에서는 기술적 어려움이 훨씬 더 클 것으로 예상해야 한다.
방법이 단순한 대수기하학은 앞으로 해석공간 이론이 발전하는 데 일종의 형식적 모형이 될 수 있을 것이다.}.
또한 세르는 체 위의 아핀 대수 대신 임의의 가환환을 쓰면 아핀 대수다양체의 코호몰로지 이론을 어려움 없이 옮겨 적을 수 있음을 이미 알아차렸다.
따라서 이 논고의 제I장과 제II장, 그리고 제III장의 첫 두 절은 본질적으로 (FAC)와 같은 저자의 뒤이은 논문~\cite{I-15}의 주요 결과를 이 확장된 틀로 곧바로 옮긴 것으로 볼 수 있다.
우리는 C.~슈발레의 \emph{대수기하학 세미나}~\cite{I-1}에서도 큰 도움을 얻었다. 특히 그가 도입한 ``구성가능 집합''을 체계적으로 사용하는 것이 스킴 이론에서 매우 유용하다는 사실이 드러났다(제IV장 참조).
또한 차원의 관점에서 사상을 연구하는 방법(제IV장)도 그에게서 가져왔다. 이 방법은 스킴의 틀로 거의 아무 변화 없이 옮겨진다.

\oldpage[I]{8}
한편 슈발레가 도입한 ``국소환의 스킴''이라는 개념은 대수기하학을 자연스럽게 확장할 수 있게 한다. 다만 우리가 여기서 부여하려는 유연성과 일반성을 모두 갖추지는 못한다. 이 개념과 우리 이론의 관계는 제1장 \S~8을 보라.
M.~나가타는 데데킨트 환 위의 대수기하학에 관한 많은 특수한 결과를 담은 일련의 논문~\cite{I-9}에서 이러한 확장을 전개하였다\footnote{대수기하학에 관한 우리의 관점에 가까운 연구로는 E.~켈러의 중요한 저작~\cite{I-22}과 저우와 이구사의 최근 노트~\cite{I-3}을 들 수 있다. 이들은 나가타-슈발레 이론의 틀 안에서 (FAC)의 몇몇 결과를 다시 다루고 퀸네트 공식도 제시한다.}.

\bigskip
\begin{center}*\quad*\quad*\end{center}
\bigskip

끝으로 대수기하학에 관한 책, 특히 그 기초에 관한 책이 O.~자리스키와 A.~베유 같은 수학자들에게서, 설령 다른 사람들을 거쳐서라도, 필연적으로 영향을 받는다는 것은 말할 필요도 없다.
특히 자리스키의 \emph{정칙함수론}~\cite{I-20}은 코호몰로지 방법으로 적절히 유연하게 만들고 존재정리(제III장 \S\S~4, 5)로 보완하면, 제VI장에서 설명하는 강하 기법과 더불어 이 논고에서 사용하는 주요 도구 가운데 하나가 된다. 우리가 보기에 이는 대수기하학에서 쓸 수 있는 가장 강력한 도구 가운데 하나이다.

이 이론이 속하는 일반적 기법은 다음과 같이 개략적으로 설명할 수 있다. 전형적인 예는 제IX장에서 기본군을 연구할 때 나타난다.
한 대수다양체에서 다른 대수다양체로 가는 고유 사상(제II장) $f:X\to Y$가 있고(더 일반적으로는 한 스킴에서 다른 스킴으로 가는 사상), 점 $y\in Y$의 근방에 관한 문제 P를 풀기 위해 $y$의 근방에서 이 사상을 연구하려 한다.
다음 단계를 차례로 밟는다.
\begin{enumerate}
  \item[$1^\circ$] $Y$가 아핀이라고 가정할 수 있다. 그러면 $X$는 $Y$의 아핀환 $A$ 위에서 정의된 스킴이 되며, 나아가 $A$를 $y$의 국소환으로 바꿀 수 있다.
  이 축약은 실제로 언제나 쉽고(제V장), $A$가 \emph{국소환}인 경우로 문제를 돌린다.
  \item[$2^\circ$] $A$가 \emph{아르틴 국소환}인 경우에 문제를 연구한다.
  $A$가 정역이라고 가정하지 않을 때에도 문제가 실제로 의미를 유지하게 하려면 문제 P를 다시 정식화해야 할 때가 있다. 이렇게 하면 흔히 이 단계에서 ``무한소적'' 성격을 띠는 문제를 더 잘 이해하게 된다.
  \item[$3^\circ$] 형식 스킴 이론(제III장 \S\S~3, 4, 5)을 이용하면 아르틴 환의 경우에서 \emph{완비 국소환}의 경우로 넘어갈 수 있다.
  \item[$4^\circ$] 끝으로 $A$가 임의의 국소환이면, $X$ 위의 적당한 스킴에서 주어진 ``형식적'' 절단에 가까운 ``다중 절단''을 고려함으로써(제IV장), $A$의 완비화로 스칼라를 확대하여 $X$에서 얻은 스킴에 대해 알려진 결과를 $A$의 충분히 단순한 유한 확대(예를 들면 비분기 확대)에 대한 유사한 결과로 흔히 옮길 수 있다.
\end{enumerate}

이 개요는 아르틴 환 $A$ 위에서 정의된 스킴을 체계적으로 연구하는 일이 중요함을 보여 준다.
세르가 국소 유체론을 정식화할 때 취한 관점과 그린버그의 최근 연구는, 그러한 스킴 $X$에 $A$의 잉여체 $k$(완전체라고 가정한다) 위의 스킴 $X'$를 함자적으로 대응시키는 방식으로 이 연구를 수행할 수 있음을 시사하는 듯하다. 여기서 $X'$의 차원은 좋은 경우 $n\dim X$이고, $n$은 $A$의 길이이다.

A.~베유의 영향에 관해서는 다음과 같이 말하는 것으로 충분하다. ``베유 코호몰로지''의 정의를 필요한 일반성 전체를 갖추어 정식화하고, 디오판토스 기하학의 유명한 추측들~\cite{I-19}을 확립하는 데 필요한 모든 형식적 성질의 증명\footnote{오해를 피하기 위해 덧붙이면, 이 서론을 쓰는 시점에는 이 과제가 이제 막 시작되었을 뿐이며, 따라서 아직 베유 추측의 증명에 이르지 못했다.}에 착수하려면 필요한 도구를 개발해야 한다. 이것이 이 논고를 집필하게 된 주된 동기 가운데 하나였다. 대수기하학에서 통상 사용하는 개념과 방법의 자연스러운 틀을 찾고, 저자들 자신이 그 개념과 기법을 이해할 기회를 얻으려는 바람도 똑같이 중요한 동기였다.

\bigskip
\begin{center}*\quad*\quad*\end{center}
\bigskip

끝으로 독자들에게 미리 알려 두는 것이 유익하다고 생각한다. 저자들 자신이 그랬듯이, 독자들도 스킴의 언어에 익숙해지고 기하학적 직관이 제안하는 통상적 구성들을 본질적으로 단 하나의 합리적인 방식으로 이 언어에 옮길 수 있음을 확신하기까지는 어느 정도 어려움을 겪을 것이다.
현대수학의 많은 분야에서 그렇듯이, 최초의 직관은 그것을 필요한 정밀성과 일반성을 모두 갖추어 표현하는 고유한 언어에서 겉보기에는 점점 더 멀어지고 있다.
여기서 심리적 어려움은 집합에 익숙한 개념들, 곧 데카르트 곱, 군·환·가군의 연산, 다발, 주동차다발 등을 집합의 범주와 이미 상당히 다른 범주의 대상들(즉 준스킴의 범주 또는 주어진 준스킴 위의 준스킴 범주)로 옮겨야 한다는 데 있다.
미래의 수학자가 이러한 새로운 추상화의 노력을 피하기는 어려울 것이다. 그러나 집합론에 익숙해지면서 우리의 선배들이 기울였던 노력과 비교하면, 결국 이 노력은 꽤 작은 것일지도 모른다.

\bigskip
\begin{center}*\quad*\quad*\end{center}
\bigskip

참조는 십진법 체계에 따라 표시한다. 예를 들어 III, 4.9.3에서 III은 장, 4는 절, 9는 그 절의 항을 나타낸다.
같은 장 안에서는 장 표시를 생략한다.

\clearpage
% Authority: EGA I ega0-1-fr.tex, SHA-256 F86663178E620A678BF0857B18BB4AC8FDB34E44BF5D241AE853EA7171E06C16.
% Sequential translation begins at authority lines 1-63.
\section{분수환}
\label{section:0.1-ko}

\setcounter{subsection}{-1}
\subsection{환과 대수}
\label{subsection:0.1.0-ko}

\oldpage[0\textsubscript{I}]{11}
\begin{env}[1.0.1]
\label{0.1.0.1-ko}
이 논고에서 다루는 모든 환은 \emph{단위원}을 갖는다.
그러한 환 위의 모든 가군은 \emph{단위원 가군}이라고 가정한다. 환 준동형은 언제나 단위원을 단위원으로 보낸다고 가정하며, 명시적으로 달리 말하지 않는 한 환 $A$의 부분환은 $A$의 \emph{단위원을 포함}한다고 가정한다.
주로 \emph{가환환}을 다루며, 아무 수식 없이 환이라고 말할 때에는 가환환을 뜻하는 것으로 이해한다.
$A$가 반드시 가환일 필요가 없는 환이면, 명시적으로 달리 말하지 않는 한 $A$-가군은 언제나 \emph{왼쪽} 가군을 뜻한다.
\end{env}

\begin{env}[1.0.2]
\label{0.1.0.2-ko}
$A$, $B$를 반드시 가환일 필요가 없는 두 환이라 하고, $\varphi:A\to B$를 준동형이라 하자.
모든 왼쪽(각각 오른쪽) $B$-가군 $M$에는 $a.m=\varphi(a).m$(각각 $m.a=m.\varphi(a)$)로 놓아 왼쪽(각각 오른쪽) $A$-가군 구조를 줄 수 있다. $M$ 위의 $A$-가군 구조와 $B$-가군 구조를 구별해야 할 때에는 이렇게 정의한 왼쪽(각각 오른쪽) $A$-가군을 $M_{[\varphi]}$로 나타낸다.
$L$이 $A$-가군이면 준동형 $u:L\to M_{[\varphi]}$는 곧 $a\in A$, $x\in L$에 대하여 $u(a.x)=\varphi(a).u(x)$를 만족하는 가환군 준동형이다. 이를 $L\to M$인 \emph{$\varphi$-준동형}이라고도 하며, 쌍 $(\varphi,u)$(또는 언어를 남용하여 $u$)를 $(A,L)$에서 $(B,M)$로 가는 \emph{디-준동형}이라고 한다.
따라서 환 $A$와 $A$-가군 $L$로 이루어진 쌍 $(A,L)$들은 디-준동형을 사상으로 하는 하나의 \emph{범주}를 이룬다.
\end{env}

\begin{env}[1.0.3]
\label{0.1.0.3-ko}
(1.0.2)의 가정 아래에서 $\mathfrak{J}$가 $A$의 왼쪽(각각 오른쪽) 아이디얼이면, $\varphi(\mathfrak{J})$가 생성하는 $B$의 왼쪽(각각 오른쪽) 아이디얼 $B\varphi(\mathfrak{J})$(각각 $\varphi(\mathfrak{J})B$)를 $B\mathfrak{J}$(각각 $\mathfrak{J}B$)로 나타낸다. 이는 왼쪽(각각 오른쪽) $B$-가군의 표준 준동형 $B\otimes_A\mathfrak{J}\to B$(각각 $\mathfrak{J}\otimes_A B\to B$)의 상이기도 하다.
\end{env}

\begin{env}[1.0.4]
\label{0.1.0.4-ko}
$A$가 (가환)환이고 $B$가 반드시 가환일 필요가 없는 환이면, $B$에 \emph{$A$-대수} 구조를 주는 것은 $\varphi(A)$가 $B$의 중심에 포함되도록 하는 환 준동형 $\varphi:A\to B$를 주는 것과 동치이다.
$A$의 모든 아이디얼 $\mathfrak{J}$에 대하여 $\mathfrak{J}B=B\mathfrak{J}$는 $B$의 양쪽 아이디얼이고, 모든 $B$-가군 $M$에 대하여 $\mathfrak{J}M$은 $(B\mathfrak{J})M$과 같은 $B$-부분가군이다.
\end{env}

\begin{env}[1.0.5]
\label{0.1.0.5-ko}
\emph{유한 생성 가군}과 \emph{유한 생성} (가환) \emph{대수}의 개념은 다시 설명하지 않겠다. $A$-가군 $M$이 유한 생성이라는 것은 완전열

\oldpage[0\textsubscript{I}]{12}
$A^p\to M\to 0$이 존재한다는 뜻이다.
$A$-가군 $M$이 어떤 준동형 $A^p\to A^q$의 여핵과 동형일 때, 다시 말해 완전열 $A^p\to A^q\to M\to 0$이 존재할 때 $M$이 \emph{유한 표시}를 갖는다고 한다.
\emph{뇌터 환} $A$ 위의 모든 유한 생성 $A$-가군은 유한 표시를 갖는다는 점에 유의하자.

$A$-대수 $B$의 모든 원소가 $A$에 계수를 갖는 일계수 다항식의 $B$ 안에서의 근일 때 $B$가 \emph{$A$ 위에서 정수적}이라고 함을 상기하자. 이는 $B$의 모든 원소가 유한 생성 \emph{$A$-가군}인 $B$의 어떤 부분대수에 들어간다는 것과 동치이다.
이 조건이 성립하고 $B$가 가환이면, $B$의 유한 부분집합이 생성하는 $B$의 부분대수는 유한 생성 $A$-가군이다. 따라서 가환대수 $B$가 $A$ 위에서 정수적이면서 유한 생성일 필요충분조건은 $B$가 유한 생성 $A$-가군인 것이다. 이때 $B$를 \emph{유한 정수적} $A$-대수라고도 한다(혼동의 우려가 없으면 단순히 \emph{유한} $A$-대수라고 한다).
이 정의들에서는 $A$-대수 구조를 정하는 준동형 $A\to B$가 단사라고 가정하지 않는다는 점에 유의하자.
\end{env}

\begin{env}[1.0.6]
\label{0.1.0.6-ko}
\emph{정역}은 $0$이 아닌 원소들로 이루어진 유한족의 곱이 $0$이 아닌 환이다. 이는 그러한 환에서 $0\neq 1$이고, $0$이 아닌 두 원소의 곱이 $0$이 아니라는 것과 동치이다.
환 $A$의 \emph{소 아이디얼}은 $A/\mathfrak{p}$가 정역이 되게 하는 아이디얼 $\mathfrak{p}$이다. 따라서 $\mathfrak{p}\neq A$이다.
환 $A$가 적어도 하나의 소 아이디얼을 가질 필요충분조건은 $A\neq\{0\}$이다.
\end{env}

\begin{env}[1.0.7]
\label{0.1.0.7-ko}
\emph{국소환}은 유일한 극대 아이디얼을 갖는 환 $A$이다. 이 극대 아이디얼은 가역원들의 여집합이며 $A$가 아닌 모든 아이디얼을 포함한다.
$A$, $B$가 각각 극대 아이디얼 $\mathfrak{m}$, $\mathfrak{n}$을 갖는 두 국소환일 때, 준동형 $\varphi:A\to B$가 $\varphi(\mathfrak{m})\subset\mathfrak{n}$을 만족하면(또는 이와 동치로 $\varphi^{-1}(\mathfrak{n})=\mathfrak{m}$이면) $\varphi$를 \emph{국소 준동형}이라고 한다.
몫으로 내려가면 이러한 준동형은 잉여체 $A/\mathfrak{m}$에서 잉여체 $B/\mathfrak{n}$으로 가는 단사 준동형을 정한다.
두 국소 준동형의 합성은 국소 준동형이다.
\end{env}

\subsection{아이디얼의 근기. 환의 멱영근기와 근기}
\label{subsection:0.1.1-ko}

\begin{env}[1.1.1]
\label{0.1.1.1-ko}
$\mathfrak{a}$를 환 $A$의 아이디얼이라 하자. $\mathfrak{a}$의 \emph{근기} $\mathfrak{r}(\mathfrak{a})$는 어떤 정수 $n>0$에 대하여 $x^n\in\mathfrak{a}$가 되는 $x\in A$ 전체의 집합이다. 이는 $\mathfrak{a}$를 포함하는 아이디얼이다.
$\mathfrak{r}(\mathfrak{r}(\mathfrak{a}))=\mathfrak{r}(\mathfrak{a})$이고, 관계 $\mathfrak{a}\subset\mathfrak{b}$는 $\mathfrak{r}(\mathfrak{a})\subset\mathfrak{r}(\mathfrak{b})$를 함의하며, 유한개의 아이디얼의 교집합의 근기는 각 근기의 교집합이다.
$\varphi$가 환 $A'$에서 $A$로 가는 준동형이면, 모든 아이디얼 $\mathfrak{a}\subset A$에 대하여 $\mathfrak{r}(\varphi^{-1}(\mathfrak{a}))=\varphi^{-1}(\mathfrak{r}(\mathfrak{a}))$이다.
한 아이디얼이 어떤 아이디얼의 근기일 필요충분조건은 그것이 소 아이디얼들의 교집합인 것이다.
아이디얼 $\mathfrak{a}$의 근기는 $\mathfrak{a}$를 포함하는 소 아이디얼들 가운데 \emph{극소}인 것들의 교집합이다. $A$가 뇌터 환이면 이러한 극소 소 아이디얼은 유한개이다.

영 아이디얼 $(0)$의 근기를 $A$의 \emph{멱영근기}라고도 한다. 이는 $A$의 멱영원 전체의 집합 $\mathfrak{N}$이다.
$\mathfrak{N}=(0)$일 때 환 $A$를 \emph{축소환}이라고 한다. 모든 환 $A$에 대하여 멱영근기로 나눈 몫환 $A/\mathfrak{N}$은 축소환이다.
\end{env}

\begin{env}[1.1.2]
\label{0.1.1.2-ko}
환 $A$(반드시 가환일 필요는 없다)의 \emph{근기} $\mathfrak{R}(A)$는 $A$의 극대 왼쪽 아이디얼 전체의 교집합(또한 극대 오른쪽 아이디얼 전체의 교집합)임을 상기하자.
$A/\mathfrak{R}(A)$의 근기는 $(0)$이다.
\end{env}

\subsection{분수가군과 분수환}
\label{subsection:0.1.2-ko}

\oldpage[0\textsubscript{I}]{13}
\begin{env}[1.2.1]
\label{0.1.2.1-ko}
환 $A$의 부분집합 $S$가 $1\in S$를 만족하고 $S$의 두 원소의 곱이 다시 $S$에 속하면 $S$를 \emph{곱셈적}이라고 한다.
이후 가장 중요한 예는 다음과 같다. $1^\circ$ 원소 $f\in A$의 거듭제곱 $f^n$($n\geq 0$) 전체의 집합 $S_f$. $2^\circ$ \emph{소} 아이디얼 $\mathfrak{p}$의 여집합 $A-\mathfrak{p}$.
\end{env}

\begin{env}[1.2.2]
\label{0.1.2.2-ko}
$S$를 환 $A$의 곱셈적 부분집합이라 하고, $M$을 $A$-가군이라 하자. 집합 $M\times S$에서 두 쌍 $(m_1,s_1)$, $(m_2,s_2)$ 사이의 관계
\[
  \text{``$s(s_1m_2-s_2m_1)=0$이 되는 $s\in S$가 존재한다''}
\]
는 동치관계이다.
이 관계에 의한 $M\times S$의 몫집합을 $S^{-1}M$으로 나타내고, 쌍 $(m,s)$의 $S^{-1}M$에서의 표준 상을 $m/s$로 나타낸다. 사상 $i_M^S:m\mapsto m/1$($i^S$로도 쓴다)을 $M$에서 $S^{-1}M$으로 가는 \emph{표준 사상}이라고 한다.
이 사상은 일반적으로 단사도 전사도 아니다. 그 핵은 어떤 $s\in S$에 대하여 $sm=0$이 되는 $m\in M$ 전체의 집합이다.

$S^{-1}M$ 위에
\[
  (m_1/s_1)+(m_2/s_2)=(s_2m_1+s_1m_2)/(s_1s_2)
\]
로 놓아 덧셈군 연산을 정의한다(이는 선택한 $S^{-1}M$의 원소의 표현에 무관함을 확인할 수 있다).
또 $S^{-1}A$ 위에 $(a_1/s_1)(a_2/s_2)=(a_1a_2)/(s_1s_2)$로 곱셈을 정의하고, 끝으로 $(a/s)(m/s')=(am)/(ss')$로 놓아 $S^{-1}A$를 작용소 집합으로 하는 $S^{-1}M$ 위의 외부 연산을 정의한다.
이로써 $S^{-1}A$에는 환 구조가 주어지고($S$에 분모를 갖는 $A$의 \emph{분수환}이라 한다), $S^{-1}M$에는 $S^{-1}A$-가군 구조가 주어진다($S$에 분모를 갖는 $M$의 \emph{분수가군}이라 한다). 모든 $s\in S$에 대하여 $s/1$은 $S^{-1}A$에서 가역이고 그 역원은 $1/s$이다.
표준 사상 $i_A^S$(각각 $i_M^S$)은 환 준동형(각각 $A$-가군 준동형)이다. 여기서 $S^{-1}M$은 준동형 $i_A^S:A\to S^{-1}A$를 통해 $A$-가군으로 본다.
\end{env}

\begin{env}[1.2.3]
\label{0.1.2.3-ko}
$f\in A$에 대하여 $S_f=\{f^n\}_{n\geq 0}$이면, $S_f^{-1}A$와 $S_f^{-1}M$ 대신 $A_f$와 $M_f$를 쓴다. $A_f$를 $A$ 위의 대수로 볼 때에는 $A_f=A[1/f]$라고 쓸 수 있다.
$A_f$는 몫대수 $A[T]/(fT-1)A[T]$와 동형이다.
$f=1$이면 $A_f$와 $M_f$는 각각 $A$와 $M$에 표준적으로 동일시된다. $f$가 멱영원이면 $A_f$와 $M_f$는 모두 영 대상이다.

$S=A-\mathfrak{p}$이고 $\mathfrak{p}$가 $A$의 소 아이디얼이면, $S^{-1}A$와 $S^{-1}M$ 대신 $A_{\mathfrak{p}}$와 $M_{\mathfrak{p}}$를 쓴다. $A_{\mathfrak{p}}$는 $i_A^S(\mathfrak{p})$가 생성하는 극대 아이디얼 $\mathfrak{q}$를 갖는 \emph{국소환}이고 $(i_A^S)^{-1}(\mathfrak{q})=\mathfrak{p}$이다. 몫으로 내려가면 $i_A^S$는 정역 $A/\mathfrak{p}$에서 체 $A_{\mathfrak{p}}/\mathfrak{q}$로 가는 단사 준동형을 주며, 이 체는 $A/\mathfrak{p}$의 분수체와 동일시된다.
\end{env}

\begin{env}[1.2.4]
\label{0.1.2.4-ko}
분수환 $S^{-1}A$와 표준 준동형 $i_A^S$는 하나의 \emph{보편 사상 문제}의 해이다. $u(S)$가 $B$의 \emph{가역원}들로 이루어지도록 하는 $A$에서 환 $B$로 가는 모든 준동형 $u$는 유일한 방식으로
\[
  u:A\xrightarrow{i_A^S}S^{-1}A\xrightarrow{u^*}B
\]

\oldpage[0\textsubscript{I}]{14}
와 같이 분해된다. 여기서 $u^*$는 환 준동형이다.
같은 가정 아래에서 $M$을 $A$-가군, $N$을 $B$-가군이라 하고, $v:M\to N$을 $A$-가군 준동형이라 하자($N$의 $A$-가군 구조는 $u:A\to B$로 정한다). 그러면 $v$는 유일한 방식으로
\[
  v:M\xrightarrow{i_M^S}S^{-1}M\xrightarrow{v^*}N
\]
와 같이 분해된다. 여기서 $v^*$는 $S^{-1}A$-가군 준동형이다($N$의 $S^{-1}A$-가군 구조는 $u^*$로 정한다).
\end{env}

\begin{env}[1.2.5]
\label{0.1.2.5-ko}
원소 $(a/s)\otimes m$에 원소 $(am)/s$를 대응시켜 $S^{-1}A$-가군의 표준 동형 $S^{-1}A\otimes_A M\xrightarrow{\sim}S^{-1}M$을 정의한다. 그 역동형은 $m/s$를 $(1/s)\otimes m$으로 보낸다.
\end{env}

\begin{env}[1.2.6]
\label{0.1.2.6-ko}
$S^{-1}A$의 모든 아이디얼 $\mathfrak{a}'$에 대하여 $\mathfrak{a}=(i_A^S)^{-1}(\mathfrak{a}')$은 $A$의 아이디얼이고, $\mathfrak{a}'$은 $i_A^S(\mathfrak{a})$가 생성하는 $S^{-1}A$의 아이디얼이며 $S^{-1}\mathfrak{a}$와 동일시된다(1.3.2).
사상 $\mathfrak{p}'\mapsto(i_A^S)^{-1}(\mathfrak{p}')$은 $S^{-1}A$의 \emph{소} 아이디얼 전체의 순서집합에서 $\mathfrak{p}\cap S=\varnothing$을 만족하는 $A$의 소 아이디얼 $\mathfrak{p}$ 전체의 순서집합으로 가는 순서동형이다.
또한 이때 국소환 $A_{\mathfrak{p}}$와 $(S^{-1}A)_{S^{-1}\mathfrak{p}}$는 표준적으로 동형이다(1.5.1).
\end{env}

\begin{env}[1.2.7]
\label{0.1.2.7-ko}
$A$가 \emph{정역}이고 그 분수체를 $K$로 나타내면, $0$을 포함하지 않는 모든 곱셈적 부분집합 $S$에 대하여 표준 사상 $i_A^S:A\to S^{-1}A$는 단사이고, $S^{-1}A$는 $A$를 포함하는 $K$의 부분환과 표준적으로 동일시된다.
특히 $A$의 모든 소 아이디얼 $\mathfrak{p}$에 대하여 $A_{\mathfrak{p}}$는 $A$를 포함하고 극대 아이디얼 $\mathfrak{p}A_{\mathfrak{p}}$를 갖는 국소환이며 $\mathfrak{p}A_{\mathfrak{p}}\cap A=\mathfrak{p}$이다.
\end{env}

\begin{env}[1.2.8]
\label{0.1.2.8-ko}
$A$가 \emph{축소환}이면(1.1.1), $S^{-1}A$도 축소환이다. 실제로 $x\in A$, $s\in S$에 대하여 $(x/s)^n=0$이면 $s'x^n=0$이 되는 $s'\in S$가 존재한다. 따라서 $(s'x)^n=0$이고, 가정에 따라 $s'x=0$이므로 $x/s=0$이다.
\end{env}

\subsection{함자적 성질}
\label{subsection:0.1.3-ko}

\begin{env}[1.3.1]
\label{0.1.3.1-ko}
$M$, $N$을 두 $A$-가군이라 하고, $u$를 $A$-준동형 $M\to N$이라 하자.
$S$가 $A$의 곱셈적 부분집합이면, $(S^{-1}u)(m/s)=u(m)/s$로 놓아 $S^{-1}u$로 나타내는 $S^{-1}A$-준동형 $S^{-1}M\to S^{-1}N$을 정의한다. $S^{-1}M$과 $S^{-1}N$을 각각 $S^{-1}A\otimes_A M$과 $S^{-1}A\otimes_A N$에 표준적으로 동일시하면(1.2.5), $S^{-1}u$는 $1\otimes u$와 동일시된다.
$P$가 세 번째 $A$-가군이고 $v$가 $A$-준동형 $N\to P$이면 $S^{-1}(v\circ u)=(S^{-1}v)\circ(S^{-1}u)$이다. 다시 말해 $A$와 $S$를 고정하면 $S^{-1}M$은 $A$-가군의 범주에서 $S^{-1}A$-가군의 범주로 가는, \emph{$M$에 관한 공변 함자}이다.
\end{env}

\begin{env}[1.3.2]
\label{0.1.3.2-ko}
함자 $S^{-1}M$은 \emph{완전}하다. 다시 말해 열
\[
  M\xrightarrow{u}N\xrightarrow{v}P
\]
이 완전하면 열
\[
  S^{-1}M\xrightarrow{S^{-1}u}S^{-1}N\xrightarrow{S^{-1}v}S^{-1}P
\]
도 완전하다.
특히 $u:M\to N$이 단사(각각 전사)이면 $S^{-1}u$도 단사(각각 전사)이다.

\oldpage[0\textsubscript{I}]{15}
$N$과 $P$가 $M$의 두 부분가군이면, $S^{-1}N$과 $S^{-1}P$는 $S^{-1}M$의 부분가군에 표준적으로 동일시되며
\[
  S^{-1}(N+P)=S^{-1}N+S^{-1}P
  \quad\text{이고}\quad
  S^{-1}(N\cap P)=(S^{-1}N)\cap(S^{-1}P)
\]
이다.
\end{env}

\begin{env}[1.3.3]
\label{0.1.3.3-ko}
$(M_\alpha,\varphi_{\beta\alpha})$를 $A$-가군의 귀납계라 하자. 그러면 $(S^{-1}M_\alpha,S^{-1}\varphi_{\beta\alpha})$는 $S^{-1}A$-가군의 귀납계이다.
$S^{-1}M_\alpha$와 $S^{-1}\varphi_{\beta\alpha}$를 텐서곱으로 나타내고(1.2.5, 1.3.1), 텐서곱과 귀납극한을 취하는 연산들이 서로 교환한다는 사실을 쓰면 표준 동형
\[
  S^{-1}\varinjlim M_\alpha\xrightarrow{\sim}\varinjlim S^{-1}M_\alpha
\]
를 얻는다. 이를 함자 $S^{-1}M$($M$에 관한 함자)이 \emph{귀납극한과 가환한다}고도 표현한다.
\end{env}

\begin{env}[1.3.4]
\label{0.1.3.4-ko}
$M$, $N$을 두 $A$-가군이라 하자. $(M,N)$에 관해 \emph{함자적인} 표준 동형
\[
  (S^{-1}M)\otimes_{S^{-1}A}(S^{-1}N)
  \xrightarrow{\sim}
  S^{-1}(M\otimes_A N)
\]
이 존재하며, 이는 $(m/s)\otimes(n/t)$를 $(m\otimes n)/st$로 보낸다.
\end{env}

\begin{env}[1.3.5]
\label{0.1.3.5-ko}
마찬가지로 $(M,N)$에 관해 \emph{함자적인} 준동형
\[
  S^{-1}\operatorname{Hom}_A(M,N)
  \longrightarrow
  \operatorname{Hom}_{S^{-1}A}(S^{-1}M,S^{-1}N)
\]
이 존재하며, 이는 $u/s$에 준동형 $m/t\mapsto u(m)/st$를 대응시킨다.
$M$이 \emph{유한 표시}를 가지면 앞의 준동형은 \emph{동형}이다. $M$이 $A^r$ 꼴일 때에는 바로 알 수 있고, 일반적인 경우에는 완전열 $A^p\to A^q\to M\to 0$에서 출발하여 함자 $S^{-1}M$의 완전성과 $M$에 관한 함자 $\operatorname{Hom}_A(M,N)$의 왼쪽 완전성을 사용하면 된다.
$A$가 \emph{뇌터 환}이고 $A$-가군 $M$이 \emph{유한 생성}이면 언제나 이 경우에 해당한다는 점에 유의하자.
\end{env}

\subsection{곱셈적 부분집합의 변경}
\label{subsection:0.1.4-ko}

\begin{env}[1.4.1]
\label{0.1.4.1-ko}
$S$, $T$를 $S\subset T$인 환 $A$의 두 곱셈적 부분집합이라 하자. $S^{-1}A$의 $a/s$로 나타낸 원소를 $T^{-1}A$에서 역시 $a/s$로 나타낸 원소에 대응시키는 표준 준동형 $\rho_A^{T,S}$($\rho^{T,S}$로도 쓴다)이 $S^{-1}A$에서 $T^{-1}A$로 존재하며
\[
  i_A^T=\rho_A^{T,S}\circ i_A^S
\]
이다.
모든 $A$-가군 $M$에 대하여 $S^{-1}M$에서 $T^{-1}M$으로 가는 $S^{-1}A$-선형 사상도 존재한다. 여기서 $T^{-1}M$은 준동형 $\rho_A^{T,S}$를 통해 $S^{-1}A$-가군으로 보며, 이 사상은 $S^{-1}M$의 원소 $m/s$를 $T^{-1}M$의 원소 $m/s$로 보낸다. 이 사상을 $\rho_M^{T,S}$ 또는 단순히 $\rho^{T,S}$로 나타내며
\[
  i_M^T=\rho_M^{T,S}\circ i_M^S
\]
이다. 표준 동일시 (1.2.5) 아래에서 $\rho_M^{T,S}$는 $\rho_A^{T,S}\otimes 1$과 동일시된다.
준동형 $\rho_M^{T,S}$는 함자 $S^{-1}M$에서 함자 $T^{-1}M$으로 가는 \emph{함자적 사상}(또는 자연변환)이다. 다시 말해 그림
\[
\begin{tikzcd}[column sep=large,row sep=large]
S^{-1}M \arrow[r,"S^{-1}u"] \arrow[d,"\rho_M^{T,S}"'] & S^{-1}N \arrow[d,"\rho_N^{T,S}"] \\
T^{-1}M \arrow[r,"T^{-1}u"] & T^{-1}N
\end{tikzcd}
\]

\oldpage[0\textsubscript{I}]{16}
은 모든 준동형 $u:M\to N$에 대하여 가환한다. 또한 $T^{-1}u$는 $S^{-1}u$에 의해 완전히 결정된다는 점에 유의하자. 실제로 $m\in M$, $t\in T$에 대하여
\[
  (T^{-1}u)(m/t)=(t/1)^{-1}\rho^{T,S}\bigl((S^{-1}u)(m/1)\bigr)
\]
이다.
\end{env}

\begin{env}[1.4.2]
\label{0.1.4.2-ko}
같은 표기 아래에서 두 $A$-가군 $M$, $N$에 관한 다음 그림들은 가환한다(1.3.4, 1.3.5 참조).
\[
\begin{tikzcd}[column sep=large,row sep=large]
(S^{-1}M)\otimes_{S^{-1}A}(S^{-1}N) \arrow[r,"\sim"] \arrow[d]
  & S^{-1}(M\otimes_A N) \arrow[d]
  \\
(T^{-1}M)\otimes_{T^{-1}A}(T^{-1}N) \arrow[r,"\sim"]
  & T^{-1}(M\otimes_A N)
\end{tikzcd}
\]
\[
\begin{tikzcd}[column sep=large,row sep=large]
S^{-1}\operatorname{Hom}_A(M,N) \arrow[r] \arrow[d]
  & \operatorname{Hom}_{S^{-1}A}(S^{-1}M,S^{-1}N) \arrow[d] \\
T^{-1}\operatorname{Hom}_A(M,N) \arrow[r]
  & \operatorname{Hom}_{T^{-1}A}(T^{-1}M,T^{-1}N)
\end{tikzcd}
\]
\end{env}

\begin{env}[1.4.3]
\label{0.1.4.3-ko}
준동형 $\rho^{T,S}$가 \emph{전단사}인 중요한 경우가 있다. 곧 $T$의 모든 원소가 $S$의 어떤 원소의 약수인 경우이다. 이때 $\rho^{T,S}$를 통해 가군 $S^{-1}M$과 $T^{-1}M$을 동일시한다.
$S$의 원소의 $A$ 안에서의 모든 약수가 다시 $S$에 속할 때 $S$를 \emph{포화}되었다고 한다. $S$를 $S$의 원소들의 약수 전체의 집합 $T$로 바꾸면($T$는 곱셈적이고 포화되어 있다), 원한다면 언제나 $S$가 포화된 분수가군 $S^{-1}M$만을 고려해도 됨을 알 수 있다.
\end{env}

\begin{env}[1.4.4]
\label{0.1.4.4-ko}
$S$, $T$, $U$가 $S\subset T\subset U$인 $A$의 세 곱셈적 부분집합이면
\[
  \rho^{U,S}=\rho^{U,T}\circ\rho^{T,S}
\]
이다.
\end{env}

\begin{env}[1.4.5]
\label{0.1.4.5-ko}
$A$의 곱셈적 부분집합들의 \emph{증가 여과족} $(S_\alpha)$를 생각하자($S_\alpha\subset S_\beta$일 때 $\alpha\leq\beta$라고 쓴다). 곱셈적 부분집합 $S=\bigcup_\alpha S_\alpha$로 놓고
\[
  \rho_{\beta\alpha}=\rho_A^{S_\beta,S_\alpha}\quad\text{($\alpha\leq\beta$일 때)}
\]
로 놓자.
(1.4.4)에 따라 준동형 $\rho_{\beta\alpha}$들은 환의 귀납계 $(S_\alpha^{-1}A,\rho_{\beta\alpha})$의 \emph{귀납극한}인 환 $A'$을 정의한다.
$\rho_\alpha$를 표준 사상 $S_\alpha^{-1}A\to A'$이라 하고 $\varphi_\alpha=\rho_A^{S,S_\alpha}$로 놓자. (1.4.4)에 따라 $\alpha\leq\beta$이면 $\varphi_\alpha=\varphi_\beta\circ\rho_{\beta\alpha}$이므로, 그림
\[
\begin{tikzcd}[column sep=huge,row sep=large]
  & S_\alpha^{-1}A
      \arrow[ddl,"\rho_\alpha"']
      \arrow[d,"\rho_{\beta\alpha}"]
      \arrow[ddr,"\varphi_\alpha"] & \\
  & S_\beta^{-1}A
      \arrow[dl,"\rho_\beta"]
      \arrow[dr,"\varphi_\beta"'] & \\
A' \arrow[rr,"\varphi"'] & & S^{-1}A
\end{tikzcd}
\qquad(\alpha\leq\beta)
\]
이 가환하도록 하는 준동형 $\varphi:A'\to S^{-1}A$를 유일하게 정의할 수 있다.
실제로 $\varphi$는 \emph{동형}이다. 구성으로부터 $\varphi$가 전사라는 것은 바로 알 수 있다.
한편 $\varphi(\rho_\alpha(a/s_\alpha))=0$인 $\rho_\alpha(a/s_\alpha)\in A'$이 있다고 하자. 이는 $S^{-1}A$에서 $a/s_\alpha=0$, 즉 $sa=0$인 $s\in S$가 존재한다는 뜻이다. 그런데 $s\in S_\beta$인 $\beta\geq\alpha$가 존재하므로 $\rho_\alpha(a/s_\alpha)=\rho_\beta(sa/ss_\alpha)=0$이다. 따라서 $\varphi$는 단사이다.
$A$-가군 $M$의 경우도 똑같이 다루어 표준 동형
\[
  \varinjlim S_\alpha^{-1}A\xrightarrow{\sim}
  (\varinjlim S_\alpha)^{-1}A,
  \qquad
  \varinjlim S_\alpha^{-1}M\xrightarrow{\sim}
  (\varinjlim S_\alpha)^{-1}M
\]
을 얻는다. 둘째 동형은 \emph{$M$에 관해 함자적}이다.
\end{env}

\oldpage[0\textsubscript{I}]{17}
\begin{env}[1.4.6]
\label{0.1.4.6-ko}
$S_1$, $S_2$를 $A$의 두 곱셈적 부분집합이라 하자. 그러면 $S_1S_2$도 $A$의 곱셈적 부분집합이다.
$S_2'$을 환 $S_1^{-1}A$ 안에서의 $S_2$의 표준 상이라 하자. 이는 이 환의 곱셈적 부분집합이다.
모든 $A$-가군 $M$에 대하여 함자적 동형
\[
  S_2'^{-1}(S_1^{-1}M)\xrightarrow{\sim}(S_1S_2)^{-1}M
\]
이 존재하며, 이는 $(m/s_1)/(s_2/1)$을 $m/(s_1s_2)$에 대응시킨다.
\end{env}

\subsection{환의 변경}
\label{subsection:0.1.5-ko}

\begin{env}[1.5.1]
\label{0.1.5.1-ko}
$A$, $A'$을 두 환, $\varphi$를 준동형 $A'\to A$, $S$(각각 $S'$)를 $\varphi(S')\subset S$를 만족하는 $A$(각각 $A'$)의 곱셈적 부분집합이라 하자. 합성 준동형
\[
  A'\xrightarrow{\varphi}A\longrightarrow S^{-1}A
\]
은 (1.2.4)에 따라
\[
  A'\longrightarrow S'^{-1}A'\xrightarrow{\varphi^{S'}}S^{-1}A
\]
와 같이 분해되며, $\varphi^{S'}(a'/s')=\varphi(a')/\varphi(s')$이다.
$A=\varphi(A')$이고 $S=\varphi(S')$이면 $\varphi^{S'}$는 \emph{전사}이다.
$A'=A$이고 $\varphi$가 항등사상이면 $\varphi^{S'}$는 (1.4.1)에서 정의한 준동형 $\rho_A^{S,S'}$와 다르지 않다.
\end{env}

\begin{env}[1.5.2]
\label{0.1.5.2-ko}
(1.5.1)의 가정 아래에서 $M$을 $A$-가군이라 하자.
$S'^{-1}A'$-가군의 표준적인 함자적 준동형
\[
  \sigma:S'^{-1}(M_{[\varphi]})\longrightarrow
  (S^{-1}M)_{[\varphi^{S'}]}
\]
이 존재하며, 이는 $S'^{-1}(M_{[\varphi]})$의 모든 원소 $m/s'$을 $(S^{-1}M)_{[\varphi^{S'}]}$의 원소 $m/\varphi(s')$에 대응시킨다. 이 정의가 해당 원소의 표현 $m/s'$에 무관하다는 것은 바로 확인할 수 있다.
$S=\varphi(S')$이면 준동형 $\sigma$는 \emph{전단사}이다.
$A'=A$이고 $\varphi$가 항등사상이면 $\sigma$는 (1.4.1)에서 정의한 준동형 $\rho_M^{S,S'}$와 다르지 않다.

특히 $M=A$로 놓으면 준동형 $\varphi$는 $A$ 위에 $A'$-대수 구조를 정의한다. 이때 $S'^{-1}(A_{[\varphi]})$에는 $(\varphi(S'))^{-1}A$와 동일시되는 환 구조가 주어지고, 준동형
\[
  \sigma:S'^{-1}(A_{[\varphi]})\longrightarrow S^{-1}A
\]
는 $S'^{-1}A'$-대수의 준동형이다.
\end{env}

\begin{env}[1.5.3]
\label{0.1.5.3-ko}
$M$, $N$을 두 $A$-가군이라 하자. (1.3.4)와 (1.5.2)에서 정의한 준동형들을 합성하면 준동형
\[
  (S^{-1}M\otimes_{S^{-1}A}S^{-1}N)_{[\varphi^{S'}]}
  \longleftarrow
  S'^{-1}((M\otimes_A N)_{[\varphi]})
\]
을 얻으며, $\varphi(S')=S$이면 이는 동형이다.
마찬가지로 준동형 (1.3.5)와 (1.5.2)를 합성하면 준동형
\[
  S'^{-1}((\operatorname{Hom}_A(M,N))_{[\varphi]})
  \longrightarrow
  (\operatorname{Hom}_{S^{-1}A}(S^{-1}M,S^{-1}N))_{[\varphi^{S'}]}
\]
을 얻으며, $\varphi(S')=S$이고 $M$이 유한 표시를 가지면 이는 동형이다.
\end{env}

\begin{env}[1.5.4]
\label{0.1.5.4-ko}
이제 $A'$-가군 $N'$을 생각하고 텐서곱 $N'\otimes_{A'}A_{[\varphi]}$를 만들자. 이를
\[
  a.(n'\otimes b)=n'\otimes(ab)
\]
로 놓아 $A$-가군으로 볼 수 있다.
$S^{-1}A$-가군의 함자적 동형
\[
  \tau:(S'^{-1}N')\otimes_{S'^{-1}A'}(S^{-1}A)_{[\varphi^{S'}]}
  \xrightarrow{\sim}
  S^{-1}(N'\otimes_{A'}A_{[\varphi]})
\]
이 존재한다.

\oldpage[0\textsubscript{I}]{18}
이는 원소 $(n'/s')\otimes(a/s)$를 원소 $(n'\otimes a)/(\varphi(s')s)$에 대응시킨다. $n'/s'$(각각 $a/s$)를 같은 원소의 다른 표현으로 바꾸어도 $(n'\otimes a)/(\varphi(s')s)$가 변하지 않는다는 것을 각각 확인할 수 있다. 한편 $(n'\otimes a)/s$에 $(n'/1)\otimes(a/s)$를 대응시켜 $\tau$의 역준동형을 정의할 수 있다. 여기서는 $S^{-1}(N'\otimes_{A'}A_{[\varphi]})$가 $(N'\otimes_{A'}A_{[\varphi]})\otimes_A S^{-1}A$와 표준적으로 동형이고(1.2.5), 따라서 $A'$에서 $S^{-1}A$로 가는 합성 준동형 $a'\mapsto\varphi(a')/1$을 $\psi$로 나타낼 때 $N'\otimes_{A'}(S^{-1}A)_{[\psi]}$와도 표준적으로 동형이라는 사실을 사용한다.
\end{env}

\begin{env}[1.5.5]
\label{0.1.5.5-ko}
$M'$, $N'$이 두 $A'$-가군이면, (1.3.4)와 (1.5.4)의 동형들을 합성하여 동형
\[
  S'^{-1}M'\otimes_{S'^{-1}A'}S'^{-1}N'
  \otimes_{S'^{-1}A'}S^{-1}A
  \xrightarrow{\sim}
  S^{-1}(M'\otimes_{A'}N'\otimes_{A'}A)
\]
을 얻는다.
마찬가지로 $M'$이 유한 표시를 가지면 (1.3.5)와 (1.5.4)에 따라 동형
\[
  \operatorname{Hom}_{S'^{-1}A'}(S'^{-1}M',S'^{-1}N')
  \otimes_{S'^{-1}A'}S^{-1}A
  \xrightarrow{\sim}
  S^{-1}(\operatorname{Hom}_{A'}(M',N')\otimes_{A'}A)
\]
을 얻는다.
\end{env}

\begin{env}[1.5.6]
\label{0.1.5.6-ko}
(1.5.1)의 가정 아래에서 $T$(각각 $T'$)를 $S\subset T$(각각 $S'\subset T'$)이고 $\varphi(T')\subset T$인 $A$(각각 $A'$)의 두 번째 곱셈적 부분집합이라 하자. 그러면 그림
\[
\begin{tikzcd}[column sep=huge,row sep=large]
S'^{-1}A' \arrow[r,"\varphi^{S'}"] \arrow[d,"\rho^{T',S'}"']
  & S^{-1}A \arrow[d,"\rho^{T,S}"] \\
T'^{-1}A' \arrow[r,"\varphi^{T'}"']
  & T^{-1}A
\end{tikzcd}
\]
은 가환한다. $M$이 $A$-가군이면 그림
\[
\begin{tikzcd}[column sep=huge,row sep=large]
S'^{-1}(M_{[\varphi]}) \arrow[r,"\sigma"] \arrow[d,"\rho^{T',S'}"']
  & (S^{-1}M)_{[\varphi^{S'}]} \arrow[d,"\rho^{T,S}"] \\
T'^{-1}(M_{[\varphi]}) \arrow[r,"\sigma"']
  & (T^{-1}M)_{[\varphi^{T'}]}
\end{tikzcd}
\]
도 가환한다. 끝으로 $N'$이 $A'$-가군이면 그림
\[
\begin{tikzcd}[column sep=huge,row sep=large]
(S'^{-1}N')\otimes_{S'^{-1}A'}(S^{-1}A)_{[\varphi^{S'}]}
  \arrow[r,"\tau","{\sim}"'] \arrow[d]
  & S^{-1}(N'\otimes_{A'}A_{[\varphi]}) \arrow[d,"\rho^{T,S}"] \\
(T'^{-1}N')\otimes_{T'^{-1}A'}(T^{-1}A)_{[\varphi^{T'}]}
  \arrow[r,"\tau"',"{\sim}"]
  & T^{-1}(N'\otimes_{A'}A_{[\varphi]})
\end{tikzcd}
\]
도 가환한다. 왼쪽의 수직 화살표는 $S'^{-1}N'$에 $\rho_{N'}^{T',S'}$를, $S^{-1}A$에 $\rho_A^{T,S}$를 적용하여 얻는다.
\end{env}

\oldpage[0\textsubscript{I}]{19}

\begin{env}[1.5.7]
\label{0.1.5.7-ko}
$A''$을 세 번째 환, $\varphi':A''\to A'$을 환 준동형, $S''$을 $\varphi'(S'')\subset S'$인 $A''$의 곱셈적 부분집합이라 하자. $\varphi''=\varphi\circ\varphi'$로 놓으면
\[
  \varphi''^{S''}=\varphi^{S'}\circ\varphi'^{S''}
\]
이다.
$M$을 $A$-가군이라 하자. 분명히 $M_{[\varphi'']}=(M_{[\varphi]})_{[\varphi']}$이다. $\sigma$가 (1.5.2)에서 $\varphi$로부터 정의된 것과 같은 방식으로 $\varphi'$과 $\varphi''$에서 정의된 준동형을 각각 $\sigma'$과 $\sigma''$이라 하면 추이 공식
\[
  \sigma''=\sigma\circ\sigma'
\]
이 성립한다.
끝으로 $N''$을 $A''$-가군이라 하자. $A$-가군 $N''\otimes_{A''}A_{[\varphi'']}$는
\[
  (N''\otimes_{A''}A'_{[\varphi']})\otimes_{A'}A_{[\varphi]}
\]
와 표준적으로 동일시된다.
마찬가지로 $S^{-1}A$-가군
\[
  (S''^{-1}N'')\otimes_{S''^{-1}A''}
  (S^{-1}A)_{[\varphi''^{S''}]}
\]
는
\[
  \bigl((S''^{-1}N'')\otimes_{S''^{-1}A''}
  (S'^{-1}A')_{[\varphi'^{S''}]}\bigr)
  \otimes_{S'^{-1}A'}(S^{-1}A)_{[\varphi^{S'}]}
\]
와 표준적으로 동일시된다.
이 동일시들 아래에서 $\tau$가 (1.5.4)에서 $\varphi$로부터 정의된 것과 같은 방식으로 $\varphi'$과 $\varphi''$에서 정의된 동형을 각각 $\tau'$과 $\tau''$이라 하면 추이 공식
\[
  \tau''=\tau\circ(\tau'\otimes 1)
\]
이 성립한다.
\end{env}

\begin{env}[1.5.8]
\label{0.1.5.8-ko}
$A$를 환 $B$의 부분환이라 하자. $A$의 모든 \emph{극소} 소 아이디얼 $\mathfrak p$에 대하여 $\mathfrak p=A\cap\mathfrak q$가 되는 $B$의 극소 소 아이디얼 $\mathfrak q$가 존재한다. 실제로 $A_{\mathfrak p}$는 $B_{\mathfrak p}$의 부분환이고(1.3.2), 유일한 소 아이디얼 $\mathfrak p'$을 갖는다(1.2.6). $B_{\mathfrak p}$는 영환이 아니므로 적어도 하나의 소 아이디얼 $\mathfrak q'$을 가지며, 반드시 $\mathfrak q'\cap A_{\mathfrak p}=\mathfrak p'$이다. 따라서 $\mathfrak q'$의 역상인 $B$의 소 아이디얼 $\mathfrak q_1$은 $\mathfrak q_1\cap A=\mathfrak p$를 만족한다. 그러므로 특히 $\mathfrak q_1$에 포함된 $B$의 모든 극소 소 아이디얼 $\mathfrak q$에 대하여 $\mathfrak q\cap A=\mathfrak p$이다.
\end{env}

\subsection{가군 $M_f$와 귀납극한의 동일시}
\label{subsection:0.1.6-ko}

\begin{env}[1.6.1]
\label{0.1.6.1-ko}
$M$을 $A$-가군, $f$를 $A$의 원소라 하자. 모두 $M$과 같은 $A$-가군들의 열 $(M_n)$을 생각하고, 정수의 모든 쌍 $m\leq n$에 대하여 $\varphi_{nm}$을 $M_m$에서 $M_n$으로 가는 준동형 $z\mapsto f^{n-m}z$라 하자. $((M_n),(\varphi_{nm}))$이 $A$-가군의 \emph{귀납계}임은 자명하다. 이 계의 귀납극한을 $N=\varinjlim M_n$이라 하자. 이제 $N$에서 $M_f$ 위로 가는 함자적인 표준 $A$-동형을 정의하겠다. 이를 위해 모든 $n$에 대하여
\[
  \theta_n:z\longmapsto z/f^n
\]
은 $M=M_n$에서 $M_f$로 가는 $A$-준동형이고, 정의로부터 $m\leq n$일 때 $\theta_n\circ\varphi_{nm}=\theta_m$임에 유의하자. 따라서 표준 준동형 $M_n\to N$을 $\varphi_n$이라 할 때, 모든 $n$에 대하여 $\theta_n=\theta\circ\varphi_n$을 만족하는 $A$-준동형 $\theta:N\to M_f$가 존재한다. 가정에 따라 $M_f$의 모든 원소는 어떤 $n$에 대하여 $z/f^n$의 꼴이므로 $\theta$가 전사임은 분명하다. 한편 $\theta(\varphi_n(z))=0$, 즉 $z/f^n=0$이면 $f^kz=0$인 정수 $k>0$이 존재한다. 따라서 $\varphi_{n+k,n}(z)=0$이고, 이로부터 $\varphi_n(z)=0$이 따른다. 그러므로 $\theta$를 통해 $M_f$와 $\varinjlim M_n$을 동일시할 수 있다.
\end{env}

\begin{env}[1.6.2]
\label{0.1.6.2-ko}
이제 $M_n$, $\varphi_{nm}$, $\varphi_n$ 대신 각각 $M_{f,n}$, $\varphi_{nm}^f$, $\varphi_n^f$라 쓰자. $g$를 $A$의 두 번째 원소라 하자. $f^n$이 $f^ng^n$의 약수이므로 함자적 준동형
\[
  \rho_{fg,f}:M_f\longrightarrow M_{fg}\qquad\text{(1.4.1과 1.4.3)}
\]
이 존재한다.

\oldpage[0\textsubscript{I}]{20}

$M_f$와 $M_{fg}$를 각각 $\varinjlim M_{f,n}$과 $\varinjlim M_{fg,n}$으로 동일시하면, $\rho_{fg,f}$는
\[
  \rho_{fg,f}^n:M_{f,n}\longrightarrow M_{fg,n},
  \qquad \rho_{fg,f}^n(z)=g^nz
\]
로 정의되는 사상들의 \emph{귀납극한}과 동일시된다. 실제로 이는 그림
\[
\begin{tikzcd}[column sep=huge,row sep=large]
M_{f,n} \arrow[r,"\rho_{fg,f}^n"] \arrow[d,"\varphi_n^f"']
  & M_{fg,n} \arrow[d,"\varphi_n^{fg}"] \\
M_f \arrow[r,"\rho_{fg,f}"]
  & M_{fg}
\end{tikzcd}
\]
의 가환성에서 바로 따른다.
\end{env}

\subsection{가군의 지지집합}
\label{subsection:0.1.7-ko}

\begin{env}[1.7.1]
\label{0.1.7.1-ko}
$A$-가군 $M$이 주어졌을 때, $M_{\mathfrak p}\neq0$인 $A$의 소 아이디얼 $\mathfrak p$들의 집합을 $M$의 \emph{지지집합}이라 하고 $\operatorname{Supp}(M)$으로 나타낸다. $M=0$이기 위한 필요충분조건은 $\operatorname{Supp}(M)=\varnothing$인 것이다. 실제로 모든 $\mathfrak p$에 대하여 $M_{\mathfrak p}=0$이면, 임의의 원소 $x\in M$의 소멸자는 $A$의 어떤 소 아이디얼에도 포함될 수 없으므로 $A$ 전체와 같다.
\end{env}

\begin{env}[1.7.2]
\label{0.1.7.2-ko}
$A$-가군의 완전열 $0\to N\to M\to P\to0$이 있으면
\[
  \operatorname{Supp}(M)=\operatorname{Supp}(N)\cup\operatorname{Supp}(P)
\]
이다. 실제로 $A$의 모든 소 아이디얼 $\mathfrak p$에 대하여 열
\[
  0\longrightarrow N_{\mathfrak p}\longrightarrow M_{\mathfrak p}
  \longrightarrow P_{\mathfrak p}\longrightarrow0
\]
은 완전하고(1.3.2), $M_{\mathfrak p}=0$이기 위한 필요충분조건은 $N_{\mathfrak p}=P_{\mathfrak p}=0$인 것이다.
\end{env}

\begin{env}[1.7.3]
\label{0.1.7.3-ko}
$M$이 부분가군들의 족 $(M_\lambda)$의 합이면, $A$의 모든 소 아이디얼 $\mathfrak p$에 대하여 $M_{\mathfrak p}$는 $(M_\lambda)_{\mathfrak p}$들의 합이므로(1.3.3과 1.3.2)
\[
  \operatorname{Supp}(M)=\bigcup_\lambda\operatorname{Supp}(M_\lambda)
\]
이다.
\end{env}

\begin{env}[1.7.4]
\label{0.1.7.4-ko}
$M$이 \emph{유한 생성} $A$-가군이면, $\operatorname{Supp}(M)$은 \emph{$M$의 소멸자를 포함하는} 소 아이디얼들의 집합이다. 실제로 $M$이 한 원소로 생성되고 $x$가 그 생성원이면, $M_{\mathfrak p}=0$이라는 것은 $s.x=0$인 $s\notin\mathfrak p$가 존재한다는 뜻이고, 따라서 $\mathfrak p$가 $x$의 소멸자를 포함하지 않는다는 뜻이다. 이제 $M$이 유한 생성계 $(x_i)_{1\leq i\leq n}$을 가지고 $\mathfrak a_i$가 $x_i$의 소멸자라 하자. (1.7.3)에 따라 $\operatorname{Supp}(M)$은 $\mathfrak a_i$들 가운데 하나를 포함하는 $\mathfrak p$들의 집합이며, 이는 $M$의 소멸자인 $\mathfrak a=\bigcap_i\mathfrak a_i$를 포함하는 $\mathfrak p$들의 집합과 같다.
\end{env}

\begin{env}[1.7.5]
\label{0.1.7.5-ko}
$M$, $N$이 두 \emph{유한 생성} $A$-가군이면
\[
  \operatorname{Supp}(M\otimes_A N)
  =\operatorname{Supp}(M)\cap\operatorname{Supp}(N)
\]
이다. 이를 위해서는 $A$의 소 아이디얼 $\mathfrak p$에 대하여 조건
\[
  M_{\mathfrak p}\otimes_{A_{\mathfrak p}}N_{\mathfrak p}\neq0
\]
이 ``$M_{\mathfrak p}\neq0$이고 $N_{\mathfrak p}\neq0$이다''와 동치임을 보이면 충분하다((1.3.4) 참조). 다시 말해 국소환 $B$ 위의 두 유한 생성 가군 $P$, $Q$가 영가군이 아니면 $P\otimes_BQ\neq0$임을 보이면 된다. $B$의 극대 아이디얼을 $\mathfrak m$이라 하자. 나카야마 보조정리에 따라 벡터공간 $P/\mathfrak mP$와 $Q/\mathfrak mQ$는 영공간이 아니므로, 그 텐서곱
\[
  (P/\mathfrak mP)\otimes_{B/\mathfrak m}(Q/\mathfrak mQ)
  =(P\otimes_B Q)\otimes_B(B/\mathfrak m)
\]
도 영공간이 아니다. 이로부터 결론이 따른다.

특히 $M$이 유한 생성 $A$-가군이고 $\mathfrak a$가 $A$의 아이디얼이면, $\operatorname{Supp}(M/\mathfrak aM)$은 $\mathfrak a$와 $M$의 소멸자 $\mathfrak n$을 모두 포함하는 소 아이디얼들의 집합이며(1.7.4), 다시 말해 $\mathfrak a+\mathfrak n$을 포함하는 소 아이디얼들의 집합이다.
\end{env}

\oldpage[0\textsubscript{I}]{21}

\section{기약공간과 뇌터 공간}
\label{section:0.2-ko}

\subsection{기약공간}
\label{subsection:0.2.1-ko}

\begin{env}[2.1.1]
\label{0.2.1.1-ko}
위상공간 $X$가 공집합이 아니고 $X$ 자체와 다른 두 닫힌 부분공간의 합집합으로 나타나지 않으면 $X$를 \emph{기약}이라 한다. 이는 $X\neq\varnothing$이고 $X$의 공집합이 아닌 두 열린집합(따라서 유한 개의 열린집합)의 교집합이 공집합이 아니라고 하는 것, 또는 공집합이 아닌 모든 열린집합이 모든 곳에서 조밀하다고 하는 것, 또는 $X$와 다른 모든 닫힌 부분집합이 \emph{어디에서도 조밀하지 않다}고 하는 것, 또는 끝으로 $X$의 모든 열린집합이 \emph{연결}되어 있다고 하는 것과 동치이다.
\end{env}

\begin{env}[2.1.2]
\label{0.2.1.2-ko}
위상공간 $X$의 부분공간 $Y$가 기약이기 위한 필요충분조건은 그 폐포 $\overline Y$가 기약인 것이다. 특히 한 점만으로 이루어진 부분공간의 폐포 $\overline{\{x\}}$인 모든 부분공간은 기약이다. 관계
\[
  y\in\overline{\{x\}}
  \qquad\bigl(\text{동치로 }\overline{\{y\}}\subset\overline{\{x\}}\bigr)
\]
를 $y$가 $x$의 \emph{특수화}라고 하거나 $x$가 $y$의 \emph{일반화}라고 표현하겠다. 기약공간 $X$에 $X=\overline{\{x\}}$인 점 $x$가 존재하면 $x$를 $X$의 \emph{일반점}이라 한다. 이때 $X$의 공집합이 아닌 모든 열린집합은 $x$를 포함하고, $x$를 포함하는 모든 부분공간은 $x$를 일반점으로 갖는다.
\end{env}

\begin{env}[2.1.3]
\label{0.2.1.3-ko}
다음 분리 공리를 만족하는 위상공간 $X$를 \emph{콜모고로프 공간}이라 함을 상기하자.

\smallskip
\noindent
\textup{(T$_0$)} $X$의 서로 다른 두 점 $x\neq y$가 주어지면, $x$, $y$ 중 한 점을 포함하고 다른 점은 포함하지 않는 열린집합이 존재한다.

기약 콜모고로프 공간에 일반점이 존재하면 그것은 \emph{유일}하다. 실제로 공집합이 아닌 열린집합은 모든 일반점을 포함한다.

위상공간 $X$의 모든 열린 덮개에서 $X$의 유한 열린 덮개를 뽑아낼 수 있으면 $X$를 \emph{준콤팩트}라 함을 상기하자. 이는 공집합이 아닌 닫힌집합들의 모든 감소 여과족이 공집합이 아닌 교집합을 갖는다는 것과 동치이다. $X$가 준콤팩트 공간이면 $X$의 공집합이 아닌 모든 닫힌 부분집합 $A$는 공집합이 아닌 \emph{극소} 닫힌집합 $M$을 포함한다. 실제로 $A$의 공집합이 아닌 닫힌 부분집합들의 집합은 포함관계 $\supset$에 관하여 귀납적이다. 더 나아가 $X$가 콜모고로프 공간이면 $M$은 반드시 한 점만으로 이루어진다. 이러한 $M$을 관용적으로 \emph{닫힌점}이라고도 한다.
\end{env}

\begin{env}[2.1.4]
\label{0.2.1.4-ko}
기약공간 $X$에서 공집합이 아닌 모든 열린 부분공간 $U$는 기약이고, $X$가 일반점 $x$를 가지면 $x$는 $U$의 일반점이기도 하다.

공집합이 아닌 열린집합들로 이루어진 위상공간 $X$의 덮개 $(U_\alpha)$를 생각하자. 이때 지표집합도 공집합이 아니라고 한다. $X$가 기약이기 위한 필요충분조건은 모든 $\alpha$에 대하여 $U_\alpha$가 기약이고 임의의 $\alpha$, $\beta$에 대하여 $U_\alpha\cap U_\beta\neq\varnothing$인 것이다. 이 조건이 필요함은 자명하다. 충분함을 보이려면 $X$의 공집합이 아닌 열린집합 $V$에 대하여 모든 $\alpha$에서 $V\cap U_\alpha$가 공집합이 아님을 보이면 충분하다. 그러면 모든 $\alpha$에서 $V\cap U_\alpha$가 $U_\alpha$에 조밀하고, 따라서 $V$가 $X$에 조밀하기 때문이다. 그런데 $V\cap U_\gamma\neq\varnothing$인 지표 $\gamma$가 적어도 하나 존재하므로 $V\cap U_\gamma$는 $U_\gamma$에 조밀하다. 또한 모든 $\alpha$에 대하여 $U_\alpha\cap U_\gamma\neq\varnothing$이므로 $V\cap U_\alpha\cap U_\gamma\neq\varnothing$이다.
\end{env}

\oldpage[0\textsubscript{I}]{22}

\begin{env}[2.1.5]
\label{0.2.1.5-ko}
$X$를 기약공간, $f$를 $X$에서 위상공간 $Y$로 가는 연속사상이라 하자. 그러면 $f(X)$는 기약이고, $x$가 $X$의 일반점이면 $f(x)$는 $f(X)$의 일반점이며 따라서 $\overline{f(X)}$의 일반점이기도 하다. 특히 $Y$도 기약이고 유일한 일반점 $y$를 가지면, $f(X)$가 모든 곳에서 조밀하기 위한 필요충분조건은 $f(x)=y$인 것이다.
\end{env}

\begin{env}[2.1.6]
\label{0.2.1.6-ko}
위상공간 $X$의 모든 기약 부분공간은 \emph{극대} 기약 부분공간에 포함되고, 이러한 극대 기약 부분공간은 반드시 닫혀 있다. $X$의 극대 기약 부분공간들을 $X$의 \emph{기약 성분}이라 한다. $Z_1$, $Z_2$가 $X$의 서로 다른 두 기약 성분이면 $Z_1\cap Z_2$는 부분공간 $Z_1$, $Z_2$ 각각에서 닫혀 있고 \emph{어디에서도 조밀하지 않은} 집합이다. 특히 $X$의 한 기약 성분이 일반점을 가지면(2.1.2), 그 일반점은 다른 어떤 기약 성분에도 속할 수 없다. $X$가 유한 개의 기약 성분 $Z_i$ $(1\leq i\leq n)$만을 가지고 각 $i$에 대하여
\[
  U_i=Z_i\cap\complement\Bigl(\bigcup_{j\neq i}Z_j\Bigr)
\]
로 놓으면, $U_i$들은 열려 있고 기약이며 서로소이고 그 합집합은 $X$에 조밀하다.

$U$를 위상공간 $X$의 열린 부분집합이라 하자. $Z$가 $U$와 만나는 $X$의 기약 부분집합이면 $Z\cap U$는 $Z$에서 열려 있고 조밀하며 따라서 기약이다. 역으로 $Y$가 $U$의 닫힌 기약 부분집합이면, $X$에서의 $Y$의 폐포 $\overline Y$는 기약이고 $\overline Y\cap U=Y$이다. 따라서 $U$의 기약 성분들과 $U$와 만나는 $X$의 기약 성분들 사이에는 \emph{전단사 대응}이 있다.
\end{env}

\begin{env}[2.1.7]
\label{0.2.1.7-ko}
위상공간 $X$가 유한 개의 닫힌 기약 부분공간 $Y_i$의 합집합이면, $X$의 기약 성분들은 $Y_i$들의 집합에서 극대인 원소들이다. 실제로 $Z$가 $X$의 닫힌 기약 부분집합이면 $Z$는 $Z\cap Y_i$들의 합집합이고, 이로부터 $Z$가 어떤 $Y_i$에 포함되어야 함이 곧바로 따른다. $Y$를 위상공간 $X$의 부분공간이라 하고 $Y$가 유한 개의 기약 성분 $Y_i$ $(1\leq i\leq n)$만을 가진다고 하자. 그러면 $X$에서의 폐포 $\overline{Y_i}$들은 $\overline Y$의 기약 성분들이다.
\end{env}

\begin{env}[2.1.8]
\label{0.2.1.8-ko}
$Y$를 유일한 일반점 $y$를 갖는 기약공간이라 하고, $X$를 위상공간, $f$를 $X$에서 $Y$로 가는 연속사상이라 하자. 그러면 $f^{-1}(y)$와 만나는 $X$의 모든 기약 성분 $Z$에 대하여 $f(Z)$는 $Y$에 조밀하다. 그 역은 반드시 참인 것은 아니다. 그러나 $Z$가 일반점 $z$를 가지고 $f(Z)$가 $Y$에 조밀하면 반드시 $f(z)=y$이다(2.1.5). 또한 이때 $Z\cap f^{-1}(y)$는 $f^{-1}(y)$에서 $\{z\}$의 폐포이므로 기약이고, $z$를 포함하는 $f^{-1}(y)$의 모든 기약 부분집합은 반드시 $Z$에 포함되므로(2.1.6) $z$는 $Z\cap f^{-1}(y)$의 일반점이다. $f^{-1}(y)$의 모든 기약 성분은 $X$의 한 기약 성분에 포함되므로, $f^{-1}(y)$와 만나는 $X$의 모든 기약 성분 $Z$가 일반점을 가진다면 이 기약 성분들의 집합과 $f^{-1}(y)$의 기약 성분들 $Z\cap f^{-1}(y)$의 집합 사이에는 \emph{전단사 대응}이 있다. 이때 $Z$의 일반점과 $Z\cap f^{-1}(y)$의 일반점은 같다.
\end{env}

\oldpage[0\textsubscript{I}]{23}

\subsection{뇌터 공간}
\label{subsection:0.2.2-ko}

\begin{env}[2.2.1]
\label{0.2.2.1-ko}
위상공간 $X$의 열린집합들의 집합이 \emph{극대 조건}을 만족하면, 또는 이와 동치로 $X$의 닫힌집합들의 집합이 \emph{극소 조건}을 만족하면 $X$를 \emph{뇌터}라 한다. 모든 $x\in X$가 뇌터 부분공간인 근방을 가지면 $X$를 \emph{국소 뇌터}라 한다.
\end{env}

\begin{env}[2.2.2]
\label{0.2.2.2-ko}
$E$를 \emph{극소 조건}을 만족하는 순서집합이라 하고, $P$를 $E$의 원소들에 관한 다음 조건을 만족하는 성질이라 하자. $a\in E$에 대하여 모든 $x<a$에서 $P(x)$가 참이면 $P(a)$도 참이다. 그러면 \emph{모든 $x\in E$에 대하여 $P(x)$가 참이다}. 이를 ``뇌터 귀납법의 원리''라 한다. 실제로 $P(x)$가 거짓인 $x\in E$들의 집합을 $F$라 하자. $F$가 공집합이 아니면 극소 원소 $a$를 갖는다. 그러면 모든 $x<a$에 대하여 $P(x)$가 참이므로 $P(a)$도 참이어야 하며, 이는 모순이다.

특히 $E$가 \emph{뇌터 공간의 닫힌 부분집합들의 집합}인 경우에 이 원리를 적용하겠다.
\end{env}

\begin{env}[2.2.3]
\label{0.2.2.3-ko}
뇌터 공간의 모든 부분공간은 뇌터이다. 역으로 유한 개의 뇌터 부분공간의 합집합인 모든 위상공간은 뇌터이다.
\end{env}

\begin{env}[2.2.4]
\label{0.2.2.4-ko}
모든 뇌터 공간은 준콤팩트이다. 역으로 모든 열린집합이 준콤팩트인 위상공간은 뇌터이다.
\end{env}

\begin{env}[2.2.5]
\label{0.2.2.5-ko}
뇌터 공간은 유한 개의 기약 성분만을 가진다. 이는 뇌터 귀납법으로 알 수 있다.
\end{env}

\section{층에 관한 보충}
\label{section:0.3-ko}

\subsection{범주에 값을 갖는 층}
\label{subsection:0.3.1-ko}

\begin{env}[3.1.1]
\label{0.3.1.1-ko}
$\mathbf K$를 범주라 하고, $(A_\alpha)_{\alpha\in I}$와 $(A_{\alpha\beta})_{(\alpha,\beta)\in I\times I}$를 $A_{\beta\alpha}=A_{\alpha\beta}$를 만족하는 $\mathbf K$의 두 대상족이라 하자. 또한 $(\rho_{\alpha\beta})_{(\alpha,\beta)\in I\times I}$를 사상
\[
  \rho_{\alpha\beta}:A_\alpha\longrightarrow A_{\alpha\beta}
\]
들의 족이라 하자. $\mathbf K$의 대상 $A$와 사상족 $\rho_\alpha:A\to A_\alpha$로 이루어진 쌍이 다음 조건을 만족하면, 이 쌍을 족 $(A_\alpha)$, $(A_{\alpha\beta})$, $(\rho_{\alpha\beta})$가 정의하는 \emph{보편 문제의 해}라 한다. $\mathbf K$의 모든 대상 $B$에 대하여 각 $f\in\operatorname{Hom}(B,A)$에 족
\[
  (\rho_\alpha\circ f)\in
  \prod_\alpha\operatorname{Hom}(B,A_\alpha)
\]
를 대응시키는 사상은 $\operatorname{Hom}(B,A)$에서 모든 지표쌍 $(\alpha,\beta)$에 대하여
\[
  \rho_{\alpha\beta}\circ f_\alpha
  =\rho_{\beta\alpha}\circ f_\beta
\]
를 만족하는 족 $(f_\alpha)$들의 집합 위로 가는 \emph{전단사}이다. 이러한 해가 존재하면 동형을 제외하고 유일함은 곧바로 알 수 있다.
\end{env}

\begin{env}[3.1.2]
\label{0.3.1.2-ko}
위상공간 $X$ 위에서 범주 $\mathbf K$에 값을 갖는 \emph{준층} $U\mapsto\mathcal F(U)$의 정의는 상기하지 않겠다(G, I, 1.9). 이러한 준층이 다음 공리를 만족하면 \emph{$\mathbf K$에 값을 갖는 층}이라 한다.

\smallskip
\noindent
\textup{(F)} \emph{$X$의 열린집합 $U$를 $U$에 포함된 열린집합 $U_\alpha$들로 덮는 임의의 덮개 $(U_\alpha)$에 대하여, $\rho_\alpha$(각각 $\rho_{\alpha\beta}$)를 제한사상}
\[
  \mathcal F(U)\longrightarrow\mathcal F(U_\alpha)
  \qquad
  \bigl(\text{각각 }\mathcal F(U_\alpha)
  \longrightarrow\mathcal F(U_\alpha\cap U_\beta)\bigr)
\]

\oldpage[0\textsubscript{I}]{24}

\emph{이라 하면, $\mathcal F(U)$와 족 $(\rho_\alpha)$로 이루어진 쌍은 $(\mathcal F(U_\alpha))$, $(\mathcal F(U_\alpha\cap U_\beta))$, $(\rho_{\alpha\beta})$에 관한 보편 문제의 해이다(3.1.1).}%
\footnote{이는 여과되지 않아도 되는 일반적인 \emph{사영극한} 개념의 특수한 경우이다((T, I, 1.8) 및 서론에서 예고한 집필 중인 책 참조).}

이는 $\mathbf K$의 모든 대상 $T$에 대하여 족
\[
  U\longmapsto\operatorname{Hom}(T,\mathcal F(U))
\]
이 \emph{집합의 층}이라고 하는 것과 동치이다.
\end{env}

\begin{env}[3.1.3]
\label{0.3.1.3-ko}
$\mathbf K$가 ``사상을 갖는 구조의 한 종류'' $\Sigma$로 정의되는 범주라고 하자. 따라서 $\mathbf K$의 대상은 $\Sigma$형 구조가 주어진 집합들이고 사상은 $\Sigma$의 사상들이다. 더 나아가 범주 $\mathbf K$가 다음 조건을 만족한다고 하자.

\smallskip
\noindent
\textup{(E)} $(A,(\rho_\alpha))$가 족 $(A_\alpha)$, $(A_{\alpha\beta})$, $(\rho_{\alpha\beta})$에 관하여 \emph{범주 $\mathbf K$ 안에서} 보편 사상 문제의 해이면, 같은 족에 관하여 \emph{집합의 범주 안에서도} 보편 사상 문제의 해이다. 여기서는 $A$, $A_\alpha$, $A_{\alpha\beta}$를 집합으로, $\rho_\alpha$, $\rho_{\alpha\beta}$를 사상으로 본다.%
\footnote{이는 표준 함자 $\mathbf K\to(\mathrm{Ens})$가 여과될 필요가 없는 \emph{사영극한과 교환한다}는 뜻이기도 함을 증명할 수 있다.}

이 조건 아래에서 (F)는 $U\mapsto\mathcal F(U)$를 \emph{집합의 준층}으로 보았을 때 \emph{층}이 되게 한다. 또한 (F)에 따라 사상 $u:T\to\mathcal F(U)$가 $\mathbf K$의 사상이기 위한 필요충분조건은 각 합성 $\rho_\alpha\circ u$가 $T\to\mathcal F(U_\alpha)$인 $\mathbf K$의 사상인 것이다. 이는 $\mathcal F(U)$ 위의 $\Sigma$형 구조가 사상 $\rho_\alpha$들에 관한 \emph{시초 구조}라는 뜻이다. 역으로 $X$ 위에서 $\mathbf K$에 값을 갖는 준층 $U\mapsto\mathcal F(U)$가 \emph{집합의 층}이고 앞의 조건을 만족한다고 하자. 그러면 이 준층은 (F)를 만족하므로 \emph{$\mathbf K$에 값을 갖는 층}이다.
\end{env}

\begin{env}[3.1.4]
\label{0.3.1.4-ko}
$\Sigma$가 군 또는 환 구조의 종류이면, $\mathbf K$에 값을 갖는 준층 $U\mapsto\mathcal F(U)$가 \emph{집합의 층}이라는 사실만으로 곧바로 \emph{$\mathbf K$에 값을 갖는 층}, 즉 (G)의 의미에서 군의 층 또는 환의 층이 된다.%
\footnote{이는 범주 $\mathbf K$에서 집합의 사상으로서 \emph{전단사}인 모든 사상이 \emph{동형}이기 때문이다. 예를 들어 $\mathbf K$가 위상공간의 범주이면 더 이상 그렇지 않다.}
그러나 예를 들어 $\mathbf K$가 연속 준동형을 사상으로 하는 \emph{위상환}의 범주이면 더 이상 그렇지 않다. $\mathbf K$에 값을 갖는 층은 다음 조건을 만족하는 환의 층 $U\mapsto\mathcal F(U)$이다. 모든 열린집합 $U$와 $U$에 포함된 열린집합 $U_\alpha$들에 의한 모든 덮개에 대하여, 환 $\mathcal F(U)$의 위상은 준동형
\[
  \mathcal F(U)\longrightarrow\mathcal F(U_\alpha)
\]
들을 연속이 되게 하는 \emph{가장 약한} 위상이다. 이 경우 위상을 잊고 환의 층으로 본 $U\mapsto\mathcal F(U)$를 위상환의 층 $U\mapsto\mathcal F(U)$의 \emph{바탕} 층이라 한다. 따라서 위상환의 층 사이의 사상 $u_V:\mathcal F(V)\to\mathcal G(V)$($V$는 $X$의 임의의 열린집합)은 바탕 환의 층 사이의 준동형으로서, 모든 열린집합 $V\subset X$에 대하여 $u_V$가 \emph{연속}인 것들이다. 이를 바탕 환의 층 사이의 임의의 준동형과 구별하기 위해 위상환의 층 사이의 \emph{연속} 준동형이라 부르겠다. 위상공간의 층이나 위상군의 층에 대해서도 비슷한 정의와 규약을 사용한다.
\end{env}

\oldpage[0\textsubscript{I}]{25}

\begin{env}[3.1.5]
\label{0.3.1.5-ko}
임의의 범주 $\mathbf K$에 대하여 $\mathcal F$가 $X$ 위에서 $\mathbf K$에 값을 갖는 준층(각각 층)이고 $U$가 $X$의 열린집합이면, 열린집합 $V\subset U$에 대한 $\mathcal F(V)$들은 $\mathbf K$에 값을 갖는 준층(각각 층)을 이룬다. 이를 $U$ 위에서 $\mathcal F$가 \emph{유도하는} 준층(각각 층)이라 하고 $\mathcal F|U$로 나타낸다.

$X$ 위에서 $\mathbf K$에 값을 갖는 준층 사이의 모든 사상 $u:\mathcal F\to\mathcal G$에 대하여, $V\subset U$인 $u_V$들로 이루어진 사상 $\mathcal F|U\to\mathcal G|U$를 $u|U$로 나타낸다.
\end{env}

\begin{env}[3.1.6]
\label{0.3.1.6-ko}
이제 범주 $\mathbf K$가 \emph{귀납극한}을 갖는다고 하자(T, 1.8). 그러면 $X$ 위에서 $\mathbf K$에 값을 갖는 모든 준층(특히 모든 층) $\mathcal F$와 모든 $x\in X$에 대하여, $\mathbf K$의 대상인 \emph{줄기} $\mathcal F_x$를 다음 귀납극한으로 정의할 수 있다. $x$의 열린 근방 $U$들의 집합을 $\supset$에 관한 여과집합으로 보고, 사상 $\rho_U^V:\mathcal F(V)\to\mathcal F(U)$에 따른 $\mathcal F(U)$들의 귀납극한을 취한다. $u:\mathcal F\to\mathcal G$가 $\mathbf K$에 값을 갖는 준층 사이의 사상이면, 모든 $x\in X$에 대하여 사상
\[
  u_x:\mathcal F_x\longrightarrow\mathcal G_x
\]
를 $x$의 열린 근방들의 집합을 따라 $u_U:\mathcal F(U)\to\mathcal G(U)$들의 귀납극한으로 정의한다. 이렇게 하여 모든 $x\in X$에 대하여 $\mathcal F_x$는 $\mathcal F$에 관한 공변함자로 정의되며 $\mathbf K$에 값을 갖는다.

더 나아가 $\mathbf K$가 사상을 갖는 구조의 종류 $\Sigma$로 정의되는 경우, $\mathbf K$에 값을 갖는 \emph{층 $\mathcal F$}의 $U$ \emph{위의 단면}은 여전히 $\mathcal F(U)$의 원소를 뜻하며, 이때 $\mathcal F(U)$ 대신 $\Gamma(U,\mathcal F)$라 쓴다. $s\in\Gamma(U,\mathcal F)$이고 $V$가 $U$에 포함된 열린집합이면 $\rho_V^U(s)$ 대신 $s|V$라 쓴다. 모든 $x\in U$에 대하여 $\mathcal F_x$에서 $s$의 표준 상을 점 $x$에서 $s$의 \emph{싹}이라 하고 $s_x$로 나타낸다. 이 의미로는 \emph{표기 $s(x)$를 결코 사용하지 않겠다}. 이 표기는 이 논고에서 뒤에 다룰 특정 층에 관한 다른 개념을 위해 남겨 둔다(5.5.1).

$u:\mathcal F\to\mathcal G$가 $\mathbf K$에 값을 갖는 층 사이의 사상이면, 모든 $s\in\Gamma(U,\mathcal F)$에 대하여 $u_V(s)$ 대신 $u(s)$라 쓰겠다.

$\mathcal F$가 가환군, 환 또는 가군의 층이면 $\mathcal F_x\neq\{0\}$인 $x\in X$들의 집합을 $\mathcal F$의 \emph{지지집합}이라 하고 $\operatorname{Supp}(\mathcal F)$로 나타낸다. 이 집합은 $X$에서 반드시 닫혀 있지는 않다.

$\mathbf K$가 사상을 갖는 구조의 한 종류로 정의되는 경우, $\mathbf K$에 값을 갖는 층에 관해서는 \emph{``에탈레 공간''의 관점을 체계적으로 사용하지 않겠다}. 다시 말해 층을 위상공간으로, 심지어 그 줄기들의 합집합인 집합으로도 보지 않을 것이며, $X$ 위의 이러한 층 사이의 사상 $u:\mathcal F\to\mathcal G$를 위상공간 사이의 연속사상으로 보지도 않을 것이다.
\end{env}

\subsection{열린집합의 기저 위의 준층}
\label{subsection:0.3.2-ko}

\begin{env}[3.2.1]
\label{0.3.2.1-ko}
앞으로는 (일반화된, 즉 반드시 여과될 필요가 없는 준순서집합에 대응하는) \emph{사영극한}을 갖는 범주 $\mathbf K$만을 다루겠다((T, 1.8) 참조). $X$를 위상공간, $\mathfrak B$를 $X$의 위상의 열린집합 기저라 하자. 각 $U\in\mathfrak B$에 대응하는 대상 $\mathcal F(U)\in\mathbf K$의 족과, $U\subset V$인 $\mathfrak B$의 모든 원소쌍 $(U,V)$에 대하여 정의된 사상
\[
  \rho_U^V:\mathcal F(V)\longrightarrow\mathcal F(U)
\]
들의 족을 \emph{$\mathbf K$에 값을 갖는 $\mathfrak B$ 위의 준층}이라 한다.

\oldpage[0\textsubscript{I}]{26}

여기에는 $\rho_U^U=\text{항등사상}$이라는 조건과, $U\subset V\subset W$인 $\mathfrak B$의 원소 $U,V,W$에 대하여
\[
  \rho_U^W=\rho_U^V\circ\rho_V^W
\]
라는 조건을 부과한다. 이 준층에는 보통 의미에서 $\mathbf K$에 값을 갖는 준층 $U\mapsto\mathcal F'(U)$를 다음과 같이 대응시킬 수 있다. 모든 열린집합 $U$에 대하여
\[
  \mathcal F'(U)=\varprojlim_{V\subset U}\mathcal F(V)
\]
로 놓는다. 여기서 $V$는 $V\subset U$인 $V\in\mathfrak B$들의 $\subset$에 관한 순서집합을 움직이며, 이 순서집합은 일반적으로 \emph{여과되지 않는다}. 실제로 $\mathcal F(V)$들은 $V\subset W\subset U$, $V,W\in\mathfrak B$일 때의 $\rho_V^W$에 관하여 사영계를 이룬다. $U\subset U'$인 $X$의 두 열린집합 $U,U'$가 주어지면, $\rho'_U{}^{U'}$를 $V\subset U$에 관한 표준 사상 $\mathcal F'(U')\to\mathcal F(V)$들의 사영극한으로 정의한다. 다시 말해 이는 표준 사상 $\mathcal F'(U)\to\mathcal F(V)$들과 합성했을 때 표준 사상 $\mathcal F'(U')\to\mathcal F(V)$가 되는 유일한 사상
\[
  \mathcal F'(U')\longrightarrow\mathcal F'(U)
\]
이다. 그러면 $\rho'_U{}^{U'}$의 추이성은 곧바로 확인된다. 더 나아가 $U\in\mathfrak B$이면 표준 사상 $\mathcal F'(U)\to\mathcal F(U)$는 동형이므로 두 대상을 동일시할 수 있다.%
\footnote{$X$가 \emph{뇌터 공간}이면, $\mathbf K$가 \emph{유한} 사영계의 사영극한만을 갖는다고 가정해도 $\mathcal F'(U)$를 정의하고 이것이 보통 의미의 준층임을 보일 수 있다. 실제로 $U$가 $X$의 임의의 열린집합이면 $\mathfrak B$의 원소들로 이루어진 $U$의 \emph{유한} 덮개 $(V_i)$가 존재한다. 각 지표쌍 $(i,j)$에 대하여 $(V_{ijk})$를 $\mathfrak B$의 원소들로 이루어진 $V_i\cap V_j$의 유한 덮개라 하자. $I$를 지표 $i$와 삼중항 $(i,j,k)$들로 이루어진 집합으로 하고, 오직 $i>(i,j,k)$와 $j>(i,j,k)$라는 관계로 순서를 주자. 이때 $\mathcal F'(U)$를 $\mathcal F(V_i)$와 $\mathcal F(V_{ijk})$로 이루어진 계의 사영극한으로 취한다. 이것이 덮개 $(V_i)$와 $(V_{ijk})$에 의존하지 않으며 $U\mapsto\mathcal F'(U)$가 준층임은 쉽게 확인된다.}
\end{env}

\begin{env}[3.2.2]
\label{0.3.2.2-ko}
이렇게 정의한 준층 $\mathcal F'$가 \emph{층}이기 위한 필요충분조건은 $\mathfrak B$ 위의 준층 $\mathcal F$가 다음 조건을 만족하는 것이다.

\smallskip
\noindent
\textup{(F\textsubscript{0})} \emph{$U\in\mathfrak B$를 $U$에 포함되는 $U_\alpha\in\mathfrak B$들로 덮는 모든 덮개 $(U_\alpha)$와 모든 대상 $T\in\mathbf K$에 대하여, 각 $f\in\operatorname{Hom}(T,\mathcal F(U))$에 족}
\[
  (\rho_{U_\alpha}^{U}\circ f)
  \in\prod_\alpha\operatorname{Hom}(T,\mathcal F(U_\alpha))
\]
\emph{을 대응시키는 사상은 $\operatorname{Hom}(T,\mathcal F(U))$에서 다음 조건을 만족하는 족 $(f_\alpha)$들의 집합 위로 가는 전단사이다. 모든 지표쌍 $(\alpha,\beta)$와 $V\subset U_\alpha\cap U_\beta$인 모든 $V\in\mathfrak B$에 대하여}
\[
  \rho_V^{U_\alpha}\circ f_\alpha
  =\rho_V^{U_\beta}\circ f_\beta.
\]%
\footnote{이는 $\mathcal F(U)$와 $\rho_\alpha=\rho_{U_\alpha}^{U}$들로 이루어진 쌍이 (3.1.1)에서 다음 자료로 정의한 보편 문제의 해라는 뜻이기도 하다. $A_\alpha=\mathcal F(U_\alpha)$이고,
$A_{\alpha\beta}=\prod\mathcal F(V)$이다. 여기서 곱은 $V\subset U_\alpha\cap U_\beta$인 $V\in\mathfrak B$ 전체에 걸친다. 또한
$\rho_{\alpha\beta}=(\rho_V''):\mathcal F(U_\alpha)\to\prod\mathcal F(V)$이며, 이는 $V,V',W\in\mathfrak B$, $V\cup V'\subset U_\alpha\cap U_\beta$, $W\subset V\cap V'$일 때
$\rho_W^V\circ\rho_V''=\rho_W^{V'}\circ\rho_{V'}''$가 되도록 정의한다.}

이 조건이 필요함은 자명하다. 충분함을 보이기 위해 먼저 $\mathfrak B$에 \emph{포함되는} $X$의 위상의 두 번째 기저 $\mathfrak B'$를 생각하자. 부분족 $(\mathcal F(V))_{V\in\mathfrak B'}$에서 얻은 준층을 $\mathcal F''$라 하면 $\mathcal F''$는 $\mathcal F'$와 \emph{표준적으로 동형}임을 보이겠다. 우선 $V\in\mathfrak B'$, $V\subset U$에 관한 표준 사상 $\mathcal F'(U)\to\mathcal F(V)$들의 사영극한은 모든 열린집합 $U$에 대하여 사상 $\mathcal F'(U)\to\mathcal F''(U)$를 준다. $U\in\mathfrak B$이면 이 사상은 동형이다. 실제로 가정에 따라 $V\in\mathfrak B'$, $V\subset U$에 대한 표준 사상 $\mathcal F''(U)\to\mathcal F(V)$들은
\[
  \mathcal F''(U)\longrightarrow\mathcal F(U)
  \longrightarrow\mathcal F(V)
\]
로 분해되며, 이렇게 정의된 사상 $\mathcal F(U)\to\mathcal F''(U)$와 $\mathcal F''(U)\to\mathcal F(U)$의 두 합성이 항등사상임은 곧바로 알 수 있다. 이제 모든 열린집합 $U$에 대하여, $W\in\mathfrak B$, $W\subset U$일 때의 사상 $\mathcal F''(U)\to\mathcal F''(W)=\mathcal F(W)$들은 $\mathcal F(W)$들의 사영극한을 특징짓는 조건을 만족한다. 따라서 사영극한이 동형을 제외하고 유일하다는 사실에 의해 주장이 증명된다.

이제 $U$를 $X$의 임의의 열린집합, $(U_\alpha)$를 $U$에 포함되는 열린집합들에 의한 $U$의 덮개라 하고, $\mathfrak B'$를 다음 원소들로 이루어진 부분족이라 하자.

\oldpage[0\textsubscript{I}]{27}

그 원소들은 $\mathfrak B$에 속하며 적어도 하나의 $U_\alpha$에 포함된다. $\mathfrak B'$도 여전히 $X$의 위상의 기저임은 명백하다. 따라서 $\mathcal F'(U)$는 $V\in\mathfrak B'$, $V\subset U$인 $\mathcal F(V)$들의 사영극한이고, $\mathcal F''(U_\alpha)$는 $V\in\mathfrak B'$, $V\subset U_\alpha$인 $\mathcal F(V)$들의 사영극한이다. 그러므로 사영극한의 정의에 의해 공리 (F)가 곧바로 성립한다.

(F\textsubscript{0})가 성립할 때에는 용어를 남용하여 $\mathfrak B$ 위의 준층 $\mathcal F$를 \emph{층}이라 하겠다.
\end{env}

\begin{env}[3.2.3]
\label{0.3.2.3-ko}
$\mathcal F$, $\mathcal G$를 $\mathbf K$에 값을 갖는 $\mathfrak B$ 위의 두 준층이라 하자. \emph{사상} $u:\mathcal F\to\mathcal G$를 제한사상 $\rho_V^W$와의 통상적인 호환조건을 만족하는 사상
\[
  u_V:\mathcal F(V)\longrightarrow\mathcal G(V)
  \qquad(V\in\mathfrak B)
\]
들의 족 $(u_V)_{V\in\mathfrak B}$로 정의한다. (3.2.1)의 표기에 따라, $V\in\mathfrak B$, $V\subset U$인 $u_V$들의 사영극한을 $u_U'$로 취하면 대응하는 보통 준층 사이의 사상 $u':\mathcal F'\to\mathcal G'$를 얻는다. $\rho'_U{}^{U'}$와의 호환조건은 사영극한의 함자성에서 따른다.
\end{env}

\begin{env}[3.2.4]
\label{0.3.2.4-ko}
범주 $\mathbf K$가 귀납극한을 갖고 $\mathcal F$가 $\mathbf K$에 값을 갖는 $\mathfrak B$ 위의 준층이라고 하자. 모든 $x\in X$에 대하여 $\mathfrak B$에 속하는 $x$의 근방들은 $x$의 모든 근방들의 집합에서 $\supset$에 관한 공종집합을 이룬다. 따라서 $\mathcal F'$가 $\mathcal F$에 대응하는 보통 준층이면 줄기 $\mathcal F_x'$는
\[
  \varinjlim_{\mathfrak B}\mathcal F(V)
\]
와 같다. 여기서 극한은 $x$를 포함하는 $V\in\mathfrak B$들의 집합을 따라 취한다. $u:\mathcal F\to\mathcal G$가 $\mathbf K$에 값을 갖는 $\mathfrak B$ 위의 준층 사이의 사상이고 $u':\mathcal F'\to\mathcal G'$가 대응하는 보통 준층 사이의 사상이면, $u_x'$도 마찬가지로 $V\in\mathfrak B$, $x\in V$인 사상 $u_V:\mathcal F(V)\to\mathcal G(V)$들의 귀납극한이다.
\end{env}

\begin{env}[3.2.5]
\label{0.3.2.5-ko}
(3.2.1)의 일반적인 조건으로 돌아가자. $\mathcal F$가 $\mathbf K$에 값을 갖는 보통 \emph{층}이고 $\mathcal F_1$이 $\mathcal F$를 $\mathfrak B$에 제한하여 얻은 \emph{$\mathfrak B$ 위의 층}이면, (3.2.1)의 절차로 $\mathcal F_1$에서 얻는 보통 층 $\mathcal F_1'$은 조건 (F)와 사영극한의 유일성에 따라 $\mathcal F$와 표준적으로 동형이다. 보통 $\mathcal F$와 $\mathcal F_1'$을 동일시하겠다.

$\mathcal G$가 $X$ 위에서 $\mathbf K$에 값을 갖는 두 번째 보통 층이고 $u:\mathcal F\to\mathcal G$가 사상이라 하자. 앞의 설명에 따르면 \emph{$V\in\mathfrak B$인 $u_V:\mathcal F(V)\to\mathcal G(V)$들만} 주어도 $u$가 완전히 결정된다. 역으로 $V\in\mathfrak B$인 $u_V$들이 주어졌을 때, $V,W\in\mathfrak B$, $V\subset W$인 제한사상 $\rho_V^W$들과의 가환도만 확인하면, 모든 $V\in\mathfrak B$에 대하여 $u_V'=u_V$인 $\mathcal F$에서 $\mathcal G$로 가는 사상 $u'$가 유일하게 존재한다(3.2.3).
\end{env}

\begin{env}[3.2.6]
\label{0.3.2.6-ko}
계속해서 $\mathbf K$가 사영극한을 갖는다고 하자. 그러면 \emph{$\mathbf K$에 값을 갖는 $X$ 위의 층}들의 범주도 \emph{사영극한}을 갖는다. 실제로 $(\mathcal F_\lambda)$가 $X$ 위에서 $\mathbf K$에 값을 갖는 층들의 사영계이면
\[
  \mathcal F(U)=\varprojlim_\lambda\mathcal F_\lambda(U)
\]
들은 $\mathbf K$에 값을 갖는 준층을 정의한다. 공리 (F)는 사영극한의 추이성으로부터 성립하며, 이때 $\mathcal F$가 $\mathcal F_\lambda$들의 사영극한이라는 사실은 자명하다.

$\mathbf K$가 집합의 범주일 때, 다음 조건을 만족하는 모든 사영계 $(\mathcal H_\lambda)$에 대하여

\oldpage[0\textsubscript{I}]{28}

모든 $\lambda$에 대하여 $\mathcal H_\lambda$는 $\mathcal F_\lambda$의 \emph{부분층}이고, $\varprojlim_\lambda\mathcal H_\lambda$는 $\varprojlim_\lambda\mathcal F_\lambda$의 \emph{부분층}과 표준적으로 동일시된다. $\mathbf K$가 가환군의 범주이면 공변함자 $\varprojlim_\lambda\mathcal F_\lambda$는 \emph{가법적이고 왼쪽 완전}이다.
\end{env}

\subsection{층의 붙이기}
\label{subsection:0.3.3-ko}

\begin{env}[3.3.1]
\label{0.3.3.1-ko}
계속해서 범주 $\mathbf K$가 일반화된 사영극한을 갖는다고 하자. $X$를 위상공간, $\mathfrak U=(U_\lambda)_{\lambda\in L}$를 $X$의 열린 덮개라 하고, 각 $\lambda\in L$에 대하여 $\mathcal F_\lambda$를 $U_\lambda$ 위에서 $\mathbf K$에 값을 갖는 층이라 하자. 모든 지표쌍 $(\lambda,\mu)$에 대하여 \emph{동형}
\[
  \theta_{\lambda\mu}:
  \mathcal F_\mu|(U_\lambda\cap U_\mu)
  \xrightarrow{\sim}
  \mathcal F_\lambda|(U_\lambda\cap U_\mu)
\]
이 주어졌다고 하자. 더 나아가 모든 삼중항 $(\lambda,\mu,\nu)$에 대하여 $\theta'_{\lambda\mu}$, $\theta'_{\mu\nu}$, $\theta'_{\lambda\nu}$를 각각 $\theta_{\lambda\mu}$, $\theta_{\mu\nu}$, $\theta_{\lambda\nu}$의 $U_\lambda\cap U_\mu\cap U_\nu$로의 제한이라 할 때
\[
  \theta'_{\lambda\nu}
  =\theta'_{\lambda\mu}\circ\theta'_{\mu\nu}
\]
가 성립한다고 하자. 이를 $\theta_{\lambda\mu}$들에 관한 \emph{붙이기 조건}이라 한다. 그러면 $X$ 위에서 $\mathbf K$에 값을 갖는 층 $\mathcal F$와 각 $\lambda$에 대한 동형
\[
  \eta_\lambda:\mathcal F|U_\lambda
  \xrightarrow{\sim}\mathcal F_\lambda
\]
가 존재하여 다음 조건을 만족한다. 모든 지표쌍 $(\lambda,\mu)$에 대하여 $\eta'_\lambda$, $\eta'_\mu$를 각각 $\eta_\lambda$, $\eta_\mu$의 $U_\lambda\cap U_\mu$로의 제한이라 하면
\[
  \theta_{\lambda\mu}
  =\eta'_\lambda\circ{\eta'_\mu}^{-1}.
\]
더 나아가 이 조건들에 의해 $\mathcal F$와 $\eta_\lambda$들은 유일한 동형을 제외하고 결정된다. 유일성은 (3.2.5)에서 곧바로 따른다. $\mathcal F$의 존재를 보이기 위해, 적어도 하나의 $U_\lambda$에 포함되는 열린집합들로 이루어진 열린집합 기저를 $\mathfrak B$라 하자. 모든 $U\in\mathfrak B$에 대하여 $U\subset U_\lambda$인 $\lambda$ 하나를 골라, 그 $\mathcal F_\lambda(U)$ 가운데 하나를 힐베르트의 $\tau$ 함수로 선택한다. 이 대상을 $\mathcal F(U)$라 하면 $U\subset V$, $U,V\in\mathfrak B$인 $\rho_U^V$들은 $\theta_{\lambda\mu}$를 사용하여 명백한 방식으로 정의된다. 추이성 조건은 붙이기 조건에서 따르고, (F\textsubscript{0})도 곧바로 확인된다. 따라서 이렇게 정의한 $\mathfrak B$ 위의 준층은 층이며, 일반적인 절차 (3.2.1)에 의해 요구한 성질을 갖는 보통 층을 얻는다. 이 층도 $\mathcal F$라 쓰겠다. $\mathcal F$는 \emph{$\theta_{\lambda\mu}$들을 사용하여 $\mathcal F_\lambda$들을 붙여 얻었다}고 말하며, 보통 $\eta_\lambda$를 통해 $\mathcal F_\lambda$와 $\mathcal F|U_\lambda$를 동일시하겠다.

$X$ 위에서 $\mathbf K$에 값을 갖는 모든 층 $\mathcal F$는 다음과 같이 붙여 얻은 것으로 볼 수 있음이 명백하다. $X$의 임의의 열린 덮개 $(U_\lambda)$에 대하여 $\mathcal F_\lambda=\mathcal F|U_\lambda$로 놓고, $\theta_{\lambda\mu}$를 항등사상으로 취한다.
\end{env}

\begin{env}[3.3.2]
\label{0.3.3.2-ko}
같은 표기 아래에서, 모든 $\lambda\in L$에 대하여 $\mathcal G_\lambda$를 $U_\lambda$ 위에서 $\mathbf K$에 값을 갖는 두 번째 층이라 하자. 모든 지표쌍 $(\lambda,\mu)$에 대하여 붙이기 조건을 만족하는 동형
\[
  \omega_{\lambda\mu}:
  \mathcal G_\mu|(U_\lambda\cap U_\mu)
  \xrightarrow{\sim}
  \mathcal G_\lambda|(U_\lambda\cap U_\mu)
\]
이 주어졌다고 하자. 끝으로 모든 $\lambda$에 대하여 사상 $u_\lambda:\mathcal F_\lambda\to\mathcal G_\lambda$가 주어지고, 가환도
\begin{equation}
\begin{tikzcd}[column sep=huge,row sep=large]
\mathcal F_\mu|(U_\lambda\cap U_\mu)
  \arrow[r,"u_\mu"] \arrow[d]
  & \mathcal G_\mu|(U_\lambda\cap U_\mu) \arrow[d] \\
\mathcal F_\lambda|(U_\lambda\cap U_\mu)
  \arrow[r,"u_\lambda"']
  & \mathcal G_\lambda|(U_\lambda\cap U_\mu)
\end{tikzcd}
\tag{3.3.2.1}\label{0.3.3.2.1-ko}
\end{equation}
들이 가환이라고 하자. $\mathcal G$가 $\omega_{\lambda\mu}$들을 사용하여 $\mathcal G_\lambda$들을 붙여 얻은 층이면, 사상 $u:\mathcal F\to\mathcal G$가 유일하게 존재한다. 그 조건은 다음 가환도들이

\oldpage[0\textsubscript{I}]{29}

\[
\begin{tikzcd}[column sep=huge,row sep=large]
\mathcal F|U_\lambda
  \arrow[r,"u|U_\lambda"] \arrow[d]
  & \mathcal G|U_\lambda \arrow[d] \\
\mathcal F_\lambda
  \arrow[r,"u_\lambda"']
  & \mathcal G_\lambda
\end{tikzcd}
\]
가환인 것이다. 이는 (3.2.3)에서 곧바로 따른다. 족 $(u_\lambda)$와 $u$ 사이의 대응은 조건 (3.3.2.1)을 만족하는
\[
  \prod_\lambda
  \operatorname{Hom}(\mathcal F_\lambda,\mathcal G_\lambda)
\]
의 부분에서 $\operatorname{Hom}(\mathcal F,\mathcal G)$ 위로 가는 함자적 전단사이다.
\end{env}

\begin{env}[3.3.3]
\label{0.3.3.3-ko}
(3.3.1)의 표기에서 $V$를 $X$의 열린집합이라 하자. $\theta_{\lambda\mu}$들의 $V\cap U_\lambda\cap U_\mu$로의 제한들은 유도된 층 $\mathcal F_\lambda|(V\cap U_\lambda)$들에 관한 붙이기 조건을 만족한다. 이 층들을 붙여 얻은 $V$ 위의 층은 $\mathcal F|V$와 표준적으로 동일시됨이 곧바로 따른다.
\end{env}

\subsection{준층의 직상}
\label{subsection:0.3.4-ko}

\begin{env}[3.4.1]
\label{0.3.4.1-ko}
$X$, $Y$를 두 위상공간, $\psi:X\to Y$를 연속사상이라 하자. $\mathcal F$를 $X$ 위에서 범주 $\mathbf K$에 값을 갖는 준층이라 하자. 모든 열린집합 $U\subset Y$에 대하여
\[
  \mathcal G(U)=\mathcal F(\psi^{-1}(U))
\]
로 놓고, $U\subset V$인 $Y$의 두 열린 부분집합 $U,V$에 대하여 $\rho_U^V$를 사상
\[
  \mathcal F(\psi^{-1}(V))
  \longrightarrow\mathcal F(\psi^{-1}(U))
\]
이라 하자. $\mathcal G(U)$들과 $\rho_U^V$들은 $Y$ 위에서 $\mathbf K$에 값을 갖는 \emph{준층}을 정의함이 곧바로 따른다. 이를 \emph{$\psi$에 의한 $\mathcal F$의 직상}이라 하고 $\psi_*(\mathcal F)$로 나타낸다. $\mathcal F$가 층이면 준층 $\mathcal G$에 관한 공리 (F)가 곧바로 확인되므로 $\psi_*(\mathcal F)$도 층이다.
\end{env}

\begin{env}[3.4.2]
\label{0.3.4.2-ko}
$\mathcal F_1$, $\mathcal F_2$를 $X$ 위에서 $\mathbf K$에 값을 갖는 두 준층이라 하고, $u:\mathcal F_1\to\mathcal F_2$를 사상이라 하자. $U$가 $Y$의 열린 부분집합 전체를 움직일 때 사상족
\[
  u_{\psi^{-1}(U)}:
  \mathcal F_1(\psi^{-1}(U))
  \longrightarrow\mathcal F_2(\psi^{-1}(U))
\]
은 제한사상과의 호환조건을 만족하므로 사상
\[
  \psi_*(u):\psi_*(\mathcal F_1)
  \longrightarrow\psi_*(\mathcal F_2)
\]
를 정의한다. $\mathcal F_3$가 $X$ 위에서 $\mathbf K$에 값을 갖는 세 번째 준층이고 $v:\mathcal F_2\to\mathcal F_3$가 사상이면
\[
  \psi_*(v\circ u)=\psi_*(v)\circ\psi_*(u).
\]
다시 말해 $\psi_*(\mathcal F)$는 $\mathcal F$에 관한 \emph{공변함자}이며, $X$ 위에서 $\mathbf K$에 값을 갖는 준층(각각 층)의 범주에서 $Y$ 위에서 $\mathbf K$에 값을 갖는 준층(각각 층)의 범주로 간다.
\end{env}

\begin{env}[3.4.3]
\label{0.3.4.3-ko}
$Z$를 세 번째 위상공간, $\psi':Y\to Z$를 연속사상이라 하고 $\psi''=\psi'\circ\psi$로 놓자. $X$ 위에서 $\mathbf K$에 값을 갖는 모든 준층 $\mathcal F$에 대하여
\[
  \psi_*''(\mathcal F)
  =\psi_*'(\psi_*(\mathcal F))
\]
임은 명백하다. 또한 이러한 준층 사이의 모든 사상 $u:\mathcal F\to\mathcal G$에 대하여
\[
  \psi_*''(u)=\psi_*'(\psi_*(u)).
\]
다시 말해 $\psi_*''$은 함자 $\psi_*'$과 $\psi_*$의 \emph{합성}이며, 이를
\[
  (\psi'\circ\psi)_*=\psi_*'\circ\psi_*
\]
로 쓸 수 있다.

더 나아가 $Y$의 모든 열린집합 $U$에 대하여, 유도된 준층 $\mathcal F|\psi^{-1}(U)$의 제한사상 $\psi|\psi^{-1}(U)$에 의한 직상은 바로 유도된 준층 $\psi_*(\mathcal F)|U$이다.
\end{env}

\begin{env}[3.4.4]
\label{0.3.4.4-ko}
범주 $\mathbf K$가 귀납극한을 갖는다고 하고, $\mathcal F$를 $X$ 위에서 $\mathbf K$에 값을 갖는 준층이라 하자. 모든 $x\in X$에 대하여 사상
\[
  \Gamma(\psi^{-1}(U),\mathcal F)
  \longrightarrow\mathcal F_x
\]
들은 귀납계를 이룬다. 여기서 $U$는 $Y$에서 $\psi(x)$의 열린 근방을 움직인다. 이 귀납계에서

\oldpage[0\textsubscript{I}]{30}

귀납극한으로 넘어가면 줄기 사이의 사상
\[
  \psi_x:(\psi_*(\mathcal F))_{\psi(x)}
  \longrightarrow\mathcal F_x
\]
을 얻는다. 일반적으로 이 사상은 \emph{단사도 전사도 아니다}. 이 사상은 함자적이다. 실제로 $u:\mathcal F_1\to\mathcal F_2$가 $X$ 위에서 $\mathbf K$에 값을 갖는 준층 사이의 사상이면 가환도
\[
\begin{tikzcd}[column sep=huge,row sep=large]
(\psi_*(\mathcal F_1))_{\psi(x)}
  \arrow[r,"\psi_x"] \arrow[d,"{(\psi_*(u))_{\psi(x)}}"']
  & (\mathcal F_1)_x \arrow[d,"u_x"] \\
(\psi_*(\mathcal F_2))_{\psi(x)}
  \arrow[r,"\psi_x"']
  & (\mathcal F_2)_x
\end{tikzcd}
\]
는 가환이다. $Z$가 세 번째 위상공간이고 $\psi':Y\to Z$가 연속사상이며 $\psi''=\psi'\circ\psi$이면, 모든 $x\in X$에 대하여
\[
  \psi_x''=\psi_x\circ\psi'_{\psi(x)}.
\]
\end{env}

\begin{env}[3.4.5]
\label{0.3.4.5-ko}
(3.4.4)의 가정 아래에서, 더 나아가 $\psi$가 $X$에서 $Y$의 부분공간 $\psi(X)$ 위로 가는 \emph{위상동형}이라고 하자. 그러면 모든 $x\in X$에 대하여 $\psi_x$는 \emph{동형}이다. 이는 특히 $Y$의 부분집합 $X$에서 $Y$로 가는 표준 단사사상 $j$에 적용된다.
\end{env}

\begin{env}[3.4.6]
\label{0.3.4.6-ko}
$\mathbf K$가 군, 환 등의 범주라고 하자. $\mathcal F$가 $X$ 위에서 $\mathbf K$에 값을 갖고 지지집합이 $S$인 층이며 $y\notin\overline{\psi(S)}$이면, $\psi_*(\mathcal F)$의 정의로부터
\[
  (\psi_*(\mathcal F))_y=\{0\}
\]
을 얻는다. 다시 말해 $\psi_*(\mathcal F)$의 지지집합은 $\overline{\psi(S)}$에 포함되지만 반드시 $\psi(S)$에 포함되는 것은 아니다. 같은 가정 아래에서 $j$가 $Y$의 부분집합 $X$에서 $Y$로 가는 표준 단사사상이면, 층 $j_*(\mathcal F)$가 $X$ 위에 유도하는 층은 $\mathcal F$이다. 더 나아가 $X$가 $Y$에서 \emph{닫혀} 있으면 $j_*(\mathcal F)$는 $X$ 위에 $\mathcal F$를, $Y-X$ 위에 $0$을 유도하는 $Y$ 위의 층이다(G, II, 2.9.2). 그러나 $X$가 국소 닫혔지만 닫히지는 않았다고 가정하면 일반적으로 $j_*(\mathcal F)$는 이 마지막 층과 다르다.
\end{env}

\subsection{준층의 역상}
\label{subsection:0.3.5-ko}

\begin{env}[3.5.1]
\label{0.3.5.1-ko}
(3.4.1)의 가정 아래에서 $\mathcal F$가 $X$ 위의 준층이고 $\mathcal G$가 $Y$ 위의 준층이며 둘 다 $\mathbf K$에 값을 갖는다고 하자. $Y$ 위의 준층 사이의 모든 사상
\[
  u:\mathcal G\longrightarrow\psi_*(\mathcal F)
\]
을 역시 $\mathcal G$에서 $\mathcal F$로 가는 $\psi$-\emph{사상}이라 하고 $\mathcal G\to\mathcal F$로도 쓴다. 또한 $\mathcal G$에서 $\mathcal F$로 가는 $\psi$-사상들의 집합
\[
  \operatorname{Hom}_Y(\mathcal G,\psi_*(\mathcal F))
\]
을 $\operatorname{Hom}_\psi(\mathcal G,\mathcal F)$로 나타낸다.

$U$가 $X$의 열린집합이고 $V$가 $\psi(U)\subset V$를 만족하는 $Y$의 열린집합인 모든 쌍 $(U,V)$에 대하여, 제한사상
\[
  \mathcal F(\psi^{-1}(V))\longrightarrow\mathcal F(U)
\]
과 사상
\[
  u_V:\mathcal G(V)\longrightarrow
  \psi_*(\mathcal F)(V)=\mathcal F(\psi^{-1}(V))
\]
을 합성하여 사상 $u_{U,V}:\mathcal G(V)\to\mathcal F(U)$를 얻는다. 이 사상들은 $U'\subset U$, $V'\subset V$, $\psi(U')\subset V'$일 때 가환도
\begin{equation*}
\begin{tikzcd}[column sep=huge,row sep=large]
\mathcal G(V) \arrow[r,"u_{U,V}"] \arrow[d]
  & \mathcal F(U) \arrow[d] \\
\mathcal G(V') \arrow[r,"u_{U',V'}"']
  & \mathcal F(U')
\end{tikzcd}
\tag{3.5.1.1}\label{0.3.5.1.1-ko}
\end{equation*}
들을 가환이 되게 함이 곧바로 따른다. 역으로 가환도 (3.5.1.1)을 가환이 되게 하는 사상족 $(u_{U,V})$가 주어지면
\[
  u_V=u_{\psi^{-1}(V),V}
\]
로 취함으로써 $\psi$-사상 $u$가 정의된다.

\oldpage[0\textsubscript{I}]{31}

범주 $\mathbf K$가 일반화된 사영극한을 갖고 $\mathfrak B$, $\mathfrak B'$이 각각 $X$, $Y$의 위상의 기저라고 하자. 층 사이의 $\psi$-사상 $u$를 정의하려면 $U\in\mathfrak B$, $V\in\mathfrak B'$, $\psi(U)\subset V$인 $u_{U,V}$들만 주고, $U,U'\in\mathfrak B$, $V,V'\in\mathfrak B'$에 대하여 호환조건 (3.5.1.1)을 확인하면 된다. 실제로 모든 열린집합 $W\subset Y$에 대하여 $u_W$를 $V\in\mathfrak B'$, $V\subset W$, $U\in\mathfrak B$, $\psi(U)\subset V$인 $u_{U,V}$들의 사영극한으로 정의하면 된다.

범주 $\mathbf K$가 귀납극한을 가질 때에는 모든 $x\in X$와 $Y$에서 $\psi(x)$의 모든 열린 근방 $V$에 대하여 사상
\[
  \mathcal G(V)\longrightarrow
  \mathcal F(\psi^{-1}(V))\longrightarrow\mathcal F_x
\]
을 얻는다. 이 사상들은 귀납계를 이루며, 귀납극한으로 넘어가면 사상
\[
  \mathcal G_{\psi(x)}\longrightarrow\mathcal F_x
\]
을 얻는다.
\end{env}

\begin{env}[3.5.2]
\label{0.3.5.2-ko}
(3.4.3)의 가정 아래에서 $\mathcal F$, $\mathcal G$, $\mathcal H$를 각각 $X$, $Y$, $Z$ 위에서 $\mathbf K$에 값을 갖는 준층이라 하고,
\[
  u:\mathcal G\longrightarrow\psi_*(\mathcal F),
  \qquad
  v:\mathcal H\longrightarrow\psi'_*(\mathcal G)
\]
를 각각 $\psi$-사상과 $\psi'$-사상이라 하자. 이들로부터 $\psi''$-사상
\[
  w:\mathcal H\xrightarrow{v}\psi'_*(\mathcal G)
  \xrightarrow{\psi'_*(u)}
  \psi'_*(\psi_*(\mathcal F))=\psi''_*(\mathcal F)
\]
을 얻으며, 이를 정의에 따라 $u$와 $v$의 \emph{합성}이라 한다. 따라서 위상공간 $X$와 그 위에서 $\mathbf K$에 값을 갖는 준층 $\mathcal F$의 쌍 $(X,\mathcal F)$들을 하나의 \emph{범주}를 이루는 것으로 볼 수 있다. 여기서 사상
\[
  (\psi,\theta):(X,\mathcal F)\longrightarrow(Y,\mathcal G)
\]
은 연속사상 $\psi:X\to Y$와 $\psi$-사상 $\theta:\mathcal G\to\mathcal F$의 쌍이다.
\end{env}

\begin{env}[3.5.3]
\label{0.3.5.3-ko}
$\psi:X\to Y$를 연속사상, $\mathcal G$를 $Y$ 위에서 $\mathbf K$에 값을 갖는 \emph{준층}이라 하자. $\psi$에 의한 $\mathcal G$의 \emph{역상}을 쌍 $(\mathcal G',\rho)$로 정의한다. 여기서 $\mathcal G'$은 $X$ 위에서 $\mathbf K$에 값을 갖는 \emph{층}이고, $\rho:\mathcal G\to\mathcal G'$은 $\psi$-사상, 다시 말해 준층의 사상 $\mathcal G\to\psi_*(\mathcal G')$이며, 다음 조건을 만족한다. $X$ 위에서 $\mathbf K$에 값을 갖는 모든 \emph{층} $\mathcal F$에 대하여 $v$를 $\psi_*(v)\circ\rho$로 보내는 사상
\begin{equation*}
  \operatorname{Hom}_X(\mathcal G',\mathcal F)
  \longrightarrow \operatorname{Hom}_\psi(\mathcal G,\mathcal F)
  =\operatorname{Hom}_Y(\mathcal G,\psi_*(\mathcal F))
\tag{3.5.3.1}\label{0.3.5.3.1-ko}
\end{equation*}
이 \emph{전단사}이다. 이 사상은 $\mathcal F$에 관하여 함자적이므로 $\mathcal F$에 관한 함자들의 동형을 정의한다. 쌍 $(\mathcal G',\rho)$은 보편문제의 해이므로, 존재할 때 \emph{유일한 동형을 제외하고 결정된다}. 이때
\[
  \mathcal G'=\psi^*(\mathcal G),
  \qquad
  \rho=\rho_{\mathcal G}
\]
로 쓰고, 표현을 약간 남용하여 $\psi^*(\mathcal G)$를 $\psi$에 의한 $\mathcal G$의 \emph{역상층}이라 하겠다. 여기서는 $\psi^*(\mathcal G)$에 표준 $\psi$-사상
\[
  \rho_{\mathcal G}:\mathcal G\longrightarrow\psi^*(\mathcal G),
\]
즉 $Y$ 위의 준층 사이의 \emph{표준 사상}
\begin{equation*}
  \rho_{\mathcal G}:\mathcal G
  \longrightarrow\psi_*(\psi^*(\mathcal G))
\tag{3.5.3.2}\label{0.3.5.3.2-ko}
\end{equation*}
이 갖추어져 있다고 이해한다.

$v:\psi^*(\mathcal G)\to\mathcal F$가 $X$ 위에서 $\mathbf K$에 값을 갖는 층 사이의 사상이면
\[
  v^b=\psi_*(v)\circ\rho_{\mathcal G}:
  \mathcal G\longrightarrow\psi_*(\mathcal F)
\]
로 놓는다. 정의에 따라 준층 사이의 \emph{모든} 사상 $u:\mathcal G\to\psi_*(\mathcal F)$는 유일한 하나의 $v$에 대하여 $v^b$의 꼴이고, 이 $v$를 $u^\sharp$이라 쓴다. 다시 말해 준층 사이의 모든 사상 $u:\mathcal G\to\psi_*(\mathcal F)$는 유일하게
\begin{equation*}
  u:\mathcal G\xrightarrow{\rho_{\mathcal G}}
  \psi_*(\psi^*(\mathcal G))
  \xrightarrow{\psi_*(u^\sharp)}\psi_*(\mathcal F)
\tag{3.5.3.3}\label{0.3.5.3.3-ko}
\end{equation*}
로 분해된다.
\end{env}

\oldpage[0\textsubscript{I}]{32}

\begin{env}[3.5.4]
\label{0.3.5.4-ko}
이제 범주 $\mathbf K$가 다음 성질을 갖는다고 하자.\footnote{서론에서 인용한 책에서, $\mathbf K$에 값을 갖는 준층의 역상이 존재함을 보장하는 범주 $\mathbf K$에 대한 매우 일반적인 조건을 제시할 것이다.} $Y$ 위에서 $\mathbf K$에 값을 갖는 \emph{모든} 준층 $\mathcal G$는 $\psi$에 의한 역상을 가지며, 이를 $\psi^*(\mathcal G)$로 쓴다.

$\psi^*(\mathcal G)$를 $\mathcal G$에 관한 \emph{공변함자}, 즉 $Y$ 위에서 $\mathbf K$에 값을 갖는 준층의 범주에서 $X$ 위에서 $\mathbf K$에 값을 갖는 층의 범주로 가는 함자로 정의할 수 있음을 보이겠다. 이때 동형 $v\mapsto v^b$는 $\mathcal G$와 $\mathcal F$에 관한 \emph{쌍함자의 동형}
\begin{equation*}
  \operatorname{Hom}_X(\psi^*(\mathcal G),\mathcal F)
  \xrightarrow{\sim}
  \operatorname{Hom}_Y(\mathcal G,\psi_*(\mathcal F))
\tag{3.5.4.1}\label{0.3.5.4.1-ko}
\end{equation*}
이 된다.

실제로 $w:\mathcal G_1\to\mathcal G_2$가 $Y$ 위에서 $\mathbf K$에 값을 갖는 준층 사이의 사상이면 합성사상
\[
  \mathcal G_1\xrightarrow{w}\mathcal G_2
  \xrightarrow{\rho_{\mathcal G_2}}
  \psi_*(\psi^*(\mathcal G_2))
\]
을 생각하자. 이에 대응하는 사상
\[
  (\rho_{\mathcal G_2}\circ w)^\sharp:
  \psi^*(\mathcal G_1)\longrightarrow\psi^*(\mathcal G_2)
\]
을 $\psi^*(w)$라 쓰겠다. (3.5.3.3)에 따라
\begin{equation*}
  \psi_*(\psi^*(w))\circ\rho_{\mathcal G_1}
  =\rho_{\mathcal G_2}\circ w
\tag{3.5.4.2}\label{0.3.5.4.2-ko}
\end{equation*}
이다. 여기서 $\mathcal F$가 $X$ 위에서 $\mathbf K$에 값을 갖는 층이고 $u:\mathcal G_2\to\psi_*(\mathcal F)$가 사상이면, (3.5.3.3), (3.5.4.2), 그리고 $v^b$의 정의로부터
\[
  (u^\sharp\circ\psi^*(w))^b
  =\psi_*(u^\sharp)\circ\psi_*(\psi^*(w))\circ\rho_{\mathcal G_1}
  =\psi_*(u^\sharp)\circ\rho_{\mathcal G_2}\circ w
  =u\circ w,
\]
즉
\begin{equation*}
  (u\circ w)^\sharp=u^\sharp\circ\psi^*(w)
\tag{3.5.4.3}\label{0.3.5.4.3-ko}
\end{equation*}
을 얻는다. 특히 $u$를 사상
\[
  \mathcal G_2\xrightarrow{w'}\mathcal G_3
  \xrightarrow{\rho_{\mathcal G_3}}
  \psi_*(\psi^*(\mathcal G_3))
\]
으로 취하면
\[
  \psi^*(w'\circ w)
  =(\rho_{\mathcal G_3}\circ w'\circ w)^\sharp
  =(\rho_{\mathcal G_3}\circ w')^\sharp\circ\psi^*(w)
  =\psi^*(w')\circ\psi^*(w)
\]
을 얻는다. 이로써 주장이 증명된다.

끝으로 $X$ 위에서 $\mathbf K$에 값을 갖는 모든 층 $\mathcal F$에 대하여 $i_{\mathcal F}$를 $\psi_*(\mathcal F)$의 항등사상이라 하고,
\[
  \sigma_{\mathcal F}:
  \psi^*(\psi_*(\mathcal F))\longrightarrow\mathcal F
\]
를 사상 $(i_{\mathcal F})^\sharp$이라 쓰자. 식 (3.5.4.3)은 모든 사상 $u:\mathcal G\to\psi_*(\mathcal F)$에 대하여 특히 분해
\begin{equation*}
  u^\sharp:\psi^*(\mathcal G)
  \xrightarrow{\psi^*(u)}
  \psi^*(\psi_*(\mathcal F))
  \xrightarrow{\sigma_{\mathcal F}}\mathcal F
\tag{3.5.4.4}\label{0.3.5.4.4-ko}
\end{equation*}
를 준다. 사상 $\sigma_{\mathcal F}$를 \emph{표준적}이라고 하겠다.
\end{env}

\begin{env}[3.5.5]
\label{0.3.5.5-ko}
$\psi':Y\to Z$를 연속사상이라 하고, $Z$ 위에서 $\mathbf K$에 값을 갖는 모든 준층 $\mathcal H$가 $\psi'$에 의한 역상 $\psi'^*(\mathcal H)$을 갖는다고 하자. 그러면 (3.5.4)의 가정 아래에서 $Z$ 위에서 $\mathbf K$에 값을 갖는 모든 준층 $\mathcal H$는 $\psi''=\psi'\circ\psi$에 의한 역상을 가지며, 함자적인 표준 동형
\begin{equation*}
  \psi''^*(\mathcal H)\xrightarrow{\sim}
  \psi^*(\psi'^*(\mathcal H))
\tag{3.5.5.1}\label{0.3.5.5.1-ko}
\end{equation*}
이 존재한다.

\oldpage[0\textsubscript{I}]{33}

이는 $\psi''_*=\psi'_*\circ\psi_*$임을 고려하면 정의에서 곧바로 따른다. 또한 $u:\mathcal G\to\psi_*(\mathcal F)$가 $\psi$-사상이고, $v:\mathcal H\to\psi'_*(\mathcal G)$가 $\psi'$-사상이며, $w=\psi'_*(u)\circ v$가 이들의 합성이라 하자(3.5.2). 그러면 $w^\sharp$은 합성사상
\[
  w^\sharp:\psi^*(\psi'^*(\mathcal H))
  \xrightarrow{\psi^*(v^\sharp)}
  \psi^*(\mathcal G)
  \xrightarrow{u^\sharp}\mathcal F
\]
임이 곧바로 확인된다.
\end{env}

\begin{env}[3.5.6]
\label{0.3.5.6-ko}
특히 $\psi$를 항등사상 $1_X:X\to X$로 취하자. $X$ 위에서 $\mathbf K$에 값을 갖는 준층 $\mathcal F$의 $\psi$에 의한 역상이 존재하면, 이 역상을 \emph{준층 $\mathcal F$의 층화}라 한다. $\mathcal F$에서 $\mathbf K$에 값을 갖는 \emph{층} $\mathcal F'$으로 가는 모든 사상 $u:\mathcal F\to\mathcal F'$은 유일하게
\[
  \mathcal F\xrightarrow{\rho_{\mathcal F}}
  1_X^*(\mathcal F)\xrightarrow{u^\sharp}\mathcal F'
\]
으로 분해된다.
\end{env}

\subsection{상수층과 국소 상수층}
\label{subsection:0.3.6-ko}

\begin{env}[3.6.1]
\label{0.3.6.1-ko}
$X$ 위에서 $\mathbf K$에 값을 갖는 \emph{준층} $\mathcal F$를 생각하자. 모든 공집합이 아닌 열린집합 $U\subset X$에 대하여 표준 사상
\[
  \mathcal F(X)\longrightarrow\mathcal F(U)
\]
이 \emph{동형}이면 $\mathcal F$를 \emph{상수 준층}이라 하겠다. $\mathcal F$가 반드시 층인 것은 아님에 유의하자. 상수 준층의 층화인 \emph{층}을 \emph{상수층}이라 한다(원문의 «단순층»). 층 $\mathcal F$가 \emph{국소 상수}라는 것은 모든 $x\in X$가 $\mathcal F|U$가 상수층이 되는 열린 근방 $U$를 갖는다는 뜻이다.
\end{env}

\begin{env}[3.6.2]
\label{0.3.6.2-ko}
$X$가 \emph{기약}이라고 하자(2.1.1). 그러면 다음 성질들은 서로 동치이다.
\begin{enumerate}
  \item[a)] \emph{$\mathcal F$는 $X$ 위의 상수 준층이다.}
  \item[b)] \emph{$\mathcal F$는 $X$ 위의 상수층이다.}
  \item[c)] \emph{$\mathcal F$는 $X$ 위의 국소 상수층이다.}
\end{enumerate}

실제로 $\mathcal F$를 $X$ 위의 상수 준층이라 하자. $U$, $V$가 $X$의 공집합이 아닌 두 열린집합이면 $U\cap V$도 공집합이 아니다. 따라서
\[
  \mathcal F(X)\longrightarrow\mathcal F(U)
  \longrightarrow\mathcal F(U\cap V)
\]
와 $\mathcal F(X)\to\mathcal F(U)$가 동형이므로 $\mathcal F(U)\to\mathcal F(U\cap V)$도 동형이고, 마찬가지로 $\mathcal F(V)\to\mathcal F(U\cap V)$도 동형이다. 이로부터 (3.1.2)의 공리 (F)가 성립함이 곧바로 따른다. 따라서 $\mathcal F$는 그 층화와 동형이고, a)는 b)를 함의한다.

이제 $(U_\alpha)$를 공집합이 아닌 열린집합들로 이루어진 $X$의 열린 덮개라 하고, $\mathcal F$를 모든 $\alpha$에 대하여 $\mathcal F|U_\alpha$가 상수층인 $X$ 위의 층이라 하자. $U_\alpha$는 기약이므로 앞의 결과에 따라 $\mathcal F|U_\alpha$는 상수 준층이다. $U_\alpha\cap U_\beta$가 공집합이 아니므로
\[
  \mathcal F(U_\alpha)\longrightarrow
  \mathcal F(U_\alpha\cap U_\beta),
  \qquad
  \mathcal F(U_\beta)\longrightarrow
  \mathcal F(U_\alpha\cap U_\beta)
\]
는 동형이다. 따라서 모든 지표쌍에 대하여 표준 동형
\[
  \theta_{\alpha\beta}:
  \mathcal F(U_\alpha)\xrightarrow{\sim}\mathcal F(U_\beta)
\]
을 얻는다. 그런데 $U=X$에 조건 (F)를 적용하면, 모든 지표 $\alpha_0$에 대하여 $\mathcal F(U_{\alpha_0})$와 $\theta_{\alpha_0\alpha}$들이 보편문제의 해임을 알 수 있다. 유일성에 따라
\[
  \mathcal F(X)\longrightarrow\mathcal F(U_{\alpha_0})
\]
는 동형이다. 그러므로 c)는 a)를 함의한다.
\end{env}

\subsection{군 또는 환 준층의 역상}
\label{subsection:0.3.7-ko}

\oldpage[0\textsubscript{I}]{34}

\begin{env}[3.7.1]
\label{0.3.7.1-ko}
$\mathbf K$를 집합의 범주로 취할 때, $\mathbf K$에 값을 갖는 모든 준층 $\mathcal G$는 $\psi$에 의한 역상을 \emph{항상 가짐}을 보이겠다. $X$, $Y$, $\psi$에 관한 표기와 가정은 (3.5.3)의 것을 사용한다. 모든 열린집합 $U\subset X$에 대하여 $\mathcal G'(U)$를 다음과 같이 정의하자. $\mathcal G'(U)$의 원소 $s'$은 족 $(s'_x)_{x\in U}$이고, 모든 $x\in U$에 대하여
\[
  s'_x\in\mathcal G_{\psi(x)}
\]
이며 다음 조건을 만족한다. 모든 $x\in U$에 대하여 $Y$에서 $\psi(x)$의 열린 근방 $V$, $X$에서 $x$의 근방 $W\subset\psi^{-1}(V)\cap U$, 그리고 원소 $s\in\mathcal G(V)$가 존재하여 모든 $z\in W$에 대하여
\[
  s'_z=s_{\psi(z)}
\]
이다. $U\mapsto\mathcal G'(U)$가 \emph{층}의 공리들을 만족함은 곧바로 확인된다.

이제 $\mathcal F$를 $X$ 위의 집합의 층이라 하고
\[
  u:\mathcal G\longrightarrow\psi_*(\mathcal F),
  \qquad
  v:\mathcal G'\longrightarrow\mathcal F
\]
를 사상이라 하자. $u^\sharp$과 $v^b$를 다음과 같이 정의한다. $s'$이 $x\in X$의 근방 $U$ 위의 $\mathcal G'$의 단면이고, $V$가 $\psi(x)$의 열린 근방이며, $s\in\mathcal G(V)$가 $x$의 어떤 근방 $W\subset\psi^{-1}(V)\cap U$에서 모든 $z\in W$에 대하여 $s'_z=s_{\psi(z)}$를 만족한다고 하자. 이때
\[
  u_x^\sharp(s'_x)=u_{\psi(x)}(s_{\psi(x)})
\]
로 취한다. 마찬가지로 $V$가 $Y$의 열린집합이고 $s\in\mathcal G(V)$이면, $v^b(s)$는 $\psi^{-1}(V)$ 위의 $\mathcal G'$의 단면
\[
  s'=(s_{\psi(x)})_{x\in\psi^{-1}(V)}
\]
의 $v$에 의한 상인 $\psi^{-1}(V)$ 위의 $\mathcal F$의 단면이다.

더 나아가 표준 사상 (3.5.3)
\[
  \rho:\mathcal G\longrightarrow\psi_*(\psi^*(\mathcal G))
\]
은 다음과 같이 정의된다. 모든 열린집합 $V\subset Y$와 모든 단면 $s\in\Gamma(V,\mathcal G)$에 대하여 $\rho(s)$는 $\psi^{-1}(V)$ 위의 $\psi^*(\mathcal G)$의 단면
\[
  (s_{\psi(x)})_{x\in\psi^{-1}(V)}
\]
이다. 관계식
\[
  (u^\sharp)^b=u,\qquad
  (v^b)^\sharp=v,\qquad
  v^b=\psi_*(v)\circ\rho
\]
은 곧바로 확인되며, 이로써 주장이 증명된다.

$w:\mathcal G_1\to\mathcal G_2$가 $Y$ 위의 집합의 준층 사이의 사상이면 $\psi^*(w)$는 다음과 같이 구체적으로 주어진다. $s'=(s'_x)_{x\in U}$가 $X$의 열린집합 $U$ 위의 $\psi^*(\mathcal G_1)$의 단면이면
\[
  (\psi^*(w))(s')
  =(w_{\psi(x)}(s'_x))_{x\in U}.
\]
끝으로 $Y$의 모든 열린집합 $V$에 대하여, $\psi$의 $\psi^{-1}(V)$로의 제한에 의한 $\mathcal G|V$의 역상은 유도된 층
\[
  \psi^*(\mathcal G)|\psi^{-1}(V)
\]
과 동일하다.

$\psi$가 항등사상 $1_X$일 때에는 준층에 대응하는 집합의 층의 정의를 다시 얻는다(G, II, 1.2). 앞의 논의는 $\mathbf K$가 군 또는 반드시 가환일 필요는 없는 환의 범주일 때에도 아무런 변경 없이 적용된다.

$X$가 위상공간 $Y$의 임의의 부분집합이고 $j:X\to Y$가 표준 단사사상이라고 하자. 범주 $\mathbf K$에 값을 갖는 $Y$ 위의 모든 층 $\mathcal G$에 대하여, 역상 $j^*(\mathcal G)$이 존재할 때 이를 $\mathcal G$가 $X$ 위에 \emph{유도하는 층}이라 한다. 집합, 군 또는 환의 층에 대해서는 통상적인 정의를 다시 얻는다(G, II, 1.5).
\end{env}

\begin{env}[3.7.2]
\label{0.3.7.2-ko}
(3.5.3)의 표기와 가정을 유지하고, $\mathcal G$가 $Y$ 위의 군의 층(각각 환의 층)이라고 하자. $\psi^*(\mathcal G)$의 단면에 관한 정의 (3.7.1)와 (3.4.4)에 따르면 줄기 사상
\[
  \psi_x\circ\rho_{\psi(x)}:
  \mathcal G_{\psi(x)}
  \longrightarrow(\psi^*(\mathcal G))_x
\]
은 $\mathcal G$에 관하여 함자적인 \emph{동형}이다. 이 동형을 통해 두 줄기를 동일시할 수 있다. 이 동일시 아래에서 $u_x^\sharp$은 (3.5.1)에서 정의한 준동형과 같고, 특히
\[
  \operatorname{Supp}(\psi^*(\mathcal G))
  =\psi^{-1}(\operatorname{Supp}(\mathcal G))
\]
이다.

이 결과로부터 가환군의 층들의 아벨 범주에서 \emph{함자 $\psi^*(\mathcal G)$는 $\mathcal G$에 관하여 완전하다}는 것이 곧바로 따른다.
\end{env}

\subsection{의사이산 공간의 층}
\label{subsection:0.3.8-ko}

\oldpage[0\textsubscript{I}]{35}

\begin{env}[3.8.1]
\label{0.3.8.1-ko}
$X$를 위상이 \emph{준콤팩트} 열린집합들로 이루어진 기저 $\mathfrak B$를 갖는 위상공간이라 하자. $\mathcal F$를 $X$ 위의 \emph{집합의 층}이라 하자. 각 $\mathcal F(U)$에 \emph{이산위상}을 주면
\[
  U\longmapsto\mathcal F(U)
\]
는 \emph{위상공간의 준층}이 된다. 모든 \emph{준콤팩트} 열린집합 $U$에 대하여 $\Gamma(U,\mathcal F')$가 이산공간 $\mathcal F(U)$인, $\mathcal F$에 대응하는 \emph{위상공간의 층 $\mathcal F'$}이 존재함을 보이겠다(3.5.6).

이를 위해서는 이산 위상공간의 준층 $U\mapsto\mathcal F(U)$가 $\mathfrak B$ \emph{위에서} (3.2.2)의 조건 $(\mathrm F_0)$를 만족함을 보이면 충분하다. 더 일반적으로 $U$가 준콤팩트 열린집합이고 $(U_\alpha)$가 $\mathfrak B$의 원소들로 이루어진 $U$의 덮개이면, 사상
\[
  \Gamma(U,\mathcal F)\longrightarrow
  \Gamma(U_\alpha,\mathcal F)
\]
들을 연속이 되게 하는 $\Gamma(U,\mathcal F)$ 위의 가장 엉성한 위상 $\mathcal T$가 \emph{이산위상}임을 보이면 충분하다. 실제로
\[
  U=\bigcup_iU_{\alpha_i}
\]
가 되게 하는 유한개의 지표 $\alpha_i$가 존재한다. $s\in\Gamma(U,\mathcal F)$라 하고 $s_i$를 $\Gamma(U_{\alpha_i},\mathcal F)$에서의 그 상이라 하자. 집합 $\{s_i\}$들의 역상들의 교집합은 정의에 따라 $\mathcal T$에서 $s$의 근방이다. 그런데 $\mathcal F$는 집합의 층이고 $U_{\alpha_i}$들이 $U$를 덮으므로 이 교집합은 $\{s\}$와 같다. 이로써 주장이 증명된다.

$U$가 $X$의 준콤팩트하지 않은 열린집합이면 위상공간 $\Gamma(U,\mathcal F')$의 바탕집합은 여전히 $\Gamma(U,\mathcal F)$이지만, 그 위상은 일반적으로 이산위상이 아니다. 그 위상은 $V\in\mathfrak B$, $V\subset U$인 사상
\[
  \Gamma(U,\mathcal F)\longrightarrow\Gamma(V,\mathcal F)
\]
들을 연속이 되게 하는 가장 엉성한 위상이다. 여기서 각 $\Gamma(V,\mathcal F)$는 이산공간이다.

앞의 논의는 반드시 가환일 필요는 없는 군 또는 환의 층에도 아무런 변경 없이 적용되어, 각각 \emph{위상군} 또는 \emph{위상환}의 층을 대응시킨다. 줄여서 $\mathcal F'$을 집합(각각 군, 환)의 층 $\mathcal F$에 대응하는 \emph{의사이산 공간}(각각 \emph{의사이산 군}, \emph{의사이산 환})의 층이라 하겠다.
\end{env}

\begin{env}[3.8.2]
\label{0.3.8.2-ko}
$\mathcal F$, $\mathcal G$를 $X$ 위의 집합(각각 군, 환)의 두 층이라 하고 $u:\mathcal F\to\mathcal G$를 사상이라 하자. $\mathcal F'$, $\mathcal G'$을 각각 $\mathcal F$, $\mathcal G$에 대응하는 의사이산 층이라 하면, $u$는 \emph{연속인} 사상
\[
  \mathcal F'\longrightarrow\mathcal G'
\]
이기도 하다. 이는 (3.2.5)에서 따른다.
\end{env}

\begin{env}[3.8.3]
\label{0.3.8.3-ko}
$\mathcal F$를 집합의 층, $\mathcal H$를 $\mathcal F$의 부분층이라 하고, $\mathcal F'$, $\mathcal H'$을 각각 $\mathcal F$, $\mathcal H$에 대응하는 의사이산 층이라 하자. 그러면 모든 열린집합 $U\subset X$에 대하여
\[
  \Gamma(U,\mathcal H')
\]
은 $\Gamma(U,\mathcal F')$에서 \emph{닫혀} 있다. 실제로 이는 $V\in\mathfrak B$, $V\subset U$인 연속사상
\[
  \Gamma(U,\mathcal F)\longrightarrow\Gamma(V,\mathcal F)
\]
에 의한 $\Gamma(V,\mathcal H)$들의 역상들의 교집합이며, $\Gamma(V,\mathcal H)$는 이산공간 $\Gamma(V,\mathcal F)$에서 닫혀 있다.
\end{env}

\section{환 달린 공간}
\label{section:0.4-ko}

\subsection{환 달린 공간, $\mathcal A$-가군, $\mathcal A$-대수}
\label{subsection:0.4.1-ko}

\begin{env}[4.1.1]
\label{0.4.1.1-ko}
\emph{환 달린 공간}(각각 \emph{위상환 달린 공간})은 위상공간 $X$와 $X$ 위의 반드시 가환일 필요는 없는 환의 층(각각 위상환의 층) $\mathcal A$로 이루어진 쌍 $(X,\mathcal A)$이다. $X$를 환 달린 공간 $(X,\mathcal A)$의 바탕
\oldpage[0\textsubscript{I}]{36}
위상공간이라 하고, $\mathcal A$를 \emph{구조층}이라 한다. 구조층은 $\mathcal O_X$로 나타내며, 한 점 $x\in X$에서의 그 줄기는 $\mathcal O_{X,x}$ 또는 혼동할 염려가 없을 때에는 간단히 $\mathcal O_x$로 나타낸다.

$X$ 위의 $\mathcal O_X$의 \emph{단위원 단면}, 곧 $\Gamma(X,\mathcal O_X)$의 단위원은 $1$ 또는 $e$로 나타낸다.

이 논고에서는 주로 \emph{가환}환의 층을 다룰 것이므로, 별도의 명시 없이 환 달린 공간 $(X,\mathcal A)$을 말할 때에는 $\mathcal A$가 가환환의 층임을 전제로 한다.

구조층이 반드시 가환일 필요가 없는 환 달린 공간들(각각 위상환 달린 공간들)은 다음과 같이 \emph{범주}를 이룬다. \emph{사상}
\[
  (X,\mathcal A)\longrightarrow(Y,\mathcal B)
\]
을 연속사상 $\psi:X\to Y$와 환의 층(각각 위상환의 층)의 $\psi$-\emph{사상} $\theta:\mathcal B\to\mathcal A$ (3.5.1)로 이루어진 쌍 $(\psi,\theta)=\Psi$로 정의한다. 두 번째 사상
\[
  \Psi'=(\psi',\theta'):(Y,\mathcal B)\longrightarrow(Z,\mathcal C)
\]
과 $\Psi$의 \emph{합성}은 $\Psi''=\Psi'\circ\Psi$로 나타내며, 이는 $\psi''=\psi'\circ\psi$이고 $\theta''$가 $\theta$와 $\theta'$의 합성인 사상 $(\psi'',\theta'')$이다. 이때 $\theta''=\psi'_*(\theta)\circ\theta'$이다(3.5.2 참조). 환 달린 공간에 대해서는
\[
  \theta''^\sharp
  =\theta^\sharp\circ\psi^*(\theta'^\sharp)
\]
임을 상기하자(3.5.5). 따라서 $\theta'^\sharp$과 $\theta^\sharp$이 \emph{단사}(각각 \emph{전사}) 준동형이면 $\theta''^\sharp$도 그러하다. 여기서는 모든 $x\in X$에 대하여 $\psi_x\circ\rho_{\psi(x)}$가 동형이라는 사실을 함께 사용한다(3.7.2). 앞의 사실들로부터 $\psi$가 \emph{단사} 연속사상이고 $\theta^\sharp$이 환의 층의 \emph{전사} 준동형이면 사상 $(\psi,\theta)$가 환 달린 공간의 범주에서 \emph{단사사상}(T, 1.1)임을 곧바로 확인할 수 있다.

혼동할 염려가 없을 때에는 표기에서 흔히 $\psi$를 $\Psi$로 바꾸어 쓰겠다. 예를 들어 $Y$의 부분집합 $U$에 대하여 $\psi^{-1}(U)$ 대신 $\Psi^{-1}(U)$라고 쓴다.
\end{env}

\begin{env}[4.1.2]
\label{0.4.1.2-ko}
$X$의 임의의 부분집합 $M$에 대하여 쌍 $(M,\mathcal A|M)$은 자명하게 환 달린 공간이다. 이를 환 달린 공간 $(X,\mathcal A)$이 $M$ 위에 \emph{유도하는} 환 달린 공간이라 하며, $(X,\mathcal A)$의 $M$으로의 \emph{제한}이라고도 한다. $j:M\to X$가 표준 단사사상이고 $\omega$가 $\mathcal A|M$의 항등사상이면,
\[
  (j,\omega^b):(M,\mathcal A|M)\longrightarrow(X,\mathcal A)
\]
은 환 달린 공간의 단사사상이며 이를 \emph{표준 단사사상}이라 한다. 사상 $\Psi:(X,\mathcal A)\to(Y,\mathcal B)$와 이 표준 단사사상의 합성을 $\Psi$의 $M$으로의 \emph{제한}이라 한다.
\end{env}

\begin{env}[4.1.3]
\label{0.4.1.3-ko}
환 달린 공간 $(X,\mathcal A)$ 위의 $\mathcal A$-가군 또는 대수적 층의 정의(G, II, 2.2)는 다시 설명하지 않겠다. $\mathcal A$가 반드시 가환일 필요가 없는 환의 층이면, 반대의 것을 명시하지 않는 한 $\mathcal A$-가군은 항상 ``왼쪽 $\mathcal A$-가군''을 뜻한다. $\mathcal A$의 부분-$\mathcal A$-가군은 $\mathcal A$ 안의 (왼쪽, 오른쪽 또는 양쪽) \emph{아이디얼층} 또는 $\mathcal A$-아이디얼이라 한다.

$\mathcal A$가 가환환의 층일 때 $\mathcal A$-가군의 정의에서 가군 구조를 모두 대수 구조로 바꾸면 $X$ 위의 $\mathcal A$-대수의 정의를 얻는다. 이는 반드시 가환일 필요는 없는 $\mathcal A$-대수가 $\mathcal A$-가군 $\mathcal C$에 $\mathcal A$-가군 준동형
\[
  \varphi:
  \mathcal C\otimes_{\mathcal A}\mathcal C
  \longrightarrow\mathcal C
\]
과 $X$ 위의 단면 $e$를 주어 다음 조건들을 만족시키는 것과 같다. 1\textsuperscript{o} 도표
\oldpage[0\textsubscript{I}]{37}
\[
\begin{tikzcd}[column sep=large,row sep=large]
\mathcal{C}\otimes_{\mathcal{A}}\mathcal{C}\otimes_{\mathcal{A}}\mathcal{C}
  \arrow[r,"\varphi\otimes 1"]
  \arrow[d,"1\otimes\varphi"']
&
\mathcal{C}\otimes_{\mathcal{A}}\mathcal{C}
  \arrow[d,"\varphi"]
\\
\mathcal{C}\otimes_{\mathcal{A}}\mathcal{C}
  \arrow[r,"\varphi"']
&
\mathcal{C}
\end{tikzcd}
\]
가 가환이다. 2\textsuperscript{o} 모든 열린집합 $U\subset X$와 모든 단면 $s\in\Gamma(U,\mathcal C)$에 대하여
\[
  \varphi((e|U)\otimes s)
  =\varphi(s\otimes(e|U))
  =s
\]
이다. $\mathcal C$가 가환 $\mathcal A$-대수라는 것은 이에 더하여 도표
\[
\begin{tikzcd}[column sep=large,row sep=large]
\mathcal{C}\otimes_{\mathcal{A}}\mathcal{C}
  \arrow[rr,"\sigma"]
  \arrow[dr,"\varphi"']
&&
\mathcal{C}\otimes_{\mathcal{A}}\mathcal{C}
  \arrow[dl,"\varphi"]
\\
& \mathcal{C} &
\end{tikzcd}
\]
가 가환이라는 뜻이다. 여기서 $\sigma$는 텐서곱 $\mathcal C\otimes_{\mathcal A}\mathcal C$의 표준 대칭이다.

$\mathcal A$-대수의 준동형도 (G, II, 2.2)의 $\mathcal A$-가군 준동형과 같은 방식으로 정의하지만, 이제 그 준동형들은 당연히 아벨군을 이루지 않는다.

$\mathcal M$이 $\mathcal A$-대수 $\mathcal C$의 부분-$\mathcal A$-가군이면, $\mathcal M$이 생성하는 $\mathcal C$의 부분-$\mathcal A$-대수는 $n\geqslant0$에 대한 준동형
\[
  \bigotimes^n\mathcal M\longrightarrow\mathcal C
\]
들의 상들의 합이다. 이는 대수의 준층 $U\mapsto\mathcal B(U)$의 층화이기도 하다. 여기서 $\mathcal B(U)$는 부분가군 $\Gamma(U,\mathcal M)$이 생성하는 $\Gamma(U,\mathcal C)$의 부분대수이다.
\end{env}

\begin{env}[4.1.4]
\label{0.4.1.4-ko}
위상공간 $X$ 위의 환의 층 $\mathcal A$에 대하여 줄기 $\mathcal A_x$가 축소환(1.1.1)이면 $\mathcal A$가 $X$의 점 $x$에서 \emph{축소}라고 한다. $\mathcal A$가 $X$의 모든 점에서 축소이면 $\mathcal A$를 \emph{축소}라고 한다. 환 $A$의 모든 소 아이디얼 $\mathfrak p$에 대하여 국소화환 $A_{\mathfrak p}$가 정칙 국소환이면 $A$를 \emph{정칙환}이라 한다는 것을 상기하자. $\mathcal A_x$가 정칙환이면 $\mathcal A$가 점 $x$에서 \emph{정칙}이라고 하고, 모든 점에서 정칙이면 $\mathcal A$를 \emph{정칙}이라고 한다. 끝으로 $\mathcal A_x$가 정수적으로 닫힌 정역이면 $\mathcal A$가 점 $x$에서 \emph{정규}라고 하고, 모든 점에서 정규이면 $\mathcal A$를 \emph{정규}라고 한다. 환 달린 공간 $(X,\mathcal A)$이 앞의 성질 가운데 하나를 갖는다는 것은 환의 층 $\mathcal A$가 그 성질을 갖는다는 뜻이다.

\emph{등급환의 층} $\mathcal A$는 정의에 따라 아벨군의 층들의 족 $(\mathcal A_n)_{n\in\mathbf Z}$의 직합(G, II, 2.7)인 환의 층으로서
\[
  \mathcal A_m\mathcal A_n
  \subset\mathcal A_{m+n}
\]
을 만족하는 것이다. \emph{등급 $\mathcal A$-가군}은 아벨군의 층들의 족 $(\mathcal F_n)_{n\in\mathbf Z}$의 직합인 $\mathcal A$-가군 $\mathcal F$로서
\[
  \mathcal A_m\mathcal F_n
  \subset\mathcal F_{m+n}
\]
을 만족하는 것이다. 이는 모든 점 $x$에서
\[
  (\mathcal A_m)_x(\mathcal A_n)_x
  \subset(\mathcal A_{m+n})_x
\]
(각각
\[
  (\mathcal A_m)_x(\mathcal F_n)_x
  \subset(\mathcal F_{m+n})_x
\]
)가 성립한다는 것과 같다.
\end{env}

\begin{env}[4.1.5]
\label{0.4.1.5-ko}
반드시 가환일 필요가 없는 환 달린 공간 $(X,\mathcal A)$을 생각하자. 경우에 따라 왼쪽 또는 오른쪽 $\mathcal A$-가군의 범주에서 정의되고 아벨군의 층의 범주(또는 더
\oldpage[0\textsubscript{I}]{38}
일반적으로 $\mathcal C$가 $\mathcal A$의 중심이면 $\mathcal C$-가군의 범주)에 값을 갖는 쌍함자
\[
  \mathcal F\otimes_{\mathcal A}\mathcal G,
  \qquad
  \mathcal H\!om_{\mathcal A}(\mathcal F,\mathcal G),
  \qquad
  \operatorname{Hom}_{\mathcal A}(\mathcal F,\mathcal G)
\]
의 정의(G, II, 2.8 및 2.2)는 여기서 되풀이하지 않겠다. 모든 점 $x\in X$에서 줄기
\[
  (\mathcal F\otimes_{\mathcal A}\mathcal G)_x
\]
는 표준적으로 $\mathcal F_x\otimes_{\mathcal A_x}\mathcal G_x$와 동일시되며, 함자적인 표준 준동형
\[
  (\mathcal H\!om_{\mathcal A}(\mathcal F,\mathcal G))_x
  \longrightarrow
  \operatorname{Hom}_{\mathcal A_x}(\mathcal F_x,\mathcal G_x)
\]
이 정의된다. 이 준동형은 일반적으로 단사도 전사도 아니다.

앞의 쌍함자들은 가법적이며, 특히 유한 직합과 가환한다. $\mathcal F\otimes_{\mathcal A}\mathcal G$는 $\mathcal F$와 $\mathcal G$ 각각에 관하여 오른쪽 완전하고 귀납극한과 가환하며,
\[
  \mathcal A\otimes_{\mathcal A}\mathcal G
  \simeq\mathcal G,
  \qquad
  \mathcal F\otimes_{\mathcal A}\mathcal A
  \simeq\mathcal F
\]
인 표준 동형이 있다. 함자 $\mathcal H\!om_{\mathcal A}(\mathcal F,\mathcal G)$와 $\operatorname{Hom}_{\mathcal A}(\mathcal F,\mathcal G)$은 $\mathcal F$와 $\mathcal G$에 관하여 \emph{왼쪽 완전}하다. 더 정확히 말하면 완전열
\[
  0\longrightarrow\mathcal G'
  \longrightarrow\mathcal G
  \longrightarrow\mathcal G''
\]
이 주어질 때 열
\[
  0\longrightarrow
  \mathcal H\!om_{\mathcal A}(\mathcal F,\mathcal G')
  \longrightarrow
  \mathcal H\!om_{\mathcal A}(\mathcal F,\mathcal G)
  \longrightarrow
  \mathcal H\!om_{\mathcal A}(\mathcal F,\mathcal G'')
\]
은 완전하고, 완전열
\[
  \mathcal F'\longrightarrow\mathcal F
  \longrightarrow\mathcal F''\longrightarrow0
\]
이 주어질 때 열
\[
  0\longrightarrow
  \mathcal H\!om_{\mathcal A}(\mathcal F'',\mathcal G)
  \longrightarrow
  \mathcal H\!om_{\mathcal A}(\mathcal F,\mathcal G)
  \longrightarrow
  \mathcal H\!om_{\mathcal A}(\mathcal F',\mathcal G)
\]
은 완전하다. $\operatorname{Hom}$ 함자에도 같은 성질들이 성립한다. 또한 $\mathcal H\!om_{\mathcal A}(\mathcal A,\mathcal G)$는 표준적으로 $\mathcal G$와 동일시되며, 모든 열린집합 $U\subset X$에 대하여
\[
  \Gamma\bigl(U,\mathcal H\!om_{\mathcal A}(\mathcal F,\mathcal G)\bigr)
  =
  \operatorname{Hom}_{\mathcal A|U}(\mathcal F|U,\mathcal G|U)
\]
이다.

모든 왼쪽(각각 오른쪽) $\mathcal A$-가군 $\mathcal F$에 대하여 오른쪽(각각 왼쪽) $\mathcal A$-가군
\[
  \check{\mathcal F}
  =\mathcal H\!om_{\mathcal A}(\mathcal F,\mathcal A)
\]
를 $\mathcal F$의 \emph{쌍대}라 한다.

끝으로 $\mathcal A$가 가환환의 층이고 $\mathcal F$가 $\mathcal A$-가군이면
\[
  U\longmapsto\bigwedge^p\Gamma(U,\mathcal F)
\]
는 준층을 이룬다. 이 준층의 층화는 $\bigwedge^p\mathcal F$로 나타내는 $\mathcal A$-가군이며, 이를 $\mathcal F$의 \emph{$p$차 외승}이라 한다. 준층 $U\mapsto\bigwedge^p\Gamma(U,\mathcal F)$에서 그 층화 $\bigwedge^p\mathcal F$로 가는 표준 사상은 \emph{단사}이고, 모든 $x\in X$에 대하여
\[
  (\bigwedge^p\mathcal F)_x
  =\bigwedge^p(\mathcal F_x)
\]
이다. $\bigwedge^p\mathcal F$가 $\mathcal F$에 관한 공변함자임은 자명하다.
\end{env}

\begin{env}[4.1.6]
\label{0.4.1.6-ko}
$\mathcal A$를 반드시 가환일 필요는 없는 환의 층, $\mathcal J$를 $\mathcal A$의 왼쪽 아이디얼층, $\mathcal F$를 왼쪽 $\mathcal A$-가군이라 하자. 이때 $\mathcal J\mathcal F$는 표준 사상
\[
  \mathcal J\otimes_{\mathbf Z}\mathcal F
  \longrightarrow\mathcal F
\]
의 상인 $\mathcal F$의 부분-$\mathcal A$-가군을 나타낸다. 여기서 $\mathbf Z$는 상수 준층 $U\mapsto\mathbf Z$의 층화이다. 모든 $x\in X$에 대하여
\[
  (\mathcal J\mathcal F)_x
  =\mathcal J_x\mathcal F_x
\]
임은 자명하다. $\mathcal A$가 가환이면 $\mathcal J\mathcal F$는 표준 사상
\[
  \mathcal J\otimes_{\mathcal A}\mathcal F
  \longrightarrow\mathcal F
\]
의 상이기도 하다. 또한 $\mathcal J\mathcal F$는 준층
\[
  U\longmapsto
  \Gamma(U,\mathcal J)\Gamma(U,\mathcal F)
\]
의 층화이다. $\mathcal J_1$, $\mathcal J_2$가 $\mathcal A$의 두 왼쪽 아이디얼층이면
\[
  \mathcal J_1(\mathcal J_2\mathcal F)
  =(\mathcal J_1\mathcal J_2)\mathcal F
\]
이다.
\end{env}

\begin{env}[4.1.7]
\label{0.4.1.7-ko}
$(X_\lambda,\mathcal A_\lambda)_{\lambda\in L}$를 환 달린 공간들의 족이라 하자. 각 $(\lambda,\mu)$에 대하여 $X_\lambda$의 열린부분집합 $V_{\lambda\mu}$와 환 달린 공간의 동형사상
\[
  \varphi_{\lambda\mu}:
  (V_{\mu\lambda},\mathcal A_\mu|V_{\mu\lambda})
  \xrightarrow{\sim}
  (V_{\lambda\mu},\mathcal A_\lambda|V_{\lambda\mu})
\]
이 주어졌다고 하자. 여기서 $V_{\lambda\lambda}=X_\lambda$이고 $\varphi_{\lambda\lambda}$는 항등사상이다. 더 나아가 모든 세쌍 $(\lambda,\mu,\nu)$에 대하여, $\varphi'_{\mu\lambda}$가 $\varphi_{\mu\lambda}$의 $V_{\lambda\mu}\cap V_{\lambda\nu}$에 대한 제한을 뜻할 때, $\varphi'_{\mu\lambda}$가
\[
  (V_{\lambda\mu}\cap V_{\lambda\nu},
   \mathcal A_\lambda|(V_{\lambda\mu}\cap V_{\lambda\nu}))
\]
에서
\[
  (V_{\mu\nu}\cap V_{\mu\lambda},
   \mathcal A_\mu|(V_{\mu\nu}\cap V_{\mu\lambda}))
\]
로 가는 동형사상이고
$\varphi'_{\lambda\nu}=\varphi'_{\lambda\mu}\circ\varphi'_{\mu\nu}$
라고 하자. 이를 $\varphi_{\lambda\mu}$들에 관한 \emph{붙이기 조건}이라 한다. 그러면 먼저 $X_\lambda$들을

\oldpage[0\textsubscript{I}]{39}
$V_{\lambda\mu}$들을 따라 $\varphi_{\lambda\mu}$들로 붙여서 얻는 위상공간 $X$를 생각할 수 있다. $X_\lambda$를 $X$ 안의 대응하는 열린부분집합 $X'_\lambda$와 동일시하면, 가정에 따라 세 집합
$V_{\lambda\mu}\cap V_{\lambda\nu}$,
$V_{\mu\nu}\cap V_{\mu\lambda}$,
$V_{\nu\lambda}\cap V_{\nu\mu}$는 모두
$X'_\lambda\cap X'_\mu\cap X'_\nu$와 동일시된다. 이제 $X_\lambda$의 환 달린 공간 구조를 $X'_\lambda$로 옮길 수 있다. $\mathcal A'_\lambda$가 $\mathcal A_\lambda$로부터 옮겨진 환의 층이면, $\mathcal A'_\lambda$들은 붙이기 조건 (3.3.1)을 만족하므로 $X$ 위의 환의 층 $\mathcal A$를 정의한다. 이때 $(X,\mathcal A)$를 $\varphi_{\lambda\mu}$들을 사용하여 $V_{\lambda\mu}$들을 따라 $(X_\lambda,\mathcal A_\lambda)$들을 \emph{붙여서 얻은 환 달린 공간}이라 한다.
\end{env}

% Authority: EGA I chapter1-frontmatter-fr.tex, SHA-256 7B2D0F8F812EBA3121202F0AE6415FFC6C281B8428DA8F0F72D89DF1CEC01708.
\clearpage
\oldpage[I]{79}
\phantomsection
\label{chapter:I-ko}

\begin{center}
{\large 제1장\par}
\vspace{1.5em}
{\Large\bfseries 스킴의 언어\par}
\vspace{1.8em}
{\bfseries 목차\par}
\end{center}

\medskip
\begin{tabular}{@{}r@{\ }l@{}}
\S~1.  & 아핀 스킴.\\
\S~2.  & 준스킴과 준스킴의 사상.\\
\S~3.  & 준스킴의 곱.\\
\S~4.  & 부분준스킴과 몰입 사상.\\
\S~5.  & 축소 준스킴; 분리 조건.\\
\S~6.  & 유한성 조건.\\
\S~7.  & 유리 사상.\\
\S~8.  & 슈발레 스킴.\\
\S~9.  & 준연접층에 관한 보충.\\
\S~10. & 형식 스킴.
\end{tabular}

\medskip
\S\S~1에서~8까지는 이후 전체에서 사용할 언어를 전개하는 것에서
크게 벗어나지 않는다. 다만 이 논고의 전반적인 정신에 따라
\S\S~7과~8은 다른 절들보다 덜 쓰이고, 쓰이더라도 덜 본질적이라는
점을 밝혀 둔다. 슈발레 스킴을 다룬 것도 슈발레~\cite{I-1}와
나가타~\cite{I-9}의 언어와 연결하기 위해서일 뿐이다. \S~9에서는
준연접층에 관한 정의와 결과들을 제시한다. 그 가운데에는 알려진
가환대수의 개념들을 ``기하학적'' 언어로 옮기는 데 그치지 않고 이미
전역적인 성격을 띠는 것도 있으며, 이러한 결과들은 바로 다음 장들부터
사상을 전역적으로 연구하는 데 반드시 필요하다. 끝으로 \S~10에서는
스킴 개념의 한 일반화를 도입한다. 이는 제III장에서 고유 사상의 코호몰로지
연구에 관한 기본 결과들을 편리하게 서술하고 증명하는 중간 도구로 쓰일
것이다. 한편 형식 스킴의 개념은 ``모듈라이 이론''의 몇몇 사실, 곧
대수다양체의 분류 문제를 표현하는 데 없어서는 안 될 것으로 보인다.
\S~10의 결과들은 제III장 \S~3 이전에는 사용하지 않으므로, 그때까지는
이 절을 건너뛰어 읽기를 권한다.

\clearpage

% Authority: EGA I ega1-1-fr.tex, whole-file SHA-256 25840D1D77C5284A158CDC78D622A458221DA3FE9ECAECC5EA9382A5C4EC595B.
% Sequential coverage in this file: authority lines 1-1624, complete subsections 1.1-1.7.
\setcounter{section}{0}
\section{아핀 스킴}
\label{section:I.1-ko}

\setcounter{subsection}{0}
\subsection{환의 소 스펙트럼}
\label{subsection:I.1.1-ko}

\oldpage[I]{80}
\begin{env}[1.1.1]
\label{I.1.1.1-ko}
\emph{표기}: $A$를 (가환)환, $M$을 $A$-가군이라 하자. 이 장과
이어지는 장들에서는 다음 표기들을 계속 사용한다.

\smallskip
\noindent
$\operatorname{Spec}(A)={}$ $A$의 \emph{소 아이디얼들의 집합}이며,
$A$의 \emph{소 스펙트럼}이라고도 한다. 흔히
$x\in X=\operatorname{Spec}(A)$ 대신 $\mathfrak{j}_x$라고 쓰는 것이
편리하다. $\operatorname{Spec}(A)$가 \emph{공집합}일 필요충분조건은
$A$가 영환인 것이다.

\smallskip
\noindent
$A_x=A_{\mathfrak{j}_x}={}$ \emph{$S^{-1}A$라는 (국소) 분수환},
여기서 $S=A-\mathfrak{j}_x$이다.

\smallskip
\noindent
$\mathfrak{m}_x=\mathfrak{j}_xA_{\mathfrak{j}_x}={}$ $A_x$의
\emph{극대 아이디얼}.

\smallskip
\noindent
$k(x)=A_x/\mathfrak{m}_x={}$ $A_x$의 \emph{잉여체}. 이는 정역
$A/\mathfrak{j}_x$의 분수체와 표준적으로 동형이며, 그 분수체와
동일시한다.

\smallskip
\noindent
$f(x)={}$ $A/\mathfrak{j}_x\subset k(x)$ 안에서 취한 \emph{$f$의
$\mathfrak{j}_x$ 법 동치류} ($f\in A$, $x\in X$). $f(x)$를
$x\in\operatorname{Spec}(A)$에서의 $f$의 \emph{값}이라고도 한다.
관계 $f(x)=0$과 $f\in\mathfrak{j}_x$는 \emph{동치}이다.

\smallskip
\noindent
$M_x=M\otimes_AA_x={}$ $A-\mathfrak{j}_x$에 분모를 갖는
\emph{분수가군}.

\smallskip
\noindent
$\mathfrak{r}(E)={}$ $A$의 부분집합 $E$가 생성하는 아이디얼의
\emph{근기}.

\smallskip
\noindent
$V(E)={}$ \emph{$E\subset\mathfrak{j}_x$를 만족하는 $x\in X$들의 집합}
(여기서 $E\subset A$)
(또는 같은 말로, 모든 $f\in E$에 대하여 $f(x)=0$인 $x\in X$들의
집합). 따라서
\begin{equation}
  \mathfrak{r}(E)=\bigcap_{x\in V(E)}\mathfrak{j}_x.
  \tag{1.1.1.1}\label{I.1.1.1.1-ko}
\end{equation}

\smallskip
\noindent
$V(f)=V(\{f\})$ ($f\in A$).

\smallskip
\noindent
$D(f)=X-V(f)={}$ \emph{$f(x)\neq0$인 $x\in X$들의 집합}.
\end{env}

\begin{proposition}[1.1.2]
\label{I.1.1.2-ko}
다음 성질들이 성립한다.
\begin{enumerate}
  \item[\rm(i)] $V(0)=X$, $V(1)=\varnothing$.
  \item[\rm(ii)] $E\subset E'$이면 $V(E)\supset V(E')$이다.
  \item[\rm(iii)] $A$의 부분집합들의 임의의 족 $(E_\lambda)$에 대하여
  \[
    V\left(\bigcup_\lambda E_\lambda\right)
      =V\left(\sum_\lambda E_\lambda\right)
      =\bigcap_\lambda V(E_\lambda).
  \]
  \item[\rm(iv)] $V(EE')=V(E)\cup V(E')$.
  \item[\rm(v)] $V(E)=V(\mathfrak{r}(E))$.
\end{enumerate}
\end{proposition}

성질 (i), (ii), (iii)은 자명하고, (v)는 (ii)와 공식
(\hyperref[I.1.1.1.1-ko]{1.1.1.1})에서 따른다.
$V(EE')\supset V(E)\cup V(E')$임은 자명하다. 역으로
$x\notin V(E)$이고 $x\notin V(E')$이면 $k(x)$ 안에서
$f(x)\neq0$, $f'(x)\neq0$을 만족하는 $f\in E$, $f'\in E'$가
존재한다. 따라서 $f(x)f'(x)\neq0$, 즉 $x\notin V(EE')$이고,
이로써 (iv)가 증명된다.

명제 (\hyperref[I.1.1.2-ko]{1.1.2})에서 특히 $E$가 $A$의 모든
부분집합을 달릴 때 $V(E)$ 꼴의 집합들이 $X$ 위의 한 위상의
\emph{닫힌집합}들이 됨을 알 수 있다. 이 위상을
\emph{스펙트럼 위상}\footnote{대수기하학에 이 위상을 도입한 사람은
자리스키이다. 따라서 흔히 $X$의 ``자리스키 위상''이라고도 한다.}이라
하겠다. 명시적으로 달리 말하지 않는 한
$X=\operatorname{Spec}(A)$에는 언제나 스펙트럼 위상을 준다.

\oldpage[I]{81}
\begin{env}[1.1.3]
\label{I.1.1.3-ko}
$X$의 임의의 부분집합 $Y$에 대하여 $\mathfrak{j}(Y)$는 모든
$y\in Y$에서 $f(y)=0$인 $f\in A$들의 집합을 뜻한다. 같은 말로
$\mathfrak{j}(Y)$는 $y\in Y$에 대한 소 아이디얼
$\mathfrak{j}_y$들의 교집합이다. $Y\subset Y'$이면
$\mathfrak{j}(Y)\supset\mathfrak{j}(Y')$임은 자명하며, $X$의
부분집합들의 임의의 족 $(Y_\lambda)$에 대하여
\begin{equation}
  \mathfrak{j}\left(\bigcup_\lambda Y_\lambda\right)
    =\bigcap_\lambda\mathfrak{j}(Y_\lambda)
  \tag{1.1.3.1}\label{I.1.1.3.1-ko}
\end{equation}
이다. 끝으로
\begin{equation}
  \mathfrak{j}(\{x\})=\mathfrak{j}_x
  \tag{1.1.3.2}\label{I.1.1.3.2-ko}
\end{equation}
이다.
\end{env}

\begin{proposition}[1.1.4]
\label{I.1.1.4-ko}
\begin{enumerate}
  \item[\rm(i)] $A$의 모든 부분집합 $E$에 대하여
  $\mathfrak{j}(V(E))=\mathfrak{r}(E)$이다.
  \item[\rm(ii)] $X$의 모든 부분집합 $Y$에 대하여
  $V(\mathfrak{j}(Y))=\overline{Y}$이며, 여기서
  $\overline{Y}$는 $X$에서의 $Y$의 폐포이다.
\end{enumerate}
\end{proposition}

(i)는 정의와 (\hyperref[I.1.1.1.1-ko]{1.1.1.1})에서 바로 따른다.
한편 $V(\mathfrak{j}(Y))$는 닫혀 있고 $Y$를 포함한다. 역으로
$Y\subset V(E)$이면 모든 $f\in E$, $y\in Y$에 대하여 $f(y)=0$이므로
$E\subset\mathfrak{j}(Y)$이고
$V(E)\supset V(\mathfrak{j}(Y))$이다. 이로써 (ii)가 증명된다.

\begin{corollary}[1.1.5]
\label{I.1.1.5-ko}
$X=\operatorname{Spec}(A)$의 닫힌 부분집합들과 자기 근기와 같은
$A$의 아이디얼들(다시 말해 소 아이디얼들의 교집합들)은 포함관계를
뒤집는 사상
$Y\mapsto\mathfrak{j}(Y)$, $\mathfrak{a}\mapsto V(\mathfrak{a})$에
의하여 서로 전단사 대응한다. 두 닫힌 부분집합의 합집합
$Y_1\cup Y_2$에는 $\mathfrak{j}(Y_1)\cap\mathfrak{j}(Y_2)$가
대응하고, 닫힌 부분집합들의 임의의 족 $(Y_\lambda)$의 교집합에는
$\mathfrak{j}(Y_\lambda)$들의 합의 근기가 대응한다.
\end{corollary}

\begin{corollary}[1.1.6]
\label{I.1.1.6-ko}
$A$가 뇌터 환이면 $X=\operatorname{Spec}(A)$는 뇌터 공간이다.
\end{corollary}

이 따름정리의 역은 거짓이다. 예를 들어 영 아이디얼이 아닌 소
아이디얼을 단 하나만 갖는 비뇌터 정역, 이를테면 랭크 1인 비이산
값매김환이 반례가 된다.

스펙트럼이 뇌터 공간이 아닌 환 $A$의 예로는 무한 콤팩트 공간
$Y$ 위의 실수값 연속함수들의 환 $\mathcal{C}(Y)$를 들 수 있다.
집합으로서 $Y$는 $A$의 극대 아이디얼들의 집합과 동일시된다는 것이
알려져 있고, $X=\operatorname{Spec}(A)$의 위상이 $Y$ 위에 유도하는
위상은 $Y$의 본래 위상임을 쉽게 알 수 있다. $Y$가 뇌터 공간이
아니므로 $X$도 뇌터 공간이 아니다.

\begin{corollary}[1.1.7]
\label{I.1.1.7-ko}
모든 $x\in X$에 대하여 $\{x\}$의 폐포는
$\mathfrak{j}_x\subset\mathfrak{j}_y$를 만족하는 $y\in X$들의
집합이다. $\{x\}$가 닫혀 있을 필요충분조건은
$\mathfrak{j}_x$가 극대인 것이다.
\end{corollary}

\begin{corollary}[1.1.8]
\label{I.1.1.8-ko}
공간 $X=\operatorname{Spec}(A)$는 콜모고로프 공간이다.
\end{corollary}

실제로 $x$, $y$가 $X$의 서로 다른 두 점이면
$\mathfrak{j}_x\not\subset\mathfrak{j}_y$이거나
$\mathfrak{j}_y\not\subset\mathfrak{j}_x$이다. 따라서 두 점
$x$, $y$ 가운데 하나는 다른 하나의 폐포에 속하지 않는다.

\begin{env}[1.1.9]
\label{I.1.1.9-ko}
명제 (\hyperref[I.1.1.2-ko]{1.1.2}, (iv))에 따라 $A$의 두 원소
$f$, $g$에 대하여
\begin{equation}
  D(fg)=D(f)\cap D(g)
  \tag{1.1.9.1}\label{I.1.1.9.1-ko}
\end{equation}
이다. 또한 명제 (\hyperref[I.1.1.4-ko]{1.1.4}, (i))와
명제 (\hyperref[I.1.1.2-ko]{1.1.2}, (v))에 따르면
$D(f)=D(g)$라는 관계는
$\mathfrak{r}(f)=\mathfrak{r}(g)$라는 뜻이며, 같은 말로
$(f)$와 $(g)$를 포함하는 극소 소 아이디얼들이 같다는 뜻이다.
특히 $u$가 가역원이고 $f=ug$이면 그러하다.
\end{env}

\oldpage[I]{82}
\begin{proposition}[1.1.10]
\label{I.1.1.10-ko}
\begin{enumerate}
  \item[\rm(i)] $f$가 $A$를 달릴 때 집합 $D(f)$들은 $X$의 위상의
  한 기저를 이룬다.
  \item[\rm(ii)] 모든 $f\in A$에 대하여 $D(f)$는 준콤팩트이다.
  특히 $X=D(1)$은 준콤팩트이다.
\end{enumerate}
\end{proposition}

(i) $U$를 $X$의 열린집합이라 하자. 정의에 따라 $A$의 한 부분집합
$E$에 대하여 $U=X-V(E)$이고,
$V(E)=\bigcap_{f\in E}V(f)$이다. 따라서
$U=\bigcup_{f\in E}D(f)$이다.

(ii) (i)에 따라 다음을 보이면 충분하다. $A$의 원소들의 족
$(f_\lambda)_{\lambda\in L}$이
$D(f)\subset\bigcup_{\lambda\in L}D(f_\lambda)$를 만족하면,
$D(f)\subset\bigcup_{\lambda\in J}D(f_\lambda)$가 되는 $L$의
유한 부분집합 $J$가 존재한다. $\mathfrak{a}$를 $f_\lambda$들이
생성하는 $A$의 아이디얼이라 하자. 가정에 따라
$V(f)\supset V(\mathfrak{a})$이므로
$\mathfrak{r}(f)\subset\mathfrak{r}(\mathfrak{a})$이다.
$f\in\mathfrak{r}(f)$이므로 어떤 정수 $n\geq0$에 대하여
$f^n\in\mathfrak{a}$이다. 그러면 $f^n$은 한 유한 부분족
$(f_\lambda)_{\lambda\in J}$이 생성하는 아이디얼
$\mathfrak{b}$에 속하며,
$V(f)=V(f^n)\supset V(\mathfrak{b})
=\bigcap_{\lambda\in J}V(f_\lambda)$이다. 다시 말해
$D(f)\subset\bigcup_{\lambda\in J}D(f_\lambda)$이다.

\begin{proposition}[1.1.11]
\label{I.1.1.11-ko}
$A$의 모든 아이디얼 $\mathfrak{a}$에 대하여
$\operatorname{Spec}(A/\mathfrak{a})$는
$\operatorname{Spec}(A)$의 닫힌 부분공간 $V(\mathfrak{a})$와
표준적으로 동일시된다.
\end{proposition}

실제로 $A/\mathfrak{a}$의 아이디얼들(각각 소 아이디얼들)과
$\mathfrak{a}$를 포함하는 $A$의 아이디얼들(각각 소 아이디얼들)
사이에는 포함관계의 순서구조를 보존하는 표준 전단사 대응이 있다.

$A$의 멱영원들의 집합 $\mathfrak{N}$, 곧 $A$의
\emph{멱영근기}는 $\mathfrak{r}(0)$과 같은 아이디얼이고
$A$의 모든 소 아이디얼의 교집합임을 상기하자
(\hyperref[0.1.1.1-ko]{0, 1.1.1}).

\begin{corollary}[1.1.12]
\label{I.1.1.12-ko}
위상공간 $\operatorname{Spec}(A)$와
$\operatorname{Spec}(A/\mathfrak{N})$은 표준적으로 동형이다.
\end{corollary}

\begin{proposition}[1.1.13]
\label{I.1.1.13-ko}
$X=\operatorname{Spec}(A)$가 기약일
(\hyperref[0.2.1.1-ko]{0, 2.1.1}) 필요충분조건은
$A/\mathfrak{N}$이 정역인 것이다. 이는 $\mathfrak{N}$이
소 아이디얼인 것과도 동치이다.
\end{proposition}

따름정리 (\hyperref[I.1.1.12-ko]{1.1.12})에 따라
$\mathfrak{N}=0$인 경우만 생각하면 된다. $X$가 기약이 아니면
$X$와 다른 두 닫힌 부분집합 $Y_1$, $Y_2$가 존재하여
$X=Y_1\cup Y_2$이고, 따라서
$\mathfrak{j}(X)=\mathfrak{j}(Y_1)\cap\mathfrak{j}(Y_2)=0$이다.
아이디얼 $\mathfrak{j}(Y_1)$과 $\mathfrak{j}(Y_2)$는 모두
$(0)$과 다르므로(\hyperref[I.1.1.5-ko]{1.1.5}) $A$는 정역이
아니다. 역으로 $A$에 $f\neq0$, $g\neq0$, $fg=0$인 두 원소가
있으면, $A$의 모든 소 아이디얼의 교집합이 $(0)$이므로
$V(f)\neq X$, $V(g)\neq X$이고
$X=V(fg)=V(f)\cup V(g)$이다.

\begin{corollary}[1.1.14]
\label{I.1.1.14-ko}
\begin{enumerate}
  \item[\rm(i)] $X=\operatorname{Spec}(A)$의 닫힌 부분집합들과 자기
  근기와 같은 $A$의 아이디얼들 사이의 전단사 대응에서, $X$의 닫힌
  기약 부분집합들은 $A$의 소 아이디얼들과 대응한다. 특히 $X$의
  기약 성분들은 $A$의 극소 소 아이디얼들과 대응한다.
  \item[\rm(ii)] 사상 $x\mapsto\overline{\{x\}}$는 $X$와 $X$의
  닫힌 기약 부분집합들의 집합 사이에 전단사 대응을 정한다.
  다시 말해 $X$의 모든 닫힌 기약 부분집합은 일반점을 하나만 갖는다.
\end{enumerate}
\end{corollary}

(i)는 (\hyperref[I.1.1.13-ko]{1.1.13})과
(\hyperref[I.1.1.11-ko]{1.1.11})에서 바로 따른다. (ii)를
증명할 때에도 (\hyperref[I.1.1.11-ko]{1.1.11})에 따라 $X$가
기약인 경우만 생각하면 된다. 그러면 명제
(\hyperref[I.1.1.13-ko]{1.1.13})에 따라 $A$에는 가장 작은 소
아이디얼 $\mathfrak{N}$이 존재하고, 이는 $X$의 한 일반점에
대응한다.
\oldpage[I]{83}
더욱이 $X$는 콜모고로프 공간이므로
((\hyperref[I.1.1.8-ko]{1.1.8})과
(\hyperref[0.2.1.3-ko]{0, 2.1.3})) 일반점을 하나만 갖는다.

\begin{proposition}[1.1.15]
\label{I.1.1.15-ko}
$\mathfrak{J}$가 $A$의 근기 $\mathfrak{R}(A)$에 포함되는
아이디얼이면, $X=\operatorname{Spec}(A)$에서
$V(\mathfrak{J})$의 유일한 근방은 $X$ 전체이다.
\end{proposition}

실제로 $A$의 모든 극대 아이디얼은 정의에 따라
$V(\mathfrak{J})$에 속한다. $A$의 모든 아이디얼
$\mathfrak{a}$는 어떤 극대 아이디얼에 포함되므로
$V(\mathfrak{a})\cap V(\mathfrak{J})\neq\varnothing$이고,
이로부터 명제가 따른다.

\subsection{환의 소 스펙트럼의 함자적 성질}
\label{subsection:I.1.2-ko}

\begin{env}[1.2.1]
\label{I.1.2.1-ko}
$A$, $A'$을 두 환이라 하고
\[
  \varphi:A'\longrightarrow A
\]
를 환 준동형이라 하자. 모든 소 아이디얼
$x=\mathfrak{j}_x\in\operatorname{Spec}(A)=X$에 대하여 환
$A'/\varphi^{-1}(\mathfrak{j}_x)$는
$A/\mathfrak{j}_x$의 한 부분환과 표준적으로 동형이므로 정역이다.
다시 말해 $\varphi^{-1}(\mathfrak{j}_x)$는 $A'$의 소 아이디얼이다.
이를 ${}^a\varphi(x)$로 나타내면 사상
\[
  {}^a\varphi:X=\operatorname{Spec}(A)\longrightarrow
  X'=\operatorname{Spec}(A')
\]
가 정의된다. 이 사상은 $\operatorname{Spec}(\varphi)$라고도 쓰며,
준동형 $\varphi$에 대응하는 사상이라고 한다. 몫으로 내려가서
$\varphi$로부터 얻는
$A'/\varphi^{-1}(\mathfrak{j}_x)$에서
$A/\mathfrak{j}_x$로 가는 단사 준동형과, 이를 체의 단사 준동형으로
표준 연장한 것을 모두 $\varphi^x$로 나타낸다. 따라서
\[
  \varphi^x:k({}^a\varphi(x))\longrightarrow k(x).
\]
정의에 따라 모든 $f'\in A'$에 대하여
\begin{equation}
  \varphi^x\bigl(f'({}^a\varphi(x))\bigr)=\bigl(\varphi(f')\bigr)(x)
  \qquad(x\in X).
  \tag{1.2.1.1}\label{I.1.2.1.1-ko}
\end{equation}
\end{env}

\begin{proposition}[1.2.2]
\label{I.1.2.2-ko}
\begin{enumerate}
  \item[\rm(i)] $A'$의 임의의 부분집합 $E'$에 대하여
  \begin{equation}
    {}^a\varphi^{-1}(V(E'))=V(\varphi(E'))
    \tag{1.2.2.1}\label{I.1.2.2.1-ko}
  \end{equation}
  이고, 특히 모든 $f'\in A'$에 대하여
  \begin{equation}
    {}^a\varphi^{-1}(D(f'))=D(\varphi(f')).
    \tag{1.2.2.2}\label{I.1.2.2.2-ko}
  \end{equation}
  \item[\rm(ii)] $A$의 모든 아이디얼 $\mathfrak{a}$에 대하여
  \begin{equation}
    \overline{{}^a\varphi(V(\mathfrak{a}))}
    =V(\varphi^{-1}(\mathfrak{a})).
    \tag{1.2.2.3}\label{I.1.2.2.3-ko}
  \end{equation}
\end{enumerate}
\end{proposition}

실제로 관계 ${}^a\varphi(x)\in V(E')$는 정의에 따라
$E'\subset\varphi^{-1}(\mathfrak{j}_x)$와 동치이고, 이는 다시
$\varphi(E')\subset\mathfrak{j}_x$, 끝으로
$x\in V(\varphi(E'))$와 동치이다. 이로써 (i)가 증명된다. (ii)를
증명할 때에는 $V(\mathfrak{r}(\mathfrak{a}))=V(\mathfrak{a})$
((\hyperref[I.1.1.2-ko]{1.1.2 (v)}))이고
\[
  \varphi^{-1}(\mathfrak{r}(\mathfrak{a}))
  =\mathfrak{r}(\varphi^{-1}(\mathfrak{a}))
\]
이므로 $\mathfrak{a}$가 자기 근기와 같다고 가정해도 된다.
$Y=V(\mathfrak{a})$와
$\mathfrak{a}'=\mathfrak{j}({}^a\varphi(Y))$로 놓으면 명제
(\hyperref[I.1.1.4-ko]{1.1.4 (ii)})에 따라
$\overline{{}^a\varphi(Y)}=V(\mathfrak{a}')$이다. 관계
$f'\in\mathfrak{a}'$은 정의에 따라 모든
$x'\in{}^a\varphi(Y)$에 대하여 $f'(x')=0$이라는 것과 동치이다.
공식 (\hyperref[I.1.2.1.1-ko]{1.2.1.1})에 따라 이는 모든
$x\in Y$에 대하여 $\varphi(f')(x)=0$이라는 것과도 동치이고,
$\mathfrak{a}$가 자기 근기와 같으므로 다시
$\varphi(f')\in\mathfrak{j}(Y)=\mathfrak{a}$와 동치이다. 이로써
(ii)가 증명된다.

\begin{corollary}[1.2.3]
\label{I.1.2.3-ko}
사상 ${}^a\varphi$는 연속이다.
\end{corollary}

$A''$이 세 번째 환이고 $\varphi':A''\to A'$이 환 준동형이면
\[
  {}^a(\varphi\circ\varphi')={}^a\varphi'\circ{}^a\varphi.
\]
이 결과와 따름정리 (\hyperref[I.1.2.3-ko]{1.2.3})은
$\operatorname{Spec}(A)$가 $A$에 관하여 환의 범주에서 위상공간의
범주로 가는 반변 함자임을 뜻한다.

\begin{corollary}[1.2.4]
\label{I.1.2.4-ko}
\oldpage[I]{84}
$\varphi$가 다음 성질을 갖는다고 하자. 모든 $f\in A$를
\[
  f=h\varphi(f')
\]
꼴로 쓸 수 있고, 여기서 $h$는 $A$의 가역원이다. 특히
$\varphi$가 전사이면 이 성질이 성립한다. 그러면 ${}^a\varphi$는
$X$에서 ${}^a\varphi(X)$로 가는 위상동형이다.
\end{corollary}

임의의 부분집합 $E\subset A$에 대하여
$V(E)=V(\varphi(E'))$가 되는 $A'$의 부분집합 $E'$이 존재함을
보이자. 공리 $(T_0)$
((\hyperref[I.1.1.8-ko]{1.1.8}))과 공식
(\hyperref[I.1.2.2.1-ko]{1.2.2.1})에 따라 이 사실은 먼저
${}^a\varphi$가 단사임을 보이고, 다시 같은 공식에 따라
${}^a\varphi$가 위상동형임을 보인다. 실제로 각 $f\in E$마다
$h\varphi(f')=f$이고 $h$가 $A$의 가역원이 되도록 $f'\in A'$을
하나 택하면 충분하다. 이렇게 얻은 원소 $f'$들의 집합 $E'$이 원하는
조건을 만족한다.

\begin{env}[1.2.5]
\label{I.1.2.5-ko}
특히 $\varphi$가 $A$에서 몫환 $A/\mathfrak{a}$로 가는 표준
준동형이면 (\hyperref[I.1.1.12-ko]{1.1.12})를 다시 얻는다. 이때
${}^a\varphi$는 $\operatorname{Spec}(A/\mathfrak{a})$와 동일시한
$V(\mathfrak{a})$에서 $X=\operatorname{Spec}(A)$로 가는 표준
단사 사상이다.
\end{env}

(\hyperref[I.1.2.4-ko]{1.2.4})의 또 다른 특수한 경우로 다음을
얻는다.

\begin{corollary}[1.2.6]
\label{I.1.2.6-ko}
$S$가 $A$의 곱셈적 부분집합이면 스펙트럼
$\operatorname{Spec}(S^{-1}A)$는, 위상까지 포함하여,
$X=\operatorname{Spec}(A)$에서
$\mathfrak{j}_x\cap S=\varnothing$을 만족하는 점 $x$들로 이루어진
부분공간과 표준적으로 동일시된다.
\end{corollary}

실제로 $S^{-1}A$의 소 아이디얼들이
$\mathfrak{j}_x\cap S=\varnothing$을 만족하는
$S^{-1}\mathfrak{j}_x$들이고
\[
  \mathfrak{j}_x=({}^S i_A)^{-1}(S^{-1}\mathfrak{j}_x)
\]
임은 (\hyperref[0.1.2.6-ko]{0, 1.2.6})에서 이미 알고 있다. 따라서
${}^S i_A$에 따름정리 (\hyperref[I.1.2.4-ko]{1.2.4})를 적용하면
된다.

\begin{corollary}[1.2.7]
\label{I.1.2.7-ko}
${}^a\varphi(X)$가 $X'$에서 모든 곳에서 조밀하기 위한 필요충분조건은
$\operatorname{Ker}\varphi$의 모든 원소가 멱영원인 것이다.
\end{corollary}

실제로 공식 (\hyperref[I.1.2.2.3-ko]{1.2.2.3})을 아이디얼
$\mathfrak{a}=(0)$에 적용하면
\[
  \overline{{}^a\varphi(X)}=V(\operatorname{Ker}\varphi)
\]
를 얻는다. $V(\operatorname{Ker}\varphi)=X'$이기 위한
필요충분조건은 $\operatorname{Ker}\varphi$가 $A'$의 모든 소
아이디얼, 곧 $A'$의 멱영근기 $\mathfrak{N}$에 포함되는 것이다.

\subsection{가군에 대응하는 층}
\label{subsection:I.1.3-ko}
\label{I.1.3-ko}

\begin{env}[1.3.1]
\label{I.1.3.1-ko}
$A$를 가환환, $M$을 $A$-가군, $f$를 $A$의 원소라 하고,
$S_f$를 $n\geq0$일 때의 $f^n$들로 이루어진 곱셈적 집합이라 하자.
$A_f=S_f^{-1}A$, $M_f=S_f^{-1}M$으로 놓는다는 것을 상기하자.
$S'_f$가 $S_f$의 한 원소를 나누는 $g\in A$들로 이루어진 $A$의
포화 곱셈적 부분집합이면, $A_f$와 $M_f$는 각각
${S'_f}^{-1}A$와 ${S'_f}^{-1}M$에 표준적으로 동일시된다
(\hyperref[0.1.4.3-ko]{0, 1.4.3}).
\end{env}

\begin{lemma}[1.3.2]
\label{I.1.3.2-ko}
다음 조건들은 서로 동치이다.
\begin{enumerate}
  \item[\rm(a)] $g\in S'_f$;
  \item[\rm(b)] $S'_g\subset S'_f$;
  \item[\rm(c)] $f\in\mathfrak{r}(g)$;
  \item[\rm(d)] $\mathfrak{r}(f)\subset\mathfrak{r}(g)$;
  \item[\rm(e)] $V(g)\subset V(f)$;
  \item[\rm(f)] $D(f)\subset D(g)$.
\end{enumerate}
\end{lemma}

이는 정의와 (\hyperref[I.1.1.5-ko]{1.1.5})에서 바로 따른다.

\begin{env}[1.3.3]
\label{I.1.3.3-ko}
$D(f)=D(g)$이면 보조정리 (\hyperref[I.1.3.2-ko]{1.3.2, b)})에
따라 $M_f=M_g$이다. 더 일반적으로 $D(f)\supset D(g)$, 따라서
$S'_f\subset S'_g$이면 함자적인 표준 준동형
\[
  \rho_{g,f}:M_f\longrightarrow M_g
\]
가 존재함을 알고 있다(\hyperref[0.1.4.1-ko]{0, 1.4.1}). 또한
$D(f)\supset D(g)\supset D(h)$이면
(\hyperref[0.1.4.4-ko]{0, 1.4.4})
\begin{equation}
  \rho_{h,g}\circ\rho_{g,f}=\rho_{h,f}.
  \tag{1.3.3.1}\label{I.1.3.3.1-ko}
\end{equation}
\end{env}

\oldpage[I]{85}
주어진 $x\in X=\operatorname{Spec}(A)$에 대하여 $f$가
$A-\mathfrak{j}_x$를 움직일 때, $S'_f$들은
$A-\mathfrak{j}_x$의 부분집합들로 이루어진 증가 여과족을 이룬다.
실제로 $A-\mathfrak{j}_x$의 두 원소 $f$, $g$에 대하여 $S'_f$와
$S'_g$는 $S'_{fg}$에 포함되고, $f\in A-\mathfrak{j}_x$일 때의
$S'_f$들의 합집합은 $A-\mathfrak{j}_x$이다. 따라서
(\hyperref[0.1.4.5-ko]{0, 1.4.5}) $A_x$-가군 $M_x$는 준동형족
$(\rho_{g,f})$에 관한 귀납극한 $\varinjlim M_f$와 표준적으로
동일시된다. $f\in A-\mathfrak{j}_x$, 또는 같은 말로
$x\in D(f)$일 때의 표준 준동형을
\[
  \rho_x^f:M_f\longrightarrow M_x
\]
로 나타내겠다.

\begin{definition}[1.3.4]
\label{I.1.3.4-ko}
$D(f)$($f\in A$)들로 이루어진 $X$의 기저 $\mathfrak{B}$ 위의
준층 $D(f)\mapsto A_f$(각각 $D(f)\mapsto M_f$)에 대응하는 환의
층(각각 $\widetilde{A}$-가군)을 소 스펙트럼
$X=\operatorname{Spec}(A)$의 \emph{구조층}(각각 \emph{$A$-가군
$M$에 대응하는 층})이라 하고, $\widetilde{A}$ 또는
$\mathscr{O}_X$(각각 $\widetilde{M}$)로 나타낸다
((\hyperref[I.1.1.10-ko]{1.1.10}),
(\hyperref[0.3.2.1-ko]{0, 3.2.1} 및
\hyperref[0.3.5.6-ko]{3.5.6})).
\end{definition}

줄기 $\widetilde{A}_x$(각각 $\widetilde{M}_x$)가 환 $A_x$(각각
$A_x$-가군 $M_x$)와 동일시됨은 이미 보았다
(\hyperref[0.3.2.4-ko]{0, 3.2.4}). 표준 사상을
\[
  \theta_f:A_f\longrightarrow\Gamma(D(f),\widetilde{A})
\]
\[
  \text{(각각 }\theta_f:M_f\longrightarrow
  \Gamma(D(f),\widetilde{M})\text{)}
\]
로 나타내겠다. 그러면 모든 $x\in D(f)$와 모든 $\xi\in M_f$에
대하여
\begin{equation}
  (\theta_f(\xi))_x=\rho_x^f(\xi).
  \tag{1.3.4.1}\label{I.1.3.4.1-ko}
\end{equation}

\begin{proposition}[1.3.5]
\label{I.1.3.5-ko}
$\widetilde{M}$은 $M$에 관한 완전 공변 함자로서, $A$-가군의
범주에서 $\widetilde{A}$-가군의 범주로 간다.
\end{proposition}

실제로 $M$, $N$을 두 $A$-가군, $u:M\to N$을 준동형이라 하자.
모든 $f\in A$에 대하여 $u$에는 $A_f$-가군 $M_f$에서
$A_f$-가군 $N_f$로 가는 준동형 $u_f$가 표준적으로 대응하고,
$D(g)\subset D(f)$일 때 그림
\[
  \xymatrix{
    M_f\ar[r]^{u_f}\ar[d]_{\rho_{g,f}} & N_f\ar[d]^{\rho_{g,f}}\\
    M_g\ar[r]_{u_g} & N_g
  }
\]
은 가환한다(\hyperref[0.1.4.1-ko]{0, 1.4.1}). 따라서 이 준동형들은
$\widetilde{A}$-가군 준동형
$\widetilde{u}:\widetilde{M}\to\widetilde{N}$을 정의한다
(\hyperref[0.3.2.3-ko]{0, 3.2.3}). 더욱이 모든 $x\in X$에 대하여
$\widetilde{u}_x$는 $x\in D(f)$($f\in A$)일 때의 $u_f$들의
귀납극한이다. 따라서 $\widetilde{M}_x$와 $\widetilde{N}_x$를 각각
$M_x$와 $N_x$에 표준적으로 동일시하면
(\hyperref[0.1.4.5-ko]{0, 1.4.5}), $\widetilde{u}_x$는 $u$에서
표준적으로 얻는 준동형 $u_x$와 동일시된다. 이제 $P$를 세 번째
$A$-가군, $v:N\to P$를 준동형이라 하고 $w=v\circ u$로 놓으면
$w_x=v_x\circ u_x$이고, 따라서
$\widetilde{w}=\widetilde{v}\circ\widetilde{u}$임은 자명하다. 이로써
$M\mapsto\widetilde{M}$을 $A$-가군의 범주에서
$\widetilde{A}$-가군의 범주로 가는 \emph{공변 함자}로 정의하였다.
이 함자는 \emph{완전}하다. 실제로 모든 $x\in X$에 대하여 $M_x$는
$M$에 관한 완전 함자이기 때문이다
(\hyperref[0.1.3.2-ko]{0, 1.3.2}). 또한 두 지지집합의 정의에 따라
(\hyperref[0.1.7.1-ko]{0, 1.7.1} 및
\hyperref[0.3.1.6-ko]{3.1.6})
$\operatorname{Supp}(M)=\operatorname{Supp}(\widetilde{M})$이다.

\oldpage[I]{86}
\begin{proposition}[1.3.6]
\label{I.1.3.6-ko}
모든 $f\in A$에 대하여 열린집합 $D(f)\subset X$는 소 스펙트럼
$\operatorname{Spec}(A_f)$와 표준적으로 동일시되고, $A_f$-가군
$M_f$에 대응하는 층 $\widetilde{M_f}$는 제한
$\widetilde{M}|D(f)$와 표준적으로 동일시된다.
\end{proposition}

첫째 명제는 (\hyperref[I.1.2.6-ko]{1.2.6})의 특수한 경우이다.
더욱이 $D(g)\subset D(f)$인 $g\in A$에 대하여 $M_g$는 분모가
$A_f$에서 $g$의 표준 상의 거듭제곱인 $M_f$의 분수가군과
표준적으로 동일시된다(\hyperref[0.1.4.6-ko]{0, 1.4.6}). 이제
$\widetilde{M_f}$와 $\widetilde{M}|D(f)$의 표준적인 동일시는
정의에서 따른다.

\begin{theorem}[1.3.7]
\label{I.1.3.7-ko}
모든 $A$-가군 $M$과 모든 $f\in A$에 대하여 준동형
\[
  \theta_f:M_f\longrightarrow\Gamma(D(f),\widetilde{M})
\]
는 전단사이다. 다시 말해 준층 $D(f)\mapsto M_f$는 \emph{층}이다.
특히 $M$은 $\theta_1$에 의해 $\Gamma(X,\widetilde{M})$와
동일시된다.
\end{theorem}

$M=A$이면 $\theta_f$가 환 구조의 준동형임에 유의하자. 따라서
정리 (\hyperref[I.1.3.7-ko]{1.3.7})에 따라 $\theta_f$를 통하여
환 $A_f$와 $\Gamma(D(f),\widetilde{A})$를 동일시하면, 준동형
$\theta_f:M_f\to\Gamma(D(f),\widetilde{M})$는 \emph{가군의 동형}이
된다.

먼저 $\theta_f$가 \emph{단사}임을 보이자. 실제로
$\theta_f(\xi)=0$인 $\xi\in M_f$에 대하여, 이는 $A_f$의 모든 소
아이디얼 $\mathfrak{p}$마다 $h\notin\mathfrak{p}$이고 $h\xi=0$인
$h$가 존재한다는 뜻이다. 따라서 $\xi$의 소멸자는 $A_f$의 어느 소
아이디얼에도 포함되지 않으므로 $A_f$ 전체이고, 결국 $\xi=0$이다.

이제 $\theta_f$가 \emph{전사}임을 보이면 된다. 일반적인 경우는
(\hyperref[I.1.3.6-ko]{1.3.6})을 써서 ``국소화''함으로써 얻으므로
$f=1$인 경우로 귀착할 수 있다. $s$를 $X$ 위의
$\widetilde{M}$의 절단이라 하자. (\hyperref[I.1.3.4-ko]{1.3.4})와
(\hyperref[I.1.1.10-ko]{1.1.10, (ii)})에 따라 $X$의 \emph{유한}
덮개 $(D(f_i))_{i\in I}$($f_i\in A$)가 존재하여, 모든 $i\in I$에
대하여 제한 $s_i=s|D(f_i)$는 어떤 $\xi_i\in M_{f_i}$에 대한
$\theta_{f_i}(\xi_i)$ 꼴이다. $i$, $j$가 $I$의 두 지표이면 $s_i$와
$s_j$의 $D(f_i)\cap D(f_j)=D(f_i f_j)$로의 제한들이 같다는 조건은
$\widetilde{M}$의 정의에 따라
\begin{equation}
  \rho_{f_i f_j,f_i}(\xi_i)=\rho_{f_i f_j,f_j}(\xi_j).
  \tag{1.3.7.1}\label{I.1.3.7.1-ko}
\end{equation}
이 된다. 정의에 따라 모든 $i\in I$에 대하여
$\xi_i=z_i/f_i^{n_i}$($z_i\in M$)로 쓸 수 있다. $I$가 유한이므로
각 $z_i$에 $f_i$의 거듭제곱을 곱하여 모든 $n_i$가 같은 $n$이라고
가정할 수 있다. 그러면 (\hyperref[I.1.3.7.1-ko]{1.3.7.1})은
정의에 따라 어떤 정수 $m_{ij}\geq0$에 대하여
$(f_i f_j)^{m_{ij}}(f_j^n z_i-f_i^n z_j)=0$이라는 뜻이다. 모든
$m_{ij}$도 같은 정수 $m$이라고 가정할 수 있다. 이제 $z_i$를
$f_i^m z_i$로 바꾸면 $m=0$인 경우, 즉 모든 $i$, $j$에 대하여
\begin{equation}
  f_j^n z_i=f_i^n z_j
  \tag{1.3.7.2}\label{I.1.3.7.2-ko}
\end{equation}
인 경우로 귀착된다. 한편 $D(f_i^n)=D(f_i)$이고 $D(f_i)$들이 $X$를
덮으므로, $f_i^n$들이 생성하는 아이디얼은 $A$이다. 다시 말해
$\sum_i g_i f_i^n=1$인 원소 $g_i\in A$들이 존재한다. 이제
$M$의 원소 $z=\sum_i g_i z_i$를 생각하자. 공식
(\hyperref[I.1.3.7.2-ko]{1.3.7.2})에 따라
$f_i^n z=\sum_j g_j f_i^n z_j=(\sum_j g_j f_j^n)z_i=z_i$이고,
따라서 정의에 따라 $M_{f_i}$에서 $\xi_i=z/1$이다.
\oldpage[I]{87}
그러므로 $s_i$는 $\theta_1(z)$의 $D(f_i)$로의 제한이다. 따라서
$s=\theta_1(z)$이고 증명이 끝난다.

\begin{corollary}[1.3.8]
\label{I.1.3.8-ko}
$M$, $N$을 두 $A$-가군이라 하자. 표준 준동형
\[
  u\longmapsto\widetilde{u}:
  \operatorname{Hom}_A(M,N)\longrightarrow
  \operatorname{Hom}_{\widetilde{A}}(\widetilde{M},\widetilde{N})
\]
는 전단사이다. 특히 $M=0$과 $\widetilde{M}=0$은 서로 동치이다.
\end{corollary}

실제로 표준 준동형
\[
  v\longmapsto\Gamma(v):
  \operatorname{Hom}_{\widetilde{A}}(\widetilde{M},\widetilde{N})
  \longrightarrow
  \operatorname{Hom}_{\Gamma(\widetilde{A})}
  (\Gamma(\widetilde{M}),\Gamma(\widetilde{N}))
\]
을 생각하자. 정리 (\hyperref[I.1.3.7-ko]{1.3.7})에 따라 끝의
가군은 $\operatorname{Hom}_A(M,N)$과 표준적으로 동일시된다.
$u\mapsto\widetilde{u}$와 $v\mapsto\Gamma(v)$가 서로 역임을 확인하면
된다. $\widetilde{u}$의 정의에 따라
$\Gamma(\widetilde{u})=u$임은 자명하다. 반대로
$v\in\operatorname{Hom}_{\widetilde{A}}(\widetilde{M},\widetilde{N})$에
대하여 $u=\Gamma(v)$로 놓으면, $v$에서 표준적으로 얻는 사상
\[
  w:\Gamma(D(f),\widetilde{M})\longrightarrow
  \Gamma(D(f),\widetilde{N})
\]
에 대하여 그림
\[
  \xymatrix{
    M\ar[r]^u\ar[d]_{\rho_{f,1}} & N\ar[d]^{\rho_{f,1}}\\
    M_f\ar[r]_w & N_f
  }
\]
은 가환한다. 따라서 모든 $f\in A$에 대하여 반드시 $w=u_f$이고
(\hyperref[0.1.2.4-ko]{0, 1.2.4}), 이로부터
$\widetilde{\Gamma(v)}=v$가 따른다.

\begin{corollary}[1.3.9]
\label{I.1.3.9-ko}
\begin{enumerate}
  \item[\rm(i)] $u:M\to N$을 $A$-가군의 준동형이라 하자. 그러면
  $\operatorname{Ker}u$, $\operatorname{Im}u$,
  $\operatorname{Coker}u$에 대응하는 층들은 각각
  $\operatorname{Ker}\widetilde{u}$,
  $\operatorname{Im}\widetilde{u}$,
  $\operatorname{Coker}\widetilde{u}$이다. 특히 $\widetilde{u}$가
  단사(각각 전사, 전단사)이기 위한 필요충분조건은 $u$가
  단사(각각 전사, 전단사)인 것이다.
  \item[\rm(ii)] $M$이 $A$-가군족 $(M_\lambda)$의
  귀납극한(각각 직합)이면, $\widetilde{M}$은 표준 동형을 제외하고
  $(\widetilde{M_\lambda})$의 귀납극한(각각 직합)이다.
\end{enumerate}
\end{corollary}

(i) $\widetilde{M}$이 $M$에 관한 완전 함자라는 사실
(\hyperref[I.1.3.5-ko]{1.3.5})을 두 $A$-가군의 완전열
\[
  0\to\operatorname{Ker}u\to M\to\operatorname{Im}u\to0
\]
과
\[
  0\to\operatorname{Im}u\to N\to\operatorname{Coker}u\to0
\]
에 적용하면 된다. 두 번째 명제는 정리
(\hyperref[I.1.3.7-ko]{1.3.7})에서 따른다.

(ii) $(M_\lambda,g_{\mu\lambda})$를 귀납극한이 $M$인 $A$-가군의
귀납계라 하고, $g_\lambda:M_\lambda\to M$을 표준 준동형이라 하자.
$\lambda\leq\mu\leq\nu$에 대하여
\[
  \widetilde{g_{\nu\mu}}\circ\widetilde{g_{\mu\lambda}}
  =\widetilde{g_{\nu\lambda}},
  \qquad
  \widetilde{g_\lambda}=\widetilde{g_\mu}\circ
  \widetilde{g_{\mu\lambda}}
\]
이므로 $(\widetilde{M_\lambda},\widetilde{g_{\mu\lambda}})$는
$X$ 위 층들의 귀납계이다. 표준 준동형
\[
  h_\lambda:\widetilde{M_\lambda}\longrightarrow
  \varinjlim\widetilde{M_\lambda}
\]
를 생각하면
$v\circ h_\lambda=\widetilde{g_\lambda}$를 만족하는 유일한 준동형
\[
  v:\varinjlim\widetilde{M_\lambda}\longrightarrow\widetilde{M}
\]
가 존재한다. $v$가 전단사임을 보이려면 모든 $x\in X$에 대하여
$v_x$가 $(\varinjlim\widetilde{M_\lambda})_x$에서
$\widetilde{M}_x$ 위로 가는 전단사임을 확인하면 된다. 그런데
$\widetilde{M}_x=M_x$이고
\[
  (\varinjlim\widetilde{M_\lambda})_x
  =\varinjlim(\widetilde{M_\lambda})_x
  =\varinjlim(M_\lambda)_x=M_x
  \quad(\hyperref[0.1.3.3-ko]{0, 1.3.3}).
\]
또한 정의에 따라 $(\widetilde{g_\lambda})_x$와 $(h_\lambda)_x$는
모두 $(M_\lambda)_x$에서 $M_x$로 가는 표준 사상이다. 따라서
$(\widetilde{g_\lambda})_x=v_x\circ(h_\lambda)_x$이므로 $v_x$는
항등사상이다.
\oldpage[I]{88}
끝으로 $M$이 두 $A$-가군 $N$, $P$의 직합이면
$\widetilde{M}=\widetilde{N}\oplus\widetilde{P}$임은 자명하다. 모든
직합은 유한 직합들의 귀납극한이므로 (ii)가 증명되었다.

(\hyperref[I.1.3.8-ko]{1.3.8})에 따라 $A$-가군에 대응하는 층들과
동형인 층들은 \emph{아벨 범주}를 이룸에 유의하자(T, I, 1.4).

또한 (\hyperref[I.1.3.9-ko]{1.3.9})에서 다음이 따른다. $M$이
\emph{유한 생성} $A$-가군, 즉 전사 준동형 $A^n\to M$이 존재하는
가군이면 전사 준동형
$\widetilde{A^n}\to\widetilde{M}$이 존재한다. 다시 말해
$\widetilde{A}$-가군 $\widetilde{M}$은 \emph{$X$ 위의 유한한
절단족에 의해 생성된다}(\hyperref[0.5.1.1-ko]{0, 5.1.1}). 그 역도
성립한다.

\begin{env}[1.3.10]
\label{I.1.3.10-ko}
$N$이 $A$-가군 $M$의 부분가군이면 표준 단사 준동형 $j:N\to M$은
(\hyperref[I.1.3.9-ko]{1.3.9})에 따라 단사 준동형
$\widetilde{N}\to\widetilde{M}$을 준다. 이를 통하여
$\widetilde{N}$을 $\widetilde{M}$의
\emph{부분-$\widetilde{A}$-가군}과 표준적으로 동일시하며, 앞으로
항상 이 동일시를 사용하겠다. $N$, $P$가 $M$의 두 부분가군이면
\begin{equation}
  \widetilde{(N+P)}=\widetilde{N}+\widetilde{P}
  \tag{1.3.10.1}\label{I.1.3.10.1-ko}
\end{equation}
이고
\begin{equation}
  \widetilde{(N\cap P)}=\widetilde{N}\cap\widetilde{P}.
  \tag{1.3.10.2}\label{I.1.3.10.2-ko}
\end{equation}
실제로 $N+P$와 $N\cap P$는 각각 표준 준동형
$N\oplus P\to M$의 상과 표준 준동형
$M\to(M/N)\oplus(M/P)$의 핵이므로
(\hyperref[I.1.3.9-ko]{1.3.9})를 적용하면 된다.

(\hyperref[I.1.3.10.1-ko]{1.3.10.1})과
(\hyperref[I.1.3.10.2-ko]{1.3.10.2})에 따라
$\widetilde{N}=\widetilde{P}$이면 $N=P$이다.
\end{env}

\begin{corollary}[1.3.11]
\label{I.1.3.11-ko}
$A$-가군에 대응하는 층들과 동형인 층들의 범주에서 함자 $\Gamma$는
완전하다.
\end{corollary}

실제로 $A$-가군 준동형 $u:M\to N$, $v:N\to P$에 대응하는 완전열
\[
  \widetilde{M}\xrightarrow{\widetilde{u}}\widetilde{N}
  \xrightarrow{\widetilde{v}}\widetilde{P}
\]
을 생각하자. $Q=\operatorname{Im}u$, $R=\operatorname{Ker}v$로
놓으면 따름정리 (\hyperref[I.1.3.9-ko]{1.3.9})에 따라
\[
  \widetilde{Q}=\operatorname{Im}\widetilde{u}
  =\operatorname{Ker}\widetilde{v}=\widetilde{R},
\]
따라서 $Q=R$이다.

\begin{corollary}[1.3.12]
\label{I.1.3.12-ko}
$M$, $N$을 두 $A$-가군이라 하자.
\begin{enumerate}
  \item[\rm(i)] $M\otimes_A N$에 대응하는 층은
  $\widetilde{M}\otimes_{\widetilde{A}}\widetilde{N}$과 표준적으로
  동일시된다.
  \item[\rm(ii)] 더욱이 $M$이 유한 표시를 가지면
  $\operatorname{Hom}_A(M,N)$에 대응하는 층은
  $\mathscr{H}\!om_{\widetilde{A}}(\widetilde{M},\widetilde{N})$과
  표준적으로 동일시된다.
\end{enumerate}
\end{corollary}

(i) 층
$\mathscr{F}=\widetilde{M}\otimes_{\widetilde{A}}\widetilde{N}$은
$X$에서 $D(f)$($f\in A$)들로 이루어진 기저
(\hyperref[I.1.1.10-ko]{1.1.10, (i)}) 위의 준층
\[
  U\longmapsto\mathscr{F}(U)=\Gamma(U,\widetilde{M})
  \otimes_{\Gamma(U,\widetilde{A})}\Gamma(U,\widetilde{N})
\]
에 대응한다. 그런데 (\hyperref[I.1.3.7-ko]{1.3.7})과
(\hyperref[I.1.3.6-ko]{1.3.6})에 따라 $\mathscr{F}(D(f))$는
$M_f\otimes_{A_f}N_f$와 표준적으로 동일시된다. 한편 이
$A_f$-가군은 $(M\otimes_A N)_f$와 표준적으로 동형이고
(\hyperref[0.1.3.4-ko]{0, 1.3.4}), 다시 이는
$\Gamma(D(f),\widetilde{(M\otimes_A N)})$과 표준적으로 동형이다
(\hyperref[I.1.3.7-ko]{1.3.7} 및
\hyperref[I.1.3.6-ko]{1.3.6}). 또한 이렇게 얻은 표준 동형
\[
  \mathscr{F}(D(f))\simeq
  \Gamma(D(f),\widetilde{(M\otimes_A N)})
\]
\oldpage[I]{89}
들은 제한 준동형과의 양립 조건을 만족한다
(\hyperref[0.1.4.2-ko]{0, 1.4.2}). 따라서 이들은 함자적인 표준 동형
\[
  \widetilde{M}\otimes_{\widetilde{A}}\widetilde{N}
  \simeq\widetilde{(M\otimes_A N)}
\]
을 정의한다.

(ii) 층
$\mathscr{G}=\mathscr{H}\!om_{\widetilde{A}}(\widetilde{M},
\widetilde{N})$은 $D(f)$들로 이루어진 $X$의 기저 위의 준층
\[
  U\longmapsto\mathscr{G}(U)
  =\operatorname{Hom}_{\widetilde{A}|U}
  (\widetilde{M}|U,\widetilde{N}|U)
\]
에 대응한다. 그런데 $\mathscr{G}(D(f))$는
$\operatorname{Hom}_{A_f}(M_f,N_f)$와 표준적으로 동일시된다
(\hyperref[I.1.3.6-ko]{1.3.6} 및
\hyperref[I.1.3.8-ko]{1.3.8}). 이는 $M$에 관한 가정에 따라 다시
$(\operatorname{Hom}_A(M,N))_f$와 표준적으로 동일시되고
(\hyperref[0.1.3.5-ko]{0, 1.3.5}), 끝으로
$(\operatorname{Hom}_A(M,N))_f$는
$\Gamma(D(f),\widetilde{(\operatorname{Hom}_A(M,N))})$과
표준적으로 동일시된다(\hyperref[I.1.3.6-ko]{1.3.6} 및
\hyperref[I.1.3.7-ko]{1.3.7}). 이렇게 얻은 표준 동형
\[
  \mathscr{G}(D(f))\simeq
  \Gamma(D(f),\widetilde{(\operatorname{Hom}_A(M,N))})
\]
들은 제한 준동형과 양립하므로
(\hyperref[0.1.4.2-ko]{0, 1.4.2}) 표준 동형
\[
  \mathscr{H}\!om_{\widetilde{A}}(\widetilde{M},\widetilde{N})
  \simeq\widetilde{(\operatorname{Hom}_A(M,N))}
\]
을 정의한다.

\begin{env}[1.3.13]
\label{I.1.3.13-ko}
이제 $B$를 (가환) $A$-대수라 하자. 이는 $B$가 $A$-가군이고,
한 원소 $e\in B$와 $A$-준동형
$\varphi:B\otimes_A B\to B$가 주어져 다음 조건들을 만족한다고
해석할 수 있다. $1^\circ$ 그림
\[
  \xymatrix{
    B\otimes_A B\otimes_A B\ar[r]^{\varphi\otimes 1}
      \ar[d]_{1\otimes\varphi} &
    B\otimes_A B\ar[d]^\varphi & &
    B\otimes_A B\ar[rr]^\sigma\ar[rd]_\varphi & &
    B\otimes_A B\ar[dl]^\varphi\\
    B\otimes_A B\ar[r]_\varphi &
    B & & & B
  }
\]
들이 가환한다($\sigma$는 표준 대칭사상). $2^\circ$
$\varphi(e\otimes x)=\varphi(x\otimes e)=x$이다. 따름정리
(\hyperref[I.1.3.12-ko]{1.3.12})에 따라 $\widetilde{A}$-가군의
준동형
\[
  \widetilde{\varphi}:
  \widetilde{B}\otimes_{\widetilde{A}}\widetilde{B}
  \longrightarrow\widetilde{B}
\]
도 같은 조건들을 만족하므로 $\widetilde{B}$에
$\widetilde{A}$-대수 구조를 정의한다. 마찬가지로 $B$-가군 $N$을
주는 것은 $A$-가군 $N$과, 그림
\[
  \xymatrix{
    B\otimes_A B\otimes_A N\ar[r]^{\varphi\otimes 1}
      \ar[d]_{1\otimes\psi} &
    B\otimes_A N\ar[d]^\psi\\
    B\otimes_A N\ar[r]_\psi & N
  }
\]
이 가환하고 $\psi(e\otimes n)=n$을 만족하는 $A$-준동형
$\psi:B\otimes_A N\to N$을 주는 것과 같다. 준동형
\[
  \widetilde{\psi}:
  \widetilde{B}\otimes_{\widetilde{A}}\widetilde{N}
  \longrightarrow\widetilde{N}
\]
도 같은 조건을 만족하므로 $\widetilde{N}$에
$\widetilde{B}$-가군 구조를 정의한다.

같은 방식으로 다음을 알 수 있다. $u:B\to B'$가 $A$-대수의
준동형(각각 $v:N\to N'$이 $B$-가군의 준동형)이면
$\widetilde{u}$는 $\widetilde{A}$-대수의 준동형(각각
$\widetilde{v}$는 $\widetilde{B}$-가군의 준동형)이고,
$\operatorname{Ker}\widetilde{u}$는 $\widetilde{B}$-아이디얼이며
(각각 $\operatorname{Ker}\widetilde{v}$,
$\operatorname{Coker}\widetilde{v}$,
$\operatorname{Im}\widetilde{v}$는 $\widetilde{B}$-가군이다).
$N$이 $B$-가군이면 $\widetilde{N}$이 유한 생성
$\widetilde{B}$-가군이기 위한 필요충분조건은 $N$이 유한 생성
$B$-가군인 것이다(\hyperref[0.5.2.3-ko]{0, 5.2.3}).
\oldpage[I]{90}
$M$, $N$이 두 $B$-가군이면 $\widetilde{B}$-가군
$\widetilde{M}\otimes_{\widetilde{B}}\widetilde{N}$은
$\widetilde{(M\otimes_B N)}$과 표준적으로 동일시된다. 또한 $M$이
유한 표시를 가지면
$\mathscr{H}\!om_{\widetilde{B}}(\widetilde{M},\widetilde{N})$은
$\widetilde{(\operatorname{Hom}_B(M,N))}$과 표준적으로 동일시된다.
증명은 (\hyperref[I.1.3.12-ko]{1.3.12})와 같다.

$\mathfrak{J}$가 $B$의 아이디얼이고 $N$이 $B$-가군이면
\[
  \widetilde{(\mathfrak{J}N)}
  =\widetilde{\mathfrak{J}}\cdot\widetilde{N}.
\]

끝으로 $B$가 부분-$A$-가군 $B_n$($n\in\mathbb{Z}$)들로 등급이
매겨진 $A$-대수이면, $\widetilde{A}$-가군
$\widetilde{B_n}$들의 직합인 $\widetilde{A}$-대수
$\widetilde{B}$(\hyperref[I.1.3.9-ko]{1.3.9})도 이
부분-$\widetilde{A}$-가군들로 등급이 매겨진다. 실제로 등급 공리는
준동형 $B_m\otimes_A B_n\to B$의 상이 $B_{m+n}$에 포함된다는
조건으로 표현된다. 마찬가지로 $M$이 부분가군 $M_n$들로 등급이
매겨진 $B$-가군이면, $\widetilde{M}$은
$\widetilde{M_n}$들로 등급이 매겨진 $\widetilde{B}$-가군이다.
\end{env}

\begin{env}[1.3.14]
\label{I.1.3.14-ko}
$B$가 $A$-대수이고 $M$이 $B$의 부분가군이면,
$\widetilde{M}$이 생성하는 $\widetilde{B}$의
부분-$\widetilde{A}$-대수(\hyperref[0.4.1.3-ko]{0, 4.1.3})는
$\widetilde{C}$이다. 여기서 $C$는 $M$이 생성하는 $B$의 부분대수이다.
실제로 $C$는 준동형 $\bigotimes_A^n M\to B$($n\geq0$)의 상인
$B$의 부분가군들의 합이고, (\hyperref[I.1.3.9-ko]{1.3.9})와
(\hyperref[I.1.3.12-ko]{1.3.12})를 적용하면 된다.
\end{env}

\subsection{소 스펙트럼 위의 준연접층}
\label{subsection:I.1.4-ko}

\begin{theorem}[1.4.1]
\label{I.1.4.1-ko}
$X$를 환 $A$의 소 스펙트럼, $V$를 $X$의 준콤팩트 열린 부분집합,
$\mathscr{F}$를 $(\mathscr{O}_X|V)$-가군이라 하자. 다음 네 조건은
서로 동치이다.
\begin{enumerate}
  \item[\textit{a})] $\mathscr{F}$가 $\widetilde{M}|V$와 동형이 되는
  $A$-가군 $M$이 존재한다.
  \item[\textit{b})] $V$에 포함되고 $D(f_i)$($f_i\in A$) 꼴인
  집합들로 이루어진 $V$의 유한 열린 덮개 $(V_i)$가 존재하여, 모든
  $i$에 대하여 $\mathscr{F}|V_i$는 어떤 $A_{f_i}$-가군 $M_i$에
  대응하는 층 $\widetilde{M_i}$와 동형이다.
  \item[\textit{c})] 층 $\mathscr{F}$는 준연접이다
  (\hyperref[0.5.1.3-ko]{0, 5.1.3}).
  \item[\textit{d})] 다음 두 성질이 성립한다.
  \begin{enumerate}
    \item[\textit{d}\,1)] $D(f)\subset V$인 모든 $f\in A$와 모든 절단
    $s\in\Gamma(D(f),\mathscr{F})$에 대하여, $f^n s$가 $V$ 위의
    $\mathscr{F}$의 절단으로 연장되는 정수 $n\geq0$이 존재한다.
    \item[\textit{d}\,2)] $D(f)\subset V$인 모든 $f\in A$와,
    $D(f)$로의 제한이 $0$인 모든 절단
    $t\in\Gamma(V,\mathscr{F})$에 대하여, $f^n t=0$인 정수
    $n\geq0$이 존재한다.
  \end{enumerate}
\end{enumerate}
\end{theorem}

(조건 \textit{d}\,1)과 \textit{d}\,2)를 쓸 때에는
(\hyperref[I.1.3.7-ko]{1.3.7})에 따라 $A$와
$\Gamma(\widetilde{A})$를 암묵적으로 동일시하였다.)

\textit{a})가 \textit{b})를 함의함은
(\hyperref[I.1.3.6-ko]{1.3.6})과 $D(f_i)$들이 $X$의 위상의 기저를
이룬다는 사실(\hyperref[I.1.1.10-ko]{1.1.10})에서 바로 따른다.
모든 $A$-가군은 어떤 준동형 $A^{(I)}\to A^{(J)}$의 여핵과
동형이므로, (\hyperref[I.1.3.9-ko]{1.3.9})에 따라 모든 $A$-가군에
대응하는 층은 준연접이다. 따라서 \textit{b})는 \textit{c})를
함의한다. 역으로 $\mathscr{F}$가 준연접이면 모든 $x\in V$는 다음
성질을 갖는 $D(f)\subset V$ 꼴의 근방을 가진다. 즉,
$\mathscr{F}|D(f)$는 어떤 준동형
$\widetilde{A_f}^{(I)}\to\widetilde{A_f}^{(J)}$의 여핵과 동형이고,
따라서 이에 대응하는 준동형 $A_f^{(I)}\to A_f^{(J)}$의 여핵인
가군 $N$에 대응하는 층 $\widetilde{N}$과 동형이다
(\hyperref[I.1.3.8-ko]{1.3.8} 및
\hyperref[I.1.3.9-ko]{1.3.9}). $V$가 준콤팩트이므로
\textit{c})가 \textit{b})를 함의함은 자명하다.
\oldpage[I]{91}
\textit{b})가 \textit{d}\,1)과 \textit{d}\,2)를 함의함을
증명하기 위해 먼저 어떤 $g\in A$에 대하여 $V=D(g)$이고,
$\mathscr{F}$가 $A_g$-가군 $N$에 대응하는 층
$\widetilde{N}$과 동형이라고 하자. $X$를 $V$로, $A$를 $A_g$로
바꾸면(\hyperref[I.1.3.6-ko]{1.3.6}) $g=1$인 경우로 귀착할 수
있다. 그러면 $\Gamma(D(f),\widetilde{N})$과 $N_f$는 표준적으로
동일시된다(\hyperref[I.1.3.6-ko]{1.3.6} 및
\hyperref[I.1.3.7-ko]{1.3.7}). 따라서 절단
$s\in\Gamma(D(f),\widetilde{N})$는 $z\in N$인 $z/f^n$ 꼴의
원소와 동일시된다. 절단 $f^n s$는 $N_f$의 원소 $z/1$과
동일시되고, 따라서 $z\in N$과 동일시되는 $X$ 위의
$\widetilde{N}$의 절단을 $D(f)$로 제한한 것이다. 이로써 이 경우의
\textit{d}\,1)을 얻는다. 마찬가지로
$t\in\Gamma(X,\widetilde{N})$는 어떤 원소 $z'\in N$과 동일시된다.
$t$의 $D(f)$로의 제한은 $N_f$에서 $z'$의 상 $z'/1$과 동일시되고,
이 상이 $0$이라는 것은 $N$에서 $f^n z'=0$인 어떤 $n\geq0$이
존재한다는 뜻이다. 이는 $f^n t=0$과 같다.

\textit{b})가 \textit{d}\,1)과 \textit{d}\,2)를 함의함을
마치려면 다음 보조정리를 증명하면 충분하다.

\begin{lemma}[1.4.1.1]
\label{I.1.4.1.1-ko}
$V$가 $D(g_i)$ 꼴의 집합들의 유한 합집합이고, 각 층
$\mathscr{F}|D(g_i)$와
$\mathscr{F}|(D(g_i)\cap D(g_j))=\mathscr{F}|D(g_i g_j)$가
\textit{d}\,1)과 \textit{d}\,2)를 만족한다고 하자. 그러면
$\mathscr{F}$는 다음 두 성질을 갖는다.
\begin{enumerate}
  \item[\textit{d}'\,1)] 모든 $f\in A$와 모든 절단
  $s\in\Gamma(D(f)\cap V,\mathscr{F})$에 대하여, $f^n s$가
  $V$ 위의 $\mathscr{F}$의 절단으로 연장되는 정수 $n\geq0$이
  존재한다.
  \item[\textit{d}'\,2)] 모든 $f\in A$와, $D(f)\cap V$로의 제한이
  $0$인 모든 절단 $t\in\Gamma(V,\mathscr{F})$에 대하여,
  $f^n t=0$인 정수 $n\geq0$이 존재한다.
\end{enumerate}
\end{lemma}

먼저 \textit{d}'\,2)를 증명하자. $D(f)\cap D(g_i)=D(fg_i)$이므로
모든 $i$에 대하여 $(fg_i)^{n_i}t$의 $D(g_i)$로의 제한이 $0$인
정수 $n_i$가 존재한다. $A_{g_i}$에서 $g_i$의 상은 가역원이므로
$f^{n_i}t$의 $D(g_i)$로의 제한도 $0$이다. $n$을 $n_i$들의
최댓값으로 취하면 \textit{d}'\,2)를 얻는다.

\textit{d}'\,1)을 증명하기 위해 층 $\mathscr{F}|D(g_i)$에
\textit{d}\,1)을 적용하자. 그러면 정수 $n_i\geq0$과
$D(g_i)$ 위의 $\mathscr{F}$의 절단 $s_i'$이 존재하여,
$s_i'$은 $(fg_i)^{n_i}s$의 $D(fg_i)$로의 제한을 연장한다.
$A_{g_i}$에서 $g_i$의 상이 가역원이므로, $s_i'=g_i^{n_i}s_i$이고
$f^{n_i}s$의 $D(fg_i)$로의 제한을 연장하는 $D(g_i)$ 위의 절단
$s_i$가 존재한다. 모든 $n_i$가 같은 정수 $n$이라고 가정할 수도
있다. 구성에 따라 $s_i-s_j$의
$D(f)\cap D(g_i)\cap D(g_j)=D(fg_i g_j)$로의 제한은 $0$이다.
층 $\mathscr{F}|D(g_i g_j)$에 \textit{d}\,2)를 적용하면,
$(fg_i g_j)^{m_{ij}}(s_i-s_j)$의 $D(g_i g_j)$로의 제한이 $0$인
정수 $m_{ij}\geq0$이 존재한다. $A_{g_i g_j}$에서 $g_i g_j$의 상이
가역원이므로 $f^{m_{ij}}(s_i-s_j)$의 $D(g_i g_j)$로의 제한도
$0$이다. 모든 $m_{ij}$가 같은 정수 $m$이라고 가정할 수 있다.
그러면 $f^m s_i$들을 연장하는 절단
$s'\in\Gamma(V,\mathscr{F})$가 존재하고, 이 절단은
$f^{n+m}s$를 연장한다. 이로써 \textit{d}'\,1)을 얻는다.

이제 \textit{d}\,1)과 \textit{d}\,2)가 \textit{a})를
함의함을 증명하면 된다. 먼저 이 두 조건이, $D(g)\subset V$인 모든
$g\in A$에 대하여 층 $\mathscr{F}|D(g)$에도 성립함을 보이자.
\textit{d}\,1)에 대해서는 자명하다. 한편
$t\in\Gamma(D(g),\mathscr{F})$이고 그 $D(f)\subset D(g)$로의
제한이 $0$이라고 하자. \textit{d}\,1)에 따라 $g^m t$가 $V$ 위의
$\mathscr{F}$의 절단 $s$로 연장되는 정수 $m\geq0$이 존재한다.
\oldpage[I]{92}
\textit{d}\,2)를 적용하면 $f^n g^m t=0$인 정수 $n\geq0$이
존재한다. $A_g$에서 $g$의 상이 가역원이므로 $f^n t=0$이다.

이제 $V$가 준콤팩트이므로 보조정리
(\hyperref[I.1.4.1.1-ko]{1.4.1.1})에 따라 조건
\textit{d}'\,1)과 \textit{d}'\,2)가 성립한다. $A$-가군
$M=\Gamma(V,\mathscr{F})$을 생각하자. $j:V\to X$를 표준
포함사상이라 하고 $\widetilde{A}$-가군 준동형
\[
  u:\widetilde{M}\longrightarrow j_*(\mathscr{F})
\]
을 정의하겠다. $D(f)$들이 $X$의 위상의 기저를 이루므로, 모든
$f\in A$에 대하여 통상적인 양립 조건
(\hyperref[0.3.2.5-ko]{0, 3.2.5})을 만족하는 준동형
\[
  u_f:M_f\longrightarrow
  \Gamma(D(f),j_*(\mathscr{F}))
  =\Gamma(D(f)\cap V,\mathscr{F})
\]
을 정의하면 충분하다. $A_f$에서 $f$의 표준 상이 가역원이므로 제한
준동형
\[
  M=\Gamma(V,\mathscr{F})\longrightarrow
  \Gamma(D(f)\cap V,\mathscr{F})
\]
은 다음과 같이 분해된다
(\hyperref[0.1.2.4-ko]{0, 1.2.4}).
\[
  M\longrightarrow M_f\xrightarrow{u_f}
  \Gamma(D(f)\cap V,\mathscr{F}).
\]
$D(g)\subset D(f)$에 대한 양립 조건은 바로 확인된다.
\textit{d}'\,1)(각각 \textit{d}'\,2))이 모든 $u_f$가
전사(각각 단사)임을 함의함을 보이면 $u$가 \emph{전단사}이고,
따라서 $\mathscr{F}$는 $\widetilde{M}|V$와 동형이다. 그런데
$s\in\Gamma(D(f)\cap V,\mathscr{F})$이면 \textit{d}'\,1)에 따라
$f^n s$가 어떤 절단 $z\in M$으로 연장되는 정수 $n\geq0$이
존재한다. 그러면 $u_f(z/f^n)=s$이므로 $u_f$는 전사이다. 마찬가지로
$u_f(z/1)=0$인 $z\in M$에 대하여, 이는 절단 $z$의
$D(f)\cap V$로의 제한이 $0$이라는 뜻이다. \textit{d}'\,2)에 따라
$f^n z=0$인 정수 $n\geq0$이 존재하므로 $M_f$에서 $z/1=0$이고,
따라서 $u_f$는 단사이다.

\par\hfill C.Q.F.D.\par

\begin{corollary}[1.4.2]
\label{I.1.4.2-ko}
$X$의 준콤팩트 열린집합 위의 모든 준연접층은 $X$ 위의 어떤
준연접층으로부터 유도된다.
\end{corollary}

\begin{corollary}[1.4.3]
\label{I.1.4.3-ko}
$X=\operatorname{Spec}(A)$ 위의 모든 준연접
$\mathscr{O}_X$-대수는 어떤 $A$-대수 $B$에 대한
$\widetilde{B}$ 꼴의 $\mathscr{O}_X$-대수와 동형이다. 모든 준연접
$\widetilde{B}$-가군은 어떤 $B$-가군 $N$에 대한
$\widetilde{N}$ 꼴의 $\widetilde{B}$-가군과 동형이다.
\end{corollary}

실제로 준연접 $\mathscr{O}_X$-대수는 준연접
$\mathscr{O}_X$-가군이므로, 어떤 $A$-가군 $B$에 대한
$\widetilde{B}$ 꼴이다. $B$가 $A$-대수라는 것은
$\widetilde{A}$-가군 준동형
\[
  \widetilde{B}\otimes_{\widetilde{A}}\widetilde{B}
  \longrightarrow\widetilde{B}
\]
을 써서 $\mathscr{O}_X$-대수 구조를 특징짓는 것과
(\hyperref[I.1.3.12-ko]{1.3.12})에서 따른다. $\mathscr{G}$가 준연접
$\widetilde{B}$-가군이면 $\mathscr{G}$가 준연접
$\widetilde{A}$-가군임을 보이기만 하면 같은 방식으로 결론을 얻는다.
문제는 국소적이므로 $X$의 $D(f)$ 꼴 열린집합으로 제한하여
$\mathscr{G}$가 $\widetilde{B}$-가군(따라서 특히
$\widetilde{A}$-가군)의 준동형
\[
  \widetilde{B}^{(I)}\longrightarrow\widetilde{B}^{(J)}
\]
의 여핵이라고 가정할 수 있다. 이제 명제는
(\hyperref[I.1.3.8-ko]{1.3.8})과
(\hyperref[I.1.3.9-ko]{1.3.9})에서 따른다.

\subsection{소 스펙트럼 위의 연접층}
\label{subsection:I.1.5-ko}

\begin{theorem}[1.5.1]
\label{I.1.5.1-ko}
$A$를 \emph{뇌터 환}, $X=\operatorname{Spec}(A)$를 그 소 스펙트럼,
$V$를 $X$의 열린 부분집합, $\mathscr{F}$를
$(\mathscr{O}_X|V)$-가군이라 하자. 다음 조건들은 서로 동치이다.
\begin{enumerate}
  \item[\textit{a})] $\mathscr{F}$는 연접이다.
  \item[\textit{b})] $\mathscr{F}$는 유한 생성이고 준연접이다.
  \item[\textit{c})] $\mathscr{F}$가 층 $\widetilde{M}|V$와 동형이
  되는 유한 생성 $A$-가군 $M$이 존재한다.
\end{enumerate}
\end{theorem}

\oldpage[I]{93}
\textit{a})가 \textit{b})를 함의함은 자명하다. \textit{b})가
\textit{c})를 함의함을 보이자. $V$는 준콤팩트이므로
(\hyperref[0.2.2.3-ko]{0, 2.2.3}), 먼저 $\mathscr{F}$는 어떤
$A$-가군 $N$에 대한 층 $\widetilde{N}|V$와 동형이다
(\hyperref[I.1.4.1-ko]{1.4.1}). 그런데 $M_\lambda$가 $N$의 유한 생성
부분-$A$-가군 전체를 달릴 때
$N=\varinjlim M_\lambda$이므로, (\hyperref[I.1.3.9-ko]{1.3.9})에 따라
\[
  \mathscr{F}=\widetilde{N}|V
  =\varinjlim\widetilde{M_\lambda}|V.
\]
한편 $\mathscr{F}$는 유한 생성이고 $V$는 준콤팩트이므로 어떤 지표
$\lambda$에 대하여
$\mathscr{F}=\widetilde{M_\lambda}|V$이다
(\hyperref[0.5.2.3-ko]{0, 5.2.3}).

끝으로 \textit{c})가 \textit{a})를 함의함을 보이자.
$\mathscr{F}$가 유한 생성임은 분명하다
(\hyperref[I.1.3.6-ko]{1.3.6} 및
\hyperref[I.1.3.9-ko]{1.3.9}). 또한 문제는 국소적이므로
$V=D(f)$($f\in A$)인 경우만 생각하면 된다. $A_f$는 뇌터 환이므로
$A$를 $A_f$로 바꾸어, 결국 $M$이 $A$-가군일 때 준동형
$\widetilde{A^n}\to\widetilde{M}$의 핵이 유한 생성임을 증명하면
충분하다. 이러한 준동형은 어떤 준동형 $u:A^n\to M$에 대한
$\widetilde{u}$ 꼴이다(\hyperref[I.1.3.8-ko]{1.3.8}).
$P=\operatorname{Ker}u$라 놓으면
\[
  \widetilde{P}=\operatorname{Ker}\widetilde{u}
\]
이다(\hyperref[I.1.3.9-ko]{1.3.9}). $A$가 뇌터 환이므로 $P$는
유한 생성이고, 이로써 증명이 끝난다.

\begin{corollary}[1.5.2]
\label{I.1.5.2-ko}
(\hyperref[I.1.5.1-ko]{1.5.1})의 가정 아래에서 층
$\mathscr{O}_X$는 연접인 환의 층이다.
\end{corollary}

\begin{corollary}[1.5.3]
\label{I.1.5.3-ko}
(\hyperref[I.1.5.1-ko]{1.5.1})의 가정 아래에서 $X$의 열린집합 위의
모든 연접층은 $X$ 위의 어떤 연접층으로부터 유도된다.
\end{corollary}

\begin{corollary}[1.5.4]
\label{I.1.5.4-ko}
(\hyperref[I.1.5.1-ko]{1.5.1})의 가정 아래에서 모든 준연접
$\mathscr{O}_X$-가군 $\mathscr{F}$는 $\mathscr{F}$의 연접
부분-$\mathscr{O}_X$-가군들의 귀납극한이다.
\end{corollary}

실제로 $\mathscr{F}=\widetilde{M}$인 $A$-가군 $M$이 존재하고,
$M$은 그 유한 생성 부분가군들의 귀납극한이다. 이제
(\hyperref[I.1.3.9-ko]{1.3.9})와
(\hyperref[I.1.5.1-ko]{1.5.1})을 적용하면 된다.

\subsection{소 스펙트럼 위의 준연접층의 함자적 성질}
\label{subsection:I.1.6-ko}

\begin{env}[1.6.1]
\label{I.1.6.1-ko}
$A$, $A'$을 두 환,
\[
  \varphi:A'\to A
\]
를 준동형,
\[
  {}^a\varphi:X=\operatorname{Spec}(A)\to
  X'=\operatorname{Spec}(A')
\]
를 $\varphi$에 대응하는 연속사상이라 하자
(\hyperref[I.1.2.1-ko]{1.2.1}). 환의 층의 \emph{표준 준동형}
\[
  \widetilde{\varphi}:\mathscr{O}_{X'}\to
  {}^a\varphi_*(\mathscr{O}_X)
\]
을 정의하겠다. 모든 $f'\in A'$에 대하여 $f=\varphi(f')$로 놓자.
그러면
${}^a\varphi^{-1}(D(f'))=D(f)$이다
(\hyperref[I.1.2.2.2-ko]{1.2.2.2}). 환
$\Gamma(D(f'),\widetilde{A'})$와
$\Gamma(D(f),\widetilde{A})$는 각각 $A'_{f'}$와 $A_f$와 동일시된다
(\hyperref[I.1.3.6-ko]{1.3.6} 및
\hyperref[I.1.3.7-ko]{1.3.7}). 한편 준동형 $\varphi$는 표준적으로
준동형 $\varphi_{f'}:A'_{f'}\to A_f$를 정의한다
(\hyperref[0.1.5.1-ko]{0, 1.5.1}). 다시 말해 환 준동형
\[
  \Gamma(D(f'),\widetilde{A'})\to
  \Gamma({}^a\varphi^{-1}(D(f')),\widetilde{A})
  =\Gamma(D(f'),{}^a\varphi_*(\widetilde{A}))
\]
이 존재한다.

\oldpage[I]{94}
또한 이 준동형들은 통상적인 양립 조건을 만족한다. 즉,
$D(f')\supset D(g')$일 때 그림
\[
  \xymatrix{
    \Gamma(D(f'),\widetilde{A'})\ar[r]\ar[d] &
    \Gamma(D(f'),{}^a\varphi_*(\widetilde{A}))\ar[d]\\
    \Gamma(D(g'),\widetilde{A'})\ar[r] &
    \Gamma(D(g'),{}^a\varphi_*(\widetilde{A}))
  }
\]
은 가환이다(\hyperref[0.1.5.1-ko]{0, 1.5.1}). $D(f')$들이
$X'$의 위상의 기저를 이루므로
(\hyperref[0.3.2.3-ko]{0, 3.2.3}), 이로써
$\mathscr{O}_{X'}$-대수의 준동형을 정의하였다. 따라서 쌍
$\Phi=({}^a\varphi,\widetilde{\varphi})$는 환 달린 공간의 사상
\[
  \Phi:(X,\mathscr{O}_X)\to(X',\mathscr{O}_{X'})
\]
이다(\hyperref[0.4.1.1-ko]{0, 4.1.1}).

더욱이 $x'={}^a\varphi(x)$로 놓으면 준동형
$\widetilde{\varphi}_x^\#$
(\hyperref[0.3.7.1-ko]{0, 3.7.1})은 바로 $\varphi:A'\to A$에서
표준적으로 유도되는 준동형
\[
  \varphi_x:A'_{x'}\to A_x
\]
이다(\hyperref[0.1.5.1-ko]{0, 1.5.1}). 실제로 모든
$z'\in A'_{x'}$은 $f',g'\in A'$이고
$f'\notin\mathfrak{j}_{x'}$인 $g'/f'$ 꼴로 쓸 수 있다. 따라서
$D(f')$는 $X'$에서 $x'$의 근방이고, $\widetilde{\varphi}$에서
유도되는 준동형
$\Gamma(D(f'),\widetilde{A'})\to
\Gamma({}^a\varphi^{-1}(D(f')),\widetilde{A})$은
$\varphi_{f'}$와 같다. $g'/f'\in A'_{f'}$에 대응하는 절단
$s'\in\Gamma(D(f'),\widetilde{A'})$을 생각하면 $A_x$에서
$\widetilde{\varphi}_x^\#(z')=\varphi(g')/\varphi(f')$을 얻는다.
\end{env}

\begin{example}[1.6.2]
\label{I.1.6.2-ko}
$S$를 $A$의 곱셈적 부분집합, $\varphi$를 표준 준동형
$A\to S^{-1}A$라 하자. 앞서
(\hyperref[I.1.2.6-ko]{1.2.6})에서 ${}^a\varphi$가
$Y=\operatorname{Spec}(S^{-1}A)$에서
$X=\operatorname{Spec}(A)$의 부분공간
\[
  \{x\in X\mid\mathfrak{j}_x\cap S=\varnothing\}
\]
으로 가는 \emph{위상동형}임을 보았다. 또한 이 부분공간의 모든
$x={}^a\varphi(y)$($y\in Y$)에 대하여 준동형
$\widetilde{\varphi}_y^\#:\mathscr{O}_x\to\mathscr{O}_y$는
\emph{전단사}이다(\hyperref[0.1.2.6-ko]{0, 1.2.6}). 다시 말해
$\mathscr{O}_Y$는 $\mathscr{O}_X$가 $Y$ 위에 유도하는 층과
동일시된다.
\end{example}

\begin{proposition}[1.6.3]
\label{I.1.6.3-ko}
모든 $A$-가군 $M$에 대하여 $\mathscr{O}_{X'}$-가군
$(M_{[\varphi]})^{\sim}$에서 직상 $\Phi_*(\widetilde{M})$으로 가는
함자적인 표준 동형이 존재한다.
\end{proposition}

간단히 $M'=M_{[\varphi]}$로 놓고, 모든 $f'\in A'$에 대하여 다시
$f=\varphi(f')$로 놓자. 절단가군
$\Gamma(D(f'),\widetilde{M'})$와
$\Gamma(D(f),\widetilde{M})$는 각각 $A'_{f'}$-가군 $M'_{f'}$와
$A_f$-가군 $M_f$와 동일시된다. 또한 $A'_{f'}$-가군
$(M_f)_{[\varphi_{f'}]}$는 $M'_{f'}$와 표준적으로 동형이다
(\hyperref[0.1.5.2-ko]{0, 1.5.2}). 따라서
$\Gamma(D(f'),\widetilde{A'})$-가군의 함자적 동형
\[
  \Gamma(D(f'),\widetilde{M'})\xrightarrow{\sim}
  \Gamma({}^a\varphi^{-1}(D(f')),\widetilde{M})_{[\varphi_{f'}]}
\]
이 존재한다. 이 동형들은 제한사상과의 통상적인 양립 조건을
만족하므로(\hyperref[0.1.5.6-ko]{0, 1.5.6}) 앞에서 말한 함자적 동형을
정의한다. 더 정확히, $u:M_1\to M_2$가 $A$-가군 준동형이면 이를
$A'$-가군 준동형
$(M_1)_{[\varphi]}\to(M_2)_{[\varphi]}$로 볼 수 있다. 이 준동형을
$u_{[\varphi]}$라 하면 $\Phi_*(\widetilde{u})$는
$(u_{[\varphi]})^{\sim}$와 동일시된다.

이 증명은 모든 $A$-대수 $B$에 대하여 함자적인 표준 동형
\oldpage[I]{95}
\[
  (B_{[\varphi]})^{\sim}\xrightarrow{\sim}\Phi_*(\widetilde{B})
\]
이 $\mathscr{O}_{X'}$-대수의 동형임도 보여 준다. $M$이
$B$-가군이면 함자적인 표준 동형
\[
  (M_{[\varphi]})^{\sim}\xrightarrow{\sim}\Phi_*(\widetilde{M})
\]
은 $\Phi_*(\widetilde{B})$-가군의 동형이다.

\begin{corollary}[1.6.4]
\label{I.1.6.4-ko}
직상함자 $\Phi_*$는 준연접 $\mathscr{O}_X$-가군의 범주에서
완전하다.
\end{corollary}

실제로 $M_{[\varphi]}$는 $M$에 관한 완전 함자이고,
$\widetilde{M'}$은 $M'$에 관한 완전 함자이다
(\hyperref[I.1.3.5-ko]{1.3.5}).

\begin{proposition}[1.6.5]
\label{I.1.6.5-ko}
$N'$을 $A'$-가군이라 하고 $N=N'\otimes_{A'}A_{[\varphi]}$를
$A$-가군이라 하자. $\mathscr{O}_X$-가군
$\Phi^*(\widetilde{N'})$에서 $\widetilde{N}$으로 가는 함자적인
표준 동형이 존재한다.
\end{proposition}

먼저 $j:z'\mapsto z'\otimes1$이 $N'$에서 $N_{[\varphi]}$으로 가는
$A'$-준동형임을 보자. 실제로 $f'\in A'$에 대하여 정의상
$(f'z')\otimes1=z'\otimes\varphi(f')=\varphi(f')(z'\otimes1)$이다.
따라서 (\hyperref[I.1.3.8-ko]{1.3.8})에 의해
$\mathscr{O}_{X'}$-가군 준동형
$\widetilde{j}:\widetilde{N'}\to(N_{[\varphi]})^{\sim}$을 얻는다.
(\hyperref[I.1.6.3-ko]{1.6.3})에 따라 $\widetilde{j}$가
$\widetilde{N'}$을 $\Phi_*(\widetilde{N})$으로 보낸다고 보아도 된다.
이 준동형 $\widetilde{j}$에는 표준적으로 준동형
$h=\widetilde{j}^{\#}:\Phi^*(\widetilde{N'})\to\widetilde{N}$이
대응한다(\hyperref[0.4.4.3-ko]{0, 4.4.3}). 모든 줄기에서 $h_x$가
\emph{전단사}임을 보이자. $x'={}^a\varphi(x)$로 놓고
$x'\in D(f')$인 $f'\in A'$를 잡은 뒤 $f=\varphi(f')$로 놓자. 환
$\Gamma(D(f),\widetilde{A})$는 $A_f$와, 가군
$\Gamma(D(f),\widetilde{N})$와
$\Gamma(D(f'),\widetilde{N'})$은 각각 $N_f$와 $N'_{f'}$와
동일시된다. 절단 $s'\in\Gamma(D(f'),\widetilde{N'})$이
$n'/{f'}^p$($n'\in N'$)와 동일시된다고 하고, $s$를
$\widetilde{j}$에 의한 그 상이라 하자. 그러면 $s$는
$(n'\otimes1)/f^p$와 동일시된다. 한편
$t\in\Gamma(D(f),\widetilde{A})$가 $g/f^q$($g\in A$)와
동일시된다고 하자. 정의에 따라
\[
  h_x(s'_{x'}\otimes t_x)=t_x\cdot s_x
\]
이다(\hyperref[0.4.4.3-ko]{0, 4.4.3}). 그런데
$N_f$는 표준적으로
\[
  N'_{f'}\otimes_{A'_{f'}}(A_f)_{[\varphi_{f'}]}
\]
와 동일시된다(\hyperref[0.1.5.4-ko]{0, 1.5.4}). 이때 $s$는
$(n'/{f'}^p)\otimes1$에 대응하고, 절단 $y\mapsto t_y\cdot s_y$는
$(n'/{f'}^p)\otimes(g/f^q)$에 대응한다.
(\hyperref[0.1.5.6-ko]{0, 1.5.6})의 양립 그림들은 $h_x$가 바로
표준 동형
\[
  \label{I.1.6.5.1-ko}
  N'_{x'}\otimes_{A'_{x'}}(A_x)_{[\varphi_x]}
  \xrightarrow{\sim}N_x=(N'\otimes_{A'}A_{[\varphi]})_x
  \tag{1.6.5.1}
\]
임을 보여 준다.

또한 $v:N'_1\to N'_2$가 $A'$-가군 준동형이면 모든 $x'\in X'$에
대하여 $\widetilde{v}_{x'}=v_{x'}$이다. 그러므로 앞의 결과에서 곧바로
$\Phi^*(\widetilde{v})$가 $(v\otimes1)^{\sim}$와 표준적으로
동일시됨을 얻으며, 이로써
(\hyperref[I.1.6.5-ko]{1.6.5})의 증명이 끝난다.

$B'$가 $A'$-대수이면 $\Phi^*(\widetilde{B'})$에서
$(B'\otimes_{A'}A_{[\varphi]})^{\sim}$로 가는 표준 동형은
$\mathscr{O}_X$-대수의 동형이다. 더욱이 $N'$이 $B'$-가군이면
$\Phi^*(\widetilde{N'})$에서
$(N'\otimes_{A'}A_{[\varphi]})^{\sim}$로 가는 표준 동형은
$\Phi^*(\widetilde{B'})$-가군의 동형이다.

\begin{corollary}[1.6.6]
\label{I.1.6.6-ko}
$s'$이 $A'$-가군 $\Gamma(\widetilde{N'})$을 달릴 때, 절단 $s'$의
$\Phi^*(\widetilde{N'})$ 안의 표준 상들은 $A$-가군
$\Gamma(X,\Phi^*(\widetilde{N'}))$을 생성한다.
\end{corollary}

실제로 $N'$과 $N$을 각각 $\Gamma(\widetilde{N'})$과
$\Gamma(\widetilde{N})$와 동일시하면
(\hyperref[I.1.3.7-ko]{1.3.7}), 이 상들은 $z'$이 $N'$을 달릴 때의
$N$의 원소 $z'\otimes1$과 동일시된다.

\begin{env}[1.6.7]
\label{I.1.6.7-ko}
(\hyperref[I.1.6.5-ko]{1.6.5})의 증명에서 표준 사상
(\hyperref[0.4.4.3.2-ko]{0, 4.4.3.2})
$\rho:\widetilde{N'}\to\Phi_*(\Phi^*(\widetilde{N'}))$이 바로
준동형 $\widetilde{j}$임을 함께 증명하였다.
\oldpage[I]{96}
여기서 $j:N'\to N'\otimes_{A'}A_{[\varphi]}$는
$z'\mapsto z'\otimes1$인 준동형이다. 마찬가지로 표준 사상
(\hyperref[0.4.4.3.3-ko]{0, 4.4.3.3})
$\sigma:\Phi^*(\Phi_*(\widetilde{M}))\to\widetilde{M}$은
$\widetilde{p}$와 같다. 여기서
$p:M_{[\varphi]}\otimes_{A'}A_{[\varphi]}\to M$은 모든 텐서곱
$z\otimes a$($z\in M$, $a\in A$)에 $a\cdot z$를 대응시키는 표준
준동형이다. 이는 정의
(\hyperref[0.3.7.1-ko]{0, 3.7.1},
\hyperref[0.4.4.3-ko]{0, 4.4.3} 및
\hyperref[I.1.3.7-ko]{I, 1.3.7})에서 곧바로 따른다.

따라서 (\hyperref[0.4.4.3-ko]{0, 4.4.3})과
(\hyperref[0.3.5.4.4-ko]{3.5.4.4})에 의해
$v:N'\to M_{[\varphi]}$가 $A'$-준동형이면
$\widetilde{v}^{\#}=(v\otimes1)^{\sim}$이다.
\end{env}

\begin{env}[1.6.8]
\label{I.1.6.8-ko}
$N'_1$, $N'_2$를 두 $A'$-가군이라 하고 $N'_1$이 \emph{유한 표시}를
갖는다고 하자. 그러면
(\hyperref[I.1.6.7-ko]{1.6.7})과
(\hyperref[I.1.3.12-ko]{1.3.12, (ii)})에 따라 표준 준동형
(\hyperref[0.4.4.6-ko]{0, 4.4.6})
\[
  \Phi^*\bigl(\mathscr{H}\!om_{\widetilde{A'}}
    (\widetilde{N'_1},\widetilde{N'_2})\bigr)
  \longrightarrow
  \mathscr{H}\!om_{\widetilde A}
    \bigl(\Phi^*(\widetilde{N'_1}),\Phi^*(\widetilde{N'_2})\bigr)
\]
은 $\widetilde{\gamma}$와 같다. 여기서 $\gamma$는 $A$-가군의
표준 준동형
\[
  \operatorname{Hom}_{A'}(N'_1,N'_2)\otimes_{A'}A
  \longrightarrow
  \operatorname{Hom}_A(N'_1\otimes_{A'}A,N'_2\otimes_{A'}A)
\]
이다.
\end{env}

\begin{env}[1.6.9]
\label{I.1.6.9-ko}
$\mathfrak{J}'$을 $A'$의 아이디얼, $M$을 $A$-가군이라 하자. 정의에
따라 $\widetilde{\mathfrak{J}'}\widetilde M$은 표준 준동형
$\Phi^*(\widetilde{\mathfrak{J}'})
\otimes_{\widetilde A}\widetilde M\to\widetilde M$의 상이다. 따라서
(\hyperref[I.1.6.5-ko]{1.6.5})와
(\hyperref[I.1.3.12-ko]{1.3.12, (i)})에 의해
$\widetilde{\mathfrak{J}'}\widetilde M$은
$(\mathfrak{J}'M)^{\sim}$와 표준적으로 동일시된다. 특히
$\Phi^*(\widetilde{\mathfrak{J}'})\widetilde A$는
$(\mathfrak{J}'A)^{\sim}$와 동일시된다. 또한 함자 $\Phi^*$의 오른쪽
완전성을 고려하면 $\widetilde A$-대수
$\Phi^*((A'/\mathfrak{J}')^{\sim})$는
$(A/\mathfrak{J}'A)^{\sim}$와 동일시된다.
\end{env}

\begin{env}[1.6.10]
\label{I.1.6.10-ko}
$A''$을 세 번째 환, $\varphi'$을 준동형 $A''\to A'$이라 하고
$\varphi''=\varphi\circ\varphi'$로 놓자. 정의에서 곧바로
\[
  {}^a\varphi''=({}^a\varphi')\circ({}^a\varphi),
  \qquad
  \widetilde{\varphi''}=\widetilde\varphi\circ\widetilde{\varphi'}
\]
를 얻는다(\hyperref[0.1.5.7-ko]{0, 1.5.7}). 따라서
$\Phi''=\Phi'\circ\Phi$이다. 다시 말해
$(\operatorname{Spec}A,\widetilde A)$는 $A$에 관하여 환의 범주에서
환 달린 공간의 범주로 가는 \emph{반변 함자}이다.
\end{env}

\subsection{아핀 스킴의 사상의 특징짓기}
\label{subsection:I.1.7-ko}

\begin{definition}[1.7.1]
\label{I.1.7.1-ko}
환 달린 공간 $(X,\mathscr{O}_X)$가 어떤 환 $A$에 대한
$(\operatorname{Spec}(A),\widetilde A)$ 꼴의 환 달린 공간과 동형이면
이를 \emph{아핀 스킴}이라 한다. 이때 $\Gamma(X,\mathscr{O}_X)$는
환 $A$와 표준적으로 동일시되며
(\hyperref[I.1.3.7-ko]{1.3.7}), 이를 아핀 스킴
$(X,\mathscr{O}_X)$의 환이라 한다. 혼동의 염려가 없을 때에는
$A(X)$로 쓴다.
\end{definition}

앞으로 \emph{아핀 스킴} $\operatorname{Spec}(A)$라고 하면 항상 환 달린
공간 $(\operatorname{Spec}(A),\widetilde A)$를 뜻한다.

\begin{env}[1.7.2]
\label{I.1.7.2-ko}
$A$, $B$를 두 환, $(X,\mathscr{O}_X)$와 $(Y,\mathscr{O}_Y)$를 각각
소 스펙트럼 $X=\operatorname{Spec}(A)$와 $Y=\operatorname{Spec}(B)$에
대응하는 아핀 스킴이라 하자. 앞서
(\hyperref[I.1.6.1-ko]{1.6.1})에서 모든 환 준동형
$\varphi:B\to A$에 사상
\[
  \Phi=({}^a\varphi,\widetilde\varphi)
  =\operatorname{Spec}(\varphi):
  (X,\mathscr{O}_X)\to(Y,\mathscr{O}_Y)
\]
이 대응함을 보았다. $\varphi$는 $\Phi$에 의해 완전히 결정된다. 실제로
정의상
\[
  \varphi=\Gamma(\widetilde\varphi):
  \Gamma(\widetilde B)\longrightarrow
  \Gamma({}^a\varphi_*(\widetilde A))=\Gamma(\widetilde A)
\]
이다.
\end{env}

\begin{theorem}[1.7.3]
\label{I.1.7.3-ko}
$(X,\mathscr{O}_X)$와 $(Y,\mathscr{O}_Y)$를 두 아핀 스킴이라 하자.
환 달린 공간의 사상
$(\psi,\theta):(X,\mathscr{O}_X)\to(Y,\mathscr{O}_Y)$가 어떤 환 준동형
$\varphi:A(Y)\to A(X)$에 대하여
$({}^a\varphi,\widetilde\varphi)$ 꼴이기 위한 필요충분조건은 모든
$x\in X$에 대하여
\[
  \theta_x^{\#}:\mathscr{O}_{\psi(x)}\to\mathscr{O}_x
\]
가 국소 준동형인 것이다.
\end{theorem}

\oldpage[I]{97}
$A=A(X)$, $B=A(Y)$로 놓자. 이 조건은 필요하다. 실제로
(\hyperref[I.1.6.1-ko]{1.6.1})에서
$\widetilde\varphi_x^{\#}$이 $\varphi$에서 표준적으로 유도되는 준동형
$B_{{}^a\varphi(x)}\to A_x$임을 보았다. 정의상
${}^a\varphi(x)=\varphi^{-1}(\mathfrak{j}_x)$이므로 이 준동형은 국소
준동형이다.

이제 조건이 충분함을 증명하자. 정의상 $\theta$는 준동형
$\mathscr{O}_Y\to\psi_*(\mathscr{O}_X)$이고, 여기서 표준적으로 환 준동형
\[
  \varphi=\Gamma(\theta):
  B=\Gamma(Y,\mathscr{O}_Y)
  \longrightarrow
  \Gamma(Y,\psi_*(\mathscr{O}_X))
  =\Gamma(X,\mathscr{O}_X)=A
\]
을 얻는다.

$\theta_x^{\#}$에 대한 가정으로부터 몫을 취하여 잉여체
$k(\psi(x))$에서 잉여체 $k(x)$로 가는 단사 준동형 $\theta^x$를 얻는다.
모든 절단 $f\in\Gamma(Y,\mathscr{O}_Y)=B$에 대하여
$\theta^x(f(\psi(x)))=\varphi(f)(x)$이다. 따라서 관계
$f(\psi(x))=0$과 $\varphi(f)(x)=0$은 동치이다. 이는
$\mathfrak{j}_{\psi(x)}=\mathfrak{j}_{{}^a\varphi(x)}$라는 뜻이므로 모든
$x\in X$에 대하여 $\psi(x)={}^a\varphi(x)$, 곧
$\psi={}^a\varphi$이다. 또한 그림
\phantomsection\label{I.1.7.3.diagram-ko}
\[
  \xymatrix{
    B=\Gamma(Y,\mathscr{O}_Y)\ar[r]^\varphi\ar[d] &
    \Gamma(X,\mathscr{O}_X)=A\ar[d]\\
    B_{\psi(x)}\ar[r]_{\theta_x^{\#}} & A_x
  }
\]
은 가환이다(\hyperref[0.3.7.2-ko]{0, 3.7.2}). 이는
$\theta_x^{\#}$이 $\varphi$에서 표준적으로 유도되는 준동형
$\varphi_x:B_{\psi(x)}\to A_x$와 같다는 뜻이다
(\hyperref[0.1.5.1-ko]{0, 1.5.1}). $\theta_x^{\#}$들의 자료가
$\theta^{\#}$, 따라서 $\theta$도 완전히 결정하므로
(\hyperref[0.3.7.1-ko]{0, 3.7.1}),
$\widetilde\varphi$의 정의
(\hyperref[I.1.6.1-ko]{1.6.1})에 의해
$\theta=\widetilde\varphi$이다.

(\hyperref[I.1.7.3-ko]{1.7.3})의 조건을 만족하는 환 달린 공간의 사상
$(\psi,\theta)$를 \emph{아핀 스킴의 사상}이라 하겠다.

\begin{corollary}[1.7.4]
\label{I.1.7.4-ko}
$(X,\mathscr{O}_X)$와 $(Y,\mathscr{O}_Y)$가 두 아핀 스킴이면 아핀
스킴의 사상들의 집합과
$\operatorname{Hom}_{\mathrm{Ring}}(B,A)$ 사이에 표준 전단사가
존재한다. 여기서 $A=\Gamma(\mathscr{O}_X)$이고
$B=\Gamma(\mathscr{O}_Y)$이다.
\end{corollary}

달리 말하면, $A$에 관한 $(\operatorname{Spec}(A),\widetilde A)$와
$(X,\mathscr{O}_X)$에 관한 $\Gamma(X,\mathscr{O}_X)$는 가환환의 범주와
아핀 스킴의 범주의 반대 범주 사이의 \emph{동치}를 정의한다
(T, I, 1.2).

\begin{corollary}[1.7.5]
\label{I.1.7.5-ko}
$\varphi:B\to A$가 전사이면 이에 대응하는 사상
$({}^a\varphi,\widetilde\varphi)$는 환 달린 공간의 단사사상이다
\emph{(4.1.7 참조)}.
\end{corollary}

실제로 ${}^a\varphi$는 단사이다
(\hyperref[I.1.2.5-ko]{1.2.5}). 또한 $\varphi$가 전사이므로 모든
$x\in X$에 대하여 $\varphi$에서 분수환을 취하여 얻는
$\varphi_x^{\#}:B_{{}^a\varphi(x)}\to A_x$도 전사이다
(\hyperref[0.1.5.1-ko]{0, 1.5.1}). 따라서 결론은
(\hyperref[0.4.1.1-ko]{0, 4.1.1})에서 따른다.

% Authority: EGA I ega1-2-fr.tex, whole-file SHA-256 AE6B128092ACBB8C1AFB4899EEA003FB966B6FF6669A264B59FD5F095AF4F029.
% Sequential coverage in this file: authority lines 1-296, complete subsections 2.1-2.2.
\section{준스킴과 준스킴의 사상}
\label{section:I.2-ko}

\subsection{준스킴의 정의}
\label{subsection:I.2.1-ko}

\begin{env}[2.1.1]
\label{I.2.1.1-ko}
환 달린 공간 $(X,\mathscr{O}_X)$가 주어졌을 때, $X$의 열린 부분집합
$V$에 유도된 환 달린 공간 $(V,\mathscr{O}_X|V)$가 아핀 스킴이면
$V$를 \emph{아핀 열린집합}이라 한다
(\hyperref[I.1.7.1-ko]{1.7.1}).
\end{env}

\begin{definition}[2.1.2]
\label{I.2.1.2-ko}
모든 점이 아핀 열린 근방을 갖는 환 달린 공간 $(X,\mathscr{O}_X)$를
\emph{준스킴}이라 한다.
\end{definition}

\oldpage[I]{98}

\begin{proposition}[2.1.3]
\label{I.2.1.3-ko}
$(X,\mathscr{O}_X)$가 준스킴이면 아핀 열린집합들은 $X$의 위상의
기저를 이룬다.
\end{proposition}

실제로 $V$가 $x\in X$의 임의의 열린 근방이면, 가정에 따라 $x$의
어떤 열린 근방 $W$에 대하여 $(W,\mathscr{O}_X|W)$가 아핀 스킴이다.
그 환을 $A$라 하자. 공간 $W$ 안에서 $V\cap W$는 $x$의 열린
근방이므로, $x$의 열린 근방이고 $V\cap W$에 포함되는 $D(f)$가
존재하도록 하는 $f\in A$가 있다
(\hyperref[I.1.1.10-ko]{1.1.10 (i)}). 이때 환 달린 공간
$(D(f),\mathscr{O}_X|D(f))$는 $A_f$와 동형인 환을 갖는 아핀
스킴이므로(\hyperref[I.1.3.6-ko]{1.3.6}), 명제가 따른다.

\begin{proposition}[2.1.4]
\label{I.2.1.4-ko}
준스킴의 바탕공간은 콜모고로프 공간이다.
\end{proposition}

실제로 준스킴 $X$의 서로 다른 두 점 $x,y$가 하나의 아핀
열린집합에 함께 들어 있지 않으면, 두 점 가운데 하나를 포함하고
다른 하나를 포함하지 않는 열린 근방이 존재함은 자명하다. 두 점이
하나의 아핀 열린집합에 함께 들어 있으면
(\hyperref[I.1.1.8-ko]{1.1.8})에서 결론이 따른다.

\begin{proposition}[2.1.5]
\label{I.2.1.5-ko}
$(X,\mathscr{O}_X)$가 준스킴이면 $X$의 모든 기약 닫힌 부분집합은
유일한 일반점을 갖는다. 따라서 사상
\[
  x\longmapsto\overline{\{x\}}
\]
은 $X$에서 그 기약 닫힌 부분집합들의 집합으로 가는 전단사이다.
\end{proposition}

실제로 $Y$를 $X$의 기약 닫힌 부분집합, $y\in Y$, $U$를 $X$에서
$y$의 아핀 열린 근방이라 하자. $U\cap Y$는 $Y$에서 조밀하고
기약이다
(\hyperref[0.2.1.1-ko]{0, 2.1.1} 및
\hyperref[0.2.1.4-ko]{2.1.4}). 따라서
(\hyperref[I.1.1.14-ko]{1.1.14})에 의해 $U\cap Y$는 어떤 점 $x$의
$U$ 안에서의 폐포이다. 그러므로 $Y\subset\overline U$는 $X$에서
$x$의 폐포이다. $Y$의 일반점의 유일성은
(\hyperref[I.2.1.4-ko]{2.1.4})와
(\hyperref[0.2.1.3-ko]{0, 2.1.3})에서 따른다.

\begin{env}[2.1.6]
\label{I.2.1.6-ko}
$Y$가 $X$의 기약 닫힌 부분집합이고 $y$가 그 일반점이면, 국소환
$\mathscr{O}_y$를 $\mathscr{O}_{X/Y}$라고도 쓰며 이를
\emph{$Y$를 따라 취한 $X$의 국소환}, 또는 \emph{$X$ 안의 $Y$의
국소환}이라 한다.

$X$ 자체가 기약이고 $x$가 그 일반점이면, $\mathscr{O}_x$를
\emph{$X$ 위의 유리함수환}이라고도 한다(§~7 참조).
\end{env}

\begin{proposition}[2.1.7]
\label{I.2.1.7-ko}
$(X,\mathscr{O}_X)$가 준스킴이면 $X$의 모든 열린 부분집합 $U$에
대하여 환 달린 공간 $(U,\mathscr{O}_X|U)$도 준스킴이다.
\end{proposition}

이는 정의 (\hyperref[I.2.1.2-ko]{2.1.2})와 명제
(\hyperref[I.2.1.3-ko]{2.1.3})에서 곧바로 따른다.

$(U,\mathscr{O}_X|U)$를 $(X,\mathscr{O}_X)$가 $U$ 위에
\emph{유도하는} 준스킴, 또는 $(X,\mathscr{O}_X)$의 $U$에 대한
\emph{제한}이라 한다.

\begin{env}[2.1.8]
\label{I.2.1.8-ko}
준스킴 $(X,\mathscr{O}_X)$의 바탕공간 $X$가 기약(각각 연결)이면
이 준스킴을 \emph{기약}(각각 \emph{연결})이라 한다. 준스킴이
\emph{기약이고 축소}이면 이를 \emph{정역 준스킴}이라 한다
\emph{(5.1.4 참조)}. 준스킴 $(X,\mathscr{O}_X)$의 모든 점
$x\in X$가, 그 위에 유도된 준스킴이 정역 준스킴인 어떤 열린 근방
$U$를 가지면 $(X,\mathscr{O}_X)$를 \emph{국소 정역 준스킴}이라
한다.
\end{env}

\subsection{준스킴의 사상}
\label{subsection:I.2.2-ko}
\label{I.2.2-ko}

\begin{definition}[2.2.1]
\label{I.2.2.1-ko}
두 준스킴 $(X,\mathscr{O}_X)$, $(Y,\mathscr{O}_Y)$가 주어졌다고 하자.
$(X,\mathscr{O}_X)$에서 $(Y,\mathscr{O}_Y)$로 가는 환 달린 공간의
사상 $(\psi,\theta)$ 가운데, 모든 $x\in X$에 대하여
\[
  \theta_x^\sharp:\mathscr{O}_{\psi(x)}\longrightarrow\mathscr{O}_x
\]
가 국소 준동형인 것을 \emph{준스킴의 사상}이라 한다.
\end{definition}

\oldpage[I]{99}

따라서 몫을 취하면
$\theta_x^\sharp:\mathscr{O}_{\psi(x)}\to\mathscr{O}_x$는 단사 준동형
$\theta^x:k(\psi(x))\to k(x)$를 유도하며, 이에 따라 $k(x)$를 체
$k(\psi(x))$의 \emph{확대체}로 볼 수 있다.

\begin{env}[2.2.2]
\label{I.2.2.2-ko}
준스킴의 두 사상 $(\psi,\theta)$와 $(\psi',\theta')$의 합성
$(\psi'',\theta'')$도 준스킴의 사상이다. 이는 공식
${\theta''}^\sharp=\theta^\sharp\circ\psi^*({\theta'}^\sharp)$
(\hyperref[0.3.5.5-ko]{0, 3.5.5})에서 따른다. 따라서 준스킴들은 하나의
\emph{범주}를 이룬다. 일반 표기에 따라 준스킴 $X$에서 준스킴 $Y$로
가는 사상들의 집합을 $\operatorname{Hom}(X,Y)$로 쓴다.
\end{env}

\begin{example}[2.2.3]
\label{I.2.2.3-ko}
$U$가 $X$의 열린 부분집합이면, 유도된 준스킴
$(U,\mathscr{O}_X|U)$에서 $(X,\mathscr{O}_X)$로 가는 표준 단사사상
(\hyperref[0.4.1.2-ko]{0, 4.1.2})은 준스킴의 사상이다. 더구나 이는
환 달린 공간의 단사사상이며(따라서 특히 준스킴의 단사사상이며),
이는 (\hyperref[0.4.1.1-ko]{0, 4.1.1})에서 곧바로 따른다.
\end{example}

\begin{proposition}[2.2.4]
\label{I.2.2.4-ko}
$(X,\mathscr{O}_X)$를 준스킴, $(S,\mathscr{O}_S)$를 환 $A$의 아핀
스킴이라 하자. 준스킴 $(X,\mathscr{O}_X)$에서 준스킴
$(S,\mathscr{O}_S)$로 가는 사상들과 $A$에서 환
$\Gamma(X,\mathscr{O}_X)$로 가는 준동형들 사이에는 표준 전단사 대응이
존재한다.
\end{proposition}

먼저 $(X,\mathscr{O}_X)$와 $(Y,\mathscr{O}_Y)$가 임의의 두 환 달린
공간이면, $(X,\mathscr{O}_X)$에서 $(Y,\mathscr{O}_Y)$로 가는 사상
$(\psi,\theta)$는 표준적으로 환 준동형
\[
  \Gamma(\theta):\Gamma(Y,\mathscr{O}_Y)\longrightarrow
  \Gamma(Y,\psi_*(\mathscr{O}_X))=\Gamma(X,\mathscr{O}_X)
\]
을 정의함을 주목하자. 지금의 경우에는 임의의 준동형
$\varphi:A\to\Gamma(X,\mathscr{O}_X)$가 유일한 $\theta$에 대하여
$\Gamma(\theta)$ 꼴임을 보이면 충분하다. 가정에 따라 $X$는 아핀
열린집합들 $(V_\alpha)$로 덮인다. $\varphi$를 제한 준동형
\[
  \Gamma(X,\mathscr{O}_X)\longrightarrow
  \Gamma(V_\alpha,\mathscr{O}_X|V_\alpha)
\]
과 합성하면 준동형
$\varphi_\alpha:A\to\Gamma(V_\alpha,\mathscr{O}_X|V_\alpha)$를 얻는다.
(\hyperref[I.1.7.3-ko]{1.7.3})에 따라 이는 준스킴
$(V_\alpha,\mathscr{O}_X|V_\alpha)$에서 $(S,\mathscr{O}_S)$로 가는
유일한 사상 $(\psi_\alpha,\theta_\alpha)$에 대응한다. 또한 모든 지표쌍
$(\alpha,\beta)$와 모든 점 $x\in V_\alpha\cap V_\beta$에 대하여,
$x$를 포함하고 $V_\alpha\cap V_\beta$에 포함되는 아핀 열린집합 $W$가
존재한다(\hyperref[I.2.1.3-ko]{2.1.3}). $\varphi_\alpha$와
$\varphi_\beta$를 $W$로 가는 제한 준동형과 합성하면 같은 준동형
\[
  \Gamma(S,\mathscr{O}_S)\longrightarrow
  \Gamma(W,\mathscr{O}_X|W)
\]
을 얻는다. 따라서 모든 $x\in V_\alpha$와 모든 $\alpha$에 대한 관계
$(\theta_\alpha^\sharp)_x=(\varphi_\alpha)_x$
(\hyperref[I.1.6.1-ko]{1.6.1})와 함께 보면,
$(\psi_\alpha,\theta_\alpha)$와 $(\psi_\beta,\theta_\beta)$의 $W$로의
제한은 일치한다. 그러므로 각 $V_\alpha$로의 제한이
$(\psi_\alpha,\theta_\alpha)$인 환 달린 공간의 사상
\[
  (\psi,\theta):(X,\mathscr{O}_X)\longrightarrow(S,\mathscr{O}_S)
\]
이 유일하게 존재한다. 이 사상이 준스킴의 사상이며
$\Gamma(\theta)=\varphi$임은 명백하다.

$u:A\to\Gamma(X,\mathscr{O}_X)$를 환 준동형이라 하고,
$v=(\psi,\theta)$를 이에 대응하는 사상
$(X,\mathscr{O}_X)\to(S,\mathscr{O}_S)$라 하자. 모든 $f\in A$에 대하여
\begin{equation}
  \psi^{-1}(D(f))=X_{u(f)}
  \tag{2.2.4.1}\label{I.2.2.4.1-ko}
\end{equation}
이다. 여기서 표기는 국소 자유층 $\mathscr{O}_X$에 관한
(\hyperref[0.5.5.2-ko]{0, 5.5.2})의 것을 사용한다. 실제로 이 공식을
$X$ 자체가 아핀인 경우에 확인하면 충분하며, 그때에는 바로
(\hyperref[I.1.2.2.2-ko]{1.2.2.2})이다.

\begin{proposition}[2.2.5]
\label{I.2.2.5-ko}
(\hyperref[I.2.2.4-ko]{2.2.4})의 가정 아래에서, 환 준동형
$\varphi:A\to\Gamma(X,\mathscr{O}_X)$와 이에 대응하는 준스킴의 사상
$f:(X,\mathscr{O}_X)\to(S,\mathscr{O}_S)$를 생각하자.
$\mathscr{G}$를 준연접 $\mathscr{O}_X$-가군,
$\mathscr{F}$를 준연접 $\mathscr{O}_S$-가군이라 하고
$M=\Gamma(S,\mathscr{F})$로 놓자.
\oldpage[I]{100}
그러면 $f$-사상
$\mathscr{F}\to\mathscr{G}$
(\hyperref[0.4.4.1-ko]{0, 4.4.1})들과 $A$-가군 준동형
\[
  M\longrightarrow(\Gamma(X,\mathscr{G}))_{[\varphi]}
\]
들 사이에는 표준 전단사 대응이 존재한다.
\end{proposition}

실제로 (\hyperref[I.2.2.4-ko]{2.2.4})에서와 같이 논하면 곧 $X$가
아핀인 경우로 환원되며, 그 경우 명제는
(\hyperref[I.1.6.3-ko]{1.6.3})과
(\hyperref[I.1.3.8-ko]{1.3.8})에서 따른다.

\begin{env}[2.2.6]
\label{I.2.2.6-ko}
준스킴의 사상
$(\psi,\theta):(X,\mathscr{O}_X)\to(Y,\mathscr{O}_Y)$가 $X$의 모든
열린 부분집합 $U$에 대하여 $\psi(U)$가 $Y$에서 열려 있으면 이 사상을
\emph{열린 사상}이라 하고, $X$의 모든 닫힌 부분집합 $F$에 대하여
$\psi(F)$가 $Y$에서 닫혀 있으면 \emph{닫힌 사상}이라 한다.
$\psi(X)$가 $Y$에서 조밀하면 \emph{우세 사상}, $\psi$가 전사이면
\emph{전사 사상}이라 한다. 이 조건들은 연속사상 $\psi$에만 관련됨을
주목하자.
\end{env}

\begin{proposition}[2.2.7]
\label{I.2.2.7-ko}
준스킴의 두 사상
\[
  f=(\psi,\theta):(X,\mathscr{O}_X)\to(Y,\mathscr{O}_Y),
  \qquad
  g=(\psi',\theta'):(Y,\mathscr{O}_Y)\to(Z,\mathscr{O}_Z)
\]
를 생각하자.
\begin{enumerate}
  \item[\rm(i)] $f$와 $g$가 모두 열린(각각 닫힌, 우세, 전사) 사상이면
    $g\circ f$도 그러하다.
  \item[\rm(ii)] $f$가 전사이고 $g\circ f$가 닫힌 사상이면 $g$도 닫힌
    사상이다.
  \item[\rm(iii)] $g\circ f$가 전사이면 $g$도 전사이다.
\end{enumerate}
\end{proposition}

(i)와 (iii)은 명백하다. $g\circ f=(\psi'',\theta'')$로 놓자. $F$가
$Y$에서 닫혀 있으면 $\psi^{-1}(F)$는 $X$에서 닫혀 있으므로
$\psi''(\psi^{-1}(F))$는 $Z$에서 닫혀 있다. 그런데 $\psi$가 전사이므로
$\psi(\psi^{-1}(F))=F$이고, 따라서
$\psi''(\psi^{-1}(F))=\psi'(F)$이다. 이로써 (ii)가 증명된다.

\begin{proposition}[2.2.8]
\label{I.2.2.8-ko}
$f=(\psi,\theta)$를
$(X,\mathscr{O}_X)\to(Y,\mathscr{O}_Y)$인 사상이라 하고,
$(U_\alpha)$를 $Y$의 열린 덮개라 하자. $f$가 열린(각각 닫힌, 전사,
우세) 사상이기 위한 필요충분조건은, 각각의 유도된 준스킴
$(\psi^{-1}(U_\alpha),\mathscr{O}_X|\psi^{-1}(U_\alpha))$에 대한 $f$의
제한을 이 준스킴에서 유도된 준스킴
$(U_\alpha,\mathscr{O}_Y|U_\alpha)$로 가는 사상으로 보았을 때, 그
제한이 열린(각각 닫힌, 전사, 우세) 사상인 것이다.
\end{proposition}

이 명제는 정의에서 곧바로 따른다. 여기서는 $Y$의 부분집합 $F$가
$Y$에서 닫힌(각각 열린, 조밀한) 부분집합이기 위한 필요충분조건이
모든 $\alpha$에 대하여 $F\cap U_\alpha$가 $U_\alpha$에서 닫힌(각각
열린, 조밀한) 부분집합인 것임을 사용한다.

\begin{env}[2.2.9]
\label{I.2.2.9-ko}
$(X,\mathscr{O}_X)$와 $(Y,\mathscr{O}_Y)$를 두 준스킴이라 하고, $X$와
$Y$가 같은 유한 개수의 기약 성분 $X_i$와 $Y_i$
$(1\leq i\leq n)$를 갖는다고 하자. $X_i$와 $Y_i$의 일반점을 각각
$\xi_i$와 $\eta_i$라 하자
(\hyperref[I.2.1.5-ko]{2.1.5}). 사상
\[
  f=(\psi,\theta):(X,\mathscr{O}_X)\longrightarrow(Y,\mathscr{O}_Y)
\]
이 모든 $i$에 대하여 $\psi^{-1}(\eta_i)=\{\xi_i\}$를 만족하고
\[
  \theta_{\xi_i}^{\sharp}:\mathscr{O}_{\eta_i}\longrightarrow
  \mathscr{O}_{\xi_i}
\]
가 동형이면 $f$를 \emph{쌍유리 사상}이라 한다. 쌍유리 사상이 우세
사상임은 명백하다(\hyperref[0.2.1.8-ko]{0, 2.1.8}). 따라서 닫힌
사상이기도 하면 전사이다.
\end{env}

\begin{notation}[2.2.10]
\label{I.2.2.10-ko}
이 책의 앞으로의 모든 부분에서는, 혼동의 염려가 없을 때 준스킴의
표기에서는 구조층을, 사상의 표기에서는 구조층의 준동형을
\emph{생략}한다. 준스킴 $X$의 바탕공간의 열린 부분집합 $U$를
준스킴이라고 할 때에는 언제나 $U$ 위에 유도된 준스킴을 뜻한다.
\end{notation}

\subsection{준스킴의 붙이기}
\label{subsection:I.2.3-ko}

\begin{env}[2.3.1]
\label{I.2.3.1-ko}
\oldpage[I]{101}
정의 (\hyperref[I.2.1.2-ko]{2.1.2})에 따라, 준스킴들을
\emph{붙여서} 얻은 모든 환 달린 공간
(\hyperref[0.4.1.6-ko]{0, 4.1.6})은 다시 준스킴이다. 특히 모든
준스킴은 정의상 아핀 열린집합들로 이루어진 덮개를 가지므로, 모든
준스킴은 \emph{아핀 스킴들을 붙여서} 얻을 수 있다.
\end{env}

\begin{example}[2.3.2]
\label{I.2.3.2-ko}
$K$를 체라 하고, $B=K[s]$, $C=K[t]$를 각각 $K$ 위의 하나의 부정원에
대한 다항식환이라 하자. 또한 $X_1=\operatorname{Spec}(B)$,
$X_2=\operatorname{Spec}(C)$라 하자. 이들은 서로 동형인 두 아핀
스킴이다. $X_1$에서(각각 $X_2$에서) $U_{12}$를(각각 $U_{21}$을)
아핀 열린집합 $D(s)$로(각각 $D(t)$로) 잡자. 그 환 $B_s$는(각각
$C_t$는) $f\in B$인(각각 $g\in C$인) 유리분수 $f(s)/s^m$의
꼴로(각각 $g(t)/t^n$의 꼴로) 이루어진다. $B_s$에서 $C_t$로 가며
$f(s)/s^m$을 유리분수 $f(1/t)/(1/t^m)$에 대응시키는 동형에
(\hyperref[I.2.2.4-ko]{2.2.4})에 따라 대응하는 준스킴의 동형
\[
  u_{12}:U_{21}\longrightarrow U_{12}
\]
를 생각하자. 붙이기 조건이 전혀 없음은 명백하므로, $u_{12}$를
사용하여 $U_{12}$와 $U_{21}$을 따라 $X_1$과 $X_2$를 붙일 수 있다.
이렇게 얻은 준스킴 $X$는 뒤에서 일반적인 구성법의 특수한 경우로
다시 만나게 된다(II, 2.4.3). 여기서는 $X$가 \emph{아핀 스킴이
아님}만을 보이자. 실제로 환 $\Gamma(X,\mathscr{O}_X)$는 $K$와
\emph{동형}이므로 그 스펙트럼은 한 점으로 이루어진다. $X$ 위의
$\mathscr{O}_X$의 한 단면을 $X$의 아핀 열린집합과 동일시되는
$X_1$에(각각 $X_2$에) 제한하면 다항식 $f(s)$를(각각 $g(t)$를)
얻는다. 정의에 따라 $g(t)=f(1/t)$이어야 하며, 이는
$f=g\in K$일 때에만 가능하다.
\end{example}

\subsection{국소 스킴}
\label{subsection:I.2.4-ko}

\begin{env}[2.4.1]
\label{I.2.4.1-ko}
환 $A$가 국소환인 아핀 스킴을 \emph{국소 스킴}이라 한다. 그러면
$X=\operatorname{Spec}(A)$에는 유일한 \emph{닫힌점} $a$가 존재하며,
다른 모든 점 $b\in X$에 대하여
$a\in\overline{\{b\}}$이다
(\hyperref[I.1.1.7-ko]{1.1.7}).
\end{env}

준스킴 $Y$의 점 $y$마다 국소 스킴
$\operatorname{Spec}(\mathscr{O}_y)$를 \emph{$Y$의 점 $y$에서의 국소
스킴}이라 한다. $V$를 $y$를 포함하는 $Y$의 아핀 열린집합, $B$를 아핀
스킴 $V$의 환이라 하자. $\mathscr{O}_y$는 $B_y$와 표준적으로
동일시되고(\hyperref[I.1.3.4-ko]{1.3.4}), 따라서 표준 준동형
$B\to B_y$는 (\hyperref[I.1.6.1-ko]{1.6.1})에 따라 준스킴의 사상
$\operatorname{Spec}(\mathscr{O}_y)\to V$에 대응한다. 이 사상을 표준
단사사상 $V\to Y$와 합성하면 사상
$\operatorname{Spec}(\mathscr{O}_y)\to Y$를 얻으며, 이는 선택한
$y$의 아핀 열린 근방 $V$에 \emph{무관하다}. 실제로 $V'$이 $y$를
포함하는 또 다른 아핀 열린집합이면, $y$를 포함하고
$W\subset V\cap V'$인 세 번째 아핀 열린집합 $W$가 존재한다
(\hyperref[I.2.1.3-ko]{2.1.3}). 따라서 $V\subset V'$인 경우만
생각하면 충분하다. $B'$을 $V'$의 환이라 하면, 이는 도식
\phantomsection\label{I.2.4.1.diagram-ko}
\[
  \xymatrix{
    B'\ar[rr]\ar[dr] & & B\ar[dl]\\
    & \mathscr{O}_y
  }
\]
이 가환임을 확인하는 것으로 귀착된다
(\hyperref[0.1.5.1-ko]{0, 1.5.1}). 이렇게 정의한 사상
\[
  \operatorname{Spec}(\mathscr{O}_y)\longrightarrow Y
\]
을 \emph{표준 사상}이라 한다.

\begin{proposition}[2.4.2]
\label{I.2.4.2-ko}
\oldpage[I]{102}
$(Y,\mathscr{O}_Y)$를 준스킴이라 하자. 각 $y\in Y$에 대하여
$(\psi,\theta)$를 표준 사상
\[
  (\operatorname{Spec}(\mathscr{O}_y),\widetilde{\mathscr{O}}_y)
  \longrightarrow(Y,\mathscr{O}_Y)
\]
이라 하자. 그러면 $\psi$는 $\operatorname{Spec}(\mathscr{O}_y)$에서
$y\in\overline{\{z\}}$를 만족하는 점 $z$들로 이루어진 $Y$의 부분공간
$S_y$ 위로 가는 위상동형이다. 다시 말해 이 점들은 $y$의
\emph{일반화들}이다(\hyperref[0.2.1.2-ko]{0, 2.1.2}). 또한
$z=\psi(\mathfrak{p})$이면
\[
  \theta_z^{\sharp}:\mathscr{O}_z\longrightarrow
  (\mathscr{O}_y)_{\mathfrak{p}}
\]
는 동형이다. 따라서 $(\psi,\theta)$는 환 달린 공간의 단사사상이다.
\end{proposition}

$\operatorname{Spec}(\mathscr{O}_y)$의 유일한 닫힌점 $a$는 이 공간의
모든 점의 폐포에 속하고 $\psi(a)=y$이므로, 연속사상 $\psi$에 의한
$\operatorname{Spec}(\mathscr{O}_y)$의 상은 $S_y$에 포함된다. 한편
$S_y$는 $y$를 포함하는 모든 아핀 열린집합에 포함되므로, $Y$가 아핀
스킴인 경우로 환원할 수 있다. 이 경우 명제는
(\hyperref[I.1.6.2-ko]{1.6.2})에서 따른다.

\emph{따라서 (\hyperref[I.2.1.5-ko]{2.1.5})에 의해
$\operatorname{Spec}(\mathscr{O}_y)$와 $y$를 포함하는 $Y$의 기약
닫힌 부분집합들의 집합 사이에는 전단사 대응이 있다.}

\begin{corollary}[2.4.3]
\label{I.2.4.3-ko}
$y\in Y$가 $Y$의 한 기약 성분의 일반점이기 위한 필요충분조건은,
국소환 $\mathscr{O}_y$의 유일한 소 아이디얼이 그 극대 아이디얼인
것이다. 다시 말해 $\mathscr{O}_y$의 차원이 $0$인 것이다.
\end{corollary}

\begin{proposition}[2.4.4]
\label{I.2.4.4-ko}
$(X,\mathscr{O}_X)$를 환이 $A$인 국소 스킴, $a$를 그 유일한 닫힌점,
$(Y,\mathscr{O}_Y)$를 준스킴이라 하자. 모든 사상
$u=(\psi,\theta):(X,\mathscr{O}_X)\to(Y,\mathscr{O}_Y)$는 유일하게
\[
  X\longrightarrow\operatorname{Spec}(\mathscr{O}_{\psi(a)})
  \longrightarrow Y
\]
로 분해된다. 여기서 두 번째 화살표는 표준 사상이고, 첫 번째 화살표는
국소 준동형 $\mathscr{O}_{\psi(a)}\to A$에 대응한다. 이로써 사상들의
집합
$\operatorname{Hom}((X,\mathscr{O}_X),(Y,\mathscr{O}_Y))$와 국소
준동형 $\mathscr{O}_y\to A$들의 집합($y\in Y$) 사이에 표준 전단사
대응이 주어진다.
\end{proposition}

실제로 모든 $x\in X$에 대하여 $a\in\overline{\{x\}}$이므로
$\psi(a)\in\overline{\{\psi(x)\}}$이다. 따라서 $\psi(X)$는
$\psi(a)$를 포함하는 모든 아핀 열린집합에 포함된다. 그러므로
$(Y,\mathscr{O}_Y)$가 환 $B$의 아핀 스킴인 경우로 환원할 수 있다.
이때 $\varphi\in\operatorname{Hom}(B,A)$에 대하여
$u=({}^a\varphi,\widetilde{\varphi})$이다
(\hyperref[I.1.7.3-ko]{1.7.3}). 또한
$\varphi^{-1}(\mathfrak{j}_a)=\mathfrak{j}_{\psi(a)}$이므로,
$B-\mathfrak{j}_{\psi(a)}$의 모든 원소의 $\varphi$에 의한 상은
국소환 $A$에서 가역이다. 따라서 명제의 분해는 분수환의 보편 성질
(\hyperref[0.1.2.4-ko]{0, 1.2.4})에서 따른다. 역으로, 모든 국소
준동형 $\mathscr{O}_y\to A$에는 $\psi(a)=y$인 유일한 사상
$(\psi,\theta):X\to\operatorname{Spec}(\mathscr{O}_y)$가 대응한다
(\hyperref[I.1.7.3-ko]{1.7.3}). 이를 표준 사상
$\operatorname{Spec}(\mathscr{O}_y)\to Y$와 합성하면 사상 $X\to Y$를
얻으며, 이로써 명제가 증명된다.

\begin{env}[2.4.5]
\label{I.2.4.5-ko}
환이 \emph{체} $K$인 아핀 스킴의 바탕공간은 한 점으로 이루어진다. $A$가
극대 아이디얼 $\mathfrak{m}$을 갖는 국소환이면, 모든 국소 준동형
$A\to K$의 핵은 $\mathfrak{m}$과 같으므로
\[
  A\longrightarrow A/\mathfrak{m}\longrightarrow K
\]
로 분해된다. 여기서 두 번째 화살표는 단사 준동형이다. 따라서 사상
$\operatorname{Spec}(K)\to\operatorname{Spec}(A)$들은 체의 단사
준동형 $A/\mathfrak{m}\to K$들과 전단사 대응한다.
\end{env}

$(Y,\mathscr{O}_Y)$를 준스킴이라 하자. 모든 $y\in Y$와
$\mathscr{O}_y$의 모든 아이디얼 $\mathfrak{a}_y$에 대하여 표준 준동형
$\mathscr{O}_y\to\mathscr{O}_y/\mathfrak{a}_y$는 사상
$\operatorname{Spec}(\mathscr{O}_y/\mathfrak{a}_y)
\to\operatorname{Spec}(\mathscr{O}_y)$를 정의한다. 이를 표준 사상
$\operatorname{Spec}(\mathscr{O}_y)\to Y$와 합성하면 사상
$\operatorname{Spec}(\mathscr{O}_y/\mathfrak{a}_y)\to Y$를 얻으며,
이 사상도 \emph{표준 사상}이라 한다. $\mathfrak{a}_y=\mathfrak{m}_y$,
즉 $\mathscr{O}_y/\mathfrak{a}_y=k(y)$인 경우에는 명제
(\hyperref[I.2.4.4-ko]{2.4.4})로부터 다음을 얻는다.

\begin{corollary}[2.4.6]
\label{I.2.4.6-ko}
\oldpage[I]{103}
$(X,\mathscr{O}_X)$를 그 환이 체 $K$인 국소 스킴, $\xi$를 $X$의
유일한 점, $(Y,\mathscr{O}_Y)$를 준스킴이라 하자. 모든 사상
$u=(\psi,\theta):(X,\mathscr{O}_X)\to(Y,\mathscr{O}_Y)$는 유일하게
\[
  X\longrightarrow\operatorname{Spec}(k(\psi(\xi)))
  \longrightarrow Y
\]
로 분해된다. 여기서 두 번째 화살표는 표준 사상이고, 첫 번째 화살표는
단사 준동형 $k(\psi(\xi))\to K$에 대응한다. 이로써 사상들의 집합
$\operatorname{Hom}((X,\mathscr{O}_X),(Y,\mathscr{O}_Y))$와 체의 단사
준동형 $k(y)\to K$들의 집합($y\in Y$) 사이에 표준 전단사 대응이
주어진다.
\end{corollary}

\begin{corollary}[2.4.7]
\label{I.2.4.7-ko}
모든 $y\in Y$에 대하여, 모든 표준 사상
$\operatorname{Spec}(\mathscr{O}_y/\mathfrak{a}_y)\to Y$는 환 달린
공간의 단사사상이다.
\end{corollary}

$\mathfrak{a}_y=0$인 경우에는 이미 이를 보았으며
(\hyperref[I.2.4.2-ko]{2.4.2}), 일반적인 경우에는
(\hyperref[I.1.7.5-ko]{1.7.5})를 적용하면 충분하다.

\begin{remark}[2.4.8]
\label{I.2.4.8-ko}
$X$를 국소 스킴, $a$를 그 유일한 닫힌점이라 하자. $a$를 포함하는
모든 아핀 열린집합은 반드시 $X$ 전체이므로, 모든 \emph{가역}
$\mathscr{O}_X$-가군(\hyperref[0.5.4.1-ko]{0, 5.4.1})은 반드시
\emph{$\mathscr{O}_X$와 동형}이다. 다시
말해 이 가군은 \emph{자명하다}. 이 성질은 임의의 아핀 스킴
$\operatorname{Spec}(A)$에 대해서는 일반적으로 성립하지 않는다.
$A$가 정규환일 때에는, $A$가 \emph{유일인수분해정역}이면 이 성질이
성립함을 제V장에서 보게 된다.
\end{remark}

\subsection{준스킴 위의 준스킴}
\label{subsection:I.2.5-ko}

\begin{definition}[2.5.1]
\label{I.2.5.1-ko}
준스킴 $S$가 주어졌을 때, 준스킴 $X$와 준스킴의 사상
$\varphi:X\to S$를 주는 것을
\emph{준스킴 $S$ 위의} 준스킴 $X$, 또는 \emph{$S$-준스킴}을
정의한다고 한다. 이때 $S$를 $S$-준스킴 $X$의 \emph{밑준스킴}이라
한다. 사상 $\varphi$를 $S$-준스킴 $X$의 \emph{구조 사상}이라 한다.
$S$가 환 $A$의 아핀 스킴이면, $\varphi$를 갖춘 $X$를 환 $A$ 위의
준스킴(또는 $A$-준스킴)이라고도 한다.
\end{definition}

(\hyperref[I.2.2.4-ko]{2.2.4})에 따르면, 환 $A$ 위의 준스킴을 주는
것은 구조층 $\mathscr{O}_X$가 $A$-\emph{대수들의 층}인 준스킴
$(X,\mathscr{O}_X)$를 주는 것과 동치이다. \emph{따라서 임의의
준스킴은 언제나 유일한 방식으로 $\mathbf{Z}$-준스킴으로 볼 수 있다.}

$\varphi:X\to S$가 $S$-준스킴 $X$의 구조 사상이면, 점 $x\in X$가
점 $s\in S$ \emph{위에 있다}는 것은 $\varphi(x)=s$라는 뜻이다.
$\varphi$가 우세 사상이면(\hyperref[I.2.2.6-ko]{2.2.6}), $X$가
$S$에 대해 \emph{우세하다}고 한다.

\begin{env}[2.5.2]
\label{I.2.5.2-ko}
$X$, $Y$를 두 $S$-준스킴이라 하자. 준스킴의 사상 $u:X\to Y$가
\emph{$S$ 위의 준스킴의 사상}(또는 \emph{$S$-사상})이라는 것은 도식
\phantomsection\label{I.2.5.2.diagram-ko}
\[
  \xymatrix{
    X\ar[rr]^u\ar[dr] & & Y\ar[dl]\\
    & S &
  }
\]
이 가환이라는 뜻이다. 여기서 빗금 화살표들은 구조 사상이다. 이
조건으로부터 모든 $s\in S$와 $s$ 위에 있는 모든 $x\in X$에 대하여
$u(x)$도 $s$ 위에 있어야 함이 따른다.
\end{env}

이 정의에서 두 $S$-사상의 합성이 $S$-사상임이 곧바로 보인다.
따라서 $S$-준스킴들은 하나의 \emph{범주}를 이룬다.

$S$-준스킴 $X$에서 $S$-준스킴 $Y$로 가는 $S$-사상들의 집합을
$\operatorname{Hom}_S(X,Y)$로 표기한다. $S$-준스킴 $X$의 항등사상은
$1_X$로 나타낸다.

$S$가 환 $A$의 아핀 스킴이면 $S$-사상을 $A$-\emph{사상}이라고도
한다.

\begin{env}[2.5.3]
\label{I.2.5.3-ko}
\oldpage[I]{104}
$X$가 $S$-준스킴이고 $v:X'\to X$가 준스킴의 사상이면, 합성 사상
$X'\xrightarrow{v}X\to S$는 $X'$를 $S$-준스킴으로 정의한다. 특히
$X$의 열린집합 $U$ 위에 유도된 모든 준스킴은 표준 단사사상을
통하여 $S$-준스킴으로 볼 수 있다.

$u:X\to Y$가 $S$-준스킴들의 $S$-사상이면, $X$의 열린집합 $U$ 위에
유도된 모든 준스킴에 대한 $u$의 제한은 $S$-사상 $U\to Y$이다.
역으로 $(U_\alpha)$가 $X$의 열린 덮개이고, 각 $\alpha$에 대하여
$u_\alpha:U_\alpha\to Y$가 $S$-사상이라고 하자. 모든 지표쌍
$(\alpha,\beta)$에 대하여 $U_\alpha\cap U_\beta$에 대한
$u_\alpha$와 $u_\beta$의 제한이 일치하면, 각 $U_\alpha$에 대한
제한이 $u_\alpha$인 $S$-사상 $X\to Y$가 유일하게 존재한다.

$u:X\to Y$가 $u(X)\subset V$를 만족하는 $S$-사상이고 $V$가 $Y$의
열린 부분집합이면, $u$를 $X$에서 $V$로 가는 사상으로 보아도 여전히
$S$-사상이다.
\end{env}

\begin{env}[2.5.4]
\label{I.2.5.4-ko}
$S'\to S$를 준스킴의 사상이라 하자. 모든 $S'$-준스킴 $X$에 대하여
합성 사상 $X\to S'\to S$는 $X$를 $S$-준스킴으로 정의한다.

역으로 $S'$이 $S$의 한 열린집합 위에 유도된 준스킴이라고 하자.
$X$가 $S$-준스킴이고 그 구조 사상 $f:X\to S$가
$f(X)\subset S'$을 만족하면, $X$를 $S'$-준스킴으로 볼 수 있다.
이 후자의 경우, $Y$가 바탕공간을 $S'$ 안으로 보내는 구조 사상을
가진 $S$-준스킴이면, $X$에서 $Y$로 가는 모든 $S$-사상은
$S'$-사상이기도 하다.
\end{env}

\begin{env}[2.5.5]
\label{I.2.5.5-ko}
$X$가 $S$-준스킴이고 $\varphi:X\to S$가 그 구조 사상이면,
$S$에서 $X$로 가는 $S$-사상, 즉 $\varphi\circ\psi$가 $S$의
항등사상인 준스킴의 사상 $\psi:S\to X$를 $X$의
\emph{$S$-단면}이라 한다. $X$의 $S$-단면들의 집합을
$\Gamma(X/S)$로 나타낸다.
\end{env}

\section{준스킴의 곱}
\label{section:I.3-ko}

\subsection{준스킴의 합}
\label{subsection:I.3.1-ko}

\phantomsection\label{I.3.1.1-ko}
$(X_\alpha)$를 임의의 준스킴족이라 하고, $X$를 그 바탕공간
$X_\alpha$들의 위상적 \emph{합}인 위상공간이라 하자. 그러면 $X$는
서로소인 열린 부분공간들 $X'_\alpha$의 합집합이고, 각 $\alpha$에
대하여 $X_\alpha$에서 $X'_\alpha$ 위로 가는 위상동형
$\varphi_\alpha$가 존재한다. 각 $X'_\alpha$에 층
$(\varphi_\alpha)_*(\mathscr{O}_{X_\alpha})$를 주면, $X$는 분명히
준스킴이 된다. 이를 준스킴족 $(X_\alpha)$의 \emph{합}이라 하며
$\bigsqcup_\alpha X_\alpha$로 표기한다. $Y$가 준스킴이면 사상
$f\mapsto(f\circ\varphi_\alpha)$는 집합 $\operatorname{Hom}(X,Y)$에서
곱집합 $\prod_\alpha\operatorname{Hom}(X_\alpha,Y)$ 위로 가는
\emph{전단사}이다. 특히 $X_\alpha$들이 구조 사상
$\psi_\alpha$를 가진 $S$-준스킴들이면, 모든 $\alpha$에 대하여
$\psi\circ\varphi_\alpha=\psi_\alpha$를 만족하는 유일한 사상
$\psi:X\to S$에 의해 $X$는 $S$-준스킴이 된다. 두 준스킴 $X$, $Y$의
합은 $X\sqcup Y$로 표기한다. $X=\operatorname{Spec}(A)$,
$Y=\operatorname{Spec}(B)$이면 $X\sqcup Y$가
$\operatorname{Spec}(A\times B)$와 표준적으로 동일시됨은 자명하다.

\subsection{준스킴의 곱}
\label{subsection:I.3.2-ko}

\begin{definition}[3.2.1]
\label{I.3.2.1-ko}
두 $S$-준스킴 $X$, $Y$가 주어졌다고 하자. $S$-준스킴 $Z$와 두
$S$-사상 $p_1:Z\to X$, $p_2:Z\to Y$로 이루어진 삼중항
$(Z,p_1,p_2)$가 다음 조건을 만족하면 이를 $S$-준스킴 $X$와 $Y$의
곱이라 한다. 모든 $S$-준스킴 $T$에 대하여 사상
\oldpage[I]{105}
$f\mapsto(p_1\circ f,p_2\circ f)$는 $T$에서 $Z$로 가는 $S$-사상들의
집합에서, $S$-사상 $T\to X$ 하나와 $S$-사상 $T\to Y$ 하나로
이루어진 쌍들의 집합 위로 가는 전단사이다. 다시 말해 다음은
전단사이다.
\[
  \operatorname{Hom}_S(T,Z)\xrightarrow{\sim}
  \operatorname{Hom}_S(T,X)\times\operatorname{Hom}_S(T,Y).
\]
\end{definition}

이는 범주의 두 대상의 \emph{곱}이라는 일반적인 개념을
$S$-준스킴들의 범주에 적용한 것이다(T, I, 1.1). 특히 두
$S$-준스킴의 곱은 유일한 $S$-동형까지 \emph{유일}하다. 이 유일성
때문에 두 $S$-준스킴 $X$, $Y$의 곱은 보통 $X\times_S Y$로
표기하고(혼동의 염려가 없으면 단순히 $X\times Y$로 표기한다),
$X\times_S Y$에서 각각 $X$와 $Y$로 가는 사상 $p_1$, $p_2$는
\emph{표준 사영}이라 부르되 표기에서는 생략한다. $g:T\to X$,
$h:T\to Y$가 두 $S$-사상이면 $p_1\circ f=g$, $p_2\circ f=h$를
만족하는 $S$-사상 $f:T\to X\times_S Y$를 $(g,h)_S$ 또는 단순히
$(g,h)$로 나타내겠다. $X'$, $Y'$이 두 $S$-준스킴이고 $p'_1$,
$p'_2$가 존재한다고 가정한 $X'\times_S Y'$의 표준 사영이며,
$u:X'\to X$, $v:Y'\to Y$가 두 $S$-사상이면, $X'\times_S Y'$에서
$X\times_S Y$로 가는 $S$-사상 $(u\circ p'_1,v\circ p'_2)_S$를
$u\times_S v$ 또는 단순히 $u\times v$로 쓰겠다.

$S$가 환 $A$의 아핀 스킴이면 앞의 표기에서 $S$를 $A$로 바꾸어
쓰는 경우가 많다.

\begin{proposition}[3.2.2]
\label{I.3.2.2-ko}
$X$, $Y$, $S$를 세 아핀 스킴이라 하고 각각의 환을 $B$, $C$, $A$라
하자. $Z=\operatorname{Spec}(B\otimes_A C)$라 하고, $p_1$, $p_2$를
$B$와 $C$에서 $B\otimes_A C$로 가는 표준 $A$-준동형
$u:b\mapsto b\otimes1$, $v:c\mapsto1\otimes c$에 대응하는
(\hyperref[I.2.2.4-ko]{2.2.4}) $S$-사상이라 하자. 그러면
$(Z,p_1,p_2)$는 $X$와 $Y$의 곱이다.
\end{proposition}

(\hyperref[I.2.2.4-ko]{2.2.4})에 따라, 모든 $A$-준동형
$f:B\otimes_A C\to L$($L$은 $A$-대수)에 쌍
$(f\circ u,f\circ v)$를 대응시키는 사상이 전단사
\[
  \operatorname{Hom}_A(B\otimes_A C,L)\xrightarrow{\sim}
  \operatorname{Hom}_A(B,L)\times\operatorname{Hom}_A(C,L)\footnotemark
\]
\footnotetext{여기서 표기 $\operatorname{Hom}_A$는 $A$-대수의
준동형들의 집합을 뜻한다.}를 정의함을 확인하면 충분하다. 이는 정의와
관계 $b\otimes c=(b\otimes1)(1\otimes c)$에서 곧바로 따른다.

\begin{corollary}[3.2.3]
\label{I.3.2.3-ko}
$T$를 환 $D$의 아핀 스킴이라 하자.
$\alpha=({}^a\rho,\widetilde{\rho})$
(각각 $\beta=({}^a\sigma,\widetilde{\sigma})$)를 $S$-사상 $T\to X$
(각각 $T\to Y$)라 하자. 여기서 $\rho$(각각 $\sigma$)는 $B$(각각
$C$)에서 $D$로 가는 $A$-준동형이다. 그러면
$(\alpha,\beta)_S=({}^a\tau,\widetilde{\tau})$이다. 여기서 $\tau$는
$\tau(b\otimes c)=\rho(b)\sigma(c)$를 만족하는 준동형
$B\otimes_A C\to D$이다.
\end{corollary}

\begin{proposition}[3.2.4]
\label{I.3.2.4-ko}
$f:S'\to S$를 준스킴의 단사사상이라 하고(T, I, 1.1), $X$, $Y$를
$f$에 의해 $S$-준스킴으로도 간주하는 두 $S'$-준스킴이라 하자.
$S$-준스킴 $X$, $Y$의 모든 곱은 $S'$-준스킴 $X$, $Y$의 곱이고,
그 역도 성립한다.
\end{proposition}

$\varphi:X\to S'$, $\psi:Y\to S'$를 구조 사상이라 하자. $T$가
$S$-준스킴이고 $u:T\to X$, $v:T\to Y$가 두 $S$-사상이면, 정의에
따라 $f\circ\varphi\circ u=f\circ\psi\circ v=\theta$이다. 여기서
$\theta$는 $T$의 구조 사상이다. $f$에 대한 가정으로부터
$\varphi\circ u=\psi\circ v=\theta'$을 얻으며, $T$를 구조 사상
$\theta'$을 가진 $S'$-준스킴으로, $u$와 $v$를 $S'$-사상으로 간주할
수 있다. 그러므로 (\hyperref[I.3.2.1-ko]{3.2.1})을 고려하면 명제의
결론이 곧바로 따른다.

\begin{corollary}[3.2.5]
\label{I.3.2.5-ko}
$X$, $Y$를 구조 사상 $\varphi:X\to S$, $\psi:Y\to S$를 가진 두
$S$-준스킴이라 하고, $S'$을 $\varphi(X)\subset S'$,
$\psi(Y)\subset S'$을 만족하는 $S$의 열린 부분집합이라 하자.
$S$-준스킴 $X$, $Y$의 모든 곱은 $S'$-준스킴 $X$, $Y$의 곱이고,
그 역도 성립한다.
\end{corollary}

\oldpage[I]{106}
표준 단사사상 $S'\to S$에 (\hyperref[I.3.2.4-ko]{3.2.4})를
적용하면 충분하다.

\begin{theorem}[3.2.6]
\label{I.3.2.6-ko}
두 $S$-준스킴 $X$, $Y$가 주어지면 곱 $X\times_S Y$가 존재한다.
\end{theorem}

여러 단계로 나누어 증명하겠다.

\begin{lemma}[3.2.6.1]
\label{I.3.2.6.1-ko}
$(Z,p,q)$를 $X$와 $Y$의 곱이라 하고, $U$, $V$를 각각 $X$, $Y$의
열린 부분집합이라 하자. $W=p^{-1}(U)\cap q^{-1}(V)$라 놓으면,
$W$와 $p$, $q$의 $W$로의 제한(각각 사상 $W\to U$, $W\to V$로
간주한다)로 이루어진 삼중항은 $U$와 $V$의 곱이다.
\end{lemma}

실제로 $T$가 $S$-준스킴이면, $S$-사상 $T\to W$를 $T$를 $W$ 안으로
보내는 $S$-사상 $T\to Z$와 동일시할 수 있다. 이제 $g:T\to U$,
$h:T\to V$가 임의의 두 $S$-사상이면, 이를 각각 $T$에서 $X$, $Y$로
가는 $S$-사상으로 간주할 수 있다. 가정에 따라 $g=p\circ f$,
$h=q\circ f$를 만족하는 유일한 $S$-사상 $f:T\to Z$가 존재한다.
$p(f(T))\subset U$, $q(f(T))\subset V$이므로
\[
  f(T)\subset p^{-1}(U)\cap q^{-1}(V)=W,
\]
이고, 주장이 성립한다.

\begin{lemma}[3.2.6.2]
\label{I.3.2.6.2-ko}
$Z$를 $S$-준스킴, $p:Z\to X$, $q:Z\to Y$를 두 $S$-사상,
$(U_\alpha)$를 $X$의 열린 덮개, $(V_\lambda)$를 $Y$의 열린 덮개라
하자. 모든 쌍 $(\alpha,\lambda)$에 대하여 $S$-준스킴
$W_{\alpha\lambda}=p^{-1}(U_\alpha)\cap q^{-1}(V_\lambda)$와 $p$, $q$의
$W_{\alpha\lambda}$로의 제한이 $U_\alpha$와 $V_\lambda$의 곱을
이룬다고 가정하자. 그러면 $(Z,p,q)$는 $X$와 $Y$의 곱이다.
\end{lemma}

먼저 $f_1$, $f_2$가 두 $S$-사상 $T\to Z$일 때, 관계
$p\circ f_1=p\circ f_2$, $q\circ f_1=q\circ f_2$로부터
$f_1=f_2$가 따름을 보이자. 실제로 $Z$는 $W_{\alpha\lambda}$들의
합집합이므로 $f_1^{-1}(W_{\alpha\lambda})$들은 $T$의 열린 덮개를
이루며, $f_2^{-1}(W_{\alpha\lambda})$들도 마찬가지이다. 더 나아가
가정에 따라
\[
\begin{aligned}
  f_1^{-1}(W_{\alpha\lambda})
  &=f_1^{-1}(p^{-1}(U_\alpha))\cap
    f_1^{-1}(q^{-1}(V_\lambda)) \\
  &=f_2^{-1}(p^{-1}(U_\alpha))\cap
    f_2^{-1}(q^{-1}(V_\lambda))
   =f_2^{-1}(W_{\alpha\lambda})
\end{aligned}
\]
이다. 따라서 모든 지표쌍에 대하여 $f_1$과 $f_2$의
$f_1^{-1}(W_{\alpha\lambda})=f_2^{-1}(W_{\alpha\lambda})$로의 제한이
같음을 보이면 충분하다. 그런데 이 제한들은
$f_1^{-1}(W_{\alpha\lambda})$에서 $W_{\alpha\lambda}$로 가는
$S$-사상으로 간주할 수 있으므로, 가정과 정의
(\hyperref[I.3.2.1-ko]{3.2.1})에서 주장이 따른다.

이제 두 $S$-사상 $g:T\to X$, $h:T\to Y$가 주어졌다고 하자.
$T_{\alpha\lambda}=g^{-1}(U_\alpha)\cap h^{-1}(V_\lambda)$라 놓으면
$T_{\alpha\lambda}$들은 $T$의 열린 덮개를 이룬다. 가정에 따라
$p\circ f_{\alpha\lambda}$와 $q\circ f_{\alpha\lambda}$가 각각
$g$와 $h$의 $T_{\alpha\lambda}$로의 제한이 되는 $S$-사상
$f_{\alpha\lambda}$가 존재한다. 또한 $f_{\alpha\lambda}$와
$f_{\beta\mu}$의 $T_{\alpha\lambda}\cap T_{\beta\mu}$로의 제한이
일치함을 보이자. 이것으로 (\hyperref[I.3.2.6.2-ko]{3.2.6.2})의
증명이 끝난다. 정의에 따라 $f_{\alpha\lambda}$와
$f_{\beta\mu}$에 의한 $T_{\alpha\lambda}\cap T_{\beta\mu}$의 상은
$W_{\alpha\lambda}\cap W_{\beta\mu}$에 포함된다. 한편
\[
  W_{\alpha\lambda}\cap W_{\beta\mu}
  =p^{-1}(U_\alpha\cap U_\beta)\cap
   q^{-1}(V_\lambda\cap V_\mu)
\]
이므로, (\hyperref[I.3.2.6.1-ko]{3.2.6.1})에 따라
$W_{\alpha\lambda}\cap W_{\beta\mu}$와 이 준스킴으로 제한한 $p$,
$q$는 $U_\alpha\cap U_\beta$와 $V_\lambda\cap V_\mu$의 곱을
이룬다. $p\circ f_{\alpha\lambda}$와 $p\circ f_{\beta\mu}$가
$T_{\alpha\lambda}\cap T_{\beta\mu}$에서 일치하고,
$q\circ f_{\alpha\lambda}$와 $q\circ f_{\beta\mu}$도 마찬가지이므로
$f_{\alpha\lambda}$와 $f_{\beta\mu}$는
$T_{\alpha\lambda}\cap T_{\beta\mu}$에서 일치한다. 증명이 끝났다.

\oldpage[I]{107}
\begin{lemma}[3.2.6.3]
\label{I.3.2.6.3-ko}
$(U_\alpha)$를 $X$의 열린 덮개, $(V_\lambda)$를 $Y$의 열린 덮개라
하고, 모든 쌍 $(\alpha,\lambda)$에 대하여 $U_\alpha$와
$V_\lambda$의 곱이 존재한다고 가정하자. 그러면 $X$와 $Y$의 곱이
존재한다.
\end{lemma}

(\hyperref[I.3.2.6.1-ko]{3.2.6.1})을 열린집합
$U_\alpha\cap U_\beta$와 $V_\lambda\cap V_\mu$에 적용하면 $X$와
$Y$가 각각 이 열린집합들 위에 유도하는 $S$-준스킴들의 곱이
존재함을 알 수 있다. 또한 곱의 유일성으로부터, $i=(\alpha,\lambda)$,
$j=(\beta,\mu)$라 놓을 때 이 곱에서
$U_\alpha\times_S V_\lambda$(각각 $U_\beta\times_S V_\mu$)가 한
열린집합 위에 유도하는 $S$-준스킴 $W_{ij}$(각각 $W_{ji}$) 위로 가는
표준 동형 $h_{ij}$(각각 $h_{ji}$)가 존재함을 알 수 있다. 따라서
$f_{ij}=h_{ij}\circ h_{ji}^{-1}$은 $W_{ji}$에서 $W_{ij}$ 위로 가는
동형이다. 더 나아가 세 번째 쌍 $k=(\gamma,\nu)$에 대하여
$W_{ki}\cap W_{kj}$에서 $f_{ik}=f_{ij}\circ f_{jk}$이다. 이는
(\hyperref[I.3.2.6.1-ko]{3.2.6.1})을 각각 $U_\beta$, $V_\mu$ 안의
열린집합 $U_\alpha\cap U_\beta\cap U_\gamma$와
$V_\lambda\cap V_\mu\cap V_\nu$에 적용하여 얻는다. 따라서 준스킴
$Z$, 그 바탕공간의 열린 덮개 $(Z_i)$, 그리고 각 $i$에 대하여 유도
준스킴 $Z_i$에서 준스킴 $U_\alpha\times_S V_\lambda$ 위로 가는
동형 $g_i$가 존재하여 모든 쌍 $(i,j)$에 대하여
$f_{ij}=g_i\circ g_j^{-1}$이 성립한다
(\hyperref[I.2.3.1-ko]{2.3.1}). 또한
$g_i(Z_i\cap Z_j)=W_{ij}$이다. $p_i$, $q_i$, $\theta_i$를
$S$-준스킴 $U_\alpha\times_S V_\lambda$의 사영들과 구조 사상이라
하면, $Z_i\cap Z_j$에서 $p_i\circ g_i=p_j\circ g_j$이고 나머지 두
사상에 대해서도 마찬가지임을 곧바로 확인할 수 있다. 따라서 각
$Z_i$에서 $p$(각각 $q$, $\theta$)가 $p_i\circ g_i$(각각
$q_i\circ g_i$, $\theta_i\circ g_i$)와 일치한다는 조건으로 준스킴의
사상 $p:Z\to X$(각각 $q:Z\to Y$, $\theta:Z\to S$)를 정의할 수
있다. $\theta$를 갖춘 $Z$는 $S$-준스킴이다. 이제
$Z'_i=p^{-1}(U_\alpha)\cap q^{-1}(V_\lambda)$가 $Z_i$와 같음을
보이자. 모든 지표 $j=(\beta,\mu)$에 대하여
$Z_j\cap Z'_i=g_j^{-1}(p_j^{-1}(U_\alpha)\cap q_j^{-1}(V_\lambda))$이다.
그런데
\[
  p_j^{-1}(U_\alpha)\cap q_j^{-1}(V_\lambda)
  =p_j^{-1}(U_\alpha\cap U_\beta)\cap
   q_j^{-1}(V_\lambda\cap V_\mu).
\]
(\hyperref[I.3.2.6.1-ko]{3.2.6.1})에 따라 $p_j$와 $q_j$의
$p_j^{-1}(U_\alpha)\cap q_j^{-1}(V_\lambda)$로의 제한은 이
$S$-준스킴에 $U_\alpha\cap U_\beta$와 $V_\lambda\cap V_\mu$의 곱
구조를 준다. 그런데 곱의 유일성으로부터
$p_j^{-1}(U_\alpha)\cap q_j^{-1}(V_\lambda)=W_{ji}$가 따른다.
그러므로 모든 $j$에 대하여 $Z_j\cap Z'_i=Z_j\cap Z_i$이고,
$Z'_i=Z_i$이다. 이제 (\hyperref[I.3.2.6.2-ko]{3.2.6.2})에 따라
$(Z,p,q)$는 $X$와 $Y$의 곱이다.

\begin{lemma}[3.2.6.4]
\label{I.3.2.6.4-ko}
$\varphi:X\to S$, $\psi:Y\to S$를 $X$, $Y$의 구조 사상이라 하고,
$(S_i)$를 $S$의 열린 덮개라 하자. $X_i=\varphi^{-1}(S_i)$,
$Y_i=\psi^{-1}(S_i)$라 놓자. 각 곱 $X_i\times_S Y_i$가 존재하면
$X\times_S Y$가 존재한다.
\end{lemma}

(\hyperref[I.3.2.6.3-ko]{3.2.6.3})에 따라 모든 $i$, $j$에 대하여
곱 $X_i\times_S Y_j$가 존재함을 보이면 충분하다.
$X_{ij}=X_i\cap X_j=\varphi^{-1}(S_i\cap S_j)$,
$Y_{ij}=Y_i\cap Y_j=\psi^{-1}(S_i\cap S_j)$라 놓자.
(\hyperref[I.3.2.6.1-ko]{3.2.6.1})에 따라 곱
$Z_{ij}=X_{ij}\times_S Y_{ij}$가 존재한다. 이제 $T$가 $S$-준스킴이고
$g:T\to X_i$, $h:T\to Y_j$가 $S$-사상이면, $S$-사상의 정의에 따라
반드시 $\varphi(g(T))=\psi(h(T))\subset S_i\cap S_j$이다. 따라서
$g(T)\subset X_{ij}$, $h(T)\subset Y_{ij}$이고, $Z_{ij}$가 $X_i$와
$Y_j$의 곱임은 자명하다.

\begin{env}[3.2.6.5]
\label{I.3.2.6.5-ko}
이제 정리 (\hyperref[I.3.2.6-ko]{3.2.6})의 증명을 끝낼 수 있다.
$S$가 \emph{아핀 스킴}이면, $X$와 $Y$에는 각각 아핀 열린집합들로
이루어진 덮개 $(U_\alpha)$, $(V_\lambda)$가 존재한다.
(\hyperref[I.3.2.2-ko]{3.2.2})에 따라
$U_\alpha\times_S V_\lambda$가 존재하므로
(\hyperref[I.3.2.6.3-ko]{3.2.6.3})에 따라 $X\times_S Y$도 존재한다.
$S$가 임의의 준스킴이면 아핀 열린집합들로 이루어진 $S$의 덮개
$(S_i)$가 존재한다. $\varphi:X\to S$, $\psi:Y\to S$가 구조 사상이고
$X_i=\varphi^{-1}(S_i)$, $Y_i=\psi^{-1}(S_i)$라 놓으면 앞에서 보인
바에 따라 곱 $X_i\times_{S_i}Y_i$가 존재한다.
\oldpage[I]{108}
그러면 (\hyperref[I.3.2.5-ko]{3.2.5})에 따라 곱
$X_i\times_S Y_i$도 존재하고, (\hyperref[I.3.2.6.4-ko]{3.2.6.4})에
따라 $X\times_S Y$도 존재한다.
\end{env}

\begin{corollary}[3.2.7]
\label{I.3.2.7-ko}
$Z=X\times_S Y$를 두 $S$-준스킴의 곱, $p$, $q$를 $Z$에서 $X$, $Y$로
가는 사영, $\varphi$(각각 $\psi$)를 $X$(각각 $Y$)의 구조 사상이라
하자. $S'$을 $S$의 열린 부분집합, $U$(각각 $V$)를
$\varphi^{-1}(S')$(각각 $\psi^{-1}(S')$)에 포함되는 $X$(각각
$Y$)의 열린 부분집합이라 하자. 그러면 곱 $U\times_{S'}V$는
$p^{-1}(U)\cap q^{-1}(V)$ 위에 $Z$가 유도하는 준스킴($S'$-준스킴으로
간주한다)과 표준적으로 동일시된다. 또한 $f:T\to X$, $g:T\to Y$가
$f(T)\subset U$, $g(T)\subset V$를 만족하는 $S$-사상이면,
$S'$-사상 $(f,g)_{S'}$는 $(f,g)_S:T\to Z$의 공역을
$p^{-1}(U)\cap q^{-1}(V)$로 제한하여 얻은 사상과 동일시된다.
\end{corollary}

이는 (\hyperref[I.3.2.5-ko]{3.2.5})와
(\hyperref[I.3.2.6.1-ko]{3.2.6.1})에서 따른다.

\begin{env}[3.2.8]
\label{I.3.2.8-ko}
$(X_\alpha)$, $(Y_\lambda)$를 두 $S$-준스킴족이라 하고, $X$(각각
$Y$)를 족 $(X_\alpha)$(각각 $(Y_\lambda)$)의 합이라 하자
(\hyperref[subsection:I.3.1-ko]{3.1}). 그러면 $X\times_S Y$는 족
$(X_\alpha\times_S Y_\lambda)$의 \emph{합}과 동일시된다. 이는
(\hyperref[I.3.2.6.3-ko]{3.2.6.3})에서 곧바로 따른다.
\end{env}

\subsection{곱의 형식적 성질~; 밑준스킴의 변경}
\label{subsection:I.3.3-ko}

\begin{env}[3.3.1]
\label{I.3.3.1-ko}
독자는 이 절에서 진술하는 모든 성질이
(\hyperref[I.3.3.13-ko]{3.3.13})과
(\hyperref[I.3.3.15-ko]{3.3.15})을 제외하면, 그 진술에 나오는
곱들이 \emph{존재}할 때 \emph{임의의 범주}에서 아무 변경 없이
성립함을 알 수 있을 것이다. 실제로 범주의 임의의 대상 $S$에
대해서도 $S$-대상과 $S$-사상의 개념을
(\hyperref[subsection:I.2.5-ko]{2.5})에서와 정확히 같은 방식으로
정의할 수 있음은 분명하다.
\end{env}

\begin{env}[3.3.2]
\label{I.3.3.2-ko}
먼저 $X\times_S Y$는 $S$-준스킴의 범주에서 $X$와 $Y$에 관한
\emph{공변 쌍함자}이다. 실제로 다음 그림식이 가환임을 확인하면
충분하다.
\[
  \xymatrix{
    X\times Y\ar[r]^{f\times 1}\ar[d] &
    X'\times Y\ar[r]^{f'\times 1}\ar[d] &
    X''\times Y\ar[d]\\
    X\ar[r]_f & X'\ar[r]_{f'} & X''
  }
\]
\end{env}

\begin{proposition}[3.3.3]
\label{I.3.3.3-ko}
모든 $S$-준스킴 $X$에 대하여 $X\times_S S$(각각 $S\times_S X$)의
첫째 사영(각각 둘째 사영)은 $X\times_S S$(각각 $S\times_S X$)에서
$X$ 위로 가는 함자적 동형이고, 그 역동형은
$(1_X,\varphi)_S$(각각 $(\varphi,1_X)_S$)이다. 여기서 $\varphi$는
구조 사상 $X\to S$이다. 따라서 표준 동형까지
\[
  X\times_S S=S\times_S X=X
\]
라고 쓸 수 있다.
\end{proposition}

삼중항 $(X,1_X,\varphi)$가 $X$와 $S$의 곱임을 보이면 충분하다.
그런데 $T$가 $S$-준스킴이면, $T$에서 $S$로 가는 유일한 $S$-사상은
반드시 구조 사상 $\psi:T\to S$이다. $f$가 $T$에서 $X$로 가는
$S$-사상이면 반드시 $\psi=\varphi\circ f$이므로 주장이 따른다.

\begin{corollary}[3.3.4]
\label{I.3.3.4-ko}
$X$, $Y$를 두 $S$-준스킴, $\varphi:X\to S$, $\psi:Y\to S$를 그
구조 사상이라 하자. $X$를 $X\times_S S$와, $Y$를 $S\times_S Y$와
표준적으로 동일시하면, 사영 $X\times_S Y\to X$와
$X\times_S Y\to Y$는 각각 $1_X\times\psi$와
$\varphi\times1_Y$와 동일시된다.
\end{corollary}

확인은 즉시 되므로 독자에게 맡긴다.

\begin{env}[3.3.5]
\label{I.3.3.5-ko}
임의의 유한개 $n$개의 $S$-준스킴의 곱도
(\hyperref[subsection:I.3.2-ko]{3.2})에서와 같은 방식으로 정의할
수 있다.
\oldpage[I]{109}
이러한 곱들의 존재는 $n$에 관한 귀납법으로
(\hyperref[I.3.2.6-ko]{3.2.6})에서 따른다. 실제로
$(X_1\times_S X_2\times_S\cdots\times_S X_{n-1})\times_S X_n$은
곱의 정의를 만족한다. 곱의 유일성에서, 임의의 범주에서와 마찬가지로
그 \emph{교환법칙}과 \emph{결합법칙}이 따른다. 예를 들어
$p_1$, $p_2$, $p_3$가 $X_1\times_S X_2\times_S X_3$의 사영들이고
이 준스킴을 $(X_1\times_S X_2)\times_S X_3$와 동일시하면,
$X_1\times_S X_2$로 가는 사영은 $(p_1,p_2)_S$와 동일시된다.
\end{env}

\begin{env}[3.3.6]
\label{I.3.3.6-ko}
$S$, $S'$을 두 준스킴, $\varphi:S'\to S$를 $S'$에 $S$-준스킴
구조를 주는 사상이라 하자. 모든 $S$-준스킴 $X$에 대하여 곱
$X\times_S S'$을 생각하고, 각각 $X$와 $S'$로 가는 그 사영을
$p$와 $\pi'$이라 하자. $\pi'$을 갖춘 이 곱은 $S'$-준스킴이다.
이를 $S'$-준스킴으로 간주할 때 $X_{(S')}$ 또는
$X_{(\varphi)}$로 나타내고, $\varphi$를 사용하여 밑준스킴을
$S$에서 $S'$으로 \emph{확장하여 얻은 준스킴}, 또는 $\varphi$에
의한 $X$의 \emph{역상}이라 한다. $\pi$가 $X$의 구조 사상이고,
$\theta$가 $X\times_S S'$을 $S$-준스킴으로 간주했을 때의 구조
사상이면, 다음 그림식이 가환임에 유의하자.
\phantomsection\label{I.3.3.6.base-change-diagram-ko}
\[
  \xymatrix{
    X\ar[d]_{\pi} &
    X_{(S')}\ar[l]_{p}\ar[ld]_{\theta}\ar[d]^{\pi'}\\
    S & S'\ar[l]^{\varphi}
  }
\]
\end{env}

\begin{env}[3.3.7]
\label{I.3.3.7-ko}
(\hyperref[I.3.3.6-ko]{3.3.6})의 표기를 사용하자. 모든 $S$-사상
$f:X\to Y$에 대하여 $S'$-사상
$f\times_S1:X_{(S')}\to Y_{(S')}$도 $f_{(S')}$로 나타내고,
$\varphi$에 의한 사상 $f$의 \emph{역상}이라 한다. 따라서
$X_{(S')}$은 $X$에 관한 \emph{공변 함자}이며, $S$-준스킴의
범주에서 $S'$-준스킴의 범주로 간다.
\end{env}

\begin{env}[3.3.8]
\label{I.3.3.8-ko}
준스킴 $X_{(S')}$은 하나의 \emph{보편 사상 문제}의 해로도 볼 수
있다. 모든 $S'$-준스킴 $T$는 $\varphi$를 통하여 $S$-준스킴이기도
하다. 모든 $S$-사상 $g:T\to X$는 $g=p\circ f$라는 꼴로 유일하게
쓰인다. 여기서 $f$는 $S'$-사상 $T\to X_{(S')}$이다. 이는 곱의
정의를 $S$-사상 $g$와 $\psi:T\to S'$($T$의 구조 사상)에 적용하면
따른다.
\end{env}

\begin{proposition}[3.3.9]
\label{I.3.3.9-ko}
\textup{(「밑준스킴 확장의 추이성」).} $S''$을 준스킴,
$\varphi':S''\to S'$을 사상이라 하자. 모든 $S$-준스킴 $X$에
대하여 $S''$-준스킴 $(X_{(\varphi)})_{(\varphi')}$에서
$S''$-준스킴 $X_{(\varphi\circ\varphi')}$ 위로 가는 표준적이고
함자적인 동형이 존재한다.
\end{proposition}

실제로 $T$를 $S''$-준스킴, $\psi$를 그 구조 사상, $g$를 $T$에서
$X$로 가는 $S$-사상이라 하자. 여기서 $T$는 구조 사상이
$\varphi\circ\varphi'\circ\psi$인 $S$-준스킴으로 간주한다.
$T$는 구조 사상이 $\varphi'\circ\psi$인 $S'$-준스킴이기도 하므로
$g=p\circ g'$으로 쓸 수 있다. 여기서 $g'$은
$S'$-사상 $T\to X_{(\varphi)}$이다. 이어서
$g'=p'\circ g''$으로 쓸 수 있다. 여기서 $g''$은
$S''$-사상 $T\to(X_{(\varphi)})_{(\varphi')}$이다.
\phantomsection\label{I.3.3.9.transitivity-diagram-ko}
\[
  \xymatrix{
    X\ar[d]_{\pi} &
    X_{(\varphi)}\ar[l]_{p}\ar[d]_{\pi'} &
    (X_{(\varphi)})_{(\varphi')}\ar[l]_{p'}\ar[d]_{\pi''}\\
    S & S'\ar[l]^{\varphi} & S''\ar[l]^{\varphi'}
  }
\]
\oldpage[I]{110}
그러므로 보편 사상 문제의 해의 유일성에 따라 명제가 성립한다.

이 결과는 혼동의 염려가 없을 때 표준 동형까지
$(X_{(S')})_{(S'')}=X_{(S'')}$라고 써도 되며, 다음과 같이 쓸
수도 있다.
\begin{equation}
  (X\times_S S')\times_{S'}S''=X\times_S S''~;
  \tag{3.3.9.1}\label{I.3.3.9.1-ko}
\end{equation}
(\hyperref[I.3.3.9-ko]{3.3.9})에서 정의한 동형의 함자성도 모든
$S$-사상 $f:X\to Y$에 대하여 성립하는
\emph{사상의 역상의 추이성} 공식
\begin{equation}
  (f_{(S')})_{(S'')}=f_{(S'')}
  \tag{3.3.9.2}\label{I.3.3.9.2-ko}
\end{equation}
으로 나타난다.

\begin{corollary}[3.3.10]
\label{I.3.3.10-ko}
$X$와 $Y$가 두 $S$-준스킴이면, $S'$-준스킴
$X_{(S')}\times_{S'}Y_{(S')}$에서 $S'$-준스킴
$(X\times_S Y)_{(S')}$ 위로 가는 표준적이고 함자적인 동형이
존재한다.
\end{corollary}

실제로 표준 동형까지
\[
  (X\times_S S')\times_{S'}(Y\times_S S')
  =X\times_S(Y\times_S S')=(X\times_S Y)\times_S S'
\]
이다. 이는 (\hyperref[I.3.3.9.1-ko]{3.3.9.1})과 $S$-준스킴의
곱의 결합법칙에서 따른다.

(\hyperref[I.3.3.10-ko]{3.3.10})에서 정의한 동형의 함자성은 모든
$S$-사상쌍 $u:T\to X$, $v:T\to Y$에 대하여 성립하는 공식
\begin{equation}
  (u_{(S')},v_{(S')})_{S'}=((u,v)_S)_{(S')}
  \tag{3.3.10.1}\label{I.3.3.10.1-ko}
\end{equation}
으로 나타난다.

다시 말해 역상함자 $X_{(S')}$은 \emph{곱의 형성과 가환}한다.
또한 합의 형성과도 가환한다(\hyperref[I.3.2.8-ko]{3.2.8}).

\begin{corollary}[3.3.11]
\label{I.3.3.11-ko}
$Y$를 $S$-준스킴, $f:X\to Y$를 $X$에 $Y$-준스킴 구조를 주는
사상이라 하자. 따라서 $X$는 $S$-준스킴이기도 하다. 그러면 준스킴
$X_{(S')}$은 곱 $X\times_Y Y_{(S')}$와 동일시되고, 사영
$X\times_Y Y_{(S')}\to Y_{(S')}$은 $f_{(S')}$와 동일시된다.
\end{corollary}

$\psi:Y\to S$를 $Y$의 구조 사상이라 하자. 다음 그림식은
가환이다.
\phantomsection\label{I.3.3.11.diagram-ko}
\[
  \xymatrix{
    S'\ar[d] &
    Y_{(S')}\ar[l]_{\psi_{(S')}}\ar[d] &
    X_{(S')}\ar[l]_{f_{(S')}}\ar[d]\\
    S & Y\ar[l]^{\psi} & X\ar[l]^{f}
  }
\]
그런데 $Y_{(S')}$은 $S'_{(\psi)}$와, $X_{(S')}$은
$S'_{(\psi\circ f)}$와 동일시된다.
(\hyperref[I.3.3.9-ko]{3.3.9})와
(\hyperref[I.3.3.4-ko]{3.3.4})를 고려하면 따름정리가 따른다.

\begin{env}[3.3.12]
\label{I.3.3.12-ko}
$f:X\to X'$, $g:Y\to Y'$을 준스킴의 \emph{단사사상}인 두
$S$-사상이라 하자(T, I, 1.1). 그러면 $f\times_S g$도
\emph{단사사상}이다. 실제로 $p$, $q$를 $X\times_S Y$의 사영,
$p'$, $q'$을 $X'\times_S Y'$의 사영, $u$, $v$를 두
$S$-사상 $T\to X\times_S Y$라 하자. 관계
$(f\times_S g)\circ u=(f\times_S g)\circ v$에서
$p'\circ(f\times_S g)\circ u=p'\circ(f\times_S g)\circ v$, 즉
$f\circ p\circ u=f\circ p\circ v$가 따른다. $f$가
단사사상이므로 $p\circ u=p\circ v$이다. 마찬가지로 $g$가
단사사상임을 사용하면 $q\circ u=q\circ v$를 얻으므로 $u=v$이다.
\oldpage[I]{111}
따라서 밑준스킴의 모든 확장 $S'\to S$에 대하여
\phantomsection\label{I.3.3.12.base-change-monomorphism-ko}
\[
  f_{(S')}:X_{(S')}\to Y_{(S')}
\]
은 단사사상이다.
\end{env}

\begin{env}[3.3.13]
\label{I.3.3.13-ko}
$S$, $S'$을 각각 환 $A$, $A'$의 두 아핀 스킴이라 하자. 따라서
사상 $S'\to S$은 환 준동형 $A\to A'$에 대응한다. $X$가
$S$-준스킴이면 $S'$-준스킴 $X_{(S')}$을
$X_{(A')}$ 또는 $X\otimes_A A'$로도 나타내겠다. $X$ 자체가 환
$B$의 아핀 스킴이면, $X_{(A')}$은 환
$B_{(A')}=B\otimes_A A'$의 아핀 스킴이다. 이 환은
$A$-대수 $B$의 스칼라환을 $A'$으로 확장하여 얻는다.
\end{env}

\begin{env}[3.3.14]
\label{I.3.3.14-ko}
(\hyperref[I.3.3.6-ko]{3.3.6})의 표기를 사용하자. 모든
\emph{$S$-사상} $f:S'\to X$에 대하여
$f'=(f,1_{S'})_S$는 $p\circ f'=f$, $\pi'\circ f'=1_{S'}$을
만족하는 $S'$-사상 $S'\to X'=X_{(S')}$이다. 다시 말해 $f'$은
\emph{$X'$의 $S'$-단면}이다. 역으로 $f'$이 이러한 $S'$-단면이면
$f=p\circ f'$은 $S$-사상 $S'\to X$이다. 이로써 표준적인
\emph{전단사 대응}
\phantomsection\label{I.3.3.14.correspondence-ko}
\[
  \operatorname{Hom}_S(S',X)\overset{\sim}{\longrightarrow}
  \operatorname{Hom}_{S'}(S',X')
\]
을 정의한다. $f'$을 $f$의 \emph{그래프 사상}이라 하고
$\Gamma_f$로 나타낸다.
\end{env}

\begin{env}[3.3.15]
\label{I.3.3.15-ko}
준스킴 $X$가 주어지면 언제나 이를 $\mathbf Z$-준스킴으로 간주할
수 있다. 특히 (\hyperref[I.3.3.14-ko]{3.3.14})에 따라
$X\otimes_{\mathbf Z}\mathbf Z[T]$의 $X$-단면들($T$는
부정원)은 $\mathbf Z[T]$에서 $X$로 가는 \emph{사상}들과
전단사로 대응한다. 이 $X$-단면들이
\emph{$X$ 위의 구조층 $\mathcal O_X$의 단면들}과도 전단사로
대응함을 보이자. $(U_\alpha)$를 아핀 열린집합들로 이루어진 $X$의
덮개라 하자. $u:X\to X\otimes_{\mathbf Z}\mathbf Z[T]$를
$X$-사상, $u_\alpha$를 그 $U_\alpha$로의 제한이라 하자.
$A_\alpha$가 아핀 스킴 $U_\alpha$의 환이면
$U_\alpha\otimes_{\mathbf Z}\mathbf Z[T]$는 환
$A_\alpha[T]$의 아핀 스킴이고
(\hyperref[I.3.2.2-ko]{3.2.2}), $u_\alpha$는 표준적으로
$A_\alpha$-준동형 $A_\alpha[T]\to A_\alpha$에 대응한다(1.7.3).
그런 준동형은 $T$의 $A_\alpha$에서의 상, 즉
$s_\alpha\in A_\alpha=\Gamma(U_\alpha,\mathcal O_X)$를 주면
완전히 결정된다. $u_\alpha$와 $u_\beta$의 제한이
$V\subset U_\alpha\cap U_\beta$인 아핀 열린집합 $V$에서
일치한다는 조건을 쓰면, $s_\alpha$와 $s_\beta$도 $V$에서
일치함을 곧바로 알 수 있다. 따라서 족 $(s_\alpha)$은
$X$ 위의 $\mathcal O_X$의 한 단면 $s$를 각 $U_\alpha$로 제한하여
얻은 족이다. 역으로 이러한 단면은, 어떤 $X$-사상
$X\to X\otimes_{\mathbf Z}\mathbf Z[T]$을 각 $U_\alpha$로
제한하여 얻은 사상족 $(u_\alpha)$를 정의함이 분명하다. 이 결과는
II, 1.7.12에서 일반화될 것이다.
\end{env}

\subsection{준스킴에 값을 갖는 준스킴의 점; 기하적 점}
\label{subsection:I.3.4-ko}

\begin{env}[3.4.1]
\label{I.3.4.1-ko}
$X$를 준스킴이라 하자. 모든 준스킴 $T$에 대하여 $X(T)$로
사상 $T\to X$들의 집합 $\operatorname{Hom}(T,X)$도 나타내며, 이 집합의
원소를 \emph{$T$에 값을 갖는 $X$의 점}이라고도 한다. 각 사상
$f:T\to T'$에 $X(T')$에서 $X(T)$으로 가는 사상
$u'\mapsto u'\circ f$를 대응시키면, $X$를 고정했을 때 $X(T)$가
준스킴의 범주에서 집합의 범주로 가는 \emph{$T$에 관한 반변함자}임을
알 수 있다. 또한 준스킴의 모든 사상 $g:X\to Y$는 $v\in X(T)$에
$g\circ v$를 대응시키는 함자적 사상 $X(T)\to Y(T)$를 정의한다.
\end{env}

\begin{env}[3.4.2]
\label{I.3.4.2-ko}
세 집합 $P$, $Q$, $R$와 두 사상 $\varphi:P\to R$, $\psi:Q\to R$가
주어졌다고 하자. $\varphi(p)=\psi(q)$를 만족하는 순서쌍 $(p,q)$들로
이루어진 곱집합 $P\times Q$의 부분집합을 ($\varphi$와 $\psi$에 관한)
\emph{$R$ 위에서의 $P$와 $Q$의 섬유곱}이라 하고
\oldpage[I]{112}
$P\times_R Q$로 나타낸다. $S$-준스킴의 곱의 정의
(\hyperref[I.3.2.1-ko]{3.2.1})는
(\hyperref[I.3.4.1-ko]{3.4.1})의 표기를 사용하면 다음 식으로도
해석할 수 있다.
\phantomsection\label{I.3.4.2.1-ko}
\[
  (X\times_S Y)(T)=X(T)\times_{S(T)}Y(T)
  \tag{3.4.2.1}
\]
여기서 사상 $X(T)\to S(T)$와 $Y(T)\to S(T)$는 각각 구조 사상
$X\to S$와 $Y\to S$에 대응한다.
\end{env}

\begin{env}[3.4.3]
\label{I.3.4.3-ko}
준스킴 $S$가 주어져 있고 $S$-준스킴과 $S$-사상만을 생각할 때,
$S$-사상 $T\to X$들의 집합 $\operatorname{Hom}_S(T,X)$을
$X(T)_S$로도 나타내며 혼동의 염려가 없으면 첨자 $S$를 생략한다.
$X(T)_S$의 원소를 \emph{점}(혼동할 염려가 있으면 \emph{$S$-점})이라고도
하며, 이는 \emph{$T$에 값을 갖는 $S$-준스킴 $X$의 점}이다. 특히 $X$의
\emph{$S$-단면}은 다름 아닌 \emph{$S$에 값을 갖는 $X$의 점}이다.
(\hyperref[I.3.4.2.1-ko]{3.4.2.1})의 식은 다음과 같이도 쓸 수 있다.
\phantomsection\label{I.3.4.3.1-ko}
\[
  (X\times_S Y)(T)_S=X(T)_S\times Y(T)_S~;
  \tag{3.4.3.1}
\]
더 일반적으로 $Z$가 $S$-준스킴이고 $X$, $Y$, $T$가
$Z$-준스킴(따라서 \emph{그 자체로} $S$-준스킴)이면
\phantomsection\label{I.3.4.3.2-ko}
\[
  (X\times_Z Y)(T)_S=X(T)_S\times_{Z(T)_S}Y(T)_S.
  \tag{3.4.3.2}
\]

$S$-준스킴 $W$와 두 $S$-사상 $r:W\to X$, $s:W\to Y$로 이루어진
삼중항 $(W,r,s)$가 ($Z$ 위에서) $X$와 $Y$의 곱임을 보이려면,
정의에 따라 \emph{모든 $S$-준스킴 $T$}에 대하여 다음 그림에서
\phantomsection\label{I.3.4.3.product-diagram-ko}
\[
  \xymatrix{
    W(T)_S\ar[r]^{r'}\ar[d]_{s'} & X(T)_S\ar[d]^{\varphi'}\\
    Y(T)_S\ar[r]_{\psi'} & Z(T)_S
  }
\]
$W(T)_S$가 $Z(T)_S$ 위에서 $X(T)_S$와 $Y(T)_S$의 섬유곱임을
확인하면 충분하다. 여기서 $r'$과 $s'$은 $r$과 $s$에,
$\varphi'$과 $\psi'$은 구조 사상 $\varphi:X\to Z$와
$\psi:Y\to Z$에 각각 대응한다.
\end{env}

\begin{env}[3.4.4]
\label{I.3.4.4-ko}
위에서 $T$(각각 $S$)가 환 $B$(각각 $A$)의 아핀 스킴이면 앞의
표기에서 $T$(각각 $S$)를 $B$(각각 $A$)로 바꾼다. 이때 $X(B)$의
원소를 \emph{환 $B$에 값을 갖는 $X$의 점}, $X(B)_A$의 원소를
\emph{$A$-대수 $B$에 값을 갖는 $A$-준스킴 $X$의 점}이라고 한다.
$X(B)$와 $X(B)_A$는 이제 \emph{$B$에 관한 공변함자}임에 유의하자.
마찬가지로 $A$-준스킴 $T$에 값을 갖는 $A$-준스킴 $X$의 점들의
집합을 $X(T)_A$로 쓴다.
\end{env}

\begin{env}[3.4.5]
\label{I.3.4.5-ko}
특히 $T=\operatorname{Spec}(A)$이고 $A$가 \emph{국소}환인 경우를
생각하자. 그러면 $X(A)$의 원소들은 $x\in X$에 대한 \emph{국소}
준동형 $\mathcal O_x\to A$들과 전단사로 대응한다
(\hyperref[I.2.4.4-ko]{2.4.4}). 이때 $X$의 바탕공간의 점 $x$를
그에 대응하는 $A$에 값을 갖는 $X$의 점의 \emph{소재점}이라 한다.

더욱 특별히 준스킴 $X$의 \emph{기하적 점}이란
\emph{어떤 체 $K$에 값을 갖는 $X$의 점}을 말한다. 따라서 이러한
점 하나를 주는 것은 그 점의 $X$의 바탕공간에서의
\oldpage[I]{113}
소재점 $x$와 $k(x)$의 \emph{확대체} $K$를 주는 것과 같다. $K$를
그 기하적 점의 \emph{값체}라고 하며, 이 기하적 점이
\emph{$x$에 위치한다}고도 한다. 이로써 $K$에 값을 갖는 기하적 점을
그 소재점에 대응시키는 사상 $X(K)\to X$를 정의한다.

$S'=\operatorname{Spec}(K)$가 $S$-준스킴이고(다시 말해
$s\in S$인 어떤 잉여체 $k(s)$의 확대체로 $K$를 보고 있고),
$X$가 $S$-준스킴이라고 하자. $X(K)_S$의 원소, 즉
\emph{$K$에 값을 갖고 $s$ 위에 있는 $X$의 기하적 점}은
$s$ \emph{위에 있는} $X$의 점 $x$에 대하여 잉여체 $k(x)$에서
$K$로 가는 $k(s)$-단사 준동형을 주는 것과 같다. 이때 $k(x)$는
$k(s)$의 확대체이다.

더욱 특별히 $S=\operatorname{Spec}(K)=\{\xi\}$이면
\emph{$K$에 값을 갖는 $X$의 기하적 점들은 $k(x)=K$인 점
$x\in X$와 동일시된다}. 이러한 점을 \emph{$K$ 위에서 유리적인
$K$-준스킴 $X$의 점}, 곧 $K$-유리점이라고도 한다. $K'$이 $K$의
확대체이면 $K'$에 값을 갖는 $X$의 기하적 점들은
$K'$ 위에서 유리적인 $X'=X_{(K')}$의 점들과 전단사로 대응한다
(\hyperref[I.3.3.14-ko]{3.3.14}).
\end{env}

\begin{lemma}[3.4.6]
\label{I.3.4.6-ko}
$X_i$ ($1\leq i\leq n$)를 $S$-준스킴들, $s$를 $S$의 한 점,
$x_i$ ($1\leq i\leq n$)를 $s$ 위에 있는 $X_i$의 한 점이라 하자.
그러면 $k(s)$의 확대체 $K$와, 곱
$Y=X_1\times_S X_2\times_S\cdots\times_S X_n$의 $K$에 값을 갖는
기하적 점이 존재하여 그 각 사영은 $x_i$에 위치한다.
\end{lemma}

실제로 어떤 하나의 $k(s)$의 확대체 $K$ 안에 $k(s)$-단사 준동형
$k(x_i)\to K$들을 잡을 수 있다(Bourbaki, \emph{Alg.}, chap. V,
\S~4, prop. 2). 합성 $k(s)\to k(x_i)\to K$들은 모두 같으므로
$k(x_i)\to K$에 대응하는 사상 $\operatorname{Spec}(K)\to X_i$들은
$S$-사상이다. 따라서 이들은 유일한 사상
$\operatorname{Spec}(K)\to Y$를 정의한다. $y$를 이에 대응하는
$Y$의 점이라 하면, $y$의 각 $X_i$로의 사영은 분명히 $x_i$이다.

\begin{proposition}[3.4.7]
\label{I.3.4.7-ko}
$X_i$ ($1\leq i\leq n$)를 $S$-준스킴들이라 하고, 각 첨자 $i$에
대하여 $x_i$를 $X_i$의 한 점이라 하자. $x_i$가
$Y=X_1\times_S X_2\times_S\cdots\times_S X_n$의 어떤 점 $y$의
$i$번째 사영이 될 필요충분조건은 모든 $x_i$가 $S$의 같은 한 점
$s$ 위에 있는 것이다.
\end{proposition}

이 조건이 필요함은 자명하고, 보조정리
(\hyperref[I.3.4.6-ko]{3.4.6})이 충분함을 보인다.

\phantomsection\label{I.3.4.7.underlying-map-ko}
다시 말해 $X$의 바탕집합을 $(X)$로 나타내면 표준적인 \emph{전사}
사상 $(X\times_S Y)\to (X)\times_{(S)}(Y)$이 있음을 알 수 있다.
이 사상은 일반적으로 \emph{단사이지 않음}에 유의해야 한다. 즉
\emph{$X\times_S Y$에서 서로 다른 여러 점 $z$가 같은 사영
$x\in X$, $y\in Y$를 가질 수 있다}. 실제로 $S$, $X$, $Y$가 각각
체 $k$, $K$, $K'$의 소 스펙트럼인 경우에도 이를 볼 수 있는데,
텐서곱 $K\otimes_k K'$은 일반적으로 서로 다른 여러 소 아이디얼을
갖기 때문이다(\hyperref[I.3.4.9-ko]{3.4.9} 참조).

\begin{corollary}[3.4.8]
\label{I.3.4.8-ko}
$f:X\to Y$를 $S$-사상, $f_{(S')}:X_{(S')}\to Y_{(S')}$를
밑준스킴의 확장 $S'\to S$로 $f$에서 얻는 $S'$-사상이라 하자.
$p$(각각 $q$)를 사영 $X_{(S')}\to X$(각각 $Y_{(S')}\to Y$)라
하자. $X$의 모든 부분집합 $M$에 대하여
\phantomsection\label{I.3.4.8.base-change-image-ko}
\[
  q^{-1}(f(M))=f_{(S')}(p^{-1}(M)).
\]
\end{corollary}

\oldpage[I]{114}
실제로 (\hyperref[I.3.3.11-ko]{3.3.11})에 따라 다음 가환 그림으로
$X_{(S')}$을 곱 $X\times_Y Y_{(S')}$과 동일시한다.
\phantomsection\label{I.3.4.8.proof-diagram-ko}
\[
  \xymatrix{
    X\ar[d]_f & X_{(S')}\ar[l]_p\ar[d]^{f_{(S')}}\\
    Y & Y_{(S')}\ar[l]^q
  }
\]

(\hyperref[I.3.4.7-ko]{3.4.7})에 따라 $x\in M$,
$y'\in Y_{(S')}$에 대한 관계 $q(y')=f(x)$은
$p(x')=x$이고 $f_{(S')}(x')=y'$인 어떤 $x'\in X_{(S')}$가
존재함과 동치이다. 이로부터 따름정리가 얻어진다.

보조정리 (\hyperref[I.3.4.6-ko]{3.4.6})을 다음과 같이 정밀화할 수
있다.

\begin{proposition}[3.4.9]
\label{I.3.4.9-ko}
$X$, $Y$를 두 $S$-준스킴, $x$를 $X$의 점, $y$를 $Y$의 점이라
하고, $x$와 $y$가 같은 점 $s\in S$ 위에 있다고 하자.
$X\times_S Y$에서 $x$와 $y$를 사영으로 갖는 점들의 집합은
$k(s)$의 확대체로 본 $k(x)$와 $k(y)$의 합성 확대체의 유형들의
집합과 표준적으로 전단사 대응한다(Bourbaki, \emph{Alg.}, chap. VIII,
\S~8, prop. 2).
\end{proposition}

$p$(각각 $q$)를 $X\times_S Y$에서 $X$(각각 $Y$)로 가는 사영이라
하고, $E$를 $X\times_S Y$의 바탕공간의 부분공간
$p^{-1}(x)\cap q^{-1}(y)$라 하자. 먼저 사상
$\operatorname{Spec}(k(x))\to S$와 $\operatorname{Spec}(k(y))\to S$가
$\operatorname{Spec}(k(x))\to\operatorname{Spec}(k(s))\to S$와
$\operatorname{Spec}(k(y))\to\operatorname{Spec}(k(s))\to S$로
각각 인수분해됨에 유의하자.
$\operatorname{Spec}(k(s))\to S$는 단사사상이므로
(\hyperref[I.2.4.7-ko]{2.4.7}),
(\hyperref[I.3.2.4-ko]{3.2.4})에 따라
\phantomsection\label{I.3.4.9.tensor-product-fibre-ko}
\[
  P=\operatorname{Spec}(k(x))\times_S\operatorname{Spec}(k(y))
  =\operatorname{Spec}(k(x))\times_{\operatorname{Spec}(k(s))}
   \operatorname{Spec}(k(y))
  =\operatorname{Spec}\bigl(k(x)\otimes_{k(s)}k(y)\bigr).
\]

서로 역인 두 사상 $\alpha:P_0\to E$, $\beta:E\to P_0$를 정의하자.
여기서 $P_0$는 준스킴 $P$의 바탕집합을 뜻한다.
$i:\operatorname{Spec}(k(x))\to X$와
$j:\operatorname{Spec}(k(y))\to Y$가 표준 사상들이면
(\hyperref[I.2.4.5-ko]{2.4.5}), $\alpha$를 사상 $i\times_S j$에
대응하는 바탕공간의 사상으로 잡는다. 한편 모든 $z\in E$는 가정에
따라 두 $k(s)$-단사 준동형 $k(x)\to k(z)$와 $k(y)\to k(z)$를
정의하고, 따라서 유도된 $k(s)$-대수 준동형
$k(x)\otimes_{k(s)}k(y)\to k(z)$와 이어서 사상
$\operatorname{Spec}(k(z))\to P$를 정의한다. $\beta(z)$는 이 사상에
의해 $P_0$에 정해지는 점으로 둔다. $\alpha\circ\beta$와
$\beta\circ\alpha$가 항등사상임은
(\hyperref[I.2.4.5-ko]{2.4.5})와 곱의 정의
(\hyperref[I.3.2.1-ko]{3.2.1})에서 따른다. 끝으로 $P_0$가
$k(x)$와 $k(y)$의 합성 확대체의 유형들의 집합과 전단사 대응함은
알려져 있다(Bourbaki, \emph{Alg.}, chap. VIII, \S~8, prop. 1).

\subsection{전사와 단사}
\label{subsection:I.3.5-ko}

\begin{env}[3.5.1]
\label{I.3.5.1-ko}
일반적으로 준스킴의 사상들이 갖는 한 성질 $\mathbf{P}$를 생각하고,
다음 두 명제를 생각하자.
\begin{enumerate}
  \item[(i)] $f:X\to X'$, $g:Y\to Y'$가 성질 $\mathbf{P}$를 갖는 두
  $S$-사상이면 $f\times_S g$도 성질 $\mathbf{P}$를 갖는다.
  \item[(ii)] $f:X\to Y$가 성질 $\mathbf{P}$를 갖는 $S$-사상이면,
  밑준스킴의 확장으로 $f$에서 얻는 모든 $S'$-사상
  $f_{(S')}:X_{(S')}\to Y_{(S')}$도 성질 $\mathbf{P}$를 갖는다.
\end{enumerate}

\oldpage[I]{115}
$f_{(S')}=f\times_S 1_{S'}$이므로 모든 준스킴 $X$에 대하여
\emph{항등사상} $1_X$가 성질 $\mathbf{P}$를 가지면 (i)이 (ii)를
함의함을 알 수 있다. 한편 $f\times_S g$는 합성사상
\phantomsection\label{I.3.5.1.product-composition-ko}
\[
  X\times_S Y\xrightarrow{f\times 1_Y}X'\times_S Y
  \xrightarrow{1_{X'}\times g}X'\times_S Y'
\]
이므로 성질 $\mathbf{P}$를 갖는 두 사상의 \emph{합성}도 이 성질을
가지면 (ii)가 (i)을 함의함을 알 수 있다.
\end{env}

이 관찰의 첫 응용은 다음 명제이다.

\begin{proposition}[3.5.2]
\label{I.3.5.2-ko}
\begin{enumerate}
  \item[(i)] $f:X\to X'$, $g:Y\to Y'$가 전사인 $S$-사상들이면
  $f\times_S g$도 전사이다.
  \item[(ii)] $f:X\to Y$가 전사인 $S$-사상이면 밑준스킴의 모든 확장
  $S'$에 대하여 $f_{(S')}$도 전사이다.
\end{enumerate}
\end{proposition}

두 전사의 합성은 전사이므로 (ii)만 증명하면 충분하다. 그런데 이
명제는 $M=X$에 (\hyperref[I.3.4.8-ko]{3.4.8})을 적용하면 곧바로
따른다.

\begin{proposition}[3.5.3]
\label{I.3.5.3-ko}
사상 $f:X\to Y$가 전사일 필요충분조건은 모든 체 $K$와 모든 사상
$\operatorname{Spec}(K)\to Y$에 대하여 $K$의 확대체 $K'$과 다음
그림을 가환이게 하는 사상 $\operatorname{Spec}(K')\to X$가 존재하는
것이다.
\phantomsection\label{I.3.5.3.surjectivity-diagram-ko}
\[
  \xymatrix{
    X\ar[d]_f & \operatorname{Spec}(K')\ar[l]\ar[d]\\
    Y & \operatorname{Spec}(K)\ar[l]
  }
\]
\end{proposition}

이 조건은 충분하다. 실제로 모든 $y\in Y$에 대하여 $K$가 $k(y)$의
확대체일 때 단사 준동형 $k(y)\to K$에 대응하는 사상
$\operatorname{Spec}(K)\to Y$에 이 조건을 적용하면 된다
(\hyperref[I.2.4.6-ko]{2.4.6}). 역으로 $f$가 전사라고 하고,
$y\in Y$를 $\operatorname{Spec}(K)$의 유일한 점의 상이라 하자.
$f(x)=y$인 $x\in X$가 존재한다. 이에 대응하는 단사 준동형
$k(y)\to k(x)$를 생각하자(\hyperref[I.2.2.1-ko]{2.2.1}).
$k(x)$와 $K$에서 $K'$으로 가는 $k(y)$-단사 준동형들이 존재하도록
$K'$을 $k(y)$의 확대체로 잡으면 충분하다(Bourbaki, \emph{Alg.},
chap. V, \S~4, prop. 2). 이때 $k(x)\to K'$에 대응하는 사상
$\operatorname{Spec}(K')\to X$가 요구한 조건을 만족한다.

(\hyperref[I.3.4.5-ko]{3.4.5})에서 도입한 말로 표현하면,
\emph{$K$에 값을 갖는 $Y$의 모든 기하적 점은 $K$의 어떤 확대체에
값을 갖는 $X$의 기하적 점에서 온다}.

\begin{definition}[3.5.4]
\label{I.3.5.4-ko}
준스킴의 사상 $f:X\to Y$가 모든 체 $K$에 대하여 대응하는 사상
$X(K)\to Y(K)$을 단사로 만들면, $f$를 \emph{보편적으로 단사}라고
하거나 \emph{라디시엘 사상}이라고 한다.
\end{definition}

정의에서 곧바로 준스킴의 모든 \emph{단사사상}(T, 1.1)은 라디시엘
사상임을 알 수 있다.

\begin{env}[3.5.5]
\label{I.3.5.5-ko}
사상 $f:X\to Y$가 라디시엘 사상임을 보이려면 정의
(\hyperref[I.3.5.4-ko]{3.5.4})의 조건을 모든 \emph{대수적으로
닫힌} 체에 대하여 확인하는 것으로 충분하다. 실제로 $K$가 임의의
체이고 $K'$이 $K$의 대수적으로 닫힌 확대체이면 그림
\phantomsection\label{I.3.5.5.algebraic-closure-diagram-ko}
\[
  \xymatrix{
    X(K)\ar[r]^\alpha\ar[d]_\varphi & Y(K)\ar[d]^{\varphi'}\\
    X(K')\ar[r]_{\alpha'} & Y(K')
  }
\]

\oldpage[I]{116}
은 가환이다. 여기서 $\varphi$와 $\varphi'$은 사상
$\operatorname{Spec}(K')\to\operatorname{Spec}(K)$에서 오고,
$\alpha$와 $\alpha'$은 $f$에 대응한다. 그런데 $\varphi$는 단사이고
가정에 따라 $\alpha'$도 단사이므로 $\alpha$도 반드시 단사이다.
\end{env}

\begin{proposition}[3.5.6]
\label{I.3.5.6-ko}
$f:X\to Y$, $g:Y\to Z$를 준스킴의 두 사상이라 하자.
\begin{enumerate}
  \item[(i)] $f$와 $g$가 라디시엘 사상이면 $g\circ f$도
  라디시엘 사상이다.
  \item[(ii)] 역으로 $g\circ f$가 라디시엘 사상이면 $f$도
  라디시엘 사상이다.
\end{enumerate}
\end{proposition}

정의 (\hyperref[I.3.5.4-ko]{3.5.4})를 고려하면 이 명제는 사상들
$X(K)\to Y(K)\to Z(K)$에 대한 대응하는 자명한 명제들로 귀착된다.

\begin{proposition}[3.5.7]
\label{I.3.5.7-ko}
\begin{enumerate}
  \item[(i)] $S$-사상 $f:X\to X'$, $g:Y\to Y'$가 라디시엘
  사상들이면 $f\times_S g$도 라디시엘 사상이다.
  \item[(ii)] $S$-사상 $f:X\to Y$가 라디시엘 사상이면 밑준스킴의
  모든 확장 $S'\to S$에 대하여
  $f_{(S')}:X_{(S')}\to Y_{(S')}$도 라디시엘 사상이다.
\end{enumerate}
\end{proposition}

(\hyperref[I.3.5.1-ko]{3.5.1})에 따라 (i)만 증명하면 충분하다.
(\hyperref[I.3.4.2.1-ko]{3.4.2.1})에서
\phantomsection\label{I.3.5.7.field-valued-product-ko}
\[
  (X\times_S Y)(K)=X(K)\times_{S(K)}Y(K),\qquad
  (X'\times_S Y')(K)=X'(K)\times_{S(K)}Y'(K)
\]
임을 보았다. 이때 $f\times_S g$에 대응하는 사상
$(X\times_S Y)(K)\to(X'\times_S Y')(K)$은
$(u,v)\mapsto(f\circ u,g\circ v)$와 동일시되므로 명제가 곧바로
따른다.

\begin{proposition}[3.5.8]
\label{I.3.5.8-ko}
사상 $f=(\psi,\theta):X\to Y$가 라디시엘 사상일 필요충분조건은
$\psi$가 단사이고, 모든 $x\in X$에 대하여 단사 준동형
$\theta^x:k(\psi(x))\to k(x)$가 $k(x)$을 $k(\psi(x))$의 순수 비분리
확대체로 만드는 것이다.
\end{proposition}

$f$가 라디시엘 사상이라고 하고, 먼저 관계
$\psi(x_1)=\psi(x_2)=y$가 반드시 $x_1=x_2$를 함의함을 보이자.
실제로 $k(y)$의 확대체인 어떤 체 $K$와 $k(y)$-단사 준동형
$k(x_1)\to K$, $k(x_2)\to K$가 존재한다(Bourbaki, \emph{Alg.},
chap. V, \S~4, prop. 2). 따라서 이에 대응하는 두 사상
$u_1:\operatorname{Spec}(K)\to X$,
$u_2:\operatorname{Spec}(K)\to X$는 $f\circ u_1=f\circ u_2$를
만족한다. 가정에 따라 $u_1=u_2$이고, 특히 $x_1=x_2$이다. 이제
$\theta^x$를 통하여 $k(x)$을 $k(\psi(x))$의 확대체로 보자. $k(x)$이
$k(\psi(x))$의 순수 비분리 확대체가 아니면 $k(\psi(x))$의 어떤
대수적으로 닫힌 확대체 $K$ 안으로 가는 서로 다른 두
$k(\psi(x))$-단사 준동형 $k(x)\to K$가 존재하고, 이에 대응하는 두
사상 $\operatorname{Spec}(K)\to X$는 가정에 모순된다. 역으로
(\hyperref[I.2.4.6-ko]{2.4.6})을 고려하면 명제의 조건들이 $f$가
라디시엘 사상일 충분조건임은 곧바로 알 수 있다.

\begin{corollary}[3.5.9]
\label{I.3.5.9-ko}
$A$를 환, $S$를 $A$의 곱셈적 부분집합이라 하면 표준 사상
$\operatorname{Spec}(S^{-1}A)\to\operatorname{Spec}(A)$는 라디시엘
사상이다.
\end{corollary}

실제로 이 사상은 단사사상이다(\hyperref[I.1.6.2-ko]{1.6.2}).

\begin{corollary}[3.5.10]
\label{I.3.5.10-ko}
$f:X\to Y$를 라디시엘 사상, $g:Y'\to Y$를 사상이라 하고
$X'=X_{(Y')}=X\times_Y Y'$라 하자. 그러면 라디시엘 사상
$f_{(Y')}$ (\hyperref[I.3.5.7-ko]{3.5.7 (ii)})은 $X'$의 바탕공간을
$g^{-1}(f(X))$에 전단사로 대응시킨다. 또한 모든 체 $K$에 대하여
$X'(K)$는 $g$에 대응하는 사상 $Y'(K)\to Y(K)$에 의한, $Y(K)$의
부분집합 $X(K)$의 역상인 $Y'(K)$의 부분집합과 동일시된다.
\end{corollary}

첫째 명제는 (\hyperref[I.3.5.8-ko]{3.5.8})과
(\hyperref[I.3.4.8-ko]{3.4.8})에서 곧바로 따르고, 둘째 명제는
다음 그림의 가환성에서 따른다.
\oldpage[I]{117}
\phantomsection\label{I.3.5.10.base-change-points-diagram-ko}
\[
  \xymatrix{
    X'(K)\ar[r]\ar[d] & Y'(K)\ar[d]\\
    X(K)\ar[r] & Y(K)
  }
\]

\begin{remark}[3.5.11]
\label{I.3.5.11-ko}
준스킴의 사상 $f=(\psi,\theta)$에서 바탕공간의 사상 $\psi$가
단사이면 $f$를
\emph{단사}라고 하겠다. 사상 $f=(\psi,\theta):X\to Y$가 라디시엘
사상일 필요충분조건은 모든 사상 $Y'\to Y$에 대하여 사상
$f_{(Y')}:X_{(Y')}\to Y'$가 단사인 것이다. 이것이 \emph{보편적으로
단사인 사상}이라는 용어를 정당화한다. 실제로 조건의 필요성은
(\hyperref[I.3.5.7-ko]{3.5.7, (ii)})와
(\hyperref[I.3.5.8-ko]{3.5.8})에서 따른다. 역으로 이 조건은 먼저
$\psi$가 단사임을 함의한다. 어떤 $x\in X$에 대하여 단사 준동형
$\theta^x:k(\psi(x))\to k(x)$가 $k(x)$을 $k(\psi(x))$의 순수
비분리 확대체로 만들지 않는다면
$k(\psi(x))$의 어떤 확대체 $K$와 같은 사상
$\operatorname{Spec}(K)\to Y$에 대응하는 서로 다른 두 사상
$\operatorname{Spec}(K)\to X$가 존재할 것이다
(\hyperref[I.3.5.8-ko]{3.5.8}). 그러나 $Y'=\operatorname{Spec}(K)$로
두면 $X_{(Y')}$의 서로 다른 두 $Y'$-단면이 생긴다
(\hyperref[I.3.3.14-ko]{3.3.14}). 이는 $f_{(Y')}$가 단사라는
가정에 모순된다.
\end{remark}

\subsection{섬유}
\label{subsection:I.3.6-ko}

\begin{proposition}[3.6.1]
\label{I.3.6.1-ko}
$f:X\to Y$를 사상, $y$를 $Y$의 점, $\mathfrak{a}_y$를
$\mathscr{O}_y$의 $\mathfrak{m}_y$-준아딕 위상에 대한 정의
아이디얼이라 하자. 사영
$p:X\times_Y\operatorname{Spec}(\mathscr{O}_y/\mathfrak{a}_y)\to X$는
준스킴 $X\times_Y\operatorname{Spec}(\mathscr{O}_y/\mathfrak{a}_y)$의
바탕공간을, $X$의 바탕공간 위상에서 유도된 위상을 준 섬유
$f^{-1}(y)$ 위로 위상동형시킨다.
\end{proposition}

$\operatorname{Spec}(\mathscr{O}_y/\mathfrak{a}_y)\to Y$는 라디시엘
사상이고(\hyperref[I.3.5.4-ko]{3.5.4} 및
\hyperref[I.2.4.7-ko]{2.4.7}), 아이디얼
$\mathfrak{m}_y/\mathfrak{a}_y$가 가정에 따라 멱영이므로
(\hyperref[I.1.1.12-ko]{1.1.12})
$\operatorname{Spec}(\mathscr{O}_y/\mathfrak{a}_y)$에는 점이 하나뿐이다.
따라서 이미 (\hyperref[I.3.5.10-ko]{3.5.10} 및
\hyperref[I.3.3.4-ko]{3.3.4})에서 $p$가
$X\times_Y\operatorname{Spec}(\mathscr{O}_y/\mathfrak{a}_y)$의
바탕공간을 $f^{-1}(y)$와 \emph{집합으로서} 동일시함을 알고 있다.
그러므로 $p$가 위상동형임을 보이면 충분하다.
(\hyperref[I.3.2.7-ko]{3.2.7})에 따라 이 문제는 $X$와 $Y$ 위에서
국소적이므로 $X=\operatorname{Spec}(B)$,
$Y=\operatorname{Spec}(A)$이고 $B$가 $A$-대수라고 가정할 수 있다.
이때 사상 $p$는 준동형
$1\otimes\varphi:B\to B\otimes_A A'$에 대응한다. 여기서
$A'=A_y/\mathfrak{a}_y$이고 $\varphi$는 $A$에서 $A'$으로 가는 표준
사상이다. 그런데 $B\otimes_A A'$의 모든 원소는
\phantomsection\label{I.3.6.1.localization-fraction-ko}
\[
  \sum_i b_i\otimes\varphi(a_i)/\varphi(s)
  =\left(\sum_i(a_i b_i\otimes 1)\right)(1\otimes\varphi(s))^{-1}
\]
꼴로 쓸 수 있다. 여기서 $s\notin\mathfrak{j}_y$이며 명제
(\hyperref[I.1.2.4-ko]{1.2.4})를 적용할 수 있다.

\begin{env}[3.6.2]
\label{I.3.6.2-ko}
이 책의 앞으로의 모든 부분에서 사상의 섬유 $f^{-1}(y)$에
$k(y)$-준스킴의 구조가 주어졌다고 할 때에는,
\emph{언제나 $X\times_Y\operatorname{Spec}(k(y))$의 구조를 $X$로의
사영을 통하여 옮겨 얻은 준스킴을 뜻한다}. 이 마지막 곱을
$X\otimes_Y k(y)$ 또는 $X\otimes_{\mathscr{O}_Y}k(y)$라고도 쓰겠다.
더 일반적으로 $B$가 $\mathscr{O}_y$-대수이면 곱
$X\times_Y\operatorname{Spec}(B)$를 $X\otimes_Y B$ 또는
$X\otimes_{\mathscr{O}_Y}B$로 나타내겠다.
\end{env}

앞의 규약 아래에서 (\hyperref[I.3.5.10-ko]{3.5.10})에 따라
$k(y)$의 확대체 $K$에 값을 갖는 $X$의 점들은
\emph{$K$에 값을 갖는 $f^{-1}(y)$의 점들}과 동일시된다.

\begin{env}[3.6.3]
\label{I.3.6.3-ko}
$f:X\to Y$, $g:Y\to Z$를 두 사상, $h=g\circ f$를 그 합성이라 하자.
모든 $z\in Z$에 대하여 섬유 $h^{-1}(z)$는 다음 준스킴과 동형이다.
\phantomsection\label{I.3.6.3.fibre-composition-ko}
\[
  X\times_Z\operatorname{Spec}(k(z))
  =(X\times_Y Y)\times_Z\operatorname{Spec}(k(z))
  =X\times_Y g^{-1}(z)
\]
\oldpage[I]{118}
특히 $U$가 $X$의 열린부분이면, 섬유 준스킴 $f^{-1}(y)$가
$U\cap f^{-1}(y)$ 위에 유도하는 준스킴은 $f_U^{-1}(y)$와 동형이다.
여기서 $f_U$는 $f$의 $U$로의 제한이다.
\end{env}

\begin{proposition}[3.6.4) (``섬유의 추이성'']
\label{I.3.6.4-ko}
$f:X\to Y$, $g:Y'\to Y$를 두 사상이라 하자. $X'=X\times_Y Y'
=X_{(Y')}$, $f'=f_{(Y')}:X'\to Y'$로 놓자. 모든 $y'\in Y'$에
대하여 $y=g(y')$로 놓으면 준스킴 $f'^{-1}(y')$는
$f^{-1}(y)\otimes_{k(y)}k(y')$와 동형이다.
\end{proposition}

실제로 두 준스킴
$(X\otimes_Y k(y))\otimes_{k(y)}k(y')$와
$(X\times_Y Y')\otimes_{Y'}k(y')$가 모두
(\hyperref[I.3.3.9.1-ko]{3.3.9.1})에 따라
$X\times_Y\operatorname{Spec}(k(y'))$와 표준적으로 동형임을
관찰하면 된다.

특히 $V$가 $Y$에서 $y$의 열린 근방이고, $f_V$가 $f$를
$f^{-1}(V)$ 위에 유도된 준스킴에 제한한 사상이라면 준스킴 $f^{-1}(y)$와
$f_V^{-1}(y)$는 표준적으로 동일시된다.

\begin{proposition}[3.6.5]
\label{I.3.6.5-ko}
$f:X\to Y$를 사상, $y$를 $Y$의 점,
$Z=\operatorname{Spec}(\mathscr{O}_y)$를 국소 스킴,
$p=(\psi,\theta):X\times_Y Z\to X$를 사영이라 하자. $Z$의
바탕공간을 $Y$의 부분공간과 동일시하면
(\hyperref[I.2.4.2-ko]{2.4.2}), $p$는 $X\times_Y Z$의 바탕공간을
$X$의 부분공간 $f^{-1}(Z)$ 위로 위상동형시킨다. 또한 모든
$t\in X\times_Y Z$에 대하여 $z=\psi(t)$로 놓으면
$\theta_t^\sharp$은 $\mathscr{O}_z$에서 $\mathscr{O}_t$로 가는
동형이다.
\end{proposition}

$Z$는 $Y$의 부분공간으로 동일시했을 때 $y$를 포함하는 모든 아핀
열린집합에 포함되므로(\hyperref[I.2.4.2-ko]{2.4.2}),
(\hyperref[I.3.6.1-ko]{3.6.1})에서와 같이
$X=\operatorname{Spec}(A)$와 $Y=\operatorname{Spec}(B)$가 아핀
스킴이고 $A$가 $B$-대수인 경우로 귀착시킬 수 있다. 그러면
$X\times_Y Z$는 $A\otimes_B B_y$의 소 스펙트럼이고, 이 환은
$S^{-1}A$와 표준적으로 동일시된다. 여기서 $S$는
$B-\mathfrak{j}_y$의 $A$ 안에서의 상이다(0, 1.5.2). 이때 $p$는
표준 준동형 $A\to S^{-1}A$에 대응하므로 명제는
(\hyperref[I.1.6.2-ko]{1.6.2})에서 따른다.

\subsection{응용: 준스킴의 mod.
\texorpdfstring{$\mathfrak{J}$ 환원\protect\footnotemark}{J 환원}}
\label{subsection:I.3.7-ko}
\label{I.3.7-ko}
\footnotetext{이 항은 제I장의 뒷부분과 제II장의 개념과 결과를 사용하지만
뒤에서는 쓰지 않을 것이며, 고전 대수기하학에 익숙한 독자만을 위한 것이다.}

\begin{env}[3.7.1]
\label{I.3.7.1-ko}
$A$를 환, $X$를 $A$-준스킴, $\mathfrak{J}$를 $A$의 아이디얼이라 하자.
그러면 $X_0=X\otimes_A(A/\mathfrak{J})$는
$(A/\mathfrak{J})$-준스킴이며, 때로는 $X$에서 mod. $\mathfrak{J}$
\emph{환원}으로 얻어진다고 한다.
\end{env}

\begin{env}[3.7.2]
\label{I.3.7.2-ko}
이 용어는 주로 $A$가 \emph{국소환}이고 $\mathfrak{J}$가 그 극대
아이디얼인 경우에 쓰인다. 이때 $X_0$는 $A$의 잉여체
$k=A/\mathfrak{J}$ 위의 준스킴이다.

더 나아가 $A$가 분수체 $K$를 갖는 정역이면 $K$-준스킴
$X'=X\otimes_A K$도 생각할 수 있다. 우리가 사용하지 않을 용어의
남용으로서, 종래에는 흔히 $X_0$가 mod. $\mathfrak{J}$ 환원으로
\emph{$X'$에서 얻어진다}고 하였다. 이 말을 쓰던 경우에는 $A$가
차원 1인 국소환이었고(대개 이산 값매김환이었다), 주어진
$K$-준스킴 $X'$가 어떤 $K$-준스킴 $P'$의 닫힌 부분준스킴이라는
사실을 다소 명시적으로 전제하였다. 실제로 $P'$은
$\mathbf{P}_K^r$형의 사영공간이었고(II, 4.1.1 참조),
$P'=P\otimes_A K$ 꼴이었다. 여기서 $P$는 주어진 $A$-준스킴,
실제로는 II, 4.1.1의 표기에서 $A$-스킴 $\mathbf{P}_A^r$이었다.
우리의 언어로 $X'$에서 $X_0$를 정의하는 방법은 다음과 같다.

아핀 스킴 $Y=\operatorname{Spec}(A)$를 생각하자. 이는 유일한
\oldpage[I]{119}
닫힌점 $y=\mathfrak{J}$와 일반점 $(0)$의 두 점으로 이루어진다. 따라서 일반점
하나로 이루어진 집합 $U$는 $Y$ 안의 열린집합
$U=\operatorname{Spec}(K)$이다.
$X$가 $A$-준스킴, 다시 말해 $Y$-준스킴이고 $\psi:X\to Y$가 구조
사상이면 $X\otimes_A K=X'$은 $X$가 $\psi^{-1}(U)$ 위에 유도하는
준스킴과 다름없다. 특히 $\varphi:P\to Y$가 구조 사상이면
$P'=\varphi^{-1}(U)$의 닫힌 부분준스킴 $X'$은 $P$의 (국소 닫힌)
부분준스킴이다. $P$가 뇌터이면(예를 들어 $A$가 뇌터이고 $P$가
$A$ 위에서 유한형이면), $X'$을 포함하는 $P$의 가장 작은 닫힌
부분준스킴 $X=\overline{X'}$가 존재한다(9.5.10). 또한 $X'$은
$X$가 열린집합 $\varphi^{-1}(U)\cap X$ 위에 유도하는 준스킴이므로
$X\otimes_A K$와 동형이다(9.5.10).
\emph{따라서 $X'$의 $P'=P\otimes_A K$ 안으로의 몰입은 표준적인
방식으로 $X'$를 $X'=X\otimes_A K$ 꼴로 볼 수 있게 한다. 여기서
$X$는 $A$-준스킴이다.} 이제 mod. $\mathfrak{J}$로 환원한 준스킴
$X_0=X\otimes_A k$를 생각할 수 있으며, 이는 닫힌점 $y$의 섬유
$\psi^{-1}(y)$와 다름없다. 종래에는 적절한 용어가 없어서
$A$-준스킴 $X$를 명시적으로 도입하기를 피하였다. 그러나 통상
``mod. $\mathfrak{J}$로 환원한'' 준스킴 $X_0$에 관하여 말하는 모든
명제는 $X$ 자체에 관한 더 완전한 명제들의 귀결로 보아야 하며,
그렇게 해석해야만 만족스럽게 정식화하고 이해할 수 있다. 더구나
가정들은 언제나 $X'$의 $\mathbf{P}_K^r$ 안으로의 몰입을 미리 주는
일과는 무관하게 $X$ 자체에 관한 가정들로 귀착되는 듯하다. 이로써
더 내재적인 명제들을 얻을 수 있다.
\end{env}

\begin{env}[3.7.3]
\label{I.3.7.3-ko}
끝으로 여기서 생각한 상황의 개념적 정리가 늦어진 까닭이 되었을 법한
매우 특수한 사실을 지적하자. $A$가 이산 값매김환이고 $X$가
\emph{$A$ 위에서 고유}이면($X$가 $\mathbf{P}_A^r$의 닫힌
부분준스킴일 때에는 실제로 그러하다. II, 5.5.4 참조), $A$에 값을 갖는
$X$의 점들과 $K$에 값을 갖는 $X'$의 점들은 전단사 대응한다
(II, 7.3.8). 이 때문에 실제로는 $X$에 관한 명제들을 증명하면서도
흔히 $X'$에 관한 결과를 증명한다고 여겨 왔다. 그 명제들은 이 형태로
바탕 국소환의 차원이 1이라는 가정을 없애도 여전히 성립한다.
\end{env}

% Authority: EGA I ega1-4-fr.tex, whole-file SHA-256 A5F5CDAC81E654E9ABA75BBFEAD24F3010B122452C689B9EBAB5A7711B454EC1.
% Sequential coverage in this file: authority lines 1-395, complete subsections 4.1-4.2.
\section{부분준스킴과 몰입 사상}
\label{section:I.4-ko}

\subsection{부분준스킴}
\label{subsection:I.4.1-ko}

\begin{env}[4.1.1]
\label{I.4.1.1-ko}
준연접층이라는 개념은 국소적이므로
(\hyperref[0.5.1.3-ko]{0, 5.1.3}), 준스킴 $X$ 위의 준연접
$\mathscr{O}_X$-가군 $\mathscr{F}$는 다음 조건으로 정의할 수 있다.
$X$의 모든 아핀 열린집합 $V$에 대하여 $\mathscr{F}|V$는 어떤
$\Gamma(V,\mathscr{O}_X)$-가군에 대응하는 층과 동형이다
(\hyperref[I.1.4.1-ko]{1.4.1}). 준스킴 $X$ 위에서 구조층
$\mathscr{O}_X$가 준연접이고, 준연접 $\mathscr{O}_X$-가군 사이의
준동형의 핵, 여핵, 상과 준연접 $\mathscr{O}_X$-가군들의 귀납극한과
직합이 모두 준연접임은 자명하다
(\hyperref[I.1.3.7-ko]{1.3.7} 및
\hyperref[I.1.3.9-ko]{1.3.9}).
\end{env}

\begin{proposition}[4.1.2]
\label{I.4.1.2-ko}
$X$를 준스킴, $\mathscr{I}$를 $\mathscr{O}_X$의 준연접 아이디얼층이라
하자. 그러면 층 $\mathscr{O}_X/\mathscr{I}$의 지지집합 $Y$는 닫혀
있다. $\mathscr{O}_Y$를 $\mathscr{O}_X/\mathscr{I}$의 $Y$로의
제한이라 하면 $(Y,\mathscr{O}_Y)$는 준스킴이다.
\end{proposition}

\oldpage[I]{120}
(\hyperref[I.2.1.3-ko]{2.1.3})에 따라 $X$가 아핀 스킴인 경우만
생각하고, 이 경우 $Y$가 $X$에서 닫혀 있으며 \emph{아핀 스킴}임을
보이면 충분하다. 실제로 $X=\operatorname{Spec}(A)$이면
$\mathscr{O}_X=\widetilde{A}$이고
$\mathscr{I}=\widetilde{\mathfrak{I}}$이다. 여기서 $\mathfrak{I}$는
$A$의 아이디얼이다(\hyperref[I.1.4.1-ko]{1.4.1}). 이때 $Y$는
$X$의 닫힌 부분집합 $V(\mathfrak{I})$와 같고, 환
$B=A/\mathfrak{I}$의 소 스펙트럼과 동일시된다
(\hyperref[I.1.1.11-ko]{1.1.11}). 또한 $\varphi$가 표준 준동형
$A\to B=A/\mathfrak{I}$이면, 직상
${}^a\varphi_*(\widetilde{B})$는 층
$\widetilde{A}/\widetilde{\mathfrak{I}}
=\mathscr{O}_X/\mathscr{I}$와 표준적으로 동일시된다
(\hyperref[I.1.6.3-ko]{1.6.3} 및
\hyperref[I.1.3.9-ko]{1.3.9}). 이로써 증명이 끝난다.

$(Y,\mathscr{O}_Y)$를 \emph{아이디얼층 $\mathscr{I}$에 의해 정의되는}
$(X,\mathscr{O}_X)$의 \emph{부분준스킴}이라 한다. 이는 다음과 같은
일반적인 \emph{부분준스킴} 개념의 특수한 경우이다.

\begin{definition}[4.1.3]
\label{I.4.1.3-ko}
환 달린 공간 $(Y,\mathscr{O}_Y)$가 다음 조건들을 만족하면 이를
준스킴 $(X,\mathscr{O}_X)$의 \emph{부분준스킴}이라 한다.
\begin{enumerate}
  \item[$1^\circ$] $Y$는 $X$의 국소 닫힌 부분공간이다.
  \item[$2^\circ$] $U$를 $Y$를 포함하며 그 안에서 $Y$가 닫혀 있는
  $X$의 가장 큰 열린집합이라 하자. 다시 말해 $U$는 $X$에서
  $\overline{Y}$ 안에서의 $Y$의 경계를 뺀 집합이다. 그러면
  $(Y,\mathscr{O}_Y)$는 $\mathscr{O}_X|U$의 어떤 준연접 아이디얼층에
  의해 정의되는 $(U,\mathscr{O}_X|U)$의 닫힌 부분준스킴이다.
\end{enumerate}
$X$의 부분준스킴 $(Y,\mathscr{O}_Y)$에서 $Y$가 $X$ 안에서 닫혀
있으면 $(Y,\mathscr{O}_Y)$를 $X$의 \emph{닫힌} 부분준스킴이라 한다.
이 경우 $U=X$이다.
\end{definition}

이 정의와 (\hyperref[I.4.1.2-ko]{4.1.2})에서 $X$의 닫힌
부분준스킴들과 $\mathscr{O}_X$의 준연접 아이디얼층
$\mathscr{J}$들 사이에 \emph{표준적인 전단사 대응}이 있음을 곧바로
알 수 있다. 실제로 그러한 두 층 $\mathscr{J}$, $\mathscr{J}'$의
지지집합이 같은 닫힌집합 $Y$이고, $\mathscr{O}_X/\mathscr{J}$와
$\mathscr{O}_X/\mathscr{J}'$의 $Y$로의 제한이 서로 같으면 곧
$\mathscr{J}'=\mathscr{J}$이다.

\begin{env}[4.1.4]
\label{I.4.1.4-ko}
$(Y,\mathscr{O}_Y)$를 $X$의 부분준스킴, $U$를 $Y$를 포함하며 그
안에서 $Y$가 닫혀 있는 $X$의 가장 큰 열린집합, $V$를 $U$에
포함되는 $X$의 열린집합이라 하자. 그러면 $V\cap Y$는 $V$ 안에서
닫혀 있다. 또한 $Y$가 $\mathscr{O}_X|U$의 준연접 아이디얼층
$\mathscr{J}$에 의해 정의되면, $\mathscr{J}|V$는
$\mathscr{O}_X|V$의 준연접 아이디얼층이고, $Y$가 $Y\cap V$ 위에
유도하는 준스킴은 아이디얼층 $\mathscr{J}|V$에 의해 정의되는
$V$의 닫힌 부분준스킴임이 자명하다. 역으로 다음이 성립한다.
\end{env}

\begin{proposition}[4.1.5]
\label{I.4.1.5-ko}
$(Y,\mathscr{O}_Y)$를 환 달린 공간이라 하자. $Y$가 $X$의 부분공간이고,
$Y$를 덮는 $X$의 열린집합족 $(V_\alpha)$가 존재하여 모든
$\alpha$에 대하여 $Y\cap V_\alpha$가 $V_\alpha$ 안에서 닫혀 있으며
환 달린 공간
$(Y\cap V_\alpha,\mathscr{O}_Y|(Y\cap V_\alpha))$가 $X$가
$V_\alpha$ 위에 유도하는 준스킴의 닫힌 부분준스킴이라고 하자.
그러면 $(Y,\mathscr{O}_Y)$는 $X$의 부분준스킴이다.
\end{proposition}

가정에 따라 $Y$는 $X$ 안에서 국소 닫혀 있다. 또한 $Y$를 포함하며
그 안에서 $Y$가 닫혀 있는 가장 큰 열린집합 $U$는 모든
$V_\alpha$를 포함한다. 따라서 $U=X$이고 $Y$가 $X$ 안에서 닫힌
경우로 귀착할 수 있다. 이제 $\mathscr{O}_X$의 준연접 아이디얼층
$\mathscr{J}$를 다음과 같이 정의한다. $\mathscr{J}|V_\alpha$는 닫힌
부분준스킴
$(Y\cap V_\alpha,\mathscr{O}_Y|(Y\cap V_\alpha))$를 정의하는
$\mathscr{O}_X|V_\alpha$의 아이디얼층으로 잡고, $Y$와 만나지 않는
$X$의 모든 열린집합 $W$에 대해서는
$\mathscr{J}|W=\mathscr{O}_X|W$로 잡는다.
(\hyperref[I.4.1.3-ko]{4.1.3})과
(\hyperref[I.4.1.4-ko]{4.1.4})에 따라 이 조건들을 만족하는
아이디얼층 $\mathscr{J}$가 하나뿐임을 곧바로 확인할 수 있으며, 이
층은 닫힌 부분준스킴 $(Y,\mathscr{O}_Y)$를 정의한다.

특히 $X$가 $X$의 어떤 \emph{열린집합} 위에 \emph{유도하는}
준스킴은 $X$의 \emph{부분준스킴}이다.

\begin{proposition}[4.1.6]
\label{I.4.1.6-ko}
$X$의 부분\oldpage[I]{121}준스킴(각각 닫힌 부분준스킴)의
부분준스킴(각각 닫힌 부분준스킴)은 $X$의 부분준스킴(각각 닫힌
부분준스킴)과 표준적으로 동일시된다.
\end{proposition}

$X$의 국소 닫힌 부분공간의 국소 닫힌 부분집합은 $X$의 국소 닫힌
부분공간이므로, 문제가 국소적임은 자명하다
(\hyperref[I.4.1.5-ko]{4.1.5}). 따라서 $X$가 아핀이라고 가정해도
된다. 이제 명제는 환 $A$의 두 아이디얼 $\mathfrak{J}$,
$\mathfrak{J}'$가 $\mathfrak{J}\subset\mathfrak{J}'$을 만족할 때
$A/\mathfrak{J}'$와
$(A/\mathfrak{J})/(\mathfrak{J}'/\mathfrak{J})$를 표준적으로
동일시할 수 있다는 사실에서 따른다.

이후에는 언제나 앞의 동일시를 사용한다.

\begin{env}[4.1.7]
\label{I.4.1.7-ko}
$Y$를 준스킴 $X$의 부분준스킴이라 하고, $Y$의 \emph{바탕공간}에서
$X$의 바탕공간으로 가는 표준 단사사상을 $\psi:Y\to X$라 하자.
역상 $\psi^*(\mathscr{O}_X)$는 제한 $\mathscr{O}_X|Y$임이 알려져
있다(\hyperref[0.3.7.1-ko]{0, 3.7.1}). 모든 $y\in Y$에 대하여
$\omega_y$를 표준 준동형
$(\mathscr{O}_X)_y\to(\mathscr{O}_Y)_y$라 하자. 이 준동형들은 환의
층 사이의 \emph{전사} 준동형
$\omega:\mathscr{O}_X|Y\to\mathscr{O}_Y$가 줄기들에 유도하는
준동형들이다. 실제로 이를 $Y$ 위에서 국소적으로 확인하면 충분하므로,
$X$가 아핀이고 부분준스킴 $Y$가 닫힌 경우만 생각해도 된다. 이 경우
$\mathscr{I}$가 $Y$를 정의하는 $\mathscr{O}_X$의 아이디얼층이면,
$\omega_y$들은 준동형
$\mathscr{O}_X|Y\to(\mathscr{O}_X/\mathscr{I})|Y$가 줄기들에
유도하는 준동형들이다.

따라서 \emph{환 달린 공간의 단사사상}
$j=(\psi,\omega^b)$를 정의하였다
(\hyperref[0.4.1.1-ko]{0, 4.1.1}). 이는 준스킴의 사상
$Y\to X$임이 자명하며(\hyperref[I.2.2.1-ko]{2.2.1}), 이를
\emph{표준 포함 사상}이라 한다.

$f:X\to Z$가 사상이면 합성사상
$Y\xrightarrow{j}X\xrightarrow{f}Z$도 부분준스킴 $Y$에 대한
$f$의 \emph{제한}이라 한다.
\end{env}

\begin{env}[4.1.8]
\label{I.4.1.8-ko}
일반적인 정의(T, I, 1.1)에 따라, 준스킴의 사상 $f:Z\to X$가
$X$의 부분준스킴 $Y$의 포함 사상 $j:Y\to X$를 \emph{통하여
인수분해된다}고 함은
$f$가 $Z\xrightarrow{g}Y\xrightarrow{j}X$로 인수분해된다는 뜻이다.
여기서 $g$는 준스킴의 사상이며, $j$가 단사사상이므로 $g$는 반드시
\emph{유일}하다.
\end{env}

\begin{proposition}[4.1.9]
\label{I.4.1.9-ko}
사상 $f:Z\to X$가 \emph{포함 사상 $j:Y\to X$를 통하여
인수분해되기} 위한 필요충분조건은 다음과 같다. $f(Z)\subset Y$이고,
모든 $z\in Z$에 대하여 $y=f(z)$로 놓으면 $f$에 대응하는 준동형
$(\mathscr{O}_X)_y\to\mathscr{O}_z$가
$(\mathscr{O}_X)_y\to(\mathscr{O}_Y)_y\to\mathscr{O}_z$로
인수분해된다. 동치로, $(\mathscr{O}_X)_y\to\mathscr{O}_z$의 핵은
$(\mathscr{O}_X)_y\to(\mathscr{O}_Y)_y$의 핵을 포함한다.
\end{proposition}

이 조건들이 필요함은 자명하다. 충분함을 보이기 위해, 필요하다면
$X$를 $Y$가 그 안에서 닫히는 열린집합 $U$로 바꾸어
(\hyperref[I.4.1.3-ko]{4.1.3}) $Y$가 $X$의 \emph{닫힌}
부분준스킴인 경우로 귀착할 수 있다. 이때 $Y$는
$\mathscr{O}_X$의 준연접 아이디얼층 $\mathscr{I}$에 의해 정의된다.
$f=(\psi,\theta)$로 놓고, $\mathscr{J}$를
$\theta^\sharp:\psi^*(\mathscr{O}_X)\to\mathscr{O}_Z$의 핵인
$\psi^*(\mathscr{O}_X)$의 아이디얼층이라 하자. 함자 $\psi^*$의
성질(\hyperref[0.3.7.2-ko]{0, 3.7.2})을 고려하면, 가정에 따라 모든
$z\in Z$에 대하여
$(\psi^*(\mathscr{I}))_z\subset\mathscr{J}_z$이고, 따라서
$\psi^*(\mathscr{I})\subset\mathscr{J}$이다. 그러므로
$\theta^\sharp$은 다음과 같이 인수분해된다.
\[
  \psi^*(\mathscr{O}_X)\longrightarrow
  \psi^*(\mathscr{O}_X)/\psi^*(\mathscr{I})
  =\psi^*(\mathscr{O}_X/\mathscr{I})
  \xrightarrow{\omega}\mathscr{O}_Z.
\]
여기서 첫째 화살표는 표준 준동형이다. $\psi'$을 $Y$의 $X$ 안으로의
포함과 합성하면 $\psi$가 되는 연속사상 $Z\to Y$라 하자. 그러면
${\psi'}^*(\mathscr{O}_Y)=\psi^*(\mathscr{O}_X/\mathscr{I})$임은
자명하다. 한편 $\omega$는 국소 준동형임이 자명하므로
$g=(\psi',\omega^b)$는 준스킴의 사상 $Z\to Y$이다
\oldpage[I]{122}
(\hyperref[I.2.2.1-ko]{2.2.1}). 앞에서 본 바에 따라 $f=j\circ g$이므로
명제가 증명되었다.

\begin{corollary}[4.1.10]
\label{I.4.1.10-ko}
포함 사상 $Z\to X$가 \emph{포함 사상 $Y\to X$를 통하여
인수분해되기} 위한 필요충분조건은 $Z$가 $Y$의 부분준스킴인 것이다.
\end{corollary}

이때 $Z\leq Y$라고 쓰며, 이 관계는 $X$의 부분준스킴들의 집합 위의
\emph{순서관계}임이 자명하다.

\subsection{몰입 사상}
\label{subsection:I.4.2-ko}

\begin{definition}[4.2.1]
\label{I.4.2.1-ko}
사상 $f:Y\to X$가
$Y\xrightarrow{g}Z\xrightarrow{j}X$로 인수분해되고, 여기서 $g$가
동형사상, $Z$가 $X$의 부분준스킴(각각 닫힌 부분준스킴, 열린집합 위에
유도된 부분준스킴), $j$가 포함 사상이면 $f$를 \emph{몰입 사상}(각각
\emph{닫힌 몰입 사상}, \emph{열린 몰입 사상})이라 한다.
\end{definition}

이때 부분준스킴 $Z$와 동형사상 $g$는 \emph{유일하게} 결정된다. 실제로
$Z'$이 $X$의 또 다른 부분준스킴이고, $j'$이 포함 사상 $Z'\to X$이며,
$g'$이 $j\circ g=j'\circ g'$을 만족하는 동형사상 $Y\to Z'$이라고 하자.
그러면 $j'=j\circ g\circ {g'}^{-1}$이므로
$Z'\leq Z$이고(\hyperref[I.4.1.10-ko]{4.1.10}), 같은 방법으로
$Z\leq Z'$를 보일 수 있다. 따라서 $Z'=Z$이고, $j$가 준스킴의
단사사상이므로 $g'=g$이다.

$f=j\circ g$를 몰입 사상 $f$의 \emph{표준 인수분해}라 하고,
부분준스킴 $Z$와 동형사상 $g$를 \emph{$f$에 대응하는 것들}이라 한다.

몰입 사상은 준스킴의 \emph{단사사상}이고
(\hyperref[I.4.1.7-ko]{4.1.7}), \emph{따라서 특히} \emph{라디시엘 사상}임은
자명하다(\hyperref[I.3.5.4-ko]{3.5.4}).

\begin{proposition}[4.2.2]
\label{I.4.2.2-ko}
\begin{enumerate}
  \item[a)] 사상 $f=(\psi,\theta):Y\to X$가 \emph{열린 몰입 사상}이기
  위한 필요충분조건은, $\psi$가 $Y$에서 $X$의 열린 부분집합으로 가는
  위상동형이고 모든 $y\in Y$에 대하여 준동형
  $\theta_y^\sharp:\mathscr{O}_{\psi(y)}\to\mathscr{O}_y$가
  전단사인 것이다.
  \item[b)] 사상 $f=(\psi,\theta):Y\to X$가 \emph{몰입 사상}(각각
  \emph{닫힌 몰입 사상})이기 위한 필요충분조건은, $\psi$가 $Y$에서
  $X$의 국소 닫힌(각각 닫힌) 부분집합으로 가는 위상동형이고 모든
  $y\in Y$에 대하여 준동형
  $\theta_y^\sharp:\mathscr{O}_{\psi(y)}\to\mathscr{O}_y$가 전사인
  것이다.
\end{enumerate}
\end{proposition}

\begin{enumerate}
  \item[a)] 조건들이 필요함은 자명하다. 역으로 이 조건들이 성립하면
  $\theta^\sharp:\psi^*(\mathscr{O}_X)\to\mathscr{O}_Y$가 동형임이
  분명하다. 또한 $\psi^*(\mathscr{O}_X)$는
  $\mathscr{O}_X|\psi(Y)$에서 $\psi^{-1}$을 써서 구조를 옮겨 얻은
  층이므로 결론이 따른다.
  \item[b)] 조건들이 필요함은 자명하므로, 이제 충분함을 증명하자.
  먼저 $X$가 아핀 스킴이고 $Z=\psi(Y)$가 $X$ 안에서 \emph{닫혀}
  있다고 가정하는 특수한 경우를 생각하자.
  (\hyperref[0.3.4.6-ko]{0, 3.4.6})에 따라
  $\psi_*(\mathscr{O}_Y)$의 지지집합은 $Z$이고, 이 층의 $Z$로의
  제한을 $\mathscr{O}'_Z$라 하면 환 달린 공간
  $(Z,\mathscr{O}'_Z)$는 $Y$에서 $Z$로 가는 위상동형으로 본 $\psi$를
  써서 $(Y,\mathscr{O}_Y)$의 구조를 옮겨 얻는다.

  이제 $f_*(\mathscr{O}_Y)=\psi_*(\mathscr{O}_Y)$가 \emph{준연접}
  $\mathscr{O}_X$-가군임을 보이자. 실제로 $x\notin Z$이면
  $\psi_*(\mathscr{O}_Y)$는 $x$의 적당한 근방에 제한했을 때 영이다.
  반대로 $x\in Z$이면 유일하게 정해지는 $y\in Y$에 대하여
  $x=\psi(y)$이다. $V$를 $Y$에서 $y$의 아핀 열린 근방이라 하자.
  그러면 $\psi(V)$는 $Z$에서 열려 있으므로, $X$의 어떤 열린집합
  $U$와 $Z$의 교집합이다. $\psi_*(\mathscr{O}_Y)$의 $U$로의 제한은
  직\oldpage[I]{123}상 $(\psi_V)_*(\mathscr{O}_Y|V)$의 $U$로의 제한과 같다. 여기서
  $\psi_V$는 $\psi$의 $V$로의 제한이다.
  한편 사상 $(\psi,\theta)$의 $(V,\mathscr{O}_Y|V)$로의 제한은 이
  준스킴에서 $(X,\mathscr{O}_X)$로 가는 사상이므로
  $({}^a\varphi,\widetilde{\varphi})$ 꼴이다. 여기서 $\varphi$는 환
  $A=\Gamma(X,\mathscr{O}_X)$에서 환
  $\Gamma(V,\mathscr{O}_Y)$로 가는 준동형이다
  (\hyperref[I.1.7.3-ko]{1.7.3}). 따라서
  $(\psi_V)_*(\mathscr{O}_Y|V)$는 준연접 $\mathscr{O}_X$-가군이고
  (\hyperref[I.1.6.3-ko]{1.6.3}), 준연접층이라는 조건이 국소적이므로
  앞의 주장이 증명된다.

  더 나아가 $\psi$가 위상동형이라는 가정에서 모든 $y\in Y$에 대하여
  $\psi_y:(\psi_*(\mathscr{O}_Y))_{\psi(y)}\to\mathscr{O}_y$가
  동형임이 따른다(\hyperref[0.3.4.5-ko]{0, 3.4.5}). 도표
  \phantomsection\label{I.4.2.2.theta-diagram-ko}
  \[
    \xymatrix{
      \mathscr{O}_{\psi(y)}\ar[r]^{\theta_{\psi(y)}}
        \ar[d]_{\psi_y\circ\rho_{\psi(y)}} &
      (\psi_*(\mathscr{O}_Y))_{\psi(y)}\ar[d]^{\psi_y}\\
      (\psi^*(\mathscr{O}_X))_y\ar[r]^{\theta_y^\sharp} &
      \mathscr{O}_y
    }
  \]
  는 가환하고 두 세로 화살표는 동형이므로
  (\hyperref[0.3.7.2-ko]{0, 3.7.2}), $\theta_y^\sharp$이 전사라는
  가정에서 $\theta_{\psi(y)}$도 전사임이 따른다. 한편
  $\psi_*(\mathscr{O}_Y)$의 지지집합이 $Z=\psi(Y)$이므로,
  $\theta$는 $\mathscr{O}_X=\widetilde{A}$에서 준연접
  $\mathscr{O}_X$-가군 $f_*(\mathscr{O}_Y)$로 가는 \emph{전사}
  준동형이다. 따라서 어떤 $A$의 아이디얼 $\mathfrak{J}$와 유일한
  동형 $\omega:\widetilde{A}/\widetilde{\mathfrak{J}}
  \to f_*(\mathscr{O}_Y)$이 존재하여, $\omega$를 표준 준동형
  $\widetilde{A}\to\widetilde{A}/\widetilde{\mathfrak{J}}$와 합성하면
  $\theta$가 된다(\hyperref[I.1.3.8-ko]{1.3.8}).
  $\mathscr{O}_Z$를
  $\widetilde{A}/\widetilde{\mathfrak{J}}$의 $Z$로의 제한이라 하면
  $(Z,\mathscr{O}_Z)$는 $(X,\mathscr{O}_X)$의 부분준스킴이다.
  또한 $f$는 이 부분준스킴의 $X$ 안으로의 표준 포함 사상과 동형사상
  $(\psi_0,\omega_0)$을 거쳐 인수분해된다. 여기서 $\psi_0$는
  $Y$에서 $Z$로 가는 사상으로 본 $\psi$이고, $\omega_0$은
  $\omega$의 $\mathscr{O}_Z$로의 제한이다.

  이제 일반적인 경우를 생각하자. $U$를 $X$의 아핀 열린집합으로서
  $U\cap\psi(Y)$가 $U$ 안에서 닫혀 있고 공집합이 아닌 것으로 잡자.
  $f$를 열린집합 $\psi^{-1}(U)$ 위에서 $Y$가 유도하는 준스킴으로
  제한하고, 이를 $U$ 위에서 $X$가 유도하는 준스킴으로 가는 사상으로
  보면 첫째 경우로 귀착된다. 따라서 $f$의 $\psi^{-1}(U)$로의 제한은
  닫힌 몰입 사상 $\psi^{-1}(U)\to U$이고, 표준적으로
  $j_U\circ g_U$로 인수분해된다. 여기서 $g_U$는
  $\psi^{-1}(U)$에서 $U$의 부분준스킴 $Z_U$로 가는 동형사상이고,
  $j_U$는 표준 포함 사상 $Z_U\to U$이다.

  $V\subset U$인 $X$의 또 다른 아핀 열린집합 $V$를 잡자. $Z_U$의
  $V$로의 제한 $Z'_V$는 준스킴 $V$의 부분준스킴이므로, $f$의
  $\psi^{-1}(V)$로의 제한은 $j'_V\circ g'_V$로 인수분해된다. 여기서
  $j'_V$는 표준 포함 사상 $Z'_V\to V$이고, $g'_V$는
  $\psi^{-1}(V)$에서 $Z'_V$로 가는 동형사상이다. 몰입 사상의 표준
  인수분해가 유일하므로(\hyperref[I.4.2.1-ko]{4.2.1}) 반드시
  $Z'_V=Z_V$이고 $g'_V=g_V$이다. 따라서
  (\hyperref[I.4.1.5-ko]{4.1.5})에 따라 바탕공간이 $\psi(Y)$이고
  모든 $U\cap\psi(Y)$로의 제한이 $Z_U$인 $X$의 부분준스킴 $Z$가
  존재한다. 이때 $g_U$들은 각 $\psi^{-1}(U)$ 위에서 어떤 동형사상
  $g:Y\to Z$의 제한이고, $f=j\circ g$이다. 여기서 $j$는 표준 포함
  사상 $Z\to X$이다.
\end{enumerate}

\begin{corollary}[4.2.3]
\label{I.4.2.3-ko}
$X$를 아핀 스킴이라 하자. 사상 $f=(\psi,\theta):Y\to X$가 닫힌
몰입 사상이기 위한 필요충분조건은 $Y$가 아핀 스킴이고 준동형
$\Gamma(\theta):\Gamma(\mathscr{O}_X)\to\Gamma(\mathscr{O}_Y)$가
전사인 것이다.
\end{corollary}

\begin{corollary}[4.2.4]
\label{I.4.2.4-ko}
\begin{enumerate}
  \item[a)] $f:Y\to X$를 사상, $(V_\lambda)$를 $X$의 열린집합들로
  이루어진 $f(Y)$의 덮개라 하자. $f$가 몰입 사상(각각 열린 몰입
  사상)이기 위한 필요충분조건은
  \oldpage[I]{124}
  각 유도 준스킴
  $f^{-1}(V_\lambda)$로의 제한이 $V_\lambda$ 안으로의 몰입 사상
  (각각 열린 몰입 사상)인 것이다.
  \item[b)] $f:Y\to X$를 사상, $(V_\lambda)$를 $X$의 열린덮개라
  하자. $f$가 닫힌 몰입 사상이기 위한 필요충분조건은 각 유도
  준스킴 $f^{-1}(V_\lambda)$로의 제한이 $V_\lambda$ 안으로의 닫힌
  몰입 사상인 것이다.
\end{enumerate}
\end{corollary}

$f=(\psi,\theta)$라 하자. a)의 경우 모든 $y\in Y$에 대하여
$\theta_y^\sharp$은 전사(각각 전단사)이고, b)의 경우 모든 $y\in Y$에
대하여 $\theta_y^\sharp$은 전사이다. 따라서 a)의 경우 $\psi$가
$Y$에서 $X$의 국소 닫힌(각각 열린) 부분집합으로 가는 위상동형이고,
b)의 경우 $Y$에서 $X$의 닫힌 부분집합으로 가는 위상동형임만 확인하면
충분하다. 그런데 가정에 따라 $\psi$는 자명하게 단사이고, 모든
$y\in Y$에 대하여 $Y$에서 $y$의 모든 근방을 $\psi(Y)$에서
$\psi(y)$의 근방으로 보낸다. a)의 경우
$\psi(Y)\cap V_\lambda$는 $V_\lambda$에서 국소 닫혀(각각 열려)
있으므로 $\psi(Y)$는 $V_\lambda$들의 합집합에서 국소 닫혀(각각 열려)
있고, \emph{따라서 특히} $X$에서도 그러하다. b)의 경우
$\psi(Y)\cap V_\lambda$는 $V_\lambda$에서 닫혀 있으므로
$X=\bigcup_\lambda V_\lambda$에서 $\psi(Y)$는 닫혀 있다.

\begin{proposition}[4.2.5]
\label{I.4.2.5-ko}
두 몰입 사상(각각 두 열린 몰입 사상, 두 닫힌 몰입 사상)의 합성은
몰입 사상(각각 열린 몰입 사상, 닫힌 몰입 사상)이다.
\end{proposition}

이는 (\hyperref[I.4.1.6-ko]{4.1.6})에서 곧바로 따른다.

\subsection{몰입 사상의 곱}
\label{subsection:I.4.3-ko}

\begin{proposition}[4.3.1]
\label{I.4.3.1-ko}
$\alpha:X'\to X$, $\beta:Y'\to Y$를 두 $S$-사상이라 하자. $\alpha$와
$\beta$가 몰입 사상(각각 열린 몰입 사상, 닫힌 몰입 사상)이면
$\alpha\times_S\beta$도 몰입 사상(각각 열린 몰입 사상, 닫힌 몰입
사상)이다. 더욱이 $\alpha$(각각 $\beta$)가 $X'$(각각 $Y'$)를
$X$(각각 $Y$)의 부분준스킴 $X''$(각각 $Y''$)와 동일시한다면,
$\alpha\times_S\beta$는 $X'\times_S Y'$의 바탕공간을
$X\times_S Y$의 바탕공간의 부분공간
$p^{-1}(X'')\cap q^{-1}(Y'')$와 동일시한다. 여기서 $p$와 $q$는 각각
$X\times_S Y$에서 $X$와 $Y$로 가는 사영이다.
\end{proposition}

정의 (\hyperref[I.4.2.1-ko]{4.2.1})을 고려하면 $X'$와 $Y'$가
부분준스킴이고 $\alpha$와 $\beta$가 포함 사상인 경우로 한정할 수
있다. 열린집합 위에 유도된 부분준스킴에 대해서는 이 명제가 이미
(\hyperref[I.3.2.7-ko]{3.2.7})에서 증명되었다. 모든 부분준스킴은 어떤
열린집합 위에 유도된 준스킴의 닫힌 부분준스킴이므로
(\hyperref[I.4.1.3-ko]{4.1.3}), $X'$와 $Y'$가 \emph{닫힌}
부분준스킴인 경우로 귀착된다.

먼저 $S$가 \emph{아핀}이라고 가정할 수 있음을 보이자. 실제로
$(S_\lambda)$를 아핀 열린집합들로 이루어진 $S$의 덮개라 하고,
$\varphi$와 $\psi$를 $X$와 $Y$의 구조 사상이라 하자. 또
$X_\lambda=\varphi^{-1}(S_\lambda)$,
$Y_\lambda=\psi^{-1}(S_\lambda)$라 하자. $X'$의
$X_\lambda\cap X'$로의 제한 $X'_\lambda$(각각 $Y'$의
$Y_\lambda\cap Y'$로의 제한 $Y'_\lambda$)은 $X_\lambda$(각각
$Y_\lambda$)의 닫힌 부분준스킴이다. 준스킴 $X_\lambda$,
$Y_\lambda$, $X'_\lambda$, $Y'_\lambda$는 $S_\lambda$-준스킴으로
볼 수 있고, 곱 $X_\lambda\times_S Y_\lambda$와
$X_\lambda\times_{S_\lambda}Y_\lambda$(각각
$X'_\lambda\times_S Y'_\lambda$와
$X'_\lambda\times_{S_\lambda}Y'_\lambda$)는 동일하다
(\hyperref[I.3.2.5-ko]{3.2.5}). $S$가 아핀일 때 명제가 참이라면,
$\alpha\times_S\beta$를 각 $X'_\lambda\times_S Y'_\lambda$에
제한한 사상은 몰입 사상이다(\hyperref[I.3.2.7-ko]{3.2.7}). 또한 곱
$X'_\lambda\times_S Y'_\mu$(각각 $X_\lambda\times_S Y_\mu$)는
$(X'_\lambda\cap X'_\mu)\times_S
(Y'_\lambda\cap Y'_\mu)$(각각
$(X_\lambda\cap X_\mu)\times_S(Y_\lambda\cap Y_\mu)$)와
동일시되므로(\hyperref[I.3.2.6.4-ko]{3.2.6.4}),
$\alpha\times_S\beta$
\oldpage[I]{125}
를 각 $X'_\lambda\times_S Y'_\mu$에 제한한 사상도 몰입 사상이다.
따라서 (\hyperref[I.4.2.4-ko]{4.2.4})에 의해
$\alpha\times_S\beta$ 자체도 몰입 사상이다.

다음으로 $X$와 $Y$도 \emph{아핀}이라고 가정할 수 있음을 증명하자.
실제로 $(U_i)$(각각 $(V_j)$)를 아핀 열린집합들로 이루어진 $X$(각각
$Y$)의 덮개라 하고, $X'_i$(각각 $Y'_j$)를 $X'$의
$X'\cap U_i$로의 제한(각각 $Y'$의 $Y'\cap V_j$로의 제한)이라 하자.
이는 $U_i$(각각 $V_j$)의 닫힌 부분준스킴이다.
$U_i\times_S V_j$는 $X\times_S Y$를
$p^{-1}(U_i)\cap q^{-1}(V_j)$에 제한하여 얻은 준스킴과 동일시된다
(\hyperref[I.3.2.7-ko]{3.2.7}). 마찬가지로 $p'$, $q'$가
$X'\times_S Y'$의 사영이면 $X'_i\times_S Y'_j$는
$X'\times_S Y'$를
${p'}^{-1}(X'_i)\cap{q'}^{-1}(Y'_j)$에 제한하여 얻은 준스킴과
동일시된다. $\gamma=\alpha\times_S\beta$라 놓자. 정의에 따라
$p\circ\gamma=\alpha\circ p'$, $q\circ\gamma=\beta\circ q'$이고,
$X'_i=\alpha^{-1}(U_i)$, $Y'_j=\beta^{-1}(V_j)$이므로
${p'}^{-1}(X'_i)=\gamma^{-1}(p^{-1}(U_i))$,
${q'}^{-1}(Y'_j)=\gamma^{-1}(q^{-1}(V_j))$도 성립한다. 따라서
\[
  {p'}^{-1}(X'_i)\cap{q'}^{-1}(Y'_j)
  =\gamma^{-1}\bigl(p^{-1}(U_i)\cap q^{-1}(V_j)\bigr)
  =\gamma^{-1}(U_i\times_S V_j).
\]
이제 첫째 부분의 논증과 같이 결론을 얻는다.

그러므로 $X$, $Y$, $S$가 아핀이라고 가정하고, 그 환을 각각 $B$,
$C$, $A$라 하자. 이때 $B$와 $C$는 $A$-대수이고, $X'$와 $Y'$는
각각 $B$와 $C$의 몫대수 $B'$, $C'$를 환으로 갖는 아핀 스킴이다.
또한 $\alpha=({}^a\rho,\widetilde{\rho})$,
$\beta=({}^a\sigma,\widetilde{\sigma})$이고, 여기서 $\rho$와
$\sigma$는 각각 표준 준동형 $B\to B'$, $C\to C'$이다
(\hyperref[I.1.7.3-ko]{1.7.3}). 한편 $X\times_S Y$(각각
$X'\times_S Y'$)는 $B\otimes_A C$(각각 $B'\otimes_A C'$)를 환으로
갖는 아핀 스킴이고,
$\alpha\times_S\beta=({}^a\tau,\widetilde{\tau})$이다. 여기서 $\tau$는
$B\otimes_A C$에서 $B'\otimes_A C'$로 가는 준동형
$\rho\otimes\sigma$이다(\hyperref[I.3.2.2-ko]{3.2.2}와
\hyperref[I.3.2.3-ko]{3.2.3}). 이 준동형은 전사이므로
$\alpha\times_S\beta$는 실제로 몰입 사상이다. 더욱이
$\mathfrak{b}$(각각 $\mathfrak{c}$)가 $\rho$(각각 $\sigma$)의
핵이면 $\tau$의 핵은
$\operatorname{Im}(\mathfrak{b}\otimes_A C\to B\otimes_A C)
+\operatorname{Im}(B\otimes_A\mathfrak{c}\to B\otimes_A C)$이며, 이는
$u(\mathfrak{b})$와 $v(\mathfrak{c})$가 생성하는 아이디얼들의 합과
같다. 여기서 $u$(각각 $v$)는 준동형 $b\mapsto b\otimes1$(각각
$c\mapsto1\otimes c$)이다. $p=({}^a u,\widetilde{u})$이고
$q=({}^a v,\widetilde{v})$이므로, 이 핵은 $B\otimes_A C$의 소
스펙트럼에서 닫힌집합 $p^{-1}(X')\cap q^{-1}(Y')$에 대응한다
(\hyperref[I.1.2.2.1-ko]{1.2.2.1}과
\hyperref[I.1.1.2-ko]{1.1.2 (iii)}). 이로써 증명이 끝난다.

\begin{corollary}[4.3.2]
\label{I.4.3.2-ko}
$f:X\to Y$가 몰입 사상(각각 열린 몰입 사상, 닫힌 몰입 사상)이면서
$S$-사상이면, 밑준스킴의 임의의 확장 $S'\to S$에 대하여
$f_{(S')}$는 몰입 사상(각각 열린 몰입 사상, 닫힌 몰입 사상)이다.
\end{corollary}

\subsection{부분준스킴의 역상}
\label{subsection:I.4.4-ko}

\begin{proposition}[4.4.1]
\label{I.4.4.1-ko}
$f:X\to Y$를 사상, $Y'$를 $Y$의 부분준스킴(각각 닫힌 부분준스킴,
열린집합 위에 유도된 준스킴), $j:Y'\to Y$를 포함 사상이라 하자.
그러면 사영 $p:X\times_Y Y'\to X$는 몰입 사상(각각 닫힌 몰입 사상,
열린 몰입 사상)이다. $p$에 대응하는 $X$의 부분준스킴은
$f^{-1}(Y')$를 바탕공간으로 갖는다. 더욱이 $j'$를 이 부분준스킴에서
$X$로 가는 포함 사상이라 하면, 사상 $h:Z\to X$에 대하여
$f\circ h:Z\to Y$가 $j$를 통하여 인수분해되기 위한 필요충분조건은
$h$가 $j'$를 통하여 인수분해되는 것이다.
\end{proposition}

$p=1_X\times_Y j$이므로(\hyperref[I.3.3.4-ko]{3.3.4}), 첫째 주장은
(\hyperref[I.4.3.1-ko]{4.3.1})에서 따른다. 둘째 주장은
(\hyperref[I.3.5.10-ko]{3.5.10})의 특수한 경우이다(여기서는 $X$와
$Y'$의 역할을 서로 바꾸어야 한다). 끝으로 $h':Z\to Y'$인 사상
$h'$에 대하여 $f\circ h=j\circ h'$이라면, 곱의 정의에서 어떤 사상
$u:Z\to X\times_Y Y'$에 대하여 $h=p\circ u$임이 따른다. 따라서
마지막 주장도 얻는다.

이렇게 정의한 $X$의 부분준스킴을 사상 $f$에 의한 $Y'$의
\emph{역상}이라 부른다. 이 용어는 앞서
\oldpage[I]{126}
(\hyperref[I.3.3.6-ko]{3.3.6})에서 더 일반적으로 도입한 용어와
일치한다. $f^{-1}(Y')$를 $X$의 부분준스킴으로 말할 때에는 언제나 이
부분준스킴을 뜻한다.

준스킴 $f^{-1}(Y')$와 $X$가 동일하면 $j'$는 항등사상이고, 따라서
모든 사상 $h:Z\to X$가 $j'$를 통하여 인수분해된다. 그러므로
\emph{사상 $f:X\to Y$는 이때 다음과 같이 인수분해된다}:
$X\xrightarrow{g}Y'\xrightarrow{j}Y$.

$y$가 $Y$의 \emph{닫힌} 점이고
$Y'=\operatorname{Spec}(k(y))$가 $\{y\}$를 바탕공간으로 갖는 $Y$의
가장 작은 닫힌 부분준스킴이면(\hyperref[I.4.1.9-ko]{4.1.9}), 닫힌
부분준스킴 $f^{-1}(Y')$는
(\hyperref[I.3.6.2-ko]{3.6.2})에서 정의한 \emph{섬유} $f^{-1}(y)$와
표준적으로 동형이며, 앞으로 둘을 동일시한다.

\begin{corollary}[4.4.2]
\label{I.4.4.2-ko}
$f:X\to Y$, $g:Y\to Z$를 두 사상, $h=g\circ f$를 그 합성이라 하자.
$Z$의 모든 부분준스킴 $Z'$에 대하여 $X$의 부분준스킴
$f^{-1}(g^{-1}(Z'))$와 $h^{-1}(Z')$는 동일하다.
\end{corollary}

이는 표준 동형사상
$X\times_Y(Y\times_Z Z')\xrightarrow{\sim}X\times_Z Z'$의 존재에서
따른다(\hyperref[I.3.3.9.1-ko]{3.3.9.1}).

\begin{corollary}[4.4.3]
\label{I.4.4.3-ko}
$X'$, $X''$를 $X$의 두 부분준스킴,
$j':X'\to X$, $j'':X''\to X$를 포함 사상이라 하자. 그러면
${j'}^{-1}(X'')$와 ${j''}^{-1}(X')$는 둘 다 부분준스킴 사이의
순서관계에 관한 $X'$와 $X''$의 하한 $\inf(X',X'')$과 같고, 이
하한은 $X'\times_X X''$와 표준적으로 동형이다.
\end{corollary}

이는 (\hyperref[I.4.4.1-ko]{4.4.1})과
(\hyperref[I.4.1.10-ko]{4.1.10})에서 곧바로 따른다.

\begin{corollary}[4.4.4]
\label{I.4.4.4-ko}
$f:X\to Y$를 사상, $Y'$, $Y''$를 $Y$의 두 부분준스킴이라 하자.
그러면
\[
  f^{-1}(\inf(Y',Y''))=\inf(f^{-1}(Y'),f^{-1}(Y'')).
\]
\end{corollary}

이는 $(X\times_Y Y')\times_X(X\times_Y Y'')$와
$X\times_Y(Y'\times_Y Y'')$ 사이의 표준 동형사상이 존재함에서
따른다(\hyperref[I.3.3.9.1-ko]{3.3.9.1}).

\begin{proposition}[4.4.5]
\label{I.4.4.5-ko}
$f:X\to Y$를 사상, $Y'$를 $\mathscr{O}_Y$의 준연접 아이디얼층
$\mathscr{K}$로 정의된 $Y$의 닫힌 부분준스킴이라 하자
(\hyperref[I.4.1.3-ko]{4.1.3}). 그러면 $X$의 닫힌 부분준스킴
$f^{-1}(Y')$는 $\mathscr{O}_X$의 준연접 아이디얼층
$f^*(\mathscr{K})\mathscr{O}_X$로 정의된다.
\end{proposition}

{\raggedright
문제는 $X$와 $Y$ 위에서 자명하게 국소적이다. 따라서 $A$가
$B$-대수이고 $\mathfrak{K}$가 $B$의 아이디얼이면
$A\otimes_B(B/\mathfrak{K})=A/\mathfrak{K}A$임을 관찰하고
(\hyperref[I.1.6.9-ko]{1.6.9})를 적용하면 충분하다.\par}

\begin{corollary}[4.4.6]
\label{I.4.4.6-ko}
$X'$를 $\mathscr{O}_X$의 준연접 아이디얼층 $\mathscr{J}$로 정의된
$X$의 닫힌 부분준스킴, $i:X'\to X$를 포함 사상이라 하자. $f$를
$X'$에 제한한 사상 $f\circ i$가 포함 사상 $j:Y'\to Y$를 통하여
인수분해되기 위한(다시 말해 어떤 사상 $g:X'\to Y'$에 대하여
$j\circ g$로 인수분해되기 위한) 필요충분조건은
$f^*(\mathscr{K})\mathscr{O}_X\subset\mathscr{J}$인 것이다.
\end{corollary}

이는 명제 (\hyperref[I.4.4.1-ko]{4.4.1})을 $i$에 적용하되
(\hyperref[I.4.4.5-ko]{4.4.5})를 고려하면 된다.

\subsection{국소 몰입 사상과 국소 동형사상}
\label{subsection:I.4.5-ko}

\begin{definition}[4.5.1]
\label{I.4.5.1-ko}
준스킴의 사상 $f:X\to Y$가 \emph{점 $x\in X$에서 국소 몰입
사상}이라는 것은 $X$에서 $x$의 열린 근방 $U$와 $Y$에서 $f(x)$의
열린 근방 $V$가 존재하여, $U$ 위에 유도된 준스킴으로 $f$를 제한한
사상이 $U$에서 $V$ 위에 유도된 준스킴으로 가는 닫힌 몰입 사상인
것을 뜻한다. $f$가 $X$의 모든 점에서 국소 몰입 사상이면 $f$를 국소
몰입 사상이라 한다.
\end{definition}

\begin{definition}[4.5.2]
\label{I.4.5.2-ko}
사상 $f:X\to Y$가
\oldpage[I]{127}
점 $x\in X$에서 국소 동형사상이라는 것은 $X$에서 $x$의 열린 근방
$U$가 존재하여, $U$ 위에 유도된 준스킴으로 $f$를 제한한 사상이
$U$에서 $Y$로 가는 열린 몰입 사상인 것을 뜻한다. $f$가 $X$의 모든
점에서 국소 동형사상이면 $f$를 국소 동형사상이라 한다.
\end{definition}

\begin{env}[4.5.3]
\label{I.4.5.3-ko}
따라서 몰입 사상(각각 닫힌 몰입 사상) $f:X\to Y$는 국소 몰입
사상이면서 $X$의 바탕공간에서 $Y$의 한 부분집합(각각 닫힌
부분집합) 위로 가는 위상동형사상인 사상으로 특징지을 수 있다. 열린
몰입 사상 $f$는 바탕사상이 \emph{단사인} 국소 동형사상으로 특징지을
수 있다.
\end{env}

\begin{proposition}[4.5.4]
\label{I.4.5.4-ko}
$X$를 기약 준스킴, $f:X\to Y$를 바탕사상이 단사인 우세 사상이라
하자. $f$가 국소 몰입 사상이면 $f$는 몰입 사상이고 $f(X)$는 $Y$에서
열려 있다.
\end{proposition}

실제로 $x\in X$라 하고, $U$를 $x$의 열린 근방, $V$를 $Y$에서
$f(x)$의 열린 근방이라 하여 $f$의 $U$로의 제한이 $V$ 안으로의 닫힌
몰입 사상이 되게 하자. $U$는 $X$에서 조밀하고 가정에 따라 $f(U)$는
$Y$에서 조밀하므로 $f(U)=V$이며, $f$는 $U$에서 $V$ 위로 가는
위상동형사상이다. $f$의 바탕사상이 단사라는 가정에서
$f^{-1}(V)=U$가 따르므로 명제를 곧바로 얻는다.

\begin{proposition}[4.5.5]
\label{I.4.5.5-ko}
\begin{enumerate}
  \item[\rm(i)] 두 국소 몰입 사상(각각 두 국소 동형사상)의 합성은
    국소 몰입 사상(각각 국소 동형사상)이다.
  \item[\rm(ii)] $f:X\to X'$, $g:Y\to Y'$를 두 $S$-사상이라 하자.
    $f$와 $g$가 국소 몰입 사상(각각 국소 동형사상)이면
    $f\times_S g$도 그러하다.
  \item[\rm(iii)] $S$-사상 $f$가 국소 몰입 사상(각각 국소
    동형사상)이면, 밑준스킴의 모든 확장 $S'\to S$에 대하여
    $f_{(S')}$도 그러하다.
\end{enumerate}
\end{proposition}

(\hyperref[I.3.5.1-ko]{3.5.1})에 따라 (i)와 (ii)만 증명하면 충분하다.

(i)는 닫힌 몰입 사상(각각 열린 몰입 사상)의 추이성
(\hyperref[I.4.2.5-ko]{4.2.5})과 다음 사실에서 곧바로 따른다. $f$가
$X$에서 $Y$의 닫힌 부분집합 위로 가는 위상동형사상이면, 모든
열린집합 $U\subset X$에 대하여 $f(U)$는 $f(X)$에서 열려 있으므로
$f(U)=V\cap f(X)$가 되는 $Y$의 열린 부분집합 $V$가 존재하고, 따라서
$f(U)$는 $V$에서 닫혀 있다.

(ii)를 증명하자. $z\in X\times_S Y$를 임의의 점으로 택하고
$z'=(f\times_S g)(z)\in X'\times_S Y'$라 하자. $p$, $q$를
$X\times_S Y$의 사영, $p'$, $q'$를 $X'\times_S Y'$의 사영이라 하자.
가정에 따라 $x=p(z)$, $x'=p'(z')$, $y=q(z)$, $y'=q'(z')$의 열린
근방 $U$, $U'$, $V$, $V'$가 각각 존재하여, $f$와 $g$를 각각 $U$와
$V$로 제한한 사상이 각각 $U'$와 $V'$ 안으로의 닫힌 몰입 사상(각각
열린 몰입 사상)이 된다. $U\times_S V$와 $U'\times_S V'$의
바탕공간은 각각 $z$와 $z'$의 열린 근방
$p^{-1}(U)\cap q^{-1}(V)$와
${p'}^{-1}(U')\cap{q'}^{-1}(V')$로 동일시되므로
(\hyperref[I.3.2.7-ko]{3.2.7}), 명제는
(\hyperref[I.4.3.1-ko]{4.3.1})에서 따른다.

\section{축소 준스킴; 분리 조건}
\label{section:I.5-ko}

\subsection{축소 준스킴}
\label{subsection:I.5.1-ko}

\begin{proposition}[5.1.1]
\label{I.5.1.1-ko}
$(X,\mathscr{O}_X)$를 준스킴, $\mathscr{B}$를 준연접
$\mathscr{O}_X$-대수라 하자. 모든 $x\in X$에서 줄기 $\mathscr{N}_x$가
환 $\mathscr{B}_x$의 멱영근기가 되는 유일한 준연접
$\mathscr{O}_X$-가군 $\mathscr{N}$이 존재한다. $X$가 아핀이고 따라서
$\mathscr{B}=\widetilde{B}$이며 $B$가 $A(X)$ 위의 대수이면,
$\mathscr{N}=\widetilde{\mathfrak{N}}$이다. 여기서 $\mathfrak{N}$은
$B$의 멱영근기이다.
\end{proposition}

\oldpage[I]{128}
문제는 국소적이므로 마지막 주장만 증명하면 된다.
$\widetilde{\mathfrak{N}}$은 준연접 $\mathscr{O}_X$-가군이고
(\hyperref[I.1.4.1-ko]{1.4.1}), 점 $x\in X$에서 그 줄기는 분수환
$B_x$의 아이디얼 $\mathfrak{N}_x$이다. 따라서 반대쪽 포함관계는
자명하므로, $B_x$의 멱영근기가 $\mathfrak{N}_x$에 포함됨을 보이면
충분하다. $z\in B$, $s\notin\mathfrak{j}_x$이고 $z/s$가 $B_x$의
멱영원이라고 하자. 가정에 따라 $(z/s)^k=0$인 정수 $k$가 존재하고,
이는 $tz^k=0$인 $t\notin\mathfrak{j}_x$가 존재한다는 뜻이다. 그러면
$(tz)^k=0$이고, 따라서 $z/s=(tz)/(ts)$는 $\mathfrak{N}_x$에 속한다.

이렇게 정의한 준연접 $\mathscr{O}_X$-가군 $\mathscr{N}$을
$\mathscr{O}_X$-대수 $\mathscr{B}$의 \emph{멱영근기}라 부른다. 특히
$\mathscr{O}_X$의 멱영근기를 $\mathscr{N}_X$로 쓴다.

\begin{corollary}[5.1.2]
\label{I.5.1.2-ko}
$X$를 준스킴이라 하자. 아이디얼층 $\mathscr{N}_X$로 정의된 $X$의
닫힌 부분준스킴은 $X$를 바탕공간으로 갖는 유일한 축소 부분준스킴
(\hyperref[0.4.1.4-ko]{0, 4.1.4})이다. 또한 이는 $X$를 바탕공간으로
갖는 $X$의 가장 작은 부분준스킴이다.
\end{corollary}

$\mathscr{N}_X$로 정의된 닫힌 부분준스킴 $Y$의 구조층은
$\mathscr{O}_X/\mathscr{N}_X$이다. 모든 $x\in X$에 대하여
$\mathscr{N}_x\ne\mathscr{O}_x$이므로 $Y$가 축소이고 $X$를
바탕공간으로 갖는다는 것은 자명하다. 나머지 주장을 증명하기 위하여,
$X$를 바탕공간으로 갖는 $X$의 부분준스킴 $Z$가 모든 $x\in X$에
대하여 $\mathscr{J}_x\ne\mathscr{O}_x$인 아이디얼층 $\mathscr{J}$로
정의됨을 주목하자(\hyperref[I.4.1.3-ko]{4.1.3}). $X$가 아핀인 경우로
환원할 수 있다. $X=\operatorname{Spec}(A)$,
$\mathscr{J}=\widetilde{\mathfrak{J}}$라 하고 $\mathfrak{J}$를 $A$의
아이디얼이라 하자. 그러면 모든 $x\in X$에 대하여
$\mathfrak{J}\subset\mathfrak{j}_x$이다. 따라서 $\mathfrak{J}$는
$A$의 모든 소 아이디얼, 곧 그 교집합인 $A$의 멱영근기
$\mathfrak{N}$에 포함된다. 이로부터 $Y$가 $X$를 바탕공간으로 갖는
$X$의 가장 작은 부분준스킴임이 따른다
(\hyperref[I.4.1.9-ko]{4.1.9}). 더욱이 $Z$가 $Y$와 다르면 적어도 한
$x\in X$에 대하여 $\mathscr{J}_x\ne\mathscr{N}_x$이고, 따라서
(\hyperref[I.5.1.1-ko]{5.1.1})에 의하여 $Z$는 축소가 아니다.

\begin{definition}[5.1.3]
\label{I.5.1.3-ko}
준스킴 $X$에 대응하며 $X$를 바탕공간으로 갖는 유일한 축소
부분준스킴을 $X$에 대응하는 축소 준스킴이라 하고
$X_{\mathrm{red}}$로 쓴다.
\end{definition}

그러므로 준스킴 $X$가 축소라는 것은 $X=X_{\mathrm{red}}$임을 뜻한다.

\begin{proposition}[5.1.4]
\label{I.5.1.4-ko}
환 $A$의 소 스펙트럼이 축소 준스킴(각각 정역 준스킴)
(\hyperref[I.2.1.8-ko]{2.1.8})이기 위한 필요충분조건은 $A$가
축소환(각각 정역)인 것이다.
\end{proposition}

실제로 (\hyperref[I.5.1.1-ko]{5.1.1})에서
$X=\operatorname{Spec}(A)$가 축소이기 위한 필요충분조건이
$\mathscr{N}=(0)$임을 곧바로 얻는다. 정역에 관한 주장은
(\hyperref[I.1.1.13-ko]{1.1.13})에서 따른다.

정역의 영환이 아닌 모든 분수환은 정역이므로,
(\hyperref[I.5.1.4-ko]{5.1.4})에서 모든 \emph{국소 정역 준스킴}
$X$와 모든 $x\in X$에 대하여 $\mathscr{O}_x$가 \emph{정역}임이
따른다. $X$의 바탕공간이 \emph{국소 뇌터}이면 그 역도 성립한다.
실제로 이때 $X$는 축소이다. $X$의 아핀 열린집합 $U$가 뇌터 공간이면
$U$의 기약 성분은 유한개뿐이고, 따라서 그 환 $A$의 극소 소
아이디얼도 유한개뿐이다(\hyperref[I.1.1.14-ko]{1.1.14}). 이 성분들
$U_i$ 가운데 둘이 공통점 $x$를 가지면 $\mathscr{O}_x$는 서로 다른
극소 소 아이디얼을 적어도 둘 가지므로 정역이 아니다. 따라서
$U_i$들은 서로소인 열린집합이고, 각각 정역 준스킴이다.

\begin{env}[5.1.5]
\label{I.5.1.5-ko}
$f=(\psi,\theta):X\to Y$를 준스킴의 사상이라 하자. 준동형
\oldpage[I]{129}
$\theta_x^\sharp:\mathscr{O}_{\psi(x)}\to\mathscr{O}_x$는
$\mathscr{O}_{\psi(x)}$의 모든 멱영원을 $\mathscr{O}_x$의 멱영원으로
보낸다. 그러므로 몫을 취하면 $\theta^\sharp$에서 준동형
\[
  \omega:\psi^*(\mathscr{O}_Y/\mathscr{N}_Y)
  \longrightarrow \mathscr{O}_X/\mathscr{N}_X
\]
를 얻는다. 모든 $x\in X$에 대하여
$\omega_x:\mathscr{O}_{\psi(x)}/\mathscr{N}_{\psi(x)}
\to\mathscr{O}_x/\mathscr{N}_x$가 국소 준동형임은 명백하다. 따라서
$(\psi,\omega^b)$는 준스킴의 사상
$X_{\mathrm{red}}\to Y_{\mathrm{red}}$이다. 이를 $f_{\mathrm{red}}$로
쓰고 $f$에 대응하는 \emph{축소} 사상이라 부른다. 두 사상
$f:X\to Y$, $g:Y\to Z$에 대하여
$(g\circ f)_{\mathrm{red}}=g_{\mathrm{red}}\circ f_{\mathrm{red}}$임은
자명하다. 따라서 $X_{\mathrm{red}}$는 $X$에 관한 \emph{공변 함자}로
정의되었다.

앞의 정의에서 도식
\[
  \xymatrix{
    X_{\mathrm{red}}\ar[r]^{f_{\mathrm{red}}}\ar[d] &
    Y_{\mathrm{red}}\ar[d]\\
    X\ar[r]_f & Y
  }
\]
이 가환임이 따른다. 여기서 세로 화살표는 포함 사상이다. 다시 말해
$X_{\mathrm{red}}\to X$는 \emph{함자적}인 사상이다. 특히 $X$가
\emph{축소}이면 모든 사상 $f:X\to Y$는
$X\xrightarrow{f_{\mathrm{red}}}Y_{\mathrm{red}}\to Y$로
인수분해된다. 다시 말해 $f$는 포함 사상
$Y_{\mathrm{red}}\to Y$를 통하여 \emph{인수분해된다}.
\end{env}

\begin{proposition}[5.1.6]
\label{I.5.1.6-ko}
$f:X\to Y$를 사상이라 하자. $f$가 전사(각각 라디시엘 사상, 몰입
사상, 닫힌 몰입 사상, 열린 몰입 사상, 국소 몰입 사상, 국소
동형사상)이면 $f_{\mathrm{red}}$도 그러하다. 역으로
$f_{\mathrm{red}}$가 전사(각각 라디시엘 사상)이면 $f$도 그러하다.
\end{proposition}

$f$가 전사인 경우 명제는 자명하다. $f$가 라디시엘 사상이면, 모든
$x\in X$에 대하여 체 $k(x)$가 준스킴 $X$와 $X_{\mathrm{red}}$에서
동일하다는 사실에서 명제가 따른다
(\hyperref[I.3.5.8-ko]{3.5.8}). 끝으로 $f=(\psi,\theta)$가 몰입 사상,
닫힌 몰입 사상 또는 국소 몰입 사상(각각 열린 몰입 사상 또는 국소
동형사상)이면, $\theta_x^\sharp$가 전사(각각 전단사)일 때
$\mathscr{O}_{\psi(x)}$와 $\mathscr{O}_x$의 멱영근기로 몫을 취하여
얻는 준동형도 전사(각각 전단사)라는 사실에서 명제가 따른다
(\hyperref[I.5.1.2-ko]{5.1.2}와
\hyperref[I.4.2.2-ko]{4.2.2}; 또한
(\agref{I.5.5.12-ko}{5.5.12}) 참조).

\begin{proposition}[5.1.7]
\label{I.5.1.7-ko}
$X$, $Y$가 두 $S$-준스킴이면, 준스킴
$X_{\mathrm{red}}\times_{S_{\mathrm{red}}}Y_{\mathrm{red}}$와
$X_{\mathrm{red}}\times_S Y_{\mathrm{red}}$는 동일하며,
$X\times_S Y$와 같은 바탕공간을 갖는 $X\times_S Y$의 부분준스킴과 표준적으로
동일시된다.
\end{proposition}

$X_{\mathrm{red}}\times_S Y_{\mathrm{red}}$가 $X\times_S Y$와 같은
바탕공간을 갖는 $X\times_S Y$의 부분준스킴과 표준적으로 동일시되는 것은
(\hyperref[I.4.3.1-ko]{4.3.1})에서 따른다. 한편 $\varphi$와 $\psi$가
구조 사상 $X_{\mathrm{red}}\to S$, $Y_{\mathrm{red}}\to S$이면, 이들은
$S_{\mathrm{red}}$를 통하여 인수분해된다
(\hyperref[I.5.1.5-ko]{5.1.5}). $S_{\mathrm{red}}\to S$는
단사사상이므로 첫째 주장은 (\hyperref[I.3.2.4-ko]{3.2.4})에서
따른다.

\begin{corollary}[5.1.8]
\label{I.5.1.8-ko}
준스킴 $(X\times_S Y)_{\mathrm{red}}$와
$(X_{\mathrm{red}}\times_{S_{\mathrm{red}}}Y_{\mathrm{red}})_{\mathrm{red}}$는
표준적으로 동일시된다.
\end{corollary}

이는 (\hyperref[I.5.1.2-ko]{5.1.2}와
\hyperref[I.5.1.7-ko]{5.1.7})에서 따른다.

$X$와 $Y$가 축소 준스킴이어도 $X\times_S Y$가 반드시 축소인 것은
아니다. 두 축소 대수의 텐서곱에 멱영원이 있을 수 있기 때문이다.

\begin{proposition}[5.1.9]
\label{I.5.1.9-ko}
\oldpage[I]{130}
$X$를 준스킴, $\mathscr{J}$를 어떤 정수 $n>0$에 대하여
$\mathscr{J}^n=0$인 $\mathscr{O}_X$의 준연접 아이디얼층이라 하자.
$X_0$를 $X$의 닫힌 부분준스킴
$(X,\mathscr{O}_X/\mathscr{J})$라 하자. $X$가 아핀 스킴이기 위한
필요충분조건은 $X_0$가 아핀 스킴인 것이다.
\end{proposition}

조건이 필요하다는 것은 자명하므로 충분함을 증명하자.
$X_k=(X,\mathscr{O}_X/\mathscr{J}^{k+1})$라 두면, $k$에 관한 귀납법으로
각 $X_k$가 아핀임을 증명하면 충분하다. 따라서 $\mathscr{J}^2=0$인
경우로 환원된다. 다음과 같이 놓자.
\[
  A=\Gamma(X,\mathscr{O}_X),\qquad
  A_0=\Gamma(X_0,\mathscr{O}_{X_0})
      =\Gamma(X,\mathscr{O}_X/\mathscr{J}).
\]
표준 준동형 $\mathscr{O}_X\to\mathscr{O}_X/\mathscr{J}$에서 환의
준동형 $\varphi:A\to A_0$를 얻는다. 아래에서 $\varphi$가
\emph{전사}임을 보일 것이다. 그러면 열
\[
  0\to\Gamma(X,\mathscr{J})\to\Gamma(X,\mathscr{O}_X)
  \to\Gamma(X,\mathscr{O}_X/\mathscr{J})\to 0
  \tag{5.1.9.1}\label{I.5.1.9.1-ko}
\]
은 \emph{완전}하다. 이 점을 안다고 하고, 그것이 명제를 함의함을
보이자. $\mathfrak{K}=\Gamma(X,\mathscr{J})$는 $A$의 제곱영
아이디얼이고, 따라서 $A_0=A/\mathfrak{K}$ 위의 가군이다. 가정에
따라 $X_0=\operatorname{Spec}(A_0)$이고, 바탕 위상공간 $X_0$와 $X$는
동일하므로 $\mathfrak{K}=\Gamma(X_0,\mathscr{J})$이다. 또한
$\mathscr{J}^2=0$이므로 $\mathscr{J}$는 준연접
$(\mathscr{O}_X/\mathscr{J})$-가군이다. 따라서
$\mathscr{J}\simeq\widetilde{\mathfrak{K}}$이고, 모든 $x\in X_0$에
대하여 $\mathfrak{K}_x=\mathscr{J}_x$이다
(\hyperref[I.1.4.1-ko]{1.4.1}).

이제 $X'=\operatorname{Spec}(A)$라 하고, 항등 사상
$A\to\Gamma(X,\mathscr{O}_X)$에 대응하는 준스킴의 사상
$f=(\psi,\theta):X\to X'$를 생각하자
(\hyperref[I.2.2.4-ko]{2.2.4}). $X$의 모든 아핀 열린집합 $V$에 대하여
도식
\[
  \xymatrix{
    A\ar[r]\ar[d] & \Gamma(V,\mathscr{O}_X|V)\ar[d]\\
    A_0=A/\mathfrak{K}\ar[r] & \Gamma(V,\mathscr{O}_{X_0}|V)
  }
\]
은 가환이다. 따라서 도식
\[
  \xymatrix{
    X' & X\ar[l]_f\\
    X'_0\ar[u]^{j'} & X_0\ar[l]^{f_0}\ar[u]_j
  }
\]
도 가환이다. 여기서 $X'_0$는 준연접 아이디얼층
$\widetilde{\mathfrak{K}}$로 정의된 $X'$의 닫힌 부분준스킴이고,
$j$, $j'$는 표준 포함 사상이다. $X_0$는 아핀이므로 $f_0$는
동형사상이다. $j$와 $j'$의 바탕 연속사상은 항등사상이므로 먼저
$\psi:X\to X'$가 위상동형사상임을 알 수 있다. 또한
$\mathfrak{K}_x=\mathscr{J}_x$라는 관계에서
$\theta^\sharp:\psi^*(\mathscr{O}_{X'})\to\mathscr{O}_X$의 제한이
$\psi^*(\widetilde{\mathfrak{K}})$에서 $\mathscr{J}$ 위로 가는
\emph{동형사상}임이 따른다. 한편 몫을 취하면 $f_0$가 동형사상이므로
$\theta^\sharp$에서 \emph{동형사상}
$\psi^*(\mathscr{O}_{X'}/\widetilde{\mathfrak{K}})
\to\mathscr{O}_X/\mathscr{J}$를 얻는다. 따라서 5-보조정리
$(\mathrm{M},\mathrm{I},1.1)$에서 $\theta^\sharp$ 자체가
동형사상임이 곧바로 따른다. 그러므로 $f$는 \emph{동형사상}이고
$X$는 아핀이다.

이제 (\hyperref[I.5.1.9.1-ko]{5.1.9.1})의 완전성을 증명하면
충분하다. 이는

\oldpage[I]{131}
$H^1(X,\mathscr{J})=0$에서 따른다. 그런데
$H^1(X,\mathscr{J})=H^1(X_0,\mathscr{J})$이고, 앞에서
$\mathscr{J}$가 준연접 $\mathscr{O}_{X_0}$-가군임을 보았다. 따라서
우리의 주장은 다음 보조정리에서 따른다.

\begin{lemma}[5.1.9.2]
\label{I.5.1.9.2-ko}
$Y$가 아핀 스킴이고 $\mathscr{F}$가 준연접
$\mathscr{O}_Y$-가군이면 $H^1(Y,\mathscr{F})=0$이다.
\end{lemma}

이 보조정리는 제III장 §1에서, 모든 $i>0$에 대하여
$H^i(Y,\mathscr{F})=0$이라는 더 일반적인 정리의 결과로 증명할
것이다. 독립적인 증명을 주기 위하여, $H^1(Y,\mathscr{F})$가
$\mathscr{F}$에 의한 $\mathscr{O}_Y$-가군 $\mathscr{O}_Y$의 확장류로
이루어진 가군
$\operatorname{Ext}^1_{\mathscr{O}_Y}
(Y;\mathscr{O}_Y,\mathscr{F})$와 동일시됨을 주목하자
$(\mathrm{T},4.2.3)$. 따라서 그러한 확장 $\mathscr{G}$가 자명함을
보이면 충분하다. 모든 $y\in Y$에 대하여 $Y$에서 $y$의 근방 $V$가
존재하여 $\mathscr{G}|V$가
$\mathscr{F}|V\oplus\mathscr{O}_Y|V$와 동형이다
(\hyperref[0.5.4.9-ko]{0, 5.4.9}). 따라서 $\mathscr{G}$는 준연접
$\mathscr{O}_Y$-가군이다. $A$를 $Y$의 환이라 하면
$\mathscr{F}=\widetilde{M}$, $\mathscr{G}=\widetilde{N}$이고, 여기서
$M$, $N$은 $A$-가군이다. 가정에 따라 $N$은 $A$-가군 $M$에 의한
$A$-가군 $A$의 확장이다(\hyperref[I.1.3.11-ko]{1.3.11}). 이 확장은
반드시 자명하므로 보조정리가 증명되고, 따라서
(\hyperref[I.5.1.9-ko]{5.1.9})도 증명된다.

\begin{corollary}[5.1.10]
\label{I.5.1.10-ko}
$X$를 $\mathscr{N}_X$가 멱영인 준스킴이라 하자. $X$가 아핀 스킴이기
위한 필요충분조건은 $X_{\mathrm{red}}$가 아핀 스킴인 것이다.
\end{corollary}

\subsection{주어진 바탕공간을 갖는 부분준스킴의 존재}
\label{subsection:I.5.2-ko}

\begin{proposition}[5.2.1]
\label{I.5.2.1-ko}
준스킴 $X$의 바탕공간의 모든 국소 닫힌 부분공간 $Y$에 대하여,
$Y$를 바탕공간으로 갖는 $X$의 축소 부분준스킴이 유일하게 존재한다.
\end{proposition}

유일성은 (\hyperref[I.5.1.2-ko]{5.1.2})에서 따르므로, 문제의
부분준스킴의 존재만 증명하면 된다.

$X$가 환 $A$의 아핀 준스킴이고 $Y$가 $X$에서 닫혀 있으면 명제는
곧바로 따른다. 실제로 $\mathfrak{j}(Y)$는
$V(\mathfrak{a})=Y$를 만족하는 아이디얼 $\mathfrak{a}\subset A$ 가운데
가장 큰 것이며 자기 근기와 같다
(\hyperref[I.1.1.4-ko]{1.1.4 (i)}). 따라서
$A/\mathfrak{j}(Y)$는 축소환이다.

일반적인 경우, $U\cap Y$가 $U$에서 닫히도록 하는 모든 아핀
열린집합 $U\subset X$에 대하여, $A(U)$의 아이디얼
$\mathfrak{j}(U\cap Y)$에 대응하는 아이디얼층으로 정의된 $U$의 닫힌
부분준스킴 $Y_U$를 생각하자. 이 부분준스킴은 축소이다. 이제 $V$가
$U$에 포함되는 $X$의 아핀 열린집합이면, $Y_V$는 $V\cap Y$ 위에서
$Y_U$에 의해 \emph{유도된} 부분준스킴임을 보이자. 그런데 이렇게
유도된 준스킴은 $V$의 축소 닫힌 부분준스킴이고 $V\cap Y$를
바탕공간으로 가지므로, $Y_V$의 유일성에 의해 주장이 따른다.

\begin{proposition}[5.2.2]
\label{I.5.2.2-ko}
$X$를 축소 준스킴이라 하고, $f:X\to Y$를 사상, $Z$를
$f(X)\subset Z$를 만족하는 $Y$의 닫힌 부분준스킴이라 하자. 그러면 $f$는
$X\xrightarrow{g}Z\xrightarrow{j}Y$로 인수분해되며, 여기서 $j$는 포함
사상이다.
\end{proposition}

가정에 따라 $X$의 닫힌 부분준스킴 $f^{-1}(Z)$의 바탕공간은 $X$
전체이다 (\hyperref[I.4.4.1-ko]{4.4.1}). $X$가 축소이므로 이 닫힌
부분준스킴은 $X$와 동일하다 (\hyperref[I.5.1.2-ko]{5.1.2}). 따라서
명제는 (\hyperref[I.4.4.1-ko]{4.4.1})에서 따른다.

\begin{corollary}[5.2.3]
\label{I.5.2.3-ko}
$X$를 준스킴 $Y$의 축소 부분준스킴이라 하자. $Z$가 $\overline{X}$를
바탕공간으로 갖는 $Y$의 축소 닫힌 부분준스킴이면, $X$는 $Z$의 어떤
열린집합 위에 유도된 부분준스킴이다.
\end{corollary}

\oldpage[I]{132}
실제로 $Y$의 어떤 열린집합 $U$에 대하여, 바탕공간으로서
$X=U\cap\overline{X}$이다. (\hyperref[I.5.2.2-ko]{5.2.2})에 따라 $X$는
$Z$의 축소 부분준스킴이고, 유일성
(\hyperref[I.5.2.1-ko]{5.2.1})에 따라 부분준스킴 $X$는 열린
부분공간 $X$ 위에서 $Z$에 의해 유도된다.

\begin{corollary}[5.2.4]
\label{I.5.2.4-ko}
$f:X\to Y$를 사상이라 하고, $X'$ (각각 $Y'$)를 $\mathscr{O}_X$의
(각각 $\mathscr{O}_Y$의) 준연접 아이디얼층 $\mathscr{J}$
(각각 $\mathscr{K}$)로 정의된 $X$의 (각각 $Y$의) 닫힌
부분준스킴이라 하자. $X'$가 축소이고 $f(X')\subset Y'$라고 하자.
그러면 $f^*(\mathscr{K})\mathscr{O}_X\subset\mathscr{J}$이다.
\end{corollary}

(\hyperref[I.5.2.2-ko]{5.2.2})에 따라 $f$를 $X'$에 제한한 사상이
$X'\to Y'\to Y$로 인수분해되므로,
(\hyperref[I.4.4.6-ko]{4.4.6})을 적용하면 충분하다.

\subsection{대각선; 사상의 그래프}
\label{subsection:I.5.3-ko}

\begin{env}[5.3.1]
\label{I.5.3.1-ko}
$X$를 $S$-준스킴이라 하자. $X$에서 $X\times_S X$로 가는
$S$-사상 $(1_X,1_X)_S$, 다시 말해
\[
  p_1\circ\Delta_X=p_2\circ\Delta_X=1_X
  \tag{5.3.1.1}\label{I.5.3.1.1-ko}
\]
을 만족하는 유일한 $S$-사상 $\Delta_X$를 $X$의 \emph{대각 사상}이라
하고 $\Delta_{X|S}$ 또는 $\Delta_X$, 혼동의 우려가 없으면 간단히
$\Delta$로 쓴다. 여기서 $p_1$, $p_2$는 $X\times_S X$의 사영이다
(정의 \hyperref[I.3.2.1-ko]{3.2.1}). $f:T\to X$, $g:T\to Y$가 두
$S$-사상이면 곧바로
\[
  (f,g)_S=(f\times_S g)\circ\Delta_{T|S}
  \tag{5.3.1.2}\label{I.5.3.1.2-ko}
\]
임을 확인할 수 있다.

독자는 앞의 정의와 (\hyperref[I.5.3.1-ko]{5.3.1})부터 (5.3.8)까지
서술한 결과가 어느 범주에서도
\emph{그 안에 나타나는 곱들이 그 범주에 존재하기만 하면} 성립함을
주목할 것이다.
\end{env}

\begin{proposition}[5.3.2]
\label{I.5.3.2-ko}
$X$, $Y$를 두 $S$-준스킴이라 하자. 곱
$(X\times Y)\times(X\times Y)$를
$(X\times X)\times(Y\times Y)$와 표준적으로 동일시하면, 사상
$\Delta_{X\times Y}$는 $\Delta_X\times\Delta_Y$와 동일시된다.
\end{proposition}

실제로 $p_1$, $q_1$이 각각 $X\times X\to X$,
$Y\times Y\to Y$인 첫째 사영이면,
$(X\times Y)\times(X\times Y)\to X\times Y$인 첫째 사영은
$p_1\times q_1$과 동일시되고
\[
  (p_1\times q_1)\circ(\Delta_X\times\Delta_Y)
  =(p_1\circ\Delta_X)\times(q_1\circ\Delta_Y)=1_{X\times Y}
\]
이다. 둘째 사영들에 대해서도 마찬가지이다.

\begin{corollary}[5.3.4]
\label{I.5.3.4-ko}
밑준스킴의 모든 확장 $S'\to S$에 대하여,
$\Delta_{X_{(S')}}$는 $(\Delta_X)_{(S')}$와 표준적으로 동일시된다.
\end{corollary}

실제로 $(X\times_S X)_{(S')}$가
$X_{(S')}\times_{S'}X_{(S')}$와 표준적으로 동일시됨을 주목하면
충분하다 (\hyperref[I.3.3.10-ko]{3.3.10}).

\begin{proposition}[5.3.5]
\label{I.5.3.5-ko}
$X$, $Y$를 두 $S$-준스킴이라 하고, $\varphi:S\to T$를 모든
$S$-준스킴에 $T$-준스킴의 구조를 주는 사상이라 하자. $f:X\to S$,
$g:Y\to S$를 구조 사상, $p$, $q$를 $X\times_S Y$의 사영,
$\pi=f\circ p=g\circ q$를 구조 사상 $X\times_S Y\to S$라 하자.
그러면 그림
\[
  \xymatrix{
    X\times_S Y\ar[r]^{(p,q)_T}\ar[d]_\pi &
    X\times_T Y\ar[d]^{f\times_T g}\\
    S\ar[r]_{\Delta_{S|T}} &
    S\times_T S
  }
  \tag{5.3.5.1}\label{I.5.3.5.1-ko}
\]
\oldpage[I]{133}
은 가환이고, $X\times_S Y$를 $(S\times_T S)$-준스킴 $S$와
$X\times_T Y$의 곱과 동일시한다. 이때 사영들은 각각 $\pi$와
$(p,q)_T$로 동일시된다.
\end{proposition}

(\hyperref[I.3.4.3-ko]{3.4.3})에 따라, $Z$를 임의의 $T$-준스킴이라
하고 $X$, $Y$, $S$를 각각 $X(Z)_T$, $Y(Z)_T$, $S(Z)_T$로 바꾸어
얻는 \emph{집합의 범주 안에서}의 대응하는 명제를 증명하면 된다.
집합의 범주에서는 확인이 즉각적이므로 독자에게 맡긴다.

\begin{corollary}[5.3.6]
\label{I.5.3.6-ko}
$P=S\times_T S$라 놓으면 사상 $(p,q)_T$는
$1_{X\times_T Y}\times_P\Delta_S$와 동일시된다.
\end{corollary}

이는 (\hyperref[I.5.3.5-ko]{5.3.5})와
(\hyperref[I.3.3.4-ko]{3.3.4})에서 따른다.

\begin{corollary}[5.3.7]
\label{I.5.3.7-ko}
$f:X\to Y$가 $S$-사상이면 그림
\[
  \xymatrix{
    X\ar[r]^{(1_X,f)_S}\ar[d]_f &
    X\times_S Y\ar[d]^{f\times_S 1_Y}\\
    Y\ar[r]_{\Delta_Y} &
    Y\times_S Y
  }
\]
은 가환이고, $X$를 $(Y\times_S Y)$-준스킴 $Y$와
$X\times_S Y$의 곱과 동일시한다.
\end{corollary}

(\hyperref[I.5.3.5-ko]{5.3.5})에서 $S$를 $Y$로, $T$를 $S$로
바꾸고 $X\times_Y Y=X$임을 주목하면 충분하다
(\hyperref[I.3.3.3-ko]{3.3.3}).

\begin{proposition}[5.3.8]
\label{I.5.3.8-ko}
$f:X\to Y$가 준스킴의 단사사상이기 위한 필요충분조건은
$\Delta_{X|Y}$가 $X$에서 $X\times_Y X$ 위로 가는 동형사상인 것이다.
\end{proposition}

실제로 $f$가 단사사상이라는 것은 모든 $Y$-준스킴 $Z$에 대하여
대응하는 사상 $f':X(Z)_Y\to Y(Z)_Y$가 단사라는 뜻이다.
$Y(Z)_Y$에는 원소가 하나뿐이므로, 이는 $X(Z)_Y$의 원소가 많아야
하나라는 뜻이다. 이는 다시 $X(Z)_Y\times X(Z)_Y$가
$X(Z)_Y$와 표준적으로 동형이라는 말과 같다. 첫째 집합은
(\hyperref[I.3.4.3.1-ko]{3.4.3.1})에 따라
$(X\times_Y X)(Z)_Y$이므로, 이는 $\Delta_{X|Y}$가 동형사상이라는
뜻이다.

\begin{proposition}[5.3.9]
\label{I.5.3.9-ko}
대각 사상 $\Delta_X$는 $X$에서 $X\times_S X$ 안으로 가는 몰입
사상이다.
\end{proposition}

(\hyperref[I.4.2.4-ko]{4.2.4 (a)})를 사용하면 충분하다. 실제로
$X$의 모든 $x$와 $x$의 모든 아핀 근방 $U$에 대하여,
$U\times_S U$는 $\Delta_X(x)$의 아핀 근방이다.
(\hyperref[I.5.3.16-ko]{5.3.16})의 증명은 정의
(\hyperref[I.5.3.1.1-ko]{5.3.1.1})만을 사용하므로, $S=\operatorname{Spec}(B)$,
$X=\operatorname{Spec}(A)$가 아핀 스킴인 경우만 증명하면 된다. 이때
(\hyperref[I.5.3.1.1-ko]{5.3.1.1})에서 $\Delta_X$가
\[
A\otimes_B A\longrightarrow A,\qquad x\otimes y\longmapsto xy
\]
인 표준 준동형에 대응함이 곧바로 따른다. 이 준동형은 전사이므로,
$\Delta_X$는 이 경우 닫힌 몰입 사상이다
(\hyperref[I.4.2.3-ko]{4.2.3}). 이것으로 증명이 끝난다.

몰입 사상 $\Delta_X$에 대응하는 $X\times_S X$의 부분준스킴
(\hyperref[I.4.2.1-ko]{4.2.1})을 $X\times_S X$의
\emph{대각선}이라 한다.

\begin{corollary}[5.3.10]
\label{I.5.3.10-ko}
(\hyperref[I.5.3.5-ko]{5.3.5})의 가정 아래에서 $(p,q)_T$는 몰입
사상이다.
\end{corollary}

이는 (\hyperref[I.5.3.6-ko]{5.3.6})과
(\hyperref[I.4.3.1-ko]{4.3.1})에서 따른다.

(\hyperref[I.5.3.5-ko]{5.3.5})의 가정 아래에서 $(p,q)_T$를
$X\times_S Y$에서 $X\times_T Y$ 안으로 가는 \emph{표준 몰입}이라
한다.

\begin{corollary}[5.3.11]
\label{I.5.3.11-ko}
$X$, $Y$를 두 $S$-준스킴, $f:X\to Y$를 $S$-사상이라 하자. 그러면
$f$의 그래프 사상 $\Gamma_f=(1_X,f)_S$
(\hyperref[I.3.3.14-ko]{3.3.14})는 $X$에서 $X\times_S Y$ 안으로 가는
몰입 사상이다.
\end{corollary}

이는 (\hyperref[I.5.3.10-ko]{5.3.10})에서 $S$를 $Y$로, $T$를 $S$로
바꾼 특별한 경우이다 (\hyperref[I.5.3.7-ko]{5.3.7} 참조).

\oldpage[I]{134}
몰입 사상 $\Gamma_f$에 대응하는 $X\times_S Y$의 부분준스킴
(\hyperref[I.4.2.1-ko]{4.2.1})을 사상 $f$의 \emph{그래프}라 한다.
사상 $X\to Y$의 그래프인 $X\times_S Y$의 부분준스킴들은 다음
성질로 특징지어진다. 첫째 사영 $p_1:X\times_S Y\to X$를 그러한
부분준스킴 $G$에 제한한 사상은 $G$에서 $X$ 위로 가는
\emph{동형사상} $g$이고, 이때 $G$는 사상 $p_2\circ g^{-1}$의
그래프이다. 여기서 $p_2$는 사영 $X\times_S Y\to Y$이다.

특히 $X=S$로 취하면, $S$-사상 $S\to Y$, 곧 $Y$의 $S$-단면
(\hyperref[I.2.5.5-ko]{2.5.5})은 자신의 그래프 사상과 같다.
$S$-단면의 그래프인 $Y$의 부분준스킴들, 다시 말해 구조 사상
$Y\to S$의 제한에 의해 $S$와 동형인 부분준스킴들을 이 단면들의
\emph{상}, 또는 관용적으로 $Y$의 \emph{$S$-단면}이라고도 한다.

\begin{corollary}[5.3.12]
\label{I.5.3.12-ko}
(\hyperref[I.5.3.11-ko]{5.3.11})의 가정과 표기를 사용하자. 모든 사상
$g:S'\to S$에 대하여, $f'$을 $g$에 의한 $f$의 역상이라 하면
(\hyperref[I.3.3.7-ko]{3.3.7}), $\Gamma_{f'}$은 $g$에 의한
$\Gamma_f$의 역상이다.
\end{corollary}

이는 공식 (\hyperref[I.3.3.10.1-ko]{3.3.10.1})의 특별한 경우이다.

\begin{corollary}[5.3.13]
\label{I.5.3.13-ko}
$f:X\to Y$, $g:Y\to Z$를 두 사상이라 하자. $g\circ f$가 몰입
사상(각각 국소 몰입 사상)이면 $f$도 그러하다.
\end{corollary}

실제로 $f$는
\[
  X\xrightarrow{\Gamma_f}X\times_Z Y\xrightarrow{p_2}Y
\]
로 인수분해된다. 한편 $p_2$는 $(g\circ f)\times_Z 1_Y$와
동일시된다 (\hyperref[I.3.3.4-ko]{3.3.4}). $g\circ f$가 몰입
사상(각각 국소 몰입 사상)이면 $p_2$도 그러하다
(\hyperref[I.4.3.1-ko]{4.3.1}과 \hyperref[I.4.5.5-ko]{4.5.5}). 또한
$\Gamma_f$는 몰입 사상이므로 (\hyperref[I.5.3.11-ko]{5.3.11}),
(\hyperref[I.4.2.5-ko]{4.2.5}) (각각
\hyperref[I.4.5.5-ko]{4.5.5})에서 결론이 따른다.

\begin{corollary}[5.3.14]
\label{I.5.3.14-ko}
$j:X\to Y$, $g:X\to Z$를 두 $S$-사상이라 하자. $j$가 몰입
사상(각각 국소 몰입 사상)이면 $(j,g)_S$도 그러하다.
\end{corollary}

실제로 $p:Y\times_S Z\to Y$가 첫째 사영이면
$j=p\circ(j,g)_S$이고, (\hyperref[I.5.3.13-ko]{5.3.13})을 적용하면
충분하다.

\begin{proposition}[5.3.15]
\label{I.5.3.15-ko}
$f:X\to Y$가 $S$-사상이면 그림
\[
  \xymatrix{
    X\ar[r]^{\Delta_X}\ar[d]_f &
    X\times_S X\ar[d]^{f\times_S f}\\
    Y\ar[r]_{\Delta_Y} &
    Y\times_S Y
  }
  \tag{5.3.15.1}\label{I.5.3.15.1-ko}
\]
은 가환이다. 다시 말해 $\Delta_X$는 준스킴의 범주에서 함자적
사상이다.
\end{proposition}

확인은 즉각적이므로 독자에게 맡긴다.

\begin{corollary}[5.3.16]
\label{I.5.3.16-ko}
$X$가 $Y$의 부분준스킴이면 대각선 $\Delta_X(X)$는
$\Delta_Y(Y)$의 부분준스킴과 동일시되며, 그 바탕공간은
\[
  \Delta_Y(Y)\cap p_1^{-1}(X)=\Delta_Y(Y)\cap p_2^{-1}(X)
\]
와 동일시된다. 여기서 $p_1$, $p_2$는 $Y\times_S Y$의 사영이다.
\end{corollary}

포함 사상 $f:X\to Y$에 (\hyperref[I.5.3.15-ko]{5.3.15})를 적용하자.
그러면 $f\times_S f$는 몰입 사상이고, $X\times_S X$의 바탕공간을
$Y\times_S Y$의 부분공간
$p_1^{-1}(X)\cap p_2^{-1}(X)$와 동일시한다
(\hyperref[I.4.3.1-ko]{4.3.1}). 또한
$z\in\Delta_Y(Y)\cap p_1^{-1}(X)$이면 $z=\Delta_Y(y)$이고

\oldpage[I]{135}
$y=p_1(z)\in X$이다. 따라서 포함 사상 아래에서 $y=f(y)$로
동일시되며, 그림 (\hyperref[I.5.3.15.1-ko]{5.3.15.1})의 가환성에
따라 $z=\Delta_Y(f(y))$는 $\Delta_X(X)$에 속한다.

\begin{corollary}[5.3.17]
\label{I.5.3.17-ko}
$f_1:Y\to X$, $f_2:Y\to X$를 두 $S$-사상이라 하고, $y$를
$f_1(y)=f_2(y)=x$이면서 $f_1$, $f_2$에 대응하는 준동형
$k(x)\to k(y)$가 서로 같은 $Y$의 점이라 하자. 그러면
$f=(f_1,f_2)_S$라 할 때, 점 $f(y)$는 대각선
$\Delta_{X|S}(X)$에 속한다.
\end{corollary}

$f_i$ ($i=1,2$)에 대응하는 두 준동형 $k(x)\to k(y)$는 두
$S$-사상
$g_i:\operatorname{Spec}(k(y))\to\operatorname{Spec}(k(x))$를 정의하며,
그림
\[
  \xymatrix{
    \operatorname{Spec}(k(y))\ar[r]^{g_i}\ar[d] &
    \operatorname{Spec}(k(x))\ar[d]\\
    Y\ar[r]_{f_i} & X
  }
\]
은 가환이다. 따라서 그림
\[
  \xymatrix{
    \operatorname{Spec}(k(y))\ar[rr]^{(g_1,g_2)_S}\ar[d] & &
    \operatorname{Spec}(k(x))\times_S\operatorname{Spec}(k(x))\ar[d]\\
    Y\ar[rr]_{(f_1,f_2)_S} & & X\times_S X
  }
\]
도 가환이다. 그런데 $g_1=g_2$이므로, $(g_1,g_2)_S$에 의한
$\operatorname{Spec}(k(y))$의 유일한 점의 상은
$\operatorname{Spec}(k(x))\times_S\operatorname{Spec}(k(x))$의 대각선에
속한다. 따라서 결론은 (\hyperref[I.5.3.15-ko]{5.3.15})에서 따른다.

\subsection{분리 사상과 분리 준스킴}
\label{subsection:I.5.4-ko}

\begin{definition}[5.4.1]
\label{I.5.4.1-ko}
준스킴의 사상 $f:X\to Y$에서 대각 사상
$X\to X\times_Y X$가 \emph{닫힌 몰입 사상}이면 $f$를
\emph{분리} 사상이라고 한다. 이때 $X$를 \emph{$Y$ 위에서 분리된
준스킴} 또는 \emph{$Y$-스킴}이라고도 한다. 준스킴 $X$가
$\operatorname{Spec}(\mathbf{Z})$ 위에서 분리되어 있으면 $X$를 분리
준스킴이라 하고, 이때 $X$를 \emph{스킴}이라고도 한다
(\agref{I.5.5.7-ko}{5.5.7} 참조).
\end{definition}

(\hyperref[I.5.3.9-ko]{5.3.9})에 따라 $X$가 $Y$ 위에서 분리되어 있기
위한 필요충분조건은 $\Delta_X(X)$가 $X\times_Y X$의 바탕공간의
\emph{닫힌 부분공간}인 것이다.

\begin{proposition}[5.4.2]
\label{I.5.4.2-ko}
$S\to T$를 분리 사상이라 하자. $X$, $Y$가 두 $S$-준스킴이면 표준
몰입 $X\times_S Y\to X\times_T Y$
(\hyperref[I.5.3.10-ko]{5.3.10})는 닫힌 몰입 사상이다.
\end{proposition}

실제로 그림 (\hyperref[I.5.3.5.1-ko]{5.3.5.1})을 보면 $(p,q)_T$는
$\Delta_{S|T}$에서, 사상
$f\times_T g:X\times_T Y\to S\times_T S$에 의한 밑준스킴
$S\times_T S$의 확장을 통해 얻어진 것으로 볼 수 있다. 따라서 명제는
(\hyperref[I.4.3.2-ko]{4.3.2})에서 따른다.

\begin{corollary}[5.4.3]
\label{I.5.4.3-ko}
$Y$를 $S$-스킴, $f:X\to Y$를 $S$-사상이라 하자. 그러면 $f$의
그래프 사상 $\Gamma_f:X\to X\times_S Y$
(\hyperref[I.5.3.11-ko]{5.3.11})는 닫힌 몰입 사상이다.
\end{corollary}

이는 (\hyperref[I.5.4.2-ko]{5.4.2})에서 $S$를 $Y$로, $T$를 $S$로
바꾼 특별한 경우이다.

\begin{corollary}[5.4.4]
\label{I.5.4.4-ko}
$f:X\to Y$, $g:Y\to Z$를 두 사상이라 하고 $g$가 분리 사상이라고
하자. $g\circ f$가 닫힌 몰입 사상이면 $f$도 닫힌 몰입 사상이다.
\end{corollary}

(\hyperref[I.5.4.3-ko]{5.4.3})에서 출발하는 증명은
(\hyperref[I.5.3.11-ko]{5.3.11})에서 출발하는
(\hyperref[I.5.3.13-ko]{5.3.13})의 증명과 같다.

\oldpage[I]{136}
\begin{corollary}[5.4.5]
\label{I.5.4.5-ko}
$Z$를 $S$-스킴, $j:X\to Y$, $g:X\to Z$를 두 $S$-사상이라 하자.
$j$가 닫힌 몰입 사상이면
$(j,g)_S:X\to Y\times_S Z$도 닫힌 몰입 사상이다.
\end{corollary}

(\hyperref[I.5.4.4-ko]{5.4.4})에서 출발하는 증명은
(\hyperref[I.5.3.13-ko]{5.3.13})에서 출발하는
(\hyperref[I.5.3.14-ko]{5.3.14})의 증명과 같다.

\begin{corollary}[5.4.6]
\label{I.5.4.6-ko}
$X$가 $S$-스킴이면 $X$의 모든 $S$-단면
(\hyperref[I.2.5.5-ko]{2.5.5})은 닫힌 몰입 사상이다.
\end{corollary}

$\varphi:X\to S$가 구조 사상이고 $\psi:S\to X$가 $X$의
$S$-단면이면, $\varphi\circ\psi=1_S$에
(\hyperref[I.5.4.5-ko]{5.4.5})를 적용하면 충분하다.

\begin{corollary}[5.4.7]
\label{I.5.4.7-ko}
$S$를 정역 준스킴, $s$를 그 일반점, $X$를 $S$-스킴이라 하자.
$X$의 두 $S$-단면 $f$, $g$가 $f(s)=g(s)$를 만족하면 $f=g$이다.
\end{corollary}

실제로 $x=f(s)=g(s)$라 하면, $f$와 $g$에 대응하는 준동형
$k(x)\to k(s)$는 반드시 서로 같다. $h=(f,g)_S$라 하면
(\hyperref[I.5.3.17-ko]{5.3.17})에 따라 $h(s)$는 대각선
$Z=\Delta_X(X)$에 속한다. 그런데 $S=\overline{\{s\}}$이고 가정에
따라 $Z$가 닫혀 있으므로 $h(S)\subset Z$이다. 따라서
(\hyperref[I.5.2.2-ko]{5.2.2})에 따라 $h$는
$S\to Z\to X\times_S X$로 인수분해되고, 대각선의 정의에서
$f=g$가 따른다.

\begin{remark}[5.4.8]
\label{I.5.4.8-ko}
역으로 (\hyperref[I.5.4.3-ko]{5.4.3})의 결론이 $f=1_Y$일 때
성립한다고 가정하면 $Y$가 $S$ 위에서 분리되어 있음을 얻는다.
마찬가지로 (\hyperref[I.5.4.5-ko]{5.4.5})의 결론이 두 사상
\[
  Y\xrightarrow{\Delta_Y}Y\times_Z Y\xrightarrow{p_1}Y
\]
에 적용된다고 가정하면 $\Delta_Y$가 닫힌 몰입 사상이고, 따라서
$Y$가 $Z$ 위에서 분리되어 있음을 얻는다. 마지막으로
(\hyperref[I.5.4.6-ko]{5.4.6})의 결론이 $Y$-준스킴
$Y\times_S Y\to Y$의 $Y$-단면 $\Delta_Y$에 대하여 성립하면,
$Y$는 $S$ 위에서 분리되어 있다.
\end{remark}

\subsection{분리 판정법}
\label{subsection:I.5.5-ko}

\begin{proposition}[5.5.1]
\label{I.5.5.1-ko}
\begin{enumerate}
  \item[\rm(i)] 준스킴의 모든 단사사상(특히 모든 몰입 사상)은 분리
    사상이다.
  \item[\rm(ii)] 두 분리 사상의 합성은 분리 사상이다.
  \item[\rm(iii)] $f:X\to X'$, $g:Y\to Y'$가 두 분리 $S$-사상이면
    $f\times_S g$도 분리 사상이다.
  \item[\rm(iv)] $f:X\to Y$가 분리 $S$-사상이면, 밑준스킴의 모든 확장
    $S'\to S$에 대하여 $S'$-사상 $f_{(S')}$도 분리 사상이다.
  \item[\rm(v)] 두 사상의 합성 $g\circ f$가 분리 사상이면 $f$도 분리
    사상이다.
  \item[\rm(vi)] 사상 $f$가 분리 사상이기 위한 필요충분조건은
    $f_{\mathrm{red}}$ (\hyperref[I.5.1.5-ko]{5.1.5})가 분리 사상인
    것이다.
\end{enumerate}
\end{proposition}

(i)는 (\hyperref[I.5.3.8-ko]{5.3.8})에서 곧바로 따른다. 두 사상
$f:X\to Y$, $g:Y\to Z$에 대하여 그림
\[
  \xymatrix{
    X\ar[rr]^{\Delta_{X|Z}}\ar[dr]_{\Delta_{X|Y}} & &
    X\times_Z X\\
    & X\times_Y X\ar[ur]_j
  }
  \tag{5.5.1.1}\label{I.5.5.1.1-ko}
\]
은 가환이다. 여기서 $j$는 표준 몰입 사상
(\hyperref[I.5.3.10-ko]{5.3.10})이며, 가환성은 곧바로 확인된다.
$f$와 $g$가 분리 사상이면 정의에 따라 $\Delta_{X|Y}$는 닫힌 몰입
사상이고, (\hyperref[I.5.4.2-ko]{5.4.2})에 따라 $j$도 닫힌 몰입
사상이다. 따라서 (\hyperref[I.4.2.5-ko]{4.2.5})에 따라
$\Delta_{X|Z}$도 닫힌 몰입 사상이며, 이로써

\oldpage[I]{137}
(ii)가 증명된다. (i)와 (ii)를 고려하면 (iii)와 (iv)는 서로
동치이고 (\hyperref[I.3.5.1-ko]{3.5.1}), 따라서 (iv)만 증명하면
충분하다. 그런데 (\hyperref[I.3.3.11-ko]{3.3.11})과
(\hyperref[I.3.3.9.1-ko]{3.3.9.1})에 따라
$X_{(S')}\times_{Y_{(S')}}X_{(S')}$은
$(X\times_Y X)\times_Y Y_{(S')}$과 표준적으로 동일시된다. 이때 대각
사상 $\Delta_{X_{(S')}}$은
$\Delta_X\times_Y 1_{Y_{(S')}}$과 동일시됨을 곧바로 확인할 수 있다.
따라서 명제는 (\hyperref[I.4.3.1-ko]{4.3.1})에서 따른다.

(v)를 증명하기 위하여 (\hyperref[I.5.3.13-ko]{5.3.13})에서와 같이
$f$의 인수분해
$X\xrightarrow{\Gamma_f}X\times_Z Y\xrightarrow{p_2}Y$를 생각하자.
여기서 $p_2=(g\circ f)\times_Z 1_Y$이다. $g\circ f$가 분리
사상이라는 가정과 (iii), (i)에 따라 $p_2$는 분리 사상이다.
$\Gamma_f$는 몰입 사상이므로 (i)에 따라 분리 사상이고, 따라서
(ii)에 따라 $f$도 분리 사상이다. 마지막으로 (vi)를 증명하기 위해,
준스킴 $X_{\mathrm{red}}\times_{Y_{\mathrm{red}}}X_{\mathrm{red}}$와
$X_{\mathrm{red}}\times_Y X_{\mathrm{red}}$가 표준적으로 동일시됨을
(\hyperref[I.5.1.7-ko]{5.1.7}) 상기하자. $j$를 포함 사상
$X_{\mathrm{red}}\to X$라 하면 그림
\[
  \xymatrix{
    X_{\mathrm{red}}\ar[r]^{\Delta_{X_{\mathrm{red}}}}\ar[d]_j &
    X_{\mathrm{red}}\times_Y X_{\mathrm{red}}\ar[d]^{j\times_Y j}\\
    X\ar[r]_{\Delta_X} & X\times_Y X
  }
\]
은 가환이다 (\hyperref[I.5.3.15-ko]{5.3.15}). 수직 화살표들이
바탕공간의 위상동형이라는 사실
(\hyperref[I.4.3.1-ko]{4.3.1})에서 명제가 따른다.

\begin{corollary}[5.5.2]
\label{I.5.5.2-ko}
$f:X\to Y$가 분리 사상이면, $X$의 모든 부분준스킴에 대한 $f$의
제한도 분리 사상이다.
\end{corollary}

이는 (\hyperref[I.5.5.1-ko]{5.5.1, (i)와 (ii)})에서 따른다.

\begin{corollary}[5.5.3]
\label{I.5.5.3-ko}
$X$, $Y$가 두 $S$-준스킴이고 $Y$가 $S$ 위에서 분리되어 있으면,
$X\times_S Y$는 $X$ 위에서 분리되어 있다.
\end{corollary}

이는 (\hyperref[I.5.5.1-ko]{5.5.1, (iv)})의 특별한 경우이다.

\begin{proposition}[5.5.4]
\label{I.5.5.4-ko}
$X$를 준스킴이라 하고, 그 바탕공간이 유한개의 닫힌 부분집합
$X_k$ $(1\leq k\leq n)$의 합집합이라고 하자. 각 $k$에 대하여
$X_k$를 바탕공간으로 갖는 $X$의 축소 부분준스킴
(\hyperref[I.5.2.1-ko]{5.2.1})을 생각하고, 이 부분준스킴도 $X_k$로
표시한다. $f:X\to Y$를 사상이라 하고, 각 $k$에 대하여
$f(X_k)\subset Y_k$를 만족하는 $Y$의 닫힌 부분집합 $Y_k$를
택하자. $Y_k$를 바탕공간으로 갖는 $Y$의 축소 부분준스킴도
$Y_k$로 표시하면, $X_k$에 대한 $f$의 제한 $X_k\to Y$는
$X_k\xrightarrow{f_k}Y_k\to Y$로 인수분해된다
(\hyperref[I.5.2.2-ko]{5.2.2}). 이때 $f$가 분리 사상이기 위한
필요충분조건은 모든 $f_k$가 분리 사상인 것이다.
\end{proposition}

필요성은 (\hyperref[I.5.5.1-ko]{5.5.1, (i), (ii), (v)})에서 따른다.
역으로 명제의 조건이 성립한다고 하자. 그러면 $X_k$에 대한 $f$의
각 제한 $X_k\to Y$는 분리 사상이다
(\hyperref[I.5.5.1-ko]{5.5.1, (i)와 (ii)}). $p_1$, $p_2$를
$X\times_Y X$의 사영이라 하면, 부분공간 $\Delta_{X_k}(X_k)$는
$X\times_Y X$의 바탕공간의 부분공간
$\Delta_X(X)\cap p_1^{-1}(X_k)$와 동일시된다
(\hyperref[I.5.3.16-ko]{5.3.16}). 이 부분공간들은
$X\times_Y X$에서 닫혀 있으므로, 그 합집합 $\Delta_X(X)$도 닫혀
있다.

특히 $X_k$들이 $X$의 \emph{기약 성분}이라고 하자. 그러면
$Y_k$들도 $Y$의 기약 성분이라고 가정할 수 있다
(\hyperref[0.2.1.5-ko]{0, 2.1.5}). 따라서 이 경우 명제
(\hyperref[I.5.5.4-ko]{5.5.4})는 분리의 개념을 \emph{정역 준스킴}
(\hyperref[I.2.1.8-ko]{2.1.8})의 경우로 환원한다.

\oldpage[I]{138}

\begin{proposition}[5.5.5]
\label{I.5.5.5-ko}
$(Y_\lambda)$를 준스킴 $Y$의 열린 덮개라 하자. 사상 $f:X\to Y$가
분리 사상이기 위한 필요충분조건은 각 제한
$f^{-1}(Y_\lambda)\to Y_\lambda$가 분리 사상인 것이다.
\end{proposition}

$X_\lambda=f^{-1}(Y_\lambda)$라 놓자. 곱
$X_\lambda\times_Y X_\lambda$와
$X_\lambda\times_{Y_\lambda}X_\lambda$가 동일하다는 사실
(\hyperref[I.3.2.5-ko]{3.2.5})과
(\hyperref[I.4.2.4-ko]{4.2.4, b)})를 고려하면, 모든 것은
$X_\lambda\times_Y X_\lambda$들이 $X\times_Y X$의 덮개를 이룸을
증명하는 것으로 귀착된다. 그런데
$Y_{\lambda\mu}=Y_\lambda\cap Y_\mu$,
$X_{\lambda\mu}=X_\lambda\cap X_\mu=f^{-1}(Y_{\lambda\mu})$라 놓으면,
$X_\lambda\times_Y X_\mu$는 곱
$X_{\lambda\mu}\times_{Y_{\lambda\mu}}X_{\lambda\mu}$와 동일시되고
(\hyperref[I.3.2.6.4-ko]{3.2.6.4}), 따라서
$X_{\lambda\mu}\times_Y X_{\lambda\mu}$와도 동일시된다
(\hyperref[I.3.2.5-ko]{3.2.5}). 결국 이것은
$X_\lambda\times_Y X_\lambda$의 열린집합과 동일시되므로, 우리의
주장이 증명된다 (\hyperref[I.3.2.7-ko]{3.2.7}).

명제 (\hyperref[I.5.5.5-ko]{5.5.5})에 따라 $Y$를 아핀 열린집합들로
덮으면, 분리 사상의 연구를 아핀 스킴에 값을 갖는 분리 사상의
연구로 환원할 수 있다.

\begin{proposition}[5.5.6]
\label{I.5.5.6-ko}
$Y$를 아핀 스킴, $X$를 준스킴, $(U_\alpha)$를 $X$의 아핀 열린
덮개라 하자. 사상 $f:X\to Y$가 분리 사상이기 위한 필요충분조건은
모든 지표쌍 $(\alpha,\beta)$에 대하여 $U_\alpha\cap U_\beta$가 아핀
열린집합이고, 환 $\Gamma(U_\alpha\cap U_\beta,\mathscr{O}_X)$가 두 환
$\Gamma(U_\alpha,\mathscr{O}_X)$와
$\Gamma(U_\beta,\mathscr{O}_X)$의 표준 상들의 합집합으로 생성되는
것이다.
\end{proposition}

$U_\alpha\times_Y U_\beta$들은 $X\times_Y X$의 열린 덮개를 이룬다
(\hyperref[I.3.2.7-ko]{3.2.7}). $p$, $q$를 $X\times_Y X$의 사영이라
하면
\[
  \begin{aligned}
    \Delta_X^{-1}(U_\alpha\times_Y U_\beta)
      &=\Delta_X^{-1}\bigl(p^{-1}(U_\alpha)\cap q^{-1}(U_\beta)\bigr)\\
      &=\Delta_X^{-1}\bigl(p^{-1}(U_\alpha)\bigr)
        \cap\Delta_X^{-1}\bigl(q^{-1}(U_\beta)\bigr)
        =U_\alpha\cap U_\beta.
  \end{aligned}
\]
따라서 모든 것은 $U_\alpha\cap U_\beta$에 대한 $\Delta_X$의 제한이
$U_\alpha\times_Y U_\beta$ 안의 닫힌 몰입 사상임을 표현하는 것으로
귀착된다. 그런데 이 제한은 정의에 따라
$(j_\alpha,j_\beta)_Y$와 다르지 않다. 여기서 $j_\alpha$ (각각
$j_\beta$)는 $U_\alpha\cap U_\beta$에서 $U_\alpha$ (각각
$U_\beta$)로 가는 포함 사상이다.
$U_\alpha\times_Y U_\beta$는 그 환이
$\Gamma(U_\alpha,\mathscr{O}_X)
 \otimes_{\Gamma(Y,\mathscr{O}_Y)}
 \Gamma(U_\beta,\mathscr{O}_X)$와 표준적으로 동형인 아핀 스킴이므로
(\hyperref[I.3.2.2-ko]{3.2.2}), $U_\alpha\cap U_\beta$는 아핀
스킴이어야 하고, 환 $A(U_\alpha\times_Y U_\beta)$에서
$\Gamma(U_\alpha\cap U_\beta,\mathscr{O}_X)$로 가는 사상
$h_\alpha\otimes h_\beta\to h_\alpha h_\beta$는 전사이어야 한다
(\hyperref[I.4.2.3-ko]{4.2.3}). 이로써 증명이 끝난다.

\begin{corollary}[5.5.7]
\label{I.5.5.7-ko}
아핀 스킴은 분리되어 있다(따라서 스킴이며, 이것이
(\hyperref[I.5.4.1-ko]{5.4.1})의 용어법을 정당화한다).
\end{corollary}

\begin{corollary}[5.5.8]
\label{I.5.5.8-ko}
$Y$를 아핀 스킴이라 하자. 사상 $f:X\to Y$가 분리 사상이기 위한
필요충분조건은 $X$가 분리되어 있는 것(다시 말해 $X$가 스킴인
것)이다.
\end{corollary}

실제로 (\hyperref[I.5.5.6-ko]{5.5.6})의 판정법은 $f$에 의존하지
않는다.

\begin{corollary}[5.5.9]
\label{I.5.5.9-ko}
사상 $f:X\to Y$가 분리 사상이기 위한 필요조건은 $Y$가 분리
준스킴을 유도하는 모든 열린집합 $U$에 대하여 유도 준스킴
$f^{-1}(U)$가 분리되어 있는 것이다. 이것이 모든 아핀 열린집합
$U\subset Y$에 대하여 성립하면 충분하다.
\end{corollary}

조건의 필요성은 (\hyperref[I.5.5.5-ko]{5.5.5})와
(\hyperref[I.5.5.1-ko]{5.5.1, (ii)})에서 따른다. 충분성은 $Y$에 아핀
열린 덮개가 존재한다는 사실을 고려하면
(\hyperref[I.5.5.5-ko]{5.5.5})와
(\hyperref[I.5.5.8-ko]{5.5.8})에서 따른다.

특히 $X$와 $Y$가 아핀 스킴이면, \emph{모든} 사상 $X\to Y$는 분리
사상이다.

\begin{proposition}[5.5.10]
\label{I.5.5.10-ko}
$Y$를 스킴, $f:X\to Y$를 사상이라 하자. $X$의 모든 아핀 열린집합
$U$와 $Y$의 모든 아핀 열린집합 $V$에 대하여
$U\cap f^{-1}(V)$는 아핀이다.
\end{proposition}

$p_1$, $p_2$를 $X\times_{\mathbf{Z}}Y$의 사영이라 하자. 부분공간
$U\cap f^{-1}(V)$는
$\Gamma_f(X)\cap p_1^{-1}(U)\cap p_2^{-1}(V)$의 $p_1$에 의한
상이다. 그런데 $p_1^{-1}(U)\cap p_2^{-1}(V)$는

\oldpage[I]{139}
준스킴 $U\times_{\mathbf{Z}}V$의 바탕공간과 동일시되므로
(\hyperref[I.3.2.7-ko]{3.2.7}) 아핀 스킴이다
(\hyperref[I.3.2.2-ko]{3.2.2}). 또한 $\Gamma_f(X)$는
$X\times_{\mathbf{Z}}Y$에서 닫혀 있으므로
(\hyperref[I.5.4.3-ko]{5.4.3}),
$\Gamma_f(X)\cap p_1^{-1}(U)\cap p_2^{-1}(V)$는
$U\times_{\mathbf{Z}}V$에서 닫혀 있다. 따라서 $\Gamma_f$에 대응하는
$X\times_{\mathbf{Z}}Y$의 부분준스킴
(\hyperref[I.4.2.1-ko]{4.2.1})이 그 바탕공간의 열린 부분집합
$\Gamma_f(X)\cap p_1^{-1}(U)\cap p_2^{-1}(V)$ 위에 유도하는 준스킴은
아핀 스킴의 닫힌 부분준스킴이고, 따라서 아핀 스킴이다
(\hyperref[I.4.2.3-ko]{4.2.3}). 이제 $\Gamma_f$가 몰입 사상이라는
사실에서 명제가 따른다.

\begin{examples}[5.5.11]
\label{I.5.5.11-ko}
예 (\hyperref[I.2.3.2-ko]{2.3.2})의 준스킴(<<체 $K$ 위의 사영
직선>>)은 \emph{분리되어 있다}. 실제로 $X$의 아핀 열린 덮개
$(X_1,X_2)$에 대하여 $X_1\cap X_2=U_{12}$는 아핀이고,
$\Gamma(U_{12},\mathscr{O}_X)$는 $f\in K[s]$인
$f(s)/s^m$ 꼴의 유리함수들로 이루어진 환으로서 $K[s]$와 $1/s$로
생성된다. 따라서 (\hyperref[I.5.5.6-ko]{5.5.6})의 조건들이
성립한다.

예 (\hyperref[I.2.3.2-ko]{2.3.2})에서와 같은 $X_1$, $X_2$,
$U_{12}$, $U_{21}$을 택하되, 이번에는 $u_{12}$로서 $f(s)$를
$f(t)$에 대응시키는 동형을 택하자. 그러면 붙이기로 \emph{분리되지
않은 정역 준스킴} $X$를 얻는다. 실제로
(\hyperref[I.5.5.6-ko]{5.5.6})의 첫째 조건은 성립하지만 둘째 조건은
성립하지 않는다. 여기서
$\Gamma(X,\mathscr{O}_X)\to\Gamma(X_1,\mathscr{O}_X)=K[s]$가
동형임은 곧바로 알 수 있다. 그 역동형은 전사 사상
$f:X\to\operatorname{Spec}(K[s])$를 정의한다. 또한
$\mathfrak{j}_y\neq(s)$인 모든
$y\in\operatorname{Spec}(K[s])$에 대하여 $f^{-1}(y)$는 한 점으로만
이루어지지만, $\mathfrak{j}_y=(s)$이면 $f^{-1}(y)$는 서로 다른
\emph{두} 점으로 이루어진다($X$를 <<점 $0$을 이중화한 $K$ 위의
아핀 직선>>이라고 한다).

(\hyperref[I.5.5.6-ko]{5.5.6})의 둘째 조건은 성립하지만
\emph{첫째 조건은 성립하지 않는} 예도 줄 수 있다. 먼저 체 $K$ 위의 두 부정원 다항식환
$A=K[s,t]$의 소 스펙트럼 $Y$에서 $D(s)$와 $D(t)$의 합집합인
열린집합 $U$는 \emph{아핀 열린집합이 아니다}. 실제로 $z$가 $U$ 위의
$\mathscr{O}_Y$의 단면이면, $s^m z$와 $t^n z$가 각각 $s$, $t$의
다항식의 $U$에 대한 제한이 되는 두 정수 $m\geq0$, $n\geq0$이
존재한다 (\hyperref[I.1.4.1-ko]{1.4.1}). 이는 단면 $z$가 $Y$ 전체
위의 단면으로 연장되어 $s$, $t$의 다항식과 동일시되는 경우에만
가능함이 분명하다. 만일 $U$가 아핀 열린집합이면 포함 사상
$U\to Y$는 동형이어야 하는데 (\hyperref[I.1.7.3-ko]{1.7.3}), 이는
$U\neq Y$이므로 모순이다.

이제 두 아핀 스킴 $Y_1$, $Y_2$를 각각 환
$A_1=K[s_1,t_1]$, $A_2=K[s_2,t_2]$의 소 스펙트럼이라 하자.
$U_{12}=D(s_1)\cup D(t_1)$,
$U_{21}=D(s_2)\cup D(t_2)$로 놓고, $u_{12}$로서
$f(s_1,t_1)$을 $f(s_2,t_2)$에 대응시키는 환의 동형에 대응하는
동형 $Y_2\to Y_1$의 $U_{21}$에 대한 제한을 택하자. 이렇게 하여
(\hyperref[I.5.5.6-ko]{5.5.6})의 첫째 조건은 성립하지 않고 둘째
조건은 성립하는 예를 얻는다(이렇게 얻은 정역 준스킴을 <<점 $0$을
이중화한 $K$ 위의 아핀 평면>>이라고 한다).
\end{examples}

\begin{remark}[5.5.12]
\label{I.5.5.12-ko}
\normalfont 준스킴의 사상에 대한 성질 \textbf{P}가 주어졌다고 하고,
다음 명제들을 생각하자.
\begin{enumerate}
  \item[\rm(i)] \emph{모든 닫힌 몰입 사상은 성질 \textbf{P}를 갖는다.}
  \item[\rm(ii)] \emph{성질 \textbf{P}를 갖는 두 사상의 합성은 성질
    \textbf{P}를 갖는다.}
  \item[\rm(iii)] \emph{$f:X\to X'$, $g:Y\to Y'$가 성질 \textbf{P}를
    갖는 두 $S$-사상이면, $f\times_S g$는 성질 \textbf{P}를 갖는다.}

\oldpage[I]{140}
  \item[\rm(iv)] \emph{$f:X\to Y$가 성질 \textbf{P}를 갖는
    $S$-사상이면, 밑준스킴의 확장 $S'\to S$로부터 얻는 모든
    $S'$-사상 $f_{(S')}$는 성질 \textbf{P}를 갖는다.}
  \item[\rm(v)] \emph{두 사상 $f:X\to Y$, $g:Y\to Z$의 합성
    $g\circ f$가 성질 \textbf{P}를 갖고 $g$가 분리 사상이면,
    $f$는 성질 \textbf{P}를 갖는다.}
  \item[\rm(vi)] \emph{사상 $f:X\to Y$가 성질 \textbf{P}를 가지면
    $f_{\mathrm{red}}$ (\hyperref[I.5.1.5-ko]{5.1.5})도 성질
    \textbf{P}를 갖는다.}
\end{enumerate}

이때 (i)와 (ii)가 성립한다고 가정하면 (iii)와 (iv)는
\emph{동치}이고, (v)와 (vi)는 (i), (ii), (iii)의
\emph{귀결}이다.

첫째 주장은 이미 증명하였다 (\hyperref[I.3.5.1-ko]{3.5.1}).
(\hyperref[I.5.3.13-ko]{5.3.13})의 $f$의 인수분해
$X\xrightarrow{\Gamma_f}X\times_Z Y\xrightarrow{p_2}Y$를 생각하자.
관계식 $p_2=(g\circ f)\times_Z 1_Y$에 따라 $g\circ f$가 성질
\textbf{P}를 가지면 (iii)에 의해 $p_2$도 성질 \textbf{P}를 갖는다.
$g$가 분리 사상이면 $\Gamma_f$는 닫힌 몰입 사상이고
(\hyperref[I.5.4.3-ko]{5.4.3}), 따라서 (i)에 의해 성질 \textbf{P}를
갖는다. 마지막으로 (ii)에 의해 $f$도 성질 \textbf{P}를 갖는다.

끝으로 가환 그림
\[
  \xymatrix{
    X_{\mathrm{red}}\ar[r]^{f_{\mathrm{red}}}\ar[d] &
    Y_{\mathrm{red}}\ar[d]\\
    X\ar[r]_f &
    Y
  }
\]
을 생각하자. 수직 화살표들은 닫힌 몰입 사상이고
(\hyperref[I.5.1.5-ko]{5.1.5}), 따라서 (i)에 의해 성질
\textbf{P}를 갖는다. $f$가 성질 \textbf{P}를 갖는다는 가정과
(ii)에 따라 합성
$X_{\mathrm{red}}\xrightarrow{f_{\mathrm{red}}}Y_{\mathrm{red}}\to Y$가
성질 \textbf{P}를 갖는다. 마지막으로 닫힌 몰입 사상은 분리
사상이므로 (\hyperref[I.5.5.1-ko]{5.5.1, (i)}), (v)에 의해
$f_{\mathrm{red}}$도 성질 \textbf{P}를 갖는다.

다음 명제들을 생각하면
\begin{enumerate}
  \item[\rm(i')] \emph{모든 몰입 사상은 성질 \textbf{P}를 갖는다.}
  \item[\rm(v')] \emph{$g\circ f$가 성질 \textbf{P}를 가지면 $f$도
    성질 \textbf{P}를 갖는다.}
\end{enumerate}
위의 논증에서 (v')은 (i'), (ii), (iii)의 귀결임을 알 수 있다.
\end{remark}

\begin{env}[5.5.13]
\label{I.5.5.13-ko}
(v)와 (vi)는 (i), (iii), 그리고 다음 명제의 귀결이기도 하다.
\begin{enumerate}
  \item[\rm(ii')] \emph{$j:X\to Y$가 닫힌 몰입 사상이고 $g:Y\to Z$가
    성질 \textbf{P}를 갖는 사상이면, $g\circ j$는 성질 \textbf{P}를
    갖는다.}
\end{enumerate}
마찬가지로 (v')은 (i'), (iii), 그리고 다음 명제의 귀결이다.
\begin{enumerate}
  \item[\rm(ii'')] \emph{$j:X\to Y$가 몰입 사상이고 $g:Y\to Z$가
    성질 \textbf{P}를 갖는 사상이면, $g\circ j$는 성질 \textbf{P}를
    갖는다.}
\end{enumerate}
실제로 이는 (\hyperref[I.5.5.12-ko]{5.5.12})의 논증에서 곧바로
따른다.
\end{env}

\section{유한성 조건}
\label{section:I.6-ko}

\subsection{뇌터 준스킴과 국소 뇌터 준스킴}
\label{subsection:I.6.1-ko}

\begin{definition}[6.1.1]
\label{I.6.1.1-ko}
준스킴 $X$가 환이 뇌터인 아핀 열린집합 $V_\alpha$들의 유한 합집합
(각각 합집합)이면 $X$를 \emph{뇌터}(각각 \emph{국소 뇌터})라 한다.
여기서 각 환은 $V_\alpha$ 위에 유도된 스킴의 환을 뜻한다.
\end{definition}

$X$가 국소 뇌터이면 구조층 $\mathscr{O}_X$가
\emph{연접인 환의 층}이라는 것은 문제가 국소적임과
(\hyperref[I.1.5.2-ko]{1.5.2})에서 곧바로 따른다.
$\mathscr{O}_X$-가군 $\mathscr{F}$가 \emph{연접}이면,
$\mathscr{F}$의 \emph{모든 준연접 $\mathscr{O}_X$-부분가군}
\oldpage[I]{141}
\emph{(각각 모든 준연접 $\mathscr{O}_X$-몫가군)}은
\emph{연접이다}.
실제로 이 문제도 국소적이며, 뇌터 환 위의 유한 생성 가군의 부분가군
(각각 몫가군)이 유한 생성이라는 사실과
(\hyperref[I.1.5.1-ko]{1.5.1}),
(\hyperref[I.1.4.1-ko]{1.4.1}),
(\hyperref[I.1.3.10-ko]{1.3.10})을 적용하면 충분하다. 특히
$\mathscr{O}_X$의 \emph{모든 준연접 아이디얼층은} \emph{연접이다}.

준스킴 $X$가 열린집합 $W_\lambda$들의 유한 합집합(각각 합집합)이고,
각 $W_\lambda$ 위에 유도된 준스킴이 뇌터(각각 국소 뇌터)이면 $X$가
뇌터(각각 국소 뇌터)임은 자명하다.

\begin{proposition}[6.1.2]
\label{I.6.1.2-ko}
준스킴 $X$가 뇌터이기 위한 필요충분조건은 $X$가 국소 뇌터이고 그
바탕공간이 준콤팩트인 것이다. 이때 $X$의 바탕공간은 뇌터이다.
\end{proposition}

\begin{proof}
첫째 주장은 정의와 (\hyperref[I.1.1.10-ko]{1.1.10 (ii)})에서 곧바로
따른다. 둘째 주장은 (\hyperref[I.1.1.6-ko]{1.1.6})과 뇌터
부분공간들의 유한 합집합인 모든 공간이 뇌터라는 사실
(\hyperref[0.2.2.3-ko]{0, 2.2.3})에서 따른다.
\end{proof}

\begin{proposition}[6.1.3]
\label{I.6.1.3-ko}
$X$를 환 $A$의 아핀 스킴이라 하자. 다음 조건들은 동치이다.
a) $X$는 뇌터이다. b) $X$는 국소 뇌터이다. c) $A$는 뇌터이다.
\end{proposition}

\begin{proof}
a)와 b)의 동치는 (\hyperref[I.6.1.2-ko]{6.1.2})와 모든 아핀 스킴의
바탕공간이 준콤팩트라는 사실
(\hyperref[I.1.1.10-ko]{1.1.10})에서 따른다. 또한 c)가 a)를
함의함은 자명하다. a)가 c)를 함의함을 보이자. 환이 뇌터인 아핀
열린집합 $V_i$들의 유한 덮개를 $X$에서 잡고, $V_i$ 위에 유도된
준스킴의 환을 $A_i$라 하자. $A$의 아이디얼들의 증가열
$(\mathfrak{a}_n)$을 잡으면 이에 표준적으로 전단사 대응하는
(\hyperref[I.1.3.7-ko]{1.3.7})
$\widetilde{A}\equiv\mathscr{O}_X$의 아이디얼층들의 증가열
$(\widetilde{\mathfrak{a}}_n)$이 존재한다. $(\mathfrak{a}_n)$이
어느 단계 이후 일정함을 보이려면 $(\widetilde{\mathfrak{a}}_n)$이
어느 단계 이후 일정함을
보이면 충분하다. 그런데 $\widetilde{\mathfrak{a}}_n|V_i$는 표준
포함 사상 $V_i\to X$에 의한 $\widetilde{\mathfrak{a}}_n$의
역상이므로 (\hyperref[0.5.1.4-ko]{0, 5.1.4})
$\mathscr{O}_X|V_i$의 준연접 아이디얼층이다. 따라서
$\widetilde{\mathfrak{a}}_n|V_i=\widetilde{\mathfrak{a}}_{ni}$이고,
여기서 $\mathfrak{a}_{ni}$는 $A_i$의 아이디얼이다
(\hyperref[I.1.3.7-ko]{1.3.7}). $A_i$가 뇌터이므로 모든 $i$에
대하여 $(\mathfrak{a}_{ni})$는 어느 단계 이후 일정하며, 이로부터
명제가 따른다.
\end{proof}

앞의 논증은 \emph{$X$가 뇌터 준스킴이면 $\mathscr{O}_X$의 연접
아이디얼층들의 모든 증가열은 어느 단계 이후 일정하다}라는 것도
증명한다.

\begin{proposition}[6.1.4]
\label{I.6.1.4-ko}
뇌터(각각 국소 뇌터) 준스킴의 모든 부분준스킴은 뇌터(각각 국소
뇌터)이다.
\end{proposition}

\begin{proof}
뇌터 준스킴 $X$에 대해서만 증명하면 충분하다. 또한 정의
(\hyperref[I.6.1.1-ko]{6.1.1})에 따라 $X$가 아핀 스킴인 경우로
곧바로 환원된다. $X$의 모든 부분준스킴은 어떤 열린집합 위에 유도된
준스킴의 닫힌 부분준스킴이므로
(\hyperref[I.4.1.3-ko]{4.1.3}), $Y$가 $X$의 닫힌 부분준스킴이거나
열린집합 위에 유도된 부분준스킴인 경우만 생각하면 된다. $Y$가
닫힌 경우는 자명하다. 실제로 $A$가 $X$의 환이면 $Y$는 $A$의 어떤
아이디얼 $\mathfrak{J}$에 대하여 환 $A/\mathfrak{J}$의 아핀
스킴이다(\hyperref[I.4.2.3-ko]{4.2.3}). $A$가 뇌터이므로
(\hyperref[I.6.1.3-ko]{6.1.3}) $A/\mathfrak{J}$도 뇌터이다.

이제 $Y$가 $X$의 열린집합이라고 하자. 바탕공간 $Y$는 뇌터이므로
(\hyperref[I.6.1.2-ko]{6.1.2}) 준콤팩트이고, 따라서
$f_i\in A$인 열린집합 $D(f_i)$들의 유한 합집합이다.
\oldpage[I]{142}
그러므로 $f\in A$에 대하여 $Y=D(f)$인 경우만 증명하면 된다. 이때
$Y$는 환이 $A_f$와 동형인 아핀 스킴이다
(\hyperref[I.1.3.6-ko]{1.3.6}). $A$가 뇌터이므로
(\hyperref[I.6.1.3-ko]{6.1.3}) $A_f$도 뇌터이다.
\end{proof}

\begin{env}[6.1.5]
\label{I.6.1.5-ko}
두 뇌터 $S$-준스킴의 \emph{곱}은 반드시 뇌터인 것은 아니다. 이는
두 뇌터 대수의 텐서곱이 반드시 뇌터 환인 것은 아니므로, 두
준스킴이 아핀인 경우에도 마찬가지이다
(\agref{I.6.3.8-ko}{6.3.8} 참조).
\end{env}

\begin{proposition}[6.1.6]
\label{I.6.1.6-ko}
$X$가 뇌터 준스킴이면 $\mathscr{O}_X$의 멱영근기
$\mathscr{N}_X$는 멱영이다.
\end{proposition}

실제로 $X$를 유한개의 아핀 열린집합 $U_i$로 덮을 수 있고,
$(\mathscr{N}_X|U_i)^{n_i}=0$인 정수 $n_i$가 존재함을 보이면
충분하다. $n$을 $n_i$들의 최댓값이라 하면
$\mathscr{N}_X^n=0$이다. 따라서 $A$가 뇌터 환이고
$X=\operatorname{Spec}(A)$가 아핀인 경우로 환원된다. 이때
(\hyperref[I.5.1.1-ko]{5.1.1})과
(\hyperref[I.1.3.13-ko]{1.3.13})에 따라 $A$의 멱영근기가 멱영이라는
사실(\cite{I-11}, p.~127, cor.~4)을 확인하면 충분하다.

\begin{corollary}[6.1.7]
\label{I.6.1.7-ko}
$X$를 뇌터 준스킴이라 하자. $X$가 아핀 스킴이기 위한
필요충분조건은 $X_{\mathrm{red}}$가 아핀 스킴인 것이다.
\end{corollary}

이는 (\hyperref[I.6.1.6-ko]{6.1.6})과
(\hyperref[I.5.1.10-ko]{5.1.10})에서 따른다.

\begin{lemma}[6.1.8]
\label{I.6.1.8-ko}
$X$를 위상공간, $x$를 $X$의 한 점, $U$를 기약 성분이 유한개뿐인
$x$의 열린 근방이라 하자. 그러면 $x$의 근방 $V$가 존재하여,
$V$에 포함되는 $x$의 모든 열린 근방이 연결이다.
\end{lemma}

$x$를 포함하지 않는 $U$의 기약 성분들을 $U_i$($1\leq i\leq m$)라
하자. 그 합집합의 $U$에서의 여집합은 $U$ 안에서, 따라서 $X$ 안에서
$x$의 열린 근방 $V$이다. 또한 이는 $U$와, $x$를 포함하지 않는
$X$의 기약 성분들의 합집합을 $X$에서 뺀 여집합의 교집합이다
(\hyperref[0.2.1.6-ko]{0, 2.1.6}). 이제 $V$에 포함되는 $x$의 열린
근방 $W$를 잡자. $W$의 기약 성분들은 $W$와 만나는 $U$의 기약
성분들과 $W$의 교집합들이므로
(\hyperref[0.2.1.6-ko]{0, 2.1.6}), 모두 $x$를 포함한다. 이 성분들은
연결이므로 $W$도 연결이다.

\begin{corollary}[6.1.9]
\label{I.6.1.9-ko}
국소 뇌터 위상공간은 국소 연결이다. 따라서 특히 그 연결 성분들은
열린집합이다.
\end{corollary}

\begin{proposition}[6.1.10]
\label{I.6.1.10-ko}
$X$를 국소 뇌터 위상공간이라 하자. 다음 조건들은 동치이다.
\begin{enumerate}
  \item[a)] $X$의 기약 성분들은 열린집합이다.
  \item[b)] $X$의 기약 성분들은 연결 성분들과 일치한다.
  \item[c)] $X$의 연결 성분들은 기약이다.
  \item[d)] $X$의 서로 다른 두 기약 성분은 만나지 않는다.
\end{enumerate}

또한 $X$가 준스킴이면 이 조건들은 다음 조건과도 동치이다.
\begin{enumerate}
  \item[e)] 모든 $x\in X$에 대하여
    $\operatorname{Spec}(\mathscr{O}_x)$는 기약이다. 다시 말해
    $\mathscr{O}_x$의 멱영근기는 소 아이디얼이다.
\end{enumerate}
\end{proposition}

기약공간은 연결이므로 a)는 b)를 함의한다. 실제로 a)에 따라 $X$의
기약 성분들은 동시에 열린집합이고 닫힌집합이다. b)가 c)를 함의함은
자명하다. 역으로 닫힌집합 $F$가
\oldpage[I]{143}
$X$의 연결 성분 $C$를 포함하면서 $C$와 다르면 $F$는 기약일 수
없다. 실제로 이 집합은 연결이 아니므로
$F$ 안에서 동시에 열리고 닫힌 서로소인 두 공집합 아닌 집합의
합집합이며, 이 두 집합은 $X$에서 닫혀 있다. 따라서 c)는 b)를
함의한다. 이로부터 서로 다른 두 연결 성분은 만나지 않으므로 c)가
d)를 함의함이 곧바로 따른다.

여기까지는 $X$가 국소 뇌터라는 사실을 사용하지 않았다. 이제 이
가정을 하고 d)가 a)를 함의함을 보이자. (\hyperref[0.2.1.6-ko]{0,
2.1.6})에 따라 $X$가 뇌터인 경우로 환원할 수 있고, 이때 $X$의 기약
성분은 유한개뿐이다. 이 성분들이 닫혀 있고 서로소이므로 열린집합이다.

끝으로 d)와 e)의 동치는 준스킴 $X$의 바탕공간이 국소 뇌터라는
가정 없이도 성립한다. (\hyperref[0.2.1.6-ko]{0, 2.1.6})에 따라
$X=\operatorname{Spec}(A)$가 아핀인 경우로 환원할 수 있다. $x$가
$X$의 기약 성분 하나에만 포함된다는 것은
$\mathfrak{j}_x$가 $A$의 극소 소 아이디얼 하나만을 포함한다는 뜻이고
(\hyperref[I.1.1.14-ko]{1.1.14}), 이는
$\mathfrak{j}_x\mathscr{O}_x$가 $\mathscr{O}_x$의 극소 소 아이디얼
하나만을 포함한다는 것과 동치이다. 이로부터 결론이 따른다.

\begin{corollary}[6.1.11]
\label{I.6.1.11-ko}
$X$를 국소 뇌터 공간이라 하자. $X$가 기약이기 위한 필요충분조건은
$X$가 연결이고 공집합이 아니며, $X$의 서로 다른 두 기약 성분이
만나지 않는 것이다. $X$가 준스킴이면 마지막 조건은 모든 $x\in X$에
대하여 $\operatorname{Spec}(\mathscr{O}_x)$가 기약이라는 것과
동치이다.
\end{corollary}

마지막 주장은 (\hyperref[I.6.1.10-ko]{6.1.10})에서 보였다. 따라서
첫째 주장의 조건들이 충분함만 증명하면 된다. 그런데
(\hyperref[I.6.1.10-ko]{6.1.10})에 따라 이 조건들은 $X$의 기약
성분들이 연결 성분들과 일치함을 함의하고, $X$가 연결이고 공집합이
아니므로 $X$는 기약이다.

\begin{corollary}[6.1.12]
\label{I.6.1.12-ko}
$X$를 국소 뇌터 준스킴이라 하자. $X$가 정역 준스킴이기 위한
필요충분조건은 $X$가 공집합이 아니고 연결이며 모든 $x\in X$에 대하여
$\mathscr{O}_x$가 정역인 것이다.
\end{corollary}

\begin{proposition}[6.1.13]
\label{I.6.1.13-ko}
$X$를 국소 뇌터 준스킴, $x\in X$를 한 점이라 하자. 국소환
$\mathscr{O}_x$의 멱영근기 $\mathscr{N}_x$가 소 아이디얼이면(각각
$\mathscr{O}_x$가 축소환이면, 정역이면), $x$의 기약인(각각 축소인,
정역인) 열린 근방 $U$가 존재한다.
\end{proposition}

$\mathscr{N}_x$가 소 아이디얼인 경우와 $\mathscr{N}_x=0$인 경우만
생각하면 충분하다. 셋째 가정은 이 두 가정을 함께 말한 것이기
때문이다. $\mathscr{N}_x$가 소 아이디얼이면 $x$는 $X$의 기약 성분
$Y$ 하나에만 속한다(\hyperref[I.6.1.10-ko]{6.1.10}). $x$를 포함하지
않는 $X$의 기약 성분들의 집합은 국소 유한이므로 그 합집합은 닫혀
있다. 따라서 이 합집합의 여집합 $U$는 열린집합이고 $Y$에 포함되므로
기약이다(\hyperref[0.2.1.6-ko]{0, 2.1.6}). $\mathscr{N}_x=0$이면
$x$의 어떤 근방에 속하는 모든 $y$에 대하여 $\mathscr{N}_y=0$이다.
실제로 $\mathscr{N}$은 준연접이고
(\hyperref[I.5.1.1-ko]{5.1.1}), $X$가 국소 뇌터이므로 연접이다.
따라서 결론은 (\hyperref[0.5.2.2-ko]{0, 5.2.2})에서 따른다.

\subsection{아르틴 준스킴}
\label{subsection:I.6.2-ko}

\begin{definition}[6.2.1]
\label{I.6.2.1-ko}
준스킴이 아핀이고 그 환이 아르틴 환이면 그 준스킴을
\emph{아르틴}이라고 한다.
\end{definition}

\oldpage[I]{144}
\begin{proposition}[6.2.2]
\label{I.6.2.2-ko}
준스킴 $X$에 대하여 다음 조건들은 동치이다.
\begin{enumerate}
  \item[a)] $X$는 아르틴 스킴이다.
  \item[b)] $X$는 뇌터이고 그 바탕공간은 이산이다.
  \item[c)] $X$는 뇌터이고 그 바탕공간의 모든 점은 닫힌점이다
    ($\mathrm{T}_1$ 조건).
\end{enumerate}

이 조건들이 성립하면 $X$의 바탕공간은 유한하고, $X$의 환 $A$는
$X$의 점들의 (아르틴) 국소환들의 직접곱이다.
\end{proposition}

a)가 마지막 주장을 함의함은 알려져 있다(\cite{I-13}, p.~205,
th.~3). 이때 $A$의 모든 소 아이디얼은 극대 아이디얼이고, $A$의 한
국소환 성분의 극대 아이디얼의 역상이다. 따라서 공간 $X$는 유한하고
이산이다. 그러므로 a)는 b)를 함의하고, b)가 c)를 함의함은 자명하다.
c)가 a)를 함의함을 보이기 위해 먼저 $X$가 유한함을 보이자. 실제로
$X$가 아핀인 경우로 환원할 수 있고, 모든 소 아이디얼이 극대
아이디얼인 뇌터 환은 아르틴 환임이 알려져 있다(\cite{I-13}, p.~203).
따라서 주장이 따른다. 이제 $X$의 바탕공간은 이산이고, 유한개의 점
$x_i$의 위상적 합이다. 국소환 $\mathscr{O}_{x_i}=A_i$들은 아르틴
환이며, $A\simeq\prod_i A_i$이고 $X$는 이 환 $A$의 소 스펙트럼인
아핀 스킴과 동형임이 자명하다(\hyperref[I.1.7.3-ko]{1.7.3}).

\subsection{유한형 사상}
\label{subsection:I.6.3-ko}

\begin{definition}[6.3.1]
\label{I.6.3.1-ko}
$Y$가 다음 성질을 갖는 아핀 열린집합들의 족 $(V_\alpha)$의 합집합이면
사상 $f:X\to Y$를 \emph{유한형}이라고 한다.
\begin{enumerate}
  \item[(P)] $f^{-1}(V_\alpha)$는 유한개의 아핀 열린집합
    $U_{\alpha i}$의 합집합이고, 각 환 $A(U_{\alpha i})$는
    $A(V_\alpha)$ 위의 유한 생성 대수이다.
\end{enumerate}
이때 $X$를 $Y$ 위의 유한형 준스킴 또는 유한형 $Y$-준스킴이라고도
한다.
\end{definition}

\begin{proposition}[6.3.2]
\label{I.6.3.2-ko}
$f:X\to Y$가 유한형 사상이면 $Y$의 모든 아핀 열린집합 $W$는
정의 (\hyperref[I.6.3.1-ko]{6.3.1})의 성질 (P)를 갖는다.
\end{proposition}

먼저 다음 보조정리를 증명하자.

\begin{lemma}[6.3.2.1]
\label{I.6.3.2.1-ko}
$T\subset Y$가 성질 (P)를 갖는 아핀 열린집합이면, 모든
$g\in A(T)$에 대하여 $D(g)$도 성질 (P)를 갖는다.
\end{lemma}

실제로 가정에 따라 $f^{-1}(T)$는 유한개의 아핀 열린집합 $Z_j$의
합집합이고, 각 $A(Z_j)$는 $A(T)$ 위의 유한 생성 대수이다. $f$를
$Z_j$에 제한한 사상에 대응하는 환 준동형을
$\varphi_j:A(T)\to A(Z_j)$라 하고
(\hyperref[I.2.2.4-ko]{2.2.4}), $g_j=\varphi_j(g)$라 놓자. 그러면
$f^{-1}(D(g))\cap Z_j=D(g_j)$이다
(\hyperref[I.1.2.2.2-ko]{1.2.2.2}). 한편
$A(D(g_j))=A(Z_j)_{g_j}=A(Z_j)[1/g_j]$는 $A(Z_j)$ 위의 유한 생성
대수이고, 따라서 가정에 의해 \emph{특히} $A(T)$ 위의 유한 생성
대수이며, 그러므로 $A(D(g))=A(T)[1/g]$ 위의 유한 생성 대수이기도 하다.
이로써 보조정리가 증명된다.

이제 $W$가 준콤팩트이므로(\hyperref[I.1.1.10-ko]{1.1.10}), 각
$g_i$가 어떤 환 $A(V_{\alpha(i)})$에 속하는
$D(g_i)\subset W$들로 이루어진 $W$의 유한 덮개가 존재한다. 각
$D(g_i)$도 준콤팩트이므로
$h_{ik}\in A(W)$인 집합 $D(h_{ik})$들의 유한 합집합이다. 표준 사상을
$\varphi_i:A(W)\to A(D(g_i))$라 하면
(\hyperref[I.1.2.2.2-ko]{1.2.2.2})에 따라
$D(h_{ik})=D(\varphi_i(h_{ik}))$이다.
(\hyperref[I.6.3.2.1-ko]{6.3.2.1})에 따라 각
$f^{-1}(D(h_{ik}))$에는 유한 아핀 열린 덮개 $(U_{ijk})$가 존재하고,
각 $A(U_{ijk})$는
$A(D(h_{ik}))=A(W)[1/h_{ik}]$ 위의 유한 생성 대수이다. 이로부터 명제가
따른다.

\oldpage[I]{145}
따라서 $Y$ 위의 유한형 준스킴이라는 개념은 \emph{$Y$ 위에서
국소적}이라고 할 수 있다.

\begin{proposition}[6.3.3]
\label{I.6.3.3-ko}
$X$, $Y$를 두 아핀 스킴이라 하자. $X$가 $Y$ 위에서 유한형이기 위한
필요충분조건은 $A(X)$가 $A(Y)$ 위의 유한 생성 대수인 것이다.
\end{proposition}

이 조건이 충분함은 자명하므로 필요함을 증명하자. $A=A(Y)$,
$B=A(X)$라 놓자. (\hyperref[I.6.3.2-ko]{6.3.2})에 따라 $X$에는 유한
아핀 열린 덮개 $(V_i)$가 존재하여, 각 환 $A(V_i)$는 $A$ 위의 유한 생성
대수이다. 또한 $V_i$들이 준콤팩트이므로 각각은 유한개의 열린집합
$D(g_{ij})\subset V_i$로 덮이며, 여기서 $g_{ij}\in B$이다. 표준 포함
사상 $V_i\to X$에 대응하는 준동형을
$\varphi_i:B\to A(V_i)$라 하면
\[
  B_{g_{ij}}=(A(V_i))_{\varphi_i(g_{ij})}
  =A(V_i)[1/\varphi_i(g_{ij})].
\]
따라서 $B_{g_{ij}}$는 $A$ 위의 유한 생성 대수이다. 그러므로
$V_i=D(g_i)$이고 $g_i\in B$인 경우로 환원할 수 있다. 가정에 따라
$B_{g_i}$가 $A$ 위에서, $b_i$가 유한집합 $F_i\subset B$를 움직일 때의
원소 $b_i/g_i^{n_i}$들로 생성되도록 하는 정수 $n_i\geq0$가 존재한다.
$g_i$들이 유한개뿐이므로 모든 $n_i$가 같은 정수 $n$이라고 해도 된다.
또한 $D(g_i)$들이 $X$를 덮으므로 $g_i$들이 $B$에서 생성하는
아이디얼은 $B$이다. 다시 말해
$\sum_i h_i g_i=1$인 $h_i\in B$들이 존재한다. 이제 $F_i$들, $g_i$들,
$h_i$들의 합집합인 $B$의 유한 부분집합을 $F$라 하자. $B$의 부분환
$B'=A[F]$가 $B$와 같음을 보이겠다. 가정에 따라 모든 $b\in B$와 모든
$i$에 대하여 $B_{g_i}$에서 $b$의 표준상은 $b'_i/g_i^{m_i}$ 꼴이고,
여기서 $b'_i\in B'$이다. $b'_i$에 $g_i$의 적당한 거듭제곱을 곱하여
모든 $m_i$가 같은 정수 $m$이라고 해도 된다. 따라서 분수환의 정의에
의해 $N\geq m$이고 모든 $i$에 대하여 $g_i^N b\in B'$인 정수
$N$이 존재한다($N$은 $b$에 의존한다). 한편 \emph{환 $B'$ 안에서}
$g_i^N$들은 아이디얼 $B'$를 생성한다. 실제로 $g_i$들이 그러하고
$h_i\in B'$이기 때문이다. 따라서
$\sum_i c_i g_i^N=1$인 $c_i\in B'$들이 존재하고,
$b=\sum_i c_i g_i^N b\in B'$이다. 이로써 증명이 끝난다.

\begin{proposition}[6.3.4]
\label{I.6.3.4-ko}
\begin{enumerate}
  \item[(i)] 모든 닫힌 몰입 사상은 유한형이다.
  \item[(ii)] 두 유한형 사상의 합성은 유한형이다.
  \item[(iii)] $f:X\to X'$, $g:Y\to Y'$가 두 유한형 $S$-사상이면
    $f\times_S g:X\times_S Y\to X'\times_S Y'$도 유한형이다.
  \item[(iv)] $f:X\to Y$가 유한형 $S$-사상이면, 밑준스킴의 모든 확장
    $g:S'\to S$에 대하여 $f_{(S')}:X_{(S')}\to Y_{(S')}$도 유한형이다.
  \item[(v)] 두 사상의 합성 $g\circ f$가 유한형이고 $g$가 분리
    사상이면, $f$는 유한형이다.
  \item[(vi)] 사상 $f$가 유한형이면 $f_{\mathrm{red}}$도 유한형이다.
\end{enumerate}
\end{proposition}

(\hyperref[I.5.5.12-ko]{5.5.12})에 따라 (i), (ii), (iv)만 증명하면
충분하다.

(i)을 보이기 위해서는 $X$가 $Y$의 닫힌 부분준스킴이고
$X\to Y$가 표준 포함 사상인 경우만 다루면 된다. 또한
(\hyperref[I.6.3.2-ko]{6.3.2})에 따라 $Y$가 아핀이라고 해도 된다.
이때 $X$도 아핀이고(\hyperref[I.4.2.3-ko]{4.2.3}), 그 환은
$A/\mathfrak{J}$ 꼴의 몫환과 동형이다. 여기서 $A$는 $Y$의 환이고
$\mathfrak{J}$는 $A$의 아이디얼이다. $A/\mathfrak{J}$는 $A$ 위의
유한 생성 대수이므로 결론이 따른다.

이제 (ii)를 증명하자. $f:X\to Y$, $g:Y\to Z$를 두 유한형 사상이라
하고, $U$를 $Z$의 아핀 열린집합이라 하자. $g^{-1}(U)$에는 유한 아핀
열린 덮개 $(V_i)$가 존재하여 각 $A(V_i)$는 $A(U)$ 위의 유한 생성
대수이다(\hyperref[I.6.3.2-ko]{6.3.2}). 마찬가지로

\oldpage[I]{146}
각 $f^{-1}(V_i)$에는 유한 아핀 열린 덮개 $(W_{ij})$가 존재하여
$A(W_{ij})$는 $A(V_i)$ 위의 유한 생성 대수이고, 따라서 $A(U)$ 위의
유한 생성 대수이기도 하다. 이로부터 결론이 따른다.

마지막으로 (iv)를 증명하기 위해 $S=Y$인 경우만 다루면 된다. 실제로
$f$를 $Y$-사상으로 보고 밑준스킴의 확장을 $Y_{(S')}\to Y$로 보면
$f_{(S')}$는 $f_{(Y_{(S')})}$와도 같다
(\hyperref[I.3.3.9-ko]{3.3.9}). 이제 $p$, $q$를 각각 사영
$X_{(S')}\to X$, $X_{(S')}\to S'$라 하자. $V$를 $S$의 아핀
열린집합이라 하자. $f^{-1}(V)$는 유한개의 아핀 열린집합 $W_i$의
합집합이고, 각 $A(W_i)$는 $A(V)$ 위의 유한 생성 대수이다
(\hyperref[I.6.3.2-ko]{6.3.2}). $V'$를 $g^{-1}(V)$에 포함되는 $S'$의
아핀 열린집합이라 하자. $f\circ p=g\circ q$이므로 $q^{-1}(V')$는
$p^{-1}(W_i)$들의 합집합에 포함된다. 한편 교집합
$p^{-1}(W_i)\cap q^{-1}(V')$는 곱 $W_i\times_V V'$와 동일시된다
(\hyperref[I.3.2.7-ko]{3.2.7}). 이것은 환
$A(W_i)\otimes_{A(V)}A(V')$와 동형인 환을 갖는 아핀 스킴이다
(\hyperref[I.3.2.2-ko]{3.2.2}). 이 환은 가정에 따라 $A(V')$ 위의
유한 생성 대수이므로 명제가 증명된다.

\begin{corollary}[6.3.5]
\label{I.6.3.5-ko}
$f:X\to Y$를 몰입 사상이라 하자. $Y$의 바탕공간이 국소 뇌터이거나
$X$의 바탕공간이 뇌터이면 $f$는 유한형이다.
\end{corollary}

(\hyperref[I.6.3.2-ko]{6.3.2})에 따라 언제나 $Y$가 아핀이라고 해도
된다. $Y$의 바탕공간이 국소 뇌터이면 더 나아가 뇌터라고 해도 되고,
그 부분공간인 $X$의 바탕공간도 뇌터이다. 따라서 $Y$가 아핀이고
$X$의 바탕공간이 뇌터인 경우만 다루면 된다. 이때 $X$에는
$g_i\in A(Y)$인 유한개의 아핀 열린집합 $D(g_i)\subset Y$로 된 덮개가
존재하여, 각 $X\cap D(g_i)$는 $D(g_i)$에서 닫혀 있고 따라서 아핀
스킴이다(\hyperref[I.4.2.3-ko]{4.2.3}). 그런 덮개가 존재하는 것은 $X$가
$Y$에서 국소 닫혀 있기 때문이다(\hyperref[I.4.1.3-ko]{4.1.3}).
(\hyperref[I.6.3.4-ko]{6.3.4, (i)})와
(\hyperref[I.6.3.3-ko]{6.3.3})에 따라 $A(X\cap D(g_i))$는
$A(D(g_i))$ 위의 유한 생성 대수이다. 또한
$A(D(g_i))=A(Y)_{g_i}=A(Y)[1/g_i]$는 $A(Y)$ 위의 유한 생성 대수이므로
증명이 끝난다.

\begin{corollary}[6.3.6]
\label{I.6.3.6-ko}
$f:X\to Y$, $g:Y\to Z$를 두 사상이라 하자. $g\circ f$가 유한형이고,
$X$가 뇌터이거나 $X\times_Z Y$가 국소 뇌터이면 $f$는 유한형이다.
\end{corollary}

이는 (\hyperref[I.5.5.12-ko]{5.5.12})의 증명과 몰입 사상
$\Gamma_f$에 적용한 (\hyperref[I.6.3.5-ko]{6.3.5})에서 곧바로 따른다.

\begin{proposition}[6.3.7]
\label{I.6.3.7-ko}
$f:X\to Y$를 유한형 사상이라 하자. $Y$가 뇌터이면 $X$도 뇌터이고,
$Y$가 국소 뇌터이면 $X$도 국소 뇌터이다.
\end{proposition}

$Y$가 뇌터인 경우만 증명하면 된다. 이때 $Y$는 환 $A(V_i)$가 뇌터
환인 유한개의 아핀 열린집합 $V_i$의 합집합이다.
(\hyperref[I.6.3.2-ko]{6.3.2})에 따라 각 $f^{-1}(V_i)$는 유한개의
아핀 열린집합 $W_{ij}$의 합집합이고, 각 $A(W_{ij})$는 $A(V_i)$ 위의
유한 생성 대수이므로 뇌터 환이다. 따라서 $X$는 뇌터이다.

\begin{corollary}[6.3.8]
\label{I.6.3.8-ko}
$X$를 $S$ 위의 유한형 준스킴이라 하자. $S'$가 뇌터인(각각 국소
뇌터인) 밑준스킴의 모든 확장 $S'\to S$에 대하여 $X_{(S')}$는
뇌터이다(각각 국소 뇌터이다).
\end{corollary}

이는 (\hyperref[I.6.3.7-ko]{6.3.7})에서 따른다. 실제로
(\hyperref[I.6.3.4-ko]{6.3.4, (iv)})에 따라 $X_{(S')}$는 $S'$ 위에서
유한형이다.

또한 $S$-준스킴의 곱 $X\times_S Y$에서, 인자 $X$, $Y$ 가운데
\emph{하나가}
\oldpage[I]{147}
$S$ 위에서 \emph{유한형}이고
\emph{다른 하나가 뇌터이면}(각각 \emph{국소 뇌터이면}),
$X\times_S Y$는 \emph{뇌터이다}(각각 \emph{국소 뇌터이다}).

\begin{corollary}[6.3.9]
\label{I.6.3.9-ko}
$X$를 국소 뇌터 준스킴 $S$ 위의 유한형 준스킴이라 하자. 그러면 모든
$S$-사상 $f:X\to Y$는 유한형이다.
\end{corollary}

실제로 $S$가 뇌터라고 해도 된다. $\varphi:X\to S$,
$\psi:Y\to S$가 구조 사상이면 $\varphi=\psi\circ f$이고,
(\hyperref[I.6.3.7-ko]{6.3.7})에 따라 $X$는 뇌터이다. 따라서
(\hyperref[I.6.3.6-ko]{6.3.6})에 따라 $f$는 유한형이다.

\begin{proposition}[6.3.10]
\label{I.6.3.10-ko}
$f:X\to Y$를 유한형 사상이라 하자. $f$가 전사이기 위한 필요충분조건은
모든 대수적으로 닫힌 체 $\Omega$에 대하여 $f$에 대응하는 사상
$X(\Omega)\to Y(\Omega)$가 전사인 것이다
(\hyperref[I.3.4.1-ko]{3.4.1}).
\end{proposition}

이 조건이 충분함은 각 $y\in Y$에 대하여 $k(y)$의 대수적으로 닫힌
확대체 $\Omega$와 다음 가환 그림을 고려하면 알 수 있다.
\[
  \xymatrix{
    X\ar[dd]_f\\
    & \operatorname{Spec}(\Omega)\ar[ul]\ar[dl]\\
    Y
  }
\]
(\emph{참조} (\hyperref[I.3.5.3-ko]{3.5.3})). 역으로 $f$가 전사라고
가정하고, $\Omega$가 대수적으로 닫힌 체인 사상
$g:\{\xi\}=\operatorname{Spec}(\Omega)\to Y$를 택하자. 다음 그림을
고려하면
\[
  \xymatrix{
    X\ar[d]_f &
    X_{(\Omega)}\ar[l]\ar[d]^{f_{(\Omega)}}\\
    Y &
    \operatorname{Spec}(\Omega)\ar[l]
  }
\]
$X_{(\Omega)}$에 \emph{$\Omega$-유리점}이 존재함을 보이면 충분하다
(\hyperref[I.3.3.14-ko]{3.3.14}, \hyperref[I.3.4.3-ko]{3.4.3},
\hyperref[I.3.4.4-ko]{3.4.4}). $f$가 전사이므로 $X_{(\Omega)}$는
공집합이 아니고(\hyperref[I.3.5.10-ko]{3.5.10}), $f$가 유한형이므로
(\hyperref[I.6.3.4-ko]{6.3.4, (iv)})에 따라 $f_{(\Omega)}$도 유한형이다.
따라서 $X_{(\Omega)}$에는 $A(Z)$가 $\Omega$ 위의 영환이 아닌 유한 생성
대수인 공집합이 아닌 아핀 열린집합 $Z$가 존재한다. 힐베르트
영점정리(\cite{I-21})에 따라 $\Omega$-대수 준동형
$A(Z)\to\Omega$가 존재하고, 따라서 $\operatorname{Spec}(\Omega)$ 위의
$X_{(\Omega)}$의 단면이 존재한다. 이로써 명제가 증명된다.

\subsection{대수적 준스킴}
\label{subsection:I.6.4-ko}

\begin{definition}[6.4.1]
\label{I.6.4.1-ko}
체 $K$가 주어졌을 때 $K$ 위의 유한형 준스킴 $X$를
\emph{$K$-대수적 준스킴}이라고 한다. $K$를 $X$의 \emph{기초체}라고 한다.
더 나아가 $X$가 스킴이면(또는 이와 동치로
(\hyperref[I.5.5.8-ko]{5.5.8}), $X$가 $K$-스킴이면) $X$를
\emph{$K$-대수적 스킴}이라고도 한다.
\end{definition}

모든 $K$-대수적 준스킴은 \emph{뇌터이다}
(\hyperref[I.6.3.7-ko]{6.3.7}).

\begin{proposition}[6.4.2]
\label{I.6.4.2-ko}
$X$를 $K$-대수적 준스킴이라 하자. 점 $x\in X$가 닫힌점이기 위한
필요충분조건은 $k(x)$가 $K$의 유한 차수 대수적 확대체인 것이다.
\end{proposition}

$X$가 아핀이고 그 환 $A$가 $K$ 위의 유한 생성 대수인 경우만 다루면
된다. 실제로 $A(U)$가 $K$ 위의 유한 생성 대수인 $X$의 아핀 열린집합
$U$들은 $X$를 덮는다(\hyperref[I.6.3.1-ko]{6.3.1}). 이때 $X$의 닫힌점들은
$\mathfrak{j}_x$가
\oldpage[I]{148}
$A$의 극대 아이디얼인 점들, 다시 말해 $A/\mathfrak{j}_x$가 체인 점들이다.
이 체는 반드시 $k(x)$와 같다. $A/\mathfrak{j}_x$가 $K$ 위의 유한 생성
대수이므로, $x$가 닫힌점이면 $k(x)$는 $K$ 위의 유한 생성 대수인 체이고,
따라서 반드시 $K$ 위에서 \emph{유한 차원}이다 \cite{I-21}. 역으로
$k(x)$가 $K$ 위에서 유한 차원이면 그 부분환
$A/\mathfrak{j}_x\subset k(x)$도 그러하다. $K$ 위의 유한 차원 대수인
정역은 모두 체이므로 $A/\mathfrak{j}_x=k(x)$이고, 따라서 $x$는 닫힌점이다.

\begin{corollary}[6.4.3]
\label{I.6.4.3-ko}
$K$를 대수적으로 닫힌 체, $X$를 $K$-대수적 준스킴이라 하자. 그러면
$X$의 닫힌점들은 $K$-유리점들이고(\hyperref[I.3.4.4-ko]{3.4.4}),
$K$에 값을 갖는 $X$의 점들과 표준적으로 동일시된다.
\end{corollary}

\begin{proposition}[6.4.4]
\label{I.6.4.4-ko}
$X$를 체 $K$ 위의 대수적 준스킴이라 하자. 다음 성질들은 동치이다.
\begin{enumerate}
  \item[\rm(a)] $X$는 아르틴이다.
  \item[\rm(b)] $X$의 바탕공간은 이산공간이다.
  \item[\rm(c)] $X$의 바탕공간에는 닫힌점이 유한개뿐이다.
  \item[\rm(c')] $X$의 바탕공간은 유한집합이다.
  \item[\rm(d)] $X$의 모든 점은 닫힌점이다.
  \item[\rm(e)] $X$는 $K$ 위의 유한 차원 대수 $A$에 대하여
    $\operatorname{Spec}(A)$와 동형이다.
\end{enumerate}
\end{proposition}

$X$가 뇌터이므로 (\hyperref[I.6.2.2-ko]{6.2.2})에 따라 조건 (a), (b),
(d)는 동치이고 (c)와 (c')를 함의한다. 또한 (e)가 (a)를 함의함은
명백하다. (c)가 (d)와 (e)를 함의함을 보이는 일만 남는다. $X$가 아핀인
경우만 다루면 된다. 이때 $A(X)$는 $K$ 위의 유한 생성 대수이고
(\hyperref[I.6.3.3-ko]{6.3.3}), 따라서 야콥슨 환이다
(\cite[p.~3-11 및 3-12]{I-1}). 가정에 따라 이 환에는 극대 아이디얼이
유한개뿐이다. 소 아이디얼들의 유한 교집합이 소 아이디얼이면 그 가운데
하나와 같아야 하므로, $A(X)$의 모든 소 아이디얼은 극대 아이디얼이다.
따라서 (d)가 따른다. 또한 (\hyperref[I.6.2.2-ko]{6.2.2})에 따라
$A(X)$는 유한 생성 아르틴 $K$-대수이고, 따라서 반드시
\emph{유한 차원}이다 \cite{I-21}.

\begin{env}[6.4.5]
\label{I.6.4.5-ko}
(\hyperref[I.6.4.4-ko]{6.4.4})의 조건들이 성립할 때 $X$를
\emph{$K$ 위의 유한 스킴}(\emph{참조}
(\agref{II.6.1.1-ko}{II, 6.1.1})) 또는 \emph{유한 $K$-스킴}이라고 한다.
그 \emph{계수(階數)}는 $[A:K]$이며 $\operatorname{rg}_K(X)$로도 쓴다.
$X$, $Y$가 $K$ 위의 두 유한 스킴이면
\[
  \operatorname{rg}_K(X\amalg Y)
  =\operatorname{rg}_K(X)+\operatorname{rg}_K(Y)
  \tag{6.4.5.1}\label{I.6.4.5.1-ko}
\]
\[
  \operatorname{rg}_K(X\times_K Y)
  =\operatorname{rg}_K(X)\operatorname{rg}_K(Y)
  \tag{6.4.5.2}\label{I.6.4.5.2-ko}
\]
이다. 이는 (\hyperref[I.3.2.2-ko]{3.2.2})에서 따른다.
\end{env}

\begin{corollary}[6.4.6]
\label{I.6.4.6-ko}
$X$를 체 $K$ 위의 유한 스킴이라 하자. $K$의 모든 확대체 $K'$에 대하여
$X\otimes_K K'$는 $K'$ 위의 유한 스킴이고, $K'$ 위의 계수는 $X$의
$K$ 위의 계수와 같다.
\end{corollary}

실제로 $A=A(X)$라 하면 $[A\otimes_K K':K']=[A:K]$이다.

\begin{corollary}[6.4.7]
\label{I.6.4.7-ko}
$X$를 체 $K$ 위의 유한 스킴이라 하고
$n=\sum_{x\in X}[k(x):K]_s$라 놓자.
\emph{(확대체 $K'/K$에 대하여 $[K':K]_s$는 $K'$에 포함되는 $K$의
가장 큰 대수적 분리확대체의 $K$ 위의 \emph{분리차수}임을 상기하자.)}
\oldpage[I]{149}
그러면 $K$의 모든 대수적으로 닫힌 확대체 $\Omega$에 대하여
$X\otimes_K\Omega$의 바탕공간에는 점이 정확히 $n$개 있고, 이 점들은
$X(\Omega)_K$의 원소들과 동일시된다.
\end{corollary}

(\hyperref[I.6.2.2-ko]{6.2.2})에 따라 환 $A=A(X)$가 국소환인 경우만
다루면 된다. $\mathfrak m$을 그 극대 아이디얼, $L=A/\mathfrak m$을
$K$의 대수적 확대체인 잉여체라 하자. 그러면 $X(\Omega)_K$의 원소들은
$X\otimes_K\Omega$의 $\Omega$-단면들과 전단사로 대응하고
(\hyperref[I.3.4.1-ko]{3.4.1}, \hyperref[I.3.3.14-ko]{3.3.14}), 또한
$L$에서 $\Omega$로 가는 $K$-준동형들과 전단사로 대응한다
(\hyperref[I.1.7.3-ko]{1.7.3}). 따라서
(Bourbaki, \emph{Alg.}, 제V장, \textsection7, n\textsuperscript{o}~5,
명제~8)과 (\hyperref[I.6.4.3-ko]{6.4.3})을 고려하면 결론이 따른다.

\begin{env}[6.4.8]
\label{I.6.4.8-ko}
(\hyperref[I.6.4.7-ko]{6.4.7})에서 정의한 수 $n$을 $A$(또는 $X$)의
$K$ 위의 \emph{분리차수}, 또는 \emph{$X$의 점의 기하적 개수}라고 한다.
따라서 이는 $X(\Omega)_K$의 원소 수와 같다. 이 정의에서 곧바로,
$K$의 모든 확대체 $K'$에 대하여 $X\otimes_K K'$의 점의 기하적 개수가
$X$의 점의 기하적 개수와 같음이 따른다. 이 수를 $n(X)$라 하자. $X$, $Y$가
$K$ 위의 두 유한 스킴이면 명백히
\[
  n(X\amalg Y)=n(X)+n(Y).
  \tag{6.4.8.1}\label{I.6.4.8.1-ko}
\]
이다.

같은 가정 아래에서
\[
  n(X\times_K Y)=n(X)n(Y)
  \tag{6.4.8.2}\label{I.6.4.8.2-ko}
\]
도 성립한다. 이는 $n(X)$를 $X(\Omega)_K$의 원소 수로 해석하고 공식
(\hyperref[I.3.4.3.1-ko]{3.4.3.1})을 적용하면 곧바로 따른다.
\end{env}

\begin{proposition}[6.4.9]
\label{I.6.4.9-ko}
$K$를 체, $X$, $Y$를 두 $K$-대수적 준스킴,
$f:X\to Y$를 $K$-사상이라 하자. 또한 $\Omega$를 $K$ 위의 초월차수가
무한인 대수적으로 닫힌 확대체라 하자. $f$가 전사이기 위한 필요충분조건은
$f$에 대응하는 사상 $X(\Omega)_K\to Y(\Omega)_K$가 전사인 것이다
(\hyperref[I.3.4.1-ko]{3.4.1}).
\end{proposition}

필요성은 (\hyperref[I.6.3.10-ko]{6.3.10})에서 따른다. 여기서 $f$가
반드시 유한형인 것은 (\hyperref[I.6.3.9-ko]{6.3.9})에 따른다.
충분성을 보이기 위해 (\hyperref[I.6.3.10-ko]{6.3.10})에서와 같이
논증한다. 실제로 모든 $y\in Y$에 대하여 $k(y)$는 $K$의 유한 생성
확대체이고, 따라서 $\Omega$의 부분체와 $K$-동형이다.

\begin{remark}[6.4.10]
\label{I.6.4.10-ko}
제IV장에서 (\hyperref[I.6.4.9-ko]{6.4.9})의 결론이 $\Omega$의 $K$ 위의
초월차수에 관한 가정 없이도 성립함을 보일 것이다.
\end{remark}

\begin{proposition}[6.4.11]
\label{I.6.4.11-ko}
$f:X\to Y$가 유한형 사상이면, 모든 $y\in Y$에 대하여 섬유
$f^{-1}(y)$는 잉여체 $k(y)$ 위의 대수적 준스킴이고, 모든
$x\in f^{-1}(y)$에 대하여 $k(x)$는 $k(y)$의 유한 생성 확대체이다.
\end{proposition}

$f^{-1}(y)=X\otimes_Y k(y)$이므로(\hyperref[I.3.6.3-ko]{3.6.3}), 이 명제는
(\hyperref[I.6.3.4-ko]{6.3.4, (iv)})와
(\hyperref[I.6.3.3-ko]{6.3.3})에서 따른다.

\begin{proposition}[6.4.12]
\label{I.6.4.12-ko}
$f:X\to Y$, $g:Y'\to Y$를 두 사상이라 하자.
$X'=X\times_Y Y'$라 놓고 $f'=f_{(Y')}:X'\to Y'$라 하자.
$y'\in Y'$, $y=g(y')$라 하자. 섬유 $f^{-1}(y)$가 $k(y)$ 위의 유한
대수적 스킴이면, 섬유 $f'^{-1}(y')$도 $k(y')$ 위의 유한 대수적
스킴이고, 그 계수와 점의 기하적 개수는 $f^{-1}(y)$의 것과 같다.
\end{proposition}

섬유의 추이성(\hyperref[I.3.6.5-ko]{3.6.5})을 고려하면 이는
(\hyperref[I.6.4.6-ko]{6.4.6})과 (\hyperref[I.6.4.8-ko]{6.4.8})에서
곧바로 따른다.

\begin{env}[6.4.13]
\label{I.6.4.13-ko}
명제 (\hyperref[I.6.4.11-ko]{6.4.11})은 유한형 사상이 직관적으로
``대수다양체들의 대수적 족''에 해당함을 보여 준다. 여기서 $Y$의 점들은
\oldpage[I]{150}
``매개변수''의 역할을 하며, 이것이 이 사상들에 ``기하적'' 의미를 준다.
유한형이 아닌 사상들은 뒤에서 주로, 예를 들어 국소화나 완비화에 의한
``밑준스킴의 변경'' 문제에 나타날 것이다.
\end{env}

\subsection{사상의 국소적 결정}
\label{subsection:I.6.5-ko}

\begin{proposition}[6.5.1]
\label{I.6.5.1-ko}
$X$, $Y$를 두 $S$-준스킴이라 하고, $Y$가 $S$ 위의 유한형이라고 하자.
또한 $x\in X$, $y\in Y$가 같은 점 $s\in S$ 위에 있다고 하자.
\begin{enumerate}
  \item[\rm(i)] $X$에서 $Y$로 가는 두 $S$-사상
    $f=(\psi,\theta)$, $f'=(\psi',\theta')$가
    $\psi(x)=\psi'(x)=y$를 만족하고, $\mathscr{O}_y$에서
    $\mathscr{O}_x$로 가는 (국소) $\mathscr{O}_s$-준동형
    $\theta_x^\sharp$와 ${\theta'}_x^\sharp$가 같으면, $f$와 $f'$는
    $x$의 어떤 열린 근방에서 일치한다.
  \item[\rm(ii)] 더 나아가 $S$가 국소 뇌터라고 하자. 모든 국소
    $\mathscr{O}_s$-준동형
    $\varphi:\mathscr{O}_y\to\mathscr{O}_x$에 대하여 $X$에서 $x$의
    열린 근방 $U$와 $U$에서 $Y$로 가는 $S$-사상
    $f=(\psi,\theta)$가 존재하여
    $\psi(x)=y$이고 $\theta_x^\sharp=\varphi$이다.
\end{enumerate}
\end{proposition}

\begin{enumerate}
  \item[\rm(i)] 문제는 $S$, $X$, $Y$ 위에서 국소적이므로 $S$, $X$, $Y$가
    각각 환 $A$, $B$, $C$의 아핀 준스킴이라고 해도 된다. 이때 $f$, $f'$는
    각각 $({}^a\varphi,\widetilde{\varphi})$,
    $({}^a\varphi',\widetilde{\varphi}')$의 꼴이며, $\varphi$, $\varphi'$는
    $C$에서 $B$로 가는 두 $A$-준동형으로서
    $\varphi^{-1}(\mathfrak{j}_x)
    ={\varphi'}^{-1}(\mathfrak{j}_x)=\mathfrak{j}_y$를 만족한다. 또한
    $\varphi$, $\varphi'$에서 얻어지는 $C_y$에서 $B_x$로
    가는 준동형 $\varphi_x$, $\varphi'_x$는 같다. 더 나아가 $C$가 $A$ 위의
    유한 생성 대수라고 해도 된다. $c_i$($1\leq i\leq n$)를 $A$-대수 $C$의
    생성원이라 하고 $b_i=\varphi(c_i)$, $b'_i=\varphi'(c_i)$라 놓자. 가정에
    따라 분수환 $B_x$에서 $b_i/1=b'_i/1$이다($1\leq i\leq n$). 이는
    $1\leq i\leq n$에 대하여 $s_i(b_i-b'_i)=0$이 되는 원소
    $s_i\in B-\mathfrak{j}_x$가 존재한다는 뜻이며, 모든 $s_i$가 같은 원소
    $g\in B-\mathfrak{j}_x$라고 해도 된다. 따라서 분수환 $B_g$에서
    $1\leq i\leq n$에 대하여 $b_i/1=b'_i/1$이다. $i_g:B\to B_g$를 표준
    준동형이라 하면 $i_g\circ\varphi=i_g\circ\varphi'$이고, 그러므로
    $f$와 $f'$의 $D(g)$로의 제한은 같다.
  \item[\rm(ii)] (i)와 같은 상황으로 환원할 수 있고, 더 나아가 환 $A$가
    뇌터라고 해도 된다. $c_i$($1\leq i\leq n$)를 $A$-대수 $C$의 생성원이라
    하고, $\alpha:A[X_1,\ldots,X_n]\to C$를 $X_i$를 $c_i$로 보내는
    다항식 대수 $A[X_1,\ldots,X_n]$에서 $C$로 가는 전사 준동형이라 하자
    ($1\leq i\leq n$). 한편 $i_y:C\to C_y$를 표준 준동형이라 하고, 합성
    준동형
    \[
      \beta:A[X_1,\ldots,X_n]
      \xrightarrow{\alpha} C
      \xrightarrow{i_y} C_y
      \xrightarrow{\varphi} B_x
    \]
    을 생각하자. $\mathfrak{a}$를 $\beta$의 핵이라 하자. $A$가 뇌터이므로
    $A[X_1,\ldots,X_n]$도 뇌터이고, 따라서 $\mathfrak{a}$는 유한 생성계
    $Q_j(X_1,\ldots,X_n)$($1\leq j\leq m$)를 갖는다. 한편 각 원소
    $\varphi(i_y(c_i))$는 $b_i\in B$, $s_i\notin\mathfrak{j}_x$에 대하여
    $b_i/s_i$로 쓸 수 있고, 모든 $s_i$가 같은 원소
    $g\in B-\mathfrak{j}_x$라고 해도 된다. 이제 가정에 따라 $B_x$에서
    $Q_j(b_1/g,\ldots,b_n/g)=0$이다. 다음과 같이 놓자.
    \[
      Q_j(X_1/T,\ldots,X_n/T)
      =P_j(X_1,\ldots,X_n,T)/T^{k_j},
    \]
    여기서 $P_j$는 $k_j$차 동차다항식이다. 또한
    $d_j=P_j(b_1,\ldots,b_n,g)\in B$라 하자. 가정에 따라
    $t_j\in B-\mathfrak{j}_x$에 대하여 $t_jd_j=0$이다
    ($1\leq j\leq m$). 모든 $t_j$가
\oldpage[I]{151}
    같은 원소 $h\in B-\mathfrak{j}_x$라고 해도 된다. 따라서
    $1\leq j\leq m$에 대하여
    $P_j(hb_1,\ldots,hb_n,hg)=0$이다. 이제 $X_i$를 $hb_i/hg$로 보내는
    $A[X_1,\ldots,X_n]$에서 분수환 $B_{hg}$로 가는 준동형 $\rho$를
    생각하자($1\leq i\leq n$). 이 준동형에 의한 $\mathfrak a$의 상은
    0이고, 따라서 $\rho$에 의한 $\alpha^{-1}(0)$의 상도
    \emph{특히} 0이다. 그러므로 $\rho$는
    $A[X_1,\ldots,X_n]\xrightarrow{\alpha}C
    \xrightarrow{\gamma}B_{hg}$로 인수분해되며,
    $\gamma(c_i)=hb_i/hg$이다. $i_x:B_{hg}\to B_x$를 표준
    준동형이라 하면 그림
    \[
      \xymatrix{
        C\ar[r]^\gamma\ar[d]_{i_y} & B_{hg}\ar[d]^{i_x}\\
        C_y\ar[r]_\varphi & B_x
      }
      \tag{6.5.1.1}\label{I.6.5.1.1-ko}
    \]
    은 가환임이 명백하다. 따라서 $\varphi=\gamma_x$이고, $\varphi$가 국소
    준동형이므로 ${}^a\gamma(x)=y$이다. 그러므로
    $f=({}^a\gamma,\widetilde{\gamma})$는 $x$의 근방 $D(hg)$에서 $Y$로
    가는, 문제의 조건을 만족하는 $S$-사상이다.
\end{enumerate}

\begin{corollary}[6.5.2]
\label{I.6.5.2-ko}
(\hyperref[I.6.5.1-ko]{6.5.1, (ii)})의 가정 아래에서, 더 나아가 $X$가
$S$ 위의 유한형이면 사상 $f$가 유한형이라고 할 수 있다.
\end{corollary}

이는 (\hyperref[I.6.3.6-ko]{6.3.6})에서 따른다.

\begin{corollary}[6.5.3]
\label{I.6.5.3-ko}
(\hyperref[I.6.5.1-ko]{6.5.1, (ii)})의 가정들이 성립한다고 하자. 더 나아가
$Y$가 정역 준스킴이고 $\varphi$가 단사 준동형이면,
$f=({}^a\gamma,\widetilde{\gamma})$이고 $\gamma$가 단사 준동형이라고 할 수
있다.
\end{corollary}

실제로 $C$가 정역이라고 해도 되고(\hyperref[I.5.1.4-ko]{5.1.4}), 따라서
$i_y$는 단사 준동형이다. 그러므로 그림
(\hyperref[I.6.5.1.1-ko]{6.5.1.1})에서 $\gamma$가 단사 준동형임이 따른다.

\begin{proposition}[6.5.4]
\label{I.6.5.4-ko}
$f=(\psi,\theta):X\to Y$를 유한형 사상, $x$를 $X$의 점,
$y=\psi(x)$라 하자.
\begin{enumerate}
  \item[\rm(i)] $f$가 점 $x$에서 국소 몰입 사상이기 위한
    (\hyperref[I.4.5.1-ko]{4.5.1}) 필요충분조건은
    $\theta_x^\sharp:\mathscr{O}_y\to\mathscr{O}_x$가 전사 준동형인 것이다.
  \item[\rm(ii)] 더 나아가 $Y$가 국소 뇌터라고 하자. $f$가 점 $x$에서
    국소 동형사상이기 위한(\hyperref[I.4.5.2-ko]{4.5.2}) 필요충분조건은
    $\theta_x^\sharp$가 동형인 것이다.
\end{enumerate}
\end{proposition}

\begin{enumerate}
  \item[\rm(ii)] (\hyperref[I.6.5.1-ko]{6.5.1})에 따라 $y$의 열린 근방
    $V$와 사상 $g:V\to X$가 존재하여 $g\circ f$(각각 $f\circ g$)가
    정의되고 $x$(각각 $y$)의 어떤 근방에서 항등사상과 일치한다. 이로부터
    $f$가 국소 동형사상임이 곧바로 따른다.
  \item[\rm(i)] 문제는 $X$, $Y$ 위에서 국소적이므로 $X$, $Y$가 각각 환
    $A$, $B$의 아핀 준스킴이라고 해도 된다. 이때
    $f=({}^a\varphi,\widetilde{\varphi})$이고, $\varphi:B\to A$는 $A$를
    유한 생성 $B$-대수로 만드는 환 준동형이다. 또한
    $\varphi^{-1}(\mathfrak{j}_x)=\mathfrak{j}_y$이고, $\varphi$에서
    얻어지는 준동형 $\varphi_x:B_y\to A_x$는 전사이다. $(t_i)$
    ($1\leq i\leq n$)를 $B$-대수 $A$의 생성계라 하자. $\varphi_x$에 관한
    가정에 따라 $b_i\in B$와 $c\in B-\mathfrak{j}_y$가 존재하여 분수환
    $A_x$에서 $1\leq i\leq n$에 대하여
    $t_i/1=\varphi(b_i)/\varphi(c)$이다. 따라서
    (\hyperref[I.1.3.3-ko]{1.3.3})에 따라 $a\in A-\mathfrak{j}_x$가
    존재하여, $g=a\varphi(c)$라 놓으면
    $t_i/1=a\varphi(b_i)/g$가 \emph{분수환 $A_g$ 안에서}도 성립한다.
    한편 가정에 따라 $\varphi(B)$를 계수환으로 하는 다항식
    $Q(X_1,\ldots,X_n)$가
\oldpage[I]{152}
    존재하여 $a=Q(t_1,\ldots,t_n)$이다. 다음과 같이 놓자.
    \[
      Q(X_1/T,\ldots,X_n/T)
      =P(X_1,\ldots,X_n,T)/T^m,
    \]
    여기서 $P$는 $m$차 동차다항식이다. 환 $A_g$에서
    \[
      a/1
      =a^mP(\varphi(b_1),\ldots,\varphi(b_n),\varphi(c))/g^m
      =a^m\varphi(d)/g^m
    \]
    이며, 여기서 $d\in B$이다. $A_g$에서
    $g/1=(a/1)(\varphi(c)/1)$은 정의에 따라 가역이므로 $a/1$과
    $\varphi(c)/1$도 가역이고, 따라서
    \[
      a/1=(\varphi(d)/1)(\varphi(c)/1)^{-m}
    \]
    으로 쓸 수 있다. 그러므로 $\varphi(d)/1$도 $A_g$에서 가역이다.
    이제 $h=cd$라 놓자. $\varphi(h)/1$이 $A_g$에서 가역이므로 합성
    준동형
    \[
      B\xrightarrow{\varphi}A\longrightarrow A_g
    \]
    은 다음과 같이 인수분해된다.
    \[
      B\longrightarrow B_h\xrightarrow{\gamma}A_g
    \]
    (\hyperref[0.1.2.4-ko]{0, 1.2.4}). $\gamma$가
    \emph{전사}임을 보이자. $A_g$에서 $B_h$의 상이 $t_i/1$과
    $(g/1)^{-1}$을 포함함을 보이면 충분하다. 그런데
    \[
      (g/1)^{-1}
      =(\varphi(c)/1)^{m-1}(\varphi(d)/1)^{-1}
      =\gamma(c^m/h),
    \]
    이고
    \[
      a/1=\gamma(d^{m+1}/h^m),
    \]
    이므로
    \[
      (a\varphi(b_i))/1=\gamma(b_i d^{m+1}/h^m)
    \]
    이다. 또한
    $t_i/1=(a\varphi(b_i)/1)(g/1)^{-1}$이므로 주장이 증명된다. $h$의
    선택에 따라 $\psi(D(g))\subset D(h)$이고, $f$를 $D(g)$에 제한한
    사상은 $({}^a\gamma,\widetilde{\gamma})$와 같다. $\gamma$가 전사이므로
    이 제한은 $D(g)$에서 $D(h)$로 가는 닫힌 몰입 사상이다
    (\hyperref[I.4.2.3-ko]{4.2.3}).
\end{enumerate}

\begin{corollary}[6.5.5]
\label{I.6.5.5-ko}
$f=(\psi,\theta):X\to Y$를 유한형 사상이라 하자. $X$가 기약이라고 하고,
그 일반점을 $x$라 하며 $y=\psi(x)$라 놓자.
\begin{enumerate}
  \item[\rm(i)] $f$가 $X$의 어떤 점에서 국소 몰입 사상이기 위한
    필요충분조건은
    $\theta_x^\sharp:\mathscr{O}_y\to\mathscr{O}_x$가 전사 준동형인 것이다.
  \item[\rm(ii)] 더 나아가 $Y$가 기약이고 국소 뇌터라고 하자. $f$가
    $X$의 어떤 점에서 국소 동형사상이기 위한 필요충분조건은 $y$가 $Y$의
    일반점인 것, 다시 말해(\hyperref[0.2.1.4-ko]{0, 2.1.4}) $f$가
    \emph{우세 사상}인 것과 $\theta_x^\sharp$가 동형인 것, 다시 말해
    $f$가 \emph{쌍유리 사상}인 것이다
    (\hyperref[I.2.2.9-ko]{2.2.9}).
\end{enumerate}
\end{corollary}

(i)는 $X$의 공집합이 아닌 모든 열린집합이 $x$를 포함한다는 사실을
고려하면 (\hyperref[I.6.5.4-ko]{6.5.4, (i)})에서 명백히 따른다. 마찬가지로
(ii)는 (\hyperref[I.6.5.4-ko]{6.5.4, (ii)})에서 따른다.

\subsection{준콤팩트 사상과 국소 유한형 사상}
\label{subsection:I.6.6-ko}

\begin{definition}[6.6.1]
\label{I.6.6.1-ko}
사상 $f:X\to Y$에 의한 $Y$의 모든 준콤팩트 열린집합의 역상이
준콤팩트이면, $f$를 \emph{준콤팩트} 사상이라 한다.
\end{definition}

$\mathfrak{B}$를 준콤팩트 열린집합들(예를 들어 아핀 열린집합들)로
이루어진 $Y$의 위상의 기저라 하자. $f$가 준콤팩트이기 위한 필요충분조건은
$\mathfrak{B}$의 모든 원소의 $f$에 의한 역상이 준콤팩트인 것, 또는 같은
말로 아핀 열린집합들의 \emph{유한} 합집합인 것이다. 실제로 $Y$의 모든
준콤팩트 열린집합은 $\mathfrak{B}$의 유한개의 원소의 합집합이다. 예를 들어
$X$가 \emph{준콤팩트}이고 $Y$가 \emph{아핀}이면 \emph{모든} 사상
$f:X\to Y$는 준콤팩트이다. 실제로 $X$는 유한개의 아핀 열린집합 $U_i$의
합집합이고, $Y$의 모든 아핀 열린집합 $V$에 대하여
$U_i\cap f^{-1}(V)$는 아핀이므로
(\hyperref[I.5.5.10-ko]{5.5.10}) 준콤팩트이다.

$f:X\to Y$가 준콤팩트 사상이면, $Y$의 모든 열린집합 $V$에 대하여
$f$를 $f^{-1}(V)$에 제한한 사상 $f^{-1}(V)\to V$가 준콤팩트임은
명백하다. 역으로 $(U_\alpha)$가 $Y$의 열린 덮개이고, 사상 $f:X\to Y$의
제한들 $f^{-1}(U_\alpha)\to U_\alpha$가 준콤팩트이면, $f$도
준콤팩트이다.

\begin{definition}[6.6.2]
\label{I.6.6.2-ko}
사상 $f:X\to Y$가 다음 조건을 만족하면 \emph{국소 유한형} 사상이라 한다.
모든 $x\in X$에 대하여 $x$의 열린 근방 $U$와
$f(U)$를 포함하는 $y=f(x)$의 열린 근방 $V$가 존재하고, $f$의
\oldpage[I]{153}
$U$로의 제한이 $U$에서 $V$로 가는 유한형 사상이다. 이때 $X$를 $Y$
위의 국소 유한형 준스킴 또는 국소 유한형 $Y$-준스킴이라고도 한다.
\end{definition}

(\hyperref[I.6.3.2-ko]{6.3.2})에 따라 $f$가 국소 유한형이면, $Y$의 모든
열린집합 $W$에 대하여 $f$를 $f^{-1}(W)$에 제한한 사상
$f^{-1}(W)\to W$도 국소 유한형임이 곧바로 따른다.

$Y$가 국소 뇌터이고 $X$가 $Y$ 위의 국소 유한형이면,
(\hyperref[I.6.3.7-ko]{6.3.7})에 따라 $X$는 국소 뇌터이다.

\begin{proposition}[6.6.3]
\label{I.6.6.3-ko}
사상 $f:X\to Y$가 유한형이기 위한 필요충분조건은 $f$가 준콤팩트이고
국소 유한형인 것이다.
\end{proposition}

조건의 필요성은 (\hyperref[I.6.3.1-ko]{6.3.1})과
(\hyperref[I.6.6.1-ko]{6.6.1}) 뒤의 주에서 곧바로 따른다. 역으로 이
조건들이 성립한다고 하고, $U$를 $Y$의 아핀 열린집합, $A$를 그 환이라
하자. 모든
$x\in f^{-1}(U)$에 대하여 가정에 따라 $x$의 근방
$V(x)\subset f^{-1}(U)$와 $y=f(x)$의 근방 $W(x)\subset U$가 존재하여,
$W(x)$는 $f(V(x))$를 포함하고 $f$를 $V(x)$에 제한한 사상
$V(x)\to W(x)$는 유한형이다. $W(x)$를 $y=f(x)$의
$D(g)$ 꼴인 근방 $W_1(x)\subset W(x)$로($g\in A$), $V(x)$를
$V(x)\cap f^{-1}(W_1(x))$로 바꾸면, $W(x)$가 $D(g)$ 꼴이라고 해도 된다.
따라서 그 환이 $A[1/g]$로 쓰이므로 $W(x)$는 $U$ 위의 유한형이고, 그러므로
$V(x)$도 $U$ 위의 유한형이다. 더 나아가 $f^{-1}(U)$는 가정에 따라
준콤팩트이므로 유한개의 열린집합 $V(x_i)$의 합집합이다. 이로써 증명이
끝난다.

\begin{proposition}[6.6.4]
\label{I.6.6.4-ko}
\begin{enumerate}
  \item[\rm(i)] 몰입 사상 $X\to Y$가 닫힌 몰입 사상이거나, $Y$의
    바탕공간이 국소 뇌터이거나, $X$의 바탕공간이 뇌터이면 이 몰입 사상은
    준콤팩트이다.
  \item[\rm(ii)] 두 준콤팩트 사상의 합성은 준콤팩트이다.
  \item[\rm(iii)] $f:X\to Y$가 준콤팩트 $S$-사상이면, 밑준스킴의 모든
    확장 $g:S'\to S$에 대하여
    $f_{(S')}:X_{(S')}\to Y_{(S')}$도 준콤팩트이다.
  \item[\rm(iv)] $f:X\to X'$와 $g:Y\to Y'$가 두 준콤팩트 $S$-사상이면
    $f\times_S g$도 준콤팩트이다.
  \item[\rm(v)] 두 사상 $f:X\to Y$, $g:Y\to Z$의 합성이 준콤팩트이고,
    또한 $g$가 분리 사상이거나 $X$의 바탕공간이 국소 뇌터이면 $f$는
    준콤팩트이다.
  \item[\rm(vi)] 사상 $f$가 준콤팩트이기 위한 필요충분조건은
    $f_{\mathrm{red}}$가 준콤팩트인 것이다.
\end{enumerate}
\end{proposition}

(vi)는 사상의 준콤팩트성이 바탕공간 사이의 대응하는 연속사상에만
의존하므로 명백하다. 마찬가지로 $X$의 바탕공간이 국소 뇌터라고 가정한
경우에 해당하는 (v)를 증명하자. $h=g\circ f$라 놓고, $U$를 $Y$의
준콤팩트 열린집합이라 하자. $g(U)$는 $Z$에서 준콤팩트이지만 반드시
열린집합인 것은 아니므로, 유한개의 준콤팩트 열린집합 $V_j$의 합집합에
포함된다(\hyperref[I.2.1.3-ko]{2.1.3}). 따라서 $f^{-1}(U)$는
$h^{-1}(V_j)$들의 합집합에 포함된다. 이들은 $X$의 준콤팩트
부분공간들이므로 뇌터 부분공간들이다. 그러므로
(\hyperref[0.2.2.3-ko]{0, 2.2.3})에 따라 $f^{-1}(U)$는 뇌터 공간이고,
\emph{따라서 특히} 준콤팩트이다.

나머지 주장들을 증명하기 위해서는 (i), (ii), (iii)를 증명하면 충분하다
(\hyperref[I.5.5.12-ko]{5.5.12}). 그런데 (ii)는 명백하고, (i)는 $Y$의
바탕공간이 국소 뇌터이거나 $X$의 바탕공간이 뇌터인 경우에는
(\hyperref[I.6.3.5-ko]{6.3.5})에서 따르며 닫힌 몰입 사상의 경우에는
명백하다. (iii)를 증명하기 위하여,
\oldpage[I]{154}
$S=Y$인 경우로 환원할 수 있다(\hyperref[I.3.3.11-ko]{3.3.11}).
$f'=f_{(S')}$라 놓고, $U'$를 $S'$의 준콤팩트 열린집합이라 하자. 모든
$s'\in U'$에 대하여 $T$를 $S$에서 $g(s')$의 아핀 열린 근방이라 하고,
$W$를 $U'\cap g^{-1}(T)$에 포함되는 $s'$의 아핀 열린 근방이라 하자.
$f'^{-1}(W)$가 준콤팩트임을 보이면 충분하다. 다시 말해 $S$와 $S'$가
아핀인 경우에 $X\times_S S'$의 바탕공간이 준콤팩트임을 보이는 문제로
환원된다. 그런데 이때 $X$는 가정에 따라 유한개의 아핀 열린집합 $V_j$의
합집합이고, $X\times_S S'$의 바탕공간은 아핀 스킴
$V_j\times_S S'$들의 바탕공간의 합집합이므로
(\hyperref[I.3.2.2-ko]{3.2.2}와 \hyperref[I.3.2.7-ko]{3.2.7})에 따라 명제의
증명이 끝난다.

또한 $X=X'\amalg X''$가 두 준스킴의 합이면, 사상 $f:X\to Y$가
준콤팩트이기 위한 필요충분조건은 $f$를 $X'$와 $X''$에 제한한 사상들이
준콤팩트인 것임을 주목하자.

\begin{proposition}[6.6.5]
\label{I.6.6.5-ko}
$f:X\to Y$를 준콤팩트 사상이라 하자. $f$가 우세 사상이기 위한
필요충분조건은 $Y$의 모든 기약 성분의 일반점 $y$에 대하여
$f^{-1}(y)$가 $X$의 어떤 기약 성분의 일반점을 포함하는 것이다.
\end{proposition}

이 조건이 충분함은($f$가 준콤팩트 사상이라고 가정하지 않아도) 명백하다.
필요성을 보이기 위하여 $y$의 아핀 열린 근방 $U$를 생각하자.
$f^{-1}(U)$는 준콤팩트이므로 아핀 열린집합 $V_i$들의 \emph{유한}
합집합이고, $f$가 우세 사상이라는 가정에 따라 $y$는 $U$에서 어떤
$f(V_i)$의 폐포에 속한다. $X$와 $Y$가 축소라고 해도 됨은 명백하다.
$V_i$의 기약 성분의 $X$에서의 폐포는 $X$의 기약 성분이므로
(\hyperref[0.2.1.6-ko]{0, 2.1.6}), $X$를 $V_i$로, $Y$를
$\overline{f(V_i)}\cap U$를 바탕공간으로 갖는 $U$의 축소 닫힌
부분준스킴으로 바꿀 수 있다
(\hyperref[I.5.2.1-ko]{5.2.1}). 따라서 $X=\operatorname{Spec}(A)$와
$Y=\operatorname{Spec}(B)$가 아핀이고 축소인 경우의 명제를 증명하면 된다.
$f$가 우세 사상이므로 $B$는 $A$의 부분환이고
(\hyperref[I.1.2.7-ko]{1.2.7}), 이제 명제는 $B$의 모든 극소 소 아이디얼이
$B$와 $A$의 어떤 극소 소 아이디얼의 교집합이라는 사실에서 따른다
(\hyperref[0.1.5.8-ko]{0, 1.5.8}).

\begin{proposition}[6.6.6]
\label{I.6.6.6-ko}
\begin{enumerate}
  \item[\rm(i)] 모든 국소 몰입 사상은 국소 유한형이다.
  \item[\rm(ii)] 두 사상 $f:X\to Y$, $g:Y\to Z$가 국소 유한형이면
    $g\circ f$도 국소 유한형이다.
  \item[\rm(iii)] $f:X\to Y$가 국소 유한형 $S$-사상이면, 밑준스킴의
    모든 확장 $S'\to S$에 대하여
    $f_{(S')}:X_{(S')}\to Y_{(S')}$도 국소 유한형이다.
  \item[\rm(iv)] $f:X\to X'$와 $g:Y\to Y'$가 두 국소 유한형
    $S$-사상이면 $f\times_S g$도 국소 유한형이다.
  \item[\rm(v)] 두 사상의 합성 $g\circ f$가 국소 유한형이면 $f$도
    국소 유한형이다.
  \item[\rm(vi)] 사상 $f$가 국소 유한형이면 $f_{\mathrm{red}}$도
    국소 유한형이다.
\end{enumerate}
\end{proposition}

(\hyperref[I.5.5.12-ko]{5.5.12})에 따라 (i), (ii), (iii)를 증명하면
충분하다. $j:X\to Y$가 국소 몰입 사상이면, 모든 $x\in X$에 대하여
$Y$에서 $j(x)$의 열린 근방 $V$와 $X$에서 $x$의 열린 근방 $U$가
존재하여 $j$를 $U$에 제한한 사상은 닫힌 몰입 사상 $U\to V$이다
(\hyperref[I.4.5.1-ko]{4.5.1}). 따라서 이 제한은 유한형이다. (ii)를
증명하기 위하여 점 $x\in X$를 생각하자. 가정에 따라 $g(f(x))$의 열린
근방 $W$와 $f(x)$의 열린 근방 $V$가 존재하여 $g(V)\subset W$이고
$V$는 $W$ 위의 유한형이다. 더 나아가 $f^{-1}(V)$는 $V$ 위의 국소
유한형이므로(\hyperref[I.6.6.2-ko]{6.6.2}), $x$의 열린 근방 $U$가
\oldpage[I]{155}
$f^{-1}(V)$에 포함되고 $V$ 위의 유한형이 되게 할 수 있다. 따라서
$g(f(U))\subset W$이고 $U$는 $W$ 위의 유한형이다
(\hyperref[I.6.3.4-ko]{6.3.4, (ii)}). 끝으로 (iii)를 증명하기 위해서는
$Y=S$인 경우로 환원할 수 있다(\hyperref[I.3.3.11-ko]{3.3.11}). 모든
$x'\in X'=X_{(S')}$에 대하여 $x$를 $X$에서 $x'$의 상, $s$를 $S$에서
$x$의 상, $T$를 $s$의 열린 근방, $T'$를 $S'$에서 $T$의 역상, $U$를
$x$의 열린 근방으로서 그 상이 $T$에 포함되고 $T$ 위의 유한형인 것으로
잡자. 그러면 $U\times_S T'=U\times_T T'$는 $x'$의 열린 근방이고
(\hyperref[I.3.2.7-ko]{3.2.7}), $T'$ 위의 유한형이다
(\hyperref[I.6.3.4-ko]{6.3.4, (iv)}).

\begin{corollary}[6.6.7]
\label{I.6.6.7-ko}
$X$, $Y$를 $S$ 위의 국소 유한형인 두 $S$-준스킴이라 하자. $S$가 국소
뇌터이면 $X\times_S Y$는 국소 뇌터이다.
\end{corollary}

실제로 $X$는 $S$ 위의 국소 유한형이므로 국소 뇌터이고,
$X\times_S Y$는 $X$ 위의 국소 유한형이므로 역시 국소 뇌터이다.

\begin{remark}[6.6.8]
\label{I.6.6.8-ko}
명제 (\hyperref[I.6.3.10-ko]{6.3.10})과 그 증명은 사상 $f$가 국소
유한형이라고만 가정하는 경우로 곧바로 확장된다. 마찬가지로 명제
(\hyperref[I.6.4.2-ko]{6.4.2})와 (\hyperref[I.6.4.9-ko]{6.4.9})는 그
명제들에 나타나는 준스킴 $X$, $Y$가 체 $K$ 위의 국소 유한형이라고만
가정할 때에도 성립한다.
\end{remark}

\section{유리 사상}
\label{section:I.7-ko}

\subsection{유리 사상과 유리함수}
\label{subsection:I.7.1-ko}

\begin{env}[7.1.1]
\label{I.7.1.1-ko}
$X$, $Y$를 두 준스킴, $U$, $V$를 $X$의 두 조밀한 열린집합,
$f$(각각 $g$)를 $U$(각각 $V$)에서 $Y$로 가는 사상이라 하자.
$f$와 $g$가 $U\cap V$ 안의 어떤 조밀한 열린집합에서 일치하면
이들을 \emph{동치}라고 하겠다. $X$의 조밀한 열린집합들의 유한
교집합은 $X$의 조밀한 열린집합이므로 이 관계가
\emph{동치관계}임은 자명하다.
\end{env}

\begin{definition}[7.1.2]
\label{I.7.1.2-ko}
두 준스킴 $X$, $Y$가 주어졌다고 하자.
(\hyperref[I.7.1.1-ko]{7.1.1})에서 정의한 관계에 따른 $X$의 조밀한
열린 부분집합에서 $Y$로 가는 사상들의 동치류를
\emph{$X$에서 $Y$로 가는 유리 사상}이라 한다. $X$와 $Y$가
$S$-준스킴일 때, 이 동치류의 어떤 대표가 $S$-사상이면 이를
\emph{$S$-유리 사상}이라 한다. $S$에서 $X$로 가는 모든
$S$-유리 사상을 \emph{$X$의 $S$-유리 단면}이라 한다. $X$-준스킴
$X\otimes_{\mathbf{Z}}\mathbf{Z}[T]$($T$는 부정원)의 모든
$X$-유리 단면을 \emph{준스킴 $X$ 위의 유리함수}라 한다.

$S$-준스킴만을 다룰 때에는 혼동의 염려가 없으면 용어를 남용하여
``$S$-유리 사상'' 대신 ``유리 사상''이라고 하겠다.

$f$를 $X$에서 $Y$로 가는 유리 사상, $U$를 $X$의 열린집합이라 하자.
$f$의 동치류에 속하는 사상 $f_1$, $f_2$가 각각 $X$의 조밀한
열린집합 $V$, $W$ 위에서 정의되어 있으면, 그 제한
$f_1|(U\cap V)$와 $f_2|(U\cap W)$는 $U\cap V\cap W$ 안의,
$U$에서 조밀한 어떤 열린집합에서 일치하므로 동치이다. 따라서
동치류 $f$는
$U$에서 $Y$로 가는 유리 사상을 정의한다. 이를
\emph{$f$의 $U$로의 제한}이라 하고 $f|U$로 쓴다.

각 $S$-사상 $f:X\to Y$에 $S$-유리 사상 하나를 대응시킬 때 그 사상을
\oldpage[I]{156}
$f$가 속하는 동치류로 정하면,
$\operatorname{Hom}_S(X,Y)$에서 $X$에서 $Y$로 가는 $S$-유리 사상들의
집합으로 가는 표준 사상이 정의된다. $X$의 $Y$-유리 단면들의 집합을
$\Gamma_{\mathrm{rat}}(X/Y)$로 나타내면 표준 사상
$\Gamma(X/Y)\to\Gamma_{\mathrm{rat}}(X/Y)$를 얻는다. 또한 $X$와 $Y$가
두 $S$-준스킴이면, $X$에서 $Y$로 가는 $S$-유리 사상들의 집합은
$\Gamma_{\mathrm{rat}}((X\times_S Y)/X)$와 표준적으로 동일시됨이
자명하다
(\hyperref[I.3.3.14-ko]{3.3.14}).
\end{definition}

\begin{env}[7.1.3]
\label{I.7.1.3-ko}
(\hyperref[I.7.1.2-ko]{7.1.2})와
(\hyperref[I.3.3.14-ko]{3.3.14})에서 곧바로 알 수 있듯이,
$X$ 위의 유리함수들은 $X$의 모든 곳에서 조밀한 열린집합 위의
\emph{구조층 $\mathscr{O}_X$의 단면들의 동치류}와 표준적으로
동일시된다. 여기서 두 단면은 그 정의역들의 교집합에 포함된 모든
곳에서 조밀한 열린집합 위에서 일치할 때 동치이다. 특히 $X$ 위의
유리함수들은 \emph{환} $R(X)$을 이룬다.
\end{env}

\begin{env}[7.1.4]
\label{I.7.1.4-ko}
$X$가 \emph{기약} 준스킴이면 공집합이 아닌 모든 열린집합은 $X$에서
조밀하다. 달리 말하면 $X$의 공집합이 아닌 열린집합들은 $X$의
\emph{일반점 $x$의 열린 근방}들이다. 그러므로 이 경우 $X$의 공집합이
아닌 열린 부분집합들에서 $Y$로 가는 두 사상이 동치라는 것은 두
사상이 $x$에서 \emph{같은 싹}을 갖는다는 뜻이다. 다시 말해,
$X\to Y$인 유리 사상(각각 $S$-유리 사상)들은 $X$의 일반점
$x$에서, $X$의 공집합이 아닌 열린 부분집합들에서 $Y$로 가는
\emph{사상의 싹}(각각 $S$-사상의 싹)들과 동일시된다. 특히 다음을
얻는다.
\end{env}

\begin{proposition}[7.1.5]
\label{I.7.1.5-ko}
$X$가 기약 준스킴이면 $X$ 위의 유리함수환 $R(X)$은 $X$의 일반점
$x$의 국소환 $\mathscr{O}_x$와 표준적으로 동일시된다. 이는 차원
0의 국소환이며, 따라서 $X$가 뇌터 준스킴이면 아르틴 국소환이다.
$X$가 정역 준스킴이면 이는 체이고, 나아가 $X$가 아핀 스킴이면 $A(X)$의
분수체와 동일시된다.
\end{proposition}

앞의 내용과 유리함수를 모든 곳에서 조밀한 열린집합 위의
$\mathscr{O}_X$의 단면과 동일시한 것을 생각하면, 첫째 주장은 한
점에서 층의 줄기를 정의한 것에 지나지 않는다. 나머지 주장들은
$X$가 환 $A$의 아핀 스킴인 경우만 생각하면 충분하다. 이때
$\mathfrak{j}_x$는 $A$의 멱영근기이므로 $\mathscr{O}_x$는 차원
0이다. $A$가 정역이면 $\mathfrak{j}_x=(0)$이므로 $\mathscr{O}_x$는
$A$의 분수체이다. 끝으로 $A$가 뇌터이면 $\mathfrak{j}_x$가
멱영이고(\cite{I-11}, p.~127, cor.~4),
$\mathscr{O}_x=A_{\mathfrak{j}_x}$가 아르틴임이 알려져 있다.

$X$가 \emph{정역 준스킴}이면 모든 $z\in X$에 대하여 환 $\mathscr{O}_z$는
정역이다. $z$를 포함하는 모든 아핀 열린집합 $U$는 $x$도 포함하며,
$A(U)$의 분수체인 $R(U)$는 따라서 $R(X)$와 동일시된다.
그러므로 $R(X)$은 \emph{$\mathscr{O}_z$의 분수체}와도 동일시된다.
$\mathscr{O}_z$를 $R(X)$의 부분환과 표준적으로 동일시하는 사상은
단면의 모든 싹 $s\in\mathscr{O}_z$에, $z$에서 싹 $s$를 갖는
$\mathscr{O}_X$의 단면(반드시 모든 곳에서 조밀한 열린집합 위에서
정의된다)의 동치류인 $X$ 위의 유일한 유리함수를 대응시킨다.

\begin{env}[7.1.6]
\label{I.7.1.6-ko}
이제 $X$가 \emph{유한} 개의 기약 성분 $X_i$($1\leq i\leq n$)를 갖는다고
하자($X$의 바탕공간이 \emph{뇌터}이면 그러하다). $X'_i$를
$X_i$에서 $j\neq i$인 $X_j\cap X_i$들의 합집합을 뺀 $X$의
열린집합이라 하자. $X'_i$는 기약이고, 그 일반점 $x_i$는 $X_i$의
일반점이며, $X'_i$들은 서로소이고 그 합집합은 $X$에서 조밀하다
(\hyperref[0.2.1.6-ko]{0, 2.1.6}).
\oldpage[I]{157}
$X$의 모든 곳에서 조밀한 임의의 열린집합 $U$에 대하여
$U_i=U\cap X'_i$는 $X'_i$의 공집합이
아닌 조밀한 열린집합이고, $U_i$들은 서로소이므로
$U'=\bigcup_i U_i$는 $X$에서 조밀하다. $U'$에서 $Y$로 가는 사상을
주는 것은 각 $U_i$에서 $Y$로 가는 사상을 임의로 하나씩 주는 것과
같다. 따라서 다음을 얻는다.
\end{env}

\begin{proposition}[7.1.7]
\label{I.7.1.7-ko}
$X$, $Y$를 두 준스킴(각각 $S$-준스킴)이라 하고, $X$가 일반점
$x_i$를 갖는 유한 개의 기약 성분 $X_i$($1\leq i\leq n$)를 갖는다고
하자. $R_i$를 $x_i$에서, $X$의 열린 부분집합에서 $Y$로 가는 사상
(각각 $S$-사상)의 싹들의 집합이라 하면, $X$에서 $Y$로 가는
유리 사상(각각 $S$-유리 사상)들의 집합은 $R_i$들의 곱
($1\leq i\leq n$)과 동일시된다.
\end{proposition}

\begin{corollary}[7.1.8]
\label{I.7.1.8-ko}
$X$를 뇌터 준스킴이라 하자. $X$ 위의 유리함수환은 아르틴 환이고,
그 국소환 성분들은 $X$의 기약 성분들의 일반점 $x_i$의 국소환
$\mathscr{O}_{x_i}$이다.
\end{corollary}

\begin{corollary}[7.1.9]
\label{I.7.1.9-ko}
$A$를 뇌터 환, $X=\operatorname{Spec}(A)$라 하자. $Q$가 $A$의
극소 소 아이디얼들의 합집합의 여집합이면 $X$ 위의 유리함수환은
분수환 $Q^{-1}A$와 표준적으로 동일시된다.
\end{corollary}

이는 다음 보조정리에서 따른다.

\begin{lemma}[7.1.9.1]
\label{I.7.1.9.1-ko}
$f\in A$에 대하여 $D(f)$가 $X$에서 모든 곳에서 조밀하기 위한
필요충분조건은 $f\in Q$인 것이다. $X$의 모든 조밀한 열린집합은
$f\in Q$인 $D(f)$ 꼴의 열린집합을 하나 포함한다.
\end{lemma}

실제로 이 보조정리가 증명되었다고 하자. 단면환
$\Gamma(D(f),\mathscr{O}_X)$는 $A_f$와 동일시된다
(\hyperref[I.1.3.6-ko]{1.3.6}과 \hyperref[I.1.3.7-ko]{1.3.7}).
$f\in Q$인 $D(f)$들은 $\supset$에 관해 순서가 주어진 $X$의 조밀한
열린집합들의 집합에서 공종집합을 이루므로, 정의
(\hyperref[I.7.1.1-ko]{7.1.1})에 따라 $X$ 위의 유리함수환은
$f\in Q$인 $A_f$들의 귀납극한과 동일시된다. 여기서 준순서관계는
``$g$가 $f$의 배수이다''이고, 이 귀납극한은 곧 $Q^{-1}A$이다
(\hyperref[0.1.4.5-ko]{0, 1.4.5}).

(\hyperref[I.7.1.9.1-ko]{7.1.9.1})을 증명하기 위하여 $X$의 기약
성분들도 $X_i$($1\leq i\leq n$)로 나타내자. $D(f)$가 $X$에서
조밀하면 모든 $1\leq i\leq n$에 대하여
$D(f)\cap X_i\neq\varnothing$이고 그 역도 성립한다. 그런데
$\mathfrak{p}_i=\mathfrak{j}(X_i)$라 놓으면 이는 모든
$1\leq i\leq n$에 대하여 $f\notin\mathfrak{p}_i$라는 뜻이다.
$\mathfrak{p}_i$들은 $A$의 극소 소 아이디얼들이므로
(\hyperref[I.1.1.14-ko]{1.1.14}), 관계들
$f\notin\mathfrak{p}_i$($1\leq i\leq n$)는 $f\in Q$와 동치이다.
이로써 보조정리의 첫째 주장이 증명되었다. 한편 $U$가 $X$의 조밀한
열린집합이면 $U$의 여집합은 $V(\mathfrak{a})$ 꼴이고, 여기서
$\mathfrak{a}$는 어느 $\mathfrak{p}_i$에도 포함되지 않는
아이디얼이다. 따라서 이 아이디얼은 그 합집합에도 포함되지 않고
(\cite{I-10}, p.~13), 어떤 $f\in\mathfrak{a}$는 $Q$에도 속한다.
그러므로 $D(f)\subset U$이고 증명이 끝난다.

\begin{env}[7.1.10]
\label{I.7.1.10-ko}
다시 $X$가 일반점 $x$를 갖는 기약 준스킴이라고 하자. $X$의
공집합이 아닌 모든 열린집합 $U$는 $x$를 포함하므로
$x\in\overline{\{z\}}$인 모든 $z\in X$도 포함한다. 따라서 모든
사상 $U\to Y$는 표준 사상 $\operatorname{Spec}(\mathscr{O}_x)\to X$
(\hyperref[I.2.4.1-ko]{2.4.1})와 합성할 수 있다. 또한 $X$의 공집합이
아닌 두 열린 부분집합에서 $Y$로 가는 두 사상이 $X$의 공집합이 아닌
어떤 열린 부분집합에서 일치하면, 이 두 사상을 각각 그 표준 사상과
합성하여 얻는 두 사상
$\operatorname{Spec}(\mathscr{O}_x)\to Y$는 같다. 다시 말해,
$X$에서 $Y$로 가는 모든 유리 사상에는 잘 정해진 사상
$\operatorname{Spec}(\mathscr{O}_x)\to Y$가 이렇게 대응한다.
\end{env}

\oldpage[I]{158}
\begin{proposition}[7.1.11]
\label{I.7.1.11-ko}
$X$, $Y$를 두 $S$-준스킴이라 하자. $X$는 일반점 $x$를 갖는
기약 준스킴이고 $Y$는 $S$ 위에서 유한형이라고 하자. 같은
$S$-사상 $\operatorname{Spec}(\mathscr{O}_x)\to Y$가 대응하는,
$X$에서 $Y$로 가는 두 $S$-유리 사상은 같다. 나아가 $S$가 국소
뇌터이면 $\operatorname{Spec}(\mathscr{O}_x)$에서 $Y$로 가는 모든
$S$-사상에는 $X$에서 $Y$로 가는 유일한 $S$-유리 사상이 대응한다.
\end{proposition}

$X$의 공집합이 아닌 모든 열린집합이 모든 곳에서 조밀함을 고려하면
이는 (\hyperref[I.6.5.1-ko]{6.5.1})에서 곧바로 따른다.

\begin{corollary}[7.1.12]
\label{I.7.1.12-ko}
$S$가 국소 뇌터이고 (\hyperref[I.7.1.11-ko]{7.1.11})의 나머지
가정들이 성립한다고 하자. 그러면 $X$에서 $Y$로 가는 $S$-유리
사상들은 $S$-준스킴 $\operatorname{Spec}(\mathscr{O}_x)$에 값을 갖는
$S$-준스킴 $Y$의 점들과 동일시된다.
\end{corollary}

이는 (\hyperref[I.7.1.11-ko]{7.1.11})을
(\hyperref[I.3.4.1-ko]{3.4.1})에서 도입한 용어로 다시 쓴 것에
지나지 않는다.

\begin{corollary}[7.1.13]
\label{I.7.1.13-ko}
(\hyperref[I.7.1.12-ko]{7.1.12})의 조건들이 성립한다고 하자.
$s$를 $S$에서 $x$의 상이라 하자. $X$에서 $Y$로 가는 $S$-유리
사상을 주는 것은 $s$ 위에 있는 $Y$의 점 $y$와 국소
$\mathscr{O}_s$-준동형
$\mathscr{O}_y\to\mathscr{O}_x=R(X)$을 주는 것과 같다.
\end{corollary}

이는 (\hyperref[I.7.1.11-ko]{7.1.11})과
(\hyperref[I.2.4.4-ko]{2.4.4})에서 따른다.

특히 다음을 얻는다.

\begin{corollary}[7.1.14]
\label{I.7.1.14-ko}
(\hyperref[I.7.1.12-ko]{7.1.12})의 조건 아래에서, $X$에서 $Y$로 가는
$S$-유리 사상들은 ($Y$가 주어졌을 때) $S$-준스킴
$\operatorname{Spec}(\mathscr{O}_x)$에만 의존한다. 특히 모든
$z\in X$에 대하여 $X$를 $\operatorname{Spec}(\mathscr{O}_z)$로
바꾸어도 이 유리 사상들은 변하지 않는다.
\end{corollary}

실제로 $z\in\overline{\{x\}}$이므로 $x$는
$Z=\operatorname{Spec}(\mathscr{O}_z)$의 일반점이고
$\mathscr{O}_{X,x}=\mathscr{O}_{Z,x}$이다.

$X$가 정역 준스킴이면 $R(X)=\mathscr{O}_x=k(x)$는 체이다
(\hyperref[I.7.1.5-ko]{7.1.5}). 따라서 앞의 따름정리들은 다음과 같이
특수화된다.

\begin{corollary}[7.1.15]
\label{I.7.1.15-ko}
(\hyperref[I.7.1.12-ko]{7.1.12})의 조건들이 성립하고, 나아가 $X$가
정역 준스킴이라고 하자. $s$를 $S$에서 $x$의 상이라 하자. 그러면 $X$에서
$Y$로 가는 $S$-유리 사상들은 $k(s)$의 확대체 $R(X)$에 값을 갖는
$Y\otimes_S k(s)$의 기하적 점들과 동일시된다. 다시 말해, 이러한
유리 사상 하나를 주는 것은 $s$ 위에 있는 점 $y\in Y$와
$k(y)$에서 $k(x)=R(X)$로 가는 $k(s)$-단사 준동형을 주는 것과 같다.
\end{corollary}

실제로 $s$ 위에 있는 $Y$의 점들은 $Y\otimes_S k(s)$의 점들과
동일시되고(\hyperref[I.3.6.3-ko]{3.6.3}), 국소
$\mathscr{O}_s$-준동형 $\mathscr{O}_y\to R(X)$들은
$k(s)$-단사 준동형 $k(y)\to R(X)$들과 동일시된다.

더욱 특별히 다음을 얻는다.

\begin{corollary}[7.1.16]
\label{I.7.1.16-ko}
$k$를 체, $X$, $Y$를 $k$ 위의 두 대수적 준스킴
(\hyperref[I.6.4.1-ko]{6.4.1})이라 하고, 나아가 $X$가 정역
준스킴이라고 하자. 그러면 $X$에서 $Y$로 가는 $k$-유리 사상들은 $k$의 확대체
$R(X)$에 값을 갖는 $Y$의 기하적 점들과 동일시된다
(\hyperref[I.3.4.4-ko]{3.4.4}).
\end{corollary}

\subsection{유리 사상의 정의역}
\label{subsection:I.7.2-ko}

\begin{env}[7.2.1]
\label{I.7.2.1-ko}
$X$, $Y$를 두 준스킴, $f$를 $X$에서 $Y$로 가는 유리 사상이라 하자.
$x$를 포함하는 모든 곳에서 조밀한 열린집합 $U$와 동치류 $f$에 속하는
사상 $U\to Y$가 존재하면 $f$가 \emph{점 $x\in X$에서 정의된다}고 한다.
$f$가 정의되는 점 $x\in X$들의 집합을 \emph{$f$의 정의역}이라 한다.
이는 $X$에서 모든 곳에서 조밀한 열린집합임이 자명하다.
\end{env}

\oldpage[I]{159}
\begin{proposition}[7.2.2]
\label{I.7.2.2-ko}
$X$, $Y$를 두 $S$-준스킴이라 하고, $X$는 축소 준스킴이며 $Y$는 $S$
위에서 분리되어 있다고 하자. $f$를 $X$에서 $Y$로 가는 $S$-유리
사상, $U_0$를 그 정의역이라 하자. 그러면 동치류 $f$에 속하는
$S$-사상 $U_0\to Y$가 유일하게 존재한다.
\end{proposition}

동치류 $f$에 속하는 모든 사상 $U\to Y$에 대하여 반드시
$U\subset U_0$이므로, 이 명제는 다음 보조정리에서 바로 따른다.

\begin{lemma}[7.2.2.1]
\label{I.7.2.2.1-ko}
(\hyperref[I.7.2.2-ko]{7.2.2})의 가정 아래에서, $U_1$, $U_2$를 $X$의
모든 곳에서 조밀한 두 열린집합, $f_i:U_i\to Y$($i=1,2$)를 두
$S$-사상이라 하자. $X$에서 조밀한 열린집합
$V\subset U_1\cap U_2$가 존재하여 $f_1$과 $f_2$가 그 위에서
일치하면, $f_1$과 $f_2$는 $U_1\cap U_2$ 전체에서 일치한다.
\end{lemma}

$X=U_1=U_2$인 경우만 생각하면 충분함이 자명하다. $X$(따라서 $V$)가
축소 준스킴이므로, $X$는 $V$를 포함하는 $X$의 가장 작은 닫힌 부분준스킴이다
(\hyperref[I.5.2.2-ko]{5.2.2}).
$g=(f_1,f_2)_S:X\to Y\times_S Y$라 하자. 가정에 따라 대각선
$T=\Delta_Y(Y)$는 $Y\times_S Y$의 닫힌 부분준스킴이므로,
$Z=g^{-1}(T)$는 $X$의 닫힌 부분준스킴이다
(\hyperref[I.4.4.1-ko]{4.4.1}). $h:V\to Y$를 $f_1$과 $f_2$의 $V$로의
공통 제한이라 하면, $g$의 $V$로의 제한은 $g'=(h,h)_S$이고, 이는
$g'=\Delta_Y\circ h$로 인수분해된다. $\Delta_Y^{-1}(T)=Y$이므로
$g'^{-1}(T)=V$이다. 따라서 $Z$는 $V$를 유도하는 $X$의 닫힌
부분준스킴이고 $V$를 포함하므로 $Z=X$이다. 관계 $g^{-1}(T)=X$에서
(\hyperref[I.4.4.1-ko]{4.4.1})에 따라
$g$는 $\Delta_Y\circ f$로 인수분해되며, 여기서 $f$는 $X\to Y$인
사상이다. 대각
사상의 정의에 따라 $f_1=f_2=f$이다.

(\hyperref[I.7.2.2-ko]{7.2.2})에서 정의한 사상 $U_0\to Y$는 $f$의
동치류에 속하면서 $U_0$을 진정하게 포함하는 $X$의 열린 부분집합으로
\emph{연장할 수 없는} 유일한 사상임이 자명하다.
\emph{(\hyperref[I.7.2.2-ko]{7.2.2})의 가정 아래에서}, $X$에서 $Y$로
가는 유리 사상들을 $X$의 모든 곳에서 조밀한 열린집합에서 $Y$로 가는,
더 큰 열린집합으로 \emph{연장할 수 없는} 사상들과 \emph{동일시}할 수
있다. 이 동일시 아래에서 명제
(\hyperref[I.7.2.2-ko]{7.2.2})로부터 다음이 따른다.

\begin{corollary}[7.2.3]
\label{I.7.2.3-ko}
$X$, $Y$에 대한 가정이 (\hyperref[I.7.2.2-ko]{7.2.2})와 같고,
$U$를 $X$의 모든 곳에서 조밀한 열린집합이라 하자. $U$에서 $Y$로
가는 $S$-사상들과 $U$의 모든 점에서 정의되는, $X$에서 $Y$로 가는
$S$-유리 사상들 사이에는 표준 전단사 대응이 존재한다.
\end{corollary}

실제로 (\hyperref[I.7.2.2-ko]{7.2.2})에 따라, $U$에서 $Y$로 가는 모든
$S$-사상 $f$에 대하여, $f$를 연장하는 $X$에서 $Y$로 가는 유일한 $S$-유리 사상
$\overline f$가 존재한다.

\begin{corollary}[7.2.4]
\label{I.7.2.4-ko}
$S$를 스킴, $X$를 축소 $S$-준스킴, $Y$를 $S$-스킴이라 하고,
$f:U\to Y$를 $X$의 조밀한 열린집합 $U$에서 $Y$로 가는 $S$-사상이라
하자. $\overline f$가 $f$를 연장하는, $X$에서 $Y$로 가는
$\mathbf{Z}$-유리 사상이면, $\overline f$는 그 정의역 위에서
$S$-사상이다. 따라서 이는 $f$를 연장하는, $X$에서 $Y$로 가는
$S$-유리 사상이다.
\end{corollary}

$\varphi:X\to S$, $\psi:Y\to S$를 구조 사상, $U_0$을 $\overline f$의
정의역, $j$를 포함 사상 $U_0\to X$라 하자. 그러면
$\psi\circ\overline f=\varphi\circ j$임을 보이면 충분하다. 이는 $f$가
$S$-사상이므로 (\hyperref[I.7.2.2.1-ko]{7.2.2.1})에서 바로 따른다.

\begin{corollary}[7.2.5]
\label{I.7.2.5-ko}
$X$, $Y$를 두 $S$-준스킴이라 하고, $X$는 축소 준스킴이며 $X$와 $Y$는
$S$ 위에서 분리되어 있다고 하자. $p:Y\to X$를 $S$-사상($Y$를
$X$-준스킴으로 만드는 사상), $U$를 $X$의 모든 곳에서 조밀한
열린집합, $f$를 $Y$의 $U$-단면이라 하자. 그러면 $f$를 연장하는,
$X$에서 $Y$로 가는 유리 사상 $\overline f$는 $Y$의 $X$-유리
단면이다.
\end{corollary}

\oldpage[I]{160}
$\overline f$의 정의역에서 $p\circ\overline f$가 항등사상임을 보이면
된다. $X$가 $S$ 위에서 분리되어 있으므로, 이 역시
(\hyperref[I.7.2.2.1-ko]{7.2.2.1})에서 따른다.

\begin{corollary}[7.2.6]
\label{I.7.2.6-ko}
$X$를 축소 준스킴, $U$를 $X$의 모든 곳에서 조밀한 열린집합이라 하자.
$U$ 위의 구조층 $\mathscr{O}_X$의 단면들과 $U$의 모든 점에서 정의되는
$X$ 위의 유리함수 $f$들 사이에는 표준 전단사 대응이 존재한다.
\end{corollary}

(\hyperref[I.7.2.3-ko]{7.2.3}), (\hyperref[I.7.1.2-ko]{7.1.2}),
(\hyperref[I.7.1.3-ko]{7.1.3})을 고려하면,
$X$-준스킴 $X\otimes_{\mathbf Z}\mathbf Z[T]$가 $X$ 위에서 분리되어
있음을 지적하면 충분하다(\hyperref[I.5.5.1-ko]{5.5.1, (iv)}).

\begin{corollary}[7.2.7]
\label{I.7.2.7-ko}
$Y$를 축소 준스킴, $f:X\to Y$를 분리 사상, $U$를 $Y$의 모든 곳에서
조밀한 열린집합, $g:U\to f^{-1}(U)$를 $f^{-1}(U)$의 $U$-단면이라
하자. $Z$를 $\overline{g(U)}$를 바탕공간으로 갖는 $X$의 축소
부분준스킴이라 하자(\hyperref[I.5.2.1-ko]{5.2.1}). $g$가 $X$의
$Y$-단면의 제한이기 위한 필요충분조건, 다시 말해
(\hyperref[I.7.2.5-ko]{7.2.5})에서 $g$를 연장하는 $Y$에서 $X$로 가는
유리 사상이 모든 곳에서 정의되기 위한 필요충분조건은 $f$의 $Z$로의
제한이 $Z$에서 $Y$로 가는 동형사상인 것이다.
\end{corollary}

$f$의 $f^{-1}(U)$로의 제한은 분리 사상이다
(\hyperref[I.5.5.1-ko]{5.5.1, (i)}). 따라서 $g$는 닫힌 몰입 사상이고
(\hyperref[I.5.4.6-ko]{5.4.6}),
$g(U)=Z\cap f^{-1}(U)$이다. 또한 $Z$의 열린집합 $g(U)$ 위에 $Z$가
유도하는 부분준스킴은 $g$에 대응하는 $f^{-1}(U)$의 닫힌 부분준스킴과
동일하다(\hyperref[I.5.2.1-ko]{5.2.1}). 이제 명제의 조건이 충분함은
자명하다. 실제로 그 조건이 성립하고 $f_Z:Z\to Y$가 $f$의 $Z$로의
제한이며 $\overline g:Y\to Z$가 그 역동형사상이면, $\overline g$는
$g$를 연장한다. 반대로 $g$가 $X$의 $Y$-단면 $h$의 $U$로의 제한이면,
$h$는 닫힌 몰입 사상이다
(\hyperref[I.5.4.6-ko]{5.4.6}). 따라서 $h(Y)$는 닫혀 있고 $Z$에
포함되므로 $Z$와 같다. 그러므로 (\hyperref[I.5.2.1-ko]{5.2.1})에 따라
$h$는 반드시 $Y$에서 $X$의 닫힌 부분준스킴 $Z$로 가는 동형사상이다.

\begin{env}[7.2.8]
\label{I.7.2.8-ko}
$X$, $Y$를 두 $S$-준스킴이라 하고, $X$는 축소 준스킴이며 $Y$는 $S$
위에서 분리되어 있다고 하자. $f$를 $X$에서 $Y$로 가는 $S$-유리
사상, $x$를 $X$의 점이라 하자. $f$를 표준 $S$-사상
$\operatorname{Spec}(\mathscr{O}_x)\to X$ 뒤에 합성할 수 있다
(\hyperref[I.2.4.1-ko]{2.4.1}). 다만 $\operatorname{Spec}(\mathscr{O}_x)$과
$f$의 정의역의 교집합이 이 국소 스펙트럼에서 조밀해야 하며, 여기서
$\operatorname{Spec}(\mathscr{O}_x)$는 $x\in\overline{\{z\}}$인 $z\in X$들의 집합과 동일시된다
(\hyperref[I.2.4.2-ko]{2.4.2}). 이 조밀성은 다음 두 경우에 성립한다.
\begin{enumerate}
  \item[1\textsuperscript{o}] $X$가 기약이다. 이때 앞의 축소 가정과
    함께 정역 준스킴이다. 실제로 $X$의 일반점 $\xi$는
    $\operatorname{Spec}(\mathscr{O}_x)$의 일반점이기도 하다. $f$의
    정의역 $U$가 $\xi$를 포함하므로
    $U\cap\operatorname{Spec}(\mathscr{O}_x)$도 $\xi$를 포함하고,
    따라서 $\operatorname{Spec}(\mathscr{O}_x)$에서 조밀하다.
  \item[2\textsuperscript{o}] $X$가 국소 뇌터 준스킴이다. 이 경우의
    주장은 다음 보조정리에서 따른다.
\end{enumerate}
\end{env}

\begin{lemma}[7.2.8.1]
\label{I.7.2.8.1-ko}
$X$를 그 바탕공간이 국소 뇌터인 준스킴, $x$를 $X$의 점이라 하자.
$\operatorname{Spec}(\mathscr{O}_x)$의 기약 성분들은 $x$를 포함하는
$X$의 기약 성분들과 $\operatorname{Spec}(\mathscr{O}_x)$의
교집합들이다. 열린집합 $U\subset X$에 대하여
$U\cap\operatorname{Spec}(\mathscr{O}_x)$가
$\operatorname{Spec}(\mathscr{O}_x)$에서 조밀하기 위한 필요충분조건은
이 열린집합이 $X$의 기약 성분들 가운데 $x$를 포함하는 모든 성분과
만나는 것이다. 특히
$U$가 $X$에서 조밀하면 이 조건이 성립한다.
\end{lemma}

둘째 주장은 첫째 주장에서 자명하게 따르므로 첫째 주장만 증명하면
된다. $\operatorname{Spec}(\mathscr{O}_x)$는 $x$를 포함하는 모든 아핀
열린집합 $U$에 포함되고, $x$를 포함하는 $U$의 기약 성분들은 $x$를
포함하는 $X$의 기약 성분들과 $U$의 교집합들이므로
(\hyperref[0.2.1.6-ko]{0, 2.1.6}), $X$가 환 $A$의 아핀 스킴인 경우만
생각하면 충분하다. $A_x$의 소 아이디얼들은
\oldpage[I]{161}
$A$의 소 아이디얼들 중 $\mathfrak{j}_x$에 포함되는 것들과 전단사로 대응하므로
(\hyperref[0.1.2.6-ko]{0, 1.2.6}), $A_x$의 극소 소 아이디얼들은
$\mathfrak{j}_x$에 포함되는 $A$의 극소 소 아이디얼들과 대응한다.
이로써 보조정리가 따른다(\hyperref[I.1.1.14-ko]{1.1.14}).

이제 앞의 두 경우 가운데 하나가 성립한다고 하자. $U$가 $S$-유리
사상 $f$의 정의역이면, $U\cap\operatorname{Spec}(\mathscr{O}_x)$에서
$f$와 일치하는(\hyperref[I.2.4.2-ko]{2.4.2})
$\operatorname{Spec}(\mathscr{O}_x)$에서 $Y$로 가는 유리 사상을
$f'$이라 하겠다. 이 유리 사상을 \emph{$f$가 유도한} 유리 사상이라
한다.

\begin{proposition}[7.2.9]
\label{I.7.2.9-ko}
$S$를 국소 뇌터 준스킴, $X$를 축소 $S$-준스킴, $Y$를 $S$ 위의
유한형 스킴이라 하자. 나아가 $X$가 기약이거나 국소 뇌터라고 하자.
$f$를 $X$에서 $Y$로 가는 $S$-유리 사상, $x$를 $X$의 점이라 하자.
$f$가 점 $x$에서 정의되기 위한 필요충분조건은 $f$가 유도한
(\hyperref[I.7.2.8-ko]{7.2.8})
$\operatorname{Spec}(\mathscr{O}_x)$에서 $Y$로 가는 유리 사상 $f'$이
사상인 것이다.
\end{proposition}

$\operatorname{Spec}(\mathscr{O}_x)$는 $x$를 포함하는 모든
열린집합에 포함되므로 이 조건이 필요함은 자명하다. 충분함을 증명하자.
(\hyperref[I.6.5.1-ko]{6.5.1})에 따라, $X$에서 $x$의 열린 근방 $U$와
$S$-사상 $g$가 존재하며, 이 사상은 $U$에서 $Y$로 가고
$\operatorname{Spec}(\mathscr{O}_x)$ 위에서 $f'$을 유도한다. $X$가 기약이면 $U$는 $X$에서
조밀하므로, (\hyperref[I.7.2.3-ko]{7.2.3})에 따라 $g$가 정하는
$S$-유리 사상으로 이를 보겠다. 또한 $X$의 일반점은
$\operatorname{Spec}(\mathscr{O}_x)$와 $f$의 정의역에 모두 속한다.
따라서 $f$와 $g$는 이 점에서 일치하고, 그러므로 $X$의 공집합이 아닌
어떤 열린집합에서 일치한다(\hyperref[I.6.5.1-ko]{6.5.1}). 그런데
$f$와 $g$는 $S$-유리 사상이므로 서로 같다
(\hyperref[I.7.2.3-ko]{7.2.3}). 따라서 $f$는 $x$에서 정의된다.

이제 $X$가 국소 뇌터라고 하자. $U$가 뇌터라고 가정할 수 있다.
그러면 $x$를 포함하는 $X$의 기약 성분 $X_i$는 유한 개뿐이고
(\hyperref[I.7.2.8.1-ko]{7.2.8.1}), 필요하면 $U$를 더 작은
열린집합으로 바꾸어 이 성분들만 $U$와 만나게 할 수 있다. 실제로
$U$가 뇌터이므로 $U$와 만나는 $X$의 기약 성분은 유한 개뿐이다.
앞에서와 같이 $f$와 $g$는 각 $X_i$의 공집합이 아닌 열린집합에서
일치한다. 각 $X_i$가 $\overline U$에 포함됨을 고려하여,
$U\cup(X-\overline U)$의 조밀한 열린 부분집합 위에서 정의되고,
$U$에서는 $g$와 같으며, $X-\overline U$와 $f$의 정의역의 교집합에서는
$f$와 같은 사상 $f_1$을 생각하자. $U\cup(X-\overline U)$가 $X$에서
조밀하므로 $f_1$과 $f$는 $X$의 조밀한 열린집합에서 일치한다.
$f$가 유리 사상이므로, $f$는 $f_1$의 연장이다
(\hyperref[I.7.2.3-ko]{7.2.3}). 따라서 원래 유리 사상은 점 $x$에서 정의된다.

\subsection{유리함수층}
\label{subsection:I.7.3-ko}

\begin{env}[7.3.1]
\label{I.7.3.1-ko}
$X$를 준스킴이라 하자. 모든 열린집합 $U\subset X$에 대하여 $R(U)$를
$U$ 위의 유리함수환이라 하자(\hyperref[I.7.1.3-ko]{7.1.3}). 이는
$\Gamma(U,\mathscr{O}_X)$-대수이다. 또한 $V\subset U$가 $X$의 또 다른
열린집합이면, $U$의 모든 곳에서 조밀한 열린 부분집합 위의
$\mathscr{O}_X$의 단면은 $V$로 제한하여 $V$의 모든 곳에서 조밀한
열린 부분집합 위의 단면을 주고, 두 단면이 $U$의 모든 곳에서 조밀한
열린 부분집합 위에서 일치하면 그들의 $V$로의 제한은 $V$의 모든 곳에서 조밀한
열린 부분집합 위에서 일치한다. 따라서 대수의 디-준동형
$\rho_{V,U}:R(U)\to R(V)$가 정의된다. $U\supset V\supset W$가 $X$의
세 열린집합이면
$\rho_{W,U}=\rho_{W,V}\circ\rho_{V,U}$임이 자명하다. 그러므로
$R(U)$들은 $X$ 위의 대수의 \emph{준층}을 정의한다.
\end{env}

\oldpage[I]{162}
\begin{definition}[7.3.2]
\label{I.7.3.2-ko}
준스킴 $X$ 위의 \emph{유리함수층}이라 함은 $R(U)$들로 이루어진
준층의 층화인 $\mathscr{O}_X$-대수를 말하며, 이를
$\mathscr{R}(X)$로 나타낸다.
\end{definition}

모든 준스킴 $X$와 모든 열린집합 $U\subset X$에 대하여 유도된 층
$\mathscr{R}(X)|U$는 바로 $\mathscr{R}(U)$이다.

\begin{proposition}[7.3.3]
\label{I.7.3.3-ko}
$X$를 그 기약 성분들의 족 $(X_\lambda)$가 국소 유한인 준스킴이라
하자. 특히 $X$의 바탕공간이 국소 뇌터이면 이 조건이 성립한다.
그러면 $\mathscr{O}_X$-가군 $\mathscr{R}(X)$는 준연접이고,
유한 개의 성분 $X_\lambda$만 만나는 $X$의 모든 열린집합 $U$에
대하여 $R(U)$는 $\Gamma(U,\mathscr{R}(X))$와 같으며, 이 환은
$U\cap X_\lambda\neq\varnothing$인 $X_\lambda$들의 일반점의
국소환들의 직접곱과 표준적으로 동일시된다.
\end{proposition}

$X$가 유한 개의 기약 성분 $X_i$만 갖고 그 일반점을 $x_i$
($1\leq i\leq n$)라 하는 경우만 생각하면 충분하다. 이때 $R(U)$가
$U\cap X_i\neq\varnothing$인 $\mathscr{O}_{x_i}=R(X_i)$들의
직접곱과 표준적으로 동일시된다는 것은
(\hyperref[I.7.1.7-ko]{7.1.7})에서 따른다. 또한 준층
$U\to R(U)$가 층의 공리를 만족함을 보이자. 그러면
$R(U)=\Gamma(U,\mathscr{R}(X))$가 증명된다. 실제로 앞의 결과에 따라
(F 1)이 성립한다. (F 2)를 보이기 위하여 $X$의 열린집합 $U$의
열린집합 $V_\alpha\subset U$들로 이루어진 열린 덮개를 생각하자. $s_\alpha\in R(V_\alpha)$들이
주어지고 모든 지표쌍에 대하여 $s_\alpha$와 $s_\beta$의
$V_\alpha\cap V_\beta$로의 제한이 일치한다고 하자. 그러면
$U\cap X_i\neq\varnothing$인 모든 지표 $i$에 대하여
$V_\alpha\cap X_i\neq\varnothing$인 모든 $s_\alpha$의
$R(X_i)$-성분은 같다. 이 성분을 $t_i$라 하면, $t_i$들을 성분으로
갖는 $R(U)$의 원소는 각 $V_\alpha$로 제한하면 $s_\alpha$가 된다.
마지막으로 $\mathscr{R}(X)$가 준연접임을 보이기 위해
$X=\operatorname{Spec}(A)$가 아핀인 경우만 생각하면 된다.
$f\in A$인 $D(f)$ 꼴의 아핀 열린집합을 $U$로 취하면 앞의 결과와
정의 (\hyperref[I.1.3.4-ko]{1.3.4})에 따라
$\mathscr{R}(X)=\widetilde M$이다. 여기서 $M$은 $A$-가군
$A_{x_i}$들의 직합이다.

\begin{corollary}[7.3.4]
\label{I.7.3.4-ko}
$X$를 유한 개의 기약 성분만 갖는 축소 준스킴이라 하고,
$X_i$ ($1\leq i\leq n$)를 $X$의 기약 성분을 바탕공간으로 갖는
$X$의 축소 닫힌 부분준스킴이라 하자
(\hyperref[I.5.2.1-ko]{5.2.1}). $h_i$가 $X_i\to X$인 표준 포함
사상이면 $\mathscr{R}(X)$는 $\mathscr{O}_X$-대수
$(h_i)_*(\mathscr{R}(X_i))$들의 직접곱이다.
\end{corollary}

\begin{corollary}[7.3.5]
\label{I.7.3.5-ko}
$X$가 기약이면 모든 준연접 $\mathscr{R}(X)$-가군
$\mathscr{F}$는 상수층이다.
\end{corollary}

모든 $x\in X$가 $\mathscr{F}|U$가 상수층이 되는 근방 $U$를
가짐을 보이면 충분하다(\hyperref[0.3.6.2-ko]{0, 3.6.2}).
다시 말해 $X$가 아핀인 경우로 환원된다. 더 나아가
$\mathscr{F}$가 어떤 준동형
$(\mathscr{R}(X))^{(I)}\to(\mathscr{R}(X))^{(J)}$의 여핵이라고
가정할 수 있다(\hyperref[0.5.1.3-ko]{0, 5.1.3}). 따라서
$\mathscr{R}(X)$가 상수층임만 보이면 된다. 그러나 $X$의 공집합이
아닌 임의의 열린집합을 $U$라 하면 그 $U$는 일반점을 포함하므로
$\Gamma(U,\mathscr{R}(X))=R(X)$이고, 이는 자명하다.

\begin{corollary}[7.3.6]
\label{I.7.3.6-ko}
$X$가 기약이면 모든 준연접 $\mathscr{O}_X$-가군
$\mathscr{F}$에 대하여
$\mathscr{F}\otimes_{\mathscr{O}_X}\mathscr{R}(X)$는 상수층이다.
더 나아가 $X$가 축소 준스킴이면(따라서 정역 준스킴이다),
$\mathscr{F}\otimes_{\mathscr{O}_X}\mathscr{R}(X)$는
$(\mathscr{R}(X))^{(I)}$ 꼴의 층과 동형이다.
\end{corollary}

두 번째 주장은 이 경우 $R(X)$가 체라는 사실에서 따른다.

\begin{proposition}[7.3.7]
\label{I.7.3.7-ko}
준스킴 $X$가 국소 정역 준스킴이거나 국소
\oldpage[I]{163}
뇌터 준스킴이라고 하자. 그러면 $\mathscr{R}(X)$는 준연접
$\mathscr{O}_X$-대수이다. 더 나아가 $X$가 축소 준스킴이면
(특히 $X$가 국소 정역 준스킴이면 그러하다), 표준 준동형
$\mathscr{O}_X\to\mathscr{R}(X)$는 단사이다.
\end{proposition}

이 문제는 국소적이므로 첫 번째 주장은
(\hyperref[I.7.3.3-ko]{7.3.3})에서 따르고, 두 번째 주장은
(\hyperref[I.7.2.3-ko]{7.2.3})에서 바로 따른다.

\begin{env}[7.3.8]
\label{I.7.3.8-ko}%
\footnote{\emph{[번역 주] 인쇄된 이 문단은 EGA II의 정오표에서
오류로 판정되어 전면 대체되었다. 여기서는 그 공식 대체문을 옮긴다.}}
$X$와 $Y$를 두 \emph{정역} 준스킴이라 하자. 그러면
$\mathscr{R}(X)$는 준연접 $\mathscr{O}_X$-가군이고,
$\mathscr{R}(Y)$는 준연접 $\mathscr{O}_Y$-가군이다
(\hyperref[I.7.3.3-ko]{7.3.3}). $f:X\to Y$를 \emph{우세} 사상이라
하자. 그러면 표준 $\mathscr{O}_X$-가군 준동형
\[
  \tau:f^*(\mathscr{R}(Y))\longrightarrow\mathscr{R}(X)
  \tag{7.3.8.1}\label{I.7.3.8.1-ko}
\]
이 존재한다.
\end{env}

먼저 $X=\operatorname{Spec}(A)$와 $Y=\operatorname{Spec}(B)$가
정역 $A$, $B$로 주어진 아핀 준스킴이라고 하자. 그러면 $f$는
단사 준동형 $B\to A$에 대응하고, 이 준동형은 $B$의 분수체
$L$에서 $A$의 분수체 $K$로 가는 단사 준동형 $L\to K$로
연장된다. 이때 (7.3.8.1)의 준동형은 표준 준동형
$L\otimes_B A\to K$에 대응한다
(\hyperref[I.1.6.5-ko]{1.6.5}).

일반적인 경우에는 $f(U)\subset V$인 공집합이 아닌 아핀
열린집합 $U\subset X$, $V\subset Y$의 각 쌍에 대하여 앞과 같이
준동형 $\tau_{U,V}$를 정의한다. $U'\subset U$, $V'\subset V$이고
$f(U')\subset V'$이면 $\tau_{U,V}$는 $\tau_{U',V'}$의 연장이다.
따라서 이 준동형들을 붙이면 앞의 표준
준동형을 얻는다. $x$, $y$가 각각 $X$, $Y$의 일반점이면
$f(x)=y$이고,
\[
  (f^*(\mathscr{R}(Y)))_x
  =\mathscr{O}_y\otimes_{\mathscr{O}_y}\mathscr{O}_x
  =\mathscr{O}_x
\]
이다(\hyperref[0.4.3.1-ko]{0, 4.3.1}). 따라서 $\tau_x$는
\emph{동형}이다.

\subsection{꼬임층과 꼬임 없는 층}
\label{subsection:I.7.4-ko}

\begin{env}[7.4.1]
\label{I.7.4.1-ko}
$X$를 \emph{정역} 준스킴이라 하자. 임의의 $\mathscr{O}_X$-가군
$\mathscr{F}$에 대하여 표준 준동형
$\mathscr{O}_X\to\mathscr{R}(X)$은 텐서곱을 취함으로써 준동형
(역시 \emph{표준}이라 한다)
\[
  \mathscr{F}\to
  \mathscr{F}\otimes_{\mathscr{O}_X}\mathscr{R}(X)
\]
을 정의한다. 각 줄기에서 이는 $\mathscr{F}_x$에서
$\mathscr{F}_x\otimes_{\mathscr{O}_x}R(X)$로 가는 준동형
$z\to z\otimes 1$이다. 이 준동형의 \emph{핵 $\mathscr{T}$}은
$\mathscr{F}$의 $\mathscr{O}_X$-부분가군이며, 이를
\emph{$\mathscr{F}$의 꼬임 부분층}이라 한다. $\mathscr{F}$가
준연접이면 이 부분층도 준연접이다
(\hyperref[I.4.1.1-ko]{4.1.1}과
\hyperref[I.7.3.6-ko]{7.3.6}). $\mathscr{T}=0$이면
$\mathscr{F}$에 \emph{꼬임이 없다}고 하고,
$\mathscr{T}=\mathscr{F}$이면 $\mathscr{F}$를 \emph{꼬임층}이라 한다.
임의의 $\mathscr{O}_X$-가군 $\mathscr{F}$에 대하여
$\mathscr{F}/\mathscr{T}$에는 꼬임이 없다.
(\hyperref[I.7.3.5-ko]{7.3.5})에서 다음을 얻는다.
\end{env}

\begin{proposition}[7.4.2]
\label{I.7.4.2-ko}
$X$가 정역 준스킴이면, 꼬임이 없는 모든 준연접
$\mathscr{O}_X$-가군 $\mathscr{F}$는 $(\mathscr{R}(X))^{(I)}$ 꼴의
상수층의 부분층 $\mathscr{G}$와 동형이다. 이 상수층은
$\mathscr{G}$에 의해 $\mathscr{R}(X)$-가군으로 생성된다.
\end{proposition}

$I$의 기수를 \emph{$\mathscr{F}$의 계수}라 한다. $X$의
공집합이 아닌 임의의 아핀 열린집합 $U$에 대하여 $\mathscr{F}$의
계수는 $\Gamma(U,\mathscr{F})$의 $\Gamma(U,\mathscr{O}_X)$-가군으로서의
계수와 같다. 이는 $U$에 들어 있는 $X$의 일반점을 고려하면 곧 알 수
있다. 특히 다음을 얻는다.

\begin{corollary}[7.4.3]
\label{I.7.4.3-ko}
정역 준스킴 $X$ 위에서 꼬임이 없고 계수가 $1$인 모든 준연접
$\mathscr{O}_X$-가군(특히 모든 가역 $\mathscr{O}_X$-가군)은
$\mathscr{R}(X)$의 $\mathscr{O}_X$-부분가군과 동형이며, 그 역도
성립한다.
\end{corollary}

\begin{corollary}[7.4.4]
\label{I.7.4.4-ko}
$X$를 정역 준스킴, $\mathscr{L}$과 $\mathscr{L}'$을 꼬임이 없는 두
$\mathscr{O}_X$-가군, $f$(각각 $f'$)를 $X$ 위의
$\mathscr{L}$(각각 $\mathscr{L}'$)의 단면이라 하자.
$f\otimes f'=0$일 필요충분조건은 두 단면 $f$, $f'$ 가운데 하나가
영인 것이다.
\end{corollary}

$x$를 $X$의 일반점이라 하자. 가정에 따라
$(f\otimes f')_x=f_x\otimes f'_x=0$이다. $\mathscr{L}_x$와
$\mathscr{L}'_x$는 체 $\mathscr{O}_x$의 $\mathscr{O}_x$-부분가군으로
동일시되므로, 앞의 관계에서 $f_x=0$ 또는 $f'_x=0$을 얻는다.
$\mathscr{L}$과 $\mathscr{L}'$에는 꼬임이 없으므로
(\hyperref[I.7.3.5-ko]{7.3.5}) 결국 $f=0$ 또는 $f'=0$이다.

\begin{proposition}[7.4.5]
\label{I.7.4.5-ko}
$X$, $Y$를 두 정역 준스킴, $f:X\to Y$를 우세 사상이라 하자.
꼬임이 없는 모든 준연접 $\mathscr{O}_X$-가군 $\mathscr{F}$에 대하여
$f_*(\mathscr{F})$는 꼬임이 없는 $\mathscr{O}_Y$-가군이다.
\end{proposition}

\oldpage[I]{164}
\begin{proof}
$f_*$는 왼쪽 완전하므로
(\hyperref[0.4.2.1-ko]{0, 4.2.1}),
(\hyperref[I.7.4.2-ko]{7.4.2})에 따라
$\mathscr{F}=(\mathscr{R}(X))^{(I)}$인 경우에 명제를 증명하면 충분하다.
그런데 $Y$의 공집합이 아닌 모든 열린집합 $U$는 $Y$의 일반점을
포함하므로, $f^{-1}(U)$는 $X$의 일반점을 포함한다
(\hyperref[0.2.1.5-ko]{0, 2.1.5}). 따라서
\[
  \Gamma(U,f_*(\mathscr{F}))
  =\Gamma(f^{-1}(U),\mathscr{F})=(R(X))^{(I)}.
\]
다시 말해 $f_*(\mathscr{F})$는 $(R(X))^{(I)}$를 값으로 하는
상수층이다. 이 층을 $\mathscr{R}(Y)$-가군으로 보면 명백히 꼬임이
없다.
\end{proof}

\begin{proposition}[7.4.6]
\label{I.7.4.6-ko}
$X$를 정역 준스킴, $x$를 그 일반점이라 하자. 유한 생성인 임의의
준연접 $\mathscr{O}_X$-가군 $\mathscr{F}$에 대하여 다음 조건들은
동치이다. a) $\mathscr{F}$는 꼬임층이다. b) $\mathscr{F}_x=0$이다.
c) $\operatorname{Supp}(\mathscr{F})\neq X$이다.
\end{proposition}

(\hyperref[I.7.3.5-ko]{7.3.5})와
(\hyperref[I.7.4.1-ko]{7.4.1})에 따라 관계 $\mathscr{F}_x=0$과
$\mathscr{F}\otimes_{\mathscr{O}_X}\mathscr{R}(X)=0$은 동치이므로
a)와 b)는 동치이다. 한편 $\operatorname{Supp}(\mathscr{F})$는
$X$에서 닫혀 있고(\hyperref[0.5.2.2-ko]{0, 5.2.2}), $X$의 공집합이
아닌 모든 열린집합은 $x$를 포함하므로 b)와 c)는 동치이다.

\begin{env}[7.4.7]
\label{I.7.4.7-ko}
$X$가 \emph{축소} 준스킴이고 기약성분을 \emph{유한}개만 갖는 경우에도
용어를 남용하여 (\hyperref[I.7.4.1-ko]{7.4.1})의 정의들을 사용한다.
그러면 (\hyperref[I.7.3.4-ko]{7.3.4})에 따라 이러한 준스킴에 대해서도
(\hyperref[I.7.4.6-ko]{7.4.6})의 a)와 c)는 동치이다.
\end{env}

\section{슈발레 스킴}
\label{section:I.8-ko}

\subsection{서로 연관된 국소환}
\label{subsection:I.8.1-ko}

임의의 국소환 $A$에 대하여 $A$의 극대 아이디얼을
$\mathfrak m(A)$로 나타낸다.

\begin{lemma}[8.1.1]
\label{I.8.1.1-ko}
$A$, $B$를 $A\subset B$인 두 국소환이라 하자. 다음 조건들은
동치이다.
\begin{enumerate}
  \item[(i)] $\mathfrak m(B)\cap A=\mathfrak m(A)$이다.
  \item[(ii)] $\mathfrak m(A)\subset\mathfrak m(B)$이다.
  \item[(iii)] $1$은 $\mathfrak m(A)$에 의해 생성되는 $B$의
    아이디얼에 속하지 않는다.
\end{enumerate}
\end{lemma}

(i)가 (ii)를 함의하고 (ii)가 (iii)을 함의함은 명백하다. 끝으로
(iii)이 성립하면 $\mathfrak m(B)\cap A$는 $\mathfrak m(A)$를 포함하고
$1$을 포함하지 않으므로 $\mathfrak m(A)$와 같다.

(\hyperref[I.8.1.1-ko]{8.1.1})의 동치인 조건들이 성립할 때 $B$가
$A$를 \emph{지배한다}고 한다. 이는 포함 준동형 $A\to B$가
\emph{국소} 준동형이라고 하는 것과 같다. 환 $R$의 국소 부분환들의
집합에서 지배관계가 순서관계임은 명백하다.

\begin{env}[8.1.2]
\label{I.8.1.2-ko}
이제 \emph{체} $R$을 생각하자. $R$의 임의의 부분환 $A$에 대하여
$\mathfrak p$가 $A$의 소 스펙트럼을 달릴 때의 국소환
$A_\mathfrak p$들의 집합을 $L(A)$로 나타낸다. 이 환들은 $A$를
포함하는 $R$의 부분환으로 동일시된다.
$\mathfrak p=(\mathfrak p A_\mathfrak p)\cap A$이므로,
$\operatorname{Spec}(A)$에서 $L(A)$로 가는 사상
$\mathfrak p\mapsto A_\mathfrak p$는 전단사이다.
\end{env}

\begin{lemma}[8.1.3]
\label{I.8.1.3-ko}
$R$을 체, $A$를 $R$의 부분환이라 하자. $R$의 국소 부분환 $M$이
환 $A_\mathfrak p\in L(A)$를 지배할 필요충분조건은
$A\subset M$인 것이다. 이때 $M$이 지배하는 국소환
$A_\mathfrak p$는 유일하며,
$\mathfrak p=\mathfrak m(M)\cap A$에 대응한다.
\end{lemma}

실제로 $M$이 $A_\mathfrak p$를 지배하면
(\hyperref[I.8.1.1-ko]{8.1.1})에 따라
$\mathfrak m(M)\cap A_\mathfrak p=\mathfrak p A_\mathfrak p$이므로
$\mathfrak p$는 유일하다. 다른 한편 $A\subset M$이면
$\mathfrak m(M)\cap A=\mathfrak p$는 $A$의 소 아이디얼이고,
$A-\mathfrak p\subset M$이므로 $A_\mathfrak p\subset M$이고
$\mathfrak p A_\mathfrak p\subset\mathfrak m(M)$이다. 따라서 $M$은
$A_\mathfrak p$를 지배한다.

\oldpage[I]{165}
\begin{lemma}[8.1.4]
\label{I.8.1.4-ko}
$R$을 체, $M$, $N$을 $R$의 두 국소 부분환, $P$를 $M\cup N$에 의해
생성되는 $R$의 부분환이라 하자. 다음 조건들은 동치이다.
\begin{enumerate}
  \item[(i)] $P$의 소 아이디얼 $\mathfrak p$가 존재하여
    $\mathfrak m(M)=\mathfrak p\cap M$,
    $\mathfrak m(N)=\mathfrak p\cap N$이다.
  \item[(ii)] $\mathfrak m(M)\cup\mathfrak m(N)$에 의해 $P$에서
    생성되는 아이디얼 $\mathfrak a$는 $P$와 다르다.
  \item[(iii)] $M$과 $N$을 모두 지배하는 $R$의 국소 부분환 $Q$가
    존재한다.
\end{enumerate}
\end{lemma}

(i)가 (ii)를 함의함은 명백하다. 역으로 $\mathfrak a\ne P$이면
$\mathfrak a$는 $P$의 어떤 극대 아이디얼 $\mathfrak n$에 포함된다.
$1\notin\mathfrak n$이므로 $\mathfrak n\cap M$은 $\mathfrak m(M)$을
포함하면서 $M$과 다르며, 따라서
$\mathfrak n\cap M=\mathfrak m(M)$이다. 마찬가지로
$\mathfrak n\cap N=\mathfrak m(N)$이다. $Q$가 $M$과 $N$을 지배하면
$P\subset Q$이고
\[
  \mathfrak m(M)=\mathfrak m(Q)\cap M
  =(\mathfrak m(Q)\cap P)\cap M,\qquad
  \mathfrak m(N)=(\mathfrak m(Q)\cap P)\cap N,
\]
이므로 (iii)은 (i)를 함의한다. 역은 $Q=P_\mathfrak p$로 두면
명백하다.

(\hyperref[I.8.1.4-ko]{8.1.4})의 조건들이 성립할 때 C.~슈발레를 따라
국소환 $M$과 $N$이 \emph{서로 연관되어 있다}고 한다.

\begin{proposition}[8.1.5]
\label{I.8.1.5-ko}
$A$, $B$를 체 $R$의 두 부분환, $C$를 $A\cup B$에 의해 생성되는
$R$의 부분환이라 하자. 다음 조건들은 동치이다.
\begin{enumerate}
  \item[(i)] $A$와 $B$를 포함하는 모든 국소환 $Q$에 대하여
    $\mathfrak p=\mathfrak m(Q)\cap A$,
    $\mathfrak q=\mathfrak m(Q)\cap B$로 놓으면
    $A_\mathfrak p=B_\mathfrak q$이다.
  \item[(ii)] $C$의 모든 소 아이디얼 $\mathfrak r$에 대하여
    $\mathfrak p=\mathfrak r\cap A$, $\mathfrak q=\mathfrak r\cap B$로
    놓으면 $A_\mathfrak p=B_\mathfrak q$이다.
  \item[(iii)] $M\in L(A)$와 $N\in L(B)$가 서로 연관되어 있으면
    두 환은 같다.
  \item[(iv)] $L(A)\cap L(B)=L(C)$이다.
\end{enumerate}
\end{proposition}

보조정리 (\hyperref[I.8.1.3-ko]{8.1.3})과
(\hyperref[I.8.1.4-ko]{8.1.4})는 (i)와 (iii)이 동치임을 보인다.
(i)를 $Q=C_\mathfrak r$에 적용하면 (i)가 (ii)를 함의함은 명백하다.
역으로 $Q$가 $A\cup B$를 포함하면 $C$도 포함한다. 이때
$\mathfrak r=\mathfrak m(Q)\cap C$로 놓으면
(\hyperref[I.8.1.3-ko]{8.1.3})에 따라
$\mathfrak p=\mathfrak r\cap A$이고 $\mathfrak q=\mathfrak r\cap B$이므로
(ii)는 (i)를 함의한다. (iv)가 (i)를 함의함도 곧 알 수 있다. 실제로
$Q$가 $A\cup B$를 포함하면 (\hyperref[I.8.1.3-ko]{8.1.3})에 따라
어떤 국소환 $C_\mathfrak r\in L(C)$를 지배한다. 가정에 따라
$C_\mathfrak r\in L(A)\cap L(B)$이고,
(\hyperref[I.8.1.1-ko]{8.1.1})과
(\hyperref[I.8.1.3-ko]{8.1.3})에 따라
$C_\mathfrak r=A_\mathfrak p=B_\mathfrak q$이다.

끝으로 (iii)이 (iv)를 함의함을 보이자. $Q\in L(C)$라 하자.
(\hyperref[I.8.1.3-ko]{8.1.3})에 따라 $Q$는 어떤 $M\in L(A)$와
$N\in L(B)$를 지배한다. 따라서 서로 연관된 $M$과 $N$은 가정에 따라
같다. 이제 $C\subset M$이므로 (\hyperref[I.8.1.3-ko]{8.1.3})에 따라
$M$은 어떤 $Q'\in L(C)$를 지배한다. 따라서 $Q$는 $Q'$을 지배한다.
그러므로 (\hyperref[I.8.1.3-ko]{8.1.3})에 따라 반드시
$Q=Q'=M$이고, $Q\in L(A)\cap L(B)$이다. 역으로
$Q\in L(A)\cap L(B)$이면 $C\subset Q$이므로
(\hyperref[I.8.1.3-ko]{8.1.3})에 따라 $Q$는 어떤
$Q''\in L(C)\subset L(A)\cap L(B)$를 지배한다. 서로 연관된 $Q$와
$Q''$은 같으므로 $Q''=Q\in L(C)$이다. 이로써 증명이 끝난다.

\subsection{정역 스킴의 국소환}
\label{subsection:I.8.2-ko}

\begin{env}[8.2.1]
\label{I.8.2.1-ko}
$X$를 \emph{정역} 준스킴이라 하자. 그 유리함수체 $R$은 $X$의 일반점
$a$의 국소환과 같다. 모든 $x\in X$에 대하여
$\mathscr{O}_x$는 $R$의 부분환과 표준적으로 동일시된다
(\hyperref[I.7.1.5-ko]{7.1.5}). 또한 임의의 유리함수 $f\in R$에
대하여 $f$의 정의역 $\delta(f)$는 $f\in\mathscr{O}_x$인 점
$x\in X$들의 열린집합이다. 따라서
(\hyperref[I.7.2.6-ko]{7.2.6}) 모든 열린집합 $U\subset X$에 대하여
\begin{equation}
  \Gamma(U,\mathscr{O}_X)=\bigcap_{x\in U}\mathscr{O}_x.
  \tag{8.2.1.1}\label{I.8.2.1.1-ko}
\end{equation}
\end{env}

\oldpage[I]{166}
\begin{proposition}[8.2.2]
\label{I.8.2.2-ko}
$X$를 정역 준스킴, $R$을 그 유리함수체라 하자. $X$가 스킴이기 위한
필요충분조건은 $X$의 두 점 $x$, $y$ 사이의 관계
``$\mathscr{O}_x$와 $\mathscr{O}_y$가 서로 연관되어 있다''
(\hyperref[I.8.1.4-ko]{8.1.4})가 $x=y$를 함의하는 것이다.
\end{proposition}

이 조건이 성립한다고 가정하고 $X$가 분리임을 보이자. $X$의 서로 다른
두 아핀 열린집합을 $U$, $V$라 하고, 그 환을 각각 $A$, $B$라 하여
이들을 $R$의 부분환과 동일시하자. 그러면 $U$와 $V$는 각각
(\hyperref[I.8.1.2-ko]{8.1.2})에 따라 $L(A)$와 $L(B)$로 동일시된다.
가정과 (\hyperref[I.8.1.5-ko]{8.1.5})에 따라 $C$가 $A\cup B$에 의해
$R$에서 생성되는 부분환이면 $W=U\cap V$는
$L(A)\cap L(B)=L(C)$와 동일시된다. 한편 $R$의 모든 부분환 $E$는
$L(E)$에 속하는 국소환들의 교집합과 같음이 알려져 있다
(\cite{I-1}, p.~5-03, 명제~4 \emph{bis}). 따라서 $C$는 $z\in W$에
대한 국소환 $\mathscr{O}_z$들의 교집합, 즉
(\hyperref[I.8.2.1.1-ko]{8.2.1.1})에 따라
$\Gamma(W,\mathscr{O}_X)$와 동일시된다. 이제 $W$ 위에 $X$가 유도하는
부분준스킴을 생각하자. 항등 준동형
$\varphi:C\to\Gamma(W,\mathscr{O}_X)$에 대응하여
(\hyperref[I.2.2.4-ko]{2.2.4}) 사상
$\Phi=(\psi,\theta):W\to\operatorname{Spec}(C)$가 존재한다. $\Phi$가
준스킴의 \emph{동형}임을 보이면 $W$가 \emph{아핀} 열린집합임이
따른다. $W$와 $L(C)=\operatorname{Spec}(C)$의 동일시에 따라
$\psi$는 \emph{전단사}이다. 또한 모든 $x\in W$에 대하여
$\mathfrak r=\mathfrak m_x\cap C$로 놓으면 $\theta_x^\sharp$는 포함
준동형 $C_\mathfrak r\to\mathscr{O}_x$이고, 정의에 따라
$C_\mathfrak r$는 $\mathscr{O}_x$와 동일시되므로
$\theta_x^\sharp$는 전단사이다. 따라서 $\psi$가 \emph{위상동형}, 즉
모든 닫힌 부분집합 $F\subset W$에 대하여 $\psi(F)$가
$\operatorname{Spec}(C)$에서 닫혔음을 보이는 일만 남는다. $F$는
$A$의 어떤 아이디얼 $\mathfrak a$에 대하여 $V(\mathfrak a)$ 꼴인
$U$의 닫힌 부분집합이 $W$ 위에 남기는 자취이다. 이제
$\psi(F)=V(\mathfrak a C)$임을 보이자. 그러면 주장이 따른다. 실제로
$\mathfrak a C$를 포함하는 $C$의 소 아이디얼들은 $\mathfrak a$를
포함하는 $C$의 소 아이디얼들, 즉 $x\in W$이고
$\mathfrak a\subset\mathfrak m_x$인
$\psi(x)=\mathfrak m_x\cap C$ 꼴의 아이디얼들이다. 이러한 점에 대하여
$\mathfrak a\subset\mathfrak m_x$인 것은
$x\in W\cap V(\mathfrak a)=F$인 것과 동치이다.\footnote{\emph{[번역 주]
인쇄 원문은 여기서 불가능한 등식
$x\in V(\mathfrak a)=W\cap F$를 적는다. 앞 문장의 정의에 맞는
등식 $F=W\cap V(\mathfrak a)$을 사용하였다.}} 따라서
$\psi(F)=V(\mathfrak a C)$이다.

$U\cap V$가 아핀이고 그 환 $C$가 $U$와 $V$의 두 환의 합집합
$A\cup B$에 의해 생성되므로 $X$는 분리이다
(\hyperref[I.5.5.6-ko]{5.5.6}).

역으로 $X$가 분리라고 하고, $\mathscr{O}_x$와 $\mathscr{O}_y$가 서로
연관되어 있는 $X$의 두 점을 $x$, $y$라 하자. $x$를 포함하는 아핀
열린집합 $U$의 환을 $A$, $y$를 포함하는 아핀 열린집합 $V$의 환을
$B$라 하자. 그러면 $U\cap V$는 아핀이고 그 환 $C$는 $A\cup B$에
의해 생성된다(\hyperref[I.5.5.6-ko]{5.5.6}).
$\mathfrak p=\mathfrak m_x\cap A$,
$\mathfrak q=\mathfrak m_y\cap B$라 놓으면
$A_\mathfrak p=\mathscr{O}_x$, $B_\mathfrak q=\mathscr{O}_y$이다.
$A_\mathfrak p$와 $B_\mathfrak q$가 서로 연관되어 있으므로
(\hyperref[I.8.1.4-ko]{8.1.4})
$\mathfrak p=\mathfrak r\cap A$,
$\mathfrak q=\mathfrak r\cap B$인 $C$의 소 아이디얼
$\mathfrak r$가 존재한다. $U\cap V$가 아핀이므로
$\mathfrak r=\mathfrak m_z\cap C$인 점 $z\in U\cap V$가 존재한다.
그러면 자명하게 $x=z$이고 $y=z$이므로 $x=y$이다.

\begin{corollary}[8.2.3]
\label{I.8.2.3-ko}
$X$를 정역 스킴, $x$, $y$를 $X$의 두 점이라 하자.
$x\in\overline{\{y\}}$이기 위한 필요충분조건은
$\mathscr{O}_x\subset\mathscr{O}_y$인 것, 즉 $x$에서 정의되는 모든
유리함수가 $y$에서도 정의되는 것이다.
\end{corollary}

유리함수 $f\in R$의 정의역 $\delta(f)$가 열려 있으므로 이 조건이
필요함은 자명하다. 충분함을 보이자. $\mathscr{O}_x\subset
\mathscr{O}_y$이면 (\hyperref[I.8.1.3-ko]{8.1.3})
$\mathscr{O}_y$가 $(\mathscr{O}_x)_\mathfrak p$를 지배하는
$\mathscr{O}_x$의 소 아이디얼 $\mathfrak p$가 존재한다. 그런데
(\hyperref[I.2.4.2-ko]{2.4.2})
$x\in\overline{\{z\}}$이고
$\mathscr{O}_z=(\mathscr{O}_x)_\mathfrak p$인 점 $z\in X$가 존재한다.
$\mathscr{O}_z$와 $\mathscr{O}_y$가 서로 연관되어 있으므로
(\hyperref[I.8.2.2-ko]{8.2.2})에 따라 $z=y$이고, 따름정리가 증명된다.

\begin{corollary}[8.2.4]
\label{I.8.2.4-ko}
$X$가 정역 스킴이면 사상 $x\mapsto\mathscr{O}_x$는 단사이다. 즉
$x$, $y$가 $X$의 서로 다른 두 점이면 그 가운데 한 점에서는 정의되고
다른 점에서는 정의되지 않는 유리함수가 존재한다.
\end{corollary}

\oldpage[I]{167}
이는 (\hyperref[I.8.2.3-ko]{8.2.3})과 공리
$(T_0)$ (\hyperref[I.2.1.4-ko]{2.1.4})에서 따른다.

\begin{corollary}[8.2.5]
\label{I.8.2.5-ko}
$X$를 바탕공간이 뇌터인 정역 스킴이라 하자. $f$가 $X$ 위의
유리함수체 $R$을 달릴 때 집합 $\delta(f)$들은 $X$의 위상을 생성한다.
\end{corollary}

실제로 $X$의 모든 닫힌 부분집합은 유한개의 기약 닫힌 부분집합, 즉
$\overline{\{y\}}$ 꼴인 집합들의 합집합이다
(\hyperref[I.2.1.5-ko]{2.1.5}). 그런데
$x\notin\overline{\{y\}}$이면 $x$에서는 정의되고 $y$에서는 정의되지
않는 유리함수 $f$가 존재한다(\hyperref[I.8.2.3-ko]{8.2.3}). 다시
말해 $x\in\delta(f)$이고 $\delta(f)$는 $\overline{\{y\}}$와 만나지
않는다. 따라서 $\overline{\{y\}}$의 여집합은 $\delta(f)$ 꼴인
집합들의 합집합이다. 첫 관찰에 따라 $X$의 모든 열린집합은
$\delta(f)$ 꼴인 열린집합들의 유한 교집합들의 합집합이다.

\begin{env}[8.2.6]
\label{I.8.2.6-ko}
따름정리 (\hyperref[I.8.2.5-ko]{8.2.5})에 따라 $X$의 위상은 $R$을
분수체로 갖는 국소환들의 족 $(\mathscr{O}_x)_{x\in X}$으로 완전히
특징지어진다. 동치로 $X$의 닫힌 부분집합들은 다음과 같이 정의된다.
$X$의 유한 부분집합 $\{x_1,\ldots,x_n\}$을 주고, 어떤 첨자 $i$에
대하여 $\mathscr{O}_y\subset\mathscr{O}_{x_i}$인 점 $y\in X$들의
집합을 생각하자. $\{x_1,\ldots,x_n\}$의 모든 선택에 대하여 얻는
이 집합들이 바로 $X$의 닫힌 부분집합들이다. 더욱이 $X$의 위상을
알고 나면 구조층 $\mathscr{O}_X$도 $\mathscr{O}_x$들의 족으로
결정된다. 실제로
$\Gamma(U,\mathscr{O}_X)=\bigcap_{x\in U}\mathscr{O}_x$이다
(\hyperref[I.8.2.1.1-ko]{8.2.1.1}). 따라서 $X$가 바탕공간이 뇌터인
정역 스킴이면 족 $(\mathscr{O}_x)_{x\in X}$가 준스킴 $X$를 완전히
결정한다.
\end{env}

\begin{proposition}[8.2.7]
\label{I.8.2.7-ko}
$X$, $Y$를 두 정역 스킴, $f:X\to Y$를 우세 사상
(\hyperref[I.2.2.6-ko]{2.2.6})이라 하자. $X$와 $Y$ 위의 유리함수체를
각각 $K$, $L$이라 하자. 그러면 $L$은 $K$의 부분체와 동일시되고,
모든 $x\in X$에 대하여 $\mathscr{O}_{f(x)}$는 $\mathscr{O}_x$가
지배하는 $Y$의 유일한 국소환이다.
\end{proposition}

실제로 $f=(\psi,\theta)$이고 $a$가 $X$의 일반점이면 $\psi(a)$는
$Y$의 일반점이다(\hyperref[0.2.1.5-ko]{0, 2.1.5}). 따라서
$\theta_a^\sharp$는 체 $L=\mathscr{O}_{\psi(a)}$에서 체
$K=\mathscr{O}_a$로 가는 단사 준동형이다. $Y$의 공집합이 아닌 모든
아핀 열린집합 $U$가 $\psi(a)$를 포함하므로
(\hyperref[I.2.2.4-ko]{2.2.4})에 따라 $f$에 대응하는 준동형
$\Gamma(U,\mathscr{O}_Y)\to
\Gamma(\psi^{-1}(U),\mathscr{O}_X)$는
$\theta_a^\sharp$를 $\Gamma(U,\mathscr{O}_Y)$에 제한한 준동형이다.
따라서 모든 $x\in X$에 대하여 $\theta_x^\sharp$는
$\theta_a^\sharp$를 $\mathscr{O}_{\psi(x)}$에 제한한 것이므로
단사이다. 또한 $\theta_x^\sharp$는 국소 준동형이다. 그러므로
$\theta_a^\sharp$를 통해 $L$을 $K$의 부분체와 동일시하면
$\mathscr{O}_{\psi(x)}$는 $\mathscr{O}_x$가 지배한다
(\hyperref[I.8.1.1-ko]{8.1.1}). 이것은 $\mathscr{O}_x$가 지배하는
$Y$의 유일한 국소환이다. 실제로 $Y$의 서로 연관된 두 국소환은
같기 때문이다(\hyperref[I.8.2.2-ko]{8.2.2}).

\begin{proposition}[8.2.8]
\label{I.8.2.8-ko}
$X$를 기약 준스킴, $f:X\to Y$를 국소 몰입 사상(각각 국소
동형사상)이라 하고, 더 나아가 사상 $f$가 분리라고 하자. 그러면
$f$는 몰입 사상(각각 열린 몰입 사상)이다.
\end{proposition}

$f=(\psi,\theta)$라 하자. 두 경우 모두 $\psi$가 $X$에서
$\psi(X)$로 가는 \emph{위상동형}임을 보이면 충분하다
(\hyperref[I.4.5.3-ko]{4.5.3}). $f$를 $f_{\mathrm{red}}$로 바꾸면
(\hyperref[I.5.1.6-ko]{5.1.6} 및
\hyperref[I.5.5.1-ko]{5.5.1, (vi)})에 따라 $X$와 $Y$가 축소라고 가정할 수
있다. $Y'$을 바탕공간이 $\overline{\psi(X)}$인 $Y$의 축소 닫힌
부분준스킴이라 하자. 그러면 $f$는
$X\xrightarrow{f'}Y'\xrightarrow{j}Y$로 인수분해되며 $j$는 표준 포함
사상이다(\hyperref[I.5.2.2-ko]{5.2.2}).
(\hyperref[I.5.5.1-ko]{5.5.1, (v)})에 따라 $f'$도 여전히 분리
사상이다. 또한 $f'$은 여전히
\oldpage[I]{168}
국소 몰입 사상(각각 국소 동형사상)이다. 실제로 이 문제는 $X$와
$Y$ 위에서 국소적이므로, 이를 보일 때 $f$가 닫힌 몰입 사상(각각
열린 몰입 사상)인 경우만 다루면 되고, 이때 주장은
(\hyperref[I.4.2.2-ko]{4.2.2})에서 곧바로 따른다.

따라서 $f$가 \emph{우세} 사상이라고 가정할 수 있다. 그러면 $Y$도
기약이다(\hyperref[0.2.1.5-ko]{0, 2.1.5}). 따라서 $X$와 $Y$는 모두
\emph{정역}이다. 또한 이 문제는 $Y$ 위에서 국소적이므로 $Y$가 아핀
스킴이라고 가정할 수 있다. $f$가 분리이므로 $X$는 스킴이다
(\hyperref[I.5.5.1-ko]{5.5.1, (ii)}). 이제
(\hyperref[I.8.2.7-ko]{8.2.7})의 가정들이 모두 성립한다. 그러면 모든
$x\in X$에 대하여 $\theta_x^\sharp$는 단사이다. 한편 $f$가 국소
몰입 사상이라는 가정에 따라 $\theta_x^\sharp$는 전사이다
(\hyperref[I.4.2.2-ko]{4.2.2}). 따라서 $\theta_x^\sharp$는
전단사이고, (\hyperref[I.8.2.7-ko]{8.2.7})의 동일시 아래
$\mathscr{O}_{\psi(x)}=\mathscr{O}_x$이다. 그러므로
(\hyperref[I.8.2.4-ko]{8.2.4})에 따라 $\psi$는 \emph{단사}이다. $f$가 국소
동형사상인 경우에는 이것으로 명제가 증명된다
(\hyperref[I.4.5.3-ko]{4.5.3}). $f$가 국소 몰입 사상이라고만
가정하는 경우, 모든 $x\in X$에 대하여 $X$에서 $x$의 열린 근방
$U$와 $Y$에서 $\psi(x)$의 열린 근방 $V$가 존재하여 $\psi$를 $U$에
제한한 사상은 $U$에서 $V$의 어떤 \emph{닫힌} 부분집합으로 가는
위상동형이다. 그런데 $U$는 $X$에서 조밀하므로 $\psi(U)$는 $Y$에서
조밀하고, \emph{따라서 특히} $V$에서 조밀하다. 그러므로
$\psi(U)=V$이다. $\psi$가 단사이므로 $\psi^{-1}(V)=U$이고, 이로써
$\psi$가 $X$에서 $\psi(X)$로 가는 위상동형임이 증명된다.

\subsection{슈발레 스킴}
\label{subsection:I.8.3-ko}

\begin{env}[8.3.1]
\label{I.8.3.1-ko}
$X$를 \emph{뇌터} 정역 스킴, $R$을 그 유리함수체라 하자. $x$가 $X$를
달릴 때 국소 부분환 $\mathscr{O}_x\subset R$들의 집합을
$X'$이라 쓰자. 집합 $X'$은 다음 세 조건을 만족한다.
\begin{enumerate}
  \item[(Sch. 1)] \emph{모든 $M\in X'$에 대하여 $R$은 $M$의
    분수체이다.}
  \item[(Sch. 2)] \emph{$R$의 뇌터 부분환 $A_i$가 유한개 존재하여
    $X'=\bigcup_i L(A_i)$이고, 모든 첨자쌍 $i$, $j$에 대하여
    $A_i\cup A_j$에 의해 $R$에서 생성되는 부분환 $A_{ij}$는 $A_i$
    위의 유한 생성 대수이다.}
  \item[(Sch. 3)] \emph{$X'$의 서로 연관된 두 원소 $M$, $N$은
    같다.}
\end{enumerate}
\end{env}

실제로 (Sch. 1)은 (\hyperref[I.8.2.1-ko]{8.2.1})에서 보았고,
(Sch. 3)은 (\hyperref[I.8.2.2-ko]{8.2.2})에서 따른다. (Sch. 2)를
보이기 위해서는 $X$를 환이 뇌터인 유한개의 아핀 열린집합 $U_i$로
덮고 $A_i=\Gamma(U_i,\mathscr{O}_X)$로 두면 충분하다. $X$가 스킴이라는
가정에 따라 $U_i\cap U_j$는 아핀이고
$\Gamma(U_i\cap U_j,\mathscr{O}_X)=A_{ij}$이다
(\hyperref[I.5.5.6-ko]{5.5.6}). 더욱이 $U_i$의 바탕공간이 뇌터이므로
몰입 사상 $U_i\cap U_j\to U_i$는 유한형이다
(\hyperref[I.6.3.5-ko]{6.3.5}). 따라서 $A_{ij}$는 $A_i$ 위의 유한
생성 대수이다(\hyperref[I.6.3.3-ko]{6.3.3}).

\begin{env}[8.3.2]
\label{I.8.3.2-ko}
공리 (Sch. 1), (Sch. 2), (Sch. 3)을 갖는 구조는 C.~슈발레가 말한
의미의 「스킴」을 일반화한다. 슈발레는 이에 더하여 $R$이 어떤 체
$K$의 유한 생성 확대체이고 $A_i$들이 $K$ 위의 유한 생성
대수라고 가정한다. 이 가정 아래에서는 (Sch. 2)의 일부가 불필요하다
\cite{I-1}. 역으로 집합 $X'$ 위에 이러한 구조가 주어지면
(\hyperref[I.8.2.6-ko]{8.2.6})의 관찰을 이용하여 이에 정역 스킴 $X$를
대응시킬 수 있다. $X$의 바탕공간은
(\hyperref[I.8.2.6-ko]{8.2.6})에서 정의한 위상을 갖춘 $X'$이고,
구조층 $\mathscr{O}_X$는

\oldpage[I]{169}
모든 열린집합 $U\subset X$에 대하여
$\Gamma(U,\mathscr{O}_X)=\bigcap_{x\in U}\mathscr{O}_x$를 만족하며
제한 준동형은 자명하게 정의된다. 이로써 국소환들이 정확히 $X'$의
원소들인 정역 스킴을 실제로 얻는다는 확인은 독자에게 맡긴다. 우리는
뒤에서 이 결과를 사용하지 않는다.
\end{env}

\section{준연접층에 관한 보충}
\label{section:I.9-ko}

\subsection{준연접층의 텐서곱}
\label{subsection:I.9.1-ko}

\begin{proposition}[9.1.1]
\label{I.9.1.1-ko}
$X$를 준스킴(각각 국소 뇌터 준스킴)이라 하자. $\mathscr{F}$와
$\mathscr{G}$를 두 준연접(각각 연접) $\mathscr{O}_X$-가군이라 하자.
그러면 $\mathscr{F}\otimes_{\mathscr{O}_X}\mathscr{G}$는 준연접(각각
연접)이고, $\mathscr{F}$와 $\mathscr{G}$가 유한 생성이면 이것도 유한
생성이다. $\mathscr{F}$가 유한 표시를 갖고 $\mathscr{G}$가 준연접(각각
연접)이면 $\mathscr{H}\!om_{\mathscr{O}_X}(\mathscr{F},\mathscr{G})$는
준연접(각각 연접)이다.
\end{proposition}

문제가 국소적이므로 $X$가 아핀(각각 뇌터 아핀)이라고 가정할 수 있다.
더욱이 $\mathscr{F}$가 연접이면 어떤 준동형
$\mathscr{O}_X^m\to\mathscr{O}_X^n$의 여핵이라고 가정할 수 있다.
그러면 준연접층에 관한 주장들은
(\hyperref[I.1.3.12-ko]{1.3.12})와
(\hyperref[I.1.3.9-ko]{1.3.9})에서 따르고, 연접층에 관한 주장들은
(\hyperref[I.1.5.1-ko]{1.5.1})과 다음 사실에서 따른다. 뇌터 환 $A$
위의 유한 생성 가군 $M$, $N$에 대하여 $M\otimes_A N$과
$\operatorname{Hom}_A(M,N)$은 유한 생성 $A$-가군이다.

\begin{definition}[9.1.2]
\label{I.9.1.2-ko}
$X$, $Y$를 두 $S$-준스킴, $p$, $q$를 $X\times_S Y$의 사영,
$\mathscr{F}$(각각 $\mathscr{G}$)를 준연접 $\mathscr{O}_X$-가군(각각
$\mathscr{O}_Y$-가군)이라 하자. $X\times_S Y$ 위의 텐서곱
$p^*(\mathscr{F})\otimes_{\mathscr{O}_{X\times_S Y}}q^*(\mathscr{G})$을
$\mathscr{F}$와 $\mathscr{G}$의 $\mathscr{O}_S$ 위의(또는 $S$
위의) 텐서곱이라 하고
$\mathscr{F}\otimes_{\mathscr{O}_S}\mathscr{G}$(또는
$\mathscr{F}\otimes_S\mathscr{G}$)로 쓴다.
\end{definition}

$X_i$($1\leq i\leq n$)가 $S$-준스킴이고 $\mathscr{F}_i$가 준연접
$\mathscr{O}_{X_i}$-가군($1\leq i\leq n$)이면, 같은 방식으로 준스킴
$Z=X_1\times_S X_2\times_S\cdots\times_S X_n$ 위의 텐서곱
$\mathscr{F}_1\otimes_S\mathscr{F}_2\otimes_S\cdots\otimes_S\mathscr{F}_n$
을 정의한다. (\hyperref[I.9.1.1-ko]{9.1.1})과
(\hyperref[0.5.1.4-ko]{0, 5.1.4})에 따라 이는
\emph{준연접} $\mathscr{O}_Z$-가군이다. $\mathscr{F}_i$들이 연접이고
$Z$가 \emph{국소 뇌터}이면 (\hyperref[I.9.1.1-ko]{9.1.1}),
(\hyperref[0.5.3.11-ko]{0, 5.3.11}),
(\hyperref[I.6.1.1-ko]{6.1.1})에 따라 이 가군은 \emph{연접}이다.

$X=Y=S$로 두면 정의 (\hyperref[I.9.1.2-ko]{9.1.2})에서
$\mathscr{O}_S$-가군의 텐서곱을 다시 얻는다. 또한
$q^*(\mathscr{O}_Y)=\mathscr{O}_{X\times_S Y}$이므로
(\hyperref[0.4.3.4-ko]{0, 4.3.4}),
$\mathscr{F}\otimes_S\mathscr{O}_Y$는 $p^*(\mathscr{F})$와 표준적으로
동형이고, 마찬가지로 $\mathscr{O}_X\otimes_S\mathscr{G}$는
$q^*(\mathscr{G})$와 표준적으로 동형이다. 특히 $Y=S$로 두고
$f:X\to Y$를 구조 사상이라 하면
$\mathscr{O}_X\otimes_Y\mathscr{G}=f^*(\mathscr{G})$이다. 따라서 보통의
텐서곱과 역상층은 일반적인 텐서곱의 특수한 경우로 나타난다.

정의 (\hyperref[I.9.1.2-ko]{9.1.2})에서 곧바로, $X$와 $Y$를 고정하면
$\mathscr{F}\otimes_S\mathscr{G}$가 $\mathscr{F}$와 $\mathscr{G}$ 각각에
관하여 \emph{공변이고 가법적이며 오른쪽 완전한 쌍함자}임을 얻는다.

\begin{proposition}[9.1.3]
\label{I.9.1.3-ko}
$S$, $X$, $Y$를 각각 환 $A$, $B$, $C$의 아핀 스킴이라 하고, $B$와
$C$를 $A$-대수라 하자. $M$(각각 $N$)을 $B$-가군(각각 $C$-가군),
$\mathscr{F}=\widetilde{M}$(각각 $\mathscr{G}=\widetilde{N}$)을 이에
대응하는 준연접층이라 하자. 그러면 $\mathscr{F}\otimes_S\mathscr{G}$는
$(B\otimes_A C)$-가군 $M\otimes_A N$에 대응하는 층과 표준적으로
동형이다.
\end{proposition}

\begin{proof}
\oldpage[I]{170}
(\hyperref[I.1.6.5-ko]{1.6.5})에 따라
$\mathscr{F}\otimes_S\mathscr{G}$는 $(B\otimes_A C)$-가군
\[
  \bigl(M\otimes_B(B\otimes_A C)\bigr)
  \otimes_{B\otimes_A C}
  \bigl((B\otimes_A C)\otimes_C N\bigr)
\]
에 대응하는 층과 표준적으로 동형이다. 텐서곱 사이의 표준 동형들에
따라 이 가군은
\[
  M\otimes_B(B\otimes_A C)\otimes_C N
  =(M\otimes_B B)\otimes_A(C\otimes_C N)=M\otimes_A N
\]
과 동형이다.
\end{proof}

\begin{proposition}[9.1.4]
\label{I.9.1.4-ko}
$f:T\to X$, $g:T\to Y$를 두 $S$-사상, $\mathscr{F}$(각각
$\mathscr{G}$)를 준연접 $\mathscr{O}_X$-가군(각각
$\mathscr{O}_Y$-가군)이라 하자. 그러면
\[
  (f,g)_S^*(\mathscr{F}\otimes_S\mathscr{G})
  =f^*(\mathscr{F})\otimes_{\mathscr{O}_T}g^*(\mathscr{G}).
\]
\end{proposition}

\begin{proof}
$p$, $q$가 $X\times_S Y$의 사영이면, 이 공식은
$(f,g)_S^*\circ p^*=f^*$, $(f,g)_S^*\circ q^*=g^*$
(\hyperref[0.3.5.5-ko]{0, 3.5.5})와 대수적 층들의 텐서곱의 역상층이
그 역상층들의 텐서곱이라는 사실
(\hyperref[0.4.3.3-ko]{0, 4.3.3})에서 따른다.
\end{proof}

\begin{corollary}[9.1.5]
\label{I.9.1.5-ko}
$f:X\to X'$, $g:Y\to Y'$을 두 $S$-사상, $\mathscr{F}'$(각각
$\mathscr{G}'$)을 준연접 $\mathscr{O}_{X'}$-가군(각각
$\mathscr{O}_{Y'}$-가군)이라 하자. 그러면
\[
  (f\times_S g)^*(\mathscr{F}'\otimes_S\mathscr{G}')
  =f^*(\mathscr{F}')\otimes_S g^*(\mathscr{G}').
\]
\end{corollary}

\begin{proof}
$p$, $q$를 $X\times_S Y$의 사영이라 하면
$f\times_S g=(f\circ p,g\circ q)_S$이므로, 따름정리는
(\hyperref[I.9.1.4-ko]{9.1.4})에서 따른다.
\end{proof}

\begin{corollary}[9.1.6]
\label{I.9.1.6-ko}
$X$, $Y$, $Z$를 세 $S$-준스킴, $\mathscr{F}$(각각 $\mathscr{G}$,
$\mathscr{H}$)를 준연접 $\mathscr{O}_X$-가군(각각
$\mathscr{O}_Y$-가군, $\mathscr{O}_Z$-가군)이라 하자. 층
$\mathscr{F}\otimes_S\mathscr{G}\otimes_S\mathscr{H}$는
$X\times_S Y\times_S Z$에서 $(X\times_S Y)\times_S Z$로 가는 표준
동형사상에 의한
$(\mathscr{F}\otimes_S\mathscr{G})\otimes_S\mathscr{H}$의 역상층이다.
\end{corollary}

\begin{proof}
$p_1$, $p_2$, $p_3$을 $X\times_S Y\times_S Z$의 사영이라 하면 이
동형사상은 $(p_1,p_2)_S\times_S p_3$으로 쓸 수 있다.

마찬가지로 $X\times_S Y$에서 $Y\times_S X$로 가는 표준 동형사상에
의한 $\mathscr{G}\otimes_S\mathscr{F}$의 역상층은
$\mathscr{F}\otimes_S\mathscr{G}$이다.
\end{proof}

\begin{corollary}[9.1.7]
\label{I.9.1.7-ko}
$X$가 $S$-준스킴이면, 모든 준연접 $\mathscr{O}_X$-가군
$\mathscr{F}$는 $X$에서 $X\times_S S$로 가는 표준 동형사상
(\hyperref[I.3.3.3-ko]{3.3.3})에 의한
$\mathscr{F}\otimes_S\mathscr{O}_S$의 역상층이다.
\end{corollary}

\begin{proof}
$\varphi:X\to S$가 구조 사상이면 이 동형사상은 $(1_X,\varphi)_S$이다.
따라서 따름정리는 (\hyperref[I.9.1.4-ko]{9.1.4})와
$\varphi^*(\mathscr{O}_S)=\mathscr{O}_X$에서 따른다.
\end{proof}

\begin{env}[9.1.8]
\label{I.9.1.8-ko}
$X$를 $S$-준스킴, $\mathscr{F}$를 준연접 $\mathscr{O}_X$-가군,
$\varphi:S'\to S$를 사상이라 하자. $X\times_S S'
=X_{(\varphi)}=X_{(S')}$ 위의 준연접층
$\mathscr{F}\otimes_S\mathscr{O}_{S'}$를
$\mathscr{F}_{(\varphi)}$ 또는 $\mathscr{F}_{(S')}$로 쓴다. 따라서
$p:X_{(S')}\to X$가 사영이면
$\mathscr{F}_{(S')}=p^*(\mathscr{F})$이다.
\end{env}

\begin{proposition}[9.1.9]
\label{I.9.1.9-ko}
$\varphi':S''\to S'$을 사상이라 하자. $S$-준스킴 $X$ 위의 모든 준연접
$\mathscr{O}_X$-가군 $\mathscr{F}$에 대하여
$(\mathscr{F}_{(\varphi)})_{(\varphi')}$은 표준 동형사상
$(X_{(\varphi)})_{(\varphi')}\xrightarrow{\sim}
X_{(\varphi\circ\varphi')}$
(\hyperref[I.3.3.9-ko]{3.3.9})에 의한
$\mathscr{F}_{(\varphi\circ\varphi')}$의 역상층이다.
\end{proposition}

\begin{proof}
이는 정의와 (\hyperref[I.3.3.9-ko]{3.3.9})에서 곧바로 따르며, 다음과
같이 쓸 수도 있다.
\[
  (\mathscr{F}\otimes_S\mathscr{O}_{S'})
  \otimes_{S'}\mathscr{O}_{S''}
  =\mathscr{F}\otimes_S\mathscr{O}_{S''}.
  \tag{9.1.9.1}\label{I.9.1.9.1-ko}
\]
\end{proof}

\begin{proposition}[9.1.10]
\label{I.9.1.10-ko}
$Y$를 $S$-준스킴, $f:X\to Y$를 $S$-사상이라 하자. 모든 준연접
$\mathscr{O}_Y$-가군 $\mathscr{G}$와 모든 사상 $S'\to S$에 대하여
\[
  (f_{(S')})^*(\mathscr{G}_{(S')})
  =(f^*(\mathscr{G}))_{(S')}
\]
이다.
\end{proposition}

\begin{proof}
이는 다음 그림의 가환성에서 곧바로 따른다.
\oldpage[I]{171}
\[
  \xymatrix{
    X_{(S')}\ar[r]^{f_{(S')}}\ar[d] &
    Y_{(S')}\ar[d]\\
    X\ar[r]_f &
    Y
  }
\]
\end{proof}

\begin{corollary}[9.1.11]
\label{I.9.1.11-ko}
$X$, $Y$를 두 $S$-준스킴, $\mathscr{F}$(각각 $\mathscr{G}$)를 준연접
$\mathscr{O}_X$-가군(각각 $\mathscr{O}_Y$-가군)이라 하자. 표준
동형사상 $(X\times_S Y)_{(S')}\xrightarrow{\sim}
(X_{(S')})\times_{S'}(Y_{(S')})$
(\hyperref[I.3.3.10-ko]{3.3.10})에 의한
$\mathscr{F}_{(S')}\otimes_{S'}\mathscr{G}_{(S')}$의 역상층은
$(\mathscr{F}\otimes_S\mathscr{G})_{(S')}$과 같다.
\end{corollary}

\begin{proof}
$p$, $q$가 $X\times_S Y$의 사영이면 위 동형사상은
$(p_{(S')},q_{(S')})_{S'}$이다. 따라서 따름정리는
(\hyperref[I.9.1.4-ko]{9.1.4})와
(\hyperref[I.9.1.10-ko]{9.1.10})에서 따른다.
\end{proof}

\begin{proposition}[9.1.12]
\label{I.9.1.12-ko}
(\hyperref[I.9.1.2-ko]{9.1.2})의 표기 아래에서
$z\in X\times_S Y$, $x=p(z)$, $y=q(z)$라 하자. 줄기
$(\mathscr{F}\otimes_S\mathscr{G})_z$는
$(\mathscr{F}_x\otimes_{\mathscr{O}_x}\mathscr{O}_z)
\otimes_{\mathscr{O}_z}
(\mathscr{G}_y\otimes_{\mathscr{O}_y}\mathscr{O}_z)
=\mathscr{F}_x\otimes_{\mathscr{O}_x}\mathscr{O}_z
\otimes_{\mathscr{O}_y}\mathscr{G}_y$와 동형이다.
\end{proposition}

\begin{proof}
아핀인 경우로 귀착할 수 있으므로 명제는 공식
(\hyperref[I.1.6.5.1-ko]{1.6.5.1})에서 따른다.
\end{proof}

\begin{corollary}[9.1.13]
\label{I.9.1.13-ko}
$\mathscr{F}$와 $\mathscr{G}$가 유한 생성이면
\[
  \operatorname{Supp}(\mathscr{F}\otimes_S\mathscr{G})
  =p^{-1}(\operatorname{Supp}(\mathscr{F}))
  \cap q^{-1}(\operatorname{Supp}(\mathscr{G})).
\]
\end{corollary}

\begin{proof}
$p^*(\mathscr{F})$와 $q^*(\mathscr{G})$는 유한 생성
$\mathscr{O}_{X\times_S Y}$-가군이므로,
(\hyperref[I.9.1.12-ko]{9.1.12})와
(\hyperref[0.1.7.5-ko]{0, 1.7.5})에 따라
$\mathscr{G}=\mathscr{O}_Y$인 경우, 다시 말해 다음 공식을 보이는
경우로 귀착한다.
\[
  \operatorname{Supp}(p^*(\mathscr{F}))
  =p^{-1}(\operatorname{Supp}(\mathscr{F})).
  \tag{9.1.13.1}\label{I.9.1.13.1-ko}
\]\footnote{\textbf{번역 주.} 인쇄 원문은 왼쪽에
$\operatorname{Supp}(p^{-1}(\mathscr{F}))$를 쓴다. 그러나 바로 앞의
귀착은 $\mathscr{F}\otimes_S\mathscr{O}_Y=p^*(\mathscr{F})$를 사용하고,
이어지는 줄기 계산도 역상층 $p^*(\mathscr{F})$에 관한 것이다. 따라서
형에 맞는 $p^*$를 사용하였다.}

(\hyperref[0.1.7.5-ko]{0, 1.7.5})에서와 같은 논증으로, 모든
$z\in X\times_S Y$에 대하여 $x=p(z)$라 할 때
$\mathscr{O}_z/\mathfrak{m}_x\mathscr{O}_z\neq0$임을 확인하면 충분하다.
이는 준동형 $\mathscr{O}_x\to\mathscr{O}_z$가 가정에 따라
\emph{국소}라는 사실에서 따른다.
\end{proof}

이 항에서 두 인자의 곱에 대하여 증명한 결과들을 임의의 유한 개수의
인자의 곱으로 확장하는 일은 독자에게 맡긴다.

\subsection{준연접층의 직상}
\label{subsection:I.9.2-ko}

\begin{proposition}[9.2.1]
\label{I.9.2.1-ko}
$f:X\to Y$를 준스킴의 사상이라 하자. $Y$가 다음 성질을 갖는 아핀
열린집합들의 덮개 $(Y_\alpha)$를 갖는다고 가정하자. 각
$f^{-1}(Y_\alpha)$는 $f^{-1}(Y_\alpha)$ 안에 포함된 아핀 열린집합들의 유한 덮개
$(X_{\alpha i})$를 가지며, 각 교집합
$X_{\alpha i}\cap X_{\alpha j}$ 자체가 아핀 열린집합들의 유한
합집합이다. 이 조건 아래에서 모든 준연접 $\mathscr{O}_X$-가군
$\mathscr{F}$에 대하여 $f_*(\mathscr{F})$는 준연접
$\mathscr{O}_Y$-가군이다.
\end{proposition}

\begin{proof}
문제는 $Y$ 위에서 국소적이므로 $Y$가 $Y_\alpha$들 가운데 하나와
같다고 가정하고 첨자 $\alpha$를 생략할 수 있다.

\begin{enumerate}
\item[a)] 먼저 $X_i\cap X_j$ 자체가 \emph{아핀} 열린집합이라고
가정하자. $\mathscr{F}_i=\mathscr{F}|X_i$,
$\mathscr{F}_{ij}=\mathscr{F}|(X_i\cap X_j)$라 두고,
$\mathscr{F}'_i$와 $\mathscr{F}'_{ij}$를 각각 $f$를 $X_i$와
$X_i\cap X_j$에 제한한 사상에 의한 $\mathscr{F}_i$와
$\mathscr{F}_{ij}$의 직상이라 하자. $\mathscr{F}'_i$와
$\mathscr{F}'_{ij}$는 준연접임이 알려져 있다
(\hyperref[I.1.6.3-ko]{1.6.3}). 이제
$\mathscr{G}=\bigoplus_i\mathscr{F}'_i$,
$\mathscr{H}=\bigoplus_{i,j}\mathscr{F}'_{ij}$라 두자.
$\mathscr{G}$와 $\mathscr{H}$는 준연접 $\mathscr{O}_Y$-가군이다.
준동형 $u:\mathscr{G}\to\mathscr{H}$를 정의하여
$f_*(\mathscr{F})$가 $u$의 \emph{핵}이 되게 할 것이다. 그러면
(\hyperref[I.1.3.9-ko]{1.3.9})에 따라 $f_*(\mathscr{F})$가 준연접임이
따른다. $u$를
\oldpage[I]{172}
준층의 준동형으로 정의하면 충분하다. 따라서 $\mathscr{G}$와
$\mathscr{H}$의 정의를 고려하면, 모든 열린집합 $W\subset Y$에 대하여
$W$가 변할 때 통상적 양립 조건들을 만족하는 준동형
\[
  u_W:\bigoplus_i\Gamma(f^{-1}(W)\cap X_i,\mathscr{F})
  \longrightarrow
  \bigoplus_{i,j}\Gamma(f^{-1}(W)\cap X_i\cap X_j,\mathscr{F})
\]
을 정의하면 된다. 각 단면
$s_i\in\Gamma(f^{-1}(W)\cap X_i,\mathscr{F})$에 대하여
$f^{-1}(W)\cap X_i\cap X_j$로의 제한을 $s_{i|j}$로 나타내고
\[
  u_W((s_i))=(s_{i|j}-s_{j|i})
\]
라 두면 양립 조건들이 자명하게 성립한다. $u$의 핵 $\mathscr{R}$이
$f_*(\mathscr{F})$임을 보이기 위하여, 모든 단면
$s\in\Gamma(f^{-1}(W),\mathscr{F})$에 그 제한들로 이루어진 족
$(s_i)$를 대응시켜 $f_*(\mathscr{F})$에서 $\mathscr{R}$로 가는
준동형을 정의하자. 여기서 $s_i$는 $s$를
$f^{-1}(W)\cap X_i$에 제한한 것이다. 층의 공리 (F 1)과 (F 2)
(G, II, 1.1)에 따라 이 준동형은 \emph{전단사}이고, 이 경우의 증명이
끝난다.

\item[b)] 일반적인 경우에는 $\mathscr{F}'_{ij}$가 준연접임을 보인
뒤에 같은 논증을 적용한다. 가정에 따라 $X_i\cap X_j$는 아핀
열린집합 $X_{ijk}$들의 유한 합집합이다. 그런데 $X_{ijk}$들은
\emph{어떤 스킴 안의} 아핀 열린집합들이므로 그 가운데 임의의 두
열린집합의 교집합도 아핀 열린집합이다
(\hyperref[I.5.5.6-ko]{5.5.6}). 따라서 첫째 경우로 귀착되고,
(9.2.1)이 증명된다.
\end{enumerate}
\end{proof}

\begin{corollary}[9.2.2]
\label{I.9.2.2-ko}
(9.2.1)의 결론은 다음 각 경우에 성립한다.
\begin{enumerate}
\item[a)] $f$가 분리 사상이고 준콤팩트이다.
\item[b)] $f$가 분리 사상이고 유한형이다.
\item[c)] $f$가 준콤팩트이고 $X$의 바탕공간이 국소 뇌터이다.
\end{enumerate}
\end{corollary}

\begin{proof}
a)의 경우에는 $X_{\alpha i}\cap X_{\alpha j}$가 아핀이다
(\hyperref[I.5.5.6-ko]{5.5.6}). b)는 a)의 특수한 경우이다
(\hyperref[I.6.6.3-ko]{6.6.3}). 마지막으로 c)의 경우에는 $Y$가
아핀이고 $X$의 바탕공간이 뇌터인 경우로 귀착할 수 있다. 이때 $X$는
아핀 열린집합들의 유한 덮개 $(X_i)$를 갖고, $X_i\cap X_j$는
준콤팩트이므로 아핀 열린집합들의 유한 합집합이다
(\hyperref[I.2.1.3-ko]{2.1.3}).
\end{proof}

\subsection{준연접층의 단면 연장}
\label{subsection:I.9.3-ko}

\begin{theorem}[9.3.1]
\label{I.9.3.1-ko}
$X$를 바탕공간이 뇌터인 준스킴이거나 바탕공간이 준콤팩트인 스킴이라
하자. $\mathscr{L}$을 가역 $\mathscr{O}_X$-가군
(\hyperref[0.5.4.1-ko]{0, 5.4.1}), $f$를 $X$ 위의 $\mathscr{L}$의
단면, $X_f$를 $f(x)\neq0$인 $x\in X$들의 열린집합
(\hyperref[0.5.5.1-ko]{0, 5.5.1}), $\mathscr{F}$를 준연접
$\mathscr{O}_X$-가군이라 하자.
\begin{enumerate}
\item[(i)] $s\in\Gamma(X,\mathscr{F})$이고 $s|X_f=0$이면,
  $s\otimes f^{\otimes n}=0$인 정수 $n>0$이 존재한다.
\item[(ii)] 모든 단면 $s\in\Gamma(X_f,\mathscr{F})$에 대하여,
  $s\otimes f^{\otimes n}$이 $X$ 위의
  $\mathscr{F}\otimes\mathscr{L}^{\otimes n}$의 단면으로 연장되는 정수
  $n>0$이 존재한다.
\end{enumerate}
\end{theorem}

\begin{proof}
\begin{enumerate}
\item[(i)] $X$의 바탕공간은 준콤팩트이므로,
  $\mathscr{L}|U_i$가 $\mathscr{O}_X|U_i$와 동형인 아핀 열린집합
  $U_i$들의 유한 합집합이다. 따라서 $X$가 아핀이고
  $\mathscr{L}=\mathscr{O}_X$인 경우로 귀착된다. 이 경우 $f$는
  $A(X)$의 한 원소와 동일시되고 $X_f=D(f)$이다. 또한 $s$는 어떤
  $A(X)$-가군 $M$의 한 원소와 동일시되며, $s|X_f$는 $M_f$의 대응하는
  원소와 동일시된다. 그러므로 결과는 분수가군의 정의에서 곧바로
  따른다.
\oldpage[I]{173}

\item[(ii)] 다시 $X$는 $\mathscr{L}|U_i\simeq\mathscr{O}_X|U_i$인 아핀
  열린집합 $U_i$들($1\leq i\leq r$)의 유한 합집합이다. 각 $i$에
  대하여, 앞의 동형에 의해
  $(s\otimes f^{\otimes n})|(U_i\cap X_f)$는
  $(f|(U_i\cap X_f))^n(s|(U_i\cap X_f))$와 동일시된다.
  그러면 (\hyperref[I.1.4.1-ko]{1.4.1})에 따라 모든 $i$에 대하여
  $(s\otimes f^{\otimes n})|(U_i\cap X_f)$가 $U_i$ 위의
  $\mathscr{F}\otimes\mathscr{L}^{\otimes n}$의 단면 $s_i$로 연장되는
  정수 $n>0$이 존재한다. $s_{i|j}$를 $s_i$의 $U_i\cap U_j$로의
  제한이라 하자. 정의에 따라 $X_f\cap U_i\cap U_j$에서
  $s_{i|j}-s_{j|i}=0$이다. 그런데 $X$의 바탕공간이 뇌터이면
  $U_i\cap U_j$는 준콤팩트이고, $X$가 스킴이면 $U_i\cap U_j$는 아핀
  열린집합이므로 (\hyperref[I.5.5.6-ko]{5.5.6}) 역시 준콤팩트이다.
  따라서 (i)에 의해 $i$와 $j$에 무관한 정수 $m$이 존재하여
  $(s_{i|j}-s_{j|i})\otimes f^{\otimes m}=0$이다. 이에 따라 $X$ 위의
  $\mathscr{F}\otimes\mathscr{L}^{\otimes(n+m)}$의 단면 $s'$가 존재하여,
  각 $U_i$ 위에서의 제한은 $s_i\otimes f^{\otimes m}$이고, 따라서
  $X_f$ 위에서의 제한은 $s\otimes f^{\otimes(n+m)}$이다.
\end{enumerate}
\end{proof}

다음 따름정리들은 정리 (9.3.1)을 더 대수적인 언어로 해석한다.

\begin{corollary}[9.3.2]
\label{I.9.3.2-ko}
(9.3.1)의 가정 아래에서 등급환 $A_* = \Gamma_*(\mathscr{L})$과
등급 $A_*$-가군 $M_* = \Gamma_*(\mathscr{L},\mathscr{F})$를 생각하자
(\hyperref[0.5.4.6-ko]{0, 5.4.6}). $f\in A_n$이고
$n\in\mathbb{Z}$이면 표준 동형
\[
  \Gamma(X_f,\mathscr{F})\xrightarrow{\sim}((M_*)_f)_0
\]
이 존재한다. 여기서 오른쪽은 분수가군 $(M_*)_f$의 차수 $0$인
원소들로 이루어진 부분군이다.
\end{corollary}

\begin{corollary}[9.3.3]
\label{I.9.3.3-ko}
(9.3.1)의 가정에 더하여 $\mathscr{L}=\mathscr{O}_X$라고 하자. 이때
$A=\Gamma(X,\mathscr{O}_X)$, $M=\Gamma(X,\mathscr{F})$로 놓으면,
$A_f$-가군 $\Gamma(X_f,\mathscr{F})$는 $M_f$와 표준적으로 동형이다.
\end{corollary}

\begin{proposition}[9.3.4]
\label{I.9.3.4-ko}
$X$를 뇌터 준스킴, $\mathscr{F}$를 연접 $\mathscr{O}_X$-가군,
$\mathscr{J}$를 $\mathscr{O}_X$의 연접 아이디얼층이라 하자.
$\mathscr{F}$의 지지집합이 $\mathscr{O}_X/\mathscr{J}$의 지지집합에
포함되면, $\mathscr{J}^n\mathscr{F}=0$인 정수 $n>0$이 존재한다.
\end{proposition}

\begin{proof}
$X$는 환이 뇌터인 아핀 열린집합들의 유한 합집합이므로, $X$가 뇌터
환 $A$의 아핀 스킴인 경우로 귀착할 수 있다. 그러면
$\mathscr{F}=\widetilde M$이고, 여기서
$M=\Gamma(X,\mathscr{F})$는 유한 생성 $A$-가군이다. 또한
$\mathscr{J}=\widetilde{\mathfrak{J}}$이고, 여기서
$\mathfrak{J}=\Gamma(X,\mathscr{J})$는 $A$의 아이디얼이다
(\hyperref[I.1.4.1-ko]{1.4.1}과 \hyperref[I.1.5.1-ko]{1.5.1}).
$A$가 뇌터이므로 $\mathfrak{J}$는 유한 생성계 $f_i$들
($1\leq i\leq m$)을 갖는다. 가정에 따라 $X$ 위의 $\mathscr{F}$의 모든
단면은 각 $D(f_i)$ 위에서 영단면이다. $s_j$들($1\leq j\leq q$)이
$M$을 생성하는 $\mathscr{F}$의 단면들이면,
(\hyperref[I.1.4.1-ko]{1.4.1})에 따라 $i$와 $j$에 무관한 정수 $h$가
존재하여 $f_i^h s_j=0$이다. 따라서 모든 $s\in M$에 대하여
$f_i^h s=0$이다. 그러므로 $n=mh$로 놓으면
$\mathfrak{J}^nM=0$이고, 이에 대응하는 $\mathscr{O}_X$-가군
$\mathscr{J}^n\mathscr{F}=\widetilde{\mathfrak{J}^nM}$
(\hyperref[I.1.3.13-ko]{1.3.13})은 영가군이다.
\end{proof}

\begin{corollary}[9.3.5]
\label{I.9.3.5-ko}
(9.3.4)의 가정 아래에서, 바탕공간이
$\mathscr{O}_X/\mathscr{J}$의 지지집합인 $X$의 닫힌 부분준스킴 $Y$가
존재한다. 또한 $j:Y\to X$가 표준 포함 사상이면
$\mathscr{F}=j_*(j^*(\mathscr{F}))$이다.
\end{corollary}

\begin{proof}
먼저 $\mathscr{O}_X/\mathscr{J}$와
$\mathscr{O}_X/\mathscr{J}^n$의 지지집합은 같다. 실제로
$\mathscr{J}_x=\mathscr{O}_x$이면
$\mathscr{J}_x^n=\mathscr{O}_x$이고, 한편 모든 $x\in X$에 대하여
$\mathscr{J}_x^n\subset\mathscr{J}_x$이다. 따라서 (9.3.4)에 의해
$\mathscr{J}\mathscr{F}=0$이라고 가정할 수 있다. 이제 $Y$를
$\mathscr{J}$로 정의되는 $X$의 닫힌 부분준스킴으로 잡으면,
$\mathscr{F}$가 $(\mathscr{O}_X/\mathscr{J})$-가군이므로 결론은
곧바로 따른다.
\end{proof}

\oldpage[I]{174}

\subsection{준연접층의 연장}
\label{subsection:I.9.4-ko}

\begin{env}[9.4.1]
\label{I.9.4.1-ko}
$X$를 위상공간, $\mathscr{F}$를 $X$ 위의 집합의 층(각각 군의 층,
환의 층), $U$를 $X$의 열린부분집합, $\psi:U\to X$를 표준 포함 사상,
$\mathscr{G}$를 $\mathscr{F}|U=\psi^*(\mathscr{F})$의 부분층이라 하자.
$\psi_*$는 왼쪽 완전하므로 $\psi_*(\mathscr{G})$는
$\psi_*(\psi^*(\mathscr{F}))$의 부분층이다. 표준 준동형
$\rho:\mathscr{F}\to\psi_*(\psi^*(\mathscr{F}))$
(\hyperref[0.3.5.3-ko]{0, 3.5.3})을 생각할 때,
$\mathscr{F}$의 부분층
$\rho^{-1}(\psi_*(\mathscr{G}))$를
$\overline{\mathscr{G}}$로 나타내겠다. 정의로부터 곧바로, $X$의 모든
열린집합 $V$에 대하여 $\Gamma(V,\overline{\mathscr{G}})$는
$V\cap U$로의 제한이 $V\cap U$ 위의 $\mathscr{G}$의 단면인
$s\in\Gamma(V,\mathscr{F})$들로 이루어짐을 알 수 있다. 따라서
$\overline{\mathscr{G}}|U=\psi^*(\overline{\mathscr{G}})=\mathscr{G}$이고,
$\overline{\mathscr{G}}$는 $U$ 위에 $\mathscr{G}$를 유도하는
$\mathscr{F}$의 \emph{가장 큰} 부분층이다. 이
$\overline{\mathscr{G}}$를 $\mathscr{F}|U$의 부분층
$\mathscr{G}$를 $\mathscr{F}$의 부분층으로 만드는
\emph{표준 연장}이라 하겠다.
\end{env}

\begin{proposition}[9.4.2]
\label{I.9.4.2-ko}
$X$를 준스킴, $U$를 표준 포함 사상 $j:U\to X$가 준콤팩트 사상인
$X$의 열린부분집합이라 하자
\emph{(이는 $X$의 바탕공간이 \emph{국소 뇌터}이면
\emph{모든} $U$에 대하여 성립한다
(\hyperref[I.6.6.4-ko]{6.6.4, (i)}))}. 그러면 다음이 성립한다.
\begin{enumerate}
\item[(i)] 모든 준연접 $(\mathscr{O}_X|U)$-가군 $\mathscr{G}$에 대하여
  $j_*(\mathscr{G})$는 준연접 $\mathscr{O}_X$-가군이고
  $j_*(\mathscr{G})|U=j^*(j_*(\mathscr{G}))=\mathscr{G}$이다.
\item[(ii)] 모든 준연접 $\mathscr{O}_X$-가군 $\mathscr{F}$와
  그 제한의 모든 준연접 부분-$(\mathscr{O}_X|U)$-가군
  $\mathscr{G}$에 대하여, $\mathscr{G}$의 표준 연장
  $\overline{\mathscr{G}}$(9.4.1)는 $\mathscr{F}$의 준연접
  부분-$\mathscr{O}_X$-가군이다.
\end{enumerate}
\end{proposition}

\begin{proof}
$j=(\psi,\theta)$라 하자. 여기서 $\psi$는 바탕공간의 포함 사상
$U\to X$이다. 모든 $(\mathscr{O}_X|U)$-가군 $\mathscr{G}$에 대하여
정의상 $j_*(\mathscr{G})=\psi_*(\mathscr{G})$이고, 열린집합 위에 유도된
준스킴의 정의에 따라 모든 $\mathscr{O}_X$-가군 $\mathscr{H}$에 대하여
$j^*(\mathscr{H})=\psi^*(\mathscr{H})=\mathscr{H}|U$이다. 따라서 (i)는
(\hyperref[I.9.2.2-ko]{9.2.2, a)})의 특수한 경우이다. 같은 이유로
$j_*(j^*(\mathscr{F}))$는 준연접이고,
$\overline{\mathscr{G}}$는 준동형
$\rho:\mathscr{F}\to j_*(j^*(\mathscr{F}))$에 의한
$j_*(\mathscr{G})$의 역상이므로 (ii)는
(\hyperref[I.4.1.1-ko]{4.1.1})에서 따른다.
\end{proof}

사상 $j:U\to X$가 준콤팩트라는 가정은 열린집합 $U$가
\emph{준콤팩트}이고 $X$가 \emph{스킴}일 때에도 성립한다. 실제로
$U$는 유한개의 아핀 열린집합 $U_i$들의 합집합이고, $X$의 모든 아핀
열린집합 $V$에 대하여 $V\cap U_i$는 아핀 열린집합
(\hyperref[I.5.5.6-ko]{5.5.6})이므로 준콤팩트이다.

\begin{corollary}[9.4.3]
\label{I.9.4.3-ko}
$X$를 준스킴, $U$를 포함 사상 $j:U\to X$가 준콤팩트인 $X$의
준콤팩트 열린집합이라 하자. 또한 모든 준연접 $\mathscr{O}_X$-가군이
그 유한 생성 준연접 부분-$\mathscr{O}_X$-가군들의 귀납극한이라고
가정하자 \emph{(이는 $X$가 \emph{아핀 스킴}이면 성립한다)}.
$\mathscr{F}$를 준연접 $\mathscr{O}_X$-가군,
$\mathscr{G}$를 $\mathscr{F}|U$의 유한 생성 준연접
부분-$(\mathscr{O}_X|U)$-가군이라 하자. 그러면 $\mathscr{F}$의
유한 생성 준연접 부분-$\mathscr{O}_X$-가군 $\mathscr{G}'$이 존재하여
$\mathscr{G}'|U=\mathscr{G}$이다.
\end{corollary}

\begin{proof}
실제로 $\mathscr{G}=\overline{\mathscr{G}}|U$이고, (9.4.2)에 따라
$\overline{\mathscr{G}}$는 준연접이므로 그 유한 생성 준연접
부분-$\mathscr{O}_X$-가군 $\mathscr{H}_\lambda$들의 귀납극한이다.
따라서 $\mathscr{G}$는 $\mathscr{H}_\lambda|U$들의 귀납극한이고,
유한 생성이므로 (\hyperref[0.5.2.3-ko]{0, 5.2.3})에 따라 어떤
$\mathscr{H}_\lambda|U$와 같다.
\end{proof}

\begin{remark}[9.4.4]
\label{I.9.4.4-ko}
\emph{모든} 아핀 열린집합 $U\subset X$에 대하여 포함 사상
$U\to X$가 준콤팩트라고 가정하자. 모든 아핀 열린집합 $U$와
$\mathscr{F}|U$의 모든 유한 생성 준연접
부분-$(\mathscr{O}_X|U)$-가군 $\mathscr{G}$에 대하여 (9.4.3)의 결론이
성립한다면,
\oldpage[I]{175}
$\mathscr{F}$는 그 유한 생성 준연접 부분-$\mathscr{O}_X$-가군들의
귀납극한이다. 실제로 모든 아핀 열린집합 $U\subset X$에 대하여
$\mathscr{F}|U=\widetilde{M}$이고, 여기서 $M$은 $A(U)$-가군이다.
이 가군은 그 유한 생성 부분가군들의 귀납극한이므로,
$\mathscr{F}|U$는 그 유한 생성 준연접 부분-$(\mathscr{O}_X|U)$-가군들의
귀납극한이다(\hyperref[I.1.3.9-ko]{1.3.9}). 그런데 가정에 따라 이러한
부분가군 각각은 $\mathscr{F}$의 유한 생성 준연접
부분-$\mathscr{O}_X$-가군 $\mathscr{G}_{\lambda,U}$가 $U$ 위에 유도하는
가군이다. $\mathscr{G}_{\lambda,U}$들의 유한합도 다시 유한 생성
준연접 $\mathscr{O}_X$-가군이다. 실제로 이 문제는 국소적이고,
$X$가 아핀인 경우는 (\hyperref[I.1.3.10-ko]{1.3.10})에서 다루었다.
이제 $\mathscr{F}$가 이러한 유한합들의 귀납극한임은 분명하며,
주장이 증명되었다.
\end{remark}

\begin{corollary}[9.4.5]
\label{I.9.4.5-ko}
(9.4.3)의 가정 아래에서 모든 유한 생성 준연접
$(\mathscr{O}_X|U)$-가군 $\mathscr{G}$에 대하여 유한 생성 준연접
$\mathscr{O}_X$-가군 $\mathscr{G}'$이 존재하여
$\mathscr{G}'|U=\mathscr{G}$이다.
\end{corollary}

\begin{proof}
$\mathscr{F}=j_*(\mathscr{G})$는 (9.4.2)에 따라 준연접이고
$\mathscr{F}|U=\mathscr{G}$이므로, $\mathscr{F}$에 (9.4.3)을 적용하면
된다.
\end{proof}

\begin{lemma}[9.4.6]
\label{I.9.4.6-ko}
$X$를 준스킴, $L$을 정렬집합,
$(V_\lambda)_{\lambda\in L}$를 $X$의 아핀 열린집합들로 이루어진 덮개,
$U$를 $X$의 열린집합이라 하자. 모든 $\lambda\in L$에 대하여
$W_\lambda=\bigcup_{\mu<\lambda}V_\mu$로 놓고, 다음을 가정하자.
\begin{enumerate}
  \item[1\textsuperscript{o}] 모든 $\lambda\in L$에 대하여
    $V_\lambda\cap W_\lambda$는 준콤팩트이다.
  \item[2\textsuperscript{o}] 몰입 사상 $U\to X$는 준콤팩트이다.
\end{enumerate}
그러면 모든 준연접 $\mathscr{O}_X$-가군 $\mathscr{F}$와
$\mathscr{F}|U$의 모든 유한 생성 준연접
부분-$(\mathscr{O}_X|U)$-가군 $\mathscr{G}$에 대하여 $\mathscr{F}$의
유한 생성 준연접 부분-$\mathscr{O}_X$-가군 $\mathscr{G}'$이 존재하여
$\mathscr{G}'|U=\mathscr{G}$이다.
\end{lemma}

\begin{proof}
$L$에 그 모든 원소보다 큰 새 원소 $\infty$를 덧붙인 정렬집합을
$\widehat L=L\cup\{\infty\}$라 하고, $\lambda\in\widehat L$에 대하여
$U_\lambda=U\cup\bigcup_{\substack{\mu\in L\\ \mu<\lambda}}V_\mu$로 놓자.
초한 귀납법으로 족 $(\mathscr{G}'_\lambda)_{\lambda\in\widehat L}$를
정의하겠다. 여기서
$\mathscr{G}'_\lambda$는 $\mathscr{F}|U_\lambda$의 유한 생성 준연접
부분-$(\mathscr{O}_X|U_\lambda)$-가군이고,
$\mu<\lambda$이면
$\mathscr{G}'_\lambda|U_\mu=\mathscr{G}'_\mu$이며
$\mathscr{G}'_\lambda|U=\mathscr{G}$이다. 이 족은
(\hyperref[0.3.3.1-ko]{0, 3.3.1})에 따라 서로 붙으며,
$U_\infty=X$이므로 $\mathscr{G}'=\mathscr{G}'_\infty$가 결론을 준다.
이제 $\widehat L$의 최소 원소 $\lambda_0$에 대해서는
$U_{\lambda_0}=U$이고 $\mathscr{G}'_{\lambda_0}=\mathscr{G}$로 둔다.
그 뒤
$\mu<\lambda$인 $\mathscr{G}'_\mu$들이 정의되어 앞의 성질들을
만족한다고 가정하자. $\lambda$가 최소 원소가 아닌 극한 원소이면,
모든 $\mu<\lambda$에 대하여
$\mathscr{G}'_\lambda|U_\mu=\mathscr{G}'_\mu$인
$\mathscr{F}|U_\lambda$의 유일한 부분-$(\mathscr{O}_X|U_\lambda)$-가군을
$\mathscr{G}'_\lambda$로 잡는다. $\mu<\lambda$인 $U_\mu$들이
$U_\lambda$를 덮으므로 이렇게 할 수 있다. 반대로
$\lambda=\mu+1$이면 $U_\lambda=U_\mu\cup V_\mu$이고,
$\mathscr{F}|V_\mu$의 유한 생성 준연접
부분-$(\mathscr{O}_X|V_\mu)$-가군 $\mathscr{G}''_\mu$를
\[
  \mathscr{G}''_\mu|(U_\mu\cap V_\mu)
  =\mathscr{G}'_\mu|(U_\mu\cap V_\mu)~;
\]
가 되도록 정의하면 충분하다. 그런 다음
$\mathscr{G}'_\lambda|U_\mu=\mathscr{G}'_\mu$이고
$\mathscr{G}'_\lambda|V_\mu=\mathscr{G}''_\mu$인
$\mathscr{F}|U_\lambda$의 부분-$(\mathscr{O}_X|U_\lambda)$-가군을
$\mathscr{G}'_\lambda$로 잡는다
(\hyperref[0.3.3.1-ko]{0, 3.3.1}). $V_\mu$가 아핀이므로,
$U_\mu\cap V_\mu$가 준콤팩트임을 보이기만 하면
$\mathscr{G}''_\mu$의 존재는 (9.4.3)에서 따른다. 그런데
$U_\mu\cap V_\mu$는 $U\cap V_\mu$와 $W_\mu\cap V_\mu$의 합집합이고,
둘 다 가정에 따라 준콤팩트이다.\footnote{\emph{[번역 주] 인쇄 원문의
귀납 지표는 $L$뿐이어서
$L$이 최대 원소를 가지면 $U_\lambda$들이 마지막 $V_\lambda$를 덮지
못하고, 최소 원소 단계에서는 앞선 $U_\mu$들이 $U$를 덮는다는 주장도
성립하지 않는다. 위에서는 새 종단 원소 $\infty$를 붙이고 최소 원소
단계를 따로 정의하여 두 지표 결함을 바로잡았다.}}
\end{proof}

\begin{theorem}[9.4.7]
\label{I.9.4.7-ko}
$X$를 준스킴, $U$를 $X$의 열린집합이라 하자. 다음 조건 중 하나가
성립한다고 가정하자.
\begin{enumerate}
  \item[(a)] $X$의 바탕공간은 국소 뇌터이다.
  \item[(b)] $X$는 준콤팩트 스킴이고 $U$는 준콤팩트 열린집합이다.
\end{enumerate}
그러면 모든 준연접 $\mathscr{O}_X$-가군 $\mathscr{F}$와
$\mathscr{F}|U$의 모든 유한 생성 준연접
부분-$(\mathscr{O}_X|U)$-가군 $\mathscr{G}$에 대하여 $\mathscr{F}$의
유한 생성 준연접 부분-$\mathscr{O}_X$-가군 $\mathscr{G}'$이 존재하여
$\mathscr{G}'|U=\mathscr{G}$이다.
\end{theorem}

\oldpage[I]{176}

\begin{proof}
$(V_\lambda)_{\lambda\in L}$를 $X$의 아핀 열린집합들로 이루어진
덮개라 하자. 경우 b)에서는 $L$이 유한하다고 가정한다. $L$에
정렬 순서를 주면, 보조정리 (9.4.6)의 조건들이 성립함을 확인하면
충분하다. 경우 a)에서는 공간 $V_\lambda$들이 뇌터이므로 이는
명백하다. 경우 b)에서는 $V_\lambda\cap V_\mu$가 아핀
(\hyperref[I.5.5.6-ko]{5.5.6})이므로 준콤팩트이고, $L$이 유한하므로
$V_\lambda\cap W_\lambda$는 준콤팩트이다. 이로써 정리가 증명되었다.
\end{proof}

\begin{corollary}[9.4.8]
\label{I.9.4.8-ko}
(9.4.7)의 가정 아래에서 모든 유한 생성 준연접
$(\mathscr{O}_X|U)$-가군 $\mathscr{G}$에 대하여 유한 생성 준연접
$\mathscr{O}_X$-가군 $\mathscr{G}'$이 존재하여
$\mathscr{G}'|U=\mathscr{G}$이다.
\end{corollary}

\begin{proof}
$\mathscr{F}=j_*(\mathscr{G})$에 (9.4.7)을 적용하면 된다. 이 가군은
(9.4.2)에 따라 준연접이고 $\mathscr{F}|U=\mathscr{G}$이다.
\end{proof}

\begin{corollary}[9.4.9]
\label{I.9.4.9-ko}
$X$를 바탕공간이 국소 뇌터인 준스킴 또는 준콤팩트 스킴이라 하자.
그러면 모든 준연접 $\mathscr{O}_X$-가군은 그 유한 생성 준연접
부분-$\mathscr{O}_X$-가군들의 귀납극한이다.
\end{corollary}

\begin{proof}
이는 (9.4.7)과 주 (9.4.4)에서 따른다.
\end{proof}

\begin{corollary}[9.4.10]
\label{I.9.4.10-ko}
(9.4.9)의 가정 아래에서 준연접 $\mathscr{O}_X$-가군
$\mathscr{F}$를 생각하자. $\mathscr{F}$의 모든 유한 생성 준연접
부분-$\mathscr{O}_X$-가군이 $X$ 위의 단면들로 생성된다고 가정하면,
$\mathscr{F}$도 $X$ 위의 단면들로 생성된다.
\end{corollary}

\begin{proof}
$U$를 점 $x\in X$의 아핀 열린 근방, $s$를 $U$ 위의
$\mathscr{F}$의 단면이라 하자. $s$로 생성되는 $\mathscr{F}|U$의
부분-$(\mathscr{O}_X|U)$-가군 $\mathscr{G}$는 준연접이고 유한 생성이므로,
$\mathscr{F}$의 유한 생성 준연접 부분-$\mathscr{O}_X$-가군
$\mathscr{G}'$이 존재하여 $\mathscr{G}'|U=\mathscr{G}$이다(9.4.7).
가정에 따라 $X$ 위의 $\mathscr{G}'$의 유한개의 단면 $t_i$와
$x$의 근방 $V\subset U$ 위의 $\mathscr{O}_X$의 단면 $a_i$들이 존재하여
$s|V=\sum_i a_i\mathbin{\cdot}(t_i|V)$이다. 이로써 따름정리가
증명되었다.
\end{proof}

\subsection{준스킴의 닫힌 상. 부분준스킴의 폐포.}
\label{subsection:I.9.5-ko}

\begin{proposition}[9.5.1]
\label{I.9.5.1-ko}
$f:X\to Y$를 $f_*(\mathscr{O}_X)$가 준연접
$\mathscr{O}_Y$-가군인 준스킴의 사상이라 하자(이는 $f$가 준콤팩트이고,
또한 $f$가 분리 사상이거나 $X$가 국소 뇌터이면 성립한다
(\hyperref[I.9.2.2-ko]{9.2.2})). 그러면 $Y$의 가장 작은
닫힌 부분준스킴\footnote{\textbf{번역 주.} 인쇄 원문은 ‘닫힌’ 없이
단순히 “가장 작은 부분준스킴”이라고 쓴다.
그러나 이어지는 (9.5.2)의 구성과 (9.5.3)의 정의는 닫힌
부분준스킴 가운데의 최소성을 증명하고 정의한다. 실제로 조밀한
준콤팩트 열린 몰입 $D(t)\to\operatorname{Spec}(k[t])$에서는 핵으로
정의되는 닫힌 부분준스킴이 $\operatorname{Spec}(k[t])$ 전체인 반면,
$f$는 진부분 열린 부분준스킴 $D(t)$를 통하여 인수분해된다. 따라서
핵으로 정의된 대상이 모든 부분준스킴 가운데 최소라는 인쇄문과
(9.5.2)의 결합된 주장은 성립하지 않으며, 여기서는 논증과 정의에
필수적인 ‘닫힌’을 보충했다. 이는 공식 정오표가 아니라 뒤의
구성·정의와 위 반례가 강제하는 편집 교정이다.}
$Y'$이 존재하여, 표준 포함 사상 $j:Y'\to Y$를 통하여 $f$가
인수분해된다(또는 동치로
(\hyperref[I.4.4.1-ko]{4.4.1}), $X$의 부분준스킴
$f^{-1}(Y')$이 $X$와 \emph{동일하다}).
\end{proposition}

좀 더 정확히 말하면 다음과 같다.

\begin{corollary}[9.5.2]
\label{I.9.5.2-ko}
(9.5.1)의 가정 아래에서 $f=(\psi,\theta)$라 하고, 준동형
$\theta:\mathscr{O}_Y\to f_*(\mathscr{O}_X)$의 준연접 핵을
$\mathscr{J}$라 하자. 그러면 $\mathscr{J}$로 정의된 $Y$의 닫힌
부분준스킴 $Y'$은 (9.5.1)의 조건을 만족한다.
\end{corollary}

\begin{proof}
함자 $\psi^*$가 완전하므로, 표준 인수분해
$\theta:\mathscr{O}_Y\to\mathscr{O}_Y/\mathscr{J}
\xrightarrow{\theta'}\psi_*(\mathscr{O}_X)$는
(\hyperref[0.3.5.4.3-ko]{0, 3.5.4.3})에 따라 다음 인수분해를 준다.
\[
  \theta^\sharp:\psi^*(\mathscr{O}_Y)
  \to\psi^*(\mathscr{O}_Y)/\psi^*(\mathscr{J})
  \xrightarrow{{\theta'}^\sharp}\mathscr{O}_X~;
\]
모든 $x\in X$에 대하여 $\theta_x^\sharp$가 국소 준동형이므로
${\theta'_x}^\sharp$도 그러하다. 연속사상 $\psi$를 $X$에서
$Y'$로 가는 사상으로 보아 $\psi_0$라 하고,\footnote{\textbf{번역 주.}
인쇄 원문은 이 증명의 해당 구간에서
$X'$, $X''$를 쓰지만, 출판된 EGA II 정오표는 이들을 전부
$Y'$, $Y''$로 고치도록 지시한다. 여기서는 그 공식 교정을
적용하였다.}
$\theta_0$를 제한 사상
$\theta'|Y':(\mathscr{O}_Y/\mathscr{J})|Y'
\to\psi_*(\mathscr{O}_X)|Y'=(\psi_0)_*(\mathscr{O}_X)$라 하자.
그러면 $f_0=(\psi_0,\theta_0)$는 준스킴의 사상 $X\to Y'$이고
(\hyperref[I.2.2.1-ko]{2.2.1}), $f=j\circ f_0$이다. 이제
$\mathscr{O}_Y$의 준연접 아이디얼층 $\mathscr{J}'$로 정의된
$Y$의 두 번째 닫힌 부분준스킴을 $Y''$
\oldpage[I]{177}
라 하고, 표준 포함 사상 $j':Y''\to Y$를 통하여 $f$가
인수분해된다고 하자. 우선 $Y''\supset\psi(X)$이어야 하므로,
$Y''$가 닫혀 있다는 데서 $Y'\subset Y''$가 따른다. 또 모든
$y\in Y''$에 대하여 $\theta$는
$\mathscr{O}_y\to\mathscr{O}_y/\mathscr{J}'_y
\to(\psi_*(\mathscr{O}_X))_y$로 인수분해되어야 한다.
따라서 정의에 의해 $\mathscr{J}'_y\subset\mathscr{J}_y$이고,
그러므로 $Y'$은 $Y''$의 닫힌 부분준스킴이다
(\hyperref[I.4.1.10-ko]{4.1.10}).
\end{proof}

\begin{definition}[9.5.3]
\label{I.9.5.3-ko}
표준 포함 사상 $j:Y'\to Y$를 통하여 $f$가 인수분해되는 $Y$의
가장 작은 닫힌 부분준스킴 $Y'$이 존재할 때, $Y'$을 사상 $f$에
의한 $X$의 닫힌 상 준스킴이라고 한다.
\end{definition}

\begin{proposition}[9.5.4]
\label{I.9.5.4-ko}
$f_*(\mathscr{O}_X)$가 준연접 $\mathscr{O}_Y$-가군이면, $f$에 의한
$X$의 닫힌 상의 바탕공간은 $Y$에서 집합론적 상의 폐포
$\overline{f(X)}$이다.
\end{proposition}

\begin{proof}
$f_*(\mathscr{O}_X)$의 지지집합은 $\overline{f(X)}$에 포함되므로,
(\hyperref[I.9.5.2-ko]{9.5.2})의 표기에서
$y\notin\overline{f(X)}$이면 $\mathscr{J}_y=\mathscr{O}_y$이다.
따라서 $\mathscr{O}_Y/\mathscr{J}$의 지지집합은
$\overline{f(X)}$에 포함된다. 한편 이 지지집합은 닫혀 있고
$f(X)$를 포함한다. 실제로 $y\in f(X)$이면 환
$(\psi_*(\mathscr{O}_X))_y$의 단위원은 영이 아닌데, 이는 단면
\[
  1\in\Gamma(X,\mathscr{O}_X)=\Gamma(Y,\psi_*(\mathscr{O}_X))~;
\]
의 $y$에서의 싹이기 때문이다. 이 단위원은 $\mathscr{O}_y$의
단위원의 $\theta$에 의한 상이므로, 후자는 $\mathscr{J}_y$에 속하지
않는다. 따라서 $\mathscr{O}_y/\mathscr{J}_y\neq0$이고, 이로써 증명이
끝난다.
\end{proof}

\begin{proposition}[9.5.5]
\label{I.9.5.5-ko}
\emph{(닫힌 상의 추이성).} $f:X\to Y$와 $g:Y\to Z$를 준스킴의
두 사상이라 하자. $f$에 의한 $X$의 닫힌 상 $Y'$이 존재하고,
$g'$이 $g$를 $Y'$에 제한한 사상일 때 $g'$에 의한 $Y'$의 닫힌 상
$Z'$이 존재한다고 가정하자. 그러면 $g\circ f$에 의한 $X$의 닫힌
상이 존재하며 $Z'$과 같다.
\end{proposition}

\begin{proof}
(\hyperref[I.9.5.1-ko]{9.5.1})에 따라, $Z'$이 다음 조건을 만족하는
$Z$의 가장 작은 닫힌 부분준스킴 $Z_1$임을 보이면 충분하다. 그
조건은 $X$의 닫힌 부분준스킴 $(g\circ f)^{-1}(Z_1)$이 $X$와 같다는
것이며, 이 부분준스킴은 (\hyperref[I.4.4.2-ko]{4.4.2})에 따라
$f^{-1}(g^{-1}(Z_1))$과 같다. 이는 $Z'$이 $f$가 포함 사상
$g^{-1}(Z_1)\to Y$를 통하여 인수분해되게 하는 $Z$의 가장 작은
닫힌 부분준스킴이라는 것과 동치이다
(\hyperref[I.4.4.1-ko]{4.4.1}). 그런데 닫힌 상 $Y'$의 존재에 따라,
이 조건을 만족하는 모든 $Z_1$에 대하여 $g^{-1}(Z_1)$은 $Y'$을
부분준스킴으로 포함한다. $j$를 포함 사상 $Y'\to Y$라 하면, 이는
$j^{-1}(g^{-1}(Z_1))=g'^{-1}(Z_1)=Y'$이라는 것과 동치이다. 이제
$Z'$의 정의로부터 $Z'$이 앞의 조건을 만족하는 $Z$의 가장 작은
닫힌 부분준스킴임이 따른다.
\end{proof}

\begin{corollary}[9.5.6]
\label{I.9.5.6-ko}
$f:X\to Y$를 $f$에 의한 $X$의 닫힌 상이 $Y$인 $S$-사상이라
하자. $Z$를 $S$-스킴이라 하자. $Y$에서 $Z$로 가는 두 $S$-사상
$g_1$, $g_2$가 $g_1\circ f=g_2\circ f$를 만족하면
$g_1=g_2$이다.
\end{corollary}

\begin{proof}
$h=(g_1,g_2)_S:Y\to Z\times_S Z$라 하자. 대각선
$T=\Delta_Z(Z)$은 $Z\times_S Z$의 닫힌 부분준스킴이므로,
$Y'=h^{-1}(T)$는 $Y$의 닫힌 부분준스킴이다
(\hyperref[I.4.4.1-ko]{4.4.1}). $u=g_1\circ f=g_2\circ f$라 두면,
곱의 정의에 따라 $h'=h\circ f=(u,u)_S$이고, 따라서
$h\circ f=\Delta_Z\circ u$이다. $\Delta_Z^{-1}(T)=Z$이므로
$h'^{-1}(T)=u^{-1}(Z)=X$이고, 따라서 $f^{-1}(Y')=X$이다. 이로부터
(\hyperref[I.4.4.1-ko]{4.4.1})에 따라 표준 포함 사상 $Y'\to Y$를 통하여
$f$가 인수분해되므로, 가정에 따라 $Y'=Y$이다. 따라서
(\hyperref[I.4.4.1-ko]{4.4.1})에 따라 $h$는 $\Delta_Z\circ v$로
인수분해된다. 여기서 $v$는 $Y\to Z$인 사상이며, 이로부터
$g_1=g_2=v$가 따른다.
\end{proof}

\begin{remark}[9.5.7]
\label{I.9.5.7-ko}
$X$와 $Y$가 $S$-스킴이면, 명제
(\hyperref[I.9.5.6-ko]{9.5.6})이 뜻하는 바는
\oldpage[I]{178}
$f$에 의한 $X$의 닫힌 상이 $Y$일 때 $f$가 $S$-스킴의 범주에서
\emph{에피사상}이라는 것이다(T, 1.1). 제V장에서 역으로, $f$에
의한 $X$의 닫힌 상 $Y'$이 존재하고 $f$가 $S$-스킴의 에피사상이면
반드시 $Y'=Y$임을 증명할 것이다.
\end{remark}

\begin{proposition}[9.5.8]
\label{I.9.5.8-ko}
(\hyperref[I.9.5.1-ko]{9.5.1})의 가정들이 성립하고, $Y'$을 $f$에
의한 $X$의 닫힌 상이라 하자. $Y$의 임의의 열린집합 $V$에 대하여,
$f_V:f^{-1}(V)\to V$를 $f$의 제한이라 하자. 그러면 $f_V$에 의한
$f^{-1}(V)$의 $V$ 안에서의 닫힌 상은 존재하며, $Y'$의 열린집합
$V\cap Y'$ 위에 $Y'$로부터 유도된 준스킴과 같다. 다시 말해 이는
$Y$의 부분준스킴 $\inf(V,Y')$와 같다
(\hyperref[I.4.4.3-ko]{4.4.3}).
\end{proposition}

\begin{proof}
$X'=f^{-1}(V)$라 하자. $f_V$에 의한 $\mathscr{O}_{X'}$의 직상은
$f_*(\mathscr{O}_X)$를 $V$에 제한한 것과 같으므로, 준동형
$\mathscr{O}_V\to(f_V)_*(\mathscr{O}_{X'})$의 핵
$\mathscr{J}'$은 $\mathscr{J}$를 $V$에 제한한 것이다. 이로부터
명제가 곧바로 따른다.
\end{proof}

이 결과는 닫힌 상을 만드는 것이 밑준스킴의 확장 $Y_1\to Y$가
\emph{열린 몰입 사상}인 경우 그 확장과 가환한다는 뜻으로 해석할
수 있다. 제IV장에서 $f$가 분리 사상이고 준콤팩트일 때에는
$Y_1\to Y$가 \emph{평탄} 사상인 경우에도 마찬가지임을 볼 것이다.

\begin{proposition}[9.5.9]
\label{I.9.5.9-ko}
$f:X\to Y$를 $f$에 의한 $X$의 닫힌 상 $Y'$이 존재하는 사상이라
하자.
\begin{enumerate}
  \item[(i)] $X$가 축소 준스킴이면 $Y'$도 축소 준스킴이다.
  \item[(ii)] (\hyperref[I.9.5.1-ko]{9.5.1})의 가정들이 성립하고
    $X$가 기약이면(각각 정역 준스킴이면), $Y'$도 기약이다
    (각각 정역 준스킴이다).
\end{enumerate}
\end{proposition}

\begin{proof}
가정에 따라 $f$는 $X\xrightarrow{g}Y'\xrightarrow{j}Y$로
인수분해되며, 여기서 $j$는 표준 포함 사상이다. $X$가 축소이므로
$g$는 $X\xrightarrow{h}Y'_{\mathrm{red}}\xrightarrow{j'}Y'$로
인수분해된다. 여기서 $j'$은 표준 포함 사상이다
(\hyperref[I.5.2.2-ko]{5.2.2}). 따라서 $Y'$의 정의로부터
$Y'_{\mathrm{red}}=Y'$이 따른다. 또한
(\hyperref[I.9.5.1-ko]{9.5.1})의 조건들이 성립하면,
(\hyperref[I.9.5.4-ko]{9.5.4})에 따라 $f(X)$는 $Y'$에서 조밀하다.
그러므로 $X$가 기약이면 $Y'$도 기약이다
(\hyperref[0.2.1.5-ko]{0, 2.1.5}). 정역 준스킴에 관한 주장은
앞의 두 주장을 함께 적용하여 얻는다.
\end{proof}

\begin{proposition}[9.5.10]
\label{I.9.5.10-ko}
$Y$를 준스킴 $X$의 부분준스킴이라 하고, 표준 포함 사상
$i:Y\to X$가 준콤팩트 사상이라고 하자. 그러면 $Y$를 부분준스킴으로
포함하는 $X$의 가장 작은 닫힌 부분준스킴 $\overline{Y}$이 존재한다.
그 바탕공간은 $Y$의 바탕공간의 폐포이고, 후자는 이
폐포에서 열려 있으며, 이 열린집합 위에서 $\overline{Y}$로부터
유도된 준스킴은 $Y$이다.
\end{proposition}

\begin{proof}
포함 사상 $i$에 (\hyperref[I.9.5.1-ko]{9.5.1})을 적용하면
된다.\footnote{\textbf{번역 주.} 인쇄 원문은 이 증명에서 앞서 정의되지
않은 $j$에 (9.5.1)을 적용한다고 쓰지만, 명제에서 정의된 표준 포함
사상은 $i:Y\to X$이다. 여기서는 명제에서 정의된 표기에 맞게 $i$로
바로잡았다.} 이 사상은
분리 사상이고 (\hyperref[I.5.5.1-ko]{5.5.1}), 가정에 따라
준콤팩트이다. 따라서 (\hyperref[I.9.5.1-ko]{9.5.1})은
$\overline{Y}$의 존재를 보이고, (\hyperref[I.9.5.4-ko]{9.5.4})는
그 바탕공간이 $X$에서 $Y$의 폐포임을 보인다. $Y$는 $X$에서 국소
닫혀 있으므로 $\overline{Y}$에서 열려 있다. 마지막 주장은 $Y$가
$X$의 어떤 열린집합 $V$ 안에서 닫혀 있도록 $V$를 택한 뒤
(\hyperref[I.9.5.8-ko]{9.5.8})을 적용하면 얻는다.
\end{proof}

이 표기 아래에서, 포함 사상 $V\to X$가 준콤팩트이고
$\mathscr{J}$가 $V$의 닫힌 부분준스킴 $Y$를 정의하는
$\mathscr{O}_X|V$의 준연접 아이디얼층이면,
(\hyperref[I.9.5.1-ko]{9.5.1})에 따라 $\overline{Y}$을 정의하는
$\mathscr{O}_X$의 준연접 아이디얼층은 $\mathscr{J}$의 표준 연장
(\hyperref[I.9.4.1-ko]{9.4.1}) $\overline{\mathscr{J}}$이다. 실제로
이는 $V$ 위에 $\mathscr{J}$를 유도하는 $\mathscr{O}_X$의 준연접
아이디얼 부분층 가운데 가장 큰 것이다.
\oldpage[I]{179}
\begin{corollary}[9.5.11]
\label{I.9.5.11-ko}
(\hyperref[I.9.5.10-ko]{9.5.10})의 가정 아래에서,
$\overline{Y}$의 열린집합 $V$ 위의 $\mathscr{O}_{\overline{Y}}$의
단면이 $V\cap Y$에서 영이면 그 단면은 영이다.
\end{corollary}

\begin{proof}
(\hyperref[I.9.5.8-ko]{9.5.8})에 따라 $V=\overline{Y}$인 경우로
환원할 수 있다. $\overline{Y}$ 위의 $\mathscr{O}_{\overline{Y}}$의
단면들이 준스킴
$\overline{Y}\otimes_{\mathbb{Z}}\mathbb{Z}[T]$의
$\overline{Y}$-단면들과 표준적으로 대응하고
(\hyperref[I.3.3.15-ko]{3.3.15}), 이 준스킴은 $\overline{Y}$ 위에서
분리되어 있으므로, 이 따름정리는
(\hyperref[I.9.5.6-ko]{9.5.6})의 특별한 경우이다.
\end{proof}

$X$의 부분준스킴 $Y$를 부분준스킴으로 포함하는 $X$의 가장 작은 닫힌
부분준스킴 $\overline{Y}$이 존재할 때, 혼동할 우려가 없으면
$\overline{Y}$을 $X$에서 $Y$의 \emph{폐포}라고 한다.

\subsection{준연접 대수; 구조층의 변경.}
\label{subsection:I.9.6-ko}

\begin{proposition}[9.6.1]
\label{I.9.6.1-ko}
$X$를 준스킴, $\mathscr{B}$를 준연접 $\mathscr{O}_X$-대수라 하자
(\hyperref[0.5.1.3-ko]{0, 5.1.3}). $\mathscr{B}$-가군 $\mathscr{F}$가
(환 달린 공간 $(X,\mathscr{B})$ 위에서) 준연접일 필요충분조건은
$\mathscr{F}$가 준연접 $\mathscr{O}_X$-가군인 것이다.
\end{proposition}

\begin{proof}
문제가 국소적이므로 $X$가 환 $A$를 갖는 아핀 스킴이라고 해도 된다.
그러면 $\mathscr{B}=\widetilde{B}$이고, 여기서 $B$는 $A$-대수이다
(\hyperref[I.1.4.3-ko]{1.4.3}). $\mathscr{F}$가 환 달린 공간
$(X,\mathscr{B})$ 위에서 준연접이면, $\mathscr{F}$가 어떤
$\mathscr{B}$-준동형 $\mathscr{B}^{(I)}\to\mathscr{B}^{(J)}$의
여핵이라고 해도 된다. 이 준동형은 $\mathscr{O}_X$-가군의
$\mathscr{O}_X$-준동형이기도 하고, $\mathscr{B}^{(I)}$와
$\mathscr{B}^{(J)}$는 준연접 $\mathscr{O}_X$-가군이므로
(\hyperref[I.1.3.9-ko]{1.3.9, (ii)}), $\mathscr{F}$도 준연접
$\mathscr{O}_X$-가군이다
(\hyperref[I.1.3.9-ko]{1.3.9, (i)}).

거꾸로 $\mathscr{F}$가 준연접 $\mathscr{O}_X$-가군이면
$\mathscr{F}=\widetilde{M}$이고, 여기서 $M$은 $B$-가군이다
(\hyperref[I.1.4.3-ko]{1.4.3}). $M$은 어떤 $B$-준동형
$B^{(I)}\to B^{(J)}$의 여핵과 동형이므로, $\mathscr{F}$는 이에
대응하는 준동형 $\mathscr{B}^{(I)}\to\mathscr{B}^{(J)}$의 여핵과
동형인 $\mathscr{B}$-가군이다
(\hyperref[I.1.3.13-ko]{1.3.13}). 이로써 증명이 끝난다.
\end{proof}

특히 $\mathscr{F}$와 $\mathscr{G}$가 두 준연접
$\mathscr{B}$-가군이면 $\mathscr{F}\otimes_{\mathscr{B}}\mathscr{G}$는
준연접 $\mathscr{B}$-가군이다. 더욱이 $\mathscr{F}$가 유한 표시를
갖는다고 가정하면 $\mathscr{H}\!om_{\mathscr{B}}(\mathscr{F},\mathscr{G})$도
그러하다(\hyperref[I.1.3.13-ko]{1.3.13}).

\begin{env}[9.6.2]
\label{I.9.6.2-ko}
준스킴 $X$가 주어졌을 때, 준연접 $\mathscr{O}_X$-대수
$\mathscr{B}$가 \emph{유한 생성}이라는 것은 모든 $x\in X$에 대하여
$x$의 \emph{아핀} 열린 근방 $U$가 존재하여
$\Gamma(U,\mathscr{B})=B$가 $\Gamma(U,\mathscr{O}_X)=A$ 위의 유한 생성
대수라는 뜻이다. 이때 $\mathscr{B}|U=\widetilde{B}$이고, 모든 $f\in A$에
대하여 그 원소에 대응하는 주 열린집합 위에 유도된
$(\mathscr{O}_X|D(f))$-대수
$\mathscr{B}|D(f)$도 유한 생성이다. 실제로 이는
$\widetilde{B_f}$와 동형이고, $B_f=B\otimes_A A_f$는 분명히 $A_f$ 위의
유한 생성 대수이다. $D(f)$들이 $U$의 위상의 기저를 이루므로, 준연접
유한 생성 $\mathscr{O}_X$-대수 $\mathscr{B}$와 $X$의 모든 열린집합
$V$에 대하여 $\mathscr{B}|V$는 준연접 유한 생성
$(\mathscr{O}_X|V)$-대수이다.
\end{env}

\begin{proposition}[9.6.3]
\label{I.9.6.3-ko}
$X$를 국소 뇌터 준스킴이라 하자. 모든 준연접 유한 생성
$\mathscr{O}_X$-대수 $\mathscr{B}$는 연접인 환의 층이다
(\hyperref[0.5.3.7-ko]{0, 5.3.7}).
\end{proposition}

\begin{proof}
다시 $X$가 뇌터 환 $A$를 갖는 아핀 스킴이고
$\mathscr{B}=\widetilde{B}$이며 $B$가 $A$ 위의 유한 생성 대수인 경우로
한정해도 된다. 그러면 $B$는 뇌터 환이다. 따라서
$\mathscr{B}$-준동형 $\mathscr{B}^m\to\mathscr{B}$의 핵
$\mathscr{N}$이 유한 생성 $\mathscr{B}$-
\oldpage[I]{180}
가군임을 보이면 된다. 그런데 이 핵은 이에 대응하는 $B$-가군의
준동형 $B^m\to B$의 핵 $N$에 대응하는 층 $\widetilde{N}$과
($\mathscr{B}$-가군으로서) 동형이다
(\hyperref[I.1.3.13-ko]{1.3.13}). $B$가 뇌터이므로 $B^m$의 부분가군
$N$은 유한 생성 $B$-가군이다. 따라서 상이 $N$인 준동형
$B^p\to B^m$이 존재한다. 열 $B^p\to B^m\to B$이 완전이므로 이에
대응하는 열 $\mathscr{B}^p\to\mathscr{B}^m\to\mathscr{B}$도 완전이고
(\hyperref[I.1.3.5-ko]{1.3.5}), $\mathscr{N}$은
$\mathscr{B}^p\to\mathscr{B}^m$의 상이므로
(\hyperref[I.1.3.9-ko]{1.3.9, (i)}) 명제가 증명된다.
\end{proof}

\begin{corollary}[9.6.4]
\label{I.9.6.4-ko}
(\hyperref[I.9.6.3-ko]{9.6.3})의 가정 아래에서, $\mathscr{B}$-가군
$\mathscr{F}$가 연접일 필요충분조건은 그 가군이 준연접
$\mathscr{O}_X$-가군이면서 유한 생성 $\mathscr{B}$-가군인 것이다.
이 조건들이 성립하고 $\mathscr{G}$가 $\mathscr{F}$의
부분-$\mathscr{B}$-가군 또는 $\mathscr{B}$-몫가군이면,
$\mathscr{G}$가 연접 $\mathscr{B}$-가군일 필요충분조건은
$\mathscr{G}$가 준연접 $\mathscr{O}_X$-가군인 것이다.
\end{corollary}

\begin{proof}
(\hyperref[I.9.6.1-ko]{9.6.1})을 고려하면 $\mathscr{F}$에 관한 조건들이
필요함은 분명하다. 충분함을 보이자. $X$가 뇌터 환 $A$를 갖는
아핀 스킴이고, $\mathscr{B}=\widetilde{B}$이며, $B$가 $A$ 위의 유한 생성
대수이고, $\mathscr{F}=\widetilde{M}$이며, 여기서 $M$은 $B$-가군이고,
전사 $\mathscr{B}$-준동형 $\mathscr{B}^m\to\mathscr{F}\to 0$이 존재하는
경우로 한정해도 된다. 그러면 이에 대응하는 완전열
$B^m\to M\to 0$이 있으므로 $M$은 유한 생성 $B$-가군이다. 또한
$B^m\to M$의 핵 $P$도 $B$가 뇌터이므로 유한 생성 $B$-가군이다.
따라서 (\hyperref[I.1.3.13-ko]{1.3.13})에 의해 $\mathscr{F}$는 어떤
$\mathscr{B}$-준동형 $\mathscr{B}^n\to\mathscr{B}^m$의 여핵이고,
$\mathscr{B}$가 연접인 환의 층이므로 연접이다
(\hyperref[0.5.3.4-ko]{0, 5.3.4}). 같은 논증으로 $\mathscr{F}$의 준연접
부분-$\mathscr{B}$-가군(각각 $\mathscr{B}$-몫가군)은 유한 생성임을
알 수 있으며, 이로부터 따름정리의 두 번째 주장이 따른다.
\end{proof}

\begin{proposition}[9.6.5]
\label{I.9.6.5-ko}
$X$를 준콤팩트 스킴 또는 바탕공간이 뇌터인 준스킴이라 하자.
모든 준연접 유한 생성 $\mathscr{O}_X$-대수 $\mathscr{B}$에 대하여,
$\mathscr{B}$의 유한 생성 준연접 부분-$\mathscr{O}_X$-가군
$\mathscr{E}$가 존재하여 $\mathscr{E}$가 $\mathscr{O}_X$-대수
$\mathscr{B}$를 생성한다(\hyperref[0.4.1.4-ko]{0, 4.1.4}).
\end{proposition}

\begin{proof}
실제로 가정에 따라 $X$는 아핀 열린집합 $(U_i)$의 유한 덮개를 가지며,
$\Gamma(U_i,\mathscr{B})=B_i$는 $\Gamma(U_i,\mathscr{O}_X)=A_i$ 위의
유한 생성 대수이다. $B_i$의 유한 생성 부분-$A_i$-가군 $E_i$를
$A_i$-대수 $B_i$를 생성하도록 잡자. (\hyperref[I.9.4.7-ko]{9.4.7})에
따르면, $\mathscr{B}$의 유한 생성 준연접 부분-$\mathscr{O}_X$-가군
$\mathscr{E}_i$가 존재하여 $\mathscr{E}_i|U_i=\widetilde{E_i}$이다.
$\mathscr{E}_i$들의 합 $\mathscr{E}$가 요구된 조건을 만족함은 분명하다.
\end{proof}

\begin{proposition}[9.6.6]
\label{I.9.6.6-ko}
$X$를 바탕공간이 국소 뇌터인 준스킴 또는 준콤팩트 스킴이라 하자.
그러면 모든 준연접 $\mathscr{O}_X$-대수 $\mathscr{B}$는 그 준연접 유한
생성 부분-$\mathscr{O}_X$-대수들의 귀납극한이다.
\end{proposition}

\begin{proof}
(\hyperref[I.9.4.9-ko]{9.4.9})에 따라 $\mathscr{B}$는
$\mathscr{O}_X$-가군으로서 그 유한 생성 준연접
부분-$\mathscr{O}_X$-가군들의 귀납극한이다. 이 부분가군들은
$\mathscr{B}$의 준연접 유한 생성 부분-$\mathscr{O}_X$-대수들을 생성하고
(\hyperref[I.1.3.14-ko]{1.3.14}), 따라서 $\mathscr{B}$는
\emph{특히} 이 부분대수들의 귀납극한이다.
\end{proof}

\section{형식적 스킴}
\label{section:I.10-ko}

\subsection{아핀 형식적 스킴.}
\label{subsection:I.10.1-ko}

\begin{env}[10.1.1]
\label{I.10.1.1-ko}
$A$를 위상환인 \emph{허용환}이라 하자
(\hyperref[0.7.1.2-ko]{0, 7.1.2}). $A$의 모든 정의 아이디얼
$\mathfrak{J}$에 대하여 $\operatorname{Spec}(A/\mathfrak{J})$는
$\operatorname{Spec}(A)$의 닫힌 부분공간 $V(\mathfrak{J})$와
동일시된다(\hyperref[I.1.1.11-ko]{1.1.11}). 이는 $A$의
\emph{열린} 소 아이디얼들의 집합이다. 이 위상공간은 고려한
\oldpage[I]{181}
정의 아이디얼 $\mathfrak{J}$의 선택과 무관하며, 이를
$\mathfrak{X}$로 나타내자. $(\mathfrak{J}_\lambda)$를 정의
아이디얼들로 이루어진 $A$에서 0의 근방들의 기본계라 하고, 모든
$\lambda$에 대하여 $\mathscr{O}_\lambda$를
$\operatorname{Spec}(A/\mathfrak{J}_\lambda)$의 구조층이라 하자. 이 층은
$\widetilde{A}/\widetilde{\mathfrak{J}}_\lambda$에 의해
$\mathfrak{X}$ 위에 유도된다(이 몫층은 $\mathfrak{X}$ 밖에서 영이다).

$\mathfrak{J}_\mu\subset\mathfrak{J}_\lambda$이면 표준 준동형
$A/\mathfrak{J}_\mu\to A/\mathfrak{J}_\lambda$는 환의 층의 준동형
$u_{\lambda\mu}:\mathscr{O}_\mu\to\mathscr{O}_\lambda$를 정의하므로
(\hyperref[I.1.6.1-ko]{1.6.1}), $(\mathscr{O}_\lambda)$는 이 준동형들에
관하여 \emph{환의 층들의 사영계}를 이룬다. $\mathfrak{X}$의 위상에는
준콤팩트 열린집합들로 이루어진 기저가 있으므로, 각
$\mathscr{O}_\lambda$에 그 바탕 환의 층(위상을 잊은 층)이
$\mathscr{O}_\lambda$인 \emph{의사이산 위상환의 층}을 대응시킬 수
있다(\hyperref[0.3.8.1-ko]{0, 3.8.1}). 이 층도 여전히
$\mathscr{O}_\lambda$로 나타내겠다. 또한 $\mathscr{O}_\lambda$들은
여전히 \emph{위상환의 층들의 사영계}를 이룬다
(\hyperref[0.3.8.2-ko]{0, 3.8.2}). 이 계
$(\mathscr{O}_\lambda)$의 사영극한인 $\mathfrak{X}$ 위의
\emph{위상환의 층}을 $\mathscr{O}_{\mathfrak{X}}$로 나타내겠다. 따라서
$\mathfrak{X}$의 모든 \emph{준콤팩트} 열린집합 $U$에 대하여
$\Gamma(U,\mathscr{O}_{\mathfrak{X}})$는 \emph{이산}환들의 계
$\Gamma(U,\mathscr{O}_\lambda)$의 사영극한인 위상환이다
(\hyperref[0.3.2.6-ko]{0, 3.2.6}).
\end{env}

\begin{definition}[10.1.2]
\label{I.10.1.2-ko}
허용환 $A$가 주어졌을 때, $A$의 열린 소 아이디얼들로 이루어진
$\operatorname{Spec}(A)$의 닫힌 부분공간 $\mathfrak{X}$를 $A$의
형식적 스펙트럼이라 하고 $\operatorname{Spf}(A)$로 나타낸다. 위상환 달린
공간이 형식적 스펙트럼 $\operatorname{Spf}(A)=\mathfrak{X}$에 다음의
구조층을 갖춘 것과 동형이면 이를 아핀 형식적 스킴이라고 한다. 여기서
$\mathfrak{J}_\lambda$가 $A$의 정의 아이디얼들의 여과집합을 달릴 때
$\mathscr{O}_{\mathfrak{X}}$는
의사이산 위상환의 층
$(\widetilde{A}/\widetilde{\mathfrak{J}}_\lambda)|\mathfrak{X}$들의
사영극한이다.
\end{definition}

\emph{형식적 스펙트럼} $\mathfrak{X}=\operatorname{Spf}(A)$를 아핀 형식적
스킴으로 말할 때에는 언제나 위에서 정의한
$\mathscr{O}_{\mathfrak{X}}$를 갖는 위상환 달린 공간
$(\mathfrak{X},\mathscr{O}_{\mathfrak{X}})$을 뜻한다.

모든 \emph{아핀 스킴} $X=\operatorname{Spec}(A)$는 $A$를 이산
위상환으로 보아 유일한 방식으로 아핀 형식적 스킴으로 볼 수 있다. 이때
$U$가 준콤팩트이면 위상환 $\Gamma(U,\mathscr{O}_X)$는 이산이지만,
$U$가 $X$의 임의의 열린집합일 때에는 일반적으로 그렇지 않다.

\begin{proposition}[10.1.3]
\label{I.10.1.3-ko}
$A$가 허용환이고 $\mathfrak{X}=\operatorname{Spf}(A)$이면,
$\Gamma(\mathfrak{X},\mathscr{O}_{\mathfrak{X}})$는 $A$와 위상적으로
동형이다.
\end{proposition}

\begin{proof}
실제로 $\mathfrak{X}$는 $\operatorname{Spec}(A)$에서 닫혀 있으므로
준콤팩트이다. 따라서
$\Gamma(\mathfrak{X},\mathscr{O}_{\mathfrak{X}})$는 이산환
$\Gamma(\mathfrak{X},\mathscr{O}_\lambda)$들의 사영극한과 위상적으로
동형이다. 그런데 $\Gamma(\mathfrak{X},\mathscr{O}_\lambda)$는
$A/\mathfrak{J}_\lambda$와 동형이다
(\hyperref[I.1.3.7-ko]{1.3.7}). $A$는 분리이고 완비이므로
$\varprojlim A/\mathfrak{J}_\lambda$와 위상적으로 동형이다
(\hyperref[0.7.2.1-ko]{0, 7.2.1}). 이로부터 명제가 따른다.
\end{proof}

\begin{proposition}[10.1.4]
\label{I.10.1.4-ko}
$A$를 허용환, $\mathfrak{X}=\operatorname{Spf}(A)$라 하고, 모든
$f\in A$에 대하여 $\mathfrak{D}(f)=D(f)\cap\mathfrak{X}$라 하자. 그러면
위상환 달린 공간
$(\mathfrak{D}(f),\mathscr{O}_{\mathfrak{X}}|\mathfrak{D}(f))$는 아핀
형식적 스펙트럼 $\operatorname{Spf}(A_{\{f\}})$와 동형이다
(\hyperref[0.7.6.15-ko]{0, 7.6.15}).
\end{proposition}

\begin{proof}
$A$의 모든 정의 아이디얼 $\mathfrak{J}$에 대하여 이산환
$S_f^{-1}A/S_f^{-1}\mathfrak{J}$는
$A_{\{f\}}/\mathfrak{J}_{\{f\}}$와 표준적으로 동일시된다
(\hyperref[0.7.6.9-ko]{0, 7.6.9}). 따라서
(\hyperref[I.1.2.5-ko]{1.2.5}와 \hyperref[I.1.2.6-ko]{1.2.6})에 따라 위상공간
$\operatorname{Spf}(A_{\{f\}})$는 $\mathfrak{D}(f)$와 표준적으로
동일시된다. 더 나아가 $\mathfrak{D}(f)$에 포함되는
$\mathfrak{X}$의 모든 준콤팩트 열린집합 $U$에 대하여
$\Gamma(U,\mathscr{O}_\lambda)$는
$\operatorname{Spec}(S_f^{-1}A/S_f^{-1}\mathfrak{J}_\lambda)$의 구조층의
$U$ 위의 단면가군과 동일시된다
(\hyperref[I.1.3.6-ko]{1.3.6}). 따라서
$\mathfrak{Y}=\operatorname{Spf}(A_{\{f\}})$로 놓으면
$\Gamma(U,\mathscr{O}_{\mathfrak{X}})$는 단면가군
$\Gamma(U,\mathscr{O}_{\mathfrak{Y}})$와 동일시된다. 이로써 명제가
증명된다.
\end{proof}

\oldpage[I]{182}
\begin{env}[10.1.5]
\label{I.10.1.5-ko}
\emph{위상을 잊은} 환의 층으로서 $\operatorname{Spf}(A)$의 구조층
$\mathscr{O}_{\mathfrak{X}}$는 모든 $x\in\mathfrak{X}$에서
(\hyperref[I.10.1.4-ko]{10.1.4})에 따라 귀납극한
$\varinjlim_{f\notin\mathfrak{j}_x}A_{\{f\}}$와 동일시되는 줄기를 갖는다.
따라서 (\hyperref[0.7.6.17-ko]{0, 7.6.17}과
\hyperref[0.7.6.18-ko]{7.6.18}):
\end{env}

\begin{proposition}[10.1.6]
\label{I.10.1.6-ko}
모든 $x\in\mathfrak{X}=\operatorname{Spf}(A)$에 대하여 줄기
$\mathscr{O}_x$는 국소환이고, 그 잉여체는
$k(x)=A_x/\mathfrak{j}_xA_x$와 동형이다. 더 나아가 $A$가 아딕이고
뇌터이면 $\mathscr{O}_x$는 뇌터 환이다.
\end{proposition}

$k(x)$가 영환이 아니므로, 이 결과로부터 환의 층
$\mathscr{O}_{\mathfrak{X}}$의 \emph{지지집합}은
\emph{$\mathfrak{X}$와 같다}는 결론을 얻는다.

\subsection{아핀 형식적 스킴의 사상.}
\label{subsection:I.10.2-ko}

\begin{env}[10.2.1]
\label{I.10.2.1-ko}
$A$, $B$를 두 허용환이라 하고,
$\varphi:B\to A$를 \emph{연속} 준동형이라 하자. 그러면 연속사상
$ {}^a\varphi:\operatorname{Spec}(A)\to\operatorname{Spec}(B)$
(\hyperref[I.1.2.1-ko]{1.2.1})은
$\mathfrak{X}=\operatorname{Spf}(A)$를
$\mathfrak{Y}=\operatorname{Spf}(B)$ 안으로 보낸다. 실제로 $A$의 열린
소 아이디얼의 $\varphi$에 의한 역상은 $B$의 열린 소 아이디얼이다.
한편 모든 $g\in B$에 대하여 $\varphi$는 연속 준동형
\[
  \Gamma(\mathfrak{D}(g),\mathscr{O}_{\mathfrak{Y}})
  \longrightarrow
  \Gamma(\mathfrak{D}(\varphi(g)),\mathscr{O}_{\mathfrak{X}})
\]
을 정의한다. 이는 (\hyperref[I.10.1.4-ko]{10.1.4}),
(\hyperref[I.10.1.3-ko]{10.1.3}),
(\hyperref[0.7.6.7-ko]{0, 7.6.7})에 따른다. 이 준동형들은 $g$를 $g$의
배수로 바꿀 때의 제한사상들과 양립하고,
$\mathfrak{D}(\varphi(g))={}^a\varphi^{-1}(\mathfrak{D}(g))$이므로,
위상환의 층 사이의 \emph{연속} 준동형
$\mathscr{O}_{\mathfrak{Y}}\to
{}^a\varphi_*(\mathscr{O}_{\mathfrak{X}})$
(\hyperref[0.3.2.5-ko]{0, 3.2.5})을 정의한다. 이를 다시
$\widetilde{\varphi}$로 나타내겠다. 이렇게 하여 위상환 달린 공간의 사상
$\Phi=({}^a\varphi,\widetilde{\varphi})$를
$\mathfrak{X}\to\mathfrak{Y}$로 얻는다. 위상을 잊은 환의 층의
준동형으로서 $\widetilde{\varphi}$는 모든 $x\in\mathfrak{X}$에서 줄기
사이의 준동형
$\widetilde{\varphi}_x^\sharp:
\mathscr{O}_{{}^a\varphi(x)}\to\mathscr{O}_x$를 정의한다.
\end{env}

\begin{proposition}[10.2.2]
\label{I.10.2.2-ko}
$A$, $B$를 두 허용환이라 하고,
$\mathfrak{X}=\operatorname{Spf}(A)$,
$\mathfrak{Y}=\operatorname{Spf}(B)$라 하자. 위상환 달린 공간의 사상
$u=(\psi,\theta):\mathfrak{X}\to\mathfrak{Y}$가 $B\to A$인 어떤 연속 환 준동형
$\varphi$로부터 얻어지는
$({}^a\varphi,\widetilde{\varphi})$ 꼴이기 위한 필요충분조건은 모든
$x\in\mathfrak{X}$에 대하여 $\theta_x^\sharp$가 국소 준동형
$\mathscr{O}_{\psi(x)}\to\mathscr{O}_x$인 것이다.
\end{proposition}

\begin{proof}
이 조건은 필요하다. 실제로
$\mathfrak{p}=\mathfrak{j}_x\in\operatorname{Spf}(A)$,
$\mathfrak{q}=\varphi^{-1}(\mathfrak{j}_x)$라 하자.
$g\notin\mathfrak{q}$이면 $\varphi(g)\notin\mathfrak{p}$이고,
$\varphi$에서 얻어지는 준동형
$B_{\{g\}}\to A_{\{\varphi(g)\}}$
(\hyperref[0.7.6.7-ko]{0, 7.6.7})은
$\mathfrak{q}_{\{g\}}$를
$\mathfrak{p}_{\{\varphi(g)\}}$의 부분집합으로 보낸다. 따라서 귀납극한을
취하고 (\hyperref[I.10.1.5-ko]{10.1.5})와
(\hyperref[0.7.6.17-ko]{0, 7.6.17})을 고려하면
$\widetilde{\varphi}_x^\sharp$가 국소 준동형임을 알 수 있다.

역으로 $(\psi,\theta)$가 명제의 조건을 만족하는 사상이라고 하자.
(\hyperref[I.10.1.3-ko]{10.1.3})에 따라 $\theta$는 다음 연속 환 준동형을
정의한다:
\[
  \varphi=\Gamma(\theta):
  B=\Gamma(\mathfrak{Y},\mathscr{O}_{\mathfrak{Y}})
  \longrightarrow
  \Gamma(\mathfrak{X},\mathscr{O}_{\mathfrak{X}})=A.
\]
$\theta$에 대한 가정에 의하여 $\mathfrak{X}$ 위의
$\mathscr{O}_{\mathfrak{X}}$의 단면 $\varphi(g)$의 $x$에서의 싹이
가역일 필요충분조건은 $g$의 $\psi(x)$에서의 싹이 가역인 것이다. 그런데
(\hyperref[0.7.6.17-ko]{0, 7.6.17})에 따라
$\mathfrak{X}$ 위의 $\mathscr{O}_{\mathfrak{X}}$의 단면(각각
$\mathfrak{Y}$ 위의 $\mathscr{O}_{\mathfrak{Y}}$의 단면) 가운데
$x$에서의 싹(각각 $\psi(x)$에서의 싹)이 가역이 아닌 것들은 정확히
$\mathfrak{j}_x$
\oldpage[I]{183}
(각각 $\mathfrak{j}_{\psi(x)}$)의 원소들이다. 따라서 앞의 관찰로부터
${}^a\varphi=\psi$를 얻는다. 마지막으로 모든 $g\in B$에 대하여 도식
\[
  \xymatrix{
    B=\Gamma(\mathfrak{Y},\mathscr{O}_{\mathfrak{Y}})
      \ar[r]^\varphi\ar[d] &
    \Gamma(\mathfrak{X},\mathscr{O}_{\mathfrak{X}})=A\ar[d]\\
    B_{\{g\}}=\Gamma(\mathfrak{D}(g),\mathscr{O}_{\mathfrak{Y}})
      \ar[r]_{\Gamma(\theta_{\mathfrak{D}(g)})} &
    \Gamma(\mathfrak{D}(\varphi(g)),\mathscr{O}_{\mathfrak{X}})
      =A_{\{\varphi(g)\}}
  }
\]
은 가환이다. 완비 분수환의 보편 성질
(\hyperref[0.7.6.6-ko]{0, 7.6.6})에 따라 모든 $g\in B$에 대하여
$\theta_{\mathfrak{D}(g)}$는
$\widetilde{\varphi}_{\mathfrak{D}(g)}$와 같고, 따라서
(\hyperref[0.3.2.5-ko]{0, 3.2.5})에 따라
$\theta=\widetilde{\varphi}$이다.
\end{proof}

(\hyperref[I.10.2.2-ko]{10.2.2})의 조건을 만족하는 위상환 달린 공간의
사상 $(\psi,\theta)$를 \emph{아핀 형식적 스킴의 사상}이라 하겠다.
$A$에 $\operatorname{Spf}(A)$를 대응시키는 함자와 $\mathfrak{X}$에
$\Gamma(\mathfrak{X},\mathscr{O}_{\mathfrak{X}})$를 대응시키는 함자는
허용환의 범주와 아핀 형식적 스킴의 범주의 반대 범주 사이의
\emph{동치}를 정의한다고 할 수 있다(T, I, 1.2).

\begin{env}[10.2.3]
\label{I.10.2.3-ko}
(\hyperref[I.10.2.2-ko]{10.2.2})의 특수한 경우로서, $f\in A$에 대하여
$\mathfrak{X}$가 $\mathfrak{D}(f)$ 위에 유도하는 아핀 형식적 스킴의
표준 포함 사상은 표준 연속 준동형 $A\to A_{\{f\}}$에 대응한다. 이제
(\hyperref[I.10.2.2-ko]{10.2.2})의 가정 아래에서 $h$를 $B$의 원소,
$g$를 $\varphi(h)$의 배수인 $A$의 원소라 하자. 그러면
$\psi(\mathfrak{D}(g))\subset\mathfrak{D}(h)$이다. $u$를
$\mathfrak{D}(g)$에 제한한 사상을 $\mathfrak{D}(g)$에서
$\mathfrak{D}(h)$로 가는 사상으로 보면, 이는 도식
\[
  \xymatrix{
    \mathfrak{D}(g)\ar[r]^v\ar[d] &
    \mathfrak{D}(h)\ar[d]\\
    \mathfrak{X}\ar[r]_u &
    \mathfrak{Y}
  }
\]
을 가환이게 하는 유일한 사상 $v$이다. 이 사상은 도식
\[
  \xymatrix{
    A\ar[d] &
    B\ar[l]_\varphi\ar[d]\\
    A_{\{g\}} &
    B_{\{h\}}\ar[l]_{\varphi'}
  }
\]
을 가환이게 하는 유일한 연속 준동형
$\varphi':B_{\{h\}}\to A_{\{g\}}$
(\hyperref[0.7.6.7-ko]{0, 7.6.7})에 대응한다.
\end{env}

\subsection{아핀 형식적 스킴의 정의 아이디얼.}
\label{subsection:I.10.3-ko}

\begin{env}[10.3.1]
\label{I.10.3.1-ko}
$A$를 허용환, $\mathfrak{J}$를 $A$의 열린 아이디얼,
$\mathfrak{X}$를 아핀 형식적 스킴 $\operatorname{Spf}(A)$라 하자.
$(\mathfrak{J}_\lambda)$를 $\mathfrak{J}$에 포함되는 $A$의 정의
아이디얼들의 집합이라 하자. 그러면
$\widetilde{\mathfrak{J}}/\widetilde{\mathfrak{J}}_\lambda$는
$\widetilde{A}/\widetilde{\mathfrak{J}}_\lambda$의 아이디얼층이다.
$\mathfrak{X}$ 위에
$\widetilde{\mathfrak{J}}/\widetilde{\mathfrak{J}}_\lambda$가 유도하는
층들의 사영극한을 $\mathfrak{J}^\Delta$로 나타내자. 이는
$\mathscr{O}_{\mathfrak{X}}$의 \emph{아이디얼층}과 동일시된다
(\hyperref[0.3.2.6-ko]{0, 3.2.6}). 모든 $f\in A$에 대하여
$\Gamma(\mathfrak{D}(f),\mathfrak{J}^\Delta)$는
$S_f^{-1}\mathfrak{J}/S_f^{-1}\mathfrak{J}_\lambda$들의 사영극한, 즉
환 $A_{\{f\}}$의 열린 아이디얼 $\mathfrak{J}_{\{f\}}$와 동일시된다
(\hyperref[0.7.6.9-ko]{0, 7.6.9}). 특히
$\Gamma(\mathfrak{X},\mathfrak{J}^\Delta)=\mathfrak{J}$이다.
$\mathfrak{D}(f)$들이 $\mathfrak{X}$의 위상의 기저를 이루므로 이로부터
다음을 얻는다:
\[
  \mathfrak{J}^\Delta|\mathfrak{D}(f)
  =(\mathfrak{J}_{\{f\}})^\Delta.
  \tag{10.3.1.1}\label{I.10.3.1.1-ko}
\]
\end{env}

\oldpage[I]{184}
\begin{env}[10.3.2]
\label{I.10.3.2-ko}
(\hyperref[I.10.3.1-ko]{10.3.1})의 표기 아래에서, 모든 $f\in A$에
대하여
$A_{\{f\}}=\Gamma(\mathfrak{D}(f),\mathscr{O}_{\mathfrak{X}})$에서
$\Gamma\bigl(\mathfrak{D}(f),
(\widetilde{A}/\widetilde{\mathfrak{J}})|\mathfrak{X}\bigr)
=S_f^{-1}A/S_f^{-1}\mathfrak{J}$로 가는 표준 사상은 \emph{전사}이고,
그 핵은
$\Gamma(\mathfrak{D}(f),\mathfrak{J}^\Delta)=\mathfrak{J}_{\{f\}}$이다
(\hyperref[0.7.6.9-ko]{0, 7.6.9}). 따라서 이 사상들은 위상환의 층
$\mathscr{O}_{\mathfrak{X}}$에서 이산환의 층
$(\widetilde{A}/\widetilde{\mathfrak{J}})|\mathfrak{X}$로 가는 연속
\emph{전사} 준동형을 정의하며, 이를 \emph{표준}이라 한다. 그 핵은
$\mathfrak{J}^\Delta$이다. 이 준동형은
$\varphi$가 연속 준동형 $A\to A/\mathfrak{J}$일 때의
$\widetilde{\varphi}$
(\hyperref[I.10.2.1-ko]{10.2.1})와 다르지 않다. 아핀 형식적 스킴의 사상
$({}^a\varphi,\widetilde{\varphi}):
\operatorname{Spec}(A/\mathfrak{J})\to\mathfrak{X}$
(여기서 ${}^a\varphi$는 닫힌 부분공간
$V(\mathfrak{J})\subset\mathfrak{X}$로의 표준 위상동형사상에 그
포함사상을 합성한 것이다)도 \emph{표준}이라 한다.\footnote{\textbf{번역 주.} 인쇄 원문은
${}^a\varphi$가 $\mathfrak{X}$에서 자기 자신으로 가는 항등
위상동형사상이라고 쓰지만, 이는 $\mathfrak{J}$가 임의의 열린 아이디얼일
때 일반적으로 성립하지 않는다. 예를 들어 체 $k$에 대하여 이산 허용환
$A=k\times k$와 열린 아이디얼 $\mathfrak{J}=k\times\{0\}$에 대하여
$\operatorname{Spf}(A)$는 두 점이지만
$\operatorname{Spec}(A/\mathfrak{J})$는 한 점이다. 올바른 상은
$V(\mathfrak{J})\subset\mathfrak{X}$이며, $\mathfrak{J}$가 정의
아이디얼일 때에는 이 상이 $\mathfrak{X}$ 전체이다. 이는 공식 정오표에
실린 항목은 아니며, 위와 같이 형에 맞게 바로잡았다.} 따라서 앞의
결과에 의하여 \emph{표준 동형사상}
\[
  \mathscr{O}_{\mathfrak{X}}/\mathfrak{J}^\Delta
  \xrightarrow{\sim}
  (\widetilde{A}/\widetilde{\mathfrak{J}})|\mathfrak{X}.
  \tag{10.3.2.1}\label{I.10.3.2.1-ko}
\]
을 얻는다.

$\Gamma(\mathfrak{X},\mathfrak{J}^\Delta)=\mathfrak{J}$이므로
$\mathfrak{J}\to\mathfrak{J}^\Delta$가
\emph{포함관계를 엄격히 보존한다}는 것은 명백하다. 앞의 결과에 따라
$\mathfrak{J}\subset\mathfrak{J}'$이면 층
${\mathfrak{J}'}^\Delta/\mathfrak{J}^\Delta$는
$\widetilde{\mathfrak{J}'}/\widetilde{\mathfrak{J}}
=\widetilde{\mathfrak{J}'/\mathfrak{J}}$와 표준적으로 동형이다.
\end{env}

\begin{env}[10.3.3]
\label{I.10.3.3-ko}
가정과 표기는 계속
(\hyperref[I.10.3.1-ko]{10.3.1})의 것이라 하자.
$\mathscr{O}_{\mathfrak{X}}$의 아이디얼층 $\mathscr{J}$가 다음 조건을
만족할 때 이를 \emph{$\mathfrak{X}$의 정의 아이디얼층}(또는
\emph{$\mathfrak{X}$의 정의 아이디얼})이라 하겠다. 모든
$x\in\mathfrak{X}$에 대하여 $x$의 열린 근방 $\mathfrak{D}(f)$가
존재하고, 여기서 $f\in A$이며,
$\mathscr{J}|\mathfrak{D}(f)$는 $\mathfrak{H}^\Delta$ 꼴이다. 이때
$\mathfrak{H}$는 $A_{\{f\}}$의 정의 아이디얼이다.
\end{env}

\begin{proposition}[10.3.4]
\label{I.10.3.4-ko}
모든 $f\in A$에 대하여 $\mathfrak{X}$의 모든 정의 아이디얼은
$\mathfrak{D}(f)$의 정의 아이디얼을 유도한다.
\end{proposition}

\begin{proof}
이는 (\hyperref[I.10.3.1.1-ko]{10.3.1.1})에서 따른다.
\end{proof}

\begin{proposition}[10.3.5]
\label{I.10.3.5-ko}
$A$가 허용환이면 $\mathfrak{X}=\operatorname{Spf}(A)$의 모든 정의
아이디얼은 $A$의 유일하게 결정되는 정의 아이디얼 $\mathfrak{J}$에 대한
$\mathfrak{J}^\Delta$ 꼴이다.
\end{proposition}

\begin{proof}
실제로 $\mathscr{J}$를 $\mathfrak{X}$의 정의 아이디얼이라 하자.
가정과 $\mathfrak{X}$의 준콤팩트성에 따라 유한 개의 원소
$f_i\in A$가 존재하여 $\mathfrak{D}(f_i)$들이 $\mathfrak{X}$를 덮고,
$\mathscr{J}|\mathfrak{D}(f_i)=\mathfrak{H}_i^\Delta$이다. 여기서
$\mathfrak{H}_i$는 $A_{\{f_i\}}$의 정의 아이디얼이다. 모든 $i$에
대하여 $(\mathfrak{K}_i)_{\{f_i\}}=\mathfrak{H}_i$인 $A$의 열린
아이디얼 $\mathfrak{K}_i$가 존재한다
(\hyperref[0.7.6.9-ko]{0, 7.6.9}). 모든 $\mathfrak{K}_i$에 포함되는
$A$의 정의 아이디얼을 $\mathfrak{K}$라 하자.
$\mathscr{J}/\mathfrak{K}^\Delta$의
$\operatorname{Spec}(A/\mathfrak{K})$의 구조층
$\widetilde{A/\mathfrak{K}}$ 안에서의 표준 상
(\hyperref[I.10.3.2-ko]{10.3.2})은 $\mathfrak{D}(f_i)$ 위로 제한하면
$\widetilde{\mathfrak{K}_i/\mathfrak{K}}$의 제한과 같다. 따라서 이
표준 상은 $\operatorname{Spec}(A/\mathfrak{K})$ 위의 준연접층이고,
$\mathfrak{K}$를 포함하는 $A$의 아이디얼 $\mathfrak{J}$에 대하여
$\widetilde{\mathfrak{J}/\mathfrak{K}}$ 꼴이다
(\hyperref[I.1.4.1-ko]{1.4.1}). 그러므로
$\mathscr{J}=\mathfrak{J}^\Delta$이다
(\hyperref[I.10.3.2-ko]{10.3.2}). 또한 모든 $i$에 대하여
$\mathfrak{H}_i^{n_i}\subset\mathfrak{K}_{\{f_i\}}$인 정수 $n_i$가
존재한다. $n$을 $n_i$들의 최댓값이라 하면
$(\mathscr{J}/\mathfrak{K}^\Delta)^n=0$이고, 따라서
(\hyperref[I.10.3.2-ko]{10.3.2})에 의하여
$(\widetilde{\mathfrak{J}/\mathfrak{K}})^n=0$이다. 그러므로 마침내
$(\mathfrak{J}/\mathfrak{K})^n=0$이고
(\hyperref[I.1.3.13-ko]{1.3.13}), 이는 $\mathfrak{J}$가 $A$의 정의
아이디얼임을 보인다 (\hyperref[0.7.1.4-ko]{0, 7.1.4}).
\end{proof}

\begin{proposition}[10.3.6]
\label{I.10.3.6-ko}
$A$를 아딕환, $\mathfrak{J}$를 $A$의 정의 아이디얼이라 하고,
$\mathfrak{J}/\mathfrak{J}^2$가 유한 생성
$A/\mathfrak{J}$-가군이라고 하자. 그러면 모든 정수 $n>0$에 대하여
$(\mathfrak{J}^\Delta)^n=(\mathfrak{J}^n)^\Delta$이다.
\end{proposition}

\begin{proof}
실제로 모든 $f\in A$에 대하여 $\mathfrak{J}^n$은 열린 아이디얼이므로
\[
  \bigl(\Gamma(\mathfrak{D}(f),\mathfrak{J}^\Delta)\bigr)^n
  =(\mathfrak{J}_{\{f\}})^n
  =(\mathfrak{J}^n)_{\{f\}}
  =\Gamma\bigl(\mathfrak{D}(f^n),(\mathfrak{J}^n)^\Delta\bigr)
\]
\oldpage[I]{185}
이다. 이는 (\hyperref[I.10.3.1.1-ko]{10.3.1.1})과
(\hyperref[0.7.6.12-ko]{0, 7.6.12})에서 따른다.
$(\mathfrak{J}^\Delta)^n$은 준층
$U\mapsto(\Gamma(U,\mathfrak{J}^\Delta))^n$의 층화이므로
(\hyperref[0.4.1.6-ko]{0, 4.1.6}), $\mathfrak{D}(f)$들이
$\mathfrak{X}$의 위상의 기저를 이룬다는 사실로부터 결론이
따른다.\footnote{\textbf{번역 주.} 인쇄 원문은 이 명제의 증명에서
“따름정리가 따른다”라고 쓰지만, 이곳에는 따름정리가 없으므로 “결론이
따른다”로 바로잡았다.}
\end{proof}

\begin{env}[10.3.7]
\label{I.10.3.7-ko}
$\mathfrak{X}$의 정의 아이디얼들의 족 $(\mathscr{J}_\lambda)$가 다음
조건을 만족할 때 이를 \emph{정의 아이디얼들의 기본계}라고 한다.
$\mathfrak{X}$의 모든 정의 아이디얼은 어떤 $\mathscr{J}_\lambda$를
포함한다. $\mathscr{J}_\lambda=\mathfrak{J}_\lambda^\Delta$이므로 이는
$\mathfrak{J}_\lambda$들이 $A$에서 \emph{0의 근방 기본계}를 이룬다는
것과 같다. $\mathfrak{D}(f_\alpha)$들이 $\mathfrak{X}$를 덮도록
$A$의 원소들의 족 $(f_\alpha)$를 잡자.
$(\mathscr{J}_\lambda)$가 $\mathscr{O}_{\mathfrak{X}}$의 아이디얼들의
감소 여과족이고, 모든 $\alpha$에 대하여
$(\mathscr{J}_\lambda|\mathfrak{D}(f_\alpha))$가
$\mathfrak{D}(f_\alpha)$의 정의 아이디얼들의 기본계를 이룬다면,
$(\mathscr{J}_\lambda)$는 $\mathfrak{X}$의 정의 아이디얼들의 기본계이다.
실제로 $\mathfrak{X}$의 모든 정의 아이디얼 $\mathscr{J}$에 대하여
$\mathfrak{D}(f_i)$들로 이루어진 $\mathfrak{X}$의 유한 덮개가 존재하여,
모든 $i$에 대하여
$\mathscr{J}_{\lambda_i}|\mathfrak{D}(f_i)$는
$\mathfrak{D}(f_i)$의 정의 아이디얼이고
$\mathscr{J}|\mathfrak{D}(f_i)$에 포함된다. 모든 $i$에 대하여
$\mathscr{J}_\mu\subset\mathscr{J}_{\lambda_i}$인 지표 $\mu$를 잡으면,
(\hyperref[I.10.3.3-ko]{10.3.3})에 따라 $\mathscr{J}_\mu$는
$\mathfrak{X}$의 정의 아이디얼이고, 명백히 $\mathscr{J}$에 포함된다.
이로부터 주장이 따른다.
\end{env}

\subsection{형식적 준스킴과 형식적 준스킴의 사상.}
\label{subsection:I.10.4-ko}

\begin{env}[10.4.1]
\label{I.10.4.1-ko}
위상환 달린 공간 $\mathfrak{X}$가 주어졌다고 하자. 열린집합
$U\subset\mathfrak{X}$ 위에 $\mathfrak{X}$가 유도하는 위상환 달린 공간이
아핀 형식적 스킴(각각 그 환이 아딕인 아핀 형식적 스킴, 아딕이고 뇌터인
아핀 형식적 스킴)일 때, $U$를 \emph{아핀 형식적 열린집합}(각각
\emph{아딕 아핀 형식적 열린집합}, \emph{뇌터 아핀 형식적 열린집합})이라
한다.
\end{env}

\begin{definition}[10.4.2]
\label{I.10.4.2-ko}
모든 점이 아핀 형식적 열린 근방을 갖는 위상환 달린 공간
$\mathfrak{X}$를 형식적 준스킴이라 한다. 형식적 준스킴
$\mathfrak{X}$가 아딕(각각 국소 뇌터)이라는 것은 $\mathfrak{X}$의 모든
점이 아딕 아핀 형식적 열린 근방(각각 뇌터 아핀 형식적 열린 근방)을
갖는다는 뜻이다.
국소 뇌터인 $\mathfrak{X}$의 바탕공간이 준콤팩트일 때(따라서 그
바탕공간은 뇌터이다) 이를 뇌터라 한다.
\end{definition}

\begin{proposition}[10.4.3]
\label{I.10.4.3-ko}
$\mathfrak{X}$가 형식적 준스킴(각각 국소 뇌터 형식적 준스킴)이면,
아핀 형식적 열린집합들(각각 뇌터 아핀 형식적 열린집합들)은
$\mathfrak{X}$의 위상의 기저를 이룬다.
\end{proposition}

\begin{proof}
이는 (\hyperref[I.10.4.2-ko]{10.4.2})와
(\hyperref[I.10.1.4-ko]{10.1.4})에서 따른다. 여기서 $A$가 뇌터
아딕환이면 $A_{\{f\}}$도 모든 $f\in A$에 대하여 그러하다는 사실
(\hyperref[0.7.6.11-ko]{0, 7.6.11})을 사용한다.
\end{proof}

\begin{corollary}[10.4.4]
\label{I.10.4.4-ko}
$\mathfrak{X}$가 형식적 준스킴(각각 국소 뇌터 형식적 준스킴, 뇌터
형식적 준스킴)이면, $\mathfrak{X}$의 모든 열린집합 위에 유도된
위상환 달린 공간도 형식적 준스킴(각각 국소 뇌터 형식적 준스킴,
뇌터 형식적 준스킴)이다.
\end{corollary}

\begin{definition}[10.4.5]
\label{I.10.4.5-ko}
두 형식적 준스킴 $\mathfrak{X}$, $\mathfrak{Y}$가 주어졌다고 하자.
$\mathfrak{X}$에서 $\mathfrak{Y}$로 가는 위상환 달린 공간의 사상
$(\psi,\theta)$ 가운데 모든 $x\in\mathfrak{X}$에 대하여
$\theta_x^\sharp$가 국소 준동형
$\mathscr{O}_{\psi(x)}\to\mathscr{O}_x$인 것을 (형식적 준스킴의)
사상이라 한다.
\end{definition}

두 형식적 준스킴의 사상의 합성도 형식적 준스킴의 사상임은 자명하다.
따라서 형식적 준스킴들은 하나의 \emph{범주}를 이룬다.
$\operatorname{Hom}(\mathfrak{X},\mathfrak{Y})$로 형식적 준스킴
$\mathfrak{X}$에서 형식적 준스킴 $\mathfrak{Y}$로 가는 사상들의 집합을
나타낸다.

\oldpage[I]{186}
$U$가 $\mathfrak{X}$의 열린부분이면, $U$ 위에 $\mathfrak{X}$가 유도하는
형식적 준스킴에서 $\mathfrak{X}$로 가는 표준 포함 사상은 형식적
준스킴의 사상이다(더 나아가 위상환 달린 공간의 \emph{단사사상}이다
(\hyperref[0.4.1.1-ko]{0, 4.1.1})).

\begin{proposition}[10.4.6]
\label{I.10.4.6-ko}
$\mathfrak{X}$를 형식적 준스킴이라 하고,
$\mathfrak{S}=\operatorname{Spf}(A)$를 아핀 형식적 스킴이라 하자.
형식적 준스킴 $\mathfrak{X}$에서 형식적 준스킴 $\mathfrak{S}$로 가는
사상들과 환 $A$에서 위상환
$\Gamma(\mathfrak{X},\mathscr{O}_{\mathfrak{X}})$로 가는 연속
준동형들 사이에는 표준 전단사 대응이 존재한다.
\end{proposition}

증명은 (\hyperref[I.2.2.4-ko]{2.2.4})의 증명과 같다. 다만
“준동형”을 “연속 준동형”으로, “아핀 열린집합”을 “아핀 형식적
열린집합”으로 바꾸고, (\hyperref[I.10.2.2-ko]{10.2.2})를
(\hyperref[I.1.7.3-ko]{1.7.3}) 대신 사용한다. 자세한 내용은 독자에게
맡긴다.

\begin{env}[10.4.7]
\label{I.10.4.7-ko}
형식적 준스킴 $\mathfrak{S}$가 주어졌다고 하자. 형식적 준스킴
$\mathfrak{X}$와 사상 $\varphi:\mathfrak{X}\to\mathfrak{S}$의 자료가
주어졌을 때, 형식적 준스킴 $\mathfrak{X}$를
\emph{$\mathfrak{S}$ 위의 형식적 준스킴} 또는
\emph{$\mathfrak{S}$-형식적 준스킴}이라 하고,
$\varphi$를 $\mathfrak{S}$-형식적 준스킴 $\mathfrak{X}$의
\emph{구조 사상}이라 한다. $\mathfrak{S}=\operatorname{Spf}(A)$이고
$A$가 허용환이면, $\mathfrak{S}$-형식적 준스킴 $\mathfrak{X}$를
\emph{$A$-형식적 준스킴} 또는 \emph{$A$ 위의 형식적 준스킴}이라고도
한다. 모든 형식적 준스킴은 (이산 위상을 갖춘) $\mathbf{Z}$ 위의
형식적 준스킴으로 볼 수 있다.

$\mathfrak{X}$, $\mathfrak{Y}$를 두 $\mathfrak{S}$-형식적 준스킴이라
하자. 사상 $u:\mathfrak{X}\to\mathfrak{Y}$에 대하여 도식
\[
  \xymatrix{
    \mathfrak{X}\ar[rr]^u\ar[rd] & &
    \mathfrak{Y}\ar[ld]\\
    & \mathfrak{S}
  }
\]
(빗금 화살표들은 구조 사상이다)이 가환일 때, 이 사상을
\emph{$\mathfrak{S}$-사상}이라 한다. 이 정의에 따라 (고정된
$\mathfrak{S}$에 대한) $\mathfrak{S}$-형식적 준스킴들은 하나의
\emph{범주}를 이룬다.
$\operatorname{Hom}_{\mathfrak{S}}(\mathfrak{X},\mathfrak{Y})$로
$\mathfrak{S}$-형식적 준스킴 $\mathfrak{X}$에서
$\mathfrak{S}$-형식적 준스킴 $\mathfrak{Y}$로 가는
$\mathfrak{S}$-사상들의 집합을 나타낸다.
$\mathfrak{S}=\operatorname{Spf}(A)$이면 \emph{$A$-사상}이라는 말도
\emph{$\mathfrak{S}$-사상} 대신 사용한다.
\end{env}

\begin{env}[10.4.8]
\label{I.10.4.8-ko}
모든 아핀 스킴은 아핀 형식적 스킴으로 볼 수 있으므로
(\hyperref[I.10.1.2-ko]{10.1.2}), 모든 (일반적인) 준스킴은 형식적
준스킴으로 볼 수 있다. 또한 (\hyperref[I.10.4.5-ko]{10.4.5})에 따라
\emph{일반적인} 준스킴에 대해서는 \emph{형식적} 준스킴의 사상(각각
$S$-사상)이 §~\hyperref[section:I.2-ko]{2}에서 정의한 사상(각각
$S$-사상)과 일치한다.
\end{env}

\subsection{형식적 준스킴의 정의 아이디얼.}
\label{subsection:I.10.5-ko}

\begin{env}[10.5.1]
\label{I.10.5.1-ko}
형식적 준스킴 $\mathfrak{X}$가 주어졌다고 하자.
$\mathscr{O}_{\mathfrak{X}}$-아이디얼 $\mathscr{J}$가 다음 조건을 만족할
때 이를 $\mathfrak{X}$의 \emph{정의 아이디얼층}(또는
\emph{정의 아이디얼})이라 한다. 모든 $x\in\mathfrak{X}$에 대하여 아핀
형식적 열린 근방 $U$가 존재하여, $\mathscr{J}|U$는 $U$ 위에
$\mathfrak{X}$가 유도하는 아핀 형식적 스킴의 정의 아이디얼이다
(\hyperref[I.10.3.3-ko]{10.3.3}). 이때
(\hyperref[I.10.3.1.1-ko]{10.3.1.1})과
(\hyperref[I.10.4.3-ko]{10.4.3})에 따라 모든 열린집합
$V\subset\mathfrak{X}$에 대하여 $\mathscr{J}|V$는 $V$ 위에
$\mathfrak{X}$가 유도하는 형식적 준스킴의 정의 아이디얼이다.

$\mathfrak{X}$의 정의 아이디얼들로 이루어진 족
$(\mathscr{J}_\lambda)$에 대하여, $\mathfrak{X}$의 아핀 형식적
열린집합들로 이루어진 덮개 $(U_\alpha)$가 존재하여 모든 $\alpha$에
대해 $\mathscr{J}_\lambda|U_\alpha$들로 이루어진 족이 $U_\alpha$ 위에
$\mathfrak{X}$가 유도하는 아핀 형식적 스킴의 정의 아이디얼들의
기본계를 이룰 때(\hyperref[I.10.3.6-ko]{10.3.6}), 이 족을
\emph{정의 아이디얼들의}
\oldpage[I]{187}
\emph{기본계}라고 한다. (\hyperref[I.10.3.7-ko]{10.3.7})의 마지막
주에 따라, $\mathfrak{X}$가 아핀 형식적 스킴이면 이 정의는
(\hyperref[I.10.3.7-ko]{10.3.7})에서 주어진 정의와 일치한다. $V$를
$\mathfrak{X}$의 임의의 열린집합이라 하면, $\mathscr{J}_\lambda|V$들은
$V$ 위에 유도된 형식적 준스킴의 정의 아이디얼들의 기본계를 이룬다. 이는
(\hyperref[I.10.3.1.1-ko]{10.3.1.1})에서 따른다. $\mathfrak{X}$가
\emph{국소 뇌터} 형식적 준스킴이고 $\mathscr{J}$가 $\mathfrak{X}$의
정의 아이디얼이면, (\hyperref[I.10.3.6-ko]{10.3.6})에 따라 거듭제곱들
$\mathscr{J}^n$은 $\mathfrak{X}$의 정의 아이디얼들의 기본계를 이룬다.
\end{env}

\begin{env}[10.5.2]
\label{I.10.5.2-ko}
$\mathfrak{X}$를 형식적 준스킴이라 하고, $\mathscr{J}$를
$\mathfrak{X}$의 정의 아이디얼이라 하자. 그러면 환 달린 공간
$(\mathfrak{X},\mathscr{O}_{\mathfrak{X}}/\mathscr{J})$는
\emph{준스킴}(일반적인 의미에서)이다. $\mathfrak{X}$가 아핀 형식적
스킴(각각 국소 뇌터 형식적 준스킴, 뇌터 형식적 준스킴)이면 이 준스킴은
아핀(각각 국소 뇌터, 뇌터)이다. 실제로 곧바로 아핀인 경우로 환원되며,
그 경우에는 (\hyperref[I.10.3.2-ko]{10.3.2})에서 이미 이 명제를
증명했다. 또한
$\theta:\mathscr{O}_{\mathfrak{X}}\to
\mathscr{O}_{\mathfrak{X}}/\mathscr{J}$가 표준 준동형이면,
$u=(1_{\mathfrak{X}},\theta)$는 형식적 준스킴의 \emph{사상}
(\emph{표준 사상}이라 한다)
$(\mathfrak{X},\mathscr{O}_{\mathfrak{X}}/\mathscr{J})\to
(\mathfrak{X},\mathscr{O}_{\mathfrak{X}})$이다. 이 사실도 곧바로 아핀인
경우로 환원되며, 그 경우에는
(\hyperref[I.10.3.2-ko]{10.3.2})에서 이미 확인했다.
\end{env}

\begin{proposition}[10.5.3]
\label{I.10.5.3-ko}
$\mathfrak{X}$를 형식적 준스킴이라 하고, $(\mathscr{J}_\lambda)$를
$\mathfrak{X}$의 정의 아이디얼들의 기본계라 하자. 그러면 위상환의 층
$\mathscr{O}_{\mathfrak{X}}$는 의사이산 위상환의 층들
(\hyperref[0.3.8.1-ko]{0, 3.8.1})
$\mathscr{O}_{\mathfrak{X}}/\mathscr{J}_\lambda$의 사영극한이다.
\end{proposition}

\begin{proof}
$\mathfrak{X}$의 위상에는 준콤팩트 아핀 형식적 열린집합들로 이루어진
기저가 있으므로(\hyperref[I.10.4.3-ko]{10.4.3}), 아핀인 경우로
환원된다. 이 경우 명제는 (\hyperref[I.10.3.5-ko]{10.3.5}),
(\hyperref[I.10.3.2-ko]{10.3.2})와
(\hyperref[I.10.1.1-ko]{10.1.1})의 정의에서 따른다.
\end{proof}

모든 형식적 준스킴이 정의 아이디얼을 갖는지는 확실하지 않다. 그러나
다음은 성립한다.

\begin{proposition}[10.5.4]
\label{I.10.5.4-ko}
$\mathfrak{X}$를 국소 뇌터 형식적 준스킴이라 하자. 그러면
정의 아이디얼 $\mathscr{T}$가 $\mathfrak{X}$의 가장 큰 정의
아이디얼로서 존재한다. 이는 정의 아이디얼 $\mathscr{J}$ 가운데 준스킴
$(\mathfrak{X},\mathscr{O}_{\mathfrak{X}}/\mathscr{J})$가 축소인
유일한 것이다. $\mathscr{J}$가
$\mathfrak{X}$의 정의 아이디얼이면, $\mathscr{T}$는 표준 준동형
$\mathscr{O}_{\mathfrak{X}}\to
\mathscr{O}_{\mathfrak{X}}/\mathscr{J}$에 의한
$\mathscr{O}_{\mathfrak{X}}/\mathscr{J}$의 멱영근기의 역상이다.
\end{proposition}

\begin{proof}
먼저 $\mathfrak{X}=\operatorname{Spf}(A)$이고 $A$가 뇌터 아딕환이라고
하자. $\mathscr{T}$의 존재성과 성질들은
(\hyperref[I.10.3.5-ko]{10.3.5})\footnote{\textbf{번역 주.} 인쇄 원문은
이곳에서 (10.3.4)를 참조하지만, 출판된 EGA II 정오표는 이를
(10.3.5)로 고치도록 지시한다. 그 공식 교정을 적용했다.}와
(\hyperref[I.5.1.1-ko]{5.1.1})에서 바로 따른다. 이때 $A$의 가장 큰
정의 아이디얼의 존재성과 성질들
(\hyperref[0.7.1.6-ko]{0, 7.1.6}과
\hyperref[0.7.1.7-ko]{7.1.7})을 이용한다.

일반적인 경우에 $\mathscr{T}$의 존재성과 성질들을 증명하려면,
$U\supset V$를 $\mathfrak{X}$의 두 뇌터 아핀 형식적 열린집합이라 할 때
$U$의 가장 큰 정의 아이디얼 $\mathscr{T}_U$가 $V$의 가장 큰 정의
아이디얼 $\mathscr{T}_V$를 유도함을 보이면 충분하다. 그런데
$(V,(\mathscr{O}_{\mathfrak{X}}|V)/(\mathscr{T}_U|V))$는 축소이므로,
이는 앞에서 보인 결과에서 따른다.
\end{proof}

$\mathfrak{X}_{\mathrm{red}}$는 축소 준스킴(일반적인 의미에서)
$(\mathfrak{X},\mathscr{O}_{\mathfrak{X}}/\mathscr{T})$를 나타낸다.

\begin{corollary}[10.5.5]
\label{I.10.5.5-ko}
$\mathfrak{X}$를 국소 뇌터 형식적 준스킴이라 하고, $\mathscr{T}$를
$\mathfrak{X}$의 가장 큰 정의 아이디얼이라 하자. $V$가
$\mathfrak{X}$의 임의의 열린집합이면, $\mathscr{T}|V$는
$\mathfrak{X}$가 $V$ 위에 유도하는 형식적 준스킴의 가장 큰 정의
아이디얼이다.
\end{corollary}

\oldpage[I]{188}

\begin{proposition}[10.5.6]
\label{I.10.5.6-ko}
$\mathfrak{X}$, $\mathfrak{Y}$를 두 형식적 준스킴이라 하고,
$\mathscr{J}$와 $\mathscr{K}$를 각각 $\mathfrak{X}$와
$\mathfrak{Y}$의 정의 아이디얼이라 하며,
$f:\mathfrak{X}\to\mathfrak{Y}$를 형식적 준스킴의 사상이라 하자.
\begin{enumerate}
  \item[{\upshape(i)}]
    $f^*(\mathscr{K})\mathscr{O}_{\mathfrak{X}}\subset\mathscr{J}$이면,
    일반적인 의미에서 준스킴의 사상인
    $f':(\mathfrak{X},\mathscr{O}_{\mathfrak{X}}/\mathscr{J})\to
    (\mathfrak{Y},\mathscr{O}_{\mathfrak{Y}}/\mathscr{K})$가 유일하게
    존재하고 다음 도식을 가환이게 한다.
    \[
      \label{I.10.5.6.1-ko}
      \xymatrix{
        (\mathfrak{X},\mathscr{O}_{\mathfrak{X}})\ar[r]^f &
        (\mathfrak{Y},\mathscr{O}_{\mathfrak{Y}})\\
        (\mathfrak{X},\mathscr{O}_{\mathfrak{X}}/\mathscr{J})
          \ar[r]_{f'}\ar[u] &
        (\mathfrak{Y},\mathscr{O}_{\mathfrak{Y}}/\mathscr{K})\ar[u]
      }
      \tag{10.5.6.1}
    \]
    여기서 세로 화살표들은 표준 사상이다.
  \item[{\upshape(ii)}]
    $\mathfrak{X}=\operatorname{Spf}(A)$,
    $\mathfrak{Y}=\operatorname{Spf}(B)$가 아핀 형식적 스킴이고,
    $\mathscr{J}=\mathfrak{J}^{\Delta}$,
    $\mathscr{K}=\mathfrak{K}^{\Delta}$라 하자. 여기서
    $\mathfrak{J}$와 $\mathfrak{K}$는 각각 $A$와 $B$의 정의
    아이디얼이다. 또한
    $f=({}^a\varphi,\widetilde{\varphi})$이고
    $\varphi:B\to A$가 연속 준동형이라 하자. 이때
    $f^*(\mathscr{K})\mathscr{O}_{\mathfrak{X}}\subset\mathscr{J}$일
    필요충분조건은
    $\varphi(\mathfrak{K})\subset\mathfrak{J}$인 것이다. 이 경우
    $f'$은 사상 $({}^a\varphi',\widetilde{\varphi'})$이며, 여기서
    $\varphi':B/\mathfrak{K}\to A/\mathfrak{J}$은 $\varphi$에서
    몫으로 내려가 얻은 준동형이다.
\end{enumerate}
\end{proposition}

\begin{proof}
\medskip\noindent
\begin{enumerate}
  \item[{\upshape(i)}] $f=(\psi,\theta)$라 하면, 가정에 따라 준동형
    $\theta^\sharp:\psi^*(\mathscr{O}_{\mathfrak{Y}})\to
    \mathscr{O}_{\mathfrak{X}}$에 의한
    $\psi^*(\mathscr{O}_{\mathfrak{Y}})$의 아이디얼층
    $\psi^*(\mathscr{K})$의 상은 $\mathscr{J}$에 포함된다
    (\hyperref[0.4.3.5-ko]{0, 4.3.5}). 따라서 몫으로 내려가면
    $\theta^\sharp$에서 환의 층의 준동형
    \[
      \omega:\psi^*(\mathscr{O}_{\mathfrak{Y}}/\mathscr{K})
      =\psi^*(\mathscr{O}_{\mathfrak{Y}})/\psi^*(\mathscr{K})
      \to\mathscr{O}_{\mathfrak{X}}/\mathscr{J}~;
    \]
    을 얻는다. 또한 모든 $x\in\mathfrak{X}$에 대하여
    $\theta_x^\sharp$가 \emph{국소} 준동형이므로 $\omega_x$도 그러하다.
    따라서 환 달린 공간의 사상 $(\psi,\omega^b)$는
    (\hyperref[I.2.2.1-ko]{2.2.1}) 문제의 요구를 만족하는 환 달린 공간의
    유일한 사상 $f'$이다.
  \item[{\upshape(ii)}] 아핀 형식적 스킴의 사상들과 환의 연속 준동형들
    사이의 표준적이고 함자적인 대응
    (\hyperref[I.10.2.2-ko]{10.2.2})에 따르면, 지금의 경우 관계
    $f^*(\mathscr{K})\mathscr{O}_{\mathfrak{X}}\subset\mathscr{J}$로부터
    $f'=({}^a\varphi',\widetilde{\varphi'})$를 얻는다. 여기서
    $\varphi':B/\mathfrak{K}\to A/\mathfrak{J}$은 다음 도식을 가환이게
    하는 유일한 준동형이다.
    \[
      \label{I.10.5.6.2-ko}
      \xymatrix{
        B\ar[r]^\varphi\ar[d] & A\ar[d]\\
        B/\mathfrak{K}\ar[r]_{\varphi'} & A/\mathfrak{J}
      }
      \tag{10.5.6.2}
    \]
    따라서 $\varphi'$이 존재하면
    $\varphi(\mathfrak{K})\subset\mathfrak{J}$이다. 역으로 이 조건이
    성립한다고 하자. $\varphi'$을 도식
    (\hyperref[I.10.5.6.2-ko]{10.5.6.2})을 가환이게 하는 유일한
    준동형이라 하고 $f'=({}^a\varphi',\widetilde{\varphi'})$로 놓으면,
    도식 (\hyperref[I.10.5.6.1-ko]{10.5.6.1})이 가환임은 명백하다. 이제
    각각 $f$와 $f'$에 대응하는 준동형
    ${}^a\varphi^*(\mathscr{O}_{\mathfrak{Y}})\to
    \mathscr{O}_{\mathfrak{X}}$와
    ${}^a{\varphi'}^*(\mathscr{O}_{\mathfrak{Y}}/\mathscr{K})\to
    \mathscr{O}_{\mathfrak{X}}/\mathscr{J}$을 고려하면, 이 가환성으로부터
    $f^*(\mathscr{K})\mathscr{O}_{\mathfrak{X}}\subset\mathscr{J}$가
    따름을 알 수 있다.
\end{enumerate}
\end{proof}

위에서 정의한 대응 $f\mapsto f'$이 \emph{함자적}임은 명백하다.

\subsection{준스킴들의 귀납극한으로서의 형식적 준스킴.}
\label{subsection:I.10.6-ko}

\begin{env}[10.6.1]
\label{I.10.6.1-ko}
$\mathfrak{X}$를 형식적 준스킴이라 하고,
$(\mathscr{J}_\lambda)$를 $\mathfrak{X}$의 정의 아이디얼들의 기본계라
하자. 각 $\lambda$에 대하여 $f_\lambda$를 표준 사상
$(\mathfrak{X},\mathscr{O}_{\mathfrak{X}}/\mathscr{J}_\lambda)
\to\mathfrak{X}$라 하자(\hyperref[I.10.5.2-ko]{10.5.2}).
$\mathscr{J}_\mu\subset\mathscr{J}_\lambda$이면 표준 준동형
$\mathscr{O}_{\mathfrak{X}}/\mathscr{J}_\mu\to
\mathscr{O}_{\mathfrak{X}}/\mathscr{J}_\lambda$는 보통 의미의 준스킴
사이의 표준 사상

\oldpage[I]{189}

$f_{\mu\lambda}:(\mathfrak{X},
\mathscr{O}_{\mathfrak{X}}/\mathscr{J}_\lambda)\to
(\mathfrak{X},\mathscr{O}_{\mathfrak{X}}/\mathscr{J}_\mu)$를 정의하며,
$f_\lambda=f_\mu\circ f_{\mu\lambda}$가 성립한다. 따라서 준스킴들
$X_\lambda=(\mathfrak{X},
\mathscr{O}_{\mathfrak{X}}/\mathscr{J}_\lambda)$와 사상들
$f_{\mu\lambda}$는 (\hyperref[I.10.4.8-ko]{10.4.8})에 따라 형식적
준스킴의 범주에서 하나의 \emph{귀납계}를 이룬다.
\end{env}

\begin{proposition}[10.6.2]
\label{I.10.6.2-ko}
(\hyperref[I.10.6.1-ko]{10.6.1})의 표기 아래에서 형식적 준스킴
$\mathfrak{X}$와 사상들 $f_\lambda$는 형식적 준스킴의 범주에서 계
$(X_\lambda,f_{\mu\lambda})$의 귀납극한(T, I, 1.8)을 이룬다.
\end{proposition}

\begin{proof}
$\mathfrak{Y}$를 형식적 준스킴이라 하고, 각 지표 $\lambda$에 대하여
사상
\[
  g_\lambda=(\psi_\lambda,\theta_\lambda):X_\lambda\to\mathfrak{Y}
\]
가 주어졌으며, $\mathscr{J}_\mu\subset\mathscr{J}_\lambda$이면
$g_\lambda=g_\mu\circ f_{\mu\lambda}$가 성립한다고 하자. 이 조건과
$X_\lambda$의 정의로부터 우선 모든 $\psi_\lambda$가 바탕공간 사이의
동일한 연속사상 $\psi:\mathfrak{X}\to\mathfrak{Y}$와 일치함을 알 수 있다.
또한 준동형
$\theta_\lambda^\sharp:\psi^*(\mathscr{O}_{\mathfrak{Y}})\to
\mathscr{O}_{X_\lambda}=
\mathscr{O}_{\mathfrak{X}}/\mathscr{J}_\lambda$들은 환의 층의 준동형들의
\emph{사영계}를 이룬다. 따라서 사영극한을 취하면 준동형
\[
  \omega:\psi^*(\mathscr{O}_{\mathfrak{Y}})\to
  \varprojlim\mathscr{O}_{\mathfrak{X}}/\mathscr{J}_\lambda
  =\mathscr{O}_{\mathfrak{X}}
\]
을 얻는다. \emph{환 달린 공간}의 사상 $g=(\psi,\omega^b)$가 도식들
\[
  \label{I.10.6.2.1-ko}
  \xymatrix{
    X_\lambda\ar[rr]^{g_\lambda}\ar[rd]_{f_\lambda} &&
    \mathfrak{Y}\\
    & \mathfrak{X}\ar[ru]_g
  }
  \tag{10.6.2.1}
\]
을 가환이게 하는 \emph{유일한} 사상임은 명백하다. 그러므로 이제 $g$가
\emph{형식적 준스킴}의 사상임을 보이면 된다. 이 문제는
$\mathfrak{X}$와 $\mathfrak{Y}$ 위에서 국소적이므로
$\mathfrak{X}=\operatorname{Spf}(A)$,
$\mathfrak{Y}=\operatorname{Spf}(B)$라 해도 된다. 여기서 $A$와 $B$는
허용환이고 $\mathscr{J}_\lambda=\mathfrak{J}_\lambda^{\Delta}$이며,
$(\mathfrak{J}_\lambda)$는 $A$의 정의 아이디얼들의 기본계이다
(\hyperref[I.10.3.5-ko]{10.3.5}). 이제
$A=\varprojlim A/\mathfrak{J}_\lambda$이므로, 도식
(\hyperref[I.10.6.2.1-ko]{10.6.2.1})을 가환이게 하는 아핀 형식적
스킴의 사상 $g$의 존재는 아핀 형식적 스킴의 사상과 환의 연속
준동형 사이의 전단사 대응 (\hyperref[I.10.2.2-ko]{10.2.2}) 및
사영극한의 정의에서 따른다. 한편 환 달린 공간의 사상으로서 $g$가
유일하므로, 이 사상은 증명 앞부분에서 같은 기호로 나타낸 사상과
일치한다.
\end{proof}

다음 명제는 몇 가지 추가 조건 아래, 보통 의미의 준스킴들로 주어진
귀납계의 귀납극한이 형식적 준스킴의 범주에서 존재함을 보인다.

\begin{proposition}[10.6.3]
\label{I.10.6.3-ko}
$\mathfrak{X}$를 위상공간이라 하고,
$(\mathscr{O}_i,u_{ji})$를 $\mathfrak{X}$ 위의 환의 층들의 사영계라
하자. 그 지표집합은 $\mathbf{N}$이다. $\mathscr{J}_i$를
$u_{0i}:\mathscr{O}_i\to\mathscr{O}_0$의 핵이라 하자. 다음을 가정한다.
\begin{enumerate}
  \item[a)] 환 달린 공간 $(\mathfrak{X},\mathscr{O}_i)$는 준스킴
    $X_i$이다.
  \item[b)] 모든 $x\in\mathfrak{X}$와 모든 $i$에 대하여,
    $\mathfrak{X}$ 안에서 $x$의 열린 근방 $U_i$가 존재하여 제한
    $\mathscr{J}_i|U_i$는 멱영이다.\footnote{\textbf{번역 주.} 인쇄
    원문은 여기서 로만체 $X$를 쓰지만, 출판된 EGA II 정오표는 이를
    프락투어체 $\mathfrak{X}$로 고치도록 지시한다. 그 공식 교정을
    적용했다.}
  \item[c)] 준동형들 $u_{ji}$는 전사이다.
\end{enumerate}

\oldpage[I]{190}

$\mathscr{O}_{\mathfrak{X}}$를 의사이산 환의 층들 $\mathscr{O}_i$의
사영극한인 위상환의 층이라 하고,
$u_i:\mathscr{O}_{\mathfrak{X}}\to\mathscr{O}_i$를 표준 준동형이라
하자. 그러면 위상환 달린 공간
$(\mathfrak{X},\mathscr{O}_{\mathfrak{X}})$는 형식적 준스킴이다.
준동형들 $u_i$는 전사이고, 그 핵들 $\mathscr{J}^{(i)}$는
$\mathfrak{X}$의 정의 아이디얼들의 기본계를 이루며,
$\mathscr{J}^{(0)}$은 아이디얼층들 $\mathscr{J}_i$의 사영극한이다.
\end{proposition}

\begin{proof}
먼저 각 줄기에서 $u_{ji}$는 전사 준동형이며
\emph{따라서 특히} 국소 준동형임에 유의하자. 그러므로
$v_{ij}=(1_{\mathfrak{X}},u_{ji})$는 준스킴의 사상
$X_j\to X_i$ ($i\geq j$)이다
(\hyperref[I.2.2.1-ko]{2.2.1}). 우선 각 $X_i$가
환 $A_i$의 아핀 스킴이라고 가정하자.
\emph{환 준동형}
$\varphi_{ji}:A_i\to A_j$가 존재하여
$u_{ji}=\widetilde{\varphi_{ji}}$를 만족한다
(\hyperref[I.1.7.3-ko]{1.7.3}). 따라서
(\hyperref[I.1.6.3-ko]{1.6.3}), 층 $\mathscr{O}_j$는
$\mathscr{O}_i$-가군으로서 $X_i$ 위에서 준연접이다(그 스칼라
작용은 $u_{ji}$로 정의된다). 이 층은 $A_j$를 $A_i$-가군으로 보아
얻는 가군에 대응하며, 이때 그 가군 구조는 $\varphi_{ji}$로 정한다.
모든 $f\in A_i$에 대하여
$f'=\varphi_{ji}(f)$라 하자. 가정에 따라 열린집합 $D(f)$와
$D(f')$는 $\mathfrak{X}$에서 동일하며,
$\Gamma(D(f),\mathscr{O}_i)=(A_i)_f$에서
$\Gamma(D(f),\mathscr{O}_j)=(A_j)_{f'}$로 가는 준동형 가운데
$u_{ji}$에 대응하는 것은 바로
$(\varphi_{ji})_f$이다
(\hyperref[I.1.6.1-ko]{1.6.1}). 한편 $A_j$를 $A_i$-가군으로 보면,
$(A_j)_{f'}$는 $(A_i)_f$-가군 $(A_j)_f$이다. 따라서
$u_{ji}=\widetilde{\varphi_{ji}}$가 여전히 성립하는데, 이번에는
$\varphi_{ji}$를 $A_i$-가군 준동형으로 본다. 이제 $u_{ji}$가
전사이므로 $\varphi_{ji}$도 전사임을 알 수 있다
(\hyperref[I.1.3.9-ko]{1.3.9}). $\mathfrak{J}_{ji}$가
$\varphi_{ji}$의 핵이면, $u_{ji}$의 핵은 준연접
$\mathscr{O}_i$-가군이고 $\widetilde{\mathfrak{J}_{ji}}$와 같다. 특히
$\mathscr{J}_i=\widetilde{\mathfrak{J}_i}$이고, 여기서
$\mathfrak{J}_i$는 $\varphi_{0i}:A_i\to A_0$의 핵이다. 가정 b)로부터
$\mathscr{J}_i$가 \emph{멱영}임이 따른다. 실제로
$\mathfrak{X}$가 준콤팩트이므로, $\mathfrak{X}$를 유한 개의
열린집합 $U_k$로 덮을 수 있고, 각각
$(\mathscr{J}_i|U_k)^{n_k}=0$을 만족한다.
$n$을 $n_k$들 가운데 가장 큰 수로 택하면
$\mathscr{J}_i^n=0$이다. 따라서 $\mathfrak{J}_i$는 멱영이다
(\hyperref[I.1.3.13-ko]{1.3.13}). 이제 환
$A=\varprojlim A_i$는 허용환이다
(\hyperref[0.7.2.2-ko]{0, 7.2.2}). 또한 표준 준동형
$\varphi_i:A\to A_i$는 전사이다. 그 핵 $\mathfrak{J}^{(i)}$는
$\mathfrak{J}_{ik}$들($k\geq i$)의 사영극한이며,
$\mathfrak{J}^{(i)}$들은 $A$에서 0의 근방들의 기본계를 이룬다.
이 경우 (\hyperref[I.10.6.3-ko]{10.6.3})의 주장들은
(\hyperref[I.10.1.1-ko]{10.1.1})과
(\hyperref[I.10.3.2-ko]{10.3.2})에서 따른다. 실제로
$(\mathfrak{X},\mathscr{O}_{\mathfrak{X}})$는 다름 아닌
$\operatorname{Spf}(A)$이다.

여전히 이 특수한 경우에, $f=(f_i)$가 사영극한
$A=\varprojlim A_i$의 원소라고 하자. 모든 열린집합 $D(f_i)$는
$X_i$에서의 아핀 열린집합이며, $\mathfrak{D}(f)$라는
$\mathfrak{X}$의 열린집합과 동일시된다. 따라서 $X_i$가
$\mathfrak{D}(f)$에 유도하는 준스킴은 아핀 스킴
$\operatorname{Spec}((A_i)_{f_i})$와 동일시된다.

일반적인 경우에는 먼저, 임의의 준콤팩트 열린집합 $U$를
$\mathfrak{X}$에서 택하면 각 $\mathscr{J}_i|U$가 멱영임에
유의하자. 이는 위의
논법으로 알 수 있다. 이제 모든 $x\in\mathfrak{X}$에 대하여,
어떤 열린집합 $U$가 존재하여 $x$를 포함하고 $\mathfrak{X}$에서 열린 근방이며,
\emph{모든 $X_i$에 대하여 아핀 열린집합}임을 보이자. 실제로
$U$를 $X_0$의 아핀 열린집합으로 택하고
$\mathscr{O}_{X_0}=\mathscr{O}_{X_i}/\mathscr{J}_i$임에 유의하자.
앞에서 본 바와 같이 $\mathscr{J}_i|U$가 멱영이므로,
(\hyperref[I.5.1.9-ko]{5.1.9})에 따라 $U$는 각 $X_i$에 대하여도
아핀 열린집합이다. 이제 앞의 조건을 만족하는 모든 $U$에 대하여,
위에서 다룬 아핀 경우에 의해 $(U,\mathscr{O}_{\mathfrak{X}}|U)$는
형식적 준스킴이고, $\mathscr{J}^{(i)}|U$들은 정의 아이디얼들의
기본계를 이루며, $\mathscr{J}^{(0)}|U$는
$\mathscr{J}_i|U$들의 사영극한이다. 이로써 결론을 얻는다.
\end{proof}

\begin{corollary}[10.6.4]
\label{I.10.6.4-ko}
$i\geq j$일 때 $u_{ji}$의 핵이
$\mathscr{J}_i^{j+1}$이고 $\mathscr{J}_1/\mathscr{J}_1^2$가

\oldpage[I]{191}

$\mathscr{O}_0=\mathscr{O}_1/\mathscr{J}_1$-가군으로서 유한
생성이라고 가정하자. 그러면 $\mathfrak{X}$는 아딕 형식적 준스킴이다.
또한 $\mathscr{J}^{(n)}$을
$\mathscr{O}_{\mathfrak{X}}\to\mathscr{O}_n$의 핵이라 하면
$\mathscr{J}^{(n)}=\mathscr{J}^{n+1}$이고
$\mathscr{J}/\mathscr{J}^2$는 $\mathscr{J}_1$과 동형이다. 더욱이
$X_0$가 국소 뇌터(각각 뇌터)이면 $\mathfrak{X}$도 국소 뇌터(각각
뇌터)이다.
\end{corollary}

\begin{proof}
$\mathfrak{X}$와 $X_0$의 바탕공간은 같으므로 문제는 국소적이고, 모든
$X_i$가 아핀이라 해도 된다. 관계
$\mathscr{J}_{ij}=\widetilde{\mathfrak{J}_{ji}}$
((\hyperref[I.10.6.3-ko]{10.6.3})의 표기)를 고려하면
(\hyperref[0.7.2.7-ko]{0, 7.2.7}과
\hyperref[0.7.2.8-ko]{7.2.8})의 대응하는 주장들로 곧바로 환원된다.
여기서 $\mathfrak{J}_1/\mathfrak{J}_1^2$는 유한 생성
$A_0$-가군임에 유의한다
(\hyperref[I.1.3.9-ko]{1.3.9}).
\end{proof}

특히 \emph{모든 국소 뇌터 형식적 준스킴
$\mathfrak{X}$}는
(\hyperref[I.10.6.3-ko]{10.6.3})과
(\hyperref[I.10.6.4-ko]{10.6.4})의 조건을 만족하는 (보통 의미의)
국소 뇌터 준스킴들의 수열 $(X_n)$의 귀납극한이다. 실제로
$\mathscr{J}$를 $\mathfrak{X}$의 정의 아이디얼로 택하고
(\hyperref[I.10.5.4-ko]{10.5.4})
$X_n=(\mathfrak{X},
\mathscr{O}_{\mathfrak{X}}/\mathscr{J}^{n+1})$
로 두면 충분하다
((\hyperref[I.10.5.1-ko]{10.5.1})과
(\hyperref[I.10.6.2-ko]{10.6.2})).

\begin{corollary}[10.6.5]
\label{I.10.6.5-ko}
$A$를 허용환이라 하자. 아핀 형식적 스킴
$\mathfrak{X}=\operatorname{Spf}(A)$가 뇌터이기 위한 필요충분조건은
$A$가 아딕이고 뇌터인 것이다.
\end{corollary}

\begin{proof}
그 조건이 충분함은 명백하다. 역으로 $\mathfrak{X}$가 뇌터라고
가정하고, $\mathfrak{J}$를 $A$의 정의 아이디얼로,
$\mathscr{J}=\mathfrak{J}^\Delta$를 이에 대응하는
$\mathfrak{X}$의 정의 아이디얼로 택하자. 그러면 (보통 의미의) 준스킴
$X_n=(\mathfrak{X},
\mathscr{O}_{\mathfrak{X}}/\mathscr{J}^{n+1})$은 아핀이고 뇌터이다.
따라서 환 $A_n=A/\mathfrak{J}^{n+1}$도 뇌터이고
(\hyperref[I.6.1.3-ko]{6.1.3}), 이로부터
$\mathfrak{J}/\mathfrak{J}^2$가 유한 생성
$A/\mathfrak{J}$-가군임을 알 수 있다. $\mathscr{J}^n$들이
$\mathfrak{X}$의 정의 아이디얼들의 기본계를 이루므로
(\hyperref[I.10.5.1-ko]{10.5.1}),
$\mathscr{O}_{\mathfrak{X}}
=\varprojlim(\mathscr{O}_{\mathfrak{X}}/\mathscr{J}^n)$이다
(\hyperref[I.10.5.3-ko]{10.5.3}). 따라서
(\hyperref[I.10.1.3-ko]{10.1.3})에 의하여 $A$는
$\varprojlim A/\mathfrak{J}^n$과 위상적으로 동형이고, 그러므로
아딕이고 뇌터이다
(\hyperref[0.7.2.8-ko]{0, 7.2.8}).
\end{proof}

\begin{remark}[10.6.6]
\label{I.10.6.6-ko}
(\hyperref[I.10.6.3-ko]{10.6.3})의 표기에서
$\mathscr{F}_i$를 $\mathscr{O}_i$-가군이라 하자. 또한
$i\geq j$에 대하여 $v_{ij}$-사상
$\theta_{ji}:\mathscr{F}_i\to\mathscr{F}_j$가 주어졌고,
$\theta_{kj}\circ\theta_{ji}=\theta_{ki}$가
$k\leq j\leq i$에 대하여 성립한다고 하자. $v_{ij}$의 바탕 연속사상은
항등사상이므로, $\theta_{ji}$는 공간 $\mathfrak{X}$ 위의 아벨군의
층의 준동형이다. 또한 $\mathscr{F}$가 아벨군의 층들의 사영계
$(\mathscr{F}_i)$의 사영극한이면, $\theta_{ji}$들이
$v_{ij}$-사상이라는 사실로부터 사영극한을 취하여 $\mathscr{F}$에
$\mathscr{O}_{\mathfrak{X}}$-가군 구조를 정의할 수 있다. 이 구조를
갖춘 $\mathscr{F}$를, $\theta_{ji}$들에 관한
$\mathscr{O}_i$-가군들의 계 $(\mathscr{F}_i)$의
\emph{사영극한}이라 하겠다. 특별히
$v_{ij}^*(\mathscr{F}_i)=\mathscr{F}_j$이고 $\theta_{ji}$가
\emph{항등사상}인 경우에는, 줄여서 $\mathscr{F}$를 계
$(\mathscr{F}_i)$의 사영극한이라 하겠다. 이때
$v_{ij}^*(\mathscr{F}_i)=\mathscr{F}_j$는 $j\leq i$에 대하여
성립하며, $\theta_{ji}$들은 언급하지 않는다.
\end{remark}

\begin{env}[10.6.7]
\label{I.10.6.7-ko}
$\mathfrak{X}$와 $\mathfrak{Y}$를 두 형식적 준스킴이라 하고,
$\mathscr{J}$와 $\mathscr{K}$를 각각 $\mathfrak{X}$와
$\mathfrak{Y}$의 정의 아이디얼이라 하자. 또한
$f:\mathfrak{X}\to\mathfrak{Y}$를
$f^*(\mathscr{K})\mathscr{O}_{\mathfrak{X}}\subset\mathscr{J}$를
만족하는 형식적 준스킴의 사상이라 하자. 그러면 모든 정수 $n>0$에 대하여
$f^*(\mathscr{K}^n)\mathscr{O}_{\mathfrak{X}}
=(f^*(\mathscr{K})\mathscr{O}_{\mathfrak{X}})^n
\subset\mathscr{J}^n$이다. 따라서
(\hyperref[I.10.5.6-ko]{10.5.6})에 의하여 $f$로부터 (보통 의미의)
준스킴의 사상 $f_n:X_n\to Y_n$을 유도할 수 있다. 여기서
$X_n=(\mathfrak{X},
\mathscr{O}_{\mathfrak{X}}/\mathscr{J}^{n+1})$,
$Y_n=(\mathfrak{Y},
\mathscr{O}_{\mathfrak{Y}}/\mathscr{K}^{n+1})$로 둔다. 정의로부터
곧바로 다음 도식들이
\[
  \label{I.10.6.7.1-ko}
  \xymatrix{
    X_m\ar[r]^{f_m}\ar[d] &
    Y_m\ar[d]\\
    X_n\ar[r]_{f_n} &
    Y_n
  }
  \tag{10.6.7.1}
\]

\oldpage[I]{192}

$m\leq n$일 때 가환함을 알 수 있다. 다시 말해, $(f_n)$은 사상들의
\emph{귀납계}이다.
\end{env}

\begin{env}[10.6.8]
\label{I.10.6.8-ko}
역으로, $(X_n)$과 $(Y_n)$을 각각
(\hyperref[I.10.6.3-ko]{10.6.3})의 조건 b)와 c)를 만족하는 (보통
의미의) 준스킴들의 귀납계라 하고, $\mathfrak{X}$와
$\mathfrak{Y}$를 각각 그 귀납극한이라 하자. 귀납극한의 정의에 따라,
임의의 수열 $(f_n)$이 $X_n\to Y_n$인 사상들로 이루어진 귀납계이면 그
\emph{귀납극한} $f:\mathfrak{X}\to\mathfrak{Y}$가 존재한다. 이 사상은
다음 도식들을 가환이게 하는 유일한 형식적 준스킴의 사상이다.
\[
  \xymatrix{
    X_n\ar[r]^{f_n}\ar[d] &
    Y_n\ar[d]\\
    \mathfrak{X}\ar[r]_f &
    \mathfrak{Y}
  }
\]
\end{env}

\begin{proposition}[10.6.9]
\label{I.10.6.9-ko}
$\mathfrak{X}$와 $\mathfrak{Y}$를 두 국소 뇌터 형식적 준스킴이라
하고, $\mathscr{J}$와 $\mathscr{K}$를 각각 $\mathfrak{X}$와
$\mathfrak{Y}$의 정의 아이디얼이라 하자.
(\hyperref[I.10.6.7-ko]{10.6.7})에서 정의한 대응
$f\mapsto(f_n)$은
사상 $f:\mathfrak{X}\to\mathfrak{Y}$ 가운데
$f^*(\mathscr{K})\mathscr{O}_{\mathfrak{X}}\subset\mathscr{J}$를
만족하는 것들의 집합과, 도식들
(\hyperref[I.10.6.7.1-ko]{10.6.7.1})을 가환이게 하는 사상들의 수열
$(f_n)$의 집합 사이의 전단사 대응이다.
\end{proposition}

\begin{proof}
$f$가 그러한 사상 수열의 귀납극한이면
$f^*(\mathscr{K})\mathscr{O}_{\mathfrak{X}}\subset\mathscr{J}$임을
보여야 한다. 문제는 $\mathfrak{X}$와 $\mathfrak{Y}$ 위에서
국소적이므로, $\mathfrak{X}=\operatorname{Spf}(A)$,
$\mathfrak{Y}=\operatorname{Spf}(B)$인 아핀 경우로 한정해도 된다.
이때 $A$와 $B$는 아딕이고 뇌터이다. 또한
$\mathscr{J}=\mathfrak{J}^\Delta$,
$\mathscr{K}=\mathfrak{K}^\Delta$이다. 여기서 $\mathfrak{J}$와
$\mathfrak{K}$는 각각 $A$와 $B$의 정의 아이디얼이다. 그러면
$X_n=\operatorname{Spec}(A_n)$,
$Y_n=\operatorname{Spec}(B_n)$이고, 여기서
$A_n=A/\mathfrak{J}^{n+1}$,
$B_n=B/\mathfrak{K}^{n+1}$이며, 이는
(\hyperref[I.10.3.6-ko]{10.3.6})과
(\hyperref[I.10.3.2-ko]{10.3.2})에 따른다. 또한
$f_n=({}^a\varphi_n,\widetilde{\varphi_n})$이며, 환 준동형
$\varphi_n:B_n\to A_n$들은 사영계를 이룬다. 따라서
$f=({}^a\varphi,\widetilde{\varphi})$이고,
$\varphi=\varprojlim\varphi_n$이다. 이제 $m=0$일 때 도식
(\hyperref[I.10.6.7.1-ko]{10.6.7.1})의 가환성으로부터 조건
$\varphi_n(\mathfrak{K}/\mathfrak{K}^{n+1})
\subset\mathfrak{J}/\mathfrak{J}^{n+1}$이 모든 $n$에 대하여
성립한다. 따라서 사영극한을 취하면
$\varphi(\mathfrak{K})\subset
\mathfrak{J}$이고, 이로부터
$f^*(\mathscr{K})\mathscr{O}_{\mathfrak{X}}\subset\mathscr{J}$가
따른다
(\hyperref[I.10.5.6-ko]{10.5.6, (ii)}).
\end{proof}

\begin{corollary}[10.6.10]
\label{I.10.6.10-ko}
$\mathfrak{X}$와 $\mathfrak{Y}$를 두 국소 뇌터 형식적 준스킴이라 하고,
$\mathscr{T}$를 $\mathfrak{X}$의 가장 큰 정의 아이디얼이라 하자
(\hyperref[I.10.5.4-ko]{10.5.4}).
\begin{enumerate}
  \item[{\upshape(i)}] $\mathfrak{Y}$의 임의의 정의 아이디얼
    $\mathscr{K}$와 임의의 사상
    $f:\mathfrak{X}\to\mathfrak{Y}$에 대하여,
    $f^*(\mathscr{K})\mathscr{O}_{\mathfrak{X}}\subset\mathscr{T}$이다.
  \item[{\upshape(ii)}]
    $\operatorname{Hom}(\mathfrak{X},\mathfrak{Y})$과 사상 수열
    $(f_n)$ 가운데 도식들
    (\hyperref[I.10.6.7.1-ko]{10.6.7.1})을 가환이게 하는 것들의 집합
    사이에는 표준 전단사 대응이 있다. 여기서
    $X_n=(\mathfrak{X},
    \mathscr{O}_{\mathfrak{X}}/\mathscr{T}^{n+1})$,
    $Y_n=(\mathfrak{Y},
    \mathscr{O}_{\mathfrak{Y}}/\mathscr{K}^{n+1})$이다.
\end{enumerate}
\end{corollary}

\begin{proof}
(ii)는 (i)와 (\hyperref[I.10.6.9-ko]{10.6.9})에서 곧바로 따른다.
(i)를 보이기 위해서는 $\mathfrak{X}=\operatorname{Spf}(A)$,
$\mathfrak{Y}=\operatorname{Spf}(B)$이고 $A$와 $B$가 뇌터이며,
$\mathscr{T}=\mathfrak{T}^\Delta$,
$\mathscr{K}=\mathfrak{K}^\Delta$인 경우로 한정해도 된다. 여기서
$\mathfrak{T}$는 $A$의 가장 큰 정의 아이디얼이고
$\mathfrak{K}$는 $B$의 정의 아이디얼이다.
$f=({}^a\varphi,\widetilde{\varphi})$라 하자. 여기서
$\varphi:B\to A$는 연속 준동형이다. $\mathfrak{K}$의 원소들은
위상적으로 멱영이고
(\hyperref[0.7.1.4-ko]{0, 7.1.4, (ii)}), 따라서
$\varphi(\mathfrak{K})$의 원소들도 위상적으로 멱영이다. 그러므로
$\mathfrak{T}$가 $A$의 위상적으로 멱영인 원소들의 집합이므로
(\hyperref[0.7.1.6-ko]{0, 7.1.6}),
$\varphi(\mathfrak{K})\subset\mathfrak{T}$이다. 따라서
(\hyperref[I.10.5.6-ko]{10.5.6, (ii)})에 의해 결론이 따른다.
\end{proof}

\begin{corollary}[10.6.11]
\label{I.10.6.11-ko}
$\mathfrak{S}$, $\mathfrak{X}$, $\mathfrak{Y}$를 세 국소 뇌터
형식적 준스킴이라 하고,
$f:\mathfrak{X}\to\mathfrak{S}$,
$g:\mathfrak{Y}\to\mathfrak{S}$를 $\mathfrak{X}$와
$\mathfrak{Y}$를 $\mathfrak{S}$-형식적 준스킴으로 만드는 사상이라 하자.
$\mathscr{J}$ (각각 $\mathscr{K}$, $\mathscr{L}$)를
$\mathfrak{S}$ (각각 $\mathfrak{X}$, $\mathfrak{Y}$)의 정의 아이디얼이라 하고,
$f^*(\mathscr{J})\mathscr{O}_{\mathfrak{X}}\subset\mathscr{K}$,
$g^*(\mathscr{J})\mathscr{O}_{\mathfrak{Y}}=\mathscr{L}$이라 가정하자.
다음과 같이 둔다.
$S_n=(\mathfrak{S},
\mathscr{O}_{\mathfrak{S}}/\mathscr{J}^{n+1})$,
$X_n=(\mathfrak{X},
\mathscr{O}_{\mathfrak{X}}/\mathscr{K}^{n+1})$,
$Y_n=(\mathfrak{Y},
\mathscr{O}_{\mathfrak{Y}}/\mathscr{L}^{n+1})$.
그러면

\oldpage[I]{193}

$\operatorname{Hom}_{\mathfrak{S}}(\mathfrak{X},\mathfrak{Y})$과 수열
$(u_n)$ 가운데 각 항이 $S_n$-사상
$u_n:X_n\to Y_n$이고 도식들
(\hyperref[I.10.6.7.1-ko]{10.6.7.1})을 가환이게 하는 것들의 집합
사이에는 표준 전단사 대응이 있다.
\end{corollary}

\begin{proof}
임의의 $\mathfrak{S}$-사상
$u:\mathfrak{X}\to\mathfrak{Y}$에 대하여, 정의에 의해
$f=g\circ u$이므로
\[
  u^*(\mathscr{L})\mathscr{O}_{\mathfrak{X}}
  =u^*(g^*(\mathscr{J})\mathscr{O}_{\mathfrak{Y}})
    \mathscr{O}_{\mathfrak{X}}
  =f^*(\mathscr{J})\mathscr{O}_{\mathfrak{X}}
  \subset\mathscr{K}
\]
이다. 따라서 이 따름정리는
(\hyperref[I.10.6.9-ko]{10.6.9})에서 따른다.
\end{proof}

$m\leq n$일 때, 사상 $f_n:X_n\to Y_n$ 하나가 주어지면
도식 (\hyperref[I.10.6.7.1-ko]{10.6.7.1})을 가환이게 하는 사상
$f_m:X_m\to Y_m$가 유일하게 정해진다. 이는 아핀 경우로 환원하면
곧바로 알 수 있다. 이로써 사상
$\varphi_{mn}:\operatorname{Hom}_{S_n}(X_n,Y_n)
\to\operatorname{Hom}_{S_m}(X_m,Y_m)$가 정의되고,
$\operatorname{Hom}_{S_n}(X_n,Y_n)$들은 $\varphi_{mn}$들에 관하여
\emph{집합들의 사영계}를 이룬다.
(\hyperref[I.10.6.11-ko]{10.6.11})은 또한 다음과 같은 표준 전단사 사상이
존재한다고 진술할 수도 있다.
\[
  \operatorname{Hom}_{\mathfrak{S}}(\mathfrak{X},\mathfrak{Y})
  \xrightarrow{\sim}
  \varprojlim_n\operatorname{Hom}_{S_n}(X_n,Y_n)
\]

\subsection{형식적 준스킴의 곱.}
\label{subsection:I.10.7-ko}

\begin{env}[10.7.1]
\label{I.10.7.1-ko}
$\mathfrak{S}$를 형식적 준스킴이라 하자.
$\mathfrak{S}$-형식적 준스킴들은 하나의 범주를 이루므로,
$\mathfrak{S}$-형식적 준스킴의 \emph{곱}이라는 개념을 정의할 수 있다.
\end{env}

\begin{proposition}[10.7.2]
\label{I.10.7.2-ko}
$\mathfrak{X}=\operatorname{Spf}(B)$,
$\mathfrak{Y}=\operatorname{Spf}(C)$를 아핀 형식적 스킴
$\mathfrak{S}=\operatorname{Spf}(A)$ 위의 두 아핀 형식적 스킴이라 하자.
$\mathfrak{Z}=\operatorname{Spf}(B\widehat{\otimes}_A C)$라 하고,
$p_1$, $p_2$를 $B$와 $C$에서
$B\widehat{\otimes}_A C$로 가는 표준 연속 $A$-준동형
$\rho$, $\sigma$에 각각 대응하는
(\hyperref[I.10.2.2-ko]{10.2.2}) $\mathfrak{S}$-사상이라 하자.
그러면 $(\mathfrak{Z},p_1,p_2)$는 $\mathfrak{S}$ 위의 아핀 형식적 스킴
$\mathfrak{X}$와 $\mathfrak{Y}$의 곱이다.
\end{proposition}

\begin{proof}
(\hyperref[I.10.4.6-ko]{10.4.6})에 따라, 위상 $A$-대수인 허용환 $D$와
임의의 연속 $A$-준동형
$\varphi:B\widehat{\otimes}_A C\to D$에 대하여, 쌍
$(\varphi\circ\rho,\varphi\circ\sigma)$를 대응시키는 사상
\[
  \operatorname{Hom}_A(B\widehat{\otimes}_A C,D)
  \xrightarrow{\sim}
  \operatorname{Hom}_A(B,D)\times\operatorname{Hom}_A(C,D)
\]
이 전단사임을 확인하면 충분하다. 이는 다름 아닌 완비 텐서곱의 보편 성질
(\hyperref[0.7.7.6-ko]{0, 7.7.6})이다.
\end{proof}

\begin{proposition}[10.7.3]
\label{I.10.7.3-ko}
두 $\mathfrak{S}$-형식적 준스킴 $\mathfrak{X}$,
$\mathfrak{Y}$가 주어지면, 곱
$\mathfrak{X}\times_{\mathfrak{S}}\mathfrak{Y}$가 존재한다.
\end{proposition}

\begin{proof}
증명은 (\hyperref[I.3.2.6-ko]{3.2.6})의 증명에서 아핀 스킴(각각
아핀 열린집합)을 아핀 형식적 스킴(각각 아핀 형식적 열린집합)으로
바꾸고, 명제 (\hyperref[I.3.2.2-ko]{3.2.2})를
(\hyperref[I.10.7.2-ko]{10.7.2})로 바꾸면 그대로 얻어진다.
\end{proof}

준스킴의 곱에 관한 모든 형식적 성질
(\hyperref[I.3.2.7-ko]{3.2.7}과
\hyperref[I.3.2.8-ko]{3.2.8},
\hyperref[I.3.3.1-ko]{3.3.1}부터
\hyperref[I.3.3.12-ko]{3.3.12}까지)은
형식적 준스킴의 곱에도 아무 변경 없이 성립한다.

\begin{env}[10.7.4]
\label{I.10.7.4-ko}
$\mathfrak{S}$, $\mathfrak{X}$, $\mathfrak{Y}$를 세 형식적 준스킴이라
하고, $f:\mathfrak{X}\to\mathfrak{S}$,
$g:\mathfrak{Y}\to\mathfrak{S}$를 두 사상이라 하자.
$\mathfrak{S}$, $\mathfrak{X}$, $\mathfrak{Y}$ 안에 각각 정의
아이디얼들의 기본계 $(\mathscr{J}_\lambda)$,
$(\mathscr{K}_\lambda)$, $(\mathscr{L}_\lambda)$가 존재한다고 가정하자.
이 세 기본계의 지표집합은 모두 $I$이며,
$f^*(\mathscr{J}_\lambda)\mathscr{O}_{\mathfrak{X}}
\subset\mathscr{K}_\lambda$ 및
$g^*(\mathscr{J}_\lambda)\mathscr{O}_{\mathfrak{Y}}
\subset\mathscr{L}_\lambda$가 모든 $\lambda$에 대하여 성립한다고 하자.
다음과 같이 놓자.
$S_\lambda=(\mathfrak{S},
\mathscr{O}_{\mathfrak{S}}/\mathscr{J}_\lambda)$,
$X_\lambda=(\mathfrak{X},
\mathscr{O}_{\mathfrak{X}}/\mathscr{K}_\lambda)$,
$Y_\lambda=(\mathfrak{Y},
\mathscr{O}_{\mathfrak{Y}}/\mathscr{L}_\lambda)$.
$\mathscr{J}_\mu\subset\mathscr{J}_\lambda$,
$\mathscr{K}_\mu\subset\mathscr{K}_\lambda$,
$\mathscr{L}_\mu\subset\mathscr{L}_\lambda$인 경우,
$S_\lambda$, $X_\lambda$, $Y_\lambda$는 각각
$S_\mu$, $X_\mu$, $Y_\mu$의 닫힌 부분준스킴이고, 각각은 대응하는
준스킴과 같은

\oldpage[I]{194}

바탕공간을 갖는다(\hyperref[I.10.6.1-ko]{10.6.1}).
$S_\lambda\to S_\mu$가 준스킴의 단사사상이므로, 먼저 두 곱
$X_\lambda\times_{S_\lambda}Y_\lambda$와
$X_\lambda\times_{S_\mu}Y_\lambda$가 동일함을 알 수 있고
(\hyperref[I.3.2.4-ko]{3.2.4}), 이어서
$X_\lambda\times_{S_\mu}Y_\lambda$는
$X_\mu\times_{S_\mu}Y_\mu$의 닫힌 부분준스킴으로 동일시되며, 이
부분준스킴은 후자와 같은 바탕공간을 갖는다
(\hyperref[I.4.3.1-ko]{4.3.1}). 이제 곱
$\mathfrak{X}\times_{\mathfrak{S}}\mathfrak{Y}$는 (보통 의미의) 준스킴
$X_\lambda\times_{S_\lambda}Y_\lambda$들의 \emph{귀납극한}이다. 실제로
(\hyperref[I.10.6.2-ko]{10.6.2})에서와 같이
$\mathfrak{S}$, $\mathfrak{X}$, $\mathfrak{Y}$가 아핀 형식적 스킴인
경우로 환원할 수 있다. (\hyperref[I.10.5.6-ko]{10.5.6, (ii)})와
$\mathfrak{S}$, $\mathfrak{X}$, $\mathfrak{Y}$의 정의 아이디얼들의
기본계에 관한 가정을 고려하면, 이 주장은 두 대수의 완비 텐서곱의
정의(\hyperref[0.7.7.1-ko]{0, 7.7.1})에서 곧바로 따른다.

또한 $\mathfrak{Z}$를 $\mathfrak{S}$-형식적 준스킴이라 하고,
$(\mathscr{M}_\lambda)$를 $\mathfrak{Z}$의 정의 아이디얼들의 기본계라
하되 그 지표집합은 $I$라 하자. 또
$u:\mathfrak{Z}\to\mathfrak{X}$,
$v:\mathfrak{Z}\to\mathfrak{Y}$를 두 $\mathfrak{S}$-사상이라 하고,
$u^*(\mathscr{K}_\lambda)\mathscr{O}_{\mathfrak{Z}}
\subset\mathscr{M}_\lambda$ 및
$v^*(\mathscr{L}_\lambda)\mathscr{O}_{\mathfrak{Z}}
\subset\mathscr{M}_\lambda$라고 하자.
$Z_\lambda=(\mathfrak{Z},
\mathscr{O}_{\mathfrak{Z}}/\mathscr{M}_\lambda)$라 놓고,
$u_\lambda:Z_\lambda\to X_\lambda$와
$v_\lambda:Z_\lambda\to Y_\lambda$를 $S_\lambda$-사상 가운데 각각
$u$와 $v$에 대응하는 것이라 하면
(\hyperref[I.10.5.6-ko]{10.5.6}), $(u,v)_{\mathfrak{S}}$가
$S_\lambda$-사상 $(u_\lambda,v_\lambda)_{S_\lambda}$들의
귀납극한임을 곧바로 확인할 수 있다.

이 번호의 고찰은 특히 $\mathfrak{S}$, $\mathfrak{X}$,
$\mathfrak{Y}$가 국소 뇌터인 경우, 각각 하나의 정의 아이디얼의
거듭제곱들로 이루어진 계를 정의 아이디얼들의 기본계로 택하면 적용된다
(\hyperref[I.10.5.1-ko]{10.5.1}). 그러나
$\mathfrak{X}\times_{\mathfrak{S}}\mathfrak{Y}$가 반드시 국소 뇌터인
것은 아니다(다만 (\agref{I.10.13.5-ko}{10.13.5})를 보라).
\end{env}

\subsection{닫힌 부분집합을 따른 준스킴의 형식적 완비화.}
\label{subsection:I.10.8-ko}

\begin{env}[10.8.1]
\label{I.10.8.1-ko}
$X$를 (보통 의미의) 국소 뇌터 준스킴이라 하고, $X'$을 $X$의
바탕공간의 닫힌 부분집합이라 하자. $\Phi$로 다음 조건을 만족하는
아이디얼층들의 집합을 나타내자. $\mathscr{J}$는
$\mathscr{O}_X$ 안의 연접 아이디얼층이고, $\mathscr{O}_X/\mathscr{J}$의
지지집합은 $X'$이다. 집합 $\Phi$는 비어 있지
않으며(\hyperref[I.5.2.1-ko]{5.2.1},
\hyperref[I.4.1.4-ko]{4.1.4}와
\hyperref[I.6.1.1-ko]{6.1.1}), 관계 $\supset$로 순서를 주기로 한다.
\end{env}

\begin{lemma}[10.8.2]
\label{I.10.8.2-ko}
순서집합 $\Phi$는 여과집합이다. $X$가 뇌터이면, 모든
$\mathscr{J}_0\in\Phi$에 대하여 거듭제곱 $\mathscr{J}_0^n$
($n>0$)으로 이루어진 집합은 $\Phi$에서 공종이다.
\end{lemma}

\begin{proof}
실제로 $\mathscr{J}_1$과 $\mathscr{J}_2$가 $\Phi$에 속하고
$\mathscr{J}=\mathscr{J}_1\cap\mathscr{J}_2$라 놓으면,
$\mathscr{O}_X$가 연접이므로 $\mathscr{J}$도 연접이고
(\hyperref[I.6.1.1-ko]{6.1.1}과
\hyperref[0.5.3.4-ko]{0, 5.3.4}), 모든 $x\in X$에 대하여
$\mathscr{J}_x=(\mathscr{J}_1)_x\cap(\mathscr{J}_2)_x$이다. 따라서
$x\notin X'$이면 $\mathscr{J}_x=\mathscr{O}_x$이고
$x\in X'$이면 $\mathscr{J}_x\neq\mathscr{O}_x$이므로
$\mathscr{J}\in\Phi$이다. 한편 $X$가 뇌터이고
$\mathscr{J}_0$와 $\mathscr{J}$가 $\Phi$에 속하면,
$\mathscr{J}_0^n(\mathscr{O}_X/\mathscr{J})=0$이 되는 정수
$n>0$이 존재한다(\hyperref[I.9.3.4-ko]{9.3.4}). 이는
$\mathscr{J}_0^n\subset\mathscr{J}$를 뜻한다.
\end{proof}

\begin{env}[10.8.3]
\label{I.10.8.3-ko}
이제 $\mathscr{F}$를 연접 $\mathscr{O}_X$-가군이라 하자. 모든
$\mathscr{J}\in\Phi$에 대하여
$\mathscr{F}\otimes_{\mathscr{O}_X}(\mathscr{O}_X/\mathscr{J})$는
연접 $\mathscr{O}_X$-가군이고
(\hyperref[I.9.1.1-ko]{9.1.1}), 그 지지집합은 $X'$에 포함된다. 우리는
이 층을 대개 $X'$으로 제한한 층과 동일시하겠다. $\mathscr{J}$가
$\Phi$를 달릴 때 이 층들은 아벨군의 층들의
\emph{사영계}를 이룬다.
\end{env}

\begin{definition}[10.8.4]
\label{I.10.8.4-ko}
국소 뇌터 준스킴 $X$의 닫힌 부분집합 $X'$과 연접
$\mathscr{O}_X$-가군 $\mathscr{F}$가 주어졌을 때, 층

\oldpage[I]{195}

$\varprojlim_{\Phi}
(\mathscr{F}\otimes_{\mathscr{O}_X}
(\mathscr{O}_X/\mathscr{J}))$의 $X'$으로의 제한을
\emph{$X'$을 따른 $\mathscr{F}$의 완비화}라 하고
$\mathscr{F}_{/X'}$로 나타내며, 혼동의 우려가 없으면
$\widehat{\mathscr{F}}$로도 나타낸다. 이 완비화의 $X'$ 위 단면들을
\emph{$X'$을 따른 $\mathscr{F}$의 형식적 단면}이라 한다.
\end{definition}

모든 열린집합 $U\subset X$에 대하여 다음이 곧바로 성립한다.
$(\mathscr{F}|U)_{/(U\cap X')}=(\mathscr{F}_{/X'})|(U\cap X')$.

사영극한을 취하면 $(\mathscr{O}_X)_{/X'}$가 환의 층이고
$\mathscr{F}_{/X'}$를 $(\mathscr{O}_X)_{/X'}$-가군으로 볼 수 있음은
명백하다. 또한 $X'$의 위상에는 준콤팩트 열린집합들로 이루어진 기저가
있으므로, $(\mathscr{O}_X)_{/X'}$ (각각 $\mathscr{F}_{/X'}$)를
\emph{의사이산} 환의 층들 (각각 의사이산 군의 층들)
$\mathscr{O}_X/\mathscr{J}$ (각각
$\mathscr{F}\otimes_{\mathscr{O}_X}
(\mathscr{O}_X/\mathscr{J})=\mathscr{F}/\mathscr{J}\mathscr{F}$)의
사영극한인 \emph{위상환의 층} (각각 \emph{위상군의 층})으로 볼 수
있다. 이때 사영극한을 취하면 $\mathscr{F}_{/X'}$는
\emph{위상 $(\mathscr{O}_X)_{/X'}$-가군}이 된다
(\hyperref[0.3.8.1-ko]{0, 3.8.1}과
\hyperref[0.3.8.2-ko]{3.8.2}). 모든 \emph{준콤팩트} 열린집합
$U\subset X$에 대하여
$\Gamma(U\cap X',(\mathscr{O}_X)_{/X'})$ (각각
$\Gamma(U\cap X',\mathscr{F}_{/X'})$)는 이산환들 (각각 이산군들)
$\Gamma(U,\mathscr{O}_X/\mathscr{J})$ (각각
$\Gamma(U,\mathscr{F}/\mathscr{J}\mathscr{F})$)의 사영극한임을
상기하자.

이제 $u:\mathscr{F}\to\mathscr{G}$가 $\mathscr{O}_X$-가군 사이의
준동형이면, 모든 $\mathscr{J}\in\Phi$에 대하여 표준적으로 준동형
$u_{\mathscr{J}}:
\mathscr{F}\otimes_{\mathscr{O}_X}(\mathscr{O}_X/\mathscr{J})
\to
\mathscr{G}\otimes_{\mathscr{O}_X}(\mathscr{O}_X/\mathscr{J})$를 얻으며,
이 준동형들은 사영계를 이룬다. 사영극한을 취하고 $X'$에 제한하면
연속 $(\mathscr{O}_X)_{/X'}$-준동형
$\mathscr{F}_{/X'}\to\mathscr{G}_{/X'}$를 얻는다. 이를
$u_{/X'}$ 또는 $\widehat u$로 나타내며,
\emph{$X'$을 따른 준동형 $u$의 완비화}라 한다. 또한
$v:\mathscr{G}\to\mathscr{H}$가 $\mathscr{O}_X$-가군 사이의
또 하나의 준동형이면
$(v\circ u)_{/X'}=(v_{/X'})\circ(u_{/X'})$임은 명백하다. 따라서
$\mathscr{F}_{/X'}$는 $\mathscr{F}$에 관한
\emph{가법적 공변 함자}이며, 연접 $\mathscr{O}_X$-가군의 범주에서
위상 $(\mathscr{O}_X)_{/X'}$-가군의 범주에 값을 갖는다.

\begin{proposition}[10.8.5]
\label{I.10.8.5-ko}
$(\mathscr{O}_X)_{/X'}$의 지지집합은 $X'$이다. 위상환 달린 공간
$(X',(\mathscr{O}_X)_{/X'})$는 국소 뇌터 형식적 준스킴이며, 모든
$\mathscr{J}\in\Phi$에 대하여 $\mathscr{J}_{/X'}$는 이 형식적
준스킴의 정의 아이디얼이다. $X=\operatorname{Spec}(A)$가 뇌터 환을
갖는 아핀 스킴이고, $\mathscr{J}=\widetilde{\mathfrak{J}}$이며,
$\mathfrak{J}$가 $A$의 아이디얼이고 $X'=V(\mathfrak{J})$이면,
$(X',(\mathscr{O}_X)_{/X'})$는 표준적으로
$\operatorname{Spf}(\widehat A)$와 동일시된다. 여기서 $\widehat A$는
$A$에 $\mathfrak{J}$-준아딕 위상을 주어 얻는 분리 완비화이다.
\end{proposition}

\begin{proof}
마지막 주장만 증명하면 충분함은 명백하다.
(\hyperref[0.7.3.3-ko]{0, 7.3.3})에 의해, $\mathfrak{J}$의
$\mathfrak{J}$-준아딕 위상에 관한 분리 완비화
$\widehat{\mathfrak{J}}$는 $\widehat A$의 아이디얼
$\mathfrak{J}\widehat A$와 동일시된다. 또한 $\widehat A$는
뇌터 $\widehat{\mathfrak{J}}$-아딕환이며
$\widehat A/\widehat{\mathfrak{J}}^n=A/\mathfrak{J}^n$임을 안다
(\hyperref[0.7.2.6-ko]{0, 7.2.6}). 이 마지막 관계에서
$\widehat A$의 열린 소 아이디얼들은
$\widehat{\mathfrak p}=\mathfrak p\widehat A$ 꼴의 아이디얼들임을 알 수
있다. 여기서 $\mathfrak p$는 $A$의 소 아이디얼이고
$\mathfrak{J}$를 포함하며, $\widehat{\mathfrak p}\cap A=\mathfrak p$이다. 따라서
바탕공간으로서 $\operatorname{Spf}(\widehat A)=X'$이다. 또한
$\mathscr{O}_X/\mathscr{J}^n=(A/\mathfrak{J}^n)^{\sim}$이므로, 명제는
정의들에서 곧바로 따른다.
\end{proof}

이렇게 정의한 형식적 준스킴을 \emph{$X'$을 따른 $X$의 완비화}라 하고
$X_{/X'}$로 나타내며, 혼동의 우려가 없으면 $\widehat X$로도 나타낸다.
$X'=X$로 잡을 때에는 $\mathscr{J}=0$을 택할 수 있으므로
$X_{/X}=X$이다.

$U$가 $X$의 열린집합 위에 유도된 부분준스킴이면,
$U_{/(U\cap X')}$는 $X_{/X'}$가 $X'$의 열린집합 $U\cap X'$ 위에
유도하는 형식적 부분준스킴과 표준적으로 동일시됨은 명백하다.

\begin{corollary}[10.8.6]
\label{I.10.8.6-ko}
(보통 의미의) 준스킴 $\widehat X_{\mathrm{red}}$는 $X$의 유일한 축소
부분준스킴으로서, 그 바탕공간은 $X'$이다
(\hyperref[I.5.2.1-ko]{5.2.1}). $\widehat X$가 뇌터이기 위한
필요충분조건은 $\widehat X_{\mathrm{red}}$가 뇌터인 것이며, 더 나아가
$X$가 뇌터이면 전자도 뇌터이다.
\end{corollary}

\oldpage[I]{196}

$\widehat X_{\mathrm{red}}$의 결정은 국소적이므로
(\hyperref[I.10.5.4-ko]{10.5.4}), 다시 $X$가 뇌터 환을 갖는 아핀 스킴인
경우로 한정할 수 있다. (\hyperref[I.10.8.5-ko]{10.8.5})의 표기를
쓰면, 아이디얼 $\mathfrak{T}$는 $\widehat A$의 위상적으로 멱영인
원소들로 이루어지며, 표준 사상
$\widehat A\to\widehat A/\widehat{\mathfrak{J}}=A/\mathfrak{J}$에 의한
$A/\mathfrak{J}$의 멱영근기의 역상이다
(\hyperref[0.7.1.3-ko]{0, 7.1.3}). 따라서
$\widehat A/\mathfrak{T}$는 $A/\mathfrak{J}$를 그 멱영근기로 나눈
몫과 동형이다. 그러므로 첫 번째 주장은
(\hyperref[I.10.5.4-ko]{10.5.4})와
(\hyperref[I.5.1.1-ko]{5.1.1})에서 따른다.
$\widehat X_{\mathrm{red}}$가 뇌터이면 그 바탕공간 $X'$도 뇌터이므로,
$X'_n=\operatorname{Spec}(\mathscr{O}_X/\mathscr{J}^n)$은 뇌터이고
(\hyperref[I.6.1.2-ko]{6.1.2}), $\widehat X$도 뇌터이다
(\hyperref[I.10.6.4-ko]{10.6.4}). 역은
(\hyperref[I.6.1.2-ko]{6.1.2})에 의해 명백하다.

\begin{env}[10.8.7]
\label{I.10.8.7-ko}
표준 준동형 $\mathscr{O}_X\to\mathscr{O}_X/\mathscr{J}$
($\mathscr{J}\in\Phi$)들은 사영계를 이루므로, 사영극한을 취하면 환의
층의 준동형
$\theta:\mathscr{O}_X\to\psi_*((\mathscr{O}_X)_{/X'})
=\varprojlim_\Phi(\mathscr{O}_X/\mathscr{J})$를 얻는다. 여기서 $\psi$는
바탕공간 사이의 표준 포함 사상 $X'\to X$이다. 환 달린 공간의 사상
\[
  (\psi,\theta):X_{/X'}\to X
\]
을 $i$ (또는 $i_X$)로 나타내고 표준 사상이라 하겠다.

텐서곱을 취하면, 모든 연접 $\mathscr{O}_X$-가군 $\mathscr{F}$에 대하여
표준 준동형 $\mathscr{O}_X\to\mathscr{O}_X/\mathscr{J}$들로부터
$\mathscr{O}_X$-가군 준동형
$\mathscr{F}\to
\mathscr{F}\otimes_{\mathscr{O}_X}(\mathscr{O}_X/\mathscr{J})$들을 얻으며,
이들도 사영계를 이룬다. 따라서 사영극한을 취하면 함자적인 표준
$\mathscr{O}_X$-가군 준동형
$\gamma:\mathscr{F}\to\psi_*(\mathscr{F}_{/X'})$를 얻는다.
\end{env}

\begin{proposition}[10.8.8]
\label{I.10.8.8-ko}
\medskip\noindent
\begin{enumerate}
  \item[\rm{(i)}] 함자 $\mathscr{F}_{/X'}$는
    ($\mathscr{F}$에 관하여) 완전하다.
  \item[\rm{(ii)}] 함자적인 준동형
    $\gamma^\sharp:i^*(\mathscr{F})\to\mathscr{F}_{/X'}$는
    $(\mathscr{O}_X)_{/X'}$-가군의 동형이다.
\end{enumerate}
\end{proposition}

\medskip\noindent
\begin{enumerate}
  \item[\rm{(i)}] 다음을 증명하면 충분하다.
    $0\to\mathscr{F}'\to\mathscr{F}\to\mathscr{F}''\to0$이 연접
    $\mathscr{O}_X$-가군의 완전열이고, $U$가 $X$의 아핀 열린집합이며
    그 환이 뇌터 환 $A$이면, 다음 열은 완전하다.
    \[
      0\to\Gamma(U\cap X',\mathscr{F}'_{/X'})
      \to\Gamma(U\cap X',\mathscr{F}_{/X'})
      \to\Gamma(U\cap X',\mathscr{F}''_{/X'})\to0
    \]
    이때 $\mathscr{F}|U=\widetilde M$,
    $\mathscr{F}'|U=\widetilde{M'}$, $\mathscr{F}''|U=\widetilde{M''}$이고,
    여기서 $M$, $M'$, $M''$는 세 유한 생성 $A$-가군이며 열
    $0\to M'\to M\to M''\to0$은 완전하다
    (\hyperref[I.1.5.1-ko]{1.5.1}과
    \hyperref[I.1.3.11-ko]{1.3.11})~; $\mathscr{J}\in\Phi$라 하고,
    $\mathfrak{J}$를 $A$의 아이디얼로서
    $\mathscr{J}|U=\widetilde{\mathfrak{J}}$인 것으로 택하자. 그러면
    \[
      \Gamma(U\cap X',
      \mathscr{F}\otimes_{\mathscr{O}_X}(\mathscr{O}_X/\mathscr{J}^n))
      =M\otimes_A(A/\mathfrak{J}^n)
    \]
    이다(\hyperref[I.1.3.12-ko]{1.3.12})~; 따라서 사영극한의 정의에 의해
    \[
      \Gamma(U\cap X',\mathscr{F}_{/X'})
      =\varprojlim_n(M\otimes_A(A/\mathfrak{J}^n))=\widehat M
    \]
    이고, 이는 $M$의 $\mathfrak{J}$-준아딕 위상에 관한 분리
    완비화이다. 마찬가지로
    \[
      \Gamma(U\cap X',\mathscr{F}'_{/X'})=\widehat{M'},\qquad
      \Gamma(U\cap X',\mathscr{F}''_{/X'})=\widehat{M''}~;
    \]
    이제 $A$가 뇌터이면 함자 $\widehat M$이 $M$에 관하여 유한 생성
    $A$-가군의 범주에서 완전하다는 사실에서 우리의 주장이 따른다
    (\hyperref[0.7.3.3-ko]{0, 7.3.3}).

\oldpage[I]{197}
  \item[\rm{(ii)}] 문제는 국소적이므로 완전열
    $\mathscr{O}_X^m\to\mathscr{O}_X^n\to\mathscr{F}\to0$이
    있다고 가정할 수 있다
    (\hyperref[0.5.3.2-ko]{0, 5.3.2})~; $\gamma^\sharp$가 함자적이고,
    함자 $i^*(\mathscr{F})$와 $\mathscr{F}_{/X'}$가 각각
    (\hyperref[0.4.3.1-ko]{0, 4.3.1})과 (i)에 의해 오른쪽 완전하므로
    다음 가환도를 얻는다.
    \[
      \xymatrix{
        i^*(\mathscr{O}_X^m)\ar[r]\ar[d]_{\gamma^\sharp} &
        i^*(\mathscr{O}_X^n)\ar[r]\ar[d]_{\gamma^\sharp} &
        i^*(\mathscr{F})\ar[r]\ar[d]_{\gamma^\sharp} &
        0\\
        (\mathscr{O}_X^m)_{/X'}\ar[r] &
        (\mathscr{O}_X^n)_{/X'}\ar[r] &
        \mathscr{F}_{/X'}\ar[r] &
        0
      }
      \tag{10.8.8.1}\label{I.10.8.8.1-ko}
    \]
    이 가환도의 두 행은 완전하다. 또한 두 함자
    $i^*(\mathscr{F})$와 $\mathscr{F}_{/X'}$는 유한 직합과 가환하므로
    (\hyperref[0.3.2.6-ko]{0, 3.2.6}과
    \hyperref[0.4.3.2-ko]{4.3.2}), 우리의 주장을
    $\mathscr{F}=\mathscr{O}_X$인 경우에 증명하는 것으로 환원된다.
    이때
    $i^*(\mathscr{O}_X)=(\mathscr{O}_X)_{/X'}=
    \mathscr{O}_{\widehat X}$
    이고(\hyperref[0.4.3.4-ko]{0, 4.3.4}), $\gamma^\sharp$는
    $\mathscr{O}_{\widehat X}$-가군 준동형이다. 따라서
    $\gamma^\sharp$가 $\mathscr{O}_{\widehat X}$의 단위원 단면을
    $X'$의 한 열린집합 위에서 자기 자신으로 보내는지 확인하면 충분하다.
    이는 명백하며, 따라서 이 경우 $\gamma^\sharp$는 항등준동형이다.
\end{enumerate}

\begin{corollary}[10.8.9]
\label{I.10.8.9-ko}
환 달린 공간의 사상 $i:X_{/X'}\to X$는 평탄하다.
\end{corollary}

이는 실제로 (\hyperref[0.6.7.3-ko]{0, 6.7.3})과
(\hyperref[I.10.8.8-ko]{10.8.8, (i)})에서 따른다.

\begin{corollary}[10.8.10]
\label{I.10.8.10-ko}
$\mathscr{F}$와 $\mathscr{G}$가 연접 $\mathscr{O}_X$-가군이면, 다음
표준 동형사상들이 존재하며 $\mathscr{F}$와 $\mathscr{G}$에 관하여 함자적이다.
\[
  (\mathscr{F}_{/X'})\otimes_{(\mathscr{O}_X)_{/X'}}
  (\mathscr{G}_{/X'})
  \xrightarrow{\sim}
  (\mathscr{F}\otimes_{\mathscr{O}_X}\mathscr{G})_{/X'}
  \tag{10.8.10.1}\label{I.10.8.10.1-ko}
\]
\[
  (\mathscr{H}\!om_{\mathscr{O}_X}(\mathscr{F},\mathscr{G}))_{/X'}
  \xrightarrow{\sim}
  \mathscr{H}\!om_{(\mathscr{O}_X)_{/X'}}
  (\mathscr{F}_{/X'},\mathscr{G}_{/X'})
  \tag{10.8.10.2}\label{I.10.8.10.2-ko}
\]
\end{corollary}

이는 $i^*(\mathscr{F})$와 $\mathscr{F}_{/X'}$의 표준 동일시에서 따른다~;
이때 첫 번째 동형의 존재는 모든 환 달린 공간의 사상에 대하여 성립하는
결과이고 (\hyperref[0.4.3.3.1-ko]{0, 4.3.3.1}), 두 번째 동형의 존재는
모든 평탄 사상에 대하여 성립하는 결과이므로
(\hyperref[0.6.7.6-ko]{0, 6.7.6}),
(\hyperref[I.10.8.9-ko]{10.8.9})에서 따른다.

\begin{proposition}[10.8.11]
\label{I.10.8.11-ko}
모든 연접 $\mathscr{O}_X$-가군 $\mathscr{F}$에 대하여, 표준 준동형
$\Gamma(X,\mathscr{F})\to\Gamma(X',\mathscr{F}_{/X'})$는
$\mathscr{F}\to\mathscr{F}_{/X'}$에서 유도되며, 그 핵은 $X'$의 어떤
근방에서 영인 단면들로 이루어진다.
\end{proposition}

그러한 단면의 표준 상이 영임은 $\mathscr{F}_{/X'}$의 정의에서 따른다.
역으로 $s\in\Gamma(X,\mathscr{F})$의 상이
$\Gamma(X',\mathscr{F}_{/X'})$에서 영이라고 하자. 모든 $x\in X'$가
$X$ 안에 하나의 근방을 가져 그곳에서 $s$가 영임을 보이면 충분하고, 따라서
$X=\operatorname{Spec}(A)$가 아핀 스킴이고 $A$가 뇌터 환이며,
$X'=V(\mathfrak{J})$이고 $\mathfrak{J}$가 $A$의 아이디얼이며,
$\mathscr{F}=\widetilde M$이고 $M$이 유한 생성 $A$-가군인 경우로
환원할 수 있다. 이때 $\Gamma(X',\mathscr{F}_{/X'})$는
분리 완비화 $\widehat M$이고, 이는 $M$을 $\mathfrak{J}$-준아딕 위상에
관하여 분리 완비화한 것이다. 또한
준동형 $\Gamma(X,\mathscr{F})\to\Gamma(X',\mathscr{F}_{/X'})$는
표준 준동형 $M\to\widehat M$이다. 이 준동형의 핵은
$z\in M$ 가운데 $1+\mathfrak{J}$의 한 원소에 의해 소멸되는 것들의 집합임이
알려져 있다 (\hyperref[0.7.3.7-ko]{0, 7.3.7}). 따라서
$(1+f)s=0$인 어떤 $f\in\mathfrak{J}$가 존재한다~; 모든 $x\in X'$에
대하여 이로부터 $(1_x+f_x)s_x=0$을 얻고,
$1_x+f_x$는 $\mathscr{O}_x$에서 가역인데(실제로
$\mathfrak{J}_x\mathscr{O}_x$는 $\mathscr{O}_x$의 극대 아이디얼에
포함된다). 따라서
$s_x=0$이고, 이는 명제를 증명한다.

\begin{corollary}[10.8.12]
\label{I.10.8.12-ko}
$\mathscr{F}_{/X'}$의 지지집합은
$\operatorname{Supp}(\mathscr{F})\cap X'$과 같다.
\end{corollary}

$\mathscr{F}_{/X'}$가 유한 생성 $(\mathscr{O}_X)_{/X'}$-가군임은
(\hyperref[I.10.8.8-ko]{10.8.8, (ii)})와
(\hyperref[0.5.2.4-ko]{0, 5.2.4})에서 명백하므로,

\oldpage[I]{198}

그 지지집합은 닫혀 있고
(\hyperref[0.5.2.2-ko]{0, 5.2.2}), 명백히
$\operatorname{Supp}(\mathscr{F})\cap X'$에 포함된다. 이 마지막
집합과 같음을 보이기 위해서는 등식
$\Gamma(X',\mathscr{F}_{/X'})=0$이
$\operatorname{Supp}(\mathscr{F})\cap X'=\varnothing$을 함의함을
증명하는 것으로 즉시 환원된다~; 그런데 이는
(\hyperref[I.10.8.11-ko]{10.8.11})과
(\hyperref[I.1.4.1-ko]{1.4.1})에서 따른다.

\begin{corollary}[10.8.13]
\label{I.10.8.13-ko}
$u:\mathscr{F}\to\mathscr{G}$를 연접 $\mathscr{O}_X$-가군 사이의 준동형이라
하자. $u_{/X'}:\mathscr{F}_{/X'}\to\mathscr{G}_{/X'}$가 영이기 위한
필요충분조건은 $u$가 $X'$의 어떤 근방에서 영 준동형인 것이다.
\end{corollary}

실제로 (\hyperref[I.10.8.8-ko]{10.8.8, (ii)})에 따라 $u_{/X'}$는
$i^*(u)$와 동일시된다. 따라서 $u$를 $X$ 위에서 층
$\mathscr{H}=\mathscr{H}\!om_{\mathscr{O}_X}(\mathscr{F},\mathscr{G})$의
단면으로 보면, $u_{/X'}$는
$i^*(\mathscr{H})=\mathscr{H}_{/X'}$의 $X'$ 위의 단면이며 앞서 본
단면에 표준적으로 대응한다
((\hyperref[I.10.8.10.2-ko]{10.8.10.2})와
(\hyperref[0.4.4.6-ko]{0, 4.4.6})). 그러므로 연접
$\mathscr{O}_X$-가군 $\mathscr{H}$에
(\hyperref[I.10.8.11-ko]{10.8.11})을 적용하면 충분하다.

\begin{corollary}[10.8.14]
\label{I.10.8.14-ko}
$u:\mathscr{F}\to\mathscr{G}$를 연접 $\mathscr{O}_X$-가군 사이의
준동형이라 하자. $u_{/X'}$가 단사(각각 전사)이기 위한 필요충분조건은
$u$가 $X'$의 어떤 근방에서 단사(각각 전사)인 것이다.
\end{corollary}

$\mathscr{P}$와 $\mathscr{N}$을 각각 $u$의 여핵과 핵이라 하자.
그러면 다음 완전열이 성립한다.
\[
  0\to\mathscr{N}\xrightarrow{v}\mathscr{F}\xrightarrow{u}
  \mathscr{G}\xrightarrow{w}\mathscr{P}\to0,
\]
따라서 (\hyperref[I.10.8.8-ko]{10.8.8, (i)})에 의해 다음 완전열이
성립한다.
\[
  0\to\mathscr{N}_{/X'}\xrightarrow{v_{/X'}}
  \mathscr{F}_{/X'}\xrightarrow{u_{/X'}}
  \mathscr{G}_{/X'}\xrightarrow{w_{/X'}}
  \mathscr{P}_{/X'}\to0.
\]
$u_{/X'}$가 단사(각각 전사)이면 $v_{/X'}=0$ (각각
$w_{/X'}=0$)이다. 따라서 (\hyperref[I.10.8.13-ko]{10.8.13})에 의해
$X'$의 어떤 근방에서는 $v=0$ (각각 $w=0$)이다.

\subsection{완비화로의 사상의 연장.}
\label{subsection:I.10.9-ko}

\begin{env}[10.9.1]
\label{I.10.9.1-ko}
$X$, $Y$를 (보통 의미의) 국소 뇌터 준스킴이라 하고, $f:X\to Y$를
사상이라 하며, $X'$ (각각 $Y'$)를 $X$ (각각 $Y$)의 바탕공간의 닫힌
부분집합이라 하자. $f(X')\subset Y'$라고 가정한다. $\mathscr{J}$
(각각 $\mathscr{K}$)를 $\mathscr{O}_X$ (각각 $\mathscr{O}_Y$)의
아이디얼층이라 하고, $\mathscr{O}_X/\mathscr{J}$ (각각
$\mathscr{O}_Y/\mathscr{K}$)의 지지집합은 $X'$ (각각 $Y'$)이며
$f^*(\mathscr{K})\mathscr{O}_X\subset\mathscr{J}$라고 하자. 이러한
아이디얼층은 항상 존재한다는 점을 지적해 두자. 예를 들어 $\mathscr{J}$로
$\mathscr{O}_X$의 아이디얼층 가운데 $X$의 부분준스킴을 정의하고
$X'$을 바탕공간으로 갖는 가장 큰 것을 택할 수 있고
(\hyperref[I.5.2.1-ko]{5.2.1}),
$f(X')\subset Y'$라는 가정에서
$f^*(\mathscr{K})\mathscr{O}_X\subset\mathscr{J}$가 따른다
(\hyperref[I.5.2.4-ko]{5.2.4}). 따라서 모든 정수 $n>0$에 대하여
$f^*(\mathscr{K}^n)\mathscr{O}_X\subset\mathscr{J}^n$
(\hyperref[0.4.3.5-ko]{0, 4.3.5})이고, 따라서
(\hyperref[I.4.4.6-ko]{4.4.6})에 의해 다음과 같이 놓으면
\[
  X'_n=(X',\mathscr{O}_X/\mathscr{J}^{n+1}),\qquad
  Y'_n=(Y',\mathscr{O}_Y/\mathscr{K}^{n+1}),
\]
$f$로부터 사상 $f_n:X'_n\to Y'_n$을 얻는다. 또한 $f_n$들은
귀납계를 이룬다. 그 귀납극한
(\hyperref[I.10.6.8-ko]{10.6.8})을
$\widehat f:X_{/X'}\to Y_{/Y'}$로 나타내고, (표현을 다소 남용하여)
$\widehat f$를 $f$의, $X$와 $Y$를 각각 $X'$와 $Y'$를 따라 완비화한
것들 사이로의 연장이라 부르겠다. 이 사상이 위 조건을
만족하는 아이디얼층 $\mathscr{J}$, $\mathscr{K}$의 선택에 의존하지
않음은 즉시 확인할 수 있다. 실제로 이를 확인하기 위해서는 $X$, $Y$가
각각 환 $A$, $B$의 아핀 뇌터 준스킴인 경우만 보면 충분하다. 이때
$\mathscr{J}=\widetilde{\mathfrak{J}}$,
$\mathscr{K}=\widetilde{\mathfrak{K}}$이고, $\mathfrak{J}$ (각각
$\mathfrak{K}$)는 $A$ (각각 $B$)의 아이디얼이며, $f$는
$\varphi:B\to A$라는 환 준동형에 대응하고
$\varphi(\mathfrak{K})\subset\mathfrak{J}$이다
(\hyperref[I.4.4.6-ko]{4.4.6} 및
\hyperref[I.1.7.4-ko]{1.7.4}). 이때
\emph{$\widehat f$는 (\hyperref[I.10.2.2-ko]{10.2.2})에 따라 연속
준동형 $\widehat\varphi:\widehat B\to\widehat A$에 대응하는
사상이다}. 여기서 $\widehat A$ (각각 $\widehat B$)는 $A$ (각각
$B$)에 $\mathfrak{J}$-준아딕 (각각 $\mathfrak{K}$-준아딕) 위상을
주어 얻는 분리 완비화이다
(\hyperref[I.10.6.8-ko]{10.6.8}). 또한
$\mathscr{J}$를 다른
\oldpage[I]{199}
아이디얼층 $\mathscr{J}'=\widetilde{\mathfrak{J}'}$로 바꾸어도,
$\mathscr{O}_X/\mathscr{J}'$의 지지집합이 여전히 $X'$이면,
$\mathfrak{J}$-준아딕 위상과 $\mathfrak{J}'$-준아딕 위상은 $A$ 위에서
같음을 알고 있다
(\hyperref[I.10.8.2-ko]{10.8.2}).
이 정의에 따르면, $X_{/X'}$와 $Y_{/Y'}$의 바탕공간 사이의
연속사상 $X'\to Y'$는 $\widehat f$에 대응하며, 다름 아닌 $X'$에 대한
$f$의 제한임을 지적하자.
\end{env}

\begin{env}[10.9.2]
\label{I.10.9.2-ko}
앞의 정의로부터 환 달린 공간의 사상들로 이루어진 도식
\[
  \label{I.10.9.2.1-ko}
  \xymatrix{
    \widehat X\ar[r]^{\widehat f}\ar[d]_{i_X} &
    \widehat Y\ar[d]^{i_Y}\\
    X\ar[r]_f &
    Y
  }
\]
이 가환이고, 세로 화살표들은 표준 사상
(\hyperref[I.10.8.7-ko]{10.8.7})임이 곧바로 따른다.
\end{env}

\begin{env}[10.9.3]
\label{I.10.9.3-ko}
$Z$를 세 번째 준스킴, $g:Y\to Z$를 사상, $Z'$을
$g(Y')\subset Z'$인 $Z$의 닫힌 부분집합이라 하자. $Y'$과 $Z'$을
따른 사상 $g$의 완비화를 $\widehat g$로 나타내면,
(\hyperref[I.10.9.1-ko]{10.9.1})로부터
$\widehat{(g\circ f)}=\widehat g\circ\widehat f$임이 곧바로 따른다.
\end{env}

\begin{proposition}[10.9.4]
\label{I.10.9.4-ko}
$X$, $Y$를 두 국소 뇌터 $S$-준스킴이라 하고, $Y$가 $S$ 위의
유한형이라고 하자. $f$, $g$를 $X$에서 $Y$로 가는 두 $S$-사상이라
하고, $f(X')\subset Y'$, $g(X')\subset Y'$라고 하자.
$\widehat f=\widehat g$일 필요충분조건은 $f$와 $g$가 $X'$의 어떤
근방에서 일치하는 것이다.
\end{proposition}

$Y$에 대한 유한성 가정 없이도 이 조건은 자명하게 충분하다.
필요성을 보기 위해, 먼저 가정 $\widehat f=\widehat g$에서 모든
$x\in X'$에 대하여 $f(x)=g(x)$가 따름을 지적하자. 한편 문제는
국소적이므로, $X$와 $Y$가 각각 $x$와 $y=f(x)=g(x)$의 아핀 열린
근방이고 그 좌표환들이 뇌터이며, $S$가 아핀이고
$\Gamma(Y,\mathscr{O}_Y)$가 $\Gamma(S,\mathscr{O}_S)$ 위의 유한 생성
대수라고 가정할 수 있다
(\hyperref[I.6.3.3-ko]{6.3.3}). 이때 $f$와 $g$는
$\Gamma(Y,\mathscr{O}_Y)$에서 $\Gamma(X,\mathscr{O}_X)$로 가는 두
$\Gamma(S,\mathscr{O}_S)$-대수 준동형 $\rho$, $\sigma$에 대응한다
(\hyperref[I.1.7.3-ko]{1.7.3}). 또한 가정에 따라 이 준동형들을
$\Gamma(Y,\mathscr{O}_Y)$의 분리 완비화에 연속적으로 연장한
준동형들은 서로 같다. 따라서 (\hyperref[I.10.8.11-ko]{10.8.11})에서
모든 단면 $s\in\Gamma(Y,\mathscr{O}_Y)$에 대하여 단면 $\rho(s)$와
$\sigma(s)$가 $X'$의 어떤 근방에서 일치한다는 결론을 얻는다(이
근방은 $s$에 의존한다)~; $\Gamma(Y,\mathscr{O}_Y)$가
$\Gamma(S,\mathscr{O}_S)$ 위의 유한 생성 대수이므로, $X'$의 한 근방
$V$가 존재하여 \emph{모든} 단면 $s\in\Gamma(Y,\mathscr{O}_Y)$에
대하여 $\rho(s)$와 $\sigma(s)$가 $V$에서 일치함이 곧바로 따른다.
$D(h)$가 $V$에 포함되는 $x$의 근방이 되도록
$h\in\Gamma(X,\mathscr{O}_X)$를 택하면, 앞의 결론과
(\hyperref[I.1.4.1-ko]{1.4.1, d)})에서 $f$와 $g$가 $D(h)$에서
일치함을 얻는다.\footnote{\textbf{번역 주.} 인쇄 원문은 여기서 형에
맞지 않는 $h\in\Gamma(Y,\mathscr{O}_X)$라고 쓴다. 그러나
$\mathscr{O}_X$는 $X$ 위의 층이고 이어지는 $D(h)$는 $X$ 안에서
$x$의 근방이므로, $h\in\Gamma(X,\mathscr{O}_X)$로 바로잡았다.}

\begin{proposition}[10.9.5]
\label{I.10.9.5-ko}
(\hyperref[I.10.9.1-ko]{10.9.1})의 가정 아래에서, 모든 연접
$\mathscr{O}_Y$-가군 $\mathscr{G}$에 대하여 다음과 같은
$(\mathscr{O}_X)_{/X'}$-가군의 함자적인 표준 동형이 존재한다.
\[
  (f^*(\mathscr{G}))_{/X'}\xrightarrow{\sim}
  \widehat f^*(\mathscr{G}_{/Y'})
\]
\end{proposition}

$(f^*(\mathscr{G}))_{/X'}$를 $i_X^*(f^*(\mathscr{G}))$와,
$\widehat f^*(\mathscr{G}_{/Y'})$를
$\widehat f^*(i_Y^*(\mathscr{G}))$와 표준적으로 동일시하면
(\hyperref[I.10.8.8-ko]{10.8.8}), 명제는
(\hyperref[I.10.9.2-ko]{10.9.2})의 도식이 가환이라는 사실에서 곧바로
따른다.

\begin{env}[10.9.6]
\label{I.10.9.6-ko}
이제 $\mathscr{F}$를 연접 $\mathscr{O}_X$-가군,
$\mathscr{G}$를 연접 $\mathscr{O}_Y$-가군이라 하자.
$u:\mathscr{G}\to\mathscr{F}$가 $\mathscr{G}$에서 $\mathscr{F}$로
가는 $f$-사상이면, 이에 대응하는 $\mathscr{O}_X$-준동형
$u^\sharp:f^*(\mathscr{G})\to\mathscr{F}$가 있다. 따라서 이 준동형을
완비화하면 연속 $(\mathscr{O}_X)_{/X'}$-준동형
$(u^\sharp)_{/X'}:(f^*(\mathscr{G}))_{/X'}\to\mathscr{F}_{/X'}$를
얻고, (\hyperref[I.10.9.5-ko]{10.9.5})에 의해
$\widehat f$-사상 $v:\mathscr{G}_{/Y'}\to\mathscr{F}_{/X'}$가 존재하고,
\oldpage[I]{200}

유일하며, $v^\sharp=(u^\sharp)_{/X'}$를 만족한다. 삼중항
$(\mathscr{F},X,X')$들($\mathscr{F}$는 연접 $\mathscr{O}_X$-가군이고
$X'$은 $X$의 닫힌 부분집합이다)을 하나의 \emph{범주}의 대상으로
생각하고, 사상
$(\mathscr{F},X,X')\to(\mathscr{G},Y,Y')$들이
$f(X')\subset Y'$인 준스킴의 사상 $f:X\to Y$와
$f$-사상 $u:\mathscr{G}\to\mathscr{F}$로 이루어진다고 하자. 그러면
$(X_{/X'},\mathscr{F}_{/X'})$는 $(\mathscr{F},X,X')$에 관한
\emph{함자}라고 말할 수 있다. 이 함자는 국소 뇌터 형식적 준스킴
$\mathfrak{Z}$와 $\mathscr{O}_{\mathfrak{Z}}$-가군 $\mathscr{H}$로
이루어진 쌍 $(\mathfrak{Z},\mathscr{H})$들의 범주에 값을 갖는다.
후자의 범주에서 사상은 형식적 준스킴의 사상 $g$와 $g$-사상으로
이루어진 쌍이다.
\end{env}

\begin{proposition}[10.9.7]
\label{I.10.9.7-ko}
$S$, $X$, $Y$를 세 국소 뇌터 준스킴,
$g:X\to S$, $h:Y\to S$를 두 사상, $S'$을 $S$의 닫힌 부분집합,
$X'$ (각각 $Y'$)을 $g(X')\subset S'$ (각각
$h(Y')\subset S'$)을 만족하는 $X$ (각각 $Y$)의 닫힌 부분집합이라
하자~; $Z=X\times_S Y$라 하고, $Z$가 국소 뇌터라고 가정하며,
$Z'=p^{-1}(X')\cap q^{-1}(Y')$라 하자. 여기서 $p$와 $q$는
$X\times_S Y$의 사영이다. 이 조건에서 완비화 $Z_{/Z'}$는
$S_{/S'}$-형식적 준스킴들의 곱
$(X_{/X'})\times_{S_{/S'}}(Y_{/Y'})$와 동일시된다. 구조 사상들은
각각 $\widehat g$, $\widehat h$와 동일시되고, 사영들은 각각
$\widehat p$, $\widehat q$와 동일시된다.
\end{proposition}

문제가 $S$, $X$, $Y$에 관하여 국소적임은 명백하므로,
$S=\operatorname{Spec}(A)$,
$X=\operatorname{Spec}(B)$, $Y=\operatorname{Spec}(C)$,
$S'=V(\mathfrak{J})$, $X'=V(\mathfrak{K})$,
$Y'=V(\mathfrak{L})$인 경우로 환원된다. 여기서
$\mathfrak{J}$, $\mathfrak{K}$, $\mathfrak{L}$은
$\varphi(\mathfrak{J})\subset\mathfrak{K}$ 및
$\psi(\mathfrak{J})\subset\mathfrak{L}$을 만족하는 세 아이디얼이고,
$\varphi:A\to B$와 $\psi:A\to C$는 각각 $g$와 $h$에 대응하는
준동형이다. 이때 $Z=\operatorname{Spec}(B\otimes_A C)$이고
$Z'=V(\mathfrak{M})$임을 알고 있다. 여기서 $\mathfrak{M}$은
아이디얼
$\operatorname{Im}(\mathfrak{K}\otimes_A C)
+\operatorname{Im}(B\otimes_A\mathfrak{L})$이다.
(\hyperref[I.10.7.2-ko]{10.7.2})에 의해, 결론은 완비 텐서곱
$(\widehat B\otimes_{\widehat A}\widehat C)^\wedge$
(여기서 $\widehat A$, $\widehat B$, $\widehat C$는 각각
$A$, $B$, $C$에 $\mathfrak{J}$-, $\mathfrak{K}$-,
$\mathfrak{L}$-준아딕 위상을 주어 얻는 분리 완비화이다)이
텐서곱 $B\otimes_A C$의 $\mathfrak{M}$-준아딕 위상에 관한
분리 완비화라는 사실에서 따른다
(\hyperref[0.7.7.2-ko]{0, 7.7.2}).

또한 $T$가 국소 뇌터 $S$-준스킴이고,
$u:T\to X$, $v:T\to Y$가 두 $S$-사상이며, $T'$이
$u(T')\subset X'$, $v(T')\subset Y'$을 만족하는 $T$의 닫힌
부분집합이면, 완비화한 것들 사이로의 연장
$((u,v)_S)^\wedge$는
$(\widehat u,\widehat v)_{S_{/S'}}$와 동일시됨을 유의하자.

\begin{corollary}[10.9.8]
\label{I.10.9.8-ko}
$X$, $Y$를 두 국소 뇌터 $S$-준스킴이라 하고,
$X\times_S Y$도 국소 뇌터라고 하자~; $S'$을 $S$의 닫힌
부분집합이라 하고, $X'$ (각각 $Y'$)을 $X$ (각각 $Y$)의 닫힌
부분집합이라 하되 그 구조 사상에 의한 상이 $S'$에 포함된다고 하자.
$f(X')\subset Y'$을 만족하는 모든 $S$-사상 $f:X\to Y$에 대하여,
그래프 사상 $\Gamma_{\widehat f}$는 $f$의 그래프 사상의 연장
$(\Gamma_f)^\wedge$와 동일시된다.
\end{corollary}

\begin{corollary}[10.9.9]
\label{I.10.9.9-ko}
$X$, $Y$를 두 국소 뇌터 준스킴, $f:X\to Y$를 사상,
$Y'$을 $Y$의 닫힌 부분집합이라 하고, $X'=f^{-1}(Y')$라 하자.
그러면 형식적 준스킴\footnote{\textbf{번역 주.} 인쇄 원문은 여기서
\emph{préschéma}라고 쓰지만, $X_{/X'}$는 정의상 형식적 준스킴이고
우변도 형식적 준스킴들의 곱이므로 \emph{préschéma formel}에서
\emph{formel}이 빠진 것으로 보아 바로잡았다.} $X_{/X'}$는 다음 가환도식
\[
  \label{I.10.9.9.1-ko}
  \xymatrix{
    X\ar[d]_f &
    X_{/X'}\ar[l]\ar[d]^{\widehat f}\\
    Y &
    Y_{/Y'}\ar[l]
  }
\]
에 의해 형식적 준스킴들의 곱
$X\times_Y(Y_{/Y'})$와 동일시된다.
\end{corollary}

$S$와 $S'$을 모두 $Y$로, $X$와 $X'$을 모두 $X$로 바꾸어
(\hyperref[I.10.9.7-ko]{10.9.7})을 적용하면 충분하다.
\oldpage[I]{201}

\begin{remark}[10.9.10]
\label{I.10.9.10-ko}
$X$가 합 $X_1\amalg X_2$
(\hyperref[subsection:I.3.1-ko]{3.1})이고, $X'$이 합집합
$X'_1\cup X'_2$이며, 여기서 $X'_i$가 $X_i$의 닫힌 부분집합이면
($i=1,2$), 곧바로
$X_{/X'}={X_1}_{/X'_1}\amalg {X_2}_{/X'_2}$임을 알 수 있다.
\end{remark}

\subsection{아핀 형식적 스킴 위의 연접층에 대한 응용.}
\label{subsection:I.10.10-ko}

\begin{env}[10.10.1]
\label{I.10.10.1-ko}
이 절 전체에서 $A$는 \emph{뇌터 아딕환}, $\mathfrak{J}$는 $A$의
정의 아이디얼을 나타낸다. $X=\operatorname{Spec}(A)$,
$\mathfrak{X}=\operatorname{Spf}(A)$라 하자. 후자의 바탕공간은
$V(\mathfrak{J})$와 동일시되며, 이는 $X$의 닫힌 부분집합이다
(\hyperref[I.10.1.2-ko]{10.1.2}). 또한 정의
(\hyperref[I.10.1.2-ko]{10.1.2})와 정의
(\hyperref[I.10.8.4-ko]{10.8.4})에서 \emph{아핀 형식적 스킴
$\mathfrak{X}$}는 \emph{완비화 $X_{/\mathfrak{X}}$}, 곧 아핀 스킴
$X$를 그 바탕공간의 닫힌 부분집합 $\mathfrak{X}$을 따라 완비화한
것과 동일함을 알 수 있다. 따라서 각 연접 $\mathscr{O}_X$-가군 $\mathscr{F}$에는
유한 생성 $\mathscr{O}_{\mathfrak{X}}$-가군
$\mathscr{F}_{/\mathfrak{X}}$가 대응하며, 이는 또한 위상환의 층
$\mathscr{O}_{\mathfrak{X}}$ 위의 위상 가군의 층이다. 한편 모든 연접
$\mathscr{O}_X$-가군 $\mathscr{F}$는 $\widetilde M$ 꼴이다. 여기서
$M$은 유한 생성 $A$-가군이다
(\hyperref[I.1.5.1-ko]{1.5.1})~; 우리는
$(\widetilde M)_{/\mathfrak{X}}=M^\Delta$로 두겠다. 또한
$u:M\to N$이 유한 생성 $A$-가군 사이의 $A$-준동형이면, 이에
준동형 $\widetilde u:\widetilde M\to\widetilde N$이 대응하고, 따라서
연속 준동형
$\widetilde u_{/\mathfrak{X}}:(\widetilde M)_{/\mathfrak{X}}\to
(\widetilde N)_{/\mathfrak{X}}$도 대응한다. 이를 $u^\Delta$로
나타내겠다. 곧바로
$(v\circ u)^\Delta=v^\Delta\circ u^\Delta$임을 알 수 있다~; 이로써
\emph{$M^\Delta$를 주는 가법적 공변 함자}를 정의하였다. 이 함자는
유한 생성 $A$-가군의 범주에서 유한 생성
$\mathscr{O}_{\mathfrak{X}}$-가군의 범주로 간다. $A$가 \emph{이산}환이면
$M^\Delta=\widetilde M$이다.
\end{env}

\begin{proposition}[10.10.2]
\label{I.10.10.2-ko}
\medskip\noindent
\begin{enumerate}
  \item[\rm{(i)}] $M^\Delta$는 $M$에 관한 완전 함자이고, $A$-가군의
    함자적인 표준 동형
    $\Gamma(\mathfrak{X},M^\Delta)\xrightarrow{\sim}M$이 존재한다.
  \item[\rm{(ii)}] $M$과 $N$이 두 유한 생성 $A$-가군이면, 다음
    함자적인 표준 동형들이 존재한다.
    \[
      (M\otimes_A N)^\Delta
      \xrightarrow{\sim}
      M^\Delta\otimes_{\mathscr{O}_{\mathfrak{X}}}N^\Delta
      \tag{10.10.2.1}\label{I.10.10.2.1-ko}
    \]
    \[
      (\operatorname{Hom}_A(M,N))^\Delta
      \xrightarrow{\sim}
      \mathscr{H}\!om_{\mathscr{O}_{\mathfrak{X}}}
      (M^\Delta,N^\Delta)
      \tag{10.10.2.2}\label{I.10.10.2.2-ko}
    \]
  \item[\rm{(iii)}] 사상 $u\to u^\Delta$는 다음 함자적인 동형이다.
    \[
      \operatorname{Hom}_A(M,N)
      \xrightarrow{\sim}
      \operatorname{Hom}_{\mathscr{O}_{\mathfrak{X}}}
      (M^\Delta,N^\Delta)
      \tag{10.10.2.3}\label{I.10.10.2.3-ko}
    \]
\end{enumerate}
\end{proposition}

$M^\Delta$를 주는 함자의 완전성은 $\widetilde M$을 주는 함자
(\hyperref[I.1.3.5-ko]{1.3.5})와 $\mathscr{F}_{/\mathfrak{X}}$
(\hyperref[I.10.8.8-ko]{10.8.8})의 완전성에서 따른다. 정의에 따라
$\Gamma(\mathfrak{X},M^\Delta)$는 $A$-가군
$\Gamma(X,\widetilde M)=M$의 $\mathfrak{J}$-준아딕 위상에 관한 분리
완비화이다~; 그러나 $A$가 완비이고 $M$이 유한 생성인 까닭에,
$M$이 분리이고 완비임을 알고 있다
(\hyperref[0.7.3.6-ko]{0, 7.3.6}). 이로써 (i)의 증명이 끝난다. 동형
(\hyperref[I.10.10.2.1-ko]{10.10.2.1}) (각각
\hyperref[I.10.10.2.2-ko]{10.10.2.2})는 동형
(\hyperref[I.1.3.12-ko]{1.3.12, (i)})와
(\hyperref[I.10.8.10.1-ko]{10.8.10.1}) (각각
(\hyperref[I.1.3.12-ko]{1.3.12, (ii)})와
(\hyperref[I.10.8.10.2-ko]{10.8.10.2}))의 합성에서 얻어진다. 끝으로
$\operatorname{Hom}_A(M,N)$은 유한 생성 $A$-가군이므로 여기에 (i)를
적용할 수 있다. 그러면
$\Gamma(\mathfrak{X},(\operatorname{Hom}_A(M,N))^\Delta)$가
$\operatorname{Hom}_A(M,N)$과 동일시되고,
(\hyperref[I.10.10.2.2-ko]{10.10.2.2})를 이용하면 준동형
(\hyperref[I.10.10.2.3-ko]{10.10.2.3})이 동형임을 알 수 있다.

(\hyperref[I.10.10.2-ko]{10.10.2})에서
(\hyperref[I.1.3.7-ko]{1.3.7})과
(\hyperref[I.1.3.12-ko]{1.3.12})에서 얻은 것들과 유사한 일련의
귀결들을 얻으며, 그 정식화는 독자에게 맡긴다.
\oldpage[I]{202}

$M^\Delta$의 완전성을 완전열
$0\to\mathfrak{J}\to A\to A/\mathfrak{J}\to0$에 적용하면, 여기서
$\mathscr{O}_{\mathfrak{X}}$의 아이디얼층 중
$\mathfrak{J}^\Delta$로 나타낸 것이
(\hyperref[I.10.3.1-ko]{10.3.1})에서 같은 기호로 나타낸
아이디얼층과 일치함을 알 수 있다는 점에 유의하자. 이는
(\hyperref[I.10.3.2-ko]{10.3.2})에 따른다.

\begin{proposition}[10.10.3]
\label{I.10.10.3-ko}
(\hyperref[I.10.10.1-ko]{10.10.1})의 가정 아래에서
$\mathscr{O}_{\mathfrak{X}}$는 연접인 환의 층이다.
\end{proposition}

$f\in A$이면 $A_{\{f\}}$가 뇌터 아딕환임을 알고 있다
(\hyperref[0.7.6.11-ko]{0, 7.6.11}). 문제는 국소적이므로
(\hyperref[I.10.1.4-ko]{10.1.4})에 따라 준동형
$v:\mathscr{O}_{\mathfrak{X}}^n\to\mathscr{O}_{\mathfrak{X}}$의 핵이
유한 생성 $\mathscr{O}_{\mathfrak{X}}$-가군임을 보이면 충분하다.
이때 $v=u^\Delta$이다. 여기서 $u$는 $A$-준동형 $A^n\to A$이다
(\hyperref[I.10.10.2-ko]{10.10.2}). $A$가 뇌터이므로 $u$의 핵은
유한 생성이다. 다시 말해, 어떤 준동형
$A^m\xrightarrow{w}A^n$이 있어서 열
$A^m\xrightarrow{w}A^n\xrightarrow{u}A$가 완전하다. 따라서
(\hyperref[I.10.10.2-ko]{10.10.2})에 의해 열
$\mathscr{O}_{\mathfrak{X}}^m\xrightarrow{w^\Delta}
\mathscr{O}_{\mathfrak{X}}^n\xrightarrow{v}
\mathscr{O}_{\mathfrak{X}}$도 완전하다. 이로써 $v$의 핵이 유한
생성임을 보였다.

\begin{env}[10.10.4]
\label{I.10.10.4-ko}
앞의 표기 아래에서 $A_n=A/\mathfrak{J}^{n+1}$로 두고, $X_n$을 아핀
스킴
$\operatorname{Spec}(A_n)=(\mathfrak{X},
\mathscr{O}_{\mathfrak{X}}/\mathscr{J}^{n+1})$이라 하자. 여기서
$\mathscr{J}=\mathfrak{J}^\Delta$는 아이디얼 $\mathfrak{J}$에 대응하는
$\mathscr{O}_{\mathfrak{X}}$의 정의 아이디얼층이다. $m\leq n$일 때
표준 준동형 $A_n\to A_m$에 대응하는 준스킴의 사상 $X_m\to X_n$을
$u_{mn}$이라 하자~; 형식적 스킴 $\mathfrak{X}$는 $u_{mn}$들에 관한
$X_n$들의 귀납극한이다
(\hyperref[I.10.6.3-ko]{10.6.3}).
\end{env}

\begin{proposition}[10.10.5]
\label{I.10.10.5-ko}
(\hyperref[I.10.10.1-ko]{10.10.1})의 가정 아래에서
$\mathscr{F}$를 $\mathscr{O}_{\mathfrak{X}}$-가군이라 하자. 다음 조건들은
서로 동치이다~:
\begin{enumerate}
  \item[a)] $\mathscr{F}$는 연접 $\mathscr{O}_{\mathfrak{X}}$-가군이다.
  \item[b)] $\mathscr{F}$는 열 $(\mathscr{F}_n)$의 사영극한
    (\hyperref[I.10.6.6-ko]{10.6.6})과 동형이다. 이 열은 연접
    $\mathscr{O}_{X_n}$-가군들로 이루어지고
    $u_{mn}^*(\mathscr{F}_n)=\mathscr{F}_m$을 만족한다.
  \item[c)] 유한 생성 $A$-가군 $M$이 존재한다. 그 가군은
    (\hyperref[I.10.10.2-ko]{10.10.2, (i)})에 의해 표준 동형을
    제외하고 결정되며, $\mathscr{F}$는 $M^\Delta$와 동형이다.
\end{enumerate}
\end{proposition}

먼저 b)가 c)를 함의함을 보이자. $\mathscr{F}_n=\widetilde{M_n}$이고,
여기서 $M_n$은 유한 생성 $A_n$-가군이다. 가정에 의해
$M_m=M_n\otimes_{A_n}A_m$이 $m\leq n$일 때 성립한다
(\hyperref[I.1.6.5-ko]{1.6.5})~; 따라서 $M_n$들은 표준 디-준동형
$M_n\to M_m$ ($m\leq n$)에 관하여 사영계를 이룬다. $A_n$들의
정의로부터 이 사영계가
(\hyperref[0.7.2.9-ko]{0, 7.2.9})의 조건들을 만족함이 곧바로 따른다~;
그러므로 그 사영극한 $M$은 유한 생성 $A$-가군이고,
$M_n=M\otimes_A A_n$이 모든 $n$에 대하여 성립한다. 따라서
$\mathscr{F}_n$은 $X_n$ 위에서
$\widetilde M\otimes_{\mathscr{O}_X}
(\mathscr{O}_X/\widetilde{\mathfrak{J}}^{n+1})$에 의해 유도되는
가군이고, 정의 (\hyperref[I.10.8.4-ko]{10.8.4})에 의해
$\mathscr{F}=M^\Delta$이다.

역으로 c)는 b)를 함의한다. 실제로 $u_n$이 몰입 사상 $X_n\to X$이면,
$u_n^*(\widetilde M)=(M\otimes_A A_n)^{\sim}$은 $X_n$ 위에서
$\widetilde M\otimes_{\mathscr{O}_X}
(\mathscr{O}_X/\widetilde{\mathfrak{J}}^{n+1})$에 의해 유도되는
가군이고, 정의 (\hyperref[I.10.8.4-ko]{10.8.4})에 의해
$M^\Delta=\varprojlim u_n^*(\widetilde M)$이다~;
$u_m=u_n\circ u_{mn}$이 $m\leq n$일 때 성립하므로,
$\mathscr{F}_n=u_n^*(\widetilde M)$들은 b)의 조건들을 만족한다. 이로써
주장이 따른다.

이제 c)가 a)를 함의함을 보이자. 정의에 의해
$\mathscr{O}_{\mathfrak{X}}=A^\Delta$이다~; $M$은 어떤 준동형
$A^m\to A^n$의 여핵이므로, (\hyperref[I.10.10.2-ko]{10.10.2})에 의해
$M^\Delta$는 준동형
$\mathscr{O}_{\mathfrak{X}}^m\to\mathscr{O}_{\mathfrak{X}}^n$의
여핵이다. $\mathscr{O}_{\mathfrak{X}}$가 연접인 환의 층이므로
(\hyperref[I.10.10.3-ko]{10.10.3}), $M^\Delta$도 연접이다
(\hyperref[0.5.3.4-ko]{0, 5.3.4}).
\oldpage[I]{203}

끝으로 a)는 b)를 함의한다. $\mathscr{O}_{\mathfrak{X}}$-가군으로
보면,
$\mathscr{O}_{X_n}=\mathscr{O}_{\mathfrak{X}}/
\mathscr{J}^{n+1}=A_n^\Delta$이다~; $\mathscr{F}_n=\mathscr{F}
\otimes_{\mathscr{O}_{\mathfrak{X}}}\mathscr{O}_{X_n}$은 연접
$\mathscr{O}_{\mathfrak{X}}$-가군이다
(\hyperref[0.5.3.5-ko]{0, 5.3.5}). 또한 이것은
$\mathscr{O}_{X_n}$-가군이고 $\mathscr{J}^{n+1}$이 연접이므로,
$\mathscr{F}_n$은 연접 $\mathscr{O}_{X_n}$-가군이다
(\hyperref[0.5.3.10-ko]{0, 5.3.10}).
$u_{mn}^*(\mathscr{F}_n)=\mathscr{F}_m$이 $m\leq n$일 때 성립함도
곧바로 알 수 있다
($X_m\to X_n$에서 얻는 바탕공간 사이의 연속사상은
$\mathfrak{X}$의 항등사상임을 기억하면 된다). 따라서 층
$\mathscr{G}=\varprojlim\mathscr{F}_n$은 연접
$\mathscr{O}_{\mathfrak{X}}$-가군이다. 이는 b)가 a)를 함의함을 이미
보였기 때문이다. 표준 준동형 $\mathscr{F}\to\mathscr{F}_n$들은
사영계를 이루며, 사영극한을 취하면 표준 준동형
$w:\mathscr{F}\to\mathscr{G}$를 얻는다. 이제 남은 것은 $w$가 전단사임을
보이는 일뿐이다. 이 문제는 국소적이므로 $\mathscr{F}$가 어떤 준동형
$\mathscr{O}_{\mathfrak{X}}^p\to\mathscr{O}_{\mathfrak{X}}^q$의
여핵인 경우로 국한할 수 있다~; 이 준동형은 $v^\Delta$ 꼴이고, 여기서
$v$는 준동형 $A^p\to A^q$이다\footnote{\textbf{번역 주.} 인쇄 원문은
앞의 층 준동형을
$\mathscr{O}_{\mathfrak{X}}^p\to\mathscr{O}_{\mathfrak{X}}^q$로
둔 뒤, 이를 유도하는 $v$를 $A^m\to A^n$으로 쓴다. 그러나
$\Delta$ 대응에서는 $(A^r)^\Delta=\mathscr{O}_{\mathfrak{X}}^r$이므로
유한 자유 가군의 계수(階數)가 보존된다. 따라서 $v:A^p\to A^q$로
바로잡았다.}
(\hyperref[I.10.10.2-ko]{10.10.2}). 그러므로 $\mathscr{F}$는
$M^\Delta$와 동형이다. 여기서 $M=\operatorname{Coker}v$이다
(\hyperref[I.10.10.2-ko]{10.10.2}). 이제
(\hyperref[I.10.10.2-ko]{10.10.2})에 의해
$\mathscr{F}_n=M^\Delta\otimes_{\mathscr{O}_{\mathfrak{X}}}A_n^\Delta
=(M\otimes_A A_n)^\Delta$이다. 또한 $\mathfrak{J}$-아딕 위상은
$M\otimes_A A_n$ 위에서 이산이므로,
$(M\otimes_A A_n)^\Delta=(M\otimes_A A_n)^{\sim}$이다
($\mathscr{O}_{X_n}$-가군으로서). 위에서
$M^\Delta=\varprojlim\mathscr{F}_n$임을 보았으므로, 이 경우 $w$는
실제로 항등사상이다.
\par\hfill C.Q.F.D.\par

\begin{corollary}[10.10.6]
\label{I.10.10.6-ko}
$\mathscr{F}$가 (\hyperref[I.10.10.5-ko]{10.10.5})의 조건 b)를
만족하면, 사영계 $(\mathscr{F}_n)$은
$\mathscr{F}\otimes_{\mathscr{O}_{\mathfrak{X}}}\mathscr{O}_{X_n}$들로
이루어진 계와 동형이다.
\end{corollary}

\begin{env}[10.10.7]
\label{I.10.10.7-ko}
이제 $A$, $B$를 두 뇌터 아딕환이라 하고,
$\varphi:B\to A$를 연속 준동형이라 하자. $\mathfrak{J}$ (각각
$\mathfrak{K}$)를 $A$ (각각 $B$)의 정의 아이디얼이라 하되
$\varphi(\mathfrak{K})\subset\mathfrak{J}$라고 하고,
$X=\operatorname{Spec}(A)$, $Y=\operatorname{Spec}(B)$,
$\mathfrak{X}=\operatorname{Spf}(A)$,
$\mathfrak{Y}=\operatorname{Spf}(B)$라 하자. $f:X\to Y$를
$\varphi$에 대응하는 준스킴의 사상이라 하고
(\hyperref[I.1.6.1-ko]{1.6.1}), 완비화한 것들 사이로의 연장을
$\widehat f:\mathfrak{X}\to\mathfrak{Y}$라 하자
(\hyperref[I.10.9.1-ko]{10.9.1}). 이는 또한 $\varphi$에 대응하는
형식적 준스킴의 사상이다 (\hyperref[I.10.2.2-ko]{10.2.2}).
\end{env}

\begin{proposition}[10.10.8]
\label{I.10.10.8-ko}
모든 유한 생성 $B$-가군 $N$에 대하여 다음과 같은
$\mathscr{O}_{\mathfrak{X}}$-가군의 함자적인 표준 동형이 존재한다~:
\[
  \widehat f^*(N^\Delta)
  \xrightarrow{\sim}
  (N\otimes_B A)^\Delta.
\]
\end{proposition}

실제로 표준 사상들을 $i_X:\mathfrak{X}\to X$와
$i_Y:\mathfrak{Y}\to Y$로 나타내면, 함자적인 표준 동형을 통해
동일시할 때 (\hyperref[I.10.8.8-ko]{10.8.8})
$N^\Delta=i_Y^*(\widetilde N)$이고
\[
  (N\otimes_B A)^\Delta
  =i_X^*((N\otimes_B A)^{\sim})
  =i_X^*(f^*(\widetilde N))
\]
이다 (\hyperref[I.1.6.5-ko]{1.6.5})~; 따라서 명제는 도표
(\hyperref[I.10.9.2-ko]{10.9.2})의 가환성에서 따른다.

\begin{corollary}[10.10.9]
\label{I.10.10.9-ko}
$\mathfrak b$를 $B$의 임의의 아이디얼이라 하면,
$\widehat f^*(\mathfrak b^\Delta)\mathscr{O}_{\mathfrak{X}}
=(\mathfrak bA)^\Delta$이다.
\end{corollary}

실제로 $j$를 표준 단사 준동형 $\mathfrak b\to B$라 하자. 이에 대응하는
$j^\Delta:\mathfrak b^\Delta\to\mathscr{O}_{\mathfrak{Y}}$는
$\mathscr{O}_{\mathfrak{Y}}$-가군의 표준 단사 준동형이다~;
정의에 따라
$\widehat f^*(\mathfrak b^\Delta)\mathscr{O}_{\mathfrak{X}}$는 준동형
$\widehat f^*(j^\Delta):\widehat f^*(\mathfrak b^\Delta)\to
\mathscr{O}_{\mathfrak{X}}=\widehat f^*(\mathscr{O}_{\mathfrak{Y}})$의
상이다~; 그런데 (\hyperref[I.10.10.8-ko]{10.10.8})에 의해 이 준동형은
$(j\otimes1)^\Delta:(\mathfrak b\otimes_B A)^\Delta\to
\mathscr{O}_{\mathfrak{X}}=(B\otimes_B A)^\Delta$와 동일시된다.
$j\otimes1$의 상은 아이디얼 $\mathfrak bA$이고 이는 $A$의 아이디얼이므로,
(\hyperref[I.10.10.2-ko]{10.10.2})에 따라 $(j\otimes1)^\Delta$의 상은
$(\mathfrak bA)^\Delta$이다. 이로써 결론이 따른다.
\oldpage[I]{204}

\subsection{형식적 준스킴 위의 연접층.}
\label{subsection:I.10.11-ko}

\begin{proposition}[10.11.1]
\label{I.10.11.1-ko}
$\mathfrak{X}$가 국소 뇌터 형식적 준스킴이면,
$\mathscr{O}_{\mathfrak{X}}$는 연접인 환의 층이고 $\mathfrak{X}$의
모든 정의 아이디얼층은 연접이다.
\end{proposition}

문제는 국소적이므로 뇌터 아핀 형식적 스킴인 경우로 환원되며, 따라서
명제는 (\hyperref[I.10.10.3-ko]{10.10.3})과
(\hyperref[I.10.10.5-ko]{10.10.5})에서 따른다.

\begin{env}[10.11.2]
\label{I.10.11.2-ko}
$\mathfrak{X}$를 국소 뇌터 형식적 준스킴,
$\mathscr{J}$를 $\mathfrak{X}$의 정의 아이디얼층이라 하자. $X_n$을
(보통 의미의) 국소 뇌터 준스킴
$(\mathfrak{X},\mathscr{O}_{\mathfrak{X}}/\mathscr{J}^{n+1})$이라 하자.
그러면 $\mathfrak{X}$는 열 $(X_n)$의 귀납극한이고, 그 전이사상들은
표준 사상 $u_{mn}:X_m\to X_n$이다
(\hyperref[I.10.6.3-ko]{10.6.3}). 이 표기 아래에서 다음이 성립한다~:
\end{env}

\begin{theorem}[10.11.3]
\label{I.10.11.3-ko}
$\mathscr{O}_{\mathfrak{X}}$-가군 $\mathscr{F}$가 연접이기 위한
필요충분조건은 그것이 열 $(\mathscr{F}_n)$의 사영극한과 동형인 것이다.
여기서 각 $\mathscr{F}_n$은 연접
$\mathscr{O}_{X_n}$-가군이며,
$u_{mn}^*(\mathscr{F}_n)=\mathscr{F}_m$이 $m\leq n$일 때 성립한다
(\hyperref[I.10.6.6-ko]{10.6.6}). 이때 사영계 $(\mathscr{F}_n)$은
$u_n^*(\mathscr{F})=\mathscr{F}
\otimes_{\mathscr{O}_{\mathfrak{X}}}\mathscr{O}_{X_n}$들로 이루어진
계와 동형이다. 여기서 $u_n$은 표준 사상 $X_n\to\mathfrak{X}$이다.
\end{theorem}

문제는 국소적이므로 $\mathfrak{X}$가 뇌터 아핀 형식적 스킴인 경우로
환원되며, 이 경우 정리는 (\hyperref[I.10.10.5-ko]{10.10.5})와
(\hyperref[I.10.10.6-ko]{10.10.6})의 귀결이다.

따라서 \emph{연접 $\mathscr{O}_{\mathfrak{X}}$-가군을 주는 것은
사영계 $(\mathscr{F}_n)$를 주는 것과 동치이다. 여기서 각 항은 연접
$\mathscr{O}_{X_n}$-가군이고,
$u_{mn}^*(\mathscr{F}_n)=\mathscr{F}_m$이 $m\leq n$일 때 성립한다}.

\begin{corollary}[10.11.4]
\label{I.10.11.4-ko}
$\mathscr{F}$와 $\mathscr{G}$가 두 연접
$\mathscr{O}_{\mathfrak{X}}$-가군이면,
(\hyperref[I.10.11.3-ko]{10.11.3})의 표기를 쓰면 다음 함자적인 표준
동형을 정의할 수 있다~:
\[
  \label{I.10.11.4.1-ko}
  \operatorname{Hom}_{\mathscr{O}_{\mathfrak{X}}}
  (\mathscr{F},\mathscr{G})
  \xrightarrow{\sim}
  \varprojlim_n
  \operatorname{Hom}_{\mathscr{O}_{X_n}}
  (\mathscr{F}_n,\mathscr{G}_n).
  \tag{10.11.4.1}
\]
\end{corollary}

우변의 사영극한은 다음 사상들에 관하여 취한 것으로 이해해야 한다~:
$\theta_n\mapsto u_{mn}^*(\theta_n)$ ($m\leq n$). 이들은
$\operatorname{Hom}_{\mathscr{O}_{X_n}}
(\mathscr{F}_n,\mathscr{G}_n)$에서
$\operatorname{Hom}_{\mathscr{O}_{X_m}}
(\mathscr{F}_m,\mathscr{G}_m)$으로 가는 사상이다. 준동형
(\hyperref[I.10.11.4.1-ko]{10.11.4.1})은 원소
$\theta\in\operatorname{Hom}_{\mathscr{O}_{\mathfrak{X}}}
(\mathscr{F},\mathscr{G})$에 열 $(u_n^*(\theta))$을 대응시킨다~;
(\hyperref[I.10.11.3-ko]{10.11.3})을 고려하면, 사영계
$(\theta_n)\in\varprojlim_n
\operatorname{Hom}_{\mathscr{O}_{X_n}}
(\mathscr{F}_n,\mathscr{G}_n)$에
$\operatorname{Hom}_{\mathscr{O}_{\mathfrak{X}}}
(\mathscr{F},\mathscr{G})$ 안에서의 그 사영극한을 대응시킴으로써
앞의 준동형의 역준동형을 정의할 수 있음을 곧 알 수 있다.

\begin{corollary}[10.11.5]
\label{I.10.11.5-ko}
준동형 $\theta:\mathscr{F}\to\mathscr{G}$가 전사일 필요충분조건은
이에 대응하는 준동형
$\theta_0=u_0^*(\theta):\mathscr{F}_0\to\mathscr{G}_0$가 전사인 것이다.
\end{corollary}

문제는 국소적이므로
$\mathfrak{X}=\operatorname{Spf}(A)$이고 $A$가 뇌터 아딕환이며,
$\mathscr{F}=M^\Delta$, $\mathscr{G}=N^\Delta$, $\theta=u^\Delta$인
경우로 환원된다. 여기서 $M$과 $N$은 유한 생성 $A$-가군이고
$u$는 준동형 $M\to N$이다~; 또한 이때
$\theta_0=\widetilde{u_0}$이고, 여기서 $u_0$는 준동형
$u\otimes1:M\otimes_A A/\mathfrak{J}\to
N\otimes_A A/\mathfrak{J}$이다~;
결론은 $\theta$와 $u$가 동시에 전사이고, 또 $\theta_0$와
$u_0$가 동시에 전사라는 사실
(\hyperref[I.1.3.9-ko]{1.3.9} 및
\hyperref[I.10.10.2-ko]{10.10.2})과 $u$와 $u_0$가 동시에 전사라는
사실 (\hyperref[0.7.1.14-ko]{0, 7.1.14})에서 따른다.

\begin{env}[10.11.6]
\label{I.10.11.6-ko}
정리 (\hyperref[I.10.11.3-ko]{10.11.3})은 모든 연접
$\mathscr{O}_{\mathfrak{X}}$-가군 $\mathscr{F}$를
\emph{위상} $\mathscr{O}_{\mathfrak{X}}$-가군으로 볼 수 있음을
보여 준다. 즉 이를 \emph{의사이산} 군의 층 $\mathscr{F}_n$들의
사영극한으로 본다
(\hyperref[0.3.8.1-ko]{0, 3.8.1}). 따라서
(\hyperref[I.10.11.4-ko]{10.11.4})에 의해 준동형
$u:\mathscr{F}\to\mathscr{G}$가 연접
$\mathscr{O}_{\mathfrak{X}}$-가군 사이의 준동형이면 자동으로
\emph{연속}이다

\oldpage[I]{205}

(\hyperref[0.3.8.2-ko]{0, 3.8.2}). 또한 $\mathscr{H}$가 연접
$\mathscr{O}_{\mathfrak{X}}$-부분가군으로서 연접
$\mathscr{O}_{\mathfrak{X}}$-가군 $\mathscr{F}$에 들어 있으면, 모든 열린집합
$U\subset\mathfrak{X}$에 대하여 $\Gamma(U,\mathscr{H})$는 위상군
$\Gamma(U,\mathscr{F})$의 \emph{닫힌} 부분군이다. 실제로 함자
$\Gamma$가 왼쪽 완전이므로 $\Gamma(U,\mathscr{H})$는 준동형
$\Gamma(U,\mathscr{F})\to\Gamma(U,\mathscr{F}/\mathscr{H})$의
핵이다. 이 준동형은 앞서 본 바에 따라 \emph{연속}인데,
$\mathscr{F}/\mathscr{H}$가 연접이기 때문이다
(\hyperref[0.5.3.4-ko]{0, 5.3.4})~; 따라서 이 주장은
$\Gamma(U,\mathscr{F}/\mathscr{H})$가 분리 위상군이라는 사실에서
따른다.
\end{env}

\begin{proposition}[10.11.7]
\label{I.10.11.7-ko}
$\mathscr{F}$와 $\mathscr{G}$를 두 연접
$\mathscr{O}_{\mathfrak{X}}$-가군이라 하자.
(\hyperref[I.10.11.3-ko]{10.11.3})의 표기를 쓰면 다음 위상
$\mathscr{O}_{\mathfrak{X}}$-가군의 함자적인 표준 동형들을 정의할 수 있다
(\hyperref[I.10.11.6-ko]{10.11.6})~:
\[
  \label{I.10.11.7.1-ko}
  \mathscr{F}\otimes_{\mathscr{O}_{\mathfrak{X}}}\mathscr{G}
  \xrightarrow{\sim}
  \varprojlim_n
  (\mathscr{F}_n\otimes_{\mathscr{O}_{X_n}}\mathscr{G}_n)
  \tag{10.11.7.1}
\]
\[
  \label{I.10.11.7.2-ko}
  \mathscr{H}\!om_{\mathscr{O}_{\mathfrak{X}}}
  (\mathscr{F},\mathscr{G})
  \xrightarrow{\sim}
  \varprojlim_n
  \mathscr{H}\!om_{\mathscr{O}_{X_n}}(\mathscr{F}_n,\mathscr{G}_n)
  \tag{10.11.7.2}
\]
\end{proposition}

동형 (\hyperref[I.10.11.7.1-ko]{10.11.7.1})의 존재는 공식
\[
  \mathscr{F}_n\otimes_{\mathscr{O}_{X_n}}\mathscr{G}_n
  =
  (\mathscr{F}\otimes_{\mathscr{O}_{\mathfrak{X}}}
  \mathscr{O}_{X_n})
  \otimes_{\mathscr{O}_{X_n}}
  (\mathscr{G}\otimes_{\mathscr{O}_{\mathfrak{X}}}
  \mathscr{O}_{X_n})
  =
  (\mathscr{F}\otimes_{\mathscr{O}_{\mathfrak{X}}}\mathscr{G})
  \otimes_{\mathscr{O}_{\mathfrak{X}}}\mathscr{O}_{X_n}
\]
과 (\hyperref[I.10.11.3-ko]{10.11.3})에서 따른다. 두 변을 위상을
고려하지 않은 가군의 층으로 보았을 때, 동형
(\hyperref[I.10.11.7.2-ko]{10.11.7.2})은
$\mathscr{H}\!om_{\mathscr{O}_{\mathfrak{X}}}
(\mathscr{F},\mathscr{G})$와
$\mathscr{H}\!om_{\mathscr{O}_{X_n}}(\mathscr{F}_n,\mathscr{G}_n)$의
단면의 정의와 $\mathfrak{X}$의 임의의 뇌터 아핀 형식적 열린집합 위에
유도된 준스킴에 (\hyperref[I.10.11.4.1-ko]{10.11.4.1})을 적용하여 얻는
동형의 존재에서 따른다. 이제 준콤팩트 집합 위에서 동형
(\hyperref[I.10.11.7.2-ko]{10.11.7.2})이 쌍연속임을 보이는 것만 남는다.
따라서 $\mathfrak{X}=\operatorname{Spf}(A)$이고 $A$가 뇌터 아딕환인
경우로 환원되며, 이때 (\hyperref[I.10.10.5-ko]{10.10.5})에 따라
$\mathscr{F}=M^\Delta$, $\mathscr{G}=N^\Delta$이고 $M$, $N$은 유한 생성
$A$-가군이다~; (\hyperref[I.10.10.2.1-ko]{10.10.2.1}),
(\hyperref[I.10.10.2.3-ko]{10.10.2.3}) 및
(\hyperref[I.1.3.12-ko]{1.3.12, (ii)})을 고려하면, 표준 동형
$\operatorname{Hom}_A(M,N)\xrightarrow{\sim}
\varprojlim_n\operatorname{Hom}_{A_n}(M_n,N_n)$
($M_n=M\otimes_A A_n$, $N_n=N\otimes_A A_n$)이 연속임을 보이는 문제로
환원되는데, 이는 (\hyperref[0.7.8.2-ko]{0, 7.8.2})에서 증명되었다.

\begin{env}[10.11.8]
\label{I.10.11.8-ko}
$\operatorname{Hom}_{\mathscr{O}_{\mathfrak{X}}}
(\mathscr{F},\mathscr{G})$는 위상군의 층
$\mathscr{H}\!om_{\mathscr{O}_{\mathfrak{X}}}
(\mathscr{F},\mathscr{G})$의 단면군이므로 군 위상을 갖는다.
$\mathfrak{X}$가 \emph{뇌터}이면,
(\hyperref[I.10.11.7.2-ko]{10.11.7.2})에 따라 이 군에서 $0$의 근방들의
기본계는 부분군
$\operatorname{Hom}_{\mathscr{O}_{\mathfrak{X}}}
(\mathscr{F},\mathscr{J}^n\mathscr{G})$ ($n$은 임의)을 취하여 얻는다.
\end{env}

\begin{proposition}[10.11.9]
\label{I.10.11.9-ko}
$\mathfrak{X}$를 뇌터 형식적 준스킴,
$\mathscr{F}$와 $\mathscr{G}$를 두 연접
$\mathscr{O}_{\mathfrak{X}}$-가군이라 하자. 위상군
$\operatorname{Hom}_{\mathscr{O}_{\mathfrak{X}}}
(\mathscr{F},\mathscr{G})$에서 전사인 준동형들(각각 단사인 준동형들,
전단사인 준동형들)은 열린 부분집합을 이룬다.
\end{proposition}

(\hyperref[I.10.11.5-ko]{10.11.5})에 따라
$\operatorname{Hom}_{\mathscr{O}_{\mathfrak{X}}}
(\mathscr{F},\mathscr{G})$ 안의 전사 준동형들의 집합은 다음 연속사상
$\operatorname{Hom}_{\mathscr{O}_{\mathfrak{X}}}
(\mathscr{F},\mathscr{G})\to\allowbreak
\operatorname{Hom}_{\mathscr{O}_{X_0}}(\mathscr{F}_0,\mathscr{G}_0)$에
의한 이산군
$\operatorname{Hom}_{\mathscr{O}_{X_0}}(\mathscr{F}_0,\mathscr{G}_0)$의
어떤 부분집합의 역상이므로 첫 번째 주장이 따른다. 두 번째 주장을
증명하기 위해 $\mathfrak{X}$를 유한 개의 뇌터 아핀 형식적 열린집합
$U_i$로 덮자. $\theta\in
\operatorname{Hom}_{\mathscr{O}_{\mathfrak{X}}}
(\mathscr{F},\mathscr{G})$가 단사일 필요충분조건은 각 연속 제한사상
$\operatorname{Hom}_{\mathscr{O}_{\mathfrak{X}}}
(\mathscr{F},\mathscr{G})\to
\operatorname{Hom}_{\mathscr{O}_{\mathfrak{X}}|U_i}
(\mathscr{F}|U_i,\mathscr{G}|U_i)$ 아래에서 얻는 상이 모두 단사인 것이다~;
따라서 아핀인 경우로 환원되며, 이 경우에는 이미
(\hyperref[0.7.8.3-ko]{0, 7.8.3})에서 증명되었다.

\oldpage[I]{206}

\subsection{형식적 준스킴의 아딕 사상.}
\label{subsection:I.10.12-ko}

\begin{env}[10.12.1]
\label{I.10.12.1-ko}
$\mathfrak{X}$, $\mathfrak{S}$를 두 \emph{국소 뇌터} 형식적
준스킴이라 하자~; 사상 $f:\mathfrak{X}\to\mathfrak{S}$가
\emph{아딕}이라 함은 아이디얼 $\mathscr{J}$가 존재하여 그것이
$\mathfrak{S}$의 정의 아이디얼이고
$\mathscr{K}=f^*(\mathscr{J})\mathscr{O}_{\mathfrak{X}}$가
$\mathfrak{X}$의 정의 아이디얼인 것을 뜻한다~; 이때
$\mathfrak{X}$를
\emph{$\mathfrak{S}$-아딕 형식적 준스킴}이라고도 한다
($f$에 관하여). 이 경우 $\mathscr{J}_1$을 $\mathfrak{S}$의
\emph{임의의} 정의 아이디얼이라 하면,
$\mathscr{K}_1=f^*(\mathscr{J}_1)\mathscr{O}_{\mathfrak{X}}$는
$\mathfrak{X}$의 정의 아이디얼이다. 실제로 문제는 국소적이므로
$\mathfrak{X}$와 $\mathfrak{S}$가 뇌터 아핀 형식적 준스킴이라고
가정해도 된다~; 따라서 정수 $n$이 존재하여
$\mathscr{J}^n\subset\mathscr{J}_1$이고
$\mathscr{J}_1^n\subset\mathscr{J}$이다
(\hyperref[I.10.3.6-ko]{10.3.6} 및
\hyperref[0.7.1.4-ko]{0, 7.1.4}). 그러므로
$\mathscr{K}^n\subset\mathscr{K}_1$이고
$\mathscr{K}_1^n\subset\mathscr{K}$이다. 첫 번째 관계는
$\mathscr{K}_1=\mathfrak{K}_1^\Delta$임을 보인다. 여기서
$\mathfrak{K}_1$은
$A=\Gamma(\mathfrak{X},\mathscr{O}_{\mathfrak{X}})$의 열린
아이디얼이다. 두 번째 관계는 $\mathfrak{K}_1$이 $A$의 정의
아이디얼임을 보이므로
(\hyperref[0.7.1.4-ko]{0, 7.1.4}), 이로써 주장이 따른다.
\end{env}

앞의 내용에서 $\mathfrak{X}$와 $\mathfrak{Y}$가 두
$\mathfrak{S}$-아딕 형식적 준스킴이면,
\emph{모든 $\mathfrak{S}$-사상
$u:\mathfrak{X}\to\mathfrak{Y}$은 아딕이다}라는 것이 곧바로 따른다~:
실제로 $f:\mathfrak{X}\to\mathfrak{S}$,
$g:\mathfrak{Y}\to\mathfrak{S}$를 각각 구조 사상이라 하고,
$\mathscr{J}$를 $\mathfrak{S}$의 정의 아이디얼이라 하자. 그러면
$f=g\circ u$이므로
$u^*(g^*(\mathscr{J})\mathscr{O}_{\mathfrak{Y}})
\mathscr{O}_{\mathfrak{X}}=f^*(\mathscr{J})
\mathscr{O}_{\mathfrak{X}}$는 $\mathfrak{X}$의 정의 아이디얼이고,
가정에 따라 $g^*(\mathscr{J})\mathscr{O}_{\mathfrak{Y}}$는
$\mathfrak{Y}$의 정의 아이디얼이다.

\begin{env}[10.12.2]
\label{I.10.12.2-ko}
이하에서는 국소 뇌터 형식적 준스킴 $\mathfrak{S}$와 아이디얼
$\mathscr{J}$를 고정하되, 후자는 $\mathfrak{S}$의 정의 아이디얼이라고
가정하자~; 또한
$S_n=(\mathfrak{S},\mathscr{O}_{\mathfrak{S}}/\mathscr{J}^{n+1})$로
두겠다. $\mathfrak{S}$-아딕 형식적 준스킴들(국소 뇌터인 것들)은
자명하게 하나의 \emph{범주}를 이룬다. 귀납계 $(X_n)$이 (보통 의미의)
국소 뇌터 $S_n$-준스킴들로 이루어져 있을 때, 다음 조건을 만족하면
이를 \emph{아딕 $(S_n)$-귀납계}라 하겠다. 구조 사상
$f_n:X_n\to S_n$들에 대하여 $m\leq n$이면 도식
\[
  \label{I.10.12.2.1-ko}
  \xymatrix{
    X_n \ar[d]_{f_n}
    & X_m \ar[l] \ar[d]^{f_m}\\
    S_n
    & S_m \ar[l]
  }
  \tag{10.12.2.1}
\]
은 가환이고, \emph{$X_m$을 곱
$X_n\times_{S_n}S_m=(X_n)_{(S_m)}$과 동일시한다}. 아딕 귀납계들도
하나의 \emph{범주}를 이룬다~: 실제로 이러한 계들의 사상
$(X_n)\to(Y_n)$을 \emph{$S_n$-사상 $u_n:X_n\to Y_n$들로 이루어진
귀납계}로 정의하되, $u_m$이 $(u_n)_{(S_m)}$과 동일시되는 조건을
모든 $m\leq n$에 대하여 만족하도록 하면 충분하다. 그러면~:
\end{env}

\begin{theorem}[10.12.3]
\label{I.10.12.3-ko}
$\mathfrak{S}$-아딕 형식적 준스킴들의 범주와 아딕
$(S_n)$-귀납계들의 범주 사이에는 표준 동치가 있다.
\end{theorem}

이 동치는 다음과 같이 얻는다~: $\mathfrak{X}$가
$\mathfrak{S}$-아딕 형식적 준스킴이고,
$f:\mathfrak{X}\to\mathfrak{S}$가 구조 사상이라 하면,
$\mathscr{K}=f^*(\mathscr{J})\mathscr{O}_{\mathfrak{X}}$는
$\mathfrak{X}$의 정의 아이디얼이며, $\mathfrak{X}$에 귀납계
$X_n=(\mathfrak{X},\mathscr{O}_{\mathfrak{X}}/\mathscr{K}^{n+1})$을
대응시킨다. 여기서 $f_n:X_n\to S_n$은 $f$에 대응하는 구조 사상이다
(\hyperref[I.10.5.6-ko]{10.5.6}). 먼저 $(X_n)$이
\emph{아딕 귀납계}임을 보이자~: $f=(\psi,\theta)$라 쓰면
$\psi^*(\mathscr{J})\mathscr{O}_{\mathfrak{X}}=\mathscr{K}$이고, 따라서
$\psi^*(\mathscr{J}^n)\mathscr{O}_{\mathfrak{X}}=\mathscr{K}^n$이며, 이는
모든 $n$에 대하여 성립한다.
또한 함자 $\psi^*$의 완전성에 의해
$\mathscr{K}^{m+1}/\mathscr{K}^{n+1}
=\psi^*(\mathscr{J}^{m+1}/\mathscr{J}^{n+1})
(\mathscr{O}_{\mathfrak{X}}/\mathscr{K}^{n+1})$이 $m\leq n$일 때
성립한다. 그러므로 원하는 결론은
(\hyperref[I.4.4.5-ko]{4.4.5})에서 따른다. 더욱이
$\mathfrak{S}$-사상 $u:\mathfrak{X}\to\mathfrak{Y}$이 두
$\mathfrak{S}$-아딕 형식적 준스킴 사이의 사상일 때에는, 명백한 표기를 쓰면,

\oldpage[I]{207}

$S_n$-사상 $u_n:X_n\to Y_n$들의 귀납계로서 $u_m$이
$(u_n)_{(S_m)}$와 동일시되는 조건이 $m\leq n$일 때 성립하는 것이
대응함도 곧바로 확인된다.

이렇게 동치를 올바르게 정의했음은 다음의 더 정밀한 명제에서 따른다~:

\begin{proposition}[10.12.3.1]
\label{I.10.12.3.1-ko}
$(X_n)$을 $S_n$-준스킴들의 귀납계라 하자. 구조 사상
$f_n:X_n\to S_n$들이 도식
(\hyperref[I.10.12.2.1-ko]{10.12.2.1})을 가환이게 하고, $X_m$을
$X_n\times_{S_n}S_m$와 동일시하는 조건이 $m\leq n$일 때 성립한다고
가정하자.
그러면 귀납계 $(X_n)$은 (\hyperref[I.10.6.3-ko]{10.6.3})의 조건
b)와 c)를 만족한다. 그 귀납극한을 $\mathfrak{X}$라 하고,
$f:\mathfrak{X}\to\mathfrak{S}$를 귀납계 $(f_n)$의 귀납극한 사상이라
하자. 이때 $X_0$이 국소 뇌터이면
$\mathfrak{X}$도 국소 뇌터이고 $f$는 아딕 사상이다.
\end{proposition}

$\mathscr{O}_{S_n}$에서 부분준스킴 $S_m$을 $S_n$ 안에 정의하는
아이디얼층은 멱영이다. 따라서 (\hyperref[I.4.4.5-ko]{4.4.5})에
의해 $\mathscr{O}_{X_n}$에서 부분준스킴 $X_m$을 $X_n$ 안에 정의하는
아이디얼층도 멱영이므로, (\hyperref[I.10.6.3-ko]{10.6.3})의 조건들은
실제로 만족된다. 그러므로 문제는 $\mathfrak{X}$와 $\mathfrak{S}$에서
국소적이며, $\mathfrak{S}=\operatorname{Spf}(A)$,
$\mathscr{J}=\mathfrak{J}^\Delta$라 가정해도 된다. 여기서 $A$는
뇌터 $\mathfrak{J}$-아딕환이고, $X_n=\operatorname{Spec}(B_n)$이다.
$A_n=A/\mathfrak{J}^{n+1}$이라 하면,
가정에 의해 $B_0$은 뇌터이고,
$\mathfrak{J}_n=\mathfrak{J}/\mathfrak{J}^{n+1}$이라 놓을 때
$B_m=B_n/\mathfrak{J}_n^{m+1}B_n$이다. 따라서 $B_n\to B_0$의 핵은
$\mathfrak{K}_n=\mathfrak{J}_nB_n$이고, $B_n\to B_m$의 핵은
$\mathfrak{K}_n^{m+1}$이며, 이는 $m\leq n$일 때 성립한다. 또한 $A_1$이
뇌터이므로 $\mathfrak{J}_1$은 유한 생성 $A_1$-가군이다. 따라서
$\overline{\mathfrak{K}}_1=\mathfrak{K}_1/\mathfrak{K}_1^2$은
유한 생성 $B_1$-가군이고, \emph{따라서 특히}
$B_0=B_1/\mathfrak{K}_1$-가군으로도 유한 생성이다. 이제
$\mathfrak{X}$가 뇌터임은 (\hyperref[I.10.6.4-ko]{10.6.4})에서
따른다. $B=\varprojlim B_n$이라 하면
$\mathfrak{X}=\operatorname{Spf}(B)$이고, $\mathfrak{K}$를
$B\to B_0$의 핵이라 하면 $B_n=B/\mathfrak{K}^{n+1}$이다.
$\rho_n:A/\mathfrak{J}^{n+1}\to B/\mathfrak{K}^{n+1}$을 $f_n$에
대응하는 준동형이라 하면
\[
  \mathfrak{K}/\mathfrak{K}^{n+1}
  =(B/\mathfrak{K}^{n+1})\rho_n(\mathfrak{J}/\mathfrak{J}^{n+1})
\]
이다. 한편 준동형 $\rho:A\to B$가 $f$에 대응하며
$\varprojlim\rho_n$과 같으므로, 아이디얼 $\mathfrak{J}B$는 $B$의
아이디얼로서 $\mathfrak{K}$ 안에서 조밀하다. 그런데 $B$의 모든
아이디얼은 닫혀 있으므로 (\hyperref[0.7.3.5-ko]{0, 7.3.5}),
$\mathfrak{K}=\mathfrak{J}B$이다.
$\mathscr{K}=\mathfrak{K}^\Delta$라 하면 관계
$f^*(\mathscr{J})\mathscr{O}_{\mathfrak{X}}=\mathscr{K}$는
(\hyperref[I.10.10.9-ko]{10.10.9})에서 따르며, 이것으로 증명이 끝난다.

\begin{env}[10.12.3.2]
\label{I.10.12.3.2-ko}
앞의 동치는 두 $\mathfrak{S}$-아딕 형식적 준스킴
$\mathfrak{X}$, $\mathfrak{Y}$에 대하여 다음 \emph{표준 전단사}를 준다~:
\[
  \operatorname{Hom}_{\mathfrak{S}}(\mathfrak{X},\mathfrak{Y})
  \xrightarrow{\sim}
  \varprojlim_n\operatorname{Hom}_{S_n}(X_n,Y_n)
\]
여기서 사영극한은 사상 $u_n\to(u_n)_{(S_m)}$들에 관하여 취하며,
$m\leq n$이다.
\end{env}

\subsection{유한형 사상.}
\label{subsection:I.10.13-ko}

\begin{proposition}[10.13.1]
\label{I.10.13.1-ko}
$\mathfrak{Y}$를 국소 뇌터 형식적 준스킴,
$\mathscr{K}$를 $\mathfrak{Y}$의 정의 아이디얼,
$f:\mathfrak{X}\to\mathfrak{Y}$를 형식적 준스킴의 사상이라 하자. 다음
조건들은 동치이다~:
\begin{enumerate}
  \item[a)] $\mathfrak{X}$는 국소 뇌터이고, $f$는 아딕 사상이며
    (\hyperref[I.10.12.1-ko]{10.12.1}),
    $\mathscr{J}=f^*(\mathscr{K})\mathscr{O}_{\mathfrak{X}}$라 놓을 때
    사상 $f_0:(\mathfrak{X},\mathscr{O}_{\mathfrak{X}}/\mathscr{J})
    \to(\mathfrak{Y},\mathscr{O}_{\mathfrak{Y}}/\mathscr{K})$는 $f$로부터
    유도되며 유한형이다.
  \item[b)] $\mathfrak{X}$는 국소 뇌터이고, 아딕
    $(Y_n)$-귀납계 $(X_n)$의 귀납극한이며, 사상 $X_0\to Y_0$는
    유한형이다.

\oldpage[I]{208}

  \item[c)] $\mathfrak{Y}$의 모든 점은 다음 성질을 갖는 뇌터 아핀
    형식적 열린집합인 근방 $V$를 가진다~:
    \begin{enumerate}
      \item[(\textbf{Q})] $f^{-1}(V)$는 뇌터 아핀 형식적 열린집합
        $U_i$들의 유한 합집합이고, 각각에 대하여 뇌터 아딕환
        $\Gamma(U_i,\mathscr{O}_{\mathfrak{X}})$는
        $\Gamma(V,\mathscr{O}_{\mathfrak{Y}})$ 위의 제한 형식적 멱급수
        대수를 어떤 (필연적으로 닫힌) 아이디얼로 나눈 몫과 위상적으로
        동형이다 (\hyperref[0.7.5.1-ko]{0, 7.5.1}).
    \end{enumerate}
\end{enumerate}
\end{proposition}

a)가 b)를 함의하는 것은 (\hyperref[I.10.12.3-ko]{10.12.3})에 따라
명백하다. b)가 c)를 함의함을 보이기 위해, 문제는 $\mathfrak{Y}$ 위에서
국소적이므로 $\mathfrak{Y}=\operatorname{Spf}(B)$이고 $B$가 뇌터
아딕환이라고 가정할 수 있다. $\mathscr{K}=\mathfrak{K}^\Delta$라
하되, $\mathfrak{K}$는 $B$의 정의 아이디얼이라 하자. 가정에 의해 $X_0$는
$Y_0$ 위에서 유한형이므로, $X_0$는 아핀 열린집합 $U_i$들의 유한
합집합이고, 환 $A_{i0}$, 즉 $X_0$가 $U_i$ 위에 유도하는 아핀
스킴의 환은 $B/\mathfrak{K}$, 즉 $Y_0$의 환 위의 유한 생성 대수이다
(\hyperref[I.6.3.2-ko]{6.3.2}).
(\hyperref[I.5.1.9-ko]{5.1.9})에 따라 $U_i$는 각 뇌터 준스킴
$X_n$에서도 아핀 열린집합이며, 환 $A_{in}$을 $X_n$이 $U_i$ 위에
유도하는 아핀 스킴의 환이라 하면 가정 b)에 의해 $m\leq n$일 때
$A_{im}$은 $A_{in}/\mathfrak{K}^{m+1}A_{in}$과 동형이다. 따라서
$U_i$ 위에서 $\mathfrak{X}$가 유도하는 형식적 준스킴은
$\operatorname{Spf}(A_i)$와 동형이며, 여기서
$A_i=\varprojlim_n A_{in}$이다 (\hyperref[I.10.6.4-ko]{10.6.4}).
$A_i$는 $\mathfrak{K}A_i$-아딕환이고, $A_i/\mathfrak{K}A_i$는
$A_{i0}$과 동형이며 $B/\mathfrak{K}$ 위의 유한 생성 대수이다.
그러므로 (\hyperref[0.7.5.5-ko]{0, 7.5.5})에 따라 $A_i$는 $B$ 위의
제한 형식적 멱급수 대수를 어떤 아이디얼로 나눈 몫과 위상적으로
동형이다. 이 아이디얼은 반드시 닫혀 있는데, 그러한 대수는 뇌터이기 때문이다
(\hyperref[0.7.5.4-ko]{0, 7.5.4}).

c)가 a)를 함의함을 증명하기 위해, $\mathfrak{X}=\operatorname{Spf}(A)$
도 아핀인 경우로 한정할 수 있다. 여기서 $A$는 뇌터 아딕환이고, $B$
위의 제한 형식적 멱급수 대수를 어떤 닫힌 아이디얼로 나눈 몫과
동형이다. 그러면 (\hyperref[0.7.5.5-ko]{0, 7.5.5})에 따라
$A/\mathfrak{K}A$는 $B/\mathfrak{K}$ 위의 유한 생성 대수이고,
$\mathfrak{K}A=\mathfrak{J}$는 $A$의 정의 아이디얼이다. 따라서
(\hyperref[I.10.10.9-ko]{10.10.9})에 의해 a)의 조건들이 성립한다.

명제 (\hyperref[I.10.13.1-ko]{10.13.1})의 조건들이 성립하면, a)는
\emph{모든 정의 아이디얼 $\mathscr{K}$}에 대하여 성립한다. 여기서
정의 아이디얼은 $\mathfrak{Y}$의 것이며, 이는 c)에 따른다. 따라서
b)에서는 \emph{모든 $f_n$이 유한형 사상}임에 유의하자.

\begin{corollary}[10.13.2]
\label{I.10.13.2-ko}
(\hyperref[I.10.13.1-ko]{10.13.1})의 조건들이 성립하면, 모든 뇌터
아핀 형식적 열린집합 $V$는 $\mathfrak{Y}$ 안에서 성질
(\textbf{Q})를 가지며, $\mathfrak{Y}$가 뇌터이면 $\mathfrak{X}$도
뇌터이다.
\end{corollary}

이는 (\hyperref[I.10.13.1-ko]{10.13.1})과
(\hyperref[I.6.3.2-ko]{6.3.2})에서 곧바로 따른다.

\begin{definition}[10.13.3]
\label{I.10.13.3-ko}
(\hyperref[I.10.13.1-ko]{10.13.1})의 동치인 성질 a), b), c)가
성립하면, 사상 $f$를 \emph{유한형}이라고 하거나 $\mathfrak{X}$를
\emph{유한형 $\mathfrak{Y}$-형식적 준스킴} 또는
\emph{$\mathfrak{Y}$ 위의 유한형 형식적 준스킴}이라고 한다.
\end{definition}

\begin{corollary}[10.13.4]
\label{I.10.13.4-ko}
$\mathfrak{X}=\operatorname{Spf}(A)$,
$\mathfrak{Y}=\operatorname{Spf}(B)$를 두 뇌터 아핀 형식적 스킴이라
하자. $\mathfrak{X}$가 $\mathfrak{Y}$ 위에서 유한형이기 위한
필요충분조건은 뇌터 아딕환 $A$가 $B$ 위의 제한 형식적 멱급수 대수를
어떤 닫힌 아이디얼로 나눈 몫과 동형인 것이다.
\end{corollary}

실제로 (\hyperref[I.10.13.1-ko]{10.13.1})의 표기를 사용하자.
$\mathfrak{X}$가 $\mathfrak{Y}$ 위에서 유한형이면
(\hyperref[I.6.3.3-ko]{6.3.3})에 따라 $A/\mathfrak{K}A$는
$(B/\mathfrak{K})$ 위의 유한 생성 대수이고,
$\mathfrak{K}A$는 $A$의 정의 아이디얼이다
(\hyperref[I.10.10.9-ko]{10.10.9}). 그러므로
(\hyperref[0.7.5.5-ko]{0, 7.5.5})에서 결론이 따른다.

\begin{proposition}[10.13.5]
\label{I.10.13.5-ko}
\medskip\noindent
\begin{enumerate}
  \item[\rm{(i)}] 형식적 준스킴의 두 유한형 사상의 합성은
    유한형이다.

\oldpage[I]{209}

  \item[\rm{(ii)}] $\mathfrak{X}$, $\mathfrak{S}$,
    $\mathfrak{S}'$을 세 국소 뇌터(각각 뇌터) 형식적 준스킴이라 하고,
    $f:\mathfrak{X}\to\mathfrak{S}$,
    $g:\mathfrak{S}'\to\mathfrak{S}$를 두 사상이라 하자. $f$가
    유한형이면 $\mathfrak{X}\times_{\mathfrak{S}}\mathfrak{S}'$은
    국소 뇌터(각각 뇌터)이고 $\mathfrak{S}'$ 위에서 유한형이다.
  \item[\rm{(iii)}] $\mathfrak{S}$를 국소 뇌터 형식적 준스킴,
    $\mathfrak{X}'$, $\mathfrak{Y}'$을 두 국소 뇌터
    $\mathfrak{S}$-형식적 준스킴이라 하되,
    $\mathfrak{X}'\times_{\mathfrak{S}}\mathfrak{Y}'$은 국소 뇌터라
    하자.
    $\mathfrak{X}$, $\mathfrak{Y}$가 국소 뇌터
    $\mathfrak{S}$-형식적 준스킴이고,
    $f:\mathfrak{X}\to\mathfrak{X}'$,
    $g:\mathfrak{Y}\to\mathfrak{Y}'$이 두 유한형
    $\mathfrak{S}$-사상이면,
    $\mathfrak{X}\times_{\mathfrak{S}}\mathfrak{Y}$은 국소 뇌터이고
    $f\times_{\mathfrak{S}}g$는 유한형 $\mathfrak{S}$-사상이다.
\end{enumerate}
\end{proposition}

(iii)은 (\hyperref[I.3.5.1-ko]{3.5.1})의 형식적 논증에 의해 (i)과
(ii)에서 따르므로, (i)과 (ii)만 증명하면 충분하다.

$\mathfrak{X}$, $\mathfrak{Y}$, $\mathfrak{Z}$를 세 국소 뇌터 형식적
준스킴이라 하고, $f:\mathfrak{X}\to\mathfrak{Y}$,
$g:\mathfrak{Y}\to\mathfrak{Z}$를 두 유한형 사상이라 하자.
$\mathscr{L}$이 $\mathfrak{Z}$의 정의 아이디얼이면,
$\mathscr{K}=g^*(\mathscr{L})\mathscr{O}_{\mathfrak{Y}}$는
$\mathfrak{Y}$의 정의 아이디얼이고,
$\mathscr{J}=f^*(g^*(\mathscr{L}))\mathscr{O}_{\mathfrak{X}}$는
$\mathfrak{X}$의 정의 아이디얼이다. 다음과 같이 놓자~:
$X_0=(\mathfrak{X},\mathscr{O}_{\mathfrak{X}}/\mathscr{J})$,
$Y_0=(\mathfrak{Y},\mathscr{O}_{\mathfrak{Y}}/\mathscr{K})$,
$Z_0=(\mathfrak{Z},\mathscr{O}_{\mathfrak{Z}}/\mathscr{L})$.
또 $f_0:X_0\to Y_0$, $g_0:Y_0\to Z_0$를 각각 $f$, $g$에 대응하는
사상이라 하자. 가정에 의해 $f_0$와 $g_0$가 유한형이므로,
$g_0\circ f_0$도 유한형이다 (\hyperref[I.6.3.4-ko]{6.3.4}). 이
사상은 $g\circ f$에 대응하므로, (\hyperref[I.10.13.1-ko]{10.13.1})에
따라 $g\circ f$는 유한형이다.

(ii)의 조건 아래에서 $\mathfrak{S}$ (각각 $\mathfrak{X}$,
$\mathfrak{S}'$)는 국소 뇌터 준스킴들의 열 $(S_n)$ (각각 $(X_n)$,
$(S'_n)$)의 귀납극한이고, (\hyperref[I.10.13.1-ko]{10.13.1})에 따라
$X_m=X_n\times_{S_n}S_m$라고 가정할 수 있으며, 이는 모든
$m\leq n$에 대해 성립한다. 그러면
형식적 준스킴
$\mathfrak{X}\times_{\mathfrak{S}}\mathfrak{S}'$은 준스킴
$X_n\times_{S_n}S'_n$들의 귀납극한이고
(\hyperref[I.10.7.4-ko]{10.7.4}), 다음이 성립한다~:
\[
  X_m\times_{S_m}S'_m
  =(X_n\times_{S_n}S_m)\times_{S_m}S'_m
  =(X_n\times_{S_n}S'_n)\times_{S'_n}S'_m.
\]
또한 $X_0\times_{S_0}S'_0$은 $X_0$가 $S_0$ 위에서 유한형이므로
국소 뇌터이다
(\hyperref[I.6.3.8-ko]{6.3.8}). 따라서 먼저
(\hyperref[I.10.12.3.1-ko]{10.12.3.1})에 의해
$\mathfrak{X}\times_{\mathfrak{S}}\mathfrak{S}'$은 국소 뇌터이다.
더 나아가 $X_0\times_{S_0}S'_0$은 $S'_0$ 위에서 유한형이므로
(\hyperref[I.6.3.8-ko]{6.3.8}),
(\hyperref[I.10.12.3.1-ko]{10.12.3.1})과
(\hyperref[I.10.13.1-ko]{10.13.1})에 따라
$\mathfrak{X}\times_{\mathfrak{S}}\mathfrak{S}'$은
$\mathfrak{S}'$ 위에서 유한형이다. 이것으로 (ii)의 증명이 끝난다.
뇌터 준스킴에 관한 주장은 (\hyperref[I.6.3.8-ko]{6.3.8})의 즉각적인
귀결이다.

\begin{corollary}[10.13.6]
\label{I.10.13.6-ko}
(\hyperref[I.10.9.9-ko]{10.9.9})의 가정 아래에서 $f$가 유한형
사상이면, 완비화한 것들 사이로의 연장 $\widehat f$도 유한형이다.
\end{corollary}

\subsection{형식적 준스킴의 닫힌 부분준스킴.}
\label{subsection:I.10.14-ko}

\begin{proposition}[10.14.1]
\label{I.10.14.1-ko}
$\mathfrak{X}$를 국소 뇌터 형식적 준스킴,
$\mathscr{A}$를 $\mathscr{O}_{\mathfrak{X}}$의 연접 아이디얼층이라
하자. $\mathfrak{Y}$가 $\mathscr{O}_{\mathfrak{X}}/\mathscr{A}$의
(닫힌) 지지집합이면, 위상환 달린 공간
$(\mathfrak{Y},(\mathscr{O}_{\mathfrak{X}}/\mathscr{A})|\mathfrak{Y})$은
국소 뇌터 형식적 준스킴이며, $\mathfrak{X}$가 뇌터이면 이 형식적
준스킴도 뇌터이다.
\end{proposition}

(\hyperref[I.10.10.3-ko]{10.10.3})과
(\hyperref[0.5.3.4-ko]{0, 5.3.4})에 따라
$\mathscr{O}_{\mathfrak{X}}/\mathscr{A}$는 구조층 위의 연접 가군이고, 따라서 그
지지집합 $\mathfrak{Y}$는 닫혀 있다
(\hyperref[0.5.2.2-ko]{0, 5.2.2}). $\mathscr{J}$를
$\mathfrak{X}$의 정의 아이디얼이라 하고,
$X_n=(\mathfrak{X},\mathscr{O}_{\mathfrak{X}}/\mathscr{J}^{n+1})$이라
하자~; 환의 층 $\mathscr{O}_{\mathfrak{X}}/\mathscr{A}$는 층들
$\mathscr{O}_{\mathfrak{X}}/(\mathscr{A}+\mathscr{J}^{n+1})
=(\mathscr{O}_{\mathfrak{X}}/\mathscr{A})
\otimes_{\mathscr{O}_{\mathfrak{X}}}
(\mathscr{O}_{\mathfrak{X}}/\mathscr{J}^{n+1})$의 사영극한이다
(\hyperref[I.10.11.3-ko]{10.11.3}). 이 층들은 모두 지지집합
$\mathfrak{Y}$를 갖는다. 층
$(\mathscr{A}+\mathscr{J}^{n+1})/\mathscr{J}^{n+1}$은
$\mathscr{O}_{\mathfrak{X}}$-가군으로서 연접이다. 실제로
$\mathscr{J}^{n+1}$이 연접이므로,
$(\mathscr{A}+\mathscr{J}^{n+1})/\mathscr{J}^{n+1}$은 또한 연접
$(\mathscr{O}_{\mathfrak{X}}/\mathscr{J}^{n+1})$-가군이다
(\hyperref[0.5.3.10-ko]{0, 5.3.10})~; $Y_n$을 이 아이디얼층으로
정의되는 $X_n$의 닫힌 부분준스킴이라 하면, 위상환 달린 공간
$(\mathfrak{Y},(\mathscr{O}_{\mathfrak{X}}/\mathscr{A})|\mathfrak{Y})$

\oldpage[I]{210}
은 $Y_n$들의 귀납극한인 형식적 준스킴임이 곧바로 확인된다.
(\hyperref[I.10.6.4-ko]{10.6.4})의 조건들이 만족되므로, 이 형식적
준스킴은 국소 뇌터이고, $\mathfrak{X}$가 뇌터이면 뇌터이다. 실제로
이 경우에는 (\hyperref[I.6.1.4-ko]{6.1.4})에 따라 $Y_0$가 뇌터이다.

\begin{definition}[10.14.2]
\label{I.10.14.2-ko}
형식적 준스킴 $\mathfrak{X}$의 \emph{닫힌 부분준스킴}이라 함은
$(\mathfrak{Y},(\mathscr{O}_{\mathfrak{X}}/\mathscr{A})|\mathfrak{Y})$
꼴의 모든 형식적 준스킴을 말한다. 여기서 $\mathscr{A}$는
$\mathscr{O}_{\mathfrak{X}}$의 연접 아이디얼층이다~; 이 형식적
준스킴을 $\mathscr{A}$가 정의하는 닫힌 부분준스킴이라고 한다.
\footnote{\textbf{번역 주.} 인쇄 원문은 여기와 바로 다음 문장에서
단지 연접 가군이라고 쓰지만, 몫환의 구성과 뒤따르는 대응에는 연접
아이디얼층이 필요하므로 두 곳 모두 이를 교정했다.}
\end{definition}

이와 같이 정의된 $\mathscr{O}_{\mathfrak{X}}$의 연접 아이디얼층들과
$\mathfrak{X}$의 닫힌 부분준스킴들 사이의 대응은 전단사임이 명백하다.

위상환 달린 공간의 사상
$j=(\psi,\theta):\mathfrak{Y}\to\mathfrak{X}$에서 $\psi$는 단사사상
$\mathfrak{Y}\to\mathfrak{X}$이고 $\theta^\sharp$는 표준 준동형
$\mathscr{O}_{\mathfrak{X}}\to\mathscr{O}_{\mathfrak{X}}/\mathscr{A}$이다.
이는 명백히 (\hyperref[I.10.4.5-ko]{10.4.5}) 형식적 준스킴의
사상이며, 이를 $\mathfrak{Y}$에서 $\mathfrak{X}$로 가는
\emph{표준 포함 사상}이라 한다. $\mathfrak{X}=\operatorname{Spf}(A)$이고
$A$가 뇌터 아딕환이면, $\mathscr{A}=\mathfrak{a}^\Delta$이고 여기서
$\mathfrak{a}$는 $A$의 아이디얼이다
(\hyperref[I.10.10.5-ko]{10.10.5}). 앞의 내용에서 곧바로
$\mathfrak{Y}=\operatorname{Spf}(A/\mathfrak{a})$임이 동형을 제외하면
따르고, $j$는 (\hyperref[I.10.2.2-ko]{10.2.2}) 표준 준동형
$A\to A/\mathfrak{a}$에 대응한다.

국소 뇌터 형식적 준스킴의 사상
$f:\mathfrak{Z}\to\mathfrak{X}$가 \emph{닫힌 몰입 사상}이라 함은
이 사상이 $\mathfrak{Z}\xrightarrow{g}\mathfrak{Y}\xrightarrow{j}\mathfrak{X}$로
인수분해되고, 여기서 $g$는 $\mathfrak{Z}$에서 닫힌 부분준스킴
$\mathfrak{Y}$로 가는 동형사상이며, 이 부분준스킴은 $\mathfrak{X}$의
닫힌 부분준스킴이고 $j$는 표준 포함
사상인 것을 뜻한다. $j$는 환 달린 공간의 단사사상이므로 $g$와
$\mathfrak{Y}$는 반드시 \emph{유일}하다.

\begin{proposition}[10.14.3]
\label{I.10.14.3-ko}
닫힌 몰입 사상은 유한형 사상이다.
\end{proposition}

$\mathfrak{X}$가 아핀 형식적 스킴 $\operatorname{Spf}(A)$이고
$\mathfrak{Y}=\operatorname{Spf}(A/\mathfrak{a})$인 경우로 곧바로
환원된다~; 명제는 (\hyperref[I.10.13.1-ko]{10.13.1, c)})에서 따른다.

\begin{lemma}[10.14.4]
\label{I.10.14.4-ko}
$f:\mathfrak{Y}\to\mathfrak{X}$를 국소 뇌터 형식적 준스킴의
사상이라 하고, $(U_\alpha)$를 $f(\mathfrak{Y})$의 덮개라 하자. 이
덮개의 원소들은 $\mathfrak{X}$의 뇌터 아핀 형식적 열린집합이고, 각
$f^{-1}(U_\alpha)$는
$\mathfrak{Y}$의 뇌터 아핀 형식적 열린집합이라고 하자. $f$가 닫힌
몰입 사상이기 위한 필요충분조건은 $f(\mathfrak{Y})$가
$\mathfrak{X}$의 닫힌 부분집합이고, 모든 $\alpha$에 대하여 $f$를
$f^{-1}(U_\alpha)$에 제한한 사상이
(\hyperref[I.10.4.6-ko]{10.4.6})에 따라 전사 준동형
$\Gamma(U_\alpha,\mathscr{O}_{\mathfrak{X}})\to
\Gamma(f^{-1}(U_\alpha),\mathscr{O}_{\mathfrak{Y}})$에 대응하는 것이다.
\end{lemma}

이 조건들이 필요함은 명백하다. 역으로 이 조건들이 성립한다고 하고,
$\mathfrak{a}_\alpha$를 준동형
$\Gamma(U_\alpha,\mathscr{O}_{\mathfrak{X}})\to
\Gamma(f^{-1}(U_\alpha),\mathscr{O}_{\mathfrak{Y}})$의 핵이라 하자.
연접 아이디얼층 $\mathscr{A}$를 $\mathscr{O}_{\mathfrak{X}}$ 위에
다음과 같이 정의한다. $\mathscr{A}|U_\alpha=\mathfrak{a}_\alpha^\Delta$로
두고, $\mathscr{A}$는 $U_\alpha$들의 합집합의 여집합에서 단위
아이디얼로 둔다. 실제로 $f(\mathfrak{Y})$가 닫혀 있고
$(\mathscr{O}_{\mathfrak{X}}|U_\alpha)/\mathfrak{a}_\alpha^\Delta$의
지지집합이 $U_\alpha\cap f(\mathfrak{Y})$이므로,
$\mathfrak{a}_\alpha^\Delta$와 $\mathfrak{a}_\beta^\Delta$가 뇌터
아핀 형식적 열린집합 $V\subset U_\alpha\cap U_\beta$ 위에서 같은
층을 유도함을 확인하면 충분하다.\footnote{\textbf{번역 주.} 인쇄
원문은 여집합에서 영 아이디얼로 둔다고 쓰지만, 닫힌 부분집합인
사상의 상 밖에서 몫층이 영층이 되려면 단위 아이디얼이어야 한다.
또한 그 닫힌 부분집합은 핵 아이디얼층 자체의 지지집합이 아니라 그
몫층의 지지집합이므로 두 곳을 함께 교정했다.}
$f$를 $f^{-1}(U_\alpha)$에 제한한 사상은 이 형식적 준스킴에서
$U_\alpha$로 가는 닫힌 몰입 사상이므로, $f^{-1}(V)$는
$f^{-1}(U_\alpha)$의 뇌터 아핀 형식적 열린집합이고 $f$를
$f^{-1}(V)$에 제한한 사상도 닫힌 몰입 사상이다~; 이 제한에 대응하는
전사 준동형의 핵을 $\mathfrak{b}$라 하자. 그 준동형은
$\Gamma(V,\mathscr{O}_{\mathfrak{X}})\to
\Gamma(f^{-1}(V),\mathscr{O}_{\mathfrak{Y}})$이다. 그러면
(\hyperref[I.10.10.2-ko]{10.10.2})에 따라
$\mathfrak{a}_\alpha^\Delta$가 $\mathfrak{b}^\Delta$를 $V$ 위에
유도함은 명백하다. 이렇게 아이디얼층 $\mathscr{A}$를 정의하면
$f=j\circ g$임이 명백하다. 여기서
$j:\mathfrak{Z}\to\mathfrak{X}$는 닫힌 부분준스킴 $\mathfrak{Z}$에서
$\mathfrak{X}$로 가는 표준 포함 사상이고, 이 부분준스킴은
$\mathscr{A}$가 정의하며,
$g$는 $\mathfrak{Y}$에서 $\mathfrak{Z}$로 가는 동형사상이다.
\footnote{\textbf{번역 주.} 인쇄 원문은 합성의 순서를 거꾸로 쓰지만,
선언된 정의역과 공역에 맞추어 바로잡았다.}

\begin{proposition}[10.14.5]
\label{I.10.14.5-ko}
\begin{enumerate}
  \item[\rm{(i)}] $f:\mathfrak{Z}\to\mathfrak{Y}$,
  $g:\mathfrak{Y}\to\mathfrak{X}$가 국소 뇌터 형식적 준스킴의
  닫힌 몰입 사상이면, $g\circ f$는 닫힌 몰입 사상이다.

\oldpage[I]{211}
  \item[\rm{(ii)}] $\mathfrak{X}$, $\mathfrak{Y}$,
  $\mathfrak{S}$를 세 국소 뇌터 형식적 준스킴이라 하고,
  $f:\mathfrak{X}\to\mathfrak{S}$를 닫힌 몰입 사상,
  $g:\mathfrak{Y}\to\mathfrak{S}$를 사상이라 하자. 그러면 사상
  $\mathfrak{X}\times_{\mathfrak{S}}\mathfrak{Y}\to\mathfrak{Y}$는
  닫힌 몰입 사상이다.
  \item[\rm{(iii)}] $\mathfrak{S}$를 국소 뇌터 형식적 준스킴,
  $\mathfrak{X}'$, $\mathfrak{Y}'$을 두 국소 뇌터
  $\mathfrak{S}$-형식적 준스킴이라 하되,
  $\mathfrak{X}'\times_{\mathfrak{S}}\mathfrak{Y}'$은 국소 뇌터라고
  하자. $\mathfrak{X}$, $\mathfrak{Y}$가 두 국소 뇌터
  $\mathfrak{S}$-형식적 준스킴이고,
  $f:\mathfrak{X}\to\mathfrak{X}'$, $g:\mathfrak{Y}\to\mathfrak{Y}'$이
  닫힌 몰입 사상인 두 $\mathfrak{S}$-사상이면,
  $f\times_{\mathfrak{S}}g$는 닫힌 몰입 사상이다.
\end{enumerate}
\end{proposition}

(\hyperref[I.3.5.1-ko]{3.5.1})에 따라 여기서도 (i)과 (ii)만 증명하면
충분하다.

(i)을 증명하기 위해 $\mathfrak{Y}$ (각각 $\mathfrak{Z}$)가
$\mathfrak{X}$ (각각 $\mathfrak{Y}$)의 닫힌 부분준스킴이고, 이들이
각각 연접 아이디얼층 $\mathscr{J}$ (각각 $\mathscr{K}$)로 정의된다고
가정할 수 있다. 이 두 층은 각각 $\mathscr{O}_{\mathfrak{X}}$ (각각
$\mathscr{O}_{\mathfrak{Y}}$)의 아이디얼층이다~; $\psi$가 바탕공간의 단사사상
$\mathfrak{Y}\to\mathfrak{X}$이면, $\psi_*(\mathscr{K})$는
$\psi_*(\mathscr{O}_{\mathfrak{Y}})
=\mathscr{O}_{\mathfrak{X}}/\mathscr{J}$의 연접 아이디얼층이다
(\hyperref[0.5.3.12-ko]{0, 5.3.12}). 따라서 이는 또한 연접
$\mathscr{O}_{\mathfrak{X}}$-가군이다
(\hyperref[0.5.3.10-ko]{0, 5.3.10})~; 핵 $\mathscr{K}_1$, 즉 준동형
$\mathscr{O}_{\mathfrak{X}}\to
(\mathscr{O}_{\mathfrak{X}}/\mathscr{J})/\psi_*(\mathscr{K})$의 핵은
$\mathscr{O}_{\mathfrak{X}}$의 연접 아이디얼층이고
(\hyperref[0.5.3.4-ko]{0, 5.3.4}),
$\mathscr{O}_{\mathfrak{X}}/\mathscr{K}_1$은
$\psi_*(\mathscr{O}_{\mathfrak{Y}}/\mathscr{K})$와 동형이다. 따라서
$\mathfrak{Z}$는 $\mathfrak{X}$의 닫힌 부분준스킴과 동형이다.

(ii)를 증명하기 위해 $\mathfrak{S}=\operatorname{Spf}(A)$,
$\mathfrak{X}=\operatorname{Spf}(B)$,
$\mathfrak{Y}=\operatorname{Spf}(C)$인 경우로 한정할 수 있음은
명백하다. 여기서 $A$는 뇌터 $\mathfrak{J}$-아딕환이고,
$B=A/\mathfrak{a}$이며 $\mathfrak{a}$는 $A$의 아이디얼이고,
$C$는 아딕이고 뇌터인 위상 $A$-대수이다. 이제 준동형
$C\to C\widehat{\otimes}_A(A/\mathfrak{a})$이 \emph{전사}임을
보이면 충분하다~: 실제로 $A/\mathfrak{a}$는 유한 생성 $A$-가군이고
그 위상은 $\mathfrak{J}$-아딕 위상이다~; 따라서
(\hyperref[0.7.7.8-ko]{0, 7.7.8})에 의해
$C\widehat{\otimes}_A(A/\mathfrak{a})$는
$C\otimes_A(A/\mathfrak{a})=C/\mathfrak{a}C$와 동일시되므로 주장이
따른다.

\begin{corollary}[10.14.6]
\label{I.10.14.6-ko}
(\hyperref[I.10.14.5-ko]{10.14.5, (ii)})의 가정 아래에서
$p:\mathfrak{X}\times_{\mathfrak{S}}\mathfrak{Y}\to\mathfrak{X}$,
$q:\mathfrak{X}\times_{\mathfrak{S}}\mathfrak{Y}\to\mathfrak{Y}$를
표준 사영들이라 하자. 그러면 도식
\[
  \xymatrix{
    \mathfrak{X} \ar[d]_{f}
    & \mathfrak{X}\times_{\mathfrak{S}}\mathfrak{Y}
      \ar[l]_{p} \ar[d]^{q}\\
    \mathfrak{S}
    & \mathfrak{Y} \ar[l]^{g}
  }
\]
은 가환이다. 모든 연접 $\mathscr{O}_{\mathfrak{X}}$-가군
$\mathscr{F}$에 대하여 다음 표준 $\mathscr{O}_{\mathfrak{Y}}$-가군
동형사상이 존재한다~:
\[
  \label{I.10.14.6.1-ko}
  u:g^*(f_*(\mathscr{F}))
  \xrightarrow{\sim}
  q_*(p^*(\mathscr{F})).
  \tag{10.14.6.1}
\]
\end{corollary}

준동형 $g^*(f_*(\mathscr{F}))\to q_*(p^*(\mathscr{F}))$을 정의하는
것은 준동형
$f_*(\mathscr{F})\to g_*(q_*(p^*(\mathscr{F})))
=f_*(p_*(p^*(\mathscr{F})))$을 정의하는 것과 같다
(\hyperref[0.4.4.3-ko]{0, 4.4.3})~: 앞의 수반 대응 아래에서 원문의
표기를 따라 $u=f_*(\rho)$로 쓰겠다. 여기서 $\rho$는 표준 준동형
$\mathscr{F}\to p_*(p^*(\mathscr{F}))$이다
(\hyperref[0.4.4.3-ko]{0, 4.4.3}).\footnote{\textbf{번역 주.} 이
등식은 서로 다른 준동형 집합의 원소들을 수반 전단사로 대응시켜 쓴
표기이며, 두 사상이 문자 그대로 같은 형을 갖는다는 뜻이 아니다.}
$u$가 동형사상임을 보이기 위해 $\mathfrak{S}$, $\mathfrak{X}$,
$\mathfrak{Y}$가 각각 뇌터 아딕환 $A$, $B$, $C$의 형식적
스펙트럼이고, (\hyperref[I.10.14.5-ko]{10.14.5, (ii)})에서 본
조건들이 성립하는 경우로 곧바로 환원된다~; 이때
$\mathscr{F}=M^\Delta$이고, $M$은 유한 생성
$(A/\mathfrak{a})$-가군이다
(\hyperref[I.10.10.5-ko]{10.10.5}).
(\hyperref[I.10.14.6.1-ko]{10.14.6.1})의 양변은
(\hyperref[I.10.10.8-ko]{10.10.8})에 따라 각각
$(C\otimes_A M)^\Delta$와
$((C/\mathfrak{a}C)\otimes_{A/\mathfrak{a}}M)^\Delta$로 동일시된다.
그런데
$(C/\mathfrak{a}C)\otimes_{A/\mathfrak{a}}M
=(C\otimes_A(A/\mathfrak{a}))\otimes_{A/\mathfrak{a}}M$은
$C\otimes_A M$과 표준적으로 동일시되므로 따름정리가 증명된다.

\oldpage[I]{212}
\begin{corollary}[10.14.7]
\label{I.10.14.7-ko}
$X$를 (보통 의미의) 국소 뇌터 준스킴, $Y$를 $X$의 닫힌 부분준스킴,
$j$를 표준 포함 사상 $Y\to X$, $X'$을 $X$의 닫힌 부분집합이라 하고
$Y'=Y\cap X'$이라 하자~; 그러면
$\widehat{j}:Y_{/Y'}\to X_{/X'}$은 닫힌 몰입 사상이고, 모든 연접
$\mathscr{O}_Y$-가군 $\mathscr{F}$에 대하여
\[
  \widehat{j}_*(\mathscr{F}_{/Y'})=(j_*(\mathscr{F}))_{/X'}.
\]
\end{corollary}

$Y'=j^{-1}(X')$이므로 (\hyperref[I.10.9.9-ko]{10.9.9})를 사용하고
(\hyperref[I.10.14.5-ko]{10.14.5})와
(\hyperref[I.10.14.6-ko]{10.14.6})을 적용하면 충분하다.
\subsection{분리된 형식적 준스킴.}
\label{subsection:I.10.15-ko}

\begin{definition}[10.15.1]
\label{I.10.15.1-ko}
형식적 준스킴 $\mathfrak{S}$, $\mathfrak{S}$-형식적 준스킴
$\mathfrak{X}$, 구조 사상 $f:\mathfrak{X}\to\mathfrak{S}$가 주어졌다고
하자. 대각 사상
$\Delta_{\mathfrak{X}|\mathfrak{S}}:\mathfrak{X}\to
\mathfrak{X}\times_{\mathfrak{S}}\mathfrak{X}$ (또한
$\Delta_{\mathfrak{X}}$라고도 쓴다)을 사상
$(1_{\mathfrak{X}},1_{\mathfrak{X}})_{\mathfrak{S}}$라 한다. 대각 사상
$\Delta_{\mathfrak{X}}$에 의한 바탕공간 $\mathfrak{X}$의 상이
$\mathfrak{X}\times_{\mathfrak{S}}\mathfrak{X}$의 바탕공간의 닫힌 부분집합이면,
$\mathfrak{X}$가 $\mathfrak{S}$ 위에서 분리되어 있다고, 또는
$\mathfrak{S}$-형식적 스킴이라고, 또는 $f$가 분리 사상이라고 한다. 형식적
준스킴 $\mathfrak{X}$가 분리되어 있다고, 또는 형식적 스킴이라고 하는 것은
$\mathbf{Z}$ 위에서 분리되어 있다는 뜻이다.
\end{definition}

\begin{proposition}[10.15.2]
\label{I.10.15.2-ko}
형식적 준스킴 $\mathfrak{S}$, $\mathfrak{X}$가 각각 보통 준스킴들의 열
$(S_n)$, $(X_n)$의 귀납극한이고, 사상
$f:\mathfrak{X}\to\mathfrak{S}$가 사상들의 열
$f_n:X_n\to S_n$의 귀납극한이라고 하자. $f$가 분리 사상이기 위한
필요충분조건은 $f_0:X_0\to S_0$가 분리 사상인 것이다.
\end{proposition}

실제로, $\Delta_{\mathfrak{X}|\mathfrak{S}}$는 사상들의 열
$\Delta_{X_n|S_n}$의 귀납극한이다
(\hyperref[I.10.7.4-ko]{10.7.4}). 또한 바탕공간 $\mathfrak{X}$의
$\Delta_{\mathfrak{X}|\mathfrak{S}}$에 의한 상(각각 바탕공간
$\mathfrak{X}\times_{\mathfrak{S}}\mathfrak{X}$)은 바탕공간 $X_0$의
$\Delta_{X_0|S_0}$에 의한 상(각각 바탕공간
$X_0\times_{S_0}X_0$)과 동일하다~; 따라서 결론이 따른다.

\begin{proposition}[10.15.3]
\label{I.10.15.3-ko}
이하에서는 고려하는 모든 형식적 준스킴(각각 형식적 준스킴의 모든 사상)이
보통 준스킴들의 열(각각 보통 준스킴의 사상들의 열)의 귀납극한이라고
가정한다.
\begin{enumerate}
  \item[\rm{(i)}] 두 분리 사상의 합성은 분리 사상이다.
  \item[\rm{(ii)}] $f:\mathfrak{X}\to\mathfrak{X}'$,
  $g:\mathfrak{Y}\to\mathfrak{Y}'$가 두 분리
  $\mathfrak{S}$-사상이면 $f\times_{\mathfrak{S}}g$도 분리 사상이다.
  \item[\rm{(iii)}] $f:\mathfrak{X}\to\mathfrak{Y}$가 분리
  $\mathfrak{S}$-사상이면, 밑 형식적 준스킴의 임의의 확장
  $\mathfrak{S}'\to\mathfrak{S}$에 대하여 $\mathfrak{S}'$-사상
  $f_{(\mathfrak{S}')}$도 분리 사상이다.
  \item[\rm{(iv)}] 두 사상의 합성 $g\circ f$가 분리 사상이면 $f$도 분리
  사상이다.
\end{enumerate}
\emph{(이 명제에서 동일한 형식적 준스킴 $\mathfrak{Z}$가 하나의 명제 안에
여러 번 나타나는 경우에는, 나타나는 모든 곳에서 그것을 보통 준스킴들의
동일한 열 $(Z_n)$의 귀납극한으로 간주한다는 것을 함의한다. 또한
형식적 준스킴 $\mathfrak{Z}$에서 형식적 준스킴으로 가는 사상(각각 형식적
준스킴에서 $\mathfrak{Z}$로 가는 사상)은 보통 준스킴의 사상들
 $Z_n$에서 보통 준스킴으로 가는 사상(각각 보통 준스킴에서 $Z_n$으로
가는 사상)의 귀납극한으로 간주한다.)}
\end{proposition}

실제로 (\hyperref[I.10.15.2-ko]{10.15.2})의 표기를 사용하면
$(g\circ f)_0=g_0\circ f_0$이고,
$(f\times_{\mathfrak{S}}g)_0=f_0\times_{S_0}g_0$이다~; 따라서
(\hyperref[I.10.15.3-ko]{10.15.3})의 주장들은 보통 준스킴에 대한
(\hyperref[I.10.15.2-ko]{10.15.2})와
(\hyperref[I.5.5.1-ko]{5.5.1})의 대응하는 주장들의 즉각적인
귀결이다.

독자에게는 (\hyperref[I.10.15.3-ko]{10.15.3})과 같은 종류의 형식적
준스킴 및 사상에 대하여
(\hyperref[I.5.5.5-ko]{5.5.5}),

\oldpage[I]{213}
(\hyperref[I.5.5.9-ko]{5.5.9}) 및
(\hyperref[I.5.5.10-ko]{5.5.10})에 대응하는 명제들을 같은 방식으로
서술하는 일을 맡긴다(여기서 «~아핀 열린집합~»을
«~(\hyperref[I.10.6.3-ko]{10.6.3})의 조건 b)를 만족하는 아핀 형식적
열린집합~»으로 바꾼다).

유사한 논증으로 모든 뇌터 아핀 형식적 스킴이 분리되어 있음도 보인다.
따라서 이 용어가 정당화된다.

\begin{proposition}[10.15.4]
\label{I.10.15.4-ko}
형식적 준스킴 $\mathfrak{S}$를 국소 뇌터라고 하고,
$\mathfrak{X}$, $\mathfrak{Y}$를 두 국소 뇌터
$\mathfrak{S}$-형식적 준스킴이라 하자. $\mathfrak{X}$ 또는
$\mathfrak{Y}$가 $\mathfrak{S}$ 위에서 유한형이라고 가정하자(따라서
$\mathfrak{X}\times_{\mathfrak{S}}\mathfrak{Y}$는 국소 뇌터이다).
또 $\mathfrak{Y}$가 $\mathfrak{S}$ 위에서 분리되어 있다고 하자. 그러면
$\mathfrak{S}$-사상 $f:\mathfrak{X}\to\mathfrak{Y}$의 그래프 사상
$\Gamma_f=(1_{\mathfrak{X}},f)_{\mathfrak{S}}:\mathfrak{X}\to
\mathfrak{X}\times_{\mathfrak{S}}\mathfrak{Y}$는 닫힌 몰입 사상이다.
\end{proposition}

형식적 준스킴 $\mathfrak{S}$가 국소 뇌터 준스킴들의 열
$(S_n)$의 귀납극한이고, $\mathfrak{X}$(각각 $\mathfrak{Y}$)가 열
$(X_n)$(각각 $(Y_n)$)의 $S_n$-준스킴들의 귀납극한이며, $f$가 $S_n$-사상들의 열
$f_n:X_n\to Y_n$의 귀납극한이라고 가정할 수 있다. 그러면
$\mathfrak{X}\times_{\mathfrak{S}}\mathfrak{Y}$는 열
$(X_n\times_{S_n}Y_n)$의 귀납극한이고, $\Gamma_f$는 열
$(\Gamma_{f_n})$의 귀납극한이다
(\hyperref[I.10.7.4-ko]{10.7.4}). 가정에 따라 $Y_0$는 $S_0$ 위에서
분리되어 있다(\hyperref[I.10.15.2-ko]{10.15.2}). 따라서 공간
$\Gamma_{f_0}(X_0)$는 $X_0\times_{S_0}Y_0$의 닫힌 부분공간이다.
$\mathfrak{X}\times_{\mathfrak{S}}\mathfrak{Y}$의 바탕공간(각각
$\Gamma_f(\mathfrak{X})$)과 $X_0\times_{S_0}Y_0$(각각
$\Gamma_{f_0}(X_0)$)의 바탕공간이 서로 같으므로, 이미
$\Gamma_f(\mathfrak{X})$가
$\mathfrak{X}\times_{\mathfrak{S}}\mathfrak{Y}$의 닫힌
부분공간임을 알 수 있다. 이제 $(U,V)$가 $\mathfrak{X}$의 뇌터 아핀
형식적 열린집합 $U$(각각 $\mathfrak{Y}$의 뇌터 아핀 형식적 열린집합
$V$)로 이루어진 순서쌍들 중 $f(U)\subset V$를 만족하는 것들을
취하면, 열린집합 $U\times_{\mathfrak{S}}V$들은
$\mathfrak{X}\times_{\mathfrak{S}}\mathfrak{Y}$ 안에서
$\Gamma_f(\mathfrak{X})$의 덮개를 이룬다. $f_U:U\to V$를 $f$의
$U$로의 제한이라 하면, $\Gamma_{f_U}:U\to
U\times_{\mathfrak{S}}V$는 $U$ 위에서 $\Gamma_f$를 제한한 것이다.
따라서 $\Gamma_{f_U}$가 닫힌 몰입 사상임을 보이면
(\hyperref[I.10.14.4-ko]{10.14.4})에 따라 $\Gamma_f$도 닫힌 몰입
사상이 되며, 문제는 다음 경우로 환원된다:
$\mathfrak{S}=\operatorname{Spf}(A)$,
$\mathfrak{X}=\operatorname{Spf}(B)$,
$\mathfrak{Y}=\operatorname{Spf}(C)$가 아핀이고
($A$, $B$, $C$는 뇌터 아딕환), $f$가 연속 $A$-준동형
$\varphi:C\to B$에 대응하는 경우이다. 이때 $\Gamma_f$는 유일한 연속
준동형
$\omega:B\widehat{\otimes}_A C\to B$에 대응한다. 표준 준동형
$B\to B\widehat{\otimes}_A C$ 및
$C\to B\widehat{\otimes}_A C$를 이 유일한 준동형과 각각 합성하면 항등사상과
$\varphi$를 각각 얻는다. 그런데 $\omega$가
\emph{전사}임은 명백하므로 주장이 따른다.

\begin{corollary}[10.15.5]
\label{I.10.15.5-ko}
형식적 준스킴 $\mathfrak{S}$를 국소 뇌터라고 하고,
$\mathfrak{X}$를 $\mathfrak{S}$ 위의 유한형 형식적 준스킴이라 하자.
$\mathfrak{X}$가 $\mathfrak{S}$ 위에서 분리되어 있기 위한 필요충분조건은
대각 사상
$\mathfrak{X}\to\mathfrak{X}\times_{\mathfrak{S}}\mathfrak{X}$가
닫힌 몰입 사상인 것이다.
\end{corollary}

\begin{proposition}[10.15.6]
\label{I.10.15.6-ko}
국소 뇌터 형식적 준스킴들의 닫힌 몰입 사상
$j:\mathfrak{Y}\to\mathfrak{X}$는 분리 사상이다.
\end{proposition}

(\hyperref[I.10.14.2-ko]{10.14.2})의 표기를 사용하면
$j_0:Y_0\to X_0$는 닫힌 몰입 사상이고, 따라서 분리 사상이다. 이제
(\hyperref[I.10.15.2-ko]{10.15.2})를 적용하면 충분하다.

\begin{proposition}[10.15.7]
\label{I.10.15.7-ko}
$X$를 (보통 의미의) 국소 뇌터 준스킴, $X'$을 $X$의 닫힌 부분집합,
$\widehat{X}=X_{/X'}$라 하자. $\widehat{X}$가 분리되기 위한
필요충분조건은 $\widehat{X}_{\mathrm{red}}$가 분리되는 것이며,
특히 $X$가 분리되어 있으면 이 조건이 충족된다.
\end{proposition}

실제로, (\hyperref[I.10.8.5-ko]{10.8.5})의 표기를 사용하면
$\widehat{X}$가 분리되기 위한 필요충분조건은 $X'_0$가 분리되는
것이다(\hyperref[I.10.15.2-ko]{10.15.2}). 그런데
$\widehat{X}_{\mathrm{red}}=(X'_0)_{\mathrm{red}}$이므로,
$\widehat{X}_{\mathrm{red}}$가 분리된다고 말하는 것과 동치이다
(\hyperref[I.5.5.1-ko]{5.5.1, (vi)}).

% Bibliographic data inherited verbatim from the current French authority.
\begin{thebibliography}{22}

\bibitem{I-1}
\oldpage[I]{214}
H.~\textsc{Cartan} et C.~\textsc{Chevalley},
\emph{Séminaire de l'École Normale Supérieure},
8\textsuperscript{e} année (1955-56), \emph{Géométrie algébrique}.

\bibitem{I-2}
H.~\textsc{Cartan} and S.~\textsc{Eilenberg},
\emph{Homological Algebra}, Princeton Math. Series (Princeton University
Press), 1956.

\bibitem{I-3}
W.~L.~\textsc{Chow} and J.~\textsc{Igusa},
\emph{Cohomology theory of varieties over rings},
\emph{Proc. Nat. Acad. Sci. U.S.A.}, t.~XLIV (1958), p.~1244-1248.

\bibitem{I-4}
R.~\textsc{Godement}, \emph{Théorie des faisceaux}, Actual. Scient. et Ind.,
n\textsuperscript{o}~1252, Paris (Hermann), 1958.

\bibitem{I-5}
H.~\textsc{Grauert},
\emph{Ein Theorem der analytischen Garbentheorie und die Modulräume komplexer
Strukturen}, \emph{Publ. Math. Inst. Hautes Études Scient.},
n\textsuperscript{o}~5, 1960.

\bibitem{I-6}
A.~\textsc{Grothendieck}, \emph{Sur quelques points d'algèbre homologique},
\emph{Tôhoku Math. Journ.}, t.~IX (1957), p.~119-221.

\bibitem{I-7}
A.~\textsc{Grothendieck},
\emph{Cohomology theory of abstract algebraic varieties},
\emph{Proc. Intern. Congress of Math.}, p.~103-118, Edinburgh (1958).

\bibitem{I-8}
A.~\textsc{Grothendieck},
\emph{Géométrie formelle et géométrie algébrique},
\emph{Séminaire Bourbaki}, 11\textsuperscript{e} année (1958-59), exposé 182.

\bibitem{I-9}
M.~\textsc{Nagata},
\emph{A general theory of algebraic geometry over Dedekind domains},
\emph{Amer. Math. Journ.}~: I, t.~LXXVIII, p.~78-116 (1956)~; II, t.~LXXX,
p.~382-420 (1958).

\bibitem{I-10}
D.~G.~\textsc{Northcott}, \emph{Ideal theory}, Cambridge Univ. Press, 1953.

\bibitem{I-11}
P.~\textsc{Samuel}, \emph{Commutative algebra} (Notes by D.~Herzig), Cornell
Univ., 1953.

\bibitem{I-12}
P.~\textsc{Samuel}, \emph{Algèbre locale}, Mém. Sci. Math.,
n\textsuperscript{o}~123, Paris, 1953.

\bibitem{I-13}
P.~\textsc{Samuel} and O.~\textsc{Zariski},
\emph{Commutative algebra}, 2 vol., New York (Van Nostrand), 1958-60.

\bibitem{I-14}
J.-P.~\textsc{Serre}, \emph{Faisceaux algébriques cohérents},
\emph{Ann. of Math.}, t.~LXI (1955), p.~197-278.

\bibitem{I-15}
J.-P.~\textsc{Serre},
\emph{Sur la cohomologie des variétés algébriques},
\emph{Journ. de Math.} (9), t.~XXXVI (1957), p.~1-16.

\bibitem{I-16}
J.-P.~\textsc{Serre},
\emph{Géométrie algébrique et géométrie analytique},
\emph{Ann. Inst. Fourier}, t.~VI (1955-56), p.~1-42.

\bibitem{I-17}
J.-P.~\textsc{Serre},
\emph{Sur la dimension homologique des anneaux et des modules noethériens},
\emph{Proc. Intern. Symp. on Alg. Number theory}, p.~176-189, Tokyo-Nikko,
1955.

\bibitem{I-18}
A.~\textsc{Weil}, \emph{Foundations of algebraic geometry}, Amer. Math. Soc.
Coll. Publ., n\textsuperscript{o}~29, 1946.

\bibitem{I-19}
A.~\textsc{Weil}, \emph{Numbers of solutions of equations in finite fields},
\emph{Bull. Amer. Math. Soc.}, t.~LV (1949), p.~497-508.

\bibitem{I-20}
O.~\textsc{Zariski},
\emph{Theory and applications of holomorphic functions on algebraic varieties
over arbitrary ground fields}, Mem. Amer. Math. Soc., n\textsuperscript{o}~5
(1951).

\bibitem{I-21}
O.~\textsc{Zariski}, \emph{A new proof of Hilbert's Nullstellensatz},
\emph{Bull. Amer. Math. Soc.}, t.~LIII (1947), p.~362-368.

\bibitem{I-22}
E.~\textsc{Kähler}, \emph{Geometria Arithmetica},
\emph{Ann. di Mat.} (4), t.~XLV (1958), p.~1-368.

\end{thebibliography}

% Authority: EGA II front-fr.tex, SHA-256 F199DE0C2C8C49E2BDAEF48D764086952806E77D8066795B97B3595F7EA03198.
\clearpage
\thispagestyle{plain}
\markboth{제II장}{제II장}
\oldpage[II]{5}
\phantomsection
\label{chapter:II-ko}

\begin{center}
{\large 제II장\par}
\vspace{1.5em}
{\Large\bfseries 몇몇 사상 부류에 대한\par}
{\Large\bfseries 초보적 전역 연구\par}
\vspace{1.8em}
{\bfseries 목차\par}
\end{center}

\medskip
\begin{tabular}{@{}r@{\ }p{0.82\textwidth}@{}}
\S~1. & 아핀 사상.\\
\S~2. & 동차 소 스펙트럼.\\
\S~3. & 등급 대수의 층의 동차 스펙트럼.\\
\S~4. & 사영 다발. 풍부한 층.\\
\S~5. & 준아핀 사상; 준사영 사상; 고유 사상; 사영 사상.\\
\S~6. & 정수적 사상과 유한 사상.\\
\S~7. & 값매김 판정법.\\
\S~8. & 블로업 스킴; 사영 원뿔; 사영 폐포.
\end{tabular}

\medskip
이 장에서 다루는 여러 사상 부류는 코호몰로지적 방법을 크게 사용하지
않고 연구한다. 이 방법을 사용하여 그 사상 부류들을 더 깊이 연구하는
일은 제III장에서 이루어질 것이며, 거기서는 제II장의 \S\S~2, 4, 5를
주로 사용할 것이다. \S~8은 처음 읽을 때 생략해도 된다. 이 절은
\S\S~1에서~3까지 전개한 형식 체계에 몇 가지 보충을 더하는데, 그
내용은 이 형식 체계의 쉬운 응용으로 귀착된다. 우리는 이 절의 결과들을
이 장의 다른 결과들만큼 자주 사용하지는 않을 것이다.

% Authority slice: EGA II ega2-1-fr.tex lines 1-4040, 173428 LF UTF-8 bytes, SHA-256 E6911F259382EBEFA590F25970FD3F47297E1EC78D8C4703AD1627DA0BA4C720.
\setcounter{section}{0}
\section{아핀 사상}
\label{section:II.1-ko}

이 절의 결과 대부분은 제I장 \S~1의 ``전역적'' 대응물이다. 따라서
본질적으로 새로운 것은 아니며, 이후를 위한 편리한 언어를 제공할 뿐이다.

\setcounter{subsection}{0}
\subsection{$S$-준스킴과 $\mathcal{O}_S$-대수}
\label{subsection:II.1.1-ko}

\begin{env}[1.1.1]
\label{II.1.1.1-ko}
$S$를 준스킴, $X$를 $S$-준스킴, $f:X\to S$를 그 구조 사상이라 하자.
(\hyperref[0.4.2.4-ko]{0, 4.2.4})에 따라 직상
$f_*(\mathcal{O}_X)$는 $\mathcal{O}_S$-대수이다. 혼동의 우려가 없으면
이를
\oldpage[II]{6}
$\mathcal{A}(X)$로 나타내기로 한다. $U$가 $S$의 열린집합이면
\[
  \mathcal{A}(f^{-1}(U))=\mathcal{A}(X)|U.
\]
마찬가지로 임의의 $\mathcal{O}_X$-가군 $\mathcal{F}$(각각
$\mathcal{O}_X$-대수 $\mathcal{B}$)에 대하여, 직상
$f_*(\mathcal{F})$(각각 $f_*(\mathcal{B})$)를
$\mathcal{A}(\mathcal{F})$(각각 $\mathcal{A}(\mathcal{B})$)로
나타내기로 한다. 이는 단순히 $\mathcal{O}_S$-가군(각각
$\mathcal{O}_S$-대수)일 뿐만 아니라 $\mathcal{A}(X)$-가군(각각
$\mathcal{A}(X)$-대수)이다.
\end{env}

\begin{env}[1.1.2]
\label{II.1.1.2-ko}
$Y$를 또 하나의 $S$-준스킴, $g:Y\to S$를 그 구조 사상,
$h:X\to Y$를 $S$-사상이라 하자. 그러면 다음 도식은 가환한다.
\[
  \xymatrix{
    X\ar[rr]^h\ar[rd]_f && Y\ar[ld]^g\\
    & S
  }
  \tag{1.1.2.1}\label{II.1.1.2.1-ko}
\]

정의상 $h=(\psi,\theta)$이고, 여기서
$\theta:\mathcal{O}_Y\to h_*(\mathcal{O}_X)=\psi_*(\mathcal{O}_X)$는
환의 층의 준동형이다. 이로부터
(\hyperref[0.4.2.2-ko]{0, 4.2.2})에 따라 다음
$\mathcal{O}_S$-대수 준동형을 얻는다.
\[
  g_*(\theta):g_*(\mathcal{O}_Y)\to
  g_*(h_*(\mathcal{O}_X))=f_*(\mathcal{O}_X),
\]
다시 말해 $\mathcal{O}_S$-대수 준동형
$\mathcal{A}(Y)\to\mathcal{A}(X)$를 얻으며, 이를
$\mathcal{A}(h)$로도 나타낸다. $h':Y\to Z$가 또 하나의
$S$-사상이면
$\mathcal{A}(h'\circ h)=\mathcal{A}(h)\circ\mathcal{A}(h')$임은
곧바로 알 수 있다. 따라서 $S$-준스킴의 범주 위에 정의되어
$\mathcal{O}_S$-대수의 범주에 값을 갖는 \emph{반변 함자}
$\mathcal{A}(X)$를 정의하였다.

이제 $\mathcal{F}$를 $\mathcal{O}_X$-가군, $\mathcal{G}$를
$\mathcal{O}_Y$-가군이라 하고, $u:\mathcal{G}\to\mathcal{F}$를
$h$-사상이라 하자. 다시 말해
(\hyperref[0.4.4.1-ko]{0, 4.4.1})에 따라 $u$는
$\mathcal{O}_Y$-가군 준동형
$\mathcal{G}\to h_*(\mathcal{F})$이다. 그러면
\[
  g_*(u):g_*(\mathcal{G})\to g_*(h_*(\mathcal{F}))=f_*(\mathcal{F})
\]
는 $\mathcal{O}_S$-가군 준동형
$\mathcal{A}(\mathcal{G})\to\mathcal{A}(\mathcal{F})$이며, 이를
$\mathcal{A}(u)$로 나타낸다. 더욱이 쌍
$(\mathcal{A}(h),\mathcal{A}(u))$는
$\mathcal{A}(Y)$-가군 $\mathcal{A}(\mathcal{G})$에서
$\mathcal{A}(X)$-가군 $\mathcal{A}(\mathcal{F})$로 가는
\emph{디-준동형}을 이룬다.
\end{env}

\begin{env}[1.1.3]
\label{II.1.1.3-ko}
$S$를 고정하면, $X$가 $S$-준스킴이고 $\mathcal{F}$가
$\mathcal{O}_X$-가군인 쌍 $(X,\mathcal{F})$들이 하나의
\emph{범주}를 이룬다고 볼 수 있다. 여기서
$(X,\mathcal{F})$에서 $(Y,\mathcal{G})$로 가는 \emph{사상}은 쌍
$(h,u)$로 정의한다. 이때 $h:X\to Y$는 $S$-사상이고
$u:\mathcal{G}\to\mathcal{F}$는 $h$-사상이다. 그러면
$(\mathcal{A}(X),\mathcal{A}(\mathcal{F}))$는
$\mathcal{O}_S$-대수와 그 대수 위의 가군으로 이루어진 쌍들을
대상으로 하고 디-준동형들을 사상으로 하는 범주에 값을 갖는
\emph{반변 함자}라고 할 수 있다.
\end{env}

\subsection{준스킴 위에서 아핀인 준스킴}
\label{subsection:II.1.2-ko}

\begin{env}[1.2.1]
\label{II.1.2.1-ko}
$X$를 $S$-준스킴, $f:X\to S$를 그 구조 사상이라 하자.
$S$의 아핀 열린집합들로 이루어진 덮개 $(S_\alpha)$가 존재하여,
모든 $\alpha$에 대하여 열린집합 $f^{-1}(S_\alpha)$ 위에 $X$가
유도하는 준스킴이 아핀이면, $X$가 \emph{$S$ 위에서 아핀}이라고
한다.
\end{env}

\begin{env}[1.2.2]
\label{II.1.2.2-ko}
$S$의 모든 닫힌 부분준스킴은 $S$ 위에서 아핀인
$S$-준스킴이다
(\hyperref[I.4.2.3-ko]{I, 4.2.3}과
\hyperref[I.4.2.4-ko]{4.2.4}).
\end{env}

\begin{env}[1.2.3]
\label{II.1.2.3-ko}
$S$ 위에서 아핀인 준스킴 $X$가 반드시 아핀 스킴인 것은 아니다. 예
$X=S$가 이를 보여 준다(\hyperref[II.1.2.2-ko]{1.2.2}). 한편 아핀 스킴
$X$가 $S$-준스킴이라 해도, $X$가 반드시 $S$ 위에서 아핀인
\oldpage[II]{7}
것은 아니다(\hyperref[II.1.3.3-ko]{1.3.3} 참조). 그러나 $S$가
\emph{스킴}이면, 아핀 스킴인 모든 $S$-준스킴은 $S$ 위에서 아핀임을
상기하자(\hyperref[I.5.5.10-ko]{I, 5.5.10}).
\end{env}

\begin{env}[1.2.4]
\label{II.1.2.4-ko}
$S$ 위에서 아핀인 모든 $S$-준스킴은 $S$ 위에서 분리되어 있다
(다시 말해 $S$-스킴이다).
\end{env}

이는 (\hyperref[I.5.5.5-ko]{I, 5.5.5})와
(\hyperref[I.5.5.8-ko]{I, 5.5.8})에서 곧바로 따른다.

\begin{env}[1.2.5]
\label{II.1.2.5-ko}
$X$를 $S$ 위에서 아핀인 $S$-스킴, $f:X\to S$를 그 구조 사상이라
하자. 모든 열린집합 $U\subset S$에 대하여, $f^{-1}(U)$는 $U$ 위에서
아핀이다.
\end{env}

정의 (\hyperref[II.1.2.1-ko]{1.2.1})에 따라,
$S=\operatorname{Spec}(A)$와 $X=\operatorname{Spec}(B)$가 아핀인
경우로 환원된다. 이때 $f=({}^a\varphi,\widetilde{\varphi})$이고,
여기서 $\varphi:A\to B$는 환 준동형이다. $g\in A$인 $D(g)$들이
$S$의 기저를 이루므로, $U=D(g)$인 경우로 환원된다. 그런데 이때
$f^{-1}(U)=D(\varphi(g))$임을 알고 있으므로
(\hyperref[I.1.2.2.2-ko]{I, 1.2.2.2}), 명제가 따른다.

\begin{env}[1.2.6]
\label{II.1.2.6-ko}
$X$를 $S$ 위에서 아핀인 $S$-스킴, $f:X\to S$를 그 구조 사상이라
하자. 임의의 준연접 $\mathcal{O}_X$-가군 $\mathcal{F}$에 대하여,
$f_*(\mathcal{F})$는 준연접 $\mathcal{O}_S$-가군이다.
\end{env}

(\hyperref[II.1.2.4-ko]{1.2.4})를 고려하면, 이는
(\hyperref[I.9.2.2-ko]{I, 9.2.2, a)})에서 따른다.

특히 $\mathcal{O}_S$-대수
$\mathcal{A}(X)=f_*(\mathcal{O}_X)$는 \emph{준연접}이다.

\begin{env}[1.2.7]
\label{II.1.2.7-ko}
$X$를 $S$ 위에서 아핀인 $S$-스킴이라 하자. 임의의 $S$-준스킴
$Y$에 대하여, 집합 $\operatorname{Hom}_S(Y,X)$에서 집합
$\operatorname{Hom}(\mathcal{A}(X),\mathcal{A}(Y))$로 가는 사상
$h\mapsto\mathcal{A}(h)$
(\hyperref[II.1.1.2-ko]{1.1.2})은 전단사이다.
\end{env}

$f:X\to S$와 $g:Y\to S$를 구조 사상이라 하자. 먼저
$S=\operatorname{Spec}(A)$와 $X=\operatorname{Spec}(B)$가
아핀이라고 하자. 임의의 $\mathcal{O}_S$-대수 준동형
$\omega:f_*(\mathcal{O}_X)\to g_*(\mathcal{O}_Y)$에 대하여,
$\mathcal{A}(h)=\omega$인 $S$-사상 $h:Y\to X$가 유일하게 존재함을
보여야 한다. 정의에 따라, 모든 열린집합 $U\subset S$에 대하여
$\omega$는 다음 $\Gamma(U,\mathcal{O}_S)$-대수 준동형을 정한다.
\[
  \omega_U=\Gamma(U,\omega):
  \Gamma(f^{-1}(U),\mathcal{O}_X)\to
  \Gamma(g^{-1}(U),\mathcal{O}_Y)
\]
특히 $U=S$일 때, 이는 $\Gamma(S,\mathcal{O}_S)$-대수 준동형
$\varphi:\Gamma(X,\mathcal{O}_X)\to\Gamma(Y,\mathcal{O}_Y)$를
준다. $X$가 아핀이므로, 이 준동형에는 유일하게 정해지는
$S$-사상 $h:Y\to X$가 대응한다
(\hyperref[I.2.2.4-ko]{I, 2.2.4}). 이제
$\mathcal{A}(h)=\omega$임을 보여야 한다. 다시 말해, $S$의 한 기저에
속하는 모든 열린집합 $U$에 대하여 $\omega_U$가 $h$의 제한인
$S$-사상 $g^{-1}(U)\to f^{-1}(U)$에 대응하는 대수 준동형
$\varphi_U$와 일치함을 보이면 된다. $\lambda\in A$인
$U=D(\lambda)$의 경우만 살피면 충분하다.\footnote{\textbf{번역 주.}
인쇄 원문은 $\lambda\in S$라고 쓰지만, $S=\operatorname{Spec}(A)$이고
$\rho:A\to B$이며 뒤에서 $\rho(\lambda)$를 사용하므로
EGA2-CANON-DISP-0001에 따라 $\lambda\in A$로 교정했다.} 이때
$f=({}^a\rho,\widetilde{\rho})$이고 $\rho:A\to B$가 환 준동형이면,
$f^{-1}(U)=D(\mu)$이며 여기서 $\mu=\rho(\lambda)$이다. 또한
$\Gamma(f^{-1}(U),\mathcal{O}_X)$는 분수환 $B_\mu$이다. 그런데 다음
도식은 가환한다.
\[
  \xymatrix{
    B\ar[r]^\varphi\ar[d] & \Gamma(Y,\mathcal{O}_Y)\ar[d]\\
    B_\mu\ar[r]_{\varphi_U} & \Gamma(g^{-1}(U),\mathcal{O}_Y)
  }
\]
$\varphi_U$를 $\omega_U$로 바꾼 유사한 도식도 가환한다. 따라서
분수환의 보편 성질
(\hyperref[0.1.2.4-ko]{0, 1.2.4})에 의해
$\varphi_U=\omega_U$이다.

이제 일반적인 경우를 보자. $(S_\alpha)$를 $S$의 아핀 열린집합들로
이루어진 덮개라 하자.
\oldpage[II]{8}
여기서 각 $f^{-1}(S_\alpha)$는 아핀이라고 하자. 그러면 임의의
$\mathcal{O}_S$-대수 준동형
$\omega:\mathcal{A}(X)\to\mathcal{A}(Y)$는 제한에 의해 다음
$\mathcal{O}_{S_\alpha}$-대수 준동형들의 족을 정한다.
\[
  \omega_\alpha:
  \mathcal{A}(f^{-1}(S_\alpha))\to
  \mathcal{A}(g^{-1}(S_\alpha))
\]
따라서 앞에서 보인 바에 의해 $S_\alpha$-사상들의 족
$h_\alpha:g^{-1}(S_\alpha)\to f^{-1}(S_\alpha)$가 정해진다. 이제
$S_\alpha\cap S_\beta$의 한 기저에 속하는 모든 아핀 열린집합 $U$에
대하여, $h_\alpha$와 $h_\beta$를 $g^{-1}(U)$에 제한한 것들이
일치함을 보이면 된다. 이는 명백하다. 실제로 앞에서 보인 바에 의해
이 두 제한은 모두 $\omega$의 제한인 준동형
$\mathcal{A}(X)|U\to\mathcal{A}(Y)|U$에 대응한다.

\begin{env}[1.2.8]
\label{II.1.2.8-ko}
$X$와 $Y$를 $S$ 위에서 아핀인 두 $S$-스킴이라 하자. $S$-사상
$h:Y\to X$가 동형사상이기 위한 필요충분조건은
$\mathcal{A}(h):\mathcal{A}(X)\to\mathcal{A}(Y)$가 동형사상인 것이다.
\end{env}

이는 (\hyperref[II.1.2.7-ko]{1.2.7})과 $\mathcal{A}(X)$의 함자성에서
곧바로 따른다.

\subsection{$\mathcal{O}_S$-대수에 대응하는 $S$ 위에서 아핀인 준스킴}
\label{subsection:II.1.3-ko}

\begin{env}[1.3.1]
\label{II.1.3.1-ko}
$S$를 준스킴이라 하자. 임의의 준연접 $\mathcal{O}_S$-대수
$\mathcal{B}$에 대하여, $\mathcal{A}(X)=\mathcal{B}$를 만족하는
$S$ 위에서 아핀인 준스킴 $X$가 존재하며, $X$는 유일한
$S$-동형까지 유일하게 결정된다.
\end{env}

유일성은 (\hyperref[II.1.2.8-ko]{1.2.8})에서 곧바로 따른다.
$X$의 존재를 증명하자. 임의의 아핀 열린집합 $U\subset S$에 대하여
$X_U$를 준스킴
$\operatorname{Spec}(\Gamma(U,\mathcal{B}))$라 하자.
$\Gamma(U,\mathcal{B})$는 $\Gamma(U,\mathcal{O}_S)$-대수이므로
$X_U$는 $S$-준스킴이다
(\hyperref[I.1.6.1-ko]{I, 1.6.1}). 또한 $\mathcal{B}$가
준연접이므로 $\mathcal{O}_S$-대수 $\mathcal{A}(X_U)$는
$\mathcal{B}|U$와 표준적으로 동일시된다
(\hyperref[I.1.3.7-ko]{I, 1.3.7},
\hyperref[I.1.3.13-ko]{1.3.13}과
\hyperref[I.1.6.3-ko]{1.6.3}).

$V$를 $S$의 두 번째 아핀 열린집합이라 하고, $f_U$가 구조 사상
$X_U\to S$를 나타낼 때 $X_{U,V}$를 $X_U$가
$f_U^{-1}(U\cap V)$ 위에 유도하는 준스킴이라 하자.
$X_{U,V}$와 $X_{V,U}$는 $U\cap V$ 위에서 아핀이고
(\hyperref[II.1.2.5-ko]{1.2.5}), 정의에 따라
$\mathcal{A}(X_{U,V})$와 $\mathcal{A}(X_{V,U})$는
$\mathcal{B}|(U\cap V)$와 표준적으로 동일시된다. 그러므로
(\hyperref[II.1.2.8-ko]{1.2.8})에 따라 표준 $S$-동형사상
$\theta_{U,V}:X_{V,U}\to X_{U,V}$가 존재한다. 더욱이 $W$를
$S$의 세 번째 아핀 열린집합이라 하고, 구조 사상들에 의해
$X_V$, $X_W$, $X_W$ 안에서 각각 얻는 $U\cap V\cap W$의
역상들로 $\theta_{U,V}$, $\theta_{V,W}$, $\theta_{U,W}$를
제한한 사상을 차례로 $\theta'_{U,V}$, $\theta'_{V,W}$,
$\theta'_{U,W}$라 하면
$\theta'_{U,V}\circ\theta'_{V,W}=\theta'_{U,W}$이다.

따라서 준스킴 $X$, 아핀 열린집합들로 이루어진 $X$의 덮개
$(T_U)$, 그리고 각 $U$에 대한 동형사상
$\varphi_U:X_U\to T_U$가 존재하여, $\varphi_U$는
$f_U^{-1}(U\cap V)$를 $T_U\cap T_V$ 위로 보내고
$\theta_{U,V}=\varphi_U^{-1}\circ\varphi_V$가 성립한다
(\hyperref[I.2.3.1-ko]{I, 2.3.1}). 사상
$g_U=f_U\circ\varphi_U^{-1}$는 $T_U$에 $S$-준스킴 구조를 주며,
$g_U$와 $g_V$는 $T_U\cap T_V$에서 일치한다. 따라서 $X$는
$S$-준스킴이다. 또한 정의상 $X$가 $S$ 위에서 아핀이고
$\mathcal{A}(T_U)=\mathcal{B}|U$임은 명백하므로
$\mathcal{A}(X)=\mathcal{B}$이다.

이렇게 정의한 $S$-스킴 $X$가
\emph{$\mathcal{O}_S$-대수 $\mathcal{B}$에 대응한다}고 하거나
$X$를 \emph{$\mathcal{B}$의 스펙트럼}이라 하며, 이를
$\operatorname{Spec}(\mathcal{B})$로 나타낸다.

\begin{env}[1.3.2]
\label{II.1.3.2-ko}
$X$를 $S$ 위에서 아핀인 준스킴, $f:X\to S$를 그 구조 사상이라
하자. 임의의 아핀 열린집합 $U\subset S$에 대하여
$f^{-1}(U)$ 위에 유도된 준스킴은 환
$\Gamma(U,\mathcal{A}(X))$의 아핀 스킴이다.
\end{env}

\oldpage[II]{9}
(\hyperref[II.1.2.6-ko]{1.2.6})과
(\hyperref[II.1.3.1-ko]{1.3.1})에 따라 $X$가 어떤
$\mathcal{O}_S$-대수에 대응한다고 가정해도 되므로, 이 따름정리는
(\hyperref[II.1.3.1-ko]{1.3.1})에서 서술한 $X$의 구성에서 따른다.

\begin{env}[1.3.3]
\label{II.1.3.3-ko}
$S$를 점 $0$을 이중화한 체 $K$ 위의 아핀 평면이라 하자
(\hyperref[I.5.5.11-ko]{I, 5.5.11}). 그 표기에서 $S$는 두 아핀
열린집합 $Y_1$, $Y_2$의 합집합이다
(\hyperref[I.5.5.11-ko]{I, 5.5.11}). $f$가 열린 몰입 사상
$Y_1\to S$이면
$f^{-1}(Y_2)=Y_1\cap Y_2$는 $Y_1$ 안의 아핀 열린집합이 아니다
(\emph{loc. cit.}). 따라서 이로써 아핀 스킴이지만 $S$ 위에서
아핀이 아닌 예를 얻는다.
\end{env}

\begin{env}[1.3.4]
\label{II.1.3.4-ko}
$S$를 아핀 스킴이라 하자. $S$-준스킴 $X$가 $S$ 위에서
아핀이기 위한 필요충분조건은 $X$가 아핀 스킴인 것이다.
\end{env}

\begin{env}[1.3.5]
\label{II.1.3.5-ko}
$S$를 준스킴, $X$를 $S$ 위에서 아핀인 준스킴, $Y$를
$X$-준스킴이라 하자. $Y$가 $S$ 위에서 아핀이기 위한
필요충분조건은 $Y$가 $X$ 위에서 아핀인 것이다.
\end{env}

곧바로 $S$가 아핀 스킴인 경우로 환원된다. 그러면 $X$도 아핀
스킴이다(\hyperref[II.1.3.4-ko]{1.3.4}). 이때 명제의 두 조건은
모두 $Y$가 아핀 스킴이라는 것을 나타낸다
(\hyperref[II.1.3.4-ko]{1.3.4}).

\begin{env}[1.3.6]
\label{II.1.3.6-ko}
$X$를 $S$ 위에서 아핀인 준스킴이라 하자.
(\hyperref[II.1.3.5-ko]{1.3.5})에 따라, $X$ 위에서 아핀인
준스킴 $Y$를 정하는 것은 $S$ 위에서 아핀인 준스킴 $Y$와
$S$-사상 $g:Y\to X$를 정하는 것과 같다. 다시 말해
(\hyperref[II.1.3.1-ko]{1.3.1}과
\hyperref[II.1.2.7-ko]{1.2.7}), 이는 준연접
$\mathcal{O}_S$-대수 $\mathcal{B}$와 $\mathcal{O}_S$-대수 준동형
$\mathcal{A}(X)\to\mathcal{B}$를 정하는 것과 같다(이를
$\mathcal{B}$에 $\mathcal{A}(X)$-대수 구조를 정하는 것으로 볼 수
있다). $f:X\to S$가 구조 사상이면, 이때
$\mathcal{B}=f_*(g_*(\mathcal{O}_Y))$이다.
\end{env}

\begin{env}[1.3.7]
\label{II.1.3.7-ko}
$X$를 $S$ 위에서 아핀인 준스킴이라 하자. $X$가 $S$ 위에서
유한형이기 위한 필요충분조건은 준연접 $\mathcal{O}_S$-대수
$\mathcal{A}(X)$가 유한 생성인 것이다
(\hyperref[I.9.6.2-ko]{I, 9.6.2}).
\end{env}

정의 (\hyperref[I.9.6.2-ko]{I, 9.6.2})에 따라 $S$가 아핀인
경우로 환원된다. 그러면 $X$는 아핀 스킴이고
(\hyperref[II.1.3.4-ko]{1.3.4}), $S=\operatorname{Spec}(A)$,
$X=\operatorname{Spec}(B)$이면 $\mathcal{A}(X)$는
$\mathcal{O}_S$-대수 $\widetilde{B}$이다. 그런데
$\Gamma(S,\widetilde{B})=B$이므로,\footnote{\textbf{번역 주.} 인쇄
원문은 이 증명에서 정의되지 않은 $U$를 쓰지만,
EGA2-CANON-DISP-0030에 따라 $S$로 교정했다.} 이 따름정리는
(\hyperref[I.9.6.2-ko]{I, 9.6.2})와
(\hyperref[I.6.3.3-ko]{I, 6.3.3})에서 따른다.

\begin{env}[1.3.8]
\label{II.1.3.8-ko}
$X$를 $S$ 위에서 아핀인 준스킴이라 하자. $X$가 축소이기 위한
필요충분조건은 준연접 $\mathcal{O}_S$-대수 $\mathcal{A}(X)$가
축소인 것이다
(\hyperref[0.4.1.4-ko]{0, 4.1.4}).
\end{env}

실제로 이 문제는 명백히 $S$ 위에서 국소적이다.
(\hyperref[II.1.3.2-ko]{1.3.2})에 따라, 이 따름정리는
(\hyperref[I.5.1.1-ko]{I, 5.1.1})과
(\hyperref[I.5.1.4-ko]{I, 5.1.4})에서 따른다.

\subsection{$S$ 위에서 아핀인 준스킴 위의 준연접층}
\label{subsection:II.1.4-ko}

\begin{env}[1.4.1]
\label{II.1.4.1-ko}
$X$를 $S$ 위에서 아핀인 준스킴, $Y$를 $S$-준스킴,
$\mathcal{F}$(각각 $\mathcal{G}$)를 준연접
$\mathcal{O}_X$-가군(각각 $\mathcal{O}_Y$-가군)이라 하자. 그러면
$(h,u)\mapsto(\mathcal{A}(h),\mathcal{A}(u))$로 주어지는, 쌍
$(Y,\mathcal{G})$에서 $(X,\mathcal{F})$로 가는 사상들의 집합에서 쌍
$(\mathcal{A}(X),\mathcal{A}(\mathcal{F}))$에서
$(\mathcal{A}(Y),\mathcal{A}(\mathcal{G}))$로 가는 디-준동형들의
집합으로 가는 사상
(\hyperref[II.1.1.2-ko]{1.1.2}와
\hyperref[II.1.1.3-ko]{1.1.3})은 전단사이다.
\end{env}

증명은 (\hyperref[I.2.2.5-ko]{I, 2.2.5})와
(\hyperref[I.2.2.4-ko]{I, 2.2.4})를 사용하여
(\hyperref[II.1.2.7-ko]{1.2.7})의 증명과 정확히 같은 절차를 따르며,
세부사항은 독자에게 맡긴다.

\begin{env}[1.4.2]
\label{II.1.4.2-ko}
(\hyperref[II.1.4.1-ko]{1.4.1})의 가정들에 더하여 $Y$가 $S$ 위에서
아핀이라고 가정하면, $(h,u)$가 동형이기 위한 필요충분조건은
$(\mathcal{A}(h),\mathcal{A}(u))$가 디-동형인 것이다.
\end{env}

\oldpage[II]{10}
\begin{env}[1.4.3]
\label{II.1.4.3-ko}
준연접 $\mathcal{O}_S$-대수 $\mathcal{B}$와 준연접
$\mathcal{B}$-가군 $\mathcal{M}$으로 이루어진 모든 쌍
$(\mathcal{B},\mathcal{M})$에 대하여
($\mathcal{M}$을 $\mathcal{O}_S$-가군으로 보든
$\mathcal{B}$-가군으로 보든 준연접성은 같은 조건이다
(\hyperref[I.9.6.1-ko]{I, 9.6.1})), $S$ 위에서 아핀인 준스킴 $X$와
준연접 $\mathcal{O}_X$-가군 $\mathcal{F}$로 이루어진 쌍
$(X,\mathcal{F})$가 존재하여
$\mathcal{A}(X)=\mathcal{B}$이고
$\mathcal{A}(\mathcal{F})=\mathcal{M}$이다. 더욱이 이 쌍은 유일한
동형까지 유일하게 결정된다.
\end{env}

유일성은 (\hyperref[II.1.4.1-ko]{1.4.1})과
(\hyperref[II.1.4.2-ko]{1.4.2})에서 따른다. 존재성은
(\hyperref[II.1.3.1-ko]{1.3.1})에서와 같이 증명되며, 여기서도
세부사항은 독자에게 맡긴다. $\mathcal{O}_X$-가군 $\mathcal{F}$를
$\widetilde{\mathcal{M}}$로 나타내고, 이 가군을 준연접
$\mathcal{B}$-가군 $\mathcal{M}$에 \emph{대응하는} 가군이라고 한다.
모든 아핀 열린집합 $U\subset S$에 대하여
$\widetilde{\mathcal{M}}|p^{-1}(U)$는
($p$가 구조 사상 $X\to S$를 나타낼 때)
$(\Gamma(U,\mathcal{M}))^{\sim}$와 표준적으로 동형이다.

\begin{env}[1.4.4]
\label{II.1.4.4-ko}
준연접 $\mathcal{B}$-가군들의 범주에서
$\widetilde{\mathcal{M}}$는 $\mathcal{M}$에 관한 공변 가법 완전
함자이고, 귀납극한과 직합과 가환한다.
\end{env}

실제로 곧 $S$가 아핀인 경우로 환원되며, 그 경우 따름정리는
(\hyperref[I.1.3.5-ko]{I, 1.3.5}),
(\hyperref[I.1.3.9-ko]{I, 1.3.9})와
(\hyperref[I.1.3.11-ko]{I, 1.3.11})에서 따른다.

\begin{env}[1.4.5]
\label{II.1.4.5-ko}
(\hyperref[II.1.4.3-ko]{1.4.3})의 가정 아래에서
$\widetilde{\mathcal{M}}$가 유한 생성 $\mathcal{O}_X$-가군이기 위한
필요충분조건은 $\mathcal{M}$이 유한 생성 $\mathcal{B}$-가군인 것이다.
\end{env}

곧 $S=\operatorname{Spec}(A)$가 아핀 스킴인 경우로 환원된다. 이때
$\mathcal{B}=\widetilde{B}$이고, 여기서 $B$는 $A$-대수이다.\footnote{
\textbf{번역 주.} 인쇄 원문은 여기서 $B$가 유한 생성 $A$-대수라고
하고 (I, 9.6.2)를 인용하지만, 이는 명제의 가정에 없으므로
EGA2-CANON-DISP-0031에 따라 그 구절과 인용을 삭제했다.} 또한
$\mathcal{M}=\widetilde{M}$이고, 여기서 $M$은 $B$-가군이다
(\hyperref[I.1.3.13-ko]{I, 1.3.13}). \emph{준스킴 $X$ 위에서}
$\mathcal{O}_X$는 환 $B$에 대응하고
$\widetilde{\mathcal{M}}$는 $B$-가군 $M$에 대응한다. 따라서
$\widetilde{\mathcal{M}}$가 유한 생성이기 위한 필요충분조건은 $M$이
유한 생성인 것이다(\hyperref[I.1.3.13-ko]{I, 1.3.13}). 이로써
주장이 증명된다.

\begin{env}[1.4.6]
\label{II.1.4.6-ko}
$Y$를 $S$ 위에서 아핀인 준스킴, $X$와 $X'$을 $Y$ 위에서 아핀인
두 준스킴이라 하자(따라서 이들은 $S$ 위에서도 아핀이다
(\hyperref[II.1.3.5-ko]{1.3.5})). 다음과 같이 놓자.
$\mathcal{B}=\mathcal{A}(Y)$, $\mathcal{A}=\mathcal{A}(X)$,
$\mathcal{A}'=\mathcal{A}(X')$. 그러면 $X\times_Y X'$은 $Y$ 위에서
아핀이고(따라서 $S$ 위에서도 아핀이다),
$\mathcal{A}(X\times_Y X')$는
$\mathcal{A}\otimes_{\mathcal{B}}\mathcal{A}'$과 표준적으로
동일시된다.
\end{env}

실제로 (\hyperref[I.9.6.1-ko]{I, 9.6.1})에 따라
$\mathcal{A}\otimes_{\mathcal{B}}\mathcal{A}'$은 준연접
$\mathcal{B}$-대수이고, 따라서 준연접 $\mathcal{O}_S$-대수이기도
하다(\hyperref[I.9.6.1-ko]{I, 9.6.1}).
$Z$를 $\mathcal{A}\otimes_{\mathcal{B}}\mathcal{A}'$의 스펙트럼이라
하자(\hyperref[II.1.3.1-ko]{1.3.1}). 표준
$\mathcal{B}$-대수 준동형
$\mathcal{A}\to\mathcal{A}\otimes_{\mathcal{B}}\mathcal{A}'$과
$\mathcal{A}'\to\mathcal{A}\otimes_{\mathcal{B}}\mathcal{A}'$은
(\hyperref[II.1.2.7-ko]{1.2.7})에 따라 $Y$-사상 $p:Z\to X$와
$p':Z\to X'$에 대응한다. 삼중항 $(Z,p,p')$가 곱
$X\times_Y X'$임을 보이기 위해, $S$가 환 $C$의 아핀 스킴인
경우로 환원할 수 있다(\hyperref[I.3.2.6.4-ko]{I, 3.2.6.4}).
그러면 $Y$, $X$, $X'$은 아핀 스킴이고
(\hyperref[II.1.3.4-ko]{1.3.4}), 그 환 $B$, $A$, $A'$은
$C$-대수이며
$\mathcal{B}=\widetilde{B}$, $\mathcal{A}=\widetilde{A}$,
$\mathcal{A}'=\widetilde{A'}$이다.
(\hyperref[I.1.3.13-ko]{I, 1.3.13})에 따라
$\mathcal{A}\otimes_{\mathcal{B}}\mathcal{A}'$은
$\mathcal{O}_S$-대수 $\widetilde{A\otimes_B A'}$와 표준적으로
동일시된다. 따라서 환 $A(Z)$는 $A\otimes_B A'$과 동일시되고,
사상 $p$, $p'$은 표준 준동형
$A\to A\otimes_B A'$, $A'\to A\otimes_B A'$에 대응한다.
이제 명제는 (\hyperref[I.3.2.2-ko]{I, 3.2.2})에서 따른다.

\begin{env}[1.4.7]
\label{II.1.4.7-ko}
$\mathcal{F}$(각각 $\mathcal{F}'$)를 준연접
$\mathcal{O}_X$-가군(각각 $\mathcal{O}_{X'}$-가군)이라 하자. 그러면
$\mathcal{A}(\mathcal{F}\otimes_Y\mathcal{F}')$는
$\mathcal{A}(\mathcal{F})\otimes_{\mathcal{A}(Y)}
\mathcal{A}(\mathcal{F}')$과 표준적으로 동일시된다.
\end{env}

$\mathcal{F}\otimes_Y\mathcal{F}'$은 $X\times_Y X'$ 위에서
준연접이다(\hyperref[I.9.1.2-ko]{I, 9.1.2}).
$g:Y\to S$, $f:X\to Y$, $f':X'\to Y$를 구조 사상이라 하자.
그러면 구조 사상
\oldpage[II]{11}
$h:Z\to S$는 $g\circ f\circ p$ 및 $g\circ f'\circ p'$과 같다.
표준 준동형
\[
  \mathcal{A}(\mathcal{F})\otimes_{\mathcal{A}(Y)}\mathcal{A}(\mathcal{F}')
  \longrightarrow \mathcal{A}(\mathcal{F}\otimes_Y\mathcal{F}')
\]
을 다음과 같이 정의한다. 임의의 열린집합 $U\subset S$에 대하여
표준 준동형
\[
  \begin{aligned}
  \Gamma(f^{-1}(g^{-1}(U)),\mathcal{F})
  &\longrightarrow \Gamma(h^{-1}(U),p^*(\mathcal{F}))\quad\text{및}\quad
  \Gamma(f^{\prime-1}(g^{-1}(U)),\mathcal{F}')
  \longrightarrow \Gamma(h^{-1}(U),p^{\prime*}(\mathcal{F}'))
  \end{aligned}
\]
들이 있다(\hyperref[0.4.4.3-ko]{0, 4.4.3}). 따라서 표준 준동형
\[
  \begin{aligned}
  &\Gamma(f^{-1}(g^{-1}(U)),\mathcal{F})
    \otimes_{\Gamma(g^{-1}(U),\mathcal{O}_Y)}
    \Gamma(f^{\prime-1}(g^{-1}(U)),\mathcal{F}')\longrightarrow\\
  &\hspace{17em}
    \Gamma(h^{-1}(U),p^*(\mathcal{F}))
    \otimes_{\Gamma(h^{-1}(U),\mathcal{O}_Z)}
    \Gamma(h^{-1}(U),p^{\prime*}(\mathcal{F}'))
  \end{aligned}
\]
을 얻는다.

이것이 실제로 $\mathcal{A}(Z)$-가군의 동형임을 보이기 위해,
$S$(따라서 $X$, $X'$, $Y$, $X\times_Y X'$도)가 아핀 스킴인
경우로 환원할 수 있다. 이때
(\hyperref[II.1.4.6-ko]{1.4.6})의 표기를 사용하여
$\mathcal{F}=\widetilde{M}$, $\mathcal{F}'=\widetilde{M'}$라 쓸 수
있으며, $M$(각각 $M'$)은 $A$-가군(각각 $A'$-가군)이다. 그러면
$\mathcal{F}\otimes_Y\mathcal{F}'$은 $X\times_Y X'$ 위에서
$(A\otimes_B A')$-가군 $M\otimes_B M'$에 대응하는 층과
동일시된다(\hyperref[I.9.1.3-ko]{I, 9.1.3}). 따라서 따름정리는
$\mathcal{O}_S$-가군 $\widetilde{M\otimes_B M'}$와
$\widetilde{M}\otimes_{\widetilde{B}}\widetilde{M'}$의 표준적
동일시에서 따른다. 여기서 $M$, $M'$, $B$는 $C$-가군으로 본다
(\hyperref[I.1.3.13-ko]{I, 1.3.13}과
\hyperref[I.1.6.3-ko]{I, 1.6.3}).

특히 (\hyperref[II.1.4.7-ko]{1.4.7})을 $X=Y$이고
$\mathcal{F}'=\mathcal{O}_{X'}$인 경우에 적용하면,
$\mathcal{A}'$-가군 $\mathcal{A}(f^{\prime*}(\mathcal{F}))$가
$\mathcal{A}(\mathcal{F})\otimes_{\mathcal{B}}\mathcal{A}'$과
동일시됨을 알 수 있다.

\begin{env}[1.4.8]
\label{II.1.4.8-ko}
특히 $X=X'=Y$일 때($X$는 $S$ 위에서 아핀이다),
$\mathcal{F}$와 $\mathcal{G}$가 두 준연접
$\mathcal{O}_X$-가군이면
\[
  \mathcal{A}(\mathcal{F}\otimes_{\mathcal{O}_X}\mathcal{G})
  =\mathcal{A}(\mathcal{F})\otimes_{\mathcal{A}(X)}\mathcal{A}(\mathcal{G})
  \tag{1.4.8.1}\label{II.1.4.8.1-ko}
\]
이라는 표준적이고 함자적인 동형이 있다. 더욱이 $\mathcal{F}$가
유한 표시를 가지면
(\hyperref[I.1.6.3-ko]{I, 1.6.3}과
\hyperref[I.1.3.12-ko]{I, 1.3.12})에 따라
\[
  \mathcal{A}(\mathcal{H}om_X(\mathcal{F},\mathcal{G}))
  =\mathcal{H}om_{\mathcal{A}(X)}
    (\mathcal{A}(\mathcal{F}),\mathcal{A}(\mathcal{G}))
  \tag{1.4.8.2}\label{II.1.4.8.2-ko}
\]
이라는 표준적 동형이 있다.
\end{env}

\begin{env}[1.4.9]
\label{II.1.4.9-ko}
$X$, $X'$가 $S$ 위에서 아핀인 두 준스킴이면, 합
$X\sqcup X'$도 $S$ 위에서 아핀이다. 두 아핀 스킴의 합은 아핀
스킴이기 때문이다.
\end{env}

\begin{env}[1.4.10]
\label{II.1.4.10-ko}
$S$를 준스킴, $\mathcal{B}$를 준연접 $\mathcal{O}_S$-대수라 하고,
$X=\operatorname{Spec}(\mathcal{B})$라 하자. $\mathcal{J}$를
$\mathcal{B}$의 임의의 준연접 아이디얼층이라 하자. 그러면
$\widetilde{\mathcal{J}}$는 $\mathcal{O}_X$의 준연접 아이디얼층이다.
$Y$를 $X$에서 $\widetilde{\mathcal{J}}$로 정의되는 닫힌 부분준스킴이라
하면, 이는 $\operatorname{Spec}(\mathcal{B}/\mathcal{J})$와 표준적으로
동형이다.
\end{env}

실제로 (\hyperref[I.4.2.3-ko]{I, 4.2.3})에서 $Y$가 $S$ 위에서
아핀임이 곧바로 따른다. 따라서 (\hyperref[II.1.3.1-ko]{1.3.1})에 의해
$S$가 아핀인 경우로 환원되고, 이 경우 명제는
(\hyperref[I.4.1.2-ko]{I, 4.1.2})에서 곧바로 따른다.

(\hyperref[II.1.4.10-ko]{1.4.10})의 결과는 다음과 같이도 표현할 수
있다. $h:\mathcal{B}\to\mathcal{B}'$가 준연접 $\mathcal{O}_S$-대수들의
\emph{전사} 준동형이면,
${}^a h:\operatorname{Spec}(\mathcal{B}')\to
\operatorname{Spec}(\mathcal{B})$는 \emph{닫힌 몰입 사상}이다.

\oldpage[II]{12}
\begin{env}[1.4.11]
\label{II.1.4.11-ko}
$S$를 준스킴, $\mathcal{B}$를 준연접 $\mathcal{O}_S$-대수라 하고,
$X=\operatorname{Spec}(\mathcal{B})$라 하자. $\mathcal{K}$를
$\mathcal{O}_S$의 임의의 준연접 아이디얼층이라 하자. 그러면
($f$가 구조 사상 $X\to S$를 나타낼 때)
$f^*(\mathcal{K})\mathcal{O}_X=\widetilde{\mathcal{K}\mathcal{B}}$가
표준 동형까지 성립한다.
\end{env}

이 문제는 $S$ 위에서 국소적이므로 $S=\operatorname{Spec}(A)$가 아핀인
경우로 환원된다. 이 경우 명제는
(\hyperref[I.1.6.9-ko]{I, 1.6.9})에 다름 아니다.

\subsection{밑준스킴의 변경}
\label{II.1.5-ko}

\begin{env}[1.5.1]
\label{II.1.5.1-ko}
$X$를 $S$ 위에서 아핀인 준스킴이라 하자. 밑준스킴의 임의의 확장
$g:S'\to S$에 대하여 $X'=X_{(S')}=X\times_S S'$은 $S'$ 위에서
아핀이다.
\end{env}

$f'$을 사영 $X'\to S'$이라 하자. $g(U')$가 $S$의 어떤 아핀
열린집합 $U$에 포함되는 $S'$의 모든 아핀 열린집합 $U'$에 대하여
$f^{\prime-1}(U')$가 아핀 열린집합임을 보이면 충분하다
(\hyperref[II.1.2.1-ko]{1.2.1}). 따라서 $S$와 $S'$가 아핀인 경우로
한정할 수 있고, 이 경우 $X'$이 아핀 스킴임을 보이면 충분하다
(\hyperref[II.1.3.4-ko]{1.3.4}). 그런데 이때
(\hyperref[II.1.3.4-ko]{1.3.4}) $X$는 아핀 스킴이다. $S$, $S'$, $X$의
환을 각각 $A$, $A'$, $B$라 하면, $X'$은 환 $A'\otimes_A B$의 아핀
스킴임을 안다(\hyperref[I.3.2.2-ko]{I, 3.2.2}).

\begin{env}[1.5.2]
\label{II.1.5.2-ko}
(\hyperref[II.1.5.1-ko]{1.5.1})의 가정 아래에서 $f:X\to S$를 구조
사상, $f':X'\to S'$과 $g':X'\to X$를 사영이라 하자. 그러면 다음
도식은 가환한다.
\[
  \xymatrix{
    X\ar[d] & X'\ar[l]_{g'}\ar[d]^{f'}\\
    S & S'\ar[l]_g
  }
\]
모든 준연접 $\mathcal{O}_X$-가군 $\mathcal{F}$에 대하여 다음 표준
$\mathcal{O}_{S'}$-가군 동형이 존재한다.
\[
  u:g^*(f_*(\mathcal{F}))\xrightarrow{\sim}
  f'_*(g^{\prime*}(\mathcal{F})).
  \tag{1.5.2.1}\label{II.1.5.2.1-ko}
\]
특히 $\mathcal{A}(X')$에서 $g^*(\mathcal{A}(X))$로 가는 표준 동형이
존재한다.
\end{env}

$u$를 정의하려면 준동형
\[
  v:f_*(\mathcal{F})\longrightarrow
  g_*(f'_*(g^{\prime*}(\mathcal{F})))
  =f_*(g'_*(g^{\prime*}(\mathcal{F})))
\]
를 정의하고 $u=v^\sharp$로 두면 충분하다
(\hyperref[0.4.4.3-ko]{0, 4.4.3}). $v=f_*(\rho)$로 두겠다. 여기서
$\rho$는 표준 준동형
$\mathcal{F}\to g'_*(g^{\prime*}(\mathcal{F}))$이다
(\hyperref[0.4.4.3-ko]{0, 4.4.3}). $u$가 동형임을 보이기 위해서는
$S$와 $S'$, 따라서 $X$와 $X'$도 아핀인 경우로 한정할 수 있다.
(\hyperref[II.1.5.1-ko]{1.5.1})의 표기를 사용하면 이때
$\mathcal{F}=\widetilde{M}$이고, 여기서 $M$은 $B$-가군이다. 그러면
$g^*(f_*(\mathcal{F}))$와 $f'_*(g^{\prime*}(\mathcal{F}))$는 둘 다
$A'$-가군 $A'\otimes_A M$에 대응하는 $\mathcal{O}_{S'}$-가군과 같다
(여기서 $M$은 $A$-가군으로 본다). 또한 $u$는 항등사상에 대응하는
준동형임을 곧바로 알 수 있다
(\hyperref[I.1.6.3-ko]{I, 1.6.3},
\hyperref[I.1.6.5-ko]{1.6.5}와 \hyperref[I.1.6.7-ko]{1.6.7}).

\begin{env}[1.5.3]
\label{II.1.5.3-ko}
$X$가 더 이상 $S$ 위에서 아핀이라고 가정되지 않을 때에는
(\hyperref[II.1.5.2-ko]{1.5.2})가 여전히 성립한다고 생각해서는 안
된다. 이는 심지어 $S'=\operatorname{Spec}(k(s))$ ($s\in S$)이고
$S'\to S$가 표준 사상인 경우에도 마찬가지이다
(\hyperref[I.2.4.5-ko]{I, 2.4.5}) --- 이 경우 $X'$은 바로
\emph{섬유} $f^{-1}(s)$이다(\hyperref[I.3.6.2-ko]{I, 3.6.2}). 다시
말해 $X$가 $S$ 위에서 아핀이 아닐 때에는
\oldpage[II]{13}
``준연접층의 직상'' 연산은 ``섬유 취하기'' 연산과 교환하지 않는다.
그러나 제III장(\agref{III.4.2.4-ko}{III, 4.2.4})에서 $f$가 고유
사상이고(\agref{II.5.4-ko}{5.4}) $S$가 뇌터일 때 $X$ 위의
\emph{연접층}에 대하여 성립하는, 이 방향의 ``점근적'' 성격의 결과를
보게 될 것이다.
\end{env}

\begin{env}[1.5.4]
\label{II.1.5.4-ko}
$S$ 위에서 아핀인 모든 준스킴 $X$와 모든 $s\in S$에 대하여, 섬유
$f^{-1}(s)$는 아핀 스킴이다(여기서 $f$는 구조 사상 $X\to S$를
나타낸다).
\end{env}

(\hyperref[II.1.5.1-ko]{1.5.1})을 $S'=\operatorname{Spec}(k(s))$에
적용하고 (\hyperref[II.1.3.4-ko]{1.3.4})을 사용하면 충분하다.

\begin{corollary}[1.5.5]
\label{II.1.5.5-ko}
$X$를 $S$-준스킴, $S'$을 $S$ 위에서 아핀인 준스킴이라 하자. 그러면
$X'=X_{(S')}$은 $X$ 위에서 아핀인 준스킴이다. 또한 $f:X\to S$가
구조 사상이면 다음 표준 $\mathcal{O}_X$-대수 동형이 존재한다.
$\mathcal{A}(X')\xrightarrow{\sim}f^*(\mathcal{A}(S'))$.
그리고 모든 준연접 $\mathcal{A}(S')$-가군 $\mathcal{M}$에 대하여 다음
표준 디-동형이 존재한다.
$f^*(\mathcal{M})\xrightarrow{\sim}
\mathcal{A}(f^{\prime*}(\widetilde{\mathcal{M}}))$.
여기서 $f'=f_{(S')}$은 구조 사상 $X'\to S'$을 나타낸다.
\end{corollary}

(\hyperref[II.1.5.1-ko]{1.5.1})과
(\hyperref[II.1.5.2-ko]{1.5.2})에서 $X$와 $S'$의 역할을 서로 바꾸면
충분하다.

\begin{env}[1.5.6]
\label{II.1.5.6-ko}
이제 $S$, $S'$을 두 준스킴, $q:S'\to S$를 사상,
$\mathcal{B}$(각각 $\mathcal{B}'$)를 준연접
$\mathcal{O}_S$-대수(각각 준연접 $\mathcal{O}_{S'}$-대수),
$u:\mathcal{B}\to\mathcal{B}'$를 $q$-사상, 즉 $\mathcal{O}_S$-대수
준동형 $\mathcal{B}\to q_*(\mathcal{B}')$이라 하자.
$X=\operatorname{Spec}(\mathcal{B})$,
$X'=\operatorname{Spec}(\mathcal{B}')$라 하면, 다음 사상을 표준적으로
얻는다.
\[
  v=\operatorname{Spec}(u):X'\to X
\]
이때 도식
\[
  \xymatrix{
    X'\ar[r]^v\ar[d] & X\ar[d]\\
    S'\ar[r]_q & S
  }
  \tag{1.5.6.1}\label{II.1.5.6.1-ko}
\]
은 가환한다(수직 화살표들은 구조 사상이다). 실제로 $u$를 주는 것은
준연접 $\mathcal{O}_{S'}$-대수들의 준동형
$u^\sharp:q^*(\mathcal{B})\to\mathcal{B}'$를 주는 것과 동치이다
(\hyperref[0.4.4.3-ko]{0, 4.4.3}). 따라서 이는
$\mathcal{A}(w)=u^\sharp$인 표준 $S'$-사상
\[
  w:\operatorname{Spec}(\mathcal{B}')\to
  \operatorname{Spec}(q^*(\mathcal{B}))
\]
을 정의한다(\hyperref[II.1.2.7-ko]{1.2.7}). 한편
(\hyperref[II.1.5.2-ko]{1.5.2})에 따라
$\operatorname{Spec}(q^*(\mathcal{B}))$는 $X\times_S S'$과 표준적으로
동일시된다. 사상 $v$는 $w$에 이어 첫째 사영을 취한 합성
$X'\xrightarrow{w}X\times_S S'\xrightarrow{p_1}X$이고,
(\hyperref[II.1.5.6.1-ko]{1.5.6.1})의 가환성은 정의에서 따른다.
$U$(각각 $U'$)를 $q(U')\subset U$인 $S$(각각 $S'$)의 아핀
열린집합이라 하고,
$A=\Gamma(U,\mathcal{O}_S)$,
$A'=\Gamma(U',\mathcal{O}_{S'})$를 그 환들,
$B=\Gamma(U,\mathcal{B})$,
$B'=\Gamma(U',\mathcal{B}')$라 하자. $u$의 제한인
$(q|U')$-사상 $\mathcal{B}|U\to\mathcal{B}'|U'$은 대수의
디-준동형 $B\to B'$에 대응한다. $V$와 $V'$을 각각 구조 사상
$X\to S$와 $X'\to S'$에 의한 $U$와 $U'$의 역상이라 하면, $v$의 제한인 사상
$V'\to V$는 (\hyperref[I.1.7.3-ko]{I, 1.7.3})에 따라 앞의
디-준동형에 대응한다.
\end{env}

\begin{env}[1.5.7]
\label{II.1.5.7-ko}
(\hyperref[II.1.5.6-ko]{1.5.6})과 같은 가정 아래에서 $\mathcal{M}$을
준연접 $\mathcal{B}$-가군이라 하자. 그러면 다음 표준
$\mathcal{O}_{X'}$-가군 동형이 존재한다.
\[
  v^*(\widetilde{\mathcal{M}})\xrightarrow{\sim}
  \bigl(q^*(\mathcal{M})\otimes_{q^*(\mathcal{B})}\mathcal{B}'\bigr)^{\sim}
  \tag{1.5.7.1}\label{II.1.5.7.1-ko}
\]

\oldpage[II]{14}
실제로 (\hyperref[II.1.5.2.1-ko]{1.5.2.1})의 표준 동형은
$p_1^*(\widetilde{\mathcal{M}})$에서
$\operatorname{Spec}(q^*(\mathcal{B}))$ 위의
$q^*(\mathcal{B})$-가군 $q^*(\mathcal{M})$에 대응하는 층으로 가는
표준 동형을 준다. 이제 (\hyperref[II.1.4.7-ko]{1.4.7})을 적용하면
충분하다.
\end{env}

\subsection{아핀 사상}
\label{II.1.6-ko}

\begin{env}[1.6.1]
\label{II.1.6.1-ko}
준스킴의 사상 $f:X\to Y$를 구조 사상으로 보았을 때 $X$가 $Y$ 위에서
아핀인 준스킴이면, $f$를 \emph{아핀} 사상이라 한다. 한 준스킴 위에서
아핀인 준스킴들의 성질은 이 용어로 다음과 같이 표현된다.
\end{env}

\begin{proposition}[1.6.2]
\label{II.1.6.2-ko}
\begin{enumerate}
  \item[\rm(i)] 닫힌 몰입 사상은 아핀 사상이다.
  \item[\rm(ii)] 두 아핀 사상의 합성은 아핀 사상이다.
  \item[\rm(iii)] $f:X\to Y$가 아핀 $S$-사상이면, 밑준스킴의 모든 확장
    $S'\to S$에 대하여 $f_{(S')}:X_{(S')}\to Y_{(S')}$도 아핀 사상이다.
  \item[\rm(iv)] $f:X\to Y$와 $f':X'\to Y'$이 두 아핀 $S$-사상이면,
    \[
      f\times_S f':X\times_S X'\to Y\times_S Y'
    \]
    도 아핀 사상이다.
  \item[\rm(v)] $f:X\to Y$와 $g:Y\to Z$가 두 사상이고 $g\circ f$가
    아핀 사상이며 $g$가 분리 사상이면, $f$는 아핀 사상이다.
  \item[\rm(vi)] $f$가 아핀 사상이면 $f_{\mathrm{red}}$도 아핀 사상이다.
\end{enumerate}
\end{proposition}

(\hyperref[I.5.5.12-ko]{I, 5.5.12})에 따라 (i), (ii), (iii)만 증명하면
충분하다. 그런데 (i)는 (\hyperref[II.1.2.2-ko]{1.2.2})의 예이고,
(ii)는 (\hyperref[II.1.3.5-ko]{1.3.5})이며, 마지막으로 (iii)은
$X_{(S')}$이 곱 $X\times_Y Y_{(S')}$과 동일시되므로
(\hyperref[II.1.5.1-ko]{1.5.1})에서 따른다
(\hyperref[I.3.3.11-ko]{I, 3.3.11}).

\begin{corollary}[1.6.3]
\label{II.1.6.3-ko}
$X$가 아핀 스킴이고 $Y$가 스킴이면, 모든 사상 $X\to Y$는 아핀
사상이다.
\end{corollary}

\begin{proposition}[1.6.4]
\label{II.1.6.4-ko}
$Y$를 국소 뇌터 준스킴, $f:X\to Y$를 유한형 사상이라 하자. $f$가
아핀 사상이기 위한 필요충분조건은 $f_{\mathrm{red}}$가 아핀 사상인
것이다.
\end{proposition}

(\hyperref[II.1.6.2-ko]{1.6.2, (vi)})에 따라 조건의 충분성만 증명하면
된다. $Y$가 아핀이고 뇌터일 때 $X$가 아핀임을 보이면 충분하다. 이때
$Y_{\mathrm{red}}$는 아핀이므로, 가정에 따라 $X_{\mathrm{red}}$도
아핀이다. 한편 $X$는 뇌터이므로, 결론은
(\hyperref[I.6.1.7-ko]{I, 6.1.7})에서 따른다.

\subsection{가군의 층에 대응하는 벡터 다발}
\label{II.1.7-ko}

\begin{env}[1.7.1]
\label{II.1.7.1-ko}
$A$를 환, $E$를 $A$-가군이라 하자. $E$의 \emph{대칭 대수}라 하고
$\mathbf{S}(E)$(또는 $\mathbf{S}_A(E)$)로 나타내는 것은 텐서 대수
$\mathbf{T}(E)$를 $x,y\in E$에 대한 원소
$x\otimes y-y\otimes x$들이 생성하는 양쪽 아이디얼로 나눈 몫대수임을
상기하자. 대수 $\mathbf{S}(E)$는 다음 보편 성질로 특징지어진다.
$\sigma$를 표준 사상 $E\to\mathbf{S}(E)$, 즉
$E\to\mathbf{T}(E)$와 표준 사상
$\mathbf{T}(E)\to\mathbf{S}(E)$의 합성이라 하자. $B$가 가환
$A$-대수이면, 모든 $A$-선형 사상 $E\to B$는 유일하게
$E\xrightarrow{\sigma}\mathbf{S}(E)\xrightarrow{g}B$로 인수분해된다. 여기서
$g$는 $A$-대수 준동형이다. 이 특징지음에서 두 $A$-가군 $E$, $F$에
대하여
\[
  \mathbf{S}(E\oplus F)=\mathbf{S}(E)\otimes\mathbf{S}(F)
\]

\oldpage[II]{15}
가 표준 동형을 제외하면 성립함이 곧바로 따른다. 또한
$E\mapsto\mathbf{S}(E)$는 $A$-가군의 범주에서 가환 $A$-대수의 범주로
가는 공변 함자이다. 마지막으로 앞의 특징지음에 따라
$E=\varinjlim E_\lambda$이면
$\mathbf{S}(E)=\varinjlim\mathbf{S}(E_\lambda)$가 표준 동형을 제외하면
성립한다. 말의 남용으로 $x_i\in E$일 때의 곱
$\sigma(x_1)\sigma(x_2)\cdots\sigma(x_n)$을 혼동의 우려가 없으면 흔히
$x_1x_2\cdots x_n$으로 나타낸다. 대수 $\mathbf{S}(E)$는
\emph{등급} 대수이며, $\mathbf{S}_n(E)$는 $E$의 $n$개 원소의 곱들의
$A$-선형 결합 전체이다($n\geq 0$). 대수 $\mathbf{S}(A)$는 한 부정원의
다항식 대수 $A[T]$와 표준적으로 동형이고, 대수 $\mathbf{S}(A^n)$은
$n$개 부정원의 다항식 대수 $A[T_1,\ldots,T_n]$과 표준적으로 동형이다.
\end{env}

\begin{env}[1.7.2]
\label{II.1.7.2-ko}
$\varphi:A\to B$를 환 준동형이라 하자. $F$가 $B$-가군이면, 표준 사상
$F\to\mathbf{S}(F)$는 표준 사상
$F_{[\varphi]}\to\mathbf{S}(F)_{[\varphi]}$를 주고, 따라서 이 사상은
$F_{[\varphi]}\to\mathbf{S}(F_{[\varphi]})\to
\mathbf{S}(F)_{[\varphi]}$로 인수분해된다. 표준 $A$-대수 준동형
$\mathbf{S}(F_{[\varphi]})\to\mathbf{S}(F)_{[\varphi]}$는 전사이지만
반드시 전단사인 것은 아니다. $E$가 $A$-가군이면, 모든 디-준동형
$E\to F$, 즉 모든 $A$-가군 준동형 $E\to F_{[\varphi]}$는 따라서
표준적으로 $A$-대수 준동형
$\mathbf{S}(E)\to\mathbf{S}(F_{[\varphi]})\to
\mathbf{S}(F)_{[\varphi]}$를 준다.
이는 곧 대수의 디-준동형 $\mathbf{S}(E)\to\mathbf{S}(F)$이다.

같은 표기 아래에서 모든 $A$-가군 $E$에 대하여
$\mathbf{S}(E\otimes_A B)$는 대수
$\mathbf{S}(E)\otimes_A B$와 표준적으로 동일시된다. 이는
$\mathbf{S}(E)$의 보편 성질에서 곧바로 따른다
(\hyperref[II.1.7.1-ko]{1.7.1}).
\end{env}

\begin{env}[1.7.3]
\label{II.1.7.3-ko}
$R$을 환 $A$의 곱셈적 부분집합이라 하자.
(\hyperref[II.1.7.2-ko]{1.7.2})를 환 $B=R^{-1}A$에 적용하고
$R^{-1}E=E\otimes_A R^{-1}A$임을 상기하면,
$\mathbf{S}(R^{-1}E)=R^{-1}\mathbf{S}(E)$가 표준 동형을 제외하면
성립함을 알 수 있다. 또한 $R'\supset R$이 $A$의 또 하나의 곱셈적
부분집합이면 도식
\[
  \xymatrix{
    R^{-1}E\ar[r]\ar[d] & {R'}^{-1}E\ar[d]\\
    \mathbf{S}(R^{-1}E)\ar[r] & \mathbf{S}({R'}^{-1}E)
  }
\]
은 가환한다.
\end{env}

\begin{env}[1.7.4]
\label{II.1.7.4-ko}
이제 $(S,\mathcal{A})$를 환 달린 공간, $\mathcal{E}$를 $S$ 위의
$\mathcal{A}$-가군이라 하자. 모든 열린집합 $U\subset S$에
$\Gamma(U,\mathcal{A})$-가군
$\mathbf{S}(\Gamma(U,\mathcal{E}))$을 대응시키면,
$\mathbf{S}(E)$의 함자성에 따라 대수의 준층이 정의된다
(\hyperref[II.1.7.2-ko]{1.7.2}). 이 준층의 층화를
$\mathbf{S}(\mathcal{E})$ 또는
$\mathbf{S}_{\mathcal{A}}(\mathcal{E})$로 나타내고,
$\mathcal{A}$-가군 $\mathcal{E}$의 \emph{대칭 대수}라 한다.
(\hyperref[II.1.7.1-ko]{1.7.1})에서
$\mathbf{S}(\mathcal{E})$가 보편 문제의 해임이 곧바로 따른다.
$\mathcal{B}$가 $\mathcal{A}$-대수일 때 모든 $\mathcal{A}$-가군
준동형 $\mathcal{E}\to\mathcal{B}$는 유일하게
$\mathcal{E}\to\mathbf{S}(\mathcal{E})\to\mathcal{B}$로 인수분해되며,
둘째 화살표는 $\mathcal{A}$-대수 준동형이다. 따라서
$\mathcal{A}$-가군 준동형 $\mathcal{E}\to\mathcal{B}$들과
$\mathcal{A}$-대수 준동형
$\mathbf{S}(\mathcal{E})\to\mathcal{B}$들 사이에는 전단사 대응이
있다. 특히 모든 $\mathcal{A}$-가군 준동형
$u:\mathcal{E}\to\mathcal{F}$는 $\mathcal{A}$-대수 준동형
$\mathbf{S}(u):\mathbf{S}(\mathcal{E})\to
\mathbf{S}(\mathcal{F})$를 정의하며, 따라서
$\mathcal{E}\mapsto\mathbf{S}(\mathcal{E})$는 공변 함자이다.

\oldpage[II]{16}
(\hyperref[II.1.7.2-ko]{1.7.2})와 $\mathbf{S}$가 귀납극한과 가환한다는
사실에 따라, 모든 점 $s\in S$에 대하여
$(\mathbf{S}(\mathcal{E}))_s=\mathbf{S}(\mathcal{E}_s)$이다.
$\mathcal{E}$, $\mathcal{F}$가 두 $\mathcal{A}$-가군이면
$\mathbf{S}(\mathcal{E}\oplus\mathcal{F})$는
$\mathbf{S}(\mathcal{E})\otimes_{\mathcal{A}}
\mathbf{S}(\mathcal{F})$와 표준적으로 동일시된다. 이는 대응하는
준층들에서 확인된다.

또한 $\mathbf{S}(\mathcal{E})$는 $\mathbf{S}_n(\mathcal{E})$들의 무한
직합인 등급 $\mathcal{A}$-대수이다. 여기서 $\mathcal{A}$-가군
$\mathbf{S}_n(\mathcal{E})$는 준층
$U\mapsto\mathbf{S}_n(\Gamma(U,\mathcal{E}))$의 층화이다. 특히
$\mathcal{E}=\mathcal{A}$로 두면
$\mathbf{S}_{\mathcal{A}}(\mathcal{A})$는
$\mathcal{A}[T]=\mathcal{A}\otimes_{\mathbf{Z}}\mathbf{Z}[T]$와
동일시된다($T$는 부정원이고, $\mathbf{Z}$는 상수층으로 본다).
\end{env}

\begin{env}[1.7.5]
\label{II.1.7.5-ko}
$(T,\mathcal{B})$를 또 하나의 환 달린 공간,
$f:(S,\mathcal{A})\to(T,\mathcal{B})$를 환 달린 공간의 사상이라 하자.
$\mathcal{F}$가 $\mathcal{B}$-가군이면
$\mathbf{S}(f^*(\mathcal{F}))$는
$f^*(\mathbf{S}(\mathcal{F}))$와 표준적으로 동일시된다. 실제로
$f=(\psi,\theta)$이면 정의상
(\hyperref[0.4.3.1-ko]{0, 4.3.1})
\[
  \mathbf{S}(f^*(\mathcal{F}))
  =\mathbf{S}(\psi^*(\mathcal{F})
    \otimes_{\psi^*(\mathcal{B})}\mathcal{A})
  =\mathbf{S}(\psi^*(\mathcal{F}))
    \otimes_{\psi^*(\mathcal{B})}\mathcal{A}
\]
이다(\hyperref[II.1.7.2-ko]{1.7.2}). $S$의 임의의 열린집합 $U$와
$U$ 위의 $\mathbf{S}(\psi^*(\mathcal{F}))$의 임의의 단면 $h$에
대하여, 각 점 $s\in U$의 어떤 근방 $V$에서 $h$는
$\mathbf{S}(\Gamma(V,\psi^*(\mathcal{F})))$의 한 원소와 일치한다.
$\psi^*(\mathcal{F})$의 정의
(\hyperref[0.3.7.1-ko]{0, 3.7.1})를 참조하고, 가군 $E$에 대한
$\mathbf{S}(E)$의 모든 원소가 $E$의 원소들의 곱을 유한 개만 사용한
선형 결합임을 고려하면, $T$에서 $\psi(s)$의 어떤 근방 $W$,
$W$ 위의 $\mathbf{S}(\mathcal{F})$의 어떤 단면 $h'$, 그리고 $s$의
어떤 근방 $V'\subset V\cap\psi^{-1}(W)$가 존재하여 $V'$ 위에서
$h$는 $t\mapsto h'(\psi(t))$와 일치함을 알 수 있다. 따라서 주장이
따른다.
\end{env}

\begin{proposition}[1.7.6]
\label{II.1.7.6-ko}
$A$를 환, $S=\operatorname{Spec}(A)$를 그 소 스펙트럼,
$\mathcal{E}=\widetilde{M}$을 $A$-가군 $M$에 대응하는
$\mathcal{O}_S$-가군이라 하자. 그러면 $\mathcal{O}_S$-대수
$\mathbf{S}(\mathcal{E})$는 $A$-대수 $\mathbf{S}(M)$에 대응하는
$\mathcal{O}_S$-대수이다.
\end{proposition}

\begin{proof}
실제로 모든 $f\in A$에 대하여
$\mathbf{S}(M_f)=(\mathbf{S}(M))_f$이다
(\hyperref[II.1.7.3-ko]{1.7.3}). 따라서 명제는 정의들
(\hyperref[I.1.3.4-ko]{I, 1.3.4})에서 따른다.
\end{proof}

\begin{corollary}[1.7.7]
\label{II.1.7.7-ko}
$S$가 준스킴이고 $\mathcal{E}$가 준연접 $\mathcal{O}_S$-가군이면,
$\mathcal{O}_S$-대수 $\mathbf{S}(\mathcal{E})$는 준연접이다. 또한
$\mathcal{E}$가 유한 생성 $\mathcal{O}_S$-가군이면, 각
$\mathcal{O}_S$-가군 $\mathbf{S}_n(\mathcal{E})$도 유한 생성이다.
\end{corollary}

\begin{proof}
첫째 주장은 (\hyperref[II.1.7.6-ko]{1.7.6})과
(\hyperref[I.1.4.1-ko]{I, 1.4.1})에서 곧바로 따른다. 둘째 주장은
$E$가 유한 생성 $A$-가군이면 $\mathbf{S}_n(E)$도 유한 생성
$A$-가군이라는 사실에서 따른다. 이제
(\hyperref[I.1.3.13-ko]{I, 1.3.13})을 적용하면 된다.
\end{proof}

\begin{definition}[1.7.8]
\label{II.1.7.8-ko}
$\mathcal{E}$를 준연접 $\mathcal{O}_S$-가군이라 하자. 준연접
$\mathcal{O}_S$-대수 $\mathbf{S}(\mathcal{E})$의 스펙트럼
(\hyperref[II.1.3.1-ko]{1.3.1})을 $\mathbf{V}(\mathcal{E})$로 나타내고,
이를 $\mathcal{E}$로 정의되는 $S$ 위의 \emph{벡터 다발}이라 한다.
\end{definition}

(\hyperref[II.1.2.7-ko]{1.2.7})에 따라 임의의 $S$-준스킴 $X$에 대하여,
$S$-사상 $X\to\mathbf{V}(\mathcal{E})$들과 $\mathcal{O}_S$-대수 준동형
$\mathbf{S}(\mathcal{E})\to\mathcal{A}(X)$들 사이에는 표준적인 전단사
대응이 있고, 따라서 이 $S$-사상들과 $\mathcal{O}_S$-가군 준동형
$\mathcal{E}\to\mathcal{A}(X)=f_*(\mathcal{O}_X)$들 사이에도 그러한
대응이 있다(여기서 $f$는 구조 사상 $X\to S$이다). 특히 다음을 얻는다.

\begin{env}[1.7.9]
\label{II.1.7.9-ko}
$X$로 $S$가 열린집합 $U\subset S$ 위에 유도하는 부분준스킴을 잡자.
그러면 $S$-사상 $U\to\mathbf{V}(\mathcal{E})$들은 열린집합
$p^{-1}(U)$ 위에 $\mathbf{V}(\mathcal{E})$가 유도하는 $U$-준스킴의
$U$-단면들(\hyperref[I.2.5.5-ko]{I, 2.5.5})에 지나지 않는다(여기서
$p$는 구조 사상 $\mathbf{V}(\mathcal{E})\to S$이다). 방금 본 바에
따르면 이 $U$-단면들은 $\mathcal{O}_S$-가군 준동형
$\mathcal{E}\to j_*(\mathcal{O}_S|U)$들과 전단사로 대응한다(여기서
$j$는 표준 포함 사상 $U\to S$이다). 또는, 이와 동치로

\oldpage[II]{17}
(\hyperref[0.4.4.3-ko]{0, 4.4.3}), 이들은 $(\mathcal{O}_S|U)$-가군
준동형
$j^*(\mathcal{E})=\mathcal{E}|U\to\mathcal{O}_S|U$들과 대응한다. 또한
$S$-사상 $U\to\mathbf{V}(\mathcal{E})$을 열린집합 $U'\subset U$로
제한한 것에는 대응하는 준동형
$\mathcal{E}|U\to\mathcal{O}_S|U$를 $U'$로 제한한 것이 대응함은
자명하다. 따라서 $\mathbf{V}(\mathcal{E})$의 \emph{$S$-단면의 싹으로
이루어진 층}은 $\mathcal{E}$의 \emph{쌍대 $\check{\mathcal{E}}$}와
표준적으로 동일시된다.

특히 $X=U=S$로 두면 영 준동형
$\mathcal{E}\to\mathcal{O}_S$에는 $\mathbf{V}(\mathcal{E})$의 표준
$S$-단면 하나가 대응하며, 이를 \emph{영 $S$-단면}이라 한다
(8.3.3 참조).
\end{env}

\begin{env}[1.7.10]
\label{II.1.7.10-ko}
이제 $X$로 체 $K$의 스펙트럼 $\{\xi\}$를 잡자. 그러면 구조 사상
$f:X\to S$는 $s=f(\xi)$일 때 체의 단사 준동형
$k(s)\to K$에 대응한다(\hyperref[I.2.4.6-ko]{I, 2.4.6}). 따라서
$S$-사상 $\{\xi\}\to\mathbf{V}(\mathcal{E})$들은 $k(s)$의 확대체
$K$를 값체로 하는 $\mathbf{V}(\mathcal{E})$의 \emph{기하적 점들}
(\hyperref[I.3.4.5-ko]{I, 3.4.5})에 지나지 않으며, 이 기하적 점들은
$p^{-1}(s)$의 점들에 놓인다. 이 점들의 집합은 점 $s$ 위의
$\mathbf{V}(\mathcal{E})$의 \emph{$K$-유리 기하적 섬유}라고도 할 수
있으며, (\hyperref[II.1.7.8-ko]{1.7.8})에 따라 $\mathcal{O}_S$-가군
준동형 $\mathcal{E}\to f_*(\mathcal{O}_X)$들의 집합과 동일시된다.
또는, 이와 동치로 (\hyperref[0.4.4.3-ko]{0, 4.4.3}),
$\mathcal{O}_X$-가군 준동형
$f^*(\mathcal{E})\to\mathcal{O}_X=K$들의 집합과 동일시된다. 한편 정의상
(\hyperref[0.4.3.1-ko]{0, 4.3.1})
$f^*(\mathcal{E})=\mathcal{E}_s\otimes_{\mathcal{O}_s}K
=\mathcal{E}^s\otimes_{k(s)}K$이며, 여기서
$\mathcal{E}^s=\mathcal{E}_s/\mathfrak{m}_s\mathcal{E}_s$로 둔다.
따라서 점 $s$ 위의 $\mathbf{V}(\mathcal{E})$의 $K$-유리 기하적
섬유는 \emph{$K$-벡터 공간
$\mathcal{E}^s\otimes_{k(s)}K$}의 \emph{쌍대}와 동일시된다.
$\mathcal{E}^s$ 또는 $K$가 $k(s)$ 위에서 유한 차원이면 이 쌍대는
$(\mathcal{E}^s)^\vee\otimes_{k(s)}K$와도 동일시된다. 여기서
$(\mathcal{E}^s)^\vee$는 $k(s)$-벡터 공간 $\mathcal{E}^s$의 쌍대를
나타낸다.
\end{env}

\begin{proposition}[1.7.11]
\label{II.1.7.11-ko}
\begin{enumerate}
  \item[\rm(i)] $\mathbf{V}(\mathcal{E})$는 준연접
    $\mathcal{O}_S$-가군의 범주에서 아핀 $S$-스킴의 범주로 가는,
    $\mathcal{E}$에 관한 반변 함자이다.
  \item[\rm(ii)] $\mathcal{E}$가 유한 생성 $\mathcal{O}_S$-가군이면,
    $\mathbf{V}(\mathcal{E})$는 $S$ 위에서 유한형이다.
  \item[\rm(iii)] $\mathcal{E}$와 $\mathcal{F}$가 두 준연접
    $\mathcal{O}_S$-가군이면,
    $\mathbf{V}(\mathcal{E}\oplus\mathcal{F})$는
    $\mathbf{V}(\mathcal{E})\times_S\mathbf{V}(\mathcal{F})$와
    표준적으로 동일시된다.
  \item[\rm(iv)] $g:S'\to S$를 사상이라 하자. 임의의 준연접
    $\mathcal{O}_S$-가군 $\mathcal{E}$에 대하여,
    $\mathbf{V}(g^*(\mathcal{E}))$는
    $\mathbf{V}(\mathcal{E})_{(S')}
    =\mathbf{V}(\mathcal{E})\times_S S'$와 표준적으로 동일시된다.
  \item[\rm(v)] 준연접 $\mathcal{O}_S$-가군의 전사 준동형
    $\mathcal{E}\to\mathcal{F}$에는 닫힌 몰입 사상
    $\mathbf{V}(\mathcal{F})\to\mathbf{V}(\mathcal{E})$가 대응한다.
\end{enumerate}
\end{proposition}

\begin{proof}
임의의 $\mathcal{O}_S$-가군 준동형 $\mathcal{E}\to\mathcal{F}$가
$\mathcal{O}_S$-대수 준동형
$\mathbf{S}(\mathcal{E})\to\mathbf{S}(\mathcal{F})$를 표준적으로
정한다는 점을 고려하면, (i)는 (\hyperref[II.1.2.7-ko]{1.2.7})의
직접적인 귀결이다. (ii)는 정의들
(\hyperref[I.6.3.1-ko]{I, 6.3.1})과, $E$가 유한 생성 $A$-가군이면
$\mathbf{S}(E)$가 유한 생성 $A$-대수라는 사실에서 곧바로 따른다.
(iii)를 증명하려면 표준 동형
$\mathbf{S}(\mathcal{E}\oplus\mathcal{F})
\simeq\mathbf{S}(\mathcal{E})
\otimes_{\mathcal{O}_S}\mathbf{S}(\mathcal{F})$
(\hyperref[II.1.7.4-ko]{1.7.4})에서 출발하여
(\hyperref[II.1.4.6-ko]{1.4.6})을 적용하면 충분하다. 마찬가지로
(iv)를 증명하려면 표준 동형
$\mathbf{S}(g^*(\mathcal{E}))
\simeq g^*(\mathbf{S}(\mathcal{E}))$
(\hyperref[II.1.7.5-ko]{1.7.5})에서 출발하여
(\hyperref[II.1.5.2-ko]{1.5.2})를 적용하면 충분하다. 끝으로 (v)를
보이려면 준동형 $\mathcal{E}\to\mathcal{F}$가 전사일 때 이에 대응하는
$\mathcal{O}_S$-대수 준동형
$\mathbf{S}(\mathcal{E})\to\mathbf{S}(\mathcal{F})$도 전사임을
관찰하면 충분하며, 결론은 (\hyperref[II.1.4.10-ko]{1.4.10})에서
따른다.
\end{proof}

\begin{env}[1.7.12]
\label{II.1.7.12-ko}
특히 $\mathcal{E}=\mathcal{O}_S$로 잡자. 준스킴
$\mathbf{V}(\mathcal{O}_S)$는 $\mathcal{O}_S$-대수
$\mathbf{S}(\mathcal{O}_S)$의 스펙트럼인 아핀 $S$-스킴이며, 이
대수는 $\mathcal{O}_S$-대수
$\mathcal{O}_S[T]=\mathcal{O}_S\otimes_{\mathbf{Z}}\mathbf{Z}[T]$와
\oldpage[II]{18}
동일시된다($T$는 부정원이다). $S=\operatorname{Spec}(\mathbf{Z})$일
때에는 (\hyperref[II.1.7.6-ko]{1.7.6})에 따라 자명하며, 구조 사상
$S\to\operatorname{Spec}(\mathbf{Z})$를 고려하고
(\hyperref[II.1.7.11-ko]{1.7.11, (iv)})를 사용하면 여기서 일반적인
경우로 넘어갈 수 있다. 이 결과에 따라
$\mathbf{V}(\mathcal{O}_S)=S[T]$로도 쓰며, 따라서 다음 공식을 얻는다.
\[
  S[T]=S\otimes_{\mathbf{Z}}\mathbf{Z}[T].
  \tag{1.7.12.1}\label{II.1.7.12.1-ko}
\]

(\hyperref[I.3.3.15-ko]{I, 3.3.15})에서 이미 본, $S[T]$의
$S$-단면의 싹으로 이루어진 층과 $\mathcal{O}_S$의 동일시는 여기에서
(\hyperref[II.1.7.9-ko]{1.7.9})의 특수한 경우로서 더 일반적인 맥락
안에서 다시 얻어진다.
\end{env}

\begin{env}[1.7.13]
\label{II.1.7.13-ko}
임의의 $S$-준스킴 $X$에 대하여,
(\hyperref[II.1.7.8-ko]{1.7.8})에서
$\operatorname{Hom}_S(X,S[T])$가
$\operatorname{Hom}_{\mathcal{O}_S}(\mathcal{O}_S,\mathcal{A}(X))$와
표준적으로 동일시됨을 보았다. 이 집합은
$\Gamma(S,\mathcal{A}(X))$와 표준적으로 동형이므로 환 구조가 주어진다.
또한 임의의 $S$-사상 $h:X\to Y$에는 이 환 구조들에 관한 사상
$\Gamma(\mathcal{A}(h)):\Gamma(S,\mathcal{A}(Y))
\to\Gamma(S,\mathcal{A}(X))$가 대응한다
(\hyperref[II.1.1.2-ko]{1.1.2}).
$\operatorname{Hom}_S(X,S[T])$에 이와 같이 정의된 환 구조를 주면,
따라서 $\operatorname{Hom}_S(X,S[T])$를 $S$-준스킴의 범주에서 환의
범주로 가는 \emph{$X$에 관한 반변 함자}로 볼 수 있다.
한편 $\operatorname{Hom}_S(X,\mathbf{V}(\mathcal{E}))$도 마찬가지로
$\operatorname{Hom}_{\mathcal{O}_S}(\mathcal{E},\mathcal{A}(X))$와
동일시된다(여기서는 $\mathcal{A}(X)$를
\emph{$\mathcal{O}_S$-가군}으로 본다). 따라서 여기에 환
$\operatorname{Hom}_S(X,S[T])$ 위의 \emph{가군} 구조를 표준적으로
줄 수 있으며, 위와 마찬가지로 쌍
\[
  \bigl(\operatorname{Hom}_S(X,S[T]),
  \operatorname{Hom}_S(X,\mathbf{V}(\mathcal{E}))\bigr)
\]
은 $X$에 관한 반변 함자임을 알 수 있다. 이 함자는 환 $A$와
$A$-가군 $M$으로 이루어진 쌍 $(A,M)$을 대상으로 하고 디-준동형들을
사상으로 하는 범주에 값을 갖는다.

이 사실들을 $S[T]$가 \emph{$S$ 위의 환 스킴}이고
$\mathbf{V}(\mathcal{E})$가 $S$ 위의 환 스킴 $S[T]$ 위의
\emph{$S$ 위의 가군 스킴}이라는 말로 해석하겠다
(0장, \textsection8 참조).
\end{env}

\begin{env}[1.7.14]
\label{II.1.7.14-ko}
$S$-스킴 $\mathbf{V}(\mathcal{E})$ 위에 정의된 $S$ 위의 가군 스킴
구조를 사용하여 $\mathcal{O}_S$-가군 $\mathcal{E}$를 유일한 동형을
제외하고 복원할 수 있음을 보이겠다. 이를 위해 $\mathcal{E}$가 이
구조를 사용하여 정의되는
$\mathbf{S}(\mathcal{E})=\mathcal{A}(\mathbf{V}(\mathcal{E}))$의 한
부분-$\mathcal{O}_S$-가군과 표준적으로 동형임을 보이겠다. 실제로
(\hyperref[II.1.7.4-ko]{1.7.4})에 따라
\emph{$\mathcal{O}_S$-대수} 준동형들의 집합
$\operatorname{Hom}_{\mathcal{O}_S}(\mathbf{S}(\mathcal{E}),
\mathcal{A}(X))$는 \emph{$\mathcal{O}_S$-가군} 준동형들의 집합
$\operatorname{Hom}_{\mathcal{O}_S}(\mathcal{E},\mathcal{A}(X))$와
표준적으로 동일시된다. $h$, $h'$을 이 뒤의 집합의 두 원소,
$s_i$ ($1\leq i\leq n$)를 열린집합 $U\subset S$ 위의
$\mathcal{E}$의 단면들, $t$를 $U$ 위의 $\mathcal{A}(X)$의 단면이라
하면, 정의에 따라
\[
  (h+h')(s_1s_2\cdots s_n)
  =\prod_{i=1}^n\bigl(h(s_i)+h'(s_i)\bigr)
\]
이고
\[
  (t.h)(s_1s_2\cdots s_n)=t^n\prod_{i=1}^n h(s_i).
\]

이제 $z$가 $U$ 위의 $\mathbf{S}(\mathcal{E})$의 단면이면,
$h\mapsto h(z)$는
$\operatorname{Hom}_S(X,\mathbf{V}(\mathcal{E}))
=\operatorname{Hom}_{\mathcal{O}_S}(\mathbf{S}(\mathcal{E}),
\mathcal{A}(X))$에서 $\Gamma(U,\mathcal{A}(X))$로 가는 사상이다.
이제 다음을
\oldpage[II]{19}
보이겠다. $\mathcal{E}$는 다음 성질을 갖는
$\mathbf{S}(\mathcal{E})$의 부분가군과 동일시된다.
\emph{모든 열린집합 $U\subset S$, 이 부분-$\mathcal{O}_S$-가군의
$U$ 위의 모든 단면 $z$, 모든 $S$-준스킴 $X$에 대하여,
$\operatorname{Hom}_{\mathcal{O}_S|U}(\mathbf{S}(\mathcal{E})|U,
\mathcal{A}(X)|U)$에서 $\Gamma(U,\mathcal{A}(X))$로 가는 사상
$h\mapsto h(z)$가 $\Gamma(U,\mathcal{A}(X))$-가군 준동형이다.}

실제로 $\mathcal{E}$가 이 성질을 갖는다는 것은 곧바로 알 수 있다.
그 역을 증명하려면, $S=\operatorname{Spec}(A)$,
$\mathcal{E}=\widetilde{M}$일 때 $S$ 위의
$\mathbf{S}(\mathcal{E})$의 단면 $z$가 ($U=S$에 대하여) 위의
성질을 가지면 반드시 $\mathcal{E}$의 단면임을 증명하는 것으로
충분하다. 이때 $z=\sum_{n=0}^{\infty}z_n$이고
$z_n\in\mathbf{S}_n(M)$이며, $n\ne1$이면 $z_n=0$임을 보이면 된다.
$B=\mathbf{S}(M)$으로 놓고 $X$로 준스킴
$\operatorname{Spec}(B[T])$를 택하자. 여기서 $T$는 부정원이다.
집합
$\operatorname{Hom}_{\mathcal{O}_S}(\mathbf{S}(\mathcal{E}),
\mathcal{A}(X))$는 환 준동형 $h:B\to B[T]$들의 집합과 동일시된다
(\hyperref[I.1.3.13-ko]{I, 1.3.13}). 위에서 본 바에 따라
$(T.h)(z)=\sum_{n=0}^{\infty}T^nh(z_n)$이며, $z$에 관한 가정은 모든
준동형 $h$에 대하여
$\sum_{n=0}^{\infty}T^nh(z_n)=T.\sum_{n=0}^{\infty}h(z_n)$임을
뜻한다. 특히 $h$로 표준 포함 사상을 택하면
$\sum_{n=0}^{\infty}T^nz_n=T.\sum_{n=0}^{\infty}z_n$을 얻으며,
여기서 실제로 $n\ne1$이면 $z_n=0$이라는 결론이 따른다.
\end{env}

\begin{proposition}[1.7.15]
\label{II.1.7.15-ko}
$Y$를 바탕공간이 뇌터인 준스킴 또는 준콤팩트 스킴이라 하자.
$Y$ 위에서 유한형인 임의의 아핀 $Y$-스킴 $X$는
$\mathbf{V}(\mathcal{E})$ 꼴의 $Y$-스킴의 닫힌 부분-$Y$-스킴과
$Y$-동형이다. 여기서 $\mathcal{E}$는 유한 생성 준연접
$\mathcal{O}_Y$-가군이다.
\end{proposition}

\begin{proof}
실제로 준연접 $\mathcal{O}_Y$-대수 $\mathcal{A}(X)$는 유한 생성이다
(\hyperref[II.1.3.7-ko]{1.3.7}). 가정에 따라 $\mathcal{A}(X)$는
유한 생성 준연접 부분-$\mathcal{O}_Y$-가군 $\mathcal{E}$에 의해
대수로서 생성된다(\hyperref[I.9.6.5-ko]{I, 9.6.5}). 정의에 따라
이는 포함 사상 $\mathcal{E}\to\mathcal{A}(X)$를 표준적으로 연장하는
표준 준동형 $\mathbf{S}(\mathcal{E})\to\mathcal{A}(X)$가
\emph{전사}임을 뜻한다. 그러면 결론은
(\hyperref[II.1.4.10-ko]{1.4.10})에서 따른다.
\end{proof}

\section{동차 소 스펙트럼}
\label{section:II.2-ko}

\subsection{등급환과 등급가군에 관한 일반 사항.}
\label{subsection:II.2.1-ko}

\begin{notation}[2.1.1]
\label{II.2.1.1-ko}
음이 아닌 차수로 \emph{등급}이 주어진 환 $S$에 대하여,
차수가 $n$인 동차 원소들로 이루어진 $S$의 부분을 $S_n$으로
($n\geq0$), $n>0$인 $S_n$들의 합(직합)을 $S_+$로 나타내겠다.
$1\in S_0$이고, $S_0$는 $S$의 부분환이며, $S_+$는 $S$의 동차
아이디얼이고, $S$는 $S_0$와 $S_+$의 직합이다. $M$이 $S$ 위의
\emph{등급가군}이면(음의 차수도 허용한다), 마찬가지로
$M$의 차수가 $n$인 동차 원소들로 이루어진 $S_0$-가군을
$M_n$으로 나타내겠다(여기서 $n\in\mathbf{Z}$).

임의의 정수 $d>0$에 대하여, $S_{nd}$들의 직합을 $S^{(d)}$로
나타낸다. $S_{nd}$의 원소들을 차수가 $n$인 동차 원소로 보면,
$S_{nd}$들은 $S^{(d)}$ 위에 등급환 구조를 정의한다.

$0\leq k\leq d-1$인 임의의 정수 $k$에 대하여, $M^{(d,k)}$로
다음 직합을 나타내겠다.
\oldpage[II]{20}
$M_{nd+k}$ ($n\in\mathbf{Z}$)들의 직합이다.
$M_{nd+k}$의 원소들을 차수가 $n$인 동차 원소로 보면, 이는 등급
$S^{(d)}$-가군이다. $M^{(d,0)}$ 대신 $M^{(d)}$로 쓰겠다.

앞의 표기법 아래에서, 임의의 정수 $n$(음수도 허용한다)에 대하여,
모든 $k\in\mathbf{Z}$에 대해 $(M(n))_k=M_{n+k}$로 정의되는 등급
$S$-가군을 $M(n)$으로 나타내겠다. 특히 $S(n)$은
$(S(n))_k=S_{n+k}$를 만족하는 등급 $S$-가군이며,
$n<0$이면 $S_n=0$으로 놓기로 한다. 등급 $S$-가군 $M$이
\emph{등급가군으로서} $S(n)$ 꼴의 가군들의 직합과 동형이면,
\emph{자유}라고 한다. $S(n)$은 $S$의 원소 $1$을 차수가
$-n$인 원소로 보아 이를 생성원으로 갖는
$S$-가군이므로, 이는 $M$이 \emph{동차} 원소들로 이루어진
$S$ 위의 \emph{기저}를 갖는다고 말하는 것과 같다.

등급 $S$-가군 $M$이 \emph{유한 표시를 갖는다}는 것은, $P$와 $Q$가
$S(n)$ 꼴의 가군들의 유한 직합이고 준동형들의 차수가 $0$인 완전열
$P\to Q\to M\to0$이 존재한다는 뜻이다
(\agref{II.2.1.2-ko}{2.1.2} 참조).
\end{notation}

\begin{env}[2.1.2]
\phantomsection
\label{II.2.1.2-ko}
$M$, $N$을 두 등급 $S$-가군이라 하자. $M\otimes_S N$ 위에 다음과
같이 \emph{등급} $S$-가군 구조를 정의한다. 텐서곱
$M\otimes_{\mathbf{Z}}N$ 위에는
$\mathbf{Z}$-가군의 등급 구조($\mathbf{Z}$에는
$\mathbf{Z}_0=\mathbf{Z}$, $n\ne0$일 때 $\mathbf{Z}_n=0$인 등급을
준다)를 다음과 같이 정의할 수 있다.
\[
  (M\otimes_{\mathbf{Z}}N)_q
  =\bigoplus_{m+n=q}M_m\otimes_{\mathbf{Z}}N_n
\]
($M$과 $N$은 각각 $M_m$들과 $N_n$들의 직합이므로,
$M\otimes_{\mathbf{Z}}N$을 모든 $M_m\otimes_{\mathbf{Z}}N_n$의
직합과 표준적으로 동일시할 수 있음을 알고 있다.) 이제
$M\otimes_S N=(M\otimes_{\mathbf{Z}}N)/P$이다. 여기서 $P$는
$x\in M$, $y\in N$, $s\in S$에 대하여
$(xs)\otimes y-x\otimes(sy)$ 꼴의 원소들로 생성되는
$M\otimes_{\mathbf{Z}}N$의 부분-$\mathbf{Z}$-가군이다. $P$가
$M\otimes_{\mathbf{Z}}N$의 \emph{등급} 부분-$\mathbf{Z}$-가군임은
명백하며, 몫을 취하면 실제로 $M\otimes_S N$ 위에 등급 $S$-가군
구조가 얻어짐을 곧바로 알 수 있다.

두 등급 $S$-가군 $M$, $N$에 대하여, $S$-가군 준동형
$u:M\to N$이 모든 $j\in\mathbf{Z}$에 대하여
$u(M_j)\subset N_{j+k}$를 만족하면 \emph{차수 $k$}라고 함을
상기하자. $M$에서 $N$으로 가는 차수 $n$인 모든 준동형의 집합을
$H_n$이라 하고, $M$에서 $N$으로 가는 모든 ($S$-가군) 준동형들의
$S$-가군 $H$ 안에서 $H_n$들의 \emph{합}(직합)($n\in\mathbf{Z}$)을
$\operatorname{Hom}_S(M,N)$으로 나타낸다. 일반적으로
$\operatorname{Hom}_S(M,N)$은 앞서 말한 가군 전체와 같지 않다. 다만 $M$이
\emph{유한 생성}이면 $H=\operatorname{Hom}_S(M,N)$이다. 실제로 이때
$M$이 유한 개의 동차 원소 $x_i$ ($1\leq i\leq n$)로 생성된다고
할 수 있고, 모든 준동형 $u\in H$는 유일하게
$\sum_{k\in\mathbf{Z}}u_k$로 쓰인다. 여기서 각 $k$에 대하여
$u_k(x_i)$는 $u(x_i)$의 차수가 $k+\deg(x_i)$인 동차 성분과 같고
($1\leq i\leq n$), 따라서 유한 개의 지표를 제외하면 $u_k=0$이다.
정의에 따라 $u_k\in H_k$이므로 결론을 얻는다.

$\operatorname{Hom}_S(M,N)$의 차수 $0$인 원소들을 \emph{등급
$S$-가군의 준동형}이라 한다. $S_mH_n\subset H_{m+n}$임은 명백하므로,
$H_n$들은 $\operatorname{Hom}_S(M,N)$ 위에 등급 $S$-가군 구조를
정의한다.

이 정의들로부터 두 등급 $S$-가군 $M$, $N$에 대하여 다음을 곧바로
얻는다.
\[
  \label{II.2.1.2.1-ko}
  M(m)\otimes_S N(n)=(M\otimes_S N)(m+n),
  \tag{2.1.2.1}
\]
\[
  \label{II.2.1.2.2-ko}
  \operatorname{Hom}_S(M(m),N(n))
  =(\operatorname{Hom}_S(M,N))(n-m)
  \tag{2.1.2.2}
\]
\oldpage[II]{21}
$S$, $S'$을 두 등급환이라 하자. \emph{등급환}의 준동형
$\varphi:S\to S'$은 모든 $n\in\mathbf{Z}$에 대하여
$\varphi(S_n)\subset S'_n$을 만족하는 환 준동형이다(다시 말해,
$\varphi$는 등급 $\mathbf{Z}$-가군의 \emph{차수 $0$인} 준동형이어야
한다). 이러한 준동형 하나를 주면 $S'$ 위에 \emph{등급} $S$-가군 구조가
정의된다. 이 구조와 원래의 등급환 구조를 갖춘 $S'$을
\emph{등급 $S$-대수}라고 한다.

이제 $M$이 등급 $S$-가군이면, 등급 $S$-가군의 텐서곱
$M\otimes_S S'$ 위에 \emph{등급} $S'$-가군 구조가 자연스럽게 주어지며,
그 등급은 위에서 정의한 것이다.
\end{env}

\begin{lemma}[2.1.3]
\phantomsection
\label{II.2.1.3-ko}
$S$를 음이 아닌 차수로 등급이 주어진 환이라 하자. 동차 원소들로
이루어진 $S_+$의 부분집합 $E$가 $S_+$를 $S$-가군으로서 생성할
필요충분조건은 $E$가 $S$를 $S_0$-대수로서 생성하는 것이다.
\end{lemma}

\begin{proof}
조건이 충분함은 명백하다. 필요함을 보이자. $E$의 차수가 $n$인
원소들의 집합(각각 차수가 $\leq n$인 원소들의 집합)을 $E_n$
(각각 $E^n$)이라 하자. $n>0$에 대한 귀납법으로, $S_n$이
$E^n$의 원소들의 곱인 차수 $n$의 원소들로 생성되는 $S_0$-가군임을
보이면 충분하다. 그런데 $n=1$일 때 이는 가정에 의해 명백하다.
또한 가정으로부터
$S_n=\sum_{p=0}^{n-1}S_pE_{n-p}$이고, 이제 귀납 논법은 즉시
따른다.
\end{proof}

\begin{corollary}[2.1.4]
\phantomsection
\label{II.2.1.4-ko}
$S_+$가 유한 생성 아이디얼이기 위한 필요충분조건은 $S$가
$S_0$ 위의 유한 생성 대수인 것이다.
\end{corollary}

\begin{proof}
실제로 $S_0$-대수 $S$ (각각 $S$-아이디얼 $S_+$)의 유한 생성계가
동차 원소들로 이루어져 있다고 언제나 가정할 수 있다. 주어진 각 생성원을
그 동차 성분들로 바꾸면 되기 때문이다.
\end{proof}

\begin{corollary}[2.1.5]
\phantomsection
\label{II.2.1.5-ko}
$S$가 뇌터이기 위한 필요충분조건은 $S_0$가 뇌터이고 $S$가
$S_0$ 위의 유한 생성 대수인 것이다.
\end{corollary}

\begin{proof}
조건이 충분함은 명백하다. 필요함은 $S_0$가 $S/S_+$와 동형이고
$S_+$가 유한 생성 아이디얼이어야 한다는 데서 따른다
(\hyperref[II.2.1.4-ko]{2.1.4}).
\end{proof}

\begin{lemma}[2.1.6]
\phantomsection
\label{II.2.1.6-ko}
$S$를 음이 아닌 차수로 등급이 주어진 환이며 $S_0$ 위의 유한 생성
대수라고 하자. $M$을 유한 생성 등급 $S$-가군이라 하자. 그러면 다음이
성립한다.
\begin{enumerate}
  \item[\rm(i)] $M_n$들은 유한 생성 $S_0$-가군이고, 모든
    $n\leq n_0$에 대하여 $M_n=0$이 되는 정수 $n_0$가 존재한다.
  \item[\rm(ii)] 모든 정수 $n\geq n_1$에 대하여
    $M_{n+h}=S_hM_n$이 되는 정수 $n_1$과 정수 $h>0$이 존재한다.
  \item[\rm(iii)] $d>0$, $0\leq k\leq d-1$을 만족하는 모든 정수쌍
    $(d,k)$에 대하여 $M^{(d,k)}$는 유한 생성 $S^{(d)}$-가군이다.
  \item[\rm(iv)] 모든 정수 $d>0$에 대하여 $S^{(d)}$는 $S_0$ 위의
    유한 생성 대수이다.
  \item[\rm(v)] 모든 $m>0$에 대하여 $S_{mh}=(S_h)^m$이 되는 정수
    $h>0$이 존재한다.
  \item[\rm(vi)] 모든 정수 $n>0$에 대하여, 모든 $m\geq m_0$에 대해
    $S_m\subset S_+^n$이 되는 정수 $m_0$가 존재한다.
\end{enumerate}
\end{lemma}

\begin{proof}
$S$가 $S_0$-대수로서 차수가 $h_i$인 동차 원소 $f_i$
($1\leq i\leq r$)들로 생성되고, $M$이 $S$-가군으로서 차수가 $k_j$인
동차 원소 $x_j$ ($1\leq j\leq s$)들로 생성된다고 가정할 수 있다.
$M_n$은 $\alpha_i$들이 $\geq0$인 정수이고
$k_j+\sum_i\alpha_i h_i=n$을 만족할 때의 원소
\oldpage[II]{22}
$f_1^{\alpha_1}\cdots f_r^{\alpha_r}x_j$들의, 계수가 $S_0$에 속하는
선형결합들로
이루어짐이 명백하다. 각 $j$에 대하여 이 방정식을 만족하는
$(\alpha_i)$는 유한 개뿐이다. 이는 $h_i$들이 $>0$이기 때문이며, 따라서
(i)의 첫 번째 주장이 따른다. 두 번째 주장은 명백하다. 한편 $h$를
$h_i$들의 최소공배수라 하고, 모든 $g_i$의 차수가 $h$가 되도록
$g_i=f_i^{h/h_i}$ ($1\leq i\leq r$)로 놓자. 또 $z_\mu$들을
$0\leq\alpha_i<h/h_i$ ($1\leq i\leq r$)인
$f_1^{\alpha_1}\cdots f_r^{\alpha_r}x_j$ 꼴의 $M$의 원소들이라 하자.
이 원소들은 유한 개이고, 그 차수들 가운데 가장 큰 것을 $n_1$이라 하자.
$n\geq n_1$이면 $M_{n+h}$의 모든 원소가 $z_\mu$들의 선형결합이며 그
계수들은 $g_i$들에 관한 차수가 $>0$인 단항식임이 명백하다. 따라서
$M_{n+h}=S_hM_n$이고, 이로써 (ii)가 증명된다. 같은 방식으로 (모든
$d>0$에 대하여) $M^{(d,k)}$의 원소는 $g$가 $S$의 동차 원소이고
$0\leq\alpha_i<d$일 때의
$g^df_1^{\alpha_1}\cdots f_r^{\alpha_r}x_j$ 꼴의 원소들의, 계수가
$S_0$에 속하는 선형결합임을 알 수 있다. 따라서 (iii)가 성립한다. 이제
$M=S_+$로 놓고 $(S_+)^{(d)}=(S^{(d)})_+$임을 이용하면, (iv)는
(iii)와 보조정리 (\hyperref[II.2.1.3-ko]{2.1.3})에서 따른다. (v)는
(ii)에서 $M=S$로 놓으면 얻어진다. 마지막으로 주어진 $n$에 대하여
$\alpha_i\geq0$이고 $\sum_i\alpha_i<n$인 계 $(\alpha_i)$는 유한 개뿐이다.
따라서 이 계들에 대한 합 $\sum_i\alpha_i h_i$의 최댓값을 $m_0$라 하면
$m>m_0$에 대하여 $S_m\subset S_+^n$이고, 이로써 (vi)가 증명된다.
\footnote{\textbf{번역 주.} 인쇄 원문의 (vi)는 $m\geq m_0$라고 서술하지만,
여기서 $m_0$로 정의한 바로 그 최댓값을 쓰면 증명이 직접 주는 범위는
$m>m_0$이다. 존재명제 자체는 이 최댓값에 $1$을 더한 수를 임계값으로
취하면 원문의 $m\geq m_0$ 형태로 성립한다.}
\end{proof}

\begin{corollary}[2.1.7]
\phantomsection
\label{II.2.1.7-ko}
$S$가 뇌터이면 모든 정수 $d>0$에 대하여 $S^{(d)}$도 뇌터이다.
\end{corollary}

\begin{proof}
이는 (\hyperref[II.2.1.5-ko]{2.1.5})와
(\hyperref[II.2.1.6-ko]{2.1.6, (iv)})에서 따른다.
\end{proof}

\begin{env}[2.1.8]
\phantomsection
\label{II.2.1.8-ko}
$\mathfrak{p}$를 등급환 $S$의 동차 소아이디얼이라 하자. 그러면
$\mathfrak{p}$는 부분군 $\mathfrak{p}_n=\mathfrak{p}\cap S_n$들의
직합이다. $\mathfrak{p}$가 $S_+$를 포함하지 않는다고 가정하자. 그러면
$\mathfrak{p}$에 속하지 않는 $f\in S_+$에 대하여, 관계
$f^nx\in\mathfrak{p}$는 $x\in\mathfrak{p}$와 동치이다. 특히
$f\in S_d$ ($d>0$)이면, 모든 $x\in S_{m-nd}$에 대하여 관계
$f^nx\in\mathfrak{p}_m$은 $x\in\mathfrak{p}_{m-nd}$와 동치이다.
\end{env}

\begin{proposition}[2.1.9]
\phantomsection
\label{II.2.1.9-ko}
$n_0$를 $>0$인 정수라 하고, 각 $n\geq n_0$에 대하여
$\mathfrak{p}_n$을 $S_n$의 부분군이라 하자. $S_+$를 포함하지 않으며
모든 $n\geq n_0$에 대하여
$\mathfrak{p}\cap S_n=\mathfrak{p}_n$을 만족하는 $S$의 동차
소아이디얼 $\mathfrak{p}$가 존재하기 위한 필요충분조건은 다음 조건들이
성립하는 것이다.
\begin{enumerate}
  \item[$1^\circ$] 모든 $m\geq 0$과 모든 $n\geq n_0$에 대하여
    $S_m\mathfrak{p}_n\subset\mathfrak{p}_{m+n}$이다.
  \item[$2^\circ$] $m\geq n_0$, $n\geq n_0$, $f\in S_m$, $g\in S_n$일 때,
    관계 $fg\in\mathfrak{p}_{m+n}$은 $f\in\mathfrak{p}_m$ 또는
    $g\in\mathfrak{p}_n$을 함의한다.
  \item[$3^\circ$] 적어도 하나의 $n\geq n_0$에 대하여
    $\mathfrak{p}_n\neq S_n$이다.
\end{enumerate}
또한 이때 동차 소아이디얼 $\mathfrak{p}$는 유일하다.
\end{proposition}

\begin{proof}
조건 $1^\circ$와 $2^\circ$가 필요함은 명백하다. 또한
$\mathfrak{p}\not\supset S_+$이면
$\mathfrak{p}\cap S_k\neq S_k$인 $k>0$이 적어도 하나 존재한다.
$f\in S_k$가 $\mathfrak{p}$에 속하지 않으면, 관계
$\mathfrak{p}\cap S_n=S_n$은
(\hyperref[II.2.1.8-ko]{2.1.8})에 따라
$\mathfrak{p}\cap S_{n-mk}=S_{n-mk}$를 함의한다. 따라서 어떤 값 이상의
모든 $n$에 대하여 $\mathfrak{p}\cap S_n=S_n$이라면
$\mathfrak{p}\supset S_+$가 되어 가정에 모순이다. 이로써 $3^\circ$가
필요함이 증명된다. 역으로 조건 $1^\circ$, $2^\circ$, $3^\circ$가
성립한다고 가정하자. 어떤 정수 $d\geq n_0$에 대하여
$f\in S_d$가 $\mathfrak{p}_d$에 속하지 않으면, $\mathfrak{p}$가
존재할 때 $m<n_0$에 대한 $\mathfrak{p}_m$은 유한 개의 값을 제외한 모든
$r$에 대하여 $f^rx\in\mathfrak{p}_{m+rd}$를 만족하는
$x\in S_m$들의 집합과 반드시 같음에 유의하자. 이는 $\mathfrak{p}$가 존재한다면
유일함을 이미 증명한다. 이제 앞의 조건으로 $m<n_0$에 대한
$\mathfrak{p}_m$들을 정의하면
$\mathfrak{p}=\sum_{n=0}^{\infty}\mathfrak{p}_n$이 소아이디얼임을 보이는
일만 남는다. 먼저 $2^\circ$에 따라, $m\geq n_0$에 대해서도
$\mathfrak{p}_m$은 유한 개의 값을 제외한 모든 $r$에 대하여
$f^rx\in\mathfrak{p}_{m+rd}$를 만족하는 $x\in S_m$들의 집합으로
정의됨에 유의하자. 이제
\oldpage[II]{23}
$g\in S_m$, $x\in\mathfrak{p}_n$이면, 유한 개의 값을 제외한 모든
$r$에 대하여 $f^rgx\in\mathfrak{p}_{m+n+rd}$이다. 따라서
$gx\in\mathfrak{p}_{m+n}$이고, 이는 $\mathfrak{p}$가 $S$의
아이디얼임을 증명한다. 이 아이디얼이 소아이디얼임을, 다시 말해 부분군
$S_n/\mathfrak{p}_n$들로 등급이 주어진 환 $S/\mathfrak{p}$가 정역임을
보이기 위해서는 ($S/\mathfrak{p}$의 두 원소의 최고차 동차성분들을
고려하여) $x\in S_m$, $y\in S_n$이
$x\notin\mathfrak{p}_m$, $y\notin\mathfrak{p}_n$을 만족하면
$xy\notin\mathfrak{p}_{m+n}$임을 보이면 충분하다. 그렇지 않다면 충분히
큰 $r$에 대하여 $f^{2r}xy\in\mathfrak{p}_{m+n+2rd}$일 것이다. 그러나
모든 $r>0$에 대하여 $f^ry\notin\mathfrak{p}_{n+rd}$이다. 그러면
$2^\circ$에 따라 유한 개의 값을 제외한 모든 $r$에 대하여
$f^rx\in\mathfrak{p}_{m+rd}$이고, 이로부터
$x\in\mathfrak{p}_m$을 얻어 가정에 모순이다.
\end{proof}

\begin{env}[2.1.10]
\phantomsection
\label{II.2.1.10-ko}
$S_+$의 부분집합 $\mathfrak{J}$가 $S$의 아이디얼이면 이를
\emph{$S_+$의 아이디얼}이라 한다. 또한 $\mathfrak{J}$가 $S_+$와,
\emph{$S_+$를 포함하지 않는} $S$의 동차 소아이디얼의 교집합이면 이를
\emph{$S_+$의 동차 소아이디얼}이라 한다(더욱이 이 소아이디얼은
(\hyperref[II.2.1.9-ko]{2.1.9})에 따라 유일하다). $\mathfrak{J}$가
$S_+$의 아이디얼이면, \emph{$S_+$에서의 $\mathfrak{J}$의 근기}는 어떤
거듭제곱이 $\mathfrak{J}$에 속하는 $S_+$의 원소들로 이루어진 집합이다.
다시 말해 이는 $r_+(\mathfrak{J})=r(\mathfrak{J})\cap S_+$이다. 특히
$S_+$에서 $0$의 근기는 다시 $S_+$의 \emph{멱영근기}라 하고
$\mathfrak{N}_+$로 나타낸다. 따라서 이는 $S_+$의 멱영원 전체의
집합이다. $\mathfrak{J}$가 $S_+$의 \emph{동차} 아이디얼이면 그 근기
$r_+(\mathfrak{J})$도 동차 아이디얼이다. 실제로 몫환
$S/\mathfrak{J}$로 넘어가면 $\mathfrak{J}=0$인 경우로 환원할 수 있고,
$x=x_h+x_{h+1}+\cdots+x_k$가 멱영원이면
$x_i\in S_i$ ($1\leq h\leq i\leq k$)들도 모두 멱영원임을 보이면
충분하다. $x_k\neq 0$이라고 가정할 수 있고, 이때 $x^n$의 최고차
성분은 $x_k^n$이므로 $x_k$는 멱영원이다. 이제 $k$에 대한 귀납법을
적용한다. 등급환 $S$가 \emph{본질적으로 축소}라는 것은
$\mathfrak{N}_+=0$이라는 뜻이다. 다시 말해 $S_+$가 $\neq 0$인 멱영원을
포함하지 않는다는 뜻이다.
\end{env}

\begin{env}[2.1.11]
\phantomsection
\label{II.2.1.11-ko}
등급환 $S$에서 원소 $x$가 어떤 영이 아닌 원소와 곱하여 $0$이 되는
영인자이면 그 최고차 동차성분도 영인자임에 유의하자. 환 $S$가
\emph{본질적으로 정역}이라는 것은 $S_+$가
(\emph{단위원을 갖지 않는 환으로서}) 영인자를 포함하지 않아 영이 아닌
두 원소의 곱이 $0$일 수 없고, 또한 $\neq 0$인 것을 뜻한다. 이를 위해서는
$S_+$의 동차 원소 가운데 $\neq 0$인 것은 모두 이 환에서 영인자가 아니어서
어떤 영이 아닌 원소와 곱해도 $0$이 되지 않으면 충분하다.
$\mathfrak{p}$가 $S_+$의
동차 소아이디얼이면 $S/\mathfrak{p}$가 본질적으로 정역임은 명백하다.

$S$를 본질적으로 정역인 등급환이라 하고 $x_0\in S_0$라 하자. 이때
$S_+$의 동차 원소 $f\neq 0$가 \emph{하나라도} 존재하여
$x_0f=0$이면 $x_0S_+=0$이다. 실제로 모든 $g\in S_+$에
대하여 $(x_0g)f=(x_0f)g=0$이고, 가정에 따라 $x_0g=0$이다. 따라서
$S$가 정역이기 위한 필요충분조건은 $S_0$가 정역이고 $S_0$에서
$S_+$의 소멸자가 $0$뿐인 것이다.
\end{env}

\subsection{등급환의 분수환.}
\label{subsection:II.2.2-ko}

\begin{env}[2.2.1]
\phantomsection
\label{II.2.2.1-ko}
$S$를 음이 아닌 차수로 등급이 주어진 환이라 하고, $f$를 차수가
$d>0$인 $S$의 \emph{동차} 원소라 하자. 그러면 분수환 $S'=S_f$는
$S'_n$을 $x\in S_{n+kd}$이고 $k\geq 0$일 때의 $x/f^k$ 전체의
집합으로 놓음으로써 등급환이 된다(여기서 $n$은 임의의 음의 값도 취할
수 있음에 유의하자). $S'$의 \emph{차수가 $0$인} 원소들로 이루어진
부분환 $S'_0=(S_f)_0$을 $S_{(f)}$로 나타내겠다.

$f\in S_d$이면 $S_f$의 단항식 $(f/1)^h$들($h$는 영을 포함한 임의의
정수)은 환 $S_{(f)}$ 위의 \emph{자유계}를 이루며, 그 선형결합 전체의
집합은 바로
\oldpage[II]{24}
환 $(S^{(d)})_f$이다. 따라서 이 환은
\emph{$S_{(f)}[T,T^{-1}]=S_{(f)}\otimes_{\mathbf{Z}}\mathbf{Z}[T,T^{-1}]$와
동형이다}($T$는 부정원이다). 실제로
$\sum_{h=-a}^{b}z_h(f/1)^h=0$인 관계가 있고 $z_h=x_h/f^m$이며 모든
$x_h$가 $S_{md}$에 속한다고 하자. 정의에 따라 이 관계는
어떤 $k\geq a$가 존재하여
$\sum_{h=-a}^{b}f^{h+k}x_h=0$이 되는 것과 동치이다. 이 합의 항들은
차수가 서로 다르므로 모든 $h$에 대하여 $f^{h+k}x_h=0$이고, 따라서
모든 $h$에 대하여 $z_h=0$이다.

$M$이 등급 $S$-가군이면 $M'=M_f$는 등급 $S_f$-가군이다. 여기서
$M'_n$은 $z/f^k$ 꼴의 원소들로 이루어진 집합이며, 여기서
$z\in M_{n+kd}$이고 $k\geq 0$이다.
$M'$의 차수가 $0$인 동차 원소들의 집합을 $M_{(f)}$로 나타낸다.
$M_{(f)}$가 $S_{(f)}$-가군이고
$(M^{(d)})_f=M_{(f)}\otimes_{S_{(f)}}(S^{(d)})_f$임은 바로 알 수 있다.
\end{env}

\begin{lemma}[2.2.2]
\phantomsection
\label{II.2.2.2-ko}
$d$, $e$를 $>0$인 두 정수, $f\in S_d$, $g\in S_e$라 하자. 다음과 같은
환의 표준 동형이 존재한다.
\[
  S_{(fg)}\xrightarrow{\sim}(S_{(f)})_{g^d/f^e}~;
\]
이 두 환을 표준적으로 동일시하면, 다음과 같은 가군의 표준 동형이
존재한다.
\[
  M_{(fg)}\xrightarrow{\sim}(M_{(f)})_{g^d/f^e}.
\]
\end{lemma}

\begin{proof}
실제로 $fg$는 $f^e g^d$의 약수이고, 후자는 $(fg)^{de}$의 약수이므로
등급환 $S_{fg}$와 $S_{f^e g^d}$는 표준적으로 동일시된다. 다른 한편
$S_{f^e g^d}$는 $(S_{f^e})_{g^d/1}$와도 동일시되고
(\hyperref[0.1.4.6-ko]{0, I.4.6}), $f^e/1$은 $S_{f^e}$에서 가역이므로
$S_{f^e g^d}$는 $(S_{f^e})_{g^d/f^e}$와도 동일시된다. 그런데
$g^d/f^e$는 차수가 $0$인 $S_{f^e}$의 원소이다. 따라서
$(S_{f^e})_{g^d/f^e}$에서 차수가 $0$인 원소들로 이루어진 부분환은
$(S_{(f^e)})_{g^d/f^e}$임을 곧바로 얻는다. 또 명백히
$S_{(f^e)}=S_{(f)}$이므로, 이는 명제의 첫 번째 부분을 증명한다.
두 번째 부분도 같은 방식으로 확립된다.
\end{proof}

\begin{env}[2.2.3]
\phantomsection
\label{II.2.2.3-ko}
(\hyperref[II.2.2.2-ko]{2.2.2})의 가정 아래에서, 표준 준동형
$S_f\to S_{fg}$ (\hyperref[0.1.4.1-ko]{0, I.4.1})은 $x/f^k$를
$g^k x/(fg)^k$로 보내며 차수가 $0$임은 명백하다. 따라서
제한함으로써
\emph{표준 준동형 $S_{(f)}\to S_{(fg)}$}을 얻으며, 다음 도식은
가환한다.
\[
  \xymatrix{
    & S_{(f)}\ar[dl]\ar[dr]\\
    S_{(fg)}\ar[rr]^-{\sim} && (S_{(f)})_{g^d/f^e}
  }
\]
마찬가지로 표준 준동형 $M_{(f)}\to M_{(fg)}$을 정의한다.
\end{env}

\begin{lemma}[2.2.4]
\phantomsection
\label{II.2.2.4-ko}
$f$, $g$가 $S_+$의 두 동차 원소이면, 환 $S_{(fg)}$는 $S_{(f)}$와
$S_{(g)}$의 표준 상들의 합집합에 의해 생성된다.
\end{lemma}

\begin{proof}
(\hyperref[II.2.2.2-ko]{2.2.2})에 따라
$1/(g^d/f^e)=f^{d+e}/(fg)^d$가 $S_{(g)}$의 $S_{(fg)}$ 안에서의 표준
상에 속함을 보이면 충분하며, 이는 정의상 명백하다.
\end{proof}

\begin{proposition}[2.2.5]
\phantomsection
\label{II.2.2.5-ko}
$d$를 $>0$인 정수, $f\in S_d$라 하자. 그러면 다음과 같은 환의
표준 동형이 존재한다.
$S_{(f)}\xrightarrow{\sim}S^{(d)}/(f-1)S^{(d)}$~; 이 동형으로 두 환을
표준적으로 동일시하면, 다음과 같은 가군의 표준 동형이 존재한다.
$M_{(f)}\xrightarrow{\sim}M^{(d)}/(f-1)M^{(d)}$.
\end{proposition}

\begin{proof}
이 동형들 중 첫 번째 것은 $x/f^n$($x\in S_{nd}$)에 원소
$\overline{x}$, 곧 $x$의 $(f-1)S^{(d)}$에 대한 잉여류를 대응시켜 정의한다.
이 사상은 잘 정의된다. 실제로 모든 양의 정수 지수에 대하여 합동식
$f^h x\equiv x\pmod{(f-1)S^{(d)}}$이 모든 $x\in S^{(d)}$에 대해
성립한다. 따라서 $f^h x=0$인 $h>0$가 있으면
\oldpage[II]{25}
$\overline{x}=0$이다. 한편 $x\in S_{nd}$이고 $x=(f-1)y$라 하자. 여기서
$y=y_{hd}+y_{(h+1)d}+\cdots+y_{kd}$이고 $y_{jd}\in S_{jd}$이며
$y_{hd}\neq 0$이다. 그러면 반드시 $h=n$이고 $x=-y_{hd}$이며, 또한
관계 $y_{(j+1)d}=fy_{jd}$가 $h\leq j\leq k-1$에 대해 성립하고
$fy_{kd}=0$이다.
이로부터 마침내 $f^{k-n+1}x=0$을 얻는다. 따라서 $x\in S_{nd}$인 원소가
나타내는 $(f-1)S^{(d)}$에 대한 각 잉여류 $\overline{x}$에 $S_{(f)}$의 원소
$x/f^n$을 대응시킬 수 있다.
앞의 고찰에 따라 이 사상은 잘 정의된다. 이렇게 정의된 두 사상이
환 준동형이며 서로 역임을 곧바로 알 수 있다. $M$에 대해서도 완전히
같은 방식으로 한다.
\end{proof}

\begin{corollary}[2.2.6]
\phantomsection
\label{II.2.2.6-ko}
$S$가 뇌터이면, 차수가 $>0$인 임의의 동차 원소 $f$에 대하여
$S_{(f)}$도 뇌터이다.
\end{corollary}

\begin{proof}
이는 (\hyperref[II.2.1.7-ko]{2.1.7})과
(\hyperref[II.2.2.5-ko]{2.2.5})에서 곧바로 따른다.
\end{proof}

\begin{env}[2.2.7]
\phantomsection
\label{II.2.2.7-ko}
$T$를 \emph{동차} 원소들로 이루어진 $S_+$의 곱셈적 부분집합이라 하자.
그러면 $T_0=T\cup\{1\}$은 $S$의 곱셈적 부분집합이다. $T_0$의
원소들은 동차이므로, 환 $T_0^{-1}S$에는 명백한 방식으로 여전히 등급이
주어진다. $S_{(T)}$로 $T_0^{-1}S$에서 차수가 $0$인 원소들, 즉
$x/h$ 꼴의 원소들로 이루어진 부분환을 나타내겠다. 여기서 $h\in T$이고
$x$는 $h$와 차수가 같은 동차 원소이다.
(\hyperref[0.1.4.5-ko]{0, 1.4.5})에 따라 $T_0^{-1}S$는 환
$S_f$들($f$가 $T$를 달리며 전이사상은 표준 준동형
$S_f\to S_{fg}$이다)의 귀납극한과 표준적으로 동일시됨이 알려져 있다.
이 동일시는 차수를 보존하므로 $S_{(T)}$를
$S_{(f)}$들의($f\in T$) \emph{귀납극한}과 동일시한다. 모든 등급
$S$-가군 $M$에 대하여도 마찬가지로 가군 $M_{(T)}$(환 $S_{(T)}$ 위의
가군)를 차수가 $0$인 $T_0^{-1}M$의 원소들로 이루어진 것으로 정의한다.
이 가군은 $M_{(f)}$들의($f\in T$) 귀납극한임을 알 수 있다.

$\mathfrak{p}$가 $S_+$의 동차 소아이디얼이면, 각각
$S_{(\mathfrak{p})}$와 $M_{(\mathfrak{p})}$로 환 $S_{(T)}$와 가군
$M_{(T)}$를 나타내겠다. 여기서 $T$는 $S_+$에서 $\mathfrak{p}$에
속하지 않는 \emph{동차} 원소들의 집합이다.
\end{env}

\subsection{등급환의 동차 소 스펙트럼.}
\label{subsection:II.2.3-ko}

\begin{env}[2.3.1]
\phantomsection
\label{II.2.3.1-ko}
음이 아닌 차수로 등급이 주어진 환 $S$에 대하여, $S$의
\emph{동차 소 스펙트럼}이라 하고 $\operatorname{Proj}(S)$로 나타내는 것은
$S_+$의 동차 소아이디얼들의 집합
(\hyperref[II.2.1.10-ko]{2.1.10}), 또는 같은 말로 $S$의 동차
소아이디얼 가운데 $S_+$를 \emph{포함하지 않는} 것들의 집합이다. 이제
$\operatorname{Proj}(S)$를 바탕집합으로 갖는 \emph{스킴} 구조를 정의하겠다.
\end{env}

\begin{env}[2.3.2]
\phantomsection
\label{II.2.3.2-ko}
$E$를 $S$의 임의의 부분집합이라 하고, $V_+(E)$를 $S$의 동차
소아이디얼 가운데 $E$를 포함하고 $S_+$를 포함하지 않는 것들의 집합이라
하자. 그러므로 이는 $V(E)\cap\operatorname{Proj}(S)$라는
$\operatorname{Spec}(S)$의 부분집합이다. 따라서
(\hyperref[I.1.1.2-ko]{I, 1.1.2})에서 다음을 얻는다.
\[
  \label{II.2.3.2.1-ko}
  V_+(0)=\operatorname{Proj}(S),\quad
  V_+(S)=V_+(S_+)=\varnothing
  \tag{2.3.2.1}
\]
\[
  \label{II.2.3.2.2-ko}
  V_+\left(\bigcup_\lambda E_\lambda\right)
  =\bigcap_\lambda V_+(E_\lambda)
  \tag{2.3.2.2}
\]
\[
  \label{II.2.3.2.3-ko}
  V_+(EE')=V_+(E)\cup V_+(E')
  \tag{2.3.2.3}
\]
$V_+(E)$는 집합 $E$를 $E$가 생성하는 동차 아이디얼로 바꾸어도 변하지
않는다. 또한 $\mathfrak{J}$가 $S$의 동차 아이디얼이면 다음이 성립한다.
\[
  \label{II.2.3.2.4-ko}
  V_+(\mathfrak{J})
  =V_+\left(\bigcup_{q\geq n}(\mathfrak{J}\cap S_q)\right)
  \tag{2.3.2.4}
\]
\oldpage[II]{26}
이는 모든 $n>0$에 대하여 성립한다. 실제로
$\mathfrak{p}\in\operatorname{Proj}(S)$가 $\mathfrak{J}$의 차수가
$\geq n$인 동차 원소들을 포함한다고 하자. 가정에 따라 동차 원소
$f\in S_d$ 가운데 $\mathfrak{p}$에 속하지 않는 것이 존재하므로, 모든
$m\geq 0$과 모든 $x\in S_m\cap\mathfrak{J}$에 대하여 관계
$f^r x\in\mathfrak{J}\cap S_{m+rd}$는 모든 음이 아닌 정수 $r$에 대해
성립한다. 한편 유한 개의 값을 제외한 모든 지수에 대해서는
$m+rd\geq n$이므로 위 교집합은 해당 소아이디얼에 포함되고, 따라서
$f^r x\in\mathfrak{p}\cap S_{m+rd}$이다. 앞서 고른 동차 원소가 해당
소아이디얼에 속하지 않으므로 소아이디얼의 성질로 그 거듭제곱을 소거하면
$x\in\mathfrak{p}\cap S_m$을 얻는다
(\hyperref[II.2.1.8-ko]{2.1.8}).

마지막으로 $\mathfrak{J}$가 $S$의 임의의 동차 아이디얼이면
\[
  \label{II.2.3.2.5-ko}
  V_+(\mathfrak{J})=V_+(r_+(\mathfrak{J})).
  \tag{2.3.2.5}
\]
\end{env}

\begin{env}[2.3.3]
\phantomsection
\label{II.2.3.3-ko}
정의에 따라 $V_+(E)$ 꼴의 집합들은 $X=\operatorname{Proj}(S)$ 위에서
$\operatorname{Spec}(S)$의 스펙트럼 위상으로부터 유도된 위상의
닫힌집합들이다. 이 위상도 $X$ 위의 \emph{스펙트럼 위상}이라 하겠다.
모든 $f\in S$에 대하여 다음과 같이 놓는다.
\[
  \label{II.2.3.3.1-ko}
  D_+(f)=D(f)\cap\operatorname{Proj}(S)
  =\operatorname{Proj}(S)\setminus V_+(f)
  \tag{2.3.3.1}
\]
따라서 $f$, $g$가 $S$의 두 원소이면
(\hyperref[I.1.1.9.1-ko]{I, 1.1.9.1})
\[
  \label{II.2.3.3.2-ko}
  D_+(fg)=D_+(f)\cap D_+(g).
  \tag{2.3.3.2}
\]
\end{env}

\begin{proposition}[2.3.4]
\phantomsection
\label{II.2.3.4-ko}
$f$가 $S_+$의 동차 원소 전체를 달릴 때, $D_+(f)$들은
$X=\operatorname{Proj}(S)$의 위상의 기저를 이룬다.
\end{proposition}

\begin{proof}
실제로 (\hyperref[II.2.3.2.2-ko]{2.3.2.2})와
(\hyperref[II.2.3.2.4-ko]{2.3.2.4})에 따라 $X$의 모든 닫힌 부분집합은
$V_+(f)$ 꼴 집합들의 교집합이다. 여기서 $f$는 차수가 $>0$인 동차
원소이다.
\end{proof}

\begin{env}[2.3.5]
\phantomsection
\label{II.2.3.5-ko}
$f$를 $S_+$의 차수가 $d>0$인 \emph{동차} 원소라 하자.
$\mathfrak{p}$가 $S$의 동차 소아이디얼이고 $f$를 포함하지 않을 때,
$x/f^n$ 꼴 원소
가운데 $x\in\mathfrak{p}$이고 $n\geq 0$인 것들의 집합은 분수환
$S_f$의 소아이디얼임이 알려져 있다
(\hyperref[0.1.2.6-ko]{0, 1.2.6}). 그 아이디얼과 $S_{(f)}$의 교집합도
이 환의 소아이디얼이므로 이를 $\psi_f(\mathfrak{p})$로 나타내겠다.
이는 $x/f^n$ 가운데 $n\geq 0$이고
$x\in\mathfrak{p}\cap S_{nd}$인 것들의 집합이다. 이로써 사상
\[
  \psi_f:D_+(f)\longrightarrow\operatorname{Spec}(S_{(f)})~;
\]
을 정의하였다. 또한 $g\in S_e$가 $S_+$의 또 하나의 동차 원소이면 다음
도식은 가환한다.
\[
  \label{II.2.3.5.1-ko}
  \xymatrix{
    D_+(f)\ar[r]^{\psi_f}
    & \operatorname{Spec}(S_{(f)})\\
    D_+(fg)\ar[u]\ar[r]_{\psi_{fg}}
    & \operatorname{Spec}(S_{(fg)})\ar[u]
  }
  \tag{2.3.5.1}
\]
여기서 왼쪽 수직 화살표는 포함 사상이고, 오른쪽 수직 화살표는 사상
${}^a\omega_{fg,f}$이다. 이는 표준 준동형
$\omega=\omega_{fg,f}:S_{(f)}\to S_{(fg)}$에서 유도된다
(\hyperref[I.1.2.1-ko]{I, 1.2.1}). 실제로
$x/f^n\in\omega^{-1}(\psi_{fg}(\mathfrak{p}))$이고
$fg\notin\mathfrak{p}$라 하자. 정의에 따라
$g^n x/(fg)^n\in\psi_{fg}(\mathfrak{p})$이므로
$g^n x\in\mathfrak{p}$이고, 따라서 $x\in\mathfrak{p}$이다. 역은
자명하다.
\end{env}

\begin{proposition}[2.3.6]
\phantomsection
\label{II.2.3.6-ko}
사상 $\psi_f$는 $D_+(f)$에서 $\operatorname{Spec}(S_{(f)})$로의
위상동형이다.
\end{proposition}

\begin{proof}
먼저 $\psi_f$는 연속이다. 실제로 $h\in S_{nd}$이고
$h/f^n\in\psi_f(\mathfrak{p})$이면 정의에 따라
$h\in\mathfrak{p}$이며, 그 역도 성립한다. 따라서
\[
  \psi_f^{-1}(D(h/f^n))=D_+(hf)
\]
이고, 주장은 공식
(\hyperref[II.2.3.3.2-ko]{2.3.3.2})에서 따른다. 또한 $D_+(hf)$들, 즉
$h$가 여러 집합 $S_{nd}$의 원소 전체를 달릴 때 얻는 이 집합들은,
(\hyperref[II.2.3.4-ko]{2.3.4})와 공식
(\hyperref[II.2.3.3.2-ko]{2.3.3.2})에 따라 $D_+(f)$의 위상의 기저를
이룬다.
\oldpage[II]{27}
그러므로 공리 $(T_0)$가 $D_+(f)$와 $\operatorname{Spec}(S_{(f)})$에서
성립함을 위의 논의와 함께 고려하면, $\psi_f$는 단사이고 그 역인 사상
$\psi_f(D_+(f))\longrightarrow D_+(f)$는 연속이다. 마지막으로
$\psi_f$가 전사임을 보이자. $\mathfrak{q}_0$를 $S_{(f)}$의
소아이디얼이라 하고, 모든 $n>0$에 대하여 $\mathfrak{p}_n$을
$x\in S_n$ 가운데 $x^d/f^n\in\mathfrak{q}_0$인 것들의 집합이라 하자.
그러면 $\mathfrak{p}_n$들은 (\hyperref[II.2.1.9-ko]{2.1.9})의 조건들을
만족한다. 실제로 $x\in S_n$, $y\in S_n$이고
$x^d/f^n\in\mathfrak{q}_0$, $y^d/f^n\in\mathfrak{q}_0$라 하자. 그러면
$(x+y)^{2d}/f^{2n}\in\mathfrak{q}_0$이고, 따라서
$(x+y)^d/f^n\in\mathfrak{q}_0$이다. 이는 $\mathfrak{q}_0$가
소아이디얼이기 때문이다. 이로써
$\mathfrak{p}_n$들이 $S_n$의 부분군임이 증명된다.
(\hyperref[II.2.1.9-ko]{2.1.9})의 다른 조건들도 $\mathfrak{q}_0$가
소아이디얼임을 고려하면 자명하게 확인된다. 이렇게 정의한 $\mathfrak{p}$는 $S$의
동차 소아이디얼이고 실제로
$\psi_f(\mathfrak{p})=\mathfrak{q}_0$이다. 왜냐하면 $x\in S_{nd}$일 때
두 관계 $x/f^n\in\mathfrak{q}_0$와
$x^d/f^{nd}\in\mathfrak{q}_0$는 $\mathfrak{q}_0$가 소아이디얼이므로
서로 동치이기 때문이다.
\end{proof}

\begin{corollary}[2.3.7]
\phantomsection
\label{II.2.3.7-ko}
$D_+(f)=\varnothing$이기 위한 필요충분조건은 $f$가 멱영원인 것이다.
\end{corollary}

\begin{proof}
실제로 $\operatorname{Spec}(S_{(f)})=\varnothing$이기 위한 필요충분조건은
$S_{(f)}=0$인 것이며, 이는 다시 $1=0$이 $S_f$에서 성립하는 것과 동치이다.
정의에 따라 이는 $f$가 멱영원이라는 뜻이다.
\end{proof}

\begin{corollary}[2.3.8]
\phantomsection
\label{II.2.3.8-ko}
$E$를 $S_+$의 부분집합이라 하자. 다음 조건들은 서로 동치이다.
\begin{enumerate}
  \item[\rm a)] $V_+(E)=X=\operatorname{Proj}(S)$.
  \item[\rm b)] $E$의 모든 원소는 멱영원이다.
  \item[\rm c)] $E$의 각 원소의 모든 동차 성분은 멱영원이다.
\end{enumerate}
\end{corollary}

\begin{proof}
c)가 b)를 함의하고 b)가 a)를 함의함은 명백하다. $\mathfrak{J}$가
$S$의 동차 아이디얼이고 $E$로 생성된다면, 조건 a)는
$V_+(\mathfrak{J})=X$와 동치이다. \emph{특히}, a)는 모든 동차 원소
$f\in\mathfrak{J}$가 $V_+(f)=X$를 만족함을
함의하므로, (\hyperref[II.2.3.7-ko]{2.3.7})에 따라 $f$는 멱영원이다.
\end{proof}

\begin{corollary}[2.3.9]
\phantomsection
\label{II.2.3.9-ko}
$\mathfrak{J}$가 $S_+$의 동차 아이디얼이면,
$r_+(\mathfrak{J})$는 $S_+$의 동차 소아이디얼들 가운데
$\mathfrak{J}$를 포함하는 것들의 교집합이다.
\end{corollary}

\begin{proof}
등급환 $S/\mathfrak{J}$를 고려하여 $\mathfrak{J}=0$인 경우로 환원할 수
있다. $f\in S_+$가 멱영원이 아니면 $S$의 동차 소아이디얼 가운데
$f$를 포함하지 않는 것이 존재함을 증명해야 한다. 그런데 $f$의 동차
성분 중 적어도 하나는 멱영원이 아니므로 $f$가 동차라고 가정할 수 있고,
이때 주장은 (\hyperref[II.2.3.7-ko]{2.3.7})에서 따른다.
\end{proof}

\begin{env}[2.3.10]
\phantomsection
\label{II.2.3.10-ko}
$Y$를 $X=\operatorname{Proj}(S)$의 임의의 부분집합이라 하자.
$\mathfrak{j}_+(Y)$를 모든 $f\in S_+$ 가운데 $Y\subset V_+(f)$를
만족하는 것들의 집합으로 둔다. 다시 말해,
$\mathfrak{j}_+(Y)=\mathfrak{j}(Y)\cap S_+$이다. 따라서
$\mathfrak{j}_+(Y)$는 $S_+$의 아이디얼이며, $S_+$에서 취한 그 근기와
같다.
\end{env}

\begin{proposition}[2.3.11]
\phantomsection
\label{II.2.3.11-ko}
\begin{enumerate}
  \item[\rm(i)] $E$를 $S_+$의 임의의 부분집합이라 하면,
    $\mathfrak{j}_+(V_+(E))$는 $S_+$에서 취한, $S_+$에서 $E$가 생성하는
    동차 아이디얼의 근기이다.
  \item[\rm(ii)] $Y$를 $X$의 임의의 부분집합이라 하면,
    $V_+(\mathfrak{j}_+(Y))=\overline{Y}$이다. 즉 우변은 $Y$의 폐포이며,
    이 폐포는 $X$ 안에서 취한다.
\end{enumerate}
\end{proposition}

\begin{proof}
\begin{enumerate}
  \item[\rm(i)] $\mathfrak{J}$가 $S_+$의 동차 아이디얼이고 $E$로
    생성된다면 $V_+(E)=V_+(\mathfrak{J})$이고, 이때 주장은
    (\hyperref[II.2.3.9-ko]{2.3.9})에서 따른다.
  \item[\rm(ii)] $V_+(\mathfrak{J})=\bigcap_{f\in\mathfrak{J}}V_+(f)$이므로,
    관계 $Y\subset V_+(\mathfrak{J})$는 $Y\subset V_+(f)$가 모든
    $f\in\mathfrak{J}$에 대하여 성립함을 함의한다. 따라서
    $\mathfrak{j}_+(Y)\supset\mathfrak{J}$이고, 이로부터
    $V_+(\mathfrak{j}_+(Y))\subset V_+(\mathfrak{J})$를 얻는다. 따라서
    닫힌 부분집합의 정의에 의해 (ii)가 증명된다.
\end{enumerate}
\end{proof}

\begin{corollary}[2.3.12]
\phantomsection
\label{II.2.3.12-ko}
$Y$가 $X=\operatorname{Proj}(S)$의 닫힌 부분집합일 때, 그러한 부분집합
전체와 $S_+$의 동차 아이디얼 가운데 $S_+$에서 취한 자신의 근기와 같은
것 전체는 포함관계를 역전하는 두 사상
$Y\longmapsto\mathfrak{j}_+(Y)$,
$\mathfrak{J}\longmapsto V_+(\mathfrak{J})$에 의해 서로 전단사로 대응한다.
또 $X$의 두 닫힌 부분집합의 합집합 $Y_1\cup Y_2$에는
\oldpage[II]{28}
$\mathfrak{j}_+(Y_1)\cap\mathfrak{j}_+(Y_2)$가 대응하고, 임의의 닫힌
부분집합족 $(Y_\lambda)$의 교집합에는 $S_+$에서
$\mathfrak{j}_+(Y_\lambda)$들의 합의 근기가 대응한다.
\end{corollary}

\begin{corollary}[2.3.13]
\phantomsection
\label{II.2.3.13-ko}
$\mathfrak{J}$를 $S_+$의 동차 아이디얼이라 하자.
$V_+(\mathfrak{J})=\varnothing$이기 위한 필요충분조건은 $S_+$의 각
원소마다 그 원소의 양의 정수 거듭제곱 중 하나가 $\mathfrak{J}$에 속하는
것이다.
\end{corollary}

마지막 따름정리는 다음과 같은 서로 동치인 형태들로도 표현된다.

\begin{corollary}[2.3.14]
\phantomsection
\label{II.2.3.14-ko}
$(f_\alpha)$를 $S_+$의 동차 원소족이라 하자. $D_+(f_\alpha)$들이
$X=\operatorname{Proj}(S)$의 덮개를 이루기 위한 필요충분조건은 $S_+$의
각 원소마다 그 원소의 양의 정수 거듭제곱 중 하나가 $f_\alpha$들이
생성하는 아이디얼에 속하는 것이다.
\end{corollary}

\begin{corollary}[2.3.15]
\phantomsection
\label{II.2.3.15-ko}
$(f_\alpha)$를 $S_+$의 동차 원소족, $f$를 $S_+$의 한 원소라 하자.
다음 관계들은 서로 동치이다.
\begin{enumerate}
  \item[\rm a)] $D_+(f)\subset\bigcup_\alpha D_+(f_\alpha)$.
  \item[\rm b)] $V_+(f)\supset\bigcap_\alpha V_+(f_\alpha)$.
  \item[\rm c)] $f$의 어떤 양의 정수 거듭제곱이 $f_\alpha$들이 생성하는
    아이디얼에 속한다.
\end{enumerate}
\end{corollary}

\begin{corollary}[2.3.16]
\phantomsection
\label{II.2.3.16-ko}
$X=\operatorname{Proj}(S)$가 공집합이기 위한 필요충분조건은 $S_+$의
모든 원소가 멱영원인 것이다.
\end{corollary}

\begin{corollary}[2.3.17]
\phantomsection
\label{II.2.3.17-ko}
(\hyperref[II.2.3.12-ko]{2.3.12})에서 서술한 전단사 대응에서 $X$의
기약 닫힌 부분집합들에는 $S_+$의 동차 소아이디얼들이 대응한다.
\end{corollary}

\begin{proof}
실제로 $Y=Y_1\cup Y_2$이고 $Y_1$, $Y_2$가 $Y$의 진부분집합인
닫힌집합들이면
\[
  \mathfrak{j}_+(Y)
  =\mathfrak{j}_+(Y_1)\cap\mathfrak{j}_+(Y_2)
\]
이며, 아이디얼 $\mathfrak{j}_+(Y_1)$과 $\mathfrak{j}_+(Y_2)$는 둘 다
$\mathfrak{j}_+(Y)$와 다르므로 $\mathfrak{j}_+(Y)$는 소아이디얼이
아니다. 역으로 $\mathfrak{J}$가 $S_+$의 소아이디얼이 아닌 동차
아이디얼이면 두 원소 $f$, $g$가 $S_+$에 존재하여
$f\notin\mathfrak{J}$, $g\notin\mathfrak{J}$, $fg\in\mathfrak{J}$를
만족한다. 그러면
$V_+(f)\not\supset V_+(\mathfrak{J})$,
$V_+(g)\not\supset V_+(\mathfrak{J})$이지만
(\hyperref[II.2.3.2.3-ko]{2.3.2.3})에 따라
$V_+(\mathfrak{J})\subset V_+(f)\cup V_+(g)$이다. 따라서
$V_+(\mathfrak{J})$는 닫힌집합
$V_+(f)\cap V_+(\mathfrak{J})$와
$V_+(g)\cap V_+(\mathfrak{J})$의 합집합이고, 이 두 닫힌집합은 모두
$V_+(\mathfrak{J})$와 다르다.
\end{proof}

\subsection{$\operatorname{Proj}(S)$의 스킴 구조}
\label{subsection:II.2.4-ko}

\begin{env}[2.4.1]
\phantomsection
\label{II.2.4.1-ko}
$f$, $g$를 $S_+$의 두 동차 원소라 하자. 아핀 스킴
$Y_f=\operatorname{Spec}(S_{(f)})$,
$Y_g=\operatorname{Spec}(S_{(g)})$,
$Y_{fg}=\operatorname{Spec}(S_{(fg)})$를 생각하자.
(\hyperref[II.2.2.2-ko]{2.2.2})에 따라, 사상
$w_{fg,f}=({}^a\omega_{fg,f},\widetilde{\omega}_{fg,f})$는 $Y_{fg}$에서
$Y_f$로 가며, 표준 준동형
$\omega_{fg,f}:S_{(f)}\to S_{(fg)}$에 대응하고,
\emph{열린 몰입 사상}이다
(\hyperref[I.1.3.6-ko]{I, 1.3.6}).
위상동형 $\psi_f:D_+(f)\to Y_f$의 역
(\hyperref[II.2.3.6-ko]{2.3.6})을 이용하여 $D_+(f)$에 $Y_f$의 아핀
스킴 구조를 옮길 수 있다. 도식
(\hyperref[II.2.3.5.1-ko]{2.3.5.1})의 가환성에 따라, 이렇게 얻은 아핀
스킴 $D_+(fg)$는 아핀 스킴 $D_+(f)$의 바탕공간의 열린 부분집합
$D_+(fg)$에 유도된 스킴과 동일시된다. 그러면
(\hyperref[II.2.3.4-ko]{2.3.4})를 고려할 때
$X=\operatorname{Proj}(S)$에는 유일한 \emph{준스킴} 구조가 주어지며,
그 구조를 각 $D_+(f)$에 제한하면 방금 정의한 아핀 스킴이 된다. 또한
다음이 성립한다.
\end{env}

\begin{proposition}[2.4.2]
\phantomsection
\label{II.2.4.2-ko}
준스킴 $\operatorname{Proj}(S)$는 스킴이다.
\end{proposition}

\begin{proof}
(\hyperref[I.5.5.6-ko]{I, 5.5.6})에 따라, $f$, $g$가 $S_+$의 임의의
동차 원소일 때 $D_+(f)\cap D_+(g)=D_+(fg)$가 아핀이고 그 환이
$D_+(f)$와 $D_+(g)$의 환들의 표준 상들에 의해 생성됨을 보이면
충분하다. 첫째 주장은 정의상 명백하고, 둘째 주장은
(\hyperref[II.2.2.4-ko]{2.2.4})에서 따른다.
\end{proof}
\oldpage[II]{29}
$\operatorname{Proj}(S)$라는 동차 소 스펙트럼을 \emph{스킴}이라고 할
때에는 언제나 방금 정의한 구조를 뜻한다.

\begin{example}[2.4.3]
\phantomsection
\label{II.2.4.3-ko}
$S=K[T_1,T_2]$라 하자. 여기서 $K$는 체이고, $T_1$, $T_2$는 두
부정원이며, $S$에는 전체 차수에 따른 등급이 주어져 있다.
(\hyperref[II.2.3.14-ko]{2.3.14})에 따라
$\operatorname{Proj}(S)$는 $D_+(T_1)$과 $D_+(T_2)$의 합집합이다.
이 두 아핀 스킴은 각각 $\operatorname{Spec}(K[T])$와 표준적으로
동형이며, $\operatorname{Proj}(S)$는
(\hyperref[I.2.3.2-ko]{I, 2.3.2})에서 서술한 대로 이 두 아핀 스킴을
붙여서 얻어진다((7.4.14) 참조).
\end{example}

\begin{proposition}[2.4.4]
\phantomsection
\label{II.2.4.4-ko}
$S$를 음이 아닌 차수로 등급이 주어진 환이라 하고, $X$를 스킴
$\operatorname{Proj}(S)$라 하자.
\begin{enumerate}
  \item[\rm(i)] $\mathfrak{N}_+$가 $S_+$의 멱영근기이면
    (\hyperref[II.2.1.10-ko]{2.1.10}), 스킴 $X_{\mathrm{red}}$는
    $\operatorname{Proj}(S/\mathfrak{N}_+)$와 표준적으로 동형이다. 특히
    $S$가 본질적으로 축소이면 $\operatorname{Proj}(S)$는 축소이다.
  \item[\rm(ii)] $S$가 본질적으로 축소라고 가정하자. 이때 $X$가 정역
    스킴이기 위한 필요충분조건은 $S$가 본질적으로 정역인 것이다.
\end{enumerate}
\end{proposition}

\begin{proof}
\begin{enumerate}
  \item[\rm(i)] $\overline{S}$를 등급환 $S/\mathfrak{N}_+$라 하고, 차수가
    $0$인 표준 준동형 $S\to\overline{S}$를
    $x\mapsto\overline{x}$로 나타내자. 모든 $f\in S_d$ ($d>0$)에 대하여
    표준 준동형
    $S_f\to\overline{S}_{\overline{f}}$
    (\hyperref[0.1.5.1-ko]{0, 1.5.1})는 전사이고 차수가 $0$이므로, 이를
    제한하면 전사 준동형
    $S_{(f)}\to\overline{S}_{(\overline{f})}$를 얻는다.
    $f\notin\mathfrak{N}_+$라고 가정하면
    $\overline{S}_{(\overline{f})}$가 축소환이고 앞의 준동형의 핵이
    $S_{(f)}$의 멱영근기임을 바로 확인할 수 있다. 다시 말해
    $\overline{S}_{(\overline{f})}=(S_{(f)})_{\mathrm{red}}$이다. 따라서 이
    준동형에는 $D_+(\overline{f})$를 $(D_+(f))_{\mathrm{red}}$와 동일시하는
    닫힌 몰입 사상
    $D_+(\overline{f})\to D_+(f)$가 대응한다
    (\hyperref[I.5.1.2-ko]{I, 5.1.2}). 특히 이 사상은 두 아핀 스킴의
    바탕공간 사이의 위상동형이다. 더욱이
    $g\notin\mathfrak{N}_+$가 $S_+$의 또 다른 동차 원소이면 도식
    \[
      \xymatrix{
        S_{(f)}\ar[r]\ar[d]
        & \overline{S}_{(\overline{f})}\ar[d]\\
        S_{(fg)}\ar[r]
        & \overline{S}_{(\overline{fg})}
      }
    \]
    은 가환한다. 한편 $D_+(f)$들은 $f$가 차수가 $>0$인 동차 원소이고
    $f\notin\mathfrak{N}_+$일 때
    $X=\operatorname{Proj}(S)$의 덮개를 이룬다
    (\hyperref[II.2.3.7-ko]{2.3.7}). 따라서 사상
    $D_+(\overline{f})\to D_+(f)$들은 닫힌 몰입 사상
    $\operatorname{Proj}(\overline{S})\to\operatorname{Proj}(S)$의
    제한들이고, 이 닫힌 몰입 사상은 바탕공간 사이의 위상동형이다. 이제
    (\hyperref[I.5.1.2-ko]{I, 5.1.2})에 의해 결론이 따른다.
  \item[\rm(ii)] $S$가 본질적으로 정역이라고 가정하자. 다시 말해
    $(0)$은 $S_+$와 다른 $S_+$의 등급 소아이디얼이다. 그러면 (i)에 따라
    $X$는 축소이고, (\hyperref[II.2.3.17-ko]{2.3.17})에 따라 기약이다.
    반대로 $S$가 본질적으로 축소이고 $X$가 정역 스킴이라고 가정하자.
    그러면 $S_+$의 동차 원소 $f\neq0$에 대하여
    $D_+(f)\neq\varnothing$이다
    (\hyperref[II.2.3.7-ko]{2.3.7}). $X$가 기약이라는 가정에 따라
    $S_+$의 동차 원소 $f$, $g$가 모두 $\neq0$이면
    $D_+(f)\cap D_+(g)\neq\varnothing$이다. 따라서
    (\hyperref[II.2.3.3.2-ko]{2.3.3.2})에 의해 $fg\neq0$이고, 이로부터
    $S_+$에는 영인자가 없음을 얻으므로 처음의 주장이 따른다.
\end{enumerate}
\end{proof}

\begin{env}[2.4.5]
\phantomsection
\label{II.2.4.5-ko}
$A$를 가환환이라 하자. 등급환 $S$에 각 부분군 $S_n$이 $A$-가군이
되도록 하는 $A$-대수 구조가 주어져 있으면 이를
\emph{등급 $A$-대수}라고 한다는 것을 상기하자. 이를 위해서는 $S_0$가
\oldpage[II]{30}
$A$-대수이면 충분하다. 다시 말해 $S_0$에 $A$-대수 구조를 정의하고
$\alpha\in A$, $x\in S_n$에 대하여
$\alpha\cdot x=(\alpha\cdot1)x$로 놓으면 $S$에 등급 $A$-대수 구조가
정의된다.
\end{env}

\begin{proposition}[2.4.6]
\phantomsection
\label{II.2.4.6-ko}
$S$가 등급 $A$-대수라고 가정하자. 그러면
$X=\operatorname{Proj}(S)$ 위에서 구조층 $\mathcal{O}_X$는
$A$-대수들의 층이다($A$는 $X$ 위의 상수층으로 본다). 다시 말해 $X$는
$\operatorname{Spec}(A)$ 위의 스킴이다.
\end{proposition}

\begin{proof}
$S_+$의 모든 동차 원소 $f$에 대하여 $S_{(f)}$가 $A$ 위의 대수이고,
$S_+$의 동차 원소 $f$, $g$에 대하여 도식
\[
  \xymatrix{
    S_{(f)}\ar[rr] && S_{(fg)}\\
    & A\ar[ul]\ar[ur] &
  }
\]
이 가환함에 유의하면 충분하다.
\end{proof}

\begin{proposition}[2.4.7]
\phantomsection
\label{II.2.4.7-ko}
$S$를 음이 아닌 차수로 등급이 주어진 환이라 하자.
\begin{enumerate}
  \item[\rm(i)] 모든 정수 $d>0$에 대하여 스킴
    $\operatorname{Proj}(S)$에서 스킴
    $\operatorname{Proj}(S^{(d)})$로 가는 표준 동형이 존재한다.
  \item[\rm(ii)] $S'$를 다음 조건을 만족하는 등급환이라 하자:
    $S'_0=\mathbf{Z}$이고, $S'_n=S_n$은 ($\mathbf{Z}$-가군으로서)
    $n>0$일 때 성립한다. 그러면 스킴 $\operatorname{Proj}(S)$에서 스킴
    $\operatorname{Proj}(S')$로 가는 표준 동형이 존재한다.
\end{enumerate}
\end{proposition}

\begin{proof}
\begin{enumerate}
  \item[\rm(i)] 먼저 사상
    $\mathfrak{p}\mapsto\mathfrak{p}\cap S^{(d)}$가 집합
    $\operatorname{Proj}(S)$에서
    $\operatorname{Proj}(S^{(d)})$로 가는 전단사임을 보이자. 실제로 동차
    소아이디얼 $\mathfrak{p}'\in\operatorname{Proj}(S^{(d)})$가 주어졌다고
    하고 $\mathfrak{p}_{nd}=\mathfrak{p}'\cap S_{nd}$ ($n\geq0$)로 놓자.
    $n>0$이고 $d$의 배수가 아닌 모든 경우에 $\mathfrak{p}_n$을
    $x\in S_n$이고 $x^d\in\mathfrak{p}_{nd}$인 원소들의 집합으로 정의하자.
    $x\in\mathfrak{p}_n$, $y\in\mathfrak{p}_n$이면
    $(x+y)^{2d}\in\mathfrak{p}_{2nd}$이고, 따라서
    $(x+y)^d\in\mathfrak{p}_{nd}$이다. 이는 $\mathfrak{p}'$이
    소아이디얼이기 때문이다. 이와 같이 정의된 $\mathfrak{p}_n$들이 모든
    $n\geq0$에 대하여
    (\hyperref[II.2.1.9-ko]{2.1.9})의 조건을 만족함은 자명하다. 따라서
    유일한 소아이디얼 $\mathfrak{p}\in\operatorname{Proj}(S)$가 존재하며,
    이 소아이디얼은 $\mathfrak{p}\cap S^{(d)}=\mathfrak{p}'$을 만족한다. 한편 $f$가
    $S_+$의 동차 원소일 때마다
    $V_+(f)=V_+(f^d)$이므로
    (\hyperref[II.2.3.2.3-ko]{2.3.2.3}), 앞의 전단사는 위상공간들의
    위상동형이다. 마지막으로 같은 표기 아래에서 $S_{(f)}$와
    $S^{(d)}_{(f^d)}$는 표준적으로 동일시되므로
    (\hyperref[II.2.2.2-ko]{2.2.2}), $\operatorname{Proj}(S)$와
    $\operatorname{Proj}(S^{(d)})$도 스킴으로서 표준적으로 동일시된다.
  \item[\rm(ii)] 각 $\mathfrak{p}\in\operatorname{Proj}(S)$에 유일한
    소아이디얼 $\mathfrak{p}'\in\operatorname{Proj}(S')$을 대응시키되,
    $\mathfrak{p}'\cap S_n=\mathfrak{p}\cap S_n$이 모든 $n>0$에 대하여
    성립하게 하면
    (\hyperref[II.2.1.9-ko]{2.1.9}), 바탕공간들의 표준 위상동형
    $\operatorname{Proj}(S)\xrightarrow{\sim}\operatorname{Proj}(S')$을
    정의함이 분명하다. 실제로 $V_+(f)$는 $S$에 대해서나 $S'$에 대해서나
    같은 집합이다($f$가 $S_+$의 동차 원소일 때). 더욱이
    $S_{(f)}=S'_{(f)}$이므로 $\operatorname{Proj}(S)$와
    $\operatorname{Proj}(S')$은 스킴으로서 동일시된다.
\end{enumerate}
\end{proof}

\begin{corollary}[2.4.8]
\phantomsection
\label{II.2.4.8-ko}
$S$가 등급 $A$-대수이고, $S'_A$가 다음 조건을 만족하는 등급
$A$-대수라고 하자: $(S'_A)_0=A$, $(S'_A)_n=S_n$ ($n>0$).
$\operatorname{Proj}(S)$에서 $\operatorname{Proj}(S'_A)$로 가는 표준
동형이 존재한다.
\end{corollary}

\begin{proof}
실제로 (\hyperref[II.2.4.7-ko]{2.4.7, (ii)})의 표기 아래에서 이 두
스킴은 모두 $\operatorname{Proj}(S')$와 표준적으로 동형이다.
\end{proof}

\subsection{등급가군에 대응하는 층.}
\label{subsection:II.2.5-ko}

\begin{env}[2.5.1]
\phantomsection
\label{II.2.5.1-ko}
$M$을 등급 $S$-가군이라 하자. $f$가 $S_+$의 동차 원소이면
$M_{(f)}$는 $S_{(f)}$-가군이고, 따라서 이 가군에 대응하는 준연접층
$\widetilde{M_{(f)}}$가 아핀 스킴
$\operatorname{Spec}(S_{(f)})$ 위에 있다. 이 스킴은 $D_+(f)$와
동일시된다(\hyperref[I.1.3.4-ko]{I, 1.3.4}).
\end{env}
\oldpage[II]{31}

\begin{proposition}[2.5.2]
\phantomsection
\label{II.2.5.2-ko}
$X=\operatorname{Proj}(S)$ 위에 다음 조건을 만족하는 준연접
$\mathcal{O}_X$-가군 $\widetilde{M}$이 유일하게 존재한다. 즉 $f$가
$S_+$의 동차 원소일 때마다
$\Gamma(D_+(f),\widetilde{M})=M_{(f)}$이고, 제한 준동형
$\Gamma(D_+(f),\widetilde{M})\to
\Gamma(D_+(fg),\widetilde{M})$은 $f$, $g$가 $S_+$의 동차 원소일 때
표준 준동형 $M_{(f)}\to M_{(fg)}$이다
(\hyperref[II.2.2.3-ko]{2.2.3}).
\end{proposition}

\begin{proof}
$f\in S_d$, $g\in S_e$라 하자.
(\hyperref[II.2.2.2-ko]{2.2.2})에 따라 $D_+(fg)$는
$(S_{(f)})_{g^d/f^e}$의 소 스펙트럼과 동일시된다. 따라서
$D_+(fg)$로 제한한 층 $\widetilde{M_{(f)}}$---여기서 제한 전의 층은
$D_+(f)$ 위에 있다---은 가군 $(M_{(f)})_{g^d/f^e}$에 대응하는 층과 표준적으로
동일시되고(\hyperref[I.1.3.6-ko]{I, 1.3.6}), 따라서
$\widetilde{M_{(fg)}}$와도 동일시된다
(\hyperref[II.2.2.2-ko]{2.2.2}). 그러므로 표준 동형
\[
  \theta_{g,f}:\widetilde{M_{(f)}}|D_+(fg)
  \xrightarrow{\sim}\widetilde{M_{(g)}}|D_+(fg)
\]
이 존재하며, $h$가 $S_+$의 세 번째 동차 원소이면
$\theta_{f,h}=\theta_{f,g}\circ\theta_{g,h}$가 $D_+(fgh)$에서 성립한다.
따라서
(\hyperref[0.3.3.1-ko]{0, 3.3.1}) $X$ 위에 준연접
$\mathcal{O}_X$-가군 $\mathcal{F}$가 존재하고, 각 $f$가 $S_+$의 동차
원소일 때 동형 $\eta_f$가 $\mathcal{F}|D_+(f)$에서
$\widetilde{M_{(f)}}$로 가며 $\theta_{g,f}=\eta_g\circ\eta_f^{-1}$을
만족한다.
이제 층 $\mathcal{G}$를 생각하자. 이는 $X$의 위상의 기저인
$D_+(f)$들 위에서 $D_+(f)\to M_{(f)}$로 정의되고, 표준 준동형
$M_{(f)}\to M_{(fg)}$를 제한 준동형으로 갖는 준층을 층화하여 얻는
층이다. 앞의 논의로부터 $\mathcal{F}$와 $\mathcal{G}$가 동형임을 안다
((\hyperref[I.1.3.7-ko]{I, 1.3.7})도 고려한다). 층 $\mathcal{G}$를
$\widetilde{M}$으로 나타내면 이 층은 명제의 조건들을 만족한다. 특히
$\widetilde{S}=\mathcal{O}_X$이다.
\end{proof}

\begin{definition}[2.5.3]
\phantomsection
\label{II.2.5.3-ko}
(\hyperref[II.2.5.2-ko]{2.5.2})에서 정의한 준연접
$\mathcal{O}_X$-가군 $\widetilde{M}$을 등급 $S$-가군 $M$에
\emph{대응하는} 층이라 한다.
\end{definition}

등급 $S$-가군들은 등급가군의 준동형을 \emph{차수 $0$의} 준동형으로
제한하면 하나의 범주를 이룸을 상기하자. 이 약속 아래에서 다음이
성립한다.

\begin{proposition}[2.5.4]
\phantomsection
\label{II.2.5.4-ko}
함자 $M\mapsto\widetilde{M}$은 등급 $S$-가군의 범주에서 준연접
$\mathcal{O}_X$-가군의 범주로 가는 공변 가법 완전 함자이고,
귀납극한 및 직합과 가환한다.
\end{proposition}

\begin{proof}
실제로 이 성질들은 국소적이므로 층
$\widetilde{M}|D_+(f)=\widetilde{M_{(f)}}$에 대하여 확인하면 충분하다.
그런데 함자 $M\mapsto M_f$, 함자 $N\mapsto N_0$ (등급
$S_f$-가군의 범주에서) 및 함자 $P\mapsto\widetilde{P}$
($S_{(f)}$-가군의 범주에서)는 모두 완전성과 귀납극한 및 직합과의
가환성을 갖는다(\hyperref[I.1.3.5-ko]{I, 1.3.5} 및
\hyperref[I.1.3.9-ko]{I, 1.3.9}). 이로부터 명제가 따른다.
\end{proof}

$\widetilde{u}$로 준동형 $\widetilde{M}\to\widetilde{N}$을 나타내기로
한다. 이는 차수 $0$의 준동형 $u:M\to N$에 대응하는 준동형이다.
(\hyperref[II.2.5.4-ko]{2.5.4})에서 곧바로, 등급 $S$-가군과 차수
$0$의 준동형에 대해서도 (여기서 정한 $\widetilde{M}$의 뜻으로)
(\hyperref[I.1.3.9-ko]{I, 1.3.9} 및
\hyperref[I.1.3.10-ko]{I, 1.3.10})의 결과들이 여전히 성립함을 알 수
있다. 실제로 그 증명들은 순전히 형식적이다.

\begin{proposition}[2.5.5]
\phantomsection
\label{II.2.5.5-ko}
모든 $\mathfrak{p}\in X=\operatorname{Proj}(S)$에 대하여
$\widetilde{M}_{\mathfrak{p}}=M_{(\mathfrak{p})}$이다.
\end{proposition}

\begin{proof}
실제로 정의에 따라
$\widetilde{M}_{\mathfrak{p}}=
\varinjlim\Gamma(D_+(f),\widetilde{M})$이다. 여기서 $f$의 범위는
$S_+$의 동차 원소 가운데 $f\notin\mathfrak{p}$인 원소 전체이다. 한편
$\Gamma(D_+(f),\widetilde{M})=M_{(f)}$이므로, 명제는
$M_{(\mathfrak{p})}$의 정의
(\hyperref[II.2.2.7-ko]{2.2.7})에서 따른다.
\end{proof}
\oldpage[II]{32}

특히 \emph{국소환 $(\mathcal{O}_X)_\mathfrak{p}$}은 바로 환
$S_{(\mathfrak{p})}$이다. 이는 분수 $x/f$들의 집합으로, 여기서 $f$는
$S_+$의 동차 원소이고 $\mathfrak{p}$에 속하지 않으며, $x$는 $f$와
같은 차수의 동차 원소이다.

더욱 구체적으로, $S$가 \emph{본질적으로 정역}이어서
$\operatorname{Proj}(S)=X$가 \emph{정역 스킴}이고
(\hyperref[II.2.4.4-ko]{2.4.4}), $\xi=(0)$가 $X$의 일반점이면,
\emph{유리함수체} $R(X)=\mathcal{O}_\xi$는 $fg^{-1}$ 꼴의 원소들로 이루어진
체이다. 여기서 $f$와 $g$는 $S_+$에 속하는 같은 차수의 동차 원소이고
$g\neq0$이다.

\begin{proposition}[2.5.6]
\phantomsection
\label{II.2.5.6-ko}
모든 $z\in M$에 대하여, $f$가 $S_+$의 동차 원소일 때마다 $f$의 어떤
거듭제곱이 $z$를 소멸시키면, $\widetilde{M}=0$이다. 이 충분조건은
$S_0$-대수 $S$가 집합 $S_1$에 의해 생성되고, 이 집합이 차수 $1$의
동차 원소 전체로 이루어질 때에는 필요조건이기도 하다.
\end{proposition}

\begin{proof}
실제로 조건 $\widetilde{M}=0$은 $M_{(f)}=0$이 $f$가 $S_+$의 동차
원소일 때마다 성립하는 것과 동치이다. 한편 $f\in S_d$일 때
$M_{(f)}=0$이라는 것은 모든 동차 $z\in M$ 가운데 차수가 $d$의 배수인
각 원소에 대하여 어떤 거듭제곱 $f^n$이 존재하여 $f^nz=0$임을 뜻한다.
$M_{(f)}=0$이 모든 $f\in S_1$에 대해 성립하면 명제의 조건은 모든
$f\in S_1$에 대하여 성립한다. 또한 앞의 생성 가정 아래에서는 $f$가
$S_+$의 동차 원소일 때마다 이 조건이 \emph{더 나아가}
성립한다. 이 생성 가정은 $S_1$이 $S$를 생성한다는 것이며, 실제로
$S_+$의 모든 동차 원소는 $S_1$의 원소들의 곱들의 선형결합이기 때문이다.
\end{proof}

\begin{proposition}[2.5.7]
\phantomsection
\label{II.2.5.7-ko}
$d$를 $>0$인 정수라 하고 $f\in S_d$라 하자. 그러면 모든
$n\in\mathbf{Z}$에 대하여 $(\mathcal{O}_X|D_+(f))$-가군
$\widetilde{S(nd)}|D_+(f)$는 $\mathcal{O}_X|D_+(f)$와 표준적으로
동형이다.
\end{proposition}

\begin{proof}
실제로 $(f/1)^n$은 $S_f$의 가역원이며, 이를 곱하는 사상은
$S_{(f)}=(S_f)_0$에서
$(S_f)_{nd}=(S_f(nd))_0=(S(nd)_f)_0=S(nd)_{(f)}$로 가는 전단사이다.
다시 말해 $S_{(f)}$-가군 $S_{(f)}$와 $S(nd)_{(f)}$는 표준적으로
동형이고, 따라서 명제가 성립한다.
\end{proof}

\begin{corollary}[2.5.8]
\phantomsection
\label{II.2.5.8-ko}
열린부분집합
\[
  U=\bigcup_{f\in S_d}D_+(f),
\]
위에서 $\mathcal{O}_X$-가군 $\widetilde{S(nd)}$의 제한은 가역인
$(\mathcal{O}_X|U)$-가군이다
(\hyperref[0.5.4.1-ko]{0, 5.4.1}).
\end{corollary}

\begin{corollary}[2.5.9]
\phantomsection
\label{II.2.5.9-ko}
아이디얼 $S_+$가 $S$에서 집합 $S_1$에 의해 생성되고, 이 집합이 차수
$1$의 동차 원소 전체로 이루어지면, $\mathcal{O}_X$-가군
$\widetilde{S(n)}$은 모든 $n\in\mathbf{Z}$에 대하여 가역이다.
\end{corollary}

\begin{proof}
$X=\bigcup_{f\in S_1}D_+(f)$임은 가정과
(\hyperref[II.2.3.14-ko]{2.3.14})에서 따른다. 이제 $U=X$로 두고
(\hyperref[II.2.5.8-ko]{2.5.8})을 적용하면 충분하다.
\end{proof}

\begin{env}[2.5.10]
\phantomsection
\label{II.2.5.10-ko}
이 절의 나머지 부분 전체에서
\[
\label{II.2.5.10.1-ko}
  \mathcal{O}_X(n)=\widetilde{S(n)}
\tag{2.5.10.1}
\]
로 두기로 한다. 이는 모든 $n\in\mathbf{Z}$에 대한 약속이다. 또한
$U$를 $X$의 임의의 열린부분집합이라 하고, 임의의
$(\mathcal{O}_X|U)$-가군 $\mathcal{F}$에 대하여
\[
\label{II.2.5.10.2-ko}
  \mathcal{F}(n)=
  \mathcal{F}\otimes_{\mathcal{O}_X|U}(\mathcal{O}_X(n)|U)
\tag{2.5.10.2}
\]
로 두기로 한다. 이는 모든 $n\in\mathbf{Z}$에 대한 약속이다. 아이디얼
$S_+$가 $S_1$에 의해 생성되면, $\mathcal{F}(n)$을 대응시키는 함자는
$\mathcal{F}$에 관해 모든 $n\in\mathbf{Z}$에 대하여 \emph{완전}하다. 이 경우
$\mathcal{O}_X(n)$이 가역인 $\mathcal{O}_X$-가군이기 때문이다.
\end{env}

\begin{env}[2.5.11]
\phantomsection
\label{II.2.5.11-ko}
$M$과 $N$을 두 등급 $S$-가군이라 하자. 모든 $f\in S_d$ ($d>0$)에
대하여 다음과 같은 $S_{(f)}$-가군의 함자적인 표준 준동형을 정의한다.
\[
\label{II.2.5.11.1-ko}
  \lambda_f:M_{(f)}\otimes_{S_{(f)}}N_{(f)}
  \longrightarrow(M\otimes_S N)_{(f)}
\tag{2.5.11.1}
\]
\oldpage[II]{33}

정의는 준동형
$M_{(f)}\otimes_{S_{(f)}}N_{(f)}\to M_f\otimes_{S_f}N_f$
---이는 단사 준동형들 $M_{(f)}\to M_f$, $N_{(f)}\to N_f$ 및
$S_{(f)}\to S_f$에서 유도된다---과 표준 동형
$M_f\otimes_{S_f}N_f\xrightarrow{\sim}(M\otimes_S N)_f$
(\hyperref[0.1.3.4-ko]{0, 1.3.4})을 합성하여 얻는다. 또한 두 등급가군의
텐서곱의 정의에 따라 뒤의 동형이 차수를 보존한다는 점을 사용한다. 따라서
$x\in M_{md}$, $y\in N_{nd}$ ($m\geq0$, $n\geq0$)에 대하여
\[
  \lambda_f((x/f^m)\otimes(y/f^n))=(x\otimes y)/f^{m+n}.
\]

이 정의에서 곧바로, $g\in S_e$ ($e>0$)이면 다음 도식이 가환함을 알 수
있다.
\[
  \xymatrix{
    M_{(f)}\otimes_{S_{(f)}}N_{(f)}\ar[r]^{\lambda_f}\ar[d]
    & (M\otimes_S N)_{(f)}\ar[d]\\
    M_{(fg)}\otimes_{S_{(fg)}}N_{(fg)}\ar[r]_{\lambda_{fg}}
    & (M\otimes_S N)_{(fg)}
  }
\]
여기서 오른쪽 세로 화살표는 표준 준동형이고 왼쪽 세로 화살표는 표준
준동형들에서 유도된다. 따라서 이들 $\lambda$로부터
$\mathcal{O}_X$-가군의 함자적인 표준 준동형
\[
\label{II.2.5.11.2-ko}
  \lambda:\widetilde{M}\otimes_{\mathcal{O}_X}\widetilde{N}
  \longrightarrow\widetilde{M\otimes_S N}.
\tag{2.5.11.2}
\]
을 얻는다.

특히 두 동차 아이디얼 $\mathfrak{J}$, $\mathfrak{K}$를 생각하자. 이들은
$S$의 아이디얼이다. $\widetilde{\mathfrak{J}}$와
$\widetilde{\mathfrak{K}}$는 $\mathcal{O}_X$의 아이디얼층이므로 표준 준동형
$\widetilde{\mathfrak{J}}\otimes_{\mathcal{O}_X}
\widetilde{\mathfrak{K}}\to\mathcal{O}_X$가 존재하며, 다음 도식은 가환한다.
\[
\label{II.2.5.11.3-ko}
  \xymatrix{
    \widetilde{\mathfrak{J}}\otimes_{\mathcal{O}_X}
    \widetilde{\mathfrak{K}}\ar[rr]^\lambda\ar[dr]
    && \widetilde{\mathfrak{J}\otimes_S\mathfrak{K}}\ar[dl]\\
    & \mathcal{O}_X &
  }
\tag{2.5.11.3}
\]
실제로 각 열린집합 $D_+(f)$ ($f$는 $S_+$의 동차 원소)에서 이를 확인하는
경우로 귀착할 수 있고, 이는 $\lambda_f$의 정의
(\hyperref[II.2.5.11.1-ko]{2.5.11.1})와
(\hyperref[I.1.3.13-ko]{I, 1.3.13})에서 곧바로 따른다.

끝으로 $M$, $N$, $P$를 세 등급 $S$-가군이라 하면 다음 도식을 얻는다.
\[
\label{II.2.5.11.4-ko}
  \xymatrix{
    \widetilde{M}\otimes_{\mathcal{O}_X}\widetilde{N}
    \otimes_{\mathcal{O}_X}\widetilde{P}
    \ar[r]^{\lambda\otimes1}\ar[d]_{1\otimes\lambda}
    & \widetilde{M\otimes_S N}\otimes_{\mathcal{O}_X}\widetilde{P}
    \ar[d]^\lambda\\
    \widetilde{M}\otimes_{\mathcal{O}_X}\widetilde{N\otimes_S P}
    \ar[r]_\lambda
    & \widetilde{M\otimes_S N\otimes_S P}
  }
\tag{2.5.11.4}
\]
\oldpage[II]{34}

이 도식은 가환한다. 다시 각 $D_+(f)$에서 확인하면 충분하며, 이는 정의와
(\hyperref[I.1.3.13-ko]{I, 1.3.13})에서 곧바로 따른다.
\end{env}

\begin{env}[2.5.12]
\phantomsection
\label{II.2.5.12-ko}
(\hyperref[II.2.5.11-ko]{2.5.11})의 가정 아래에서 다음과 같은
$S_{(f)}$-가군의 함자적인 표준 준동형을 정의한다.
\[
\label{II.2.5.12.1-ko}
  \mu_f:\bigl(\operatorname{Hom}_S(M,N)\bigr)_{(f)}
  \longrightarrow
  \operatorname{Hom}_{S_{(f)}}(M_{(f)},N_{(f)})
\tag{2.5.12.1}
\]
이는 $u/f^n$---여기서 $u$는 차수 $nd$인 준동형이다---에 준동형
$M_{(f)}\to N_{(f)}$을 대응시켜 정의한다. 이 준동형은 $x/f^m$
($x\in M_{md}$)을 $u(x)/f^{m+n}$으로 보낸다. $g\in S_e$ ($e>0$)에
대해서도 다음 도식은 가환한다.
\[
  \xymatrix{
    \bigl(\operatorname{Hom}_S(M,N)\bigr)_{(f)}
    \ar[r]^{\mu_f}\ar[d]
    & \operatorname{Hom}_{S_{(f)}}(M_{(f)},N_{(f)})\ar[d]\\
    \bigl(\operatorname{Hom}_S(M,N)\bigr)_{(fg)}
    \ar[r]_{\mu_{fg}}
    & \operatorname{Hom}_{S_{(fg)}}(M_{(fg)},N_{(fg)})
  }
\]
여기서 왼쪽 화살표는 표준 준동형이고 오른쪽 화살표는 표준 준동형들에서
유도된다. 따라서 (\hyperref[I.1.3.8-ko]{I, 1.3.8})을 고려하면
$\mu_f$들이 다음과 같은 $\mathcal{O}_X$-가군의 함자적인 표준 준동형을
정의함을 다시 얻는다.
\[
\label{II.2.5.12.2-ko}
  \mu:\widetilde{\operatorname{Hom}_S(M,N)}
  \longrightarrow
  \mathcal{H}om_{\mathcal{O}_X}(\widetilde{M},\widetilde{N}).
\tag{2.5.12.2}
\]
\end{env}

\begin{proposition}[2.5.13]
\phantomsection
\label{II.2.5.13-ko}
아이디얼 $S_+$가 $S_1$에 의해 생성된다고 가정하자. 그러면 준동형
$\lambda$(\hyperref[II.2.5.11.2-ko]{2.5.11.2})는 동형사상이고, 준동형
$\mu$(\hyperref[II.2.5.12.2-ko]{2.5.12.2})도 등급 $S$-가군 $M$이
유한 표시를 가지는 경우(\hyperref[II.2.1.1-ko]{2.1.1})에는 동형사상이다.
\end{proposition}

\begin{proof}
$X$가 $D_+(f)$ ($f\in S_1$)들의 합집합이므로
(\hyperref[II.2.3.14-ko]{2.3.14}), 고려하는 각각의 가정 아래에서
$\lambda_f$와 $\mu_f$가 $f$가 동차이고 \emph{차수가 $1$}일 때
동형사상임을 보이는 것으로 귀착된다. 이때 $\mathbf{Z}$-쌍선형 사상
\[
  M_m\times N_n\longrightarrow
  M_{(f)}\otimes_{S_{(f)}}N_{(f)}
\]
을 $(x,y)$에 원소 $(x/f^m)\otimes(y/f^n)$을 대응시켜 정의한다
($m<0$이면 $x/f^m$은 $f^{-m}x/1$을 뜻한다). 이 사상들로
$\mathbf{Z}$-선형 사상
\[
  M\otimes_{\mathbf{Z}}N\longrightarrow
  M_{(f)}\otimes_{S_{(f)}}N_{(f)},
\]
을 정의할 수 있다. $s\in S_q$이면 이 사상은 $(sx)\otimes y$를
\[
  (s/f^q)\bigl((x/f^m)\otimes(y/f^n)\bigr)
\]
으로 보낸다($x\in M_m$, $y\in N_n$). 따라서 가군의 쌍준동형
\[
  \gamma_f:M\otimes_S N\longrightarrow
  M_{(f)}\otimes_{S_{(f)}}N_{(f)},
\]
을 얻는다. 이는 표준 준동형 $S\to S_{(f)}$---$s\in S_q$를
$s/f^q$로 보내는 준동형---에 관한 것이다. 더 나아가 원소
$\sum_i(x_i\otimes y_i)$가 $M\otimes_S N$에 속하고, 여기서
$x_i$, $y_i$는 각각 차수가
$m_i$, $n_i$인 동차 원소이다---에 대하여
\[
  f^r\sum_i(x_i\otimes y_i)=0,
\]
즉 $\sum_i(f^rx_i\otimes y_i)=0$이라고 가정하자.
(\hyperref[0.1.3.4-ko]{0, 1.3.4})에 따라
\[
  \sum_i(f^rx_i/f^{m_i+r})\otimes(y_i/f^{n_i})=0,
\]
을 얻는다. 다시 말해 $\gamma_f(\sum_i(x_i\otimes y_i))=0$이다. 따라서
$\gamma_f$는 다음과 같이 인수분해된다.
\[
  M\otimes_S N\longrightarrow(M\otimes_S N)_f
  \xrightarrow{\gamma'_f}
  M_{(f)}\otimes_{S_{(f)}}N_{(f)}~;
\]
$\lambda'_f$를 $\gamma'_f$의
\oldpage[II]{35}
$(M\otimes_S N)_{(f)}$로의 제한이라 하면,
$\lambda_f$와 $\lambda'_f$가 서로 역인 $S_{(f)}$-준동형임을 곧바로
확인할 수 있다. 이로써 명제의 첫 번째 부분이 증명된다.

두 번째 부분을 증명하기 위해, $M$이 준동형 $P\to Q$의 여핵이라고
가정하자. 이는 등급 $S$-가군의 준동형이며, $P$와 $Q$는 $S(n)$ 꼴인
가군들의 유한 직합이다. $\operatorname{Hom}_S(L,N)$의 $L$에 관한
왼쪽 완전성과 $M_{(f)}$의 $M$에 관한 완전성을 사용하면, $\mu_f$가
$M=S(n)$일 때 동형사상임을 보이는 것으로 곧바로 귀착된다. 이제 임의의
동차 원소 $z$를 $N$에서 택하고, $u_z$를 $S(n)$에서 $N$으로 가며
$u_z(1)=z$를 만족하는 준동형이라 하자. 그러면
$\eta:z\mapsto u_z$는 차수 $0$이고 $N(-n)$에서
$\operatorname{Hom}_S(S(n),N)$으로 가는 동형사상임을 곧바로
알 수 있다. 이에 대응하는 동형사상
\[
  \eta_f:\bigl(N(-n)\bigr)_{(f)}
  \longrightarrow
  \bigl(\operatorname{Hom}_S(S(n),N)\bigr)_{(f)}.
\]
이 있다. 한편 $\eta'_f$를 다음 동형사상이라 하자.
\[
  N_{(f)}\longrightarrow
  \operatorname{Hom}_{S_{(f)}}\bigl(S(n)_{(f)},N_{(f)}\bigr)
\]
이는 각 $z'\in N_{(f)}$에 준동형 $v_{z'}$를 대응시키며, 이 준동형은
$v_{z'}(s/f^k)=sz'/f^{n+k}$를 만족한다
($s\in S_{n+k}=(S(n))_k$). 합성 사상
\[
  \bigl(N(-n)\bigr)_{(f)}
  \xrightarrow{\eta_f}
  \bigl(\operatorname{Hom}_S(S(n),N)\bigr)_{(f)}
  \xrightarrow{\mu_f}
  \operatorname{Hom}_{S_{(f)}}\bigl(S(n)_{(f)},N_{(f)}\bigr)
  \xrightarrow{{\eta'_f}^{-1}}N_{(f)}
\]
은 $z/f^h\mapsto z/f^{h-n}$으로 주어지는
$\bigl(N(-n)\bigr)_{(f)}$에서 $N_{(f)}$로 가는 동형사상임을 쉽게 확인할 수
있다. 따라서 $\mu_f$는 동형사상이다.
\end{proof}

아이디얼 $S_+$가 $S_1$에 의해 생성되면,
(\hyperref[II.2.5.13-ko]{2.5.13})에서 $S$의 모든 등급 아이디얼
$\mathfrak{J}$와 모든 등급 $S$-가군 $M$에 대하여
\[
\label{II.2.5.13.1-ko}
  \widetilde{\mathfrak{J}}\cdot\widetilde{M}
  =\widetilde{\mathfrak{J}\cdot M}
\tag{2.5.13.1}
\]
가 표준 동형까지 성립함을 얻는다. 실제로 이는 다음 도식의 가환성에서
따른다.
\[
  \xymatrix{
    \widetilde{\mathfrak{J}}\otimes_{\mathcal{O}_X}\widetilde{M}
    \ar[rr]^\lambda\ar[dr]
    && \widetilde{\mathfrak{J}\otimes_S M}\ar[dl]\\
    & \widetilde{M} &
  }
\]
이 가환성은 (\hyperref[II.2.5.11.3-ko]{2.5.11.3})의 경우와 마찬가지로
확인한다.

\begin{corollary}[2.5.14]
\phantomsection
\label{II.2.5.14-ko}
$S$가 $S_1$에 의해 생성된다고 하자. 어떠한 $m,n\in\mathbf{Z}$에
대해서도 다음 공식들이 성립한다.
\[
\label{II.2.5.14.1-ko}
  \mathcal{O}_X(m)\otimes_{\mathcal{O}_X}\mathcal{O}_X(n)
  =\mathcal{O}_X(m+n)
\tag{2.5.14.1}
\]
\[
\label{II.2.5.14.2-ko}
  \mathcal{O}_X(n)=\bigl(\mathcal{O}_X(1)\bigr)^{\otimes n}
\tag{2.5.14.2}
\]
여기서 등식들은 표준 동형까지 성립한다.
\end{corollary}

\begin{proof}
첫째 공식은 (\hyperref[II.2.5.13-ko]{2.5.13})과 차수 $0$인 표준 동형
$S(m)\otimes_S S(n)\xrightarrow{\sim}S(m+n)$의 존재에서 따른다. 이 동형은
$1\otimes1$---여기서 첫째 인자 $1$은 $(S(m))_{-m}$에, 둘째 인자 $1$은
$(S(n))_{-n}$에 속한다---을 원소 $1\in(S(m+n))_{-m-n}$에 대응시킨다.
이제 둘째 공식은 $n=-1$인 경우만 증명하면 충분하다. 그리고
(\hyperref[II.2.5.13-ko]{2.5.13})에 의하면 이는
$\operatorname{Hom}_S(S(1),S)$가 $S(-1)$과 표준적으로 동형임을 보이는
것으로 귀착한다. 정의들 (\hyperref[II.2.1.2-ko]{2.1.2})을 되짚어 보고
$S(1)$이 한 원소로 생성되는 $S$-가군임을 상기하면 곧바로 확인된다.
\end{proof}

\oldpage[II]{36}
\begin{corollary}[2.5.15]
\phantomsection
\label{II.2.5.15-ko}
$S$가 $S_1$에 의해 생성된다고 하자. 모든 등급 $S$-가군 $M$과 모든
$n\in\mathbf{Z}$에 대하여
\[
\label{II.2.5.15.1-ko}
  \widetilde{M(n)}=\widetilde{M}(n)
\tag{2.5.15.1}
\]
가 표준 동형까지 성립한다.
\end{corollary}

\begin{proof}
이는 정의들 (\hyperref[II.2.5.10.2-ko]{2.5.10.2})와
(\hyperref[II.2.5.10.1-ko]{2.5.10.1}), 명제
(\hyperref[II.2.5.13-ko]{2.5.13}), 그리고 차수 $0$인 표준 동형
$M(n)\xrightarrow{\sim}M\otimes_S S(n)$의 존재에서 따른다. 이 동형은 모든
$z\in(M(n))_h=M_{n+h}$를
$z\otimes1\in M_{n+h}\otimes(S(n))_{-n}\subset(M\otimes_S S(n))_h$에
대응시킨다.
\end{proof}

\begin{env}[2.5.16]
\phantomsection
\label{II.2.5.16-ko}
$S'_0=\mathbf{Z}$이고 $n>0$에 대하여 $S'_n=S_n$인 등급환을 $S'$라 하자.
그러면 $f\in S_d$ ($d>0$)일 때 모든 $n\in\mathbf{Z}$에 대하여
$(S(n))_{(f)}=(S'(n))_{(f)}$이다. 실제로 $(S'(n))_{(f)}$의 원소는
$x\in S'_{n+kd}$ ($k>0$)인 $x/f^k$ 꼴이며, 항상 $n+kd\neq0$이 되도록
$k$를 택할 수 있다. $X=\operatorname{Proj}(S)$와
$X'=\operatorname{Proj}(S')$가 표준적으로 동일시되므로
(\hyperref[II.2.4.7-ko]{2.4.7, (ii)}), 모든 $n\in\mathbf{Z}$에 대하여
$\mathcal{O}_X(n)$과 $\mathcal{O}_{X'}(n)$은 앞의 동일시에서 서로
대응한다.

한편 모든 $d>0$과 모든 $n\in\mathbf{Z}$에 대하여
\[
  (S^{(d)}(n))_h=S_{(n+h)d}=(S(nd))_{hd}
\]
이며, $f\in S_d$이면 따라서
$(S^{(d)}(n))_{(f)}=(S(nd))_{(f)}$이다. 스킴
$X=\operatorname{Proj}(S)$와 $X^{(d)}=\operatorname{Proj}(S^{(d)})$가
표준적으로 동일시됨을 알고 있다
(\hyperref[II.2.4.7-ko]{2.4.7, (i)}). 앞의 논의에 따르면 $S_0$-대수
$S^{(d)}$가 $S_d$에 의해 생성될 때, 모든 $n\in\mathbf{Z}$에 대하여
$\mathcal{O}_X(nd)$와 $\mathcal{O}_{X^{(d)}}(n)$은 이 동일시에서 서로
대응한다.
\end{env}

\begin{proposition}[2.5.17]
\phantomsection
\label{II.2.5.17-ko}
$d$를 $>0$인 정수라 하고
$U=\bigcup_{f\in S_d}D_+(f)$라 하자. 표준 준동형
$\mathcal{O}_X(nd)\otimes_{\mathcal{O}_X}\mathcal{O}_X(-nd)
\to\mathcal{O}_X$의 $U$에 대한 제한은 모든 정수 $n$에 대하여
동형이다.
\end{proposition}

\begin{proof}
(\hyperref[II.2.5.16-ko]{2.5.16})에 의해 $d=1$인 경우로 환원할 수 있고,
그 경우 결론은 (\hyperref[II.2.5.13-ko]{2.5.13})의 증명에서 따른다.
\end{proof}

\subsection{$\operatorname{Proj}(S)$ 위의 층에 대응하는 등급
$S$-가군.}
\label{subsection:II.2.6-ko}

\emph{이 절 전체에서 $S$의 아이디얼 $S_+$가 차수 $1$의 동차 원소들의
집합 $S_1$에 의해 생성된다고 가정한다.}

\begin{env}[2.6.1]
\phantomsection
\label{II.2.6.1-ko}
그러면 $\mathcal{O}_X$-가군 $\mathcal{O}_X(1)$은
\emph{가역}이다(\hyperref[II.2.5.9-ko]{2.5.9}). 이때 모든
$\mathcal{O}_X$-가군 $\mathcal{F}$
(\hyperref[0.5.4.6-ko]{0, 5.4.6})에 대하여
\[
\label{II.2.6.1.1-ko}
  \Gamma_*(\mathcal{F})
  =\Gamma_*(\mathcal{O}_X(1),\mathcal{F})
  =\bigoplus_{n\in\mathbf{Z}}\Gamma(X,\mathcal{F}(n))
\tag{2.6.1.1}
\]
로 둔다. 여기서는 (\hyperref[II.2.5.14.2-ko]{2.5.14.2})를 고려하였다.
(\hyperref[0.5.4.6-ko]{0, 5.4.6})에서 상기한 바와 같이
$\Gamma_*(\mathcal{O}_X)$는 \emph{등급환}의 구조를 가지며,
$\Gamma_*(\mathcal{F})$는 $\Gamma_*(\mathcal{O}_X)$ 위의
\emph{등급가군} 구조를 가진다.

$\mathcal{O}_X(n)$이 국소 자유이므로
$\mathcal{F}\mapsto\Gamma_*(\mathcal{F})$는 $\mathcal{F}$에 관한
공변 가법 \emph{왼쪽 완전} 함자이다. 특히 $\mathcal{J}$가
$\mathcal{O}_X$의 아이디얼층이면, $\Gamma_*(\mathcal{J})$는
$\Gamma_*(\mathcal{O}_X)$의 \emph{등급 아이디얼}과 표준적으로
동일시된다.
\end{env}

\begin{env}[2.6.2]
\phantomsection
\label{II.2.6.2-ko}
$M$을 등급 $S$-가군이라 하자. 모든 $f\in S_d$ ($d>0$)에 대하여
$x\mapsto x/1$은 아벨군의 준동형 $M_0\to M_{(f)}$이고,
$M_{(f)}$는 표준적으로

\oldpage[II]{37}

$\Gamma(D_+(f),\widetilde{M})$와 동일시되므로 아벨군의 준동형
\[
  \alpha_0^f:M_0\longrightarrow\Gamma(D_+(f),\widetilde{M}).
\]
을 얻는다. 모든 $g\in S_e$ ($e>0$)에 대하여 다음 도식이 가환함은
명백하다.
\[
  \xymatrix{
    & \Gamma(D_+(f),\widetilde{M})\ar[dd]\\
    M_0\ar[ur]^{\alpha_0^f}\ar[dr]_{\alpha_0^{fg}}\\
    & \Gamma(D_+(fg),\widetilde{M})
  }
\]
이는 모든 $x\in M_0$에 대하여 $\widetilde{M}$의 단면
$\alpha_0^f(x)$와 $\alpha_0^g(x)$가 $D_+(f)\cap D_+(g)$에서
일치함을 뜻한다. 따라서 각 $D_+(f)$로의 제한이 $\alpha_0^f(x)$인
유일한 단면 $\alpha_0(x)\in\Gamma(X,\widetilde{M})$이 존재한다.
이로써 ($S$가 $S_1$에 의해 생성된다는 가정을 사용하지 않고) 아벨군의
준동형
\[
\label{II.2.6.2.1-ko}
  \alpha_0:M_0\longrightarrow\Gamma(X,\widetilde{M}).
\tag{2.6.2.1}
\]
을 정의하였다.

이 결과를 모든 $n\in\mathbf{Z}$에 대하여 등급 $S$-가군 $M(n)$에
적용하면, 각 $n\in\mathbf{Z}$에 대하여 아벨군의 준동형
\[
\label{II.2.6.2.2-ko}
  \alpha_n:M_n=(M(n))_0\longrightarrow\Gamma(X,\widetilde{M}(n))
\tag{2.6.2.2}
\]
을 얻는다((\hyperref[II.2.5.15-ko]{2.5.15})를 고려한다). 따라서 등급
아벨군의 함자적인 준동형(차수 $0$)
\[
\label{II.2.6.2.3-ko}
  \alpha:M\longrightarrow\Gamma_*(\widetilde{M})
\tag{2.6.2.3}
\]
을 얻으며, 이를 $\alpha_M$으로도 나타낸다. 이 준동형은 각 $M_n$에서
$\alpha_n$과 일치한다.

특히 $M=S$로 두면, $\Gamma_*(\mathcal{O}_X)$의 곱셈의 정의
(\hyperref[0.5.4.6-ko]{0, 5.4.6})를 고려하여
$\alpha:S\to\Gamma_*(\mathcal{O}_X)$가 등급환의 준동형임을 곧바로
확인할 수 있다. 또한 모든 등급 $S$-가군 $M$에 대하여
(\hyperref[II.2.6.2.3-ko]{2.6.2.3})은 등급 가군의 쌍준동형이다.
\end{env}

\begin{proposition}[2.6.3]
\phantomsection
\label{II.2.6.3-ko}
모든 $f\in S_d$ ($d>0$)에 대하여, $D_+(f)$는 $\mathcal{O}_X(d)$의
단면 $\alpha_d(f)$가 소멸하지 않는 $\mathfrak{p}\in X$들의 집합과
일치한다(\hyperref[0.5.5.2-ko]{0, 5.5.2}).
\end{proposition}

\begin{proof}
가정에 따라 $X=\bigcup_{g\in S_1}D_+(g)$이므로, 모든 $g\in S_1$에
대하여 $\alpha_d(f)$가 소멸하지 않는 $\mathfrak{p}\in D_+(g)$들의
집합이 $D_+(fg)$와 일치함을 보이면 충분하다. 그런데 $\alpha_d(f)$의
$D_+(g)$에 대한 제한은 정의상 $(S(d))_{(g)}$의 원소 $f/1$에 대응하는
단면이다. 표준 동형
$(S(d))_{(g)}\xrightarrow{\sim}S_{(g)}$
(\hyperref[II.2.5.7-ko]{2.5.7}) 아래에서, $D_+(g)$ 위의
$\mathcal{O}_X(d)$의 이 단면은 $S_{(g)}$의 원소 $f/g^d$에 대응하는
$D_+(g)$ 위의 $\mathcal{O}_X$의 단면과 동일시된다. 이 단면이
$\mathfrak{p}\in D_+(g)$에서 소멸한다는 것은 $f/g^d\in\mathfrak{q}$임을
뜻한다. 여기서 $\mathfrak{q}$는 $\mathfrak{p}$에 대응하는 $S_{(g)}$의
소아이디얼이다(\hyperref[II.2.3.6-ko]{2.3.6}). 정의상 이는
$f\in\mathfrak{p}$라는 뜻이므로 명제가 성립한다.
\end{proof}

\begin{env}[2.6.4]
\phantomsection
\label{II.2.6.4-ko}
이제 $\mathcal{F}$를 $\mathcal{O}_X$-가군이라 하고
$M=\Gamma_*(\mathcal{F})$로 두자. 등급환 준동형
$\alpha:S\to\Gamma_*(\mathcal{O}_X)$가 존재하므로, $M$을 등급
$S$-가군으로 볼 수 있다. 모든 $f\in S_d$ ($d>0$)에 대하여
(\hyperref[II.2.6.3-ko]{2.6.3})에서 $\mathcal{O}_X(d)$의 단면
$\alpha_d(f)$의 $D_+(f)$에 대한 제한이 가역임을 알 수 있다. 따라서
모든 $n>0$에 대하여 $\mathcal{O}_X(nd)$의 단면
$\alpha_{nd}(f^n)$\footnote{\textbf{번역 주.} 인쇄 원문은 이 문단의 세
곳에서 각각 $\alpha_d(f^n)$, $\alpha_d(f^k)$, $\alpha_d(f^n)$을
쓰지만, $f^n\in S_{nd}$와 $f^k\in S_{kd}$이므로
EGA2-CANON-DISP-0002에 따라 각각 $\alpha_{nd}(f^n)$,
$\alpha_{kd}(f^k)$, $\alpha_{nd}(f^n)$으로 교정했다.}의 $D_+(f)$에
대한 제한도 가역이다. 이제
$z\in M_{nd}=\Gamma(X,\mathcal{F}(nd))$ ($n>0$)라 하자. 어떤 정수
$k\geq0$에 대하여 $f^kz$의 $D_+(f)$에 대한 제한, 즉
\oldpage[II]{38}
단면
$(z|D_+(f))(\alpha_{kd}(f^k)|D_+(f))$가
$\mathcal{F}((n+k)d)$의 단면으로서 영이면, 앞의 설명에 따라
$z|D_+(f)=0$이기도 하다. 이로써 원소 $z/f^n$에 $D_+(f)$ 위의
$\mathcal{F}$의 단면
$(z|D_+(f))(\alpha_{nd}(f^n)|D_+(f))^{-1}$을 대응시켜
$S_{(f)}$-가군 준동형
$\beta_f:M_{(f)}\to\Gamma(D_+(f),\mathcal{F})$를 정의할 수 있다.
또한 $g\in S_e$ ($e>0$)일 때 다음 도식이 가환함을 곧바로 확인할 수
있다.
\[
\label{II.2.6.4.1-ko}
  \xymatrix{
    M_{(f)}\ar[r]^{\beta_f}\ar[d]
    & \Gamma(D_+(f),\mathcal{F})\ar[d]\\
    M_{(fg)}\ar[r]_{\beta_{fg}}
    & \Gamma(D_+(fg),\mathcal{F})
  }
\tag{2.6.4.1}
\]
$M_{(f)}$가 $\Gamma(D_+(f),\widetilde{M})$과 표준적으로 동일시되고
$D_+(f)$들이 $X$의 위상의 기저를 이룬다는 사실
(\hyperref[II.2.3.4-ko]{2.3.4})을 상기하면, $\beta_f$들이 유일한 표준
$\mathcal{O}_X$-가군 준동형
\[
\label{II.2.6.4.2-ko}
  \beta:\widetilde{\Gamma_*(\mathcal{F})}
  \longrightarrow\mathcal{F}
\tag{2.6.4.2}
\]
에서 유도됨을 알 수 있다. 이 준동형은 $\beta_{\mathcal{F}}$로도
표기하며 명백히 함자적이다.
\end{env}

\begin{proposition}[2.6.5]
\phantomsection
\label{II.2.6.5-ko}
$M$을 등급 $S$-가군, $\mathcal{F}$를 $\mathcal{O}_X$-가군이라 하자.
그러면 다음 합성 준동형들은 각각 항등 준동형이다.
\[
\label{II.2.6.5.1-ko}
  \widetilde{M}\xrightarrow{\widetilde{\alpha}}
  \widetilde{\Gamma_*(\widetilde{M})}
  \xrightarrow{\beta}\widetilde{M}
\tag{2.6.5.1}
\]
\[
\label{II.2.6.5.2-ko}
  \Gamma_*(\mathcal{F})\xrightarrow{\alpha}
  \Gamma_*(\widetilde{\Gamma_*(\mathcal{F})})
  \xrightarrow{\Gamma_*(\beta)}\Gamma_*(\mathcal{F})
\tag{2.6.5.2}
\]
\end{proposition}

\begin{proof}
(\hyperref[II.2.6.5.1-ko]{2.6.5.1})의 확인은 국소적이다. 열린집합
$D_+(f)$에서는 정의들과, 준연접 층에 적용한 $\beta$가 $D_+(f)$ 위의
단면들에 대한 작용으로 정해진다는 사실
(\hyperref[I.1.3.8-ko]{I, 1.3.8})에서 곧바로 따른다.
(\hyperref[II.2.6.5.2-ko]{2.6.5.2})의 확인은 각 차수에서 따로 이루어진다.
$M=\Gamma_*(\mathcal{F})$로 두면
$M_n=\Gamma(X,\mathcal{F}(n))$이고
$(\Gamma_*(\widetilde{M}))_n=\Gamma(X,\widetilde{M}(n))
=\Gamma(X,\widetilde{M(n)})$이다. 그런데 $f\in S_1$이고 $z\in M_n$이면,
$\alpha_n^f(z)$는 $(M(n))_{(f)}$의 원소 $z/1$이며
$(f/1)^n(z/f^n)$과 같다. 이에 $\beta_f$로 대응하는 것은 $D_+(f)$ 위의
단면
\[
  ((\alpha_1(f))^n|D_+(f))
  ((z|D_+(f))((\alpha_1(f))^n|D_+(f))^{-1})
\]
이다. 이는 곧 $z$의 $D_+(f)$에 대한 제한이므로
(\hyperref[II.2.6.5.2-ko]{2.6.5.2})가 확인된다.
\end{proof}

\subsection{유한성 조건.}
\label{subsection:II.2.7-ko}

\begin{proposition}[2.7.1]
\phantomsection
\label{II.2.7.1-ko}
\begin{enumerate}
  \item[\rm(i)] $S$가 뇌터인 등급환이면
    $X=\operatorname{Proj}(S)$는 뇌터 스킴이다.
  \item[\rm(ii)] $S$가 유한형 등급 $A$-대수이면
    $X=\operatorname{Proj}(S)$는 $Y=\operatorname{Spec}(A)$ 위의
    유한형 스킴이다.
\end{enumerate}
\end{proposition}

\oldpage[II]{39}
\begin{proof}
\begin{enumerate}
  \item[\rm(i)] $S$가 뇌터이면 아이디얼 $S_+$는 동차 원소들
    $f_i$ ($1\leq i\leq p$)로 이루어진 유한 생성계를 갖는다. 따라서
    (\hyperref[II.2.3.14-ko]{2.3.14}) 바탕공간 $X$는
    $D_+(f_i)=\operatorname{Spec}(S_{(f_i)})$들의 합집합이고, 각
    $S_{(f_i)}$가 뇌터임을 보이면 충분하다. 이는
    (\hyperref[II.2.2.6-ko]{2.2.6})에서 따른다.
  \item[\rm(ii)] 가정으로부터 $S_0$가 유한형 $A$-대수이고 $S$가
    유한형 $S_0$-대수임이 따른다. 따라서 $S_+$는 유한 생성
    아이디얼이다(\hyperref[II.2.1.4-ko]{2.1.4}). 그러면 (i)에서와 같이
    $f\in S_d$일 때 $S_{(f)}$가 유한형 $A$-대수임을 증명하면
    충분하다. (\hyperref[II.2.2.5-ko]{2.2.5})에 따라 $S^{(d)}$가 유한형
    $A$-대수임을 보이면 충분하고, 이는
    (\hyperref[II.2.1.6-ko]{2.1.6})에서 따른다.
\end{enumerate}
\end{proof}

\begin{env}[2.7.2]
\phantomsection
\label{II.2.7.2-ko}
이하에서는 등급 $S$-가군 $M$에 관한 다음 유한성 조건들을 고려하겠다.
\begin{enumerate}
  \item[\rm(TF)] 어떤 정수 $n$이 존재하여 부분가군
    $\bigoplus_{k\geq n}M_k$가 유한 생성 $S$-가군이다.
  \item[\rm(TN)] 어떤 정수 $n$이 존재하여 $k\geq n$이면 $M_k=0$이다.
\end{enumerate}

$M$이 (TN)을 만족하면 $S_+$의 모든 동차 원소 $f$에 대하여
$M_{(f)}=0$이고, 따라서 $\widetilde{M}=0$이다.

$M$, $N$을 두 등급 $S$-가군이라 하자. 차수 $0$의 준동형
$u:M\to N$에 대하여, 어떤 정수 $n$이 존재하여 $k\geq n$일 때
$u_k:M_k\to N_k$가 단사(각각 전사, 전단사)이면 $u$를
(TN)-\emph{단사}(각각 (TN)-\emph{전사}, (TN)-\emph{전단사})라고 한다.
따라서 $u$가 (TN)-단사(각각 (TN)-전사)라는 것은
$\operatorname{Ker}u$(각각 $\operatorname{Coker}u$)가 (TN)을
만족한다는 것과 같다. (\hyperref[II.2.5.4-ko]{2.5.4})에 따라 $u$가
(TN)-단사(각각 (TN)-전사, (TN)-전단사)이면 $\widetilde{u}$는
단사(각각 전사, 전단사)이다. $u$가 (TN)-전단사일 때 $u$를
(TN)-\emph{동형사상}이라고도 한다.
\end{env}

\begin{proposition}[2.7.3]
\phantomsection
\label{II.2.7.3-ko}
$S$를 아이디얼 $S_+$가 유한 생성인 등급환이라 하고, $M$을 등급
$S$-가군이라 하자.
\begin{enumerate}
  \item[\rm(i)] $M$이 조건 (TF)를 만족하면 $\mathcal{O}_X$-가군
    $\widetilde{M}$은 유한 생성이다.
  \item[\rm(ii)] $M$이 (TF)를 만족한다고 가정하자. $\widetilde{M}=0$일
    필요충분조건은 $M$이 (TN)을 만족하는 것이다.
\end{enumerate}
\end{proposition}

\begin{proof}
방금 본 바와 같이 조건 (TN)은 $\widetilde{M}=0$을 함의한다. $M$이
(TF)를 만족하면, 가정에 따라 유한 생성인 등급 부분가군
$M'=\bigoplus_{k\geq n}M_k$에 대하여 $M/M'$은 (TN)을 만족한다. 따라서
$\widetilde{M/M'}=0$이고, 함자 $M\mapsto\widetilde{M}$의 완전성
(\hyperref[II.2.5.4-ko]{2.5.4})으로부터
$\widetilde{M}=\widetilde{M'}$을 얻는다. 그러므로 $\widetilde{M}$이 유한
생성임을 증명하기 위해 $M$이 유한 생성인 경우로 귀착할 수 있다. 문제는
국소적이므로 $M_{(f)}$가 유한 생성 $S_{(f)}$-가군임을 보이면 충분하다
(\hyperref[I.1.3.9-ko]{I, 1.3.9}). 그런데 $M^{(d)}$는 유한 생성
$S^{(d)}$-가군이고(\hyperref[II.2.1.6-ko]{2.1.6, (iii)}), 이제 주장은
(\hyperref[II.2.2.5-ko]{2.2.5})에서 따른다.

이제 $M$이 (TF)를 만족하고 $\widetilde{M}=0$이라고 가정하자. 위와 같은
표기에서 $\widetilde{M'}=0$이고, $M'$에 대한 조건 (TN)은 $M$에 대한
조건 (TN)과 동치이다. 따라서 $\widetilde{M}=0$이 $M$의 조건 (TN)을
함의함을 증명하기 위해, 다시 $M$이 유한 개의 동차 원소 $x_i$
($1\leq i\leq p$)에 의해 생성되는 경우만 고려하면 된다. 또
$(f_j)_{1\leq j\leq q}$를 아이디얼 $S_+$의 동차 생성계라 하자. 가정에
따라 모든 $j$에 대하여 $M_{(f_j)}=0$이므로, 어떤 정수 $n$이 존재하여
모든 $i$, $j$에 대하여 $f_j^nx_i=0$이다. $n_j$를 $f_j$의 차수라 하고,
$\sum_jr_j\leq nq$를 만족하는 정수족 $(r_j)$에 대하여
$\sum_jr_jn_j$가 취하는 최댓값을 $m$이라 하자. 그러면

\oldpage[II]{40}
$k>m$이면 모든 $i$에 대하여 $S_kx_i=0$임이 명백하다. $x_i$들의 차수 중
최댓값을 $h$라 하면 $k>h+m$일 때 $M_k=0$이고, 이것으로 증명이 끝난다.
\end{proof}

\begin{corollary}[2.7.4]
\phantomsection
\label{II.2.7.4-ko}
$S$를 아이디얼 $S_+$가 유한 생성인 등급환이라 하자. 그러면
$X=\operatorname{Proj}(S)=\varnothing$이기 위한 필요충분조건은 어떤 $n$이
존재하여 $k\geq n$일 때 $S_k=0$인 것이다.
\end{corollary}

\begin{proof}
실제로 조건 $X=\varnothing$은
$\mathcal{O}_X=\widetilde{S}=0$과 동치이고, $S$는 한 원소로 생성되는
$S$-가군이다.
\end{proof}

\begin{theorem}[2.7.5]
\phantomsection
\label{II.2.7.5-ko}
아이디얼 $S_+$가 차수 $1$인 유한 개의 동차 원소로 생성된다고 가정하고,
$X=\operatorname{Proj}(S)$라 하자. 그러면 모든 준연접
$\mathcal{O}_X$-가군 $\mathcal{F}$에 대하여 표준 준동형
$\beta:\widetilde{\Gamma_*(\mathcal{F})}
\to\mathcal{F}$ (\hyperref[II.2.6.4-ko]{2.6.4})는 동형이다.
\end{theorem}

\begin{proof}
실제로 $S_+$가 유한 개의 원소 $f_i\in S_1$로 생성되면, $X$는 준콤팩트인
부분공간 $\operatorname{Spec}(S_{(f_i)})$
(\hyperref[II.2.3.6-ko]{2.3.6})들의 합집합이므로 준콤팩트이다. 또한 $X$는
스킴이다(\hyperref[II.2.4.2-ko]{2.4.2}).
(\hyperref[I.9.3.2-ko]{I, 9.3.2}),
(\hyperref[II.2.5.14.2-ko]{2.5.14.2}) 및
(\hyperref[II.2.6.3-ko]{2.6.3})에 따라 모든 $f\in S_d$ ($d>0$)에 대하여
표준 동형
$(\Gamma_*(\mathcal{F}))_{(\alpha_d(f))}\xrightarrow{\sim}
\Gamma(D_+(f),\mathcal{F})$을 얻는다. 더욱이
$(\Gamma_*(\mathcal{F}))_{(\alpha_d(f))}$는
$\Gamma_*(\mathcal{F})$를 $\Gamma_*(\mathcal{O}_X)$-가군으로 볼 때의
가군인데, 이는 $\Gamma_*(\mathcal{F})$를 $S$-가군으로 볼 때의
$(\Gamma_*(\mathcal{F}))_{(f)}$와 다름없다. 앞의 표준 동형의 정의
(\hyperref[I.9.3.1-ko]{I, 9.3.1})를 참조하면 이 동형이 준동형
$\beta_f$와 일치함을 알 수 있고, 정리가 따른다.
\end{proof}

\begin{remark}[2.7.6]
\phantomsection
\label{II.2.7.6-ko}
등급환 $S$가 \emph{뇌터}라고 가정하면, 아이디얼 $S_+$가 차수 $1$의 동차
원소 전체의 집합 $S_1$에 의해 생성된다고 가정하는 것만으로
(\hyperref[II.2.7.5-ko]{2.7.5})의 조건이 \emph{자동으로} 성립한다.
\end{remark}

\begin{corollary}[2.7.7]
\phantomsection
\label{II.2.7.7-ko}
(\hyperref[II.2.7.5-ko]{2.7.5})의 가정 아래에서 모든 준연접
$\mathcal{O}_X$-가군 $\mathcal{F}$는 $\widetilde{M}$ 꼴의
$\mathcal{O}_X$-가군과 동형이다. 여기서 $M$은 등급 $S$-가군이다.
\end{corollary}

\begin{corollary}[2.7.8]
\phantomsection
\label{II.2.7.8-ko}
(\hyperref[II.2.7.5-ko]{2.7.5})의 가정 아래에서 모든 유한 생성 준연접
$\mathcal{O}_X$-가군 $\mathcal{F}$는 $\widetilde{N}$ 꼴의
$\mathcal{O}_X$-가군과 동형이다. 여기서 $N$은 유한 생성 등급
$S$-가군이다.
\end{corollary}

\begin{proof}
(\hyperref[II.2.7.7-ko]{2.7.7})에 따라 $\mathcal{F}=\widetilde{M}$이라고
가정할 수 있다. 여기서 $M$은 등급 $S$-가군이다.
$(f_\lambda)_{\lambda\in L}$를 $M$의 동차 생성계라 하자. $L$의 모든 유한
부분집합 $H$에 대하여, $\lambda\in H$인 $f_\lambda$들로 생성되는 $M$의
등급 부분가군을 $M_H$라 하자. $M$이 부분가군 $M_H$들의 귀납극한임은
명백하므로, $\mathcal{F}$는 그 $\mathcal{O}_X$-부분가군
$\widetilde{M_H}$들의 귀납극한이다
(\hyperref[II.2.5.4-ko]{2.5.4}). 그런데 $\mathcal{F}$가 유한 생성이고
바탕공간 $X$가 준콤팩트이므로, (\hyperref[0.5.2.3-ko]{0, 5.2.3})에서
$L$의 어떤 유한 부분집합 $H$에 대하여
$\mathcal{F}=\widetilde{M_H}$임이 따른다.
\end{proof}

\begin{corollary}[2.7.9]
\phantomsection
\label{II.2.7.9-ko}
(\hyperref[II.2.7.5-ko]{2.7.5})의 가정 아래에서 $\mathcal{F}$를 유한 생성
준연접 $\mathcal{O}_X$-가군이라 하자. 그러면 어떤 정수 $n_0$가 존재하여,
모든 $n\geq n_0$에 대하여 $\mathcal{F}(n)$은 $\mathcal{O}_X^k$의 어떤
몫가군과 동형이다. 여기서 $k>0$은 $n$에 의존한다. 따라서
$\mathcal{F}(n)$은 $X$ 위의 유한 개 단면으로 생성된다
(\hyperref[0.5.1.1-ko]{0, 5.1.1}).
\end{corollary}

\begin{proof}
(\hyperref[II.2.7.8-ko]{2.7.8})에 따라 $\mathcal{F}=\widetilde{M}$이라고
가정할 수 있다. 여기서 $M$은 $S(m_i)$ 꼴의 $S$-가군들의 유한 직합의
몫가군이다. 따라서 (\hyperref[II.2.5.4-ko]{2.5.4})에 따라 $M=S(m)$인
경우로 귀착할 수 있고, 이때
$\mathcal{F}(n)=\widetilde{S(m+n)}
=\mathcal{O}_X(m+n)$이다
(\hyperref[II.2.5.15-ko]{2.5.15}). 그러므로 다음 보조정리를 증명하면
충분하다.

\oldpage[II]{41}
\begin{lemma}[2.7.9.1]
\phantomsection
\label{II.2.7.9.1-ko}
(\hyperref[II.2.7.5-ko]{2.7.5})의 가정 아래에서 모든 $n\geq0$에 대하여,
어떤 정수 $k$가 존재하여(이는 $n$에 의존한다) 전사 준동형
$\mathcal{O}_X^k\to\mathcal{O}_X(n)$이 존재한다.
\end{lemma}

실제로 (\hyperref[II.2.7.2-ko]{2.7.2})에 따라, 적당한 $k$에 대하여 등급
$S$-가군 곱 $S^k$에서 $S(n)$으로 가는 \emph{차수 $0$}의
(TN)-전사 준동형 $u$가 존재함을 보이면 충분하다. 그런데
$(S(n))_0=S_n$이고 가정에 따라 모든 $h>0$에 대하여 $S_h=S_1^h$이므로
$SS_n=S_n+S_{n+1}+\cdots$이다. $S_n$은 유한 생성 $S_0$-가군이므로
(\hyperref[II.2.1.5-ko]{2.1.5} 및
\hyperref[II.2.1.6-ko]{2.1.6, (i)}), 이 가군의 생성계
$(a_i)_{1\leq i\leq k}$를 택하자. 준동형 $u$는 $S^k$의 $i$번째 표준 기저
원소를 $a_i$에 대응시킨다($1\leq i\leq k$). 이때
$\operatorname{Coker}u$는
$(S(n))_{-n}+\cdots+(S(n))_{-1}$과 동일시되므로, $u$는 요구한 조건을
충족한다.
\end{proof}

\begin{corollary}[2.7.10]
\phantomsection
\label{II.2.7.10-ko}
(\hyperref[II.2.7.5-ko]{2.7.5})의 가정 아래에서, $\mathcal{F}$를 유한 생성
준연접 $\mathcal{O}_X$-가군이라 하자. 그러면 어떤 정수 $n_0$가 존재하여,
모든 $n\geq n_0$에 대하여 $\mathcal{F}$는
$(\mathcal{O}_X(-n))^k$ 꼴의 $\mathcal{O}_X$-가군의 몫과 동형이다
($k$는 $n$에 의존한다).
\end{corollary}

\begin{proposition}[2.7.11]
\phantomsection
\label{II.2.7.11-ko}
(\hyperref[II.2.7.5-ko]{2.7.5})의 가정이 성립한다고 하고, $M$을 등급
$S$-가군이라 하자. 그러면 다음이 성립한다.
\begin{enumerate}
  \item[\rm(i)] 표준 준동형
    $\widetilde{\alpha}:\widetilde{M}\to
    \widetilde{\Gamma_*(\widetilde{M})}$은 동형이다.
  \item[\rm(ii)] $\mathcal{G}$를 $\widetilde{M}$의 준연접
    $\mathcal{O}_X$-부분가군이라 하고, $N$을 $\alpha$에 의한
    $\Gamma_*(\mathcal{G})$의 역상인 $M$의 등급 $S$-부분가군이라 하자.
    그러면 $\widetilde{N}=\mathcal{G}$이다. 여기서
    (\hyperref[II.2.5.4-ko]{2.5.4})에 따라 $\widetilde{N}$을
    $\widetilde{M}$의 $\mathcal{O}_X$-부분가군으로 동일시한다.
\end{enumerate}
\end{proposition}

\begin{proof}
(\hyperref[II.2.7.5-ko]{2.7.5})에 따라
$\beta:\widetilde{\Gamma_*(\widetilde{M})}\to\widetilde{M}$가 동형이므로,
(\hyperref[II.2.6.5.1-ko]{2.6.5.1})에 따라 $\widetilde{\alpha}$는 그
역동형이다. 이로써 (i)가 따른다. $P$를 $\Gamma_*(\widetilde{M})$의 등급
부분가군 $\alpha(M)$이라 하자. $M\mapsto\widetilde{M}$은 완전함자이므로
(\hyperref[II.2.5.4-ko]{2.5.4}), $\widetilde{\alpha}$에 의한
$\widetilde{M}$의 상은 $\widetilde{P}$와 같다. 따라서 (i)에 따라
$\widetilde{P}=\widetilde{\Gamma_*(\widetilde{M})}$이다. 이제
$Q=\Gamma_*(\mathcal{G})\cap P$로 놓자. 앞의 결과와
(\hyperref[II.2.5.4-ko]{2.5.4})에 따라
$\widetilde{Q}=\widetilde{\Gamma_*(\mathcal{G})}$이므로,
(\hyperref[II.2.7.5-ko]{2.7.5})에 따라 $\beta$를 $\widetilde{Q}$에 제한한
사상은 이 $\mathcal{O}_X$-가군에서 $\mathcal{G}$ 위로 가는
\emph{동형}이다. 한편 $N$의 정의와
(\hyperref[II.2.5.4-ko]{2.5.4})에 따라, 동형 $\widetilde{\alpha}$를
$\widetilde{N}$에 제한한 사상은 $\widetilde{N}$에서
$\widetilde{Q}$ 위로 가는 동형이다. 따라서
(\hyperref[II.2.6.5.1-ko]{2.6.5.1})에 의해 결론이 따른다.
\end{proof}

\subsection{함자적 성질.}
\label{subsection:II.2.8-ko}

\begin{env}[2.8.1]
\phantomsection
\label{II.2.8.1-ko}
$S$, $S'$을 음이 아닌 차수로 등급이 주어진 두 환이라 하고,
$\varphi:S'\to S$를 등급환의 준동형이라 하자. $X=\operatorname{Proj}(S)$에서
$V_+(\varphi(S'_+))$의 여집합인 열린부분집합, 또는 같은 말로 $S'_+$의
동차 원소 $f'$ 전체에 대한 $D_+(\varphi(f'))$들의 합집합을
$G(\varphi)$로 나타내겠다. $\operatorname{Spec}(S)$에서
$\operatorname{Spec}(S')$로 가는 연속사상 ${}^a\varphi$
(\hyperref[I.1.2.1-ko]{I, 1.2.1})를 $G(\varphi)$에 제한하면
$G(\varphi)$에서 $\operatorname{Proj}(S')$로 가는 연속사상이 된다. 표기의
남용으로 이 사상도 ${}^a\varphi$로 나타내겠다. $f'\in S'_+$가 동차
원소이면
\[
\label{II.2.8.1.1-ko}
  {}^a\varphi^{-1}(D_+(f'))=D_+(\varphi(f'))
\tag{2.8.1.1}
\]
이다. 이는 ${}^a\varphi$가 $G(\varphi)$를 $\operatorname{Proj}(S')$로
보낸다는 사실과 (\hyperref[I.1.2.2.2-ko]{I, 1.2.2.2})에서 따른다. 한편
같은 표기 아래에서 준동형 $\varphi$는 등급환의 준동형
$S'_{f'}\to S_f$를 표준적으로 정의하며, 이를 차수 $0$인 원소들에
제한하면 다음 준동형을 얻는다.

\oldpage[II]{42}
$S'_{(f')}\to S_{(f)}$; 이를 $\varphi_{(f)}$로 나타내겠다. 이에 대응하여
(\hyperref[I.1.6.1-ko]{I, 1.6.1}) 아핀 스킴의 사상
$({}^a\varphi_{(f)},\widetilde{\varphi}_{(f)}):
\operatorname{Spec}(S_{(f)})\to\operatorname{Spec}(S'_{(f')})$을 얻는다.
$\operatorname{Spec}(S_{(f)})$를 $D_+(f)$ 위에서
$\operatorname{Proj}(S)$로부터 유도된 스킴과 표준적으로 동일시하면
(\hyperref[II.2.3.6-ko]{2.3.6}), 이로써 사상
$\Phi_f:D_+(f)\to D_+(f')$를 정의하였고, ${}^a\varphi_{(f)}$는
${}^a\varphi$를 $D_+(f)$에 제한한 사상과 동일시된다. 한편 $g'$이
$S'_+$의 두 번째 동차 원소이고 $g=\varphi(g')$이면 다음 도식이
가환임을 곧 알 수 있다.
\[
  \xymatrix{
    D_+(f)\ar[r]^{\Phi_f} & D_+(f')\\
    D_+(fg)\ar[u]\ar[r]_{\Phi_{fg}} & D_+(f'g')\ar[u]
  }
\]
이는 다음 도식이 가환하기 때문이다.
\[
  \xymatrix{
    S'_{(f')}\ar[r]^{\varphi_{(f)}}\ar[d]_{\omega_{f'g',f'}}
    & S_{(f)}\ar[d]^{\omega_{fg,f}}\\
    S'_{(f'g')}\ar[r]_{\varphi_{(fg)}} & S_{(fg)}
  }
\]
$G(\varphi)$의 정의와 (\hyperref[II.2.3.3.2-ko]{2.3.3.2})을 고려하면
따라서 다음을 알 수 있다.
\end{env}

\begin{proposition}[2.8.2]
\phantomsection
\label{II.2.8.2-ko}
등급환 준동형 $\varphi:S'\to S$가 주어지면, 유도 준스킴
$G(\varphi)$에서 $\operatorname{Proj}(S')$로 가는 유일한 사상
$({}^a\varphi,\widetilde{\varphi})$가 존재한다. 이 사상을
\emph{$\varphi$에 대응하는 사상}이라 하고 $\operatorname{Proj}(\varphi)$로
나타낸다. 모든 동차 원소 $f'\in S'_+$에 대하여, 이 사상을
$D_+(\varphi(f'))$에 제한한 사상은 $\varphi$가 정하는 준동형
$S'_{(f')}\to S_{(\varphi(f'))}$에 대응하는 사상과 일치한다.
\end{proposition}

\begin{proof}
앞의 표기에서 $f'\in S'_d$이면 다음 도식이 가환함에 유의하자.
\[
\label{II.2.8.2.1-ko}
  \xymatrix{
    S'_{(f')}\ar[r]^{\varphi_{(f)}}\ar[d]_{\sim}
    & S_{(f)}\ar[d]^{\sim}\\
    {S'}^{(d)}/(f'-1){S'}^{(d)}\ar[r]
    & S^{(d)}/(f-1)S^{(d)}
  }
\tag{2.8.2.1}
\]
여기서 세로 화살표들은
(\hyperref[II.2.2.5-ko]{2.2.5})의 동형들이다.
\end{proof}

\begin{corollary}[2.8.3]
\phantomsection
\label{II.2.8.3-ko}
\begin{enumerate}
  \item[\rm(i)] 사상 $\operatorname{Proj}(\varphi)$는 아핀이다.
  \item[\rm(ii)] $\operatorname{Ker}(\varphi)$가 멱영이면(특히
    $\varphi$가 단사이면), 사상 $\operatorname{Proj}(\varphi)$는
    우세 사상이다.
\end{enumerate}
\end{corollary}

\begin{proof}
(i)은 (\hyperref[II.2.8.2-ko]{2.8.2})와
(\hyperref[II.2.8.1.1-ko]{2.8.1.1})에서 곧바로 따른다. 한편
$\operatorname{Ker}(\varphi)$가 멱영이면, 모든 동차 원소
$f'\in S'_+$에 대하여 $f=\varphi(f')$로 둘 때
$\operatorname{Ker}(\varphi_f)$가 멱영임을 곧바로 확인할 수 있고,
따라서 $\operatorname{Ker}(\varphi_{(f)})$도 멱영이다. 그러므로 결론은
(\hyperref[II.2.8.2-ko]{2.8.2})와
(\hyperref[I.1.2.7-ko]{I, 1.2.7})에서 따른다.
\end{proof}

\oldpage[II]{43}
일반적으로 $\operatorname{Proj}(S)$에서 $\operatorname{Proj}(S')$로 가는
사상들 가운데 아핀이 아닌 것이 있으며, 따라서 이러한 사상은 등급환 준동형
$S'\to S$에서 나오지 않는다. 예를 들어 $A$가 체일 때의 구조 사상
$\operatorname{Proj}(S)\to\operatorname{Spec}(A)$이 그러하다. 여기서는
$\operatorname{Spec}(A)$를 $\operatorname{Proj}(A[T])$와 동일시한다
(\agref{II.3.1.7-ko}{3.1.7} 참조). 실제로 이는
(\hyperref[I.2.3.2-ko]{I, 2.3.2})에서 따른다.

\begin{env}[2.8.4]
\phantomsection
\label{II.2.8.4-ko}
$S''$을 음이 아닌 차수로 등급이 주어진 세 번째 환이라 하고,
$\varphi':S''\to S'$을 등급환 준동형이라 하며,
$\varphi''=\varphi\circ\varphi'$으로 놓자.
(\hyperref[II.2.8.1.1-ko]{2.8.1.1})과 공식
${}^a\varphi''=({}^a\varphi')\circ({}^a\varphi)$에 따라
$G(\varphi'')\subset G(\varphi)$임을 곧바로 확인할 수 있다. 또한
$\Phi$, $\Phi'$, $\Phi''$이 각각 $\varphi$, $\varphi'$, $\varphi''$에
대응하는 사상이면
$\Phi''=\Phi'\circ(\Phi|G(\varphi''))$이다.
\end{env}

\begin{env}[2.8.5]
\phantomsection
\label{II.2.8.5-ko}
$S$가 등급 $A$-대수이고 $S'$이 등급 $A'$-대수라고 하자. 또
$\psi:A'\to A$를 다음 도식이 가환하도록 하는 환 준동형이라 하자.
\[
\label{II.2.8.5.1-ko}
  \xymatrix{
    A'\ar[r]^{\psi}\ar[d] & A\ar[d]\\
    S'\ar[r]_{\varphi} & S
  }
\tag{2.8.5.1}
\]
그러면 $G(\varphi)$와 $\operatorname{Proj}(S')$을 각각
$\operatorname{Spec}(A)$와 $\operatorname{Spec}(A')$ 위의 스킴으로 볼 수
있다. $\Phi$와 $\Psi$가 각각 $\varphi$와 $\psi$에 대응하는 사상이면
다음 도식은 가환한다.
\[
  \xymatrix{
    G(\varphi)\ar[r]^{\Phi}\ar[d] & \operatorname{Proj}(S')\ar[d]\\
    \operatorname{Spec}(A)\ar[r]_{\Psi} & \operatorname{Spec}(A')
  }
\]
실제로 $f=\varphi(f')$이고 $f'$이 $S'_+$의 동차 원소일 때, $\Phi$를
$D_+(f)$에 제한한 사상에 대해 이를 보이면 충분하다. 이 경우에는 다음
도식의 가환성에서 결론이 따른다.
\[
  \xymatrix{
    A'\ar[r]^{\psi}\ar[d] & A\ar[d]\\
    S'_{(f')}\ar[r]_{\varphi_{(f)}} & S_{(f)}
  }
\]
\end{env}

\begin{env}[2.8.6]
\phantomsection
\label{II.2.8.6-ko}
이제 $M$을 등급 $S$-가군이라 하고, 자명하게 등급이 매겨지는
$S'$-가군 $M_{[\varphi]}$를 생각하자. $f'$을 $S'_+$의 동차 원소라 하고
$f=\varphi(f')$이라 하자. 표준 동형
$(M_{[\varphi]})_{f'}\xrightarrow{\sim}(M_f)_{[\varphi_f]}$이 존재함을
알고 있으며(\hyperref[0.1.5.2-ko]{0, 1.5.2}), 이 동형이 차수를
보존함은 자명하다. 따라서 표준 동형
$(M_{[\varphi]})_{(f')}\xrightarrow{\sim}
(M_{(f)})_{[\varphi_{(f)}]}$이 존재한다. 이 동형에 표준적으로 대응하는
층의 동형
$\widetilde{M_{[\varphi]}}|D_+(f')
\xrightarrow{\sim}(\Phi_f)_*(\widetilde{M}|D_+(f))$
이 존재한다
(\hyperref[II.2.5.2-ko]{2.5.2} 및
\hyperref[I.1.6.3-ko]{I, 1.6.3}). 더욱이,

\oldpage[II]{44}
$g'$가 $S'_+$의 두 번째 동차 원소이고 $g=\varphi(g')$이면 다음 도식은
가환한다.
\[
  \xymatrix{
    (M_{[\varphi]})_{(f')}\ar[r]^{\sim}\ar[d]
    & (M_{(f)})_{[\varphi_{(f)}]}\ar[d]\\
    (M_{[\varphi]})_{(f'g')}\ar[r]^{\sim}
    & (M_{(fg)})_{[\varphi_{(fg)}]}
  }
\]
따라서 동형
\[
  \widetilde{M_{[\varphi]}}|D_+(f'g')
  \xrightarrow{\sim}(\Phi_{fg})_*(\widetilde{M}|D_+(fg))
\]
은 동형
$\widetilde{M_{[\varphi]}}|D_+(f')
\xrightarrow{\sim}(\Phi_f)_*(\widetilde{M}|D_+(f))$을
$D_+(f'g')$로 제한한 것임을 곧바로 알 수 있다. $\Phi_f$는 사상
$\Phi$를 $D_+(f)$로 제한한 것이므로,
(\hyperref[II.2.8.1.1-ko]{2.8.1.1})을 고려하고
$X'=\operatorname{Proj}(S')$으로 두면 다음을 얻는다.
\end{env}

\begin{proposition}[2.8.7]
\phantomsection
\label{II.2.8.7-ko}
$\mathcal{O}_{X'}$-가군 $\widetilde{M_{[\varphi]}}$에서
$\mathcal{O}_{X'}$-가군 $\Phi_*(\widetilde{M}|G(\varphi))$로 가는
함자적인 표준 동형이 존재한다.
\end{proposition}

이로부터 등급 $S'$-가군 $M'$에서 등급 $S$-가군 $M$으로 가는
$\varphi$-사상들의 집합에서 $\Phi$-사상
$\widetilde{M'}\to\widetilde{M}|G(\varphi)$들의 집합으로 가는 함자적인
표준 사상을 곧바로 얻는다. (\hyperref[II.2.8.4-ko]{2.8.4})의 표기
아래에서 $M''$가 등급 $S''$-가군이면, $\varphi$-사상 $M'\to M$과
$\varphi'$-사상 $M''\to M'$의 합성에
$\widetilde{M'}|G(\varphi')\to\widetilde{M}|G(\varphi'')$와
$\widetilde{M''}\to\widetilde{M'}|G(\varphi')$의 합성이 표준적으로
대응한다.

\begin{proposition}[2.8.8]
\phantomsection
\label{II.2.8.8-ko}
(\hyperref[II.2.8.1-ko]{2.8.1})의 가정 아래에서 $M'$을 등급
$S'$-가군이라 하자. $(\mathcal{O}_X|G(\varphi))$-가군
$\Phi^*(\widetilde{M'})$에서 $(\mathcal{O}_X|G(\varphi))$-가군
$\widetilde{M'\otimes_{S'}S}|G(\varphi)$로 가는 함자적인 표준 준동형
$\nu$가 존재한다. 아이디얼 $S'_+$가 $S'_1$에 의해 생성되면
$\nu$는 동형이다.
\end{proposition}

\begin{proof}
실제로 $f'\in S'_d$ ($d>0$)에 대하여 다음과 같은
$S_{(f)}$-가군의 함자적인 표준 준동형을 정의한다. 여기서
$f=\varphi(f')$이다.
\[
\label{II.2.8.8.1-ko}
  \nu_f:M'_{(f')}\otimes_{S'_{(f')}}S_{(f)}
  \longrightarrow(M'\otimes_{S'}S)_{(f)}
\tag{2.8.8.1}
\]
이는 준동형
$M'_{(f')}\otimes_{S'_{(f')}}S_{(f)}\to
M'_{f'}\otimes_{S'_{f'}}S_f$와 표준 동형
$M'_{f'}\otimes_{S'_{f'}}S_f\xrightarrow{\sim}(M'\otimes_{S'}S)_f$
(\hyperref[0.1.5.4-ko]{0, 1.5.4})을 합성하여 정의하며, 뒤의 동형이
차수를 보존한다는 점을 사용한다. 두 번째 원소 $g'\in S'_+$와
$g=\varphi(g')$에 대하여 $D_+(f)$에서 $D_+(fg)$로 가는 제한
준동형과 $\nu_f$들이 양립함을 곧바로 확인할 수 있다. 따라서
(\hyperref[I.1.6.5-ko]{I, 1.6.5})을 고려하면 준동형
\[
  \nu:\Phi^*(\widetilde{M'})
  \longrightarrow\widetilde{M'\otimes_{S'}S}|G(\varphi)
\]
가 정의된다. 두 번째 주장을 증명하려면 모든 $f'\in S'_1$에 대하여
$\nu_f$가 동형임을 보이면 충분하다. 실제로 이 경우 $G(\varphi)$는
$D_+(\varphi(f'))$들의 합집합이다. 먼저 $\mathbf{Z}$-쌍선형 사상
$M'_m\times S_n\to M'_{(f')}\otimes_{S'_{(f')}}S_{(f)}$을 정의하자.
이 사상은 $(x',s)$에 원소
$(x'/{f'}^m)\otimes(s/f^n)$을 대응시킨다. 여기서 $m<0$일 때에는
$x'/{f'}^m$이 ${f'}^{-m}x'/1$이라는 약속을 쓴다.

\oldpage[II]{45}
(\hyperref[II.2.5.13-ko]{2.5.13})의 증명에서와 같이 이 사상으로부터
가군의 쌍준동형
\[
  \eta_f:M'\otimes_{S'}S
  \longrightarrow M'_{(f')}\otimes_{S'_{(f')}}S_{(f)}.
\]
을 얻는다. 또한 어떤 $r>0$에 대하여
$f^r\sum_i(x'_i\otimes s_i)=0$이면, 이를
$\sum_i({f'}^rx'_i\otimes s_i)=0$으로도 쓸 수 있다. 따라서
(\hyperref[0.1.5.4-ko]{0, 1.5.4})에 의해
$\sum_i({f'}^rx'_i/{f'}^{m_i+r})\otimes(s_i/f^{n_i})=0$이다. 즉
$\eta_f(\sum_i x'_i\otimes s_i)=0$이므로 $\eta_f$는 다음과 같이
인수분해된다.
\[
  M'\otimes_{S'}S\longrightarrow(M'\otimes_{S'}S)_f
  \xrightarrow{\eta'_f}M'_{(f')}\otimes_{S'_{(f')}}S_{(f)}~;
\]
끝으로 $\eta'_f$와 $\nu_f$가 서로 역인 동형임을 확인할 수 있다.

특히 (\hyperref[II.2.1.2.1-ko]{2.1.2.1})에서 모든
$n\in\mathbf{Z}$에 대하여 표준 준동형
\[
\label{II.2.8.8.2-ko}
  \Phi^*(\mathcal{O}_{X'}(n))\xrightarrow{\sim}
  \mathcal{O}_X(n)|G(\varphi)
\tag{2.8.8.2}
\]
을 얻는다.
\end{proof}

\begin{env}[2.8.9]
\phantomsection
\label{II.2.8.9-ko}
$A$, $A'$을 두 환이라 하고, $\psi:A'\to A$를 환 준동형이라 하자. 이
준동형은 사상
$\Psi:\operatorname{Spec}(A)\to\operatorname{Spec}(A')$를 정의한다.
$S'$를 음이 아닌 차수로 등급이 주어진 $A'$-대수라 하고
$S=S'\otimes_{A'}A$로 두자. 이는 자명하게 $S'_n\otimes_{A'}A$들을
동차 성분으로 갖는 등급 $A$-대수이다. 그러면 사상
$\varphi:s'\mapsto s'\otimes1$은 도식
(\hyperref[II.2.8.5.1-ko]{2.8.5.1})을 가환하게 하는 등급환
준동형이다. 여기서는 $S_+$가 $\varphi(S'_+)$에 의해 생성되는
$A$-가군이므로 $G(\varphi)=\operatorname{Proj}(S)=X$이다. 따라서
$X'=\operatorname{Proj}(S')$로 두면 다음 가환 도식을 얻는다.
\[
\label{II.2.8.9.1-ko}
  \xymatrix{
    X\ar[r]^{\Phi}\ar[d]_p & X'\ar[d]\\
    Y\ar[r]_{\Psi} & Y'
  }
\tag{2.8.9.1}
\]

이제 $M'$을 등급 $S'$-가군이라 하고
$M=M'\otimes_{A'}A=M'\otimes_{S'}S$로 두자. 이 조건 아래에서는
다음이 성립한다.
\end{env}

\begin{proposition}[2.8.10]
\phantomsection
\label{II.2.8.10-ko}
도식 (\hyperref[II.2.8.9.1-ko]{2.8.9.1})은 스킴 $X$를 곱
$X'\times_{Y'}Y$와 동일시한다. 또한 표준 준동형
$\nu:\Phi^*(\widetilde{M'})\to\widetilde{M}$
(\hyperref[II.2.8.8-ko]{2.8.8})은 동형이다.
\end{proposition}

\begin{proof}
첫째 주장은 $S'_+$의 모든 동차 원소 $f'$에 대하여
$f=\varphi(f')$로 둘 때, $\Phi$와 $p$를 $D_+(f)$에 제한한 사상들이
이 스킴을 곱 $D_+(f')\times_{Y'}Y$와 동일시함을 보이면 증명된다
(\hyperref[I.3.2.6.2-ko]{I, 3.2.6.2}). 다시 말해
(\hyperref[I.3.2.2-ko]{I, 3.2.2})에 따라 $S_{(f)}$가
$S'_{(f')}\otimes_{A'}A$와 표준적으로 동일시됨을 보이면 충분하다.
이는 차수를 보존하는 표준 동형
$S_f\xrightarrow{\sim}S'_{f'}\otimes_{A'}A$가 존재하므로 곧바로
따른다(\hyperref[0.1.5.4-ko]{0, 1.5.4}). 둘째 주장은 앞에서 본 바에
따라 $M'_{(f')}\otimes_{S'_{(f')}}S_{(f)}$가
$M'_{(f')}\otimes_{A'}A$와 동일시되고, 후자가 $M_{(f)}$와 동일시되는
사실에서 따른다. 실제로 $M_f$는 차수를 보존하는 동형에 의해
$M'_{f'}\otimes_{A'}A$와 표준적으로 동일시된다.
\end{proof}

\begin{corollary}[2.8.11]
\phantomsection
\label{II.2.8.11-ko}
모든 정수 $n\in\mathbf{Z}$에 대하여 $\widetilde{M(n)}$은
$\Phi^*(\widetilde{M'(n)})
=\widetilde{M'(n)}\otimes_{Y'}\mathcal{O}_Y$와 동일시된다. 특히
$\mathcal{O}_X(n)=\Phi^*(\mathcal{O}_{X'}(n))
=\mathcal{O}_{X'}(n)\otimes_{Y'}\mathcal{O}_Y$이다.
\end{corollary}

\begin{proof}
이는 (\hyperref[II.2.8.10-ko]{2.8.10})과
(\hyperref[II.2.5.15-ko]{2.5.15})에서 따른다.
\end{proof}

\oldpage[II]{46}

\begin{env}[2.8.12]
\phantomsection
\label{II.2.8.12-ko}
(\hyperref[II.2.8.9-ko]{2.8.9})의 가정 아래에서
$f'\in S'_d$ ($d>0$) 및 $f=\varphi(f')$에 대하여 다음 도식은
가환한다.
\[
\label{II.2.8.12.1-ko}
  \xymatrix{
    M'_{(f')}\ar[r]^{\sim}\ar[d]
    & M'^{(d)}/(f'-1)M'^{(d)}\ar[d]\\
    M_{(f)}\ar[r]^{\sim}
    & M^{(d)}/(f-1)M^{(d)}
  }
\tag{2.8.12.1}
\]
(\hyperref[II.2.2.5-ko]{2.2.5} 참조).
\end{env}

\begin{env}[2.8.13]
\phantomsection
\label{II.2.8.13-ko}
(\hyperref[II.2.8.9-ko]{2.8.9})의 표기와 가정을 그대로 두고,
$\mathcal{F}'$를 $\mathcal{O}_{X'}$-가군이라 하자.
$\mathcal{F}=\Phi^*(\mathcal{F}')$로 두면,
(\hyperref[II.2.8.11-ko]{2.8.11})과
(\hyperref[0.4.3.3-ko]{0, 4.3.3})에 의해 모든 $n\in\mathbf{Z}$에
대하여 $\mathcal{F}(n)=\Phi^*(\mathcal{F}'(n))$이다. 따라서
(\hyperref[0.3.7.1-ko]{0, 3.7.1})에 의해 표준 준동형
\[
  \Gamma(\rho):\Gamma(X',\mathcal{F}'(n))
  \longrightarrow\Gamma(X,\mathcal{F}(n))
\]
이 존재하며, 이로부터 등급가군의 표준 쌍준동형
\[
  \Gamma_\bullet(\mathcal{F}')\longrightarrow
  \Gamma_\bullet(\mathcal{F}).
\]
을 얻는다.

아이디얼 $S_+$가 $S_1$에 의해 생성된다고 가정하고,
$\mathcal{F}'=\widetilde{M'}$라 하자. 그러면
$M=M'\otimes_{A'}A$로 둘 때 $\mathcal{F}=\widetilde{M}$이다.
$f'$가 $S'_+$의 동차 원소이고 $f=\varphi(f')$이면,
$M_{(f)}=M'_{(f')}\otimes_{A'}A$임을 이미 보았으며 다음 도식은
가환한다.
\[
  \xymatrix{
    M'_0\ar[r]\ar[d]
    & M'_{(f')}\ar[d]
    & =\Gamma(D_+(f'),\widetilde{M'})\\
    M_0\ar[r]
    & M_{(f)}
    & =\Gamma(D_+(f),\widetilde{M})
  }
\]
이 관찰과 준동형 $\alpha$의 정의
(\hyperref[II.2.6.2-ko]{2.6.2})에서 다음 도식이 가환함을 곧바로
얻는다.
\[
\label{II.2.8.13.1-ko}
  \xymatrix{
    M'\ar[r]^{\alpha_{M'}}\ar[d]
    & \Gamma_\bullet(\widetilde{M'})\ar[d]\\
    M\ar[r]_{\alpha_M}
    & \Gamma_\bullet(\widetilde{M})
  }
\tag{2.8.13.1}
\]
마찬가지로 다음 도식도 가환한다.
\[
\label{II.2.8.13.2-ko}
  \xymatrix{
    \widetilde{\Gamma_\bullet(\mathcal{F}')}\ar[r]^{\beta_{\mathcal{F}'}}
      \ar[d]
    & \mathcal{F}'\ar[d]\\
    \widetilde{\Gamma_\bullet(\mathcal{F})}\ar[r]_{\beta_{\mathcal{F}}}
    & \mathcal{F}
  }
\tag{2.8.13.2}
\]
(오른쪽 세로 화살표는 표준 $\Phi$-사상
$\mathcal{F}'\to\Phi^*(\mathcal{F}')=\mathcal{F}$이다.)
\end{env}

\oldpage[II]{47}

\begin{env}[2.8.14]
\phantomsection
\label{II.2.8.14-ko}
계속해서 (\hyperref[II.2.8.9-ko]{2.8.9})의 가정과 표기법을 유지하고,
$N'$을 또 하나의 등급 $S'$-가군이라 하며
$N=N'\otimes_{A'}A$라 하자. 표준 쌍준동형 $M'\to M$, $N'\to N$은
(표준 환 준동형 $S'\to S$에 상대적인) 쌍준동형
$M'\otimes_{S'}N'\to M\otimes_S N$을 주며, 따라서 차수~$0$의
$S$-준동형
\[
  (M'\otimes_{S'}N')\otimes_{A'}A\longrightarrow M\otimes_S N,
\]
도 준다는 것은 곧바로 알 수 있다. 이에
(\hyperref[II.2.8.10-ko]{2.8.10})을 고려하여 다음과 같은
$\mathcal{O}_X$-가군의 준동형이 대응한다.
\[
\label{II.2.8.14.1-ko}
  \Phi^*(\widetilde{M'\otimes_{S'}N'})
  \longrightarrow\widetilde{M\otimes_S N}.
\tag{2.8.14.1}
\]

더 나아가 다음 도식이 가환함을 곧바로 확인할 수 있다.
\[
\label{II.2.8.14.2-ko}
  \xymatrix{
    \Phi^*(\widetilde{M'}\otimes_{\mathcal{O}_{X'}}\widetilde{N'})
      \ar[r]^{\sim}\ar[d]_{\Phi^*(\lambda)}
    & \widetilde{M}\otimes_{\mathcal{O}_X}\widetilde{N}
      \ar[d]^{\lambda}
    & =\Phi^*(\widetilde{M'})\otimes_{\mathcal{O}_X}
      \Phi^*(\widetilde{N'})\\
    \Phi^*(\widetilde{M'\otimes_{S'}N'})\ar[r]
    & \widetilde{M\otimes_S N}
  }
\tag{2.8.14.2}
\]
여기서 첫째 가로줄은 표준 동형
(\hyperref[0.4.3.3-ko]{0, 4.3.3})이다. 아이디얼 $S'_+$가
$S'_1$에 의해 생성되면, $S_+$가 $S_1$에 의해 생성됨은 명백하다.
이 경우 (\hyperref[II.2.8.14.2-ko]{2.8.14.2})의 두 세로 화살표는
동형이며(\hyperref[II.2.5.13-ko]{2.5.13}), 따라서
(\hyperref[II.2.8.14.1-ko]{2.8.14.1})도 동형이다.

마찬가지로, 차수 $k$의 준동형 $u'$에 역시 차수 $k$인 준동형
$u'\otimes1$을 대응시키면 표준 쌍준동형
$\operatorname{Hom}_{S'}(M',N')\to\operatorname{Hom}_S(M,N)$을 얻는다.
이로부터 다시 차수~$0$의 $S$-준동형
\[
  (\operatorname{Hom}_{S'}(M',N'))\otimes_{A'}A
  \longrightarrow\operatorname{Hom}_S(M,N)
\]
을 얻으며, 따라서 $\mathcal{O}_X$-가군의 준동형
\[
\label{II.2.8.14.3-ko}
  \Phi^*(\widetilde{\operatorname{Hom}_{S'}(M',N')})
  \longrightarrow\widetilde{\operatorname{Hom}_S(M,N)}.
\tag{2.8.14.3}
\]
을 얻는다.

더 나아가 다음 도식은 가환한다.
\[
  \xymatrix{
    \Phi^*(\widetilde{\operatorname{Hom}_{S'}(M',N')})\ar[r]
      \ar[d]_{\Phi^*(\mu)}
    & \widetilde{\operatorname{Hom}_S(M,N)}\ar[d]^{\mu}\\
    \Phi^*(\mathcal{H}om_{\mathcal{O}_{X'}}
      (\widetilde{M'},\widetilde{N'}))\ar[r]
    & \mathcal{H}om_{\mathcal{O}_X}(\widetilde{M},\widetilde{N})
  }
\]
여기서 둘째 가로줄은 표준 준동형
(\hyperref[0.4.4.6-ko]{0, 4.4.6})이다.
\end{env}

\oldpage[II]{48}

\begin{env}[2.8.15]
\phantomsection
\label{II.2.8.15-ko}
표기법과 가정이 (\hyperref[II.2.8.1-ko]{2.8.1})의 것이라 하자.
(\hyperref[II.2.4.7-ko]{2.4.7})에 의하면, $S_0$과 $S'_0$을
$\mathbf{Z}$로 바꾸고 $\varphi_0$을 항등사상으로 바꾼 뒤, 각각 $S$와
$S'$을 $S^{(d)}$와 ${S'}^{(d)}$로 바꾸고($d>0$) $\varphi$를
$S^{(d)}$에 대한 제한 $\varphi^{(d)}$로 바꾸어도 사상 $\Phi$는 동형까지
변하지 않는다.
\end{env}

\subsection{$\operatorname{Proj}(S)$ 스킴의 닫힌 부분스킴.}
\label{subsection:II.2.9-ko}

\begin{env}[2.9.1]
\phantomsection
\label{II.2.9.1-ko}
$\varphi:S\to S'$이 등급환의 준동형이라고 하자. 어떤 정수 $n$이 존재하여
$k\geq n$일 때 $\varphi_k:S_k\to S'_k$가 전사이면(각각 단사이면,
전단사이면), $\varphi$를 (TN)-\emph{전사}(각각 (TN)-\emph{단사},
(TN)-\emph{전단사})라고 한다. 그러면
(\hyperref[II.2.8.15-ko]{2.8.15})에 따라 $\Phi$의 연구는 $\varphi$가
\emph{전사}(각각 \emph{단사}, \emph{전단사})인 경우로 환원된다.
$\varphi$가 (TN)-전단사라고 하는 대신, 이때 이를
(TN)-\emph{동형사상}이라고도 한다.
\end{env}

\begin{proposition}[2.9.2]
\phantomsection
\label{II.2.9.2-ko}
$S$를 음이 아닌 차수로 등급이 주어진 환이라 하고
$X=\operatorname{Proj}(S)$라 하자.
\begin{enumerate}
  \item[\rm(i)] $\varphi:S\to S'$이 등급환의 (TN)-전사 준동형이면, 이에
    대응하는 사상 $\Phi$ (\hyperref[II.2.8.1-ko]{2.8.1})는
    $\operatorname{Proj}(S')$ 전체에서 정의되며
    $\operatorname{Proj}(S')$에서 $X$로 가는 닫힌 몰입 사상이다.
    $\mathfrak{J}$가 $\varphi$의 핵이면, $\Phi$에 결부된 $X$의 닫힌
    부분스킴은 $\mathcal{O}_X$의 준연접 아이디얼층
    $\widetilde{\mathfrak{J}}$로 정의된다.
  \item[\rm(ii)] 더 나아가 아이디얼 $S_+$가 차수~$1$인 유한 개의 동차
    원소로 생성된다고 가정하자. $X'$을 $\mathcal{O}_X$의 준연접
    아이디얼층 $\mathcal{J}$로 정의되는 $X$의 닫힌 부분스킴이라 하자.
    $\mathfrak{J}$를 표준 준동형
    $\alpha:S\to\Gamma_\bullet(\mathcal{O}_X)$
    (\hyperref[II.2.6.2-ko]{2.6.2})에 의한
    $\Gamma_\bullet(\mathcal{J})$의 역상인 $S$의 등급 아이디얼이라 하고,
    $S'=S/\mathfrak{J}$라 놓자. 그러면 $X'$은 등급환의 표준 준동형
    $S\to S'$에 대응하는 닫힌 몰입 사상
    $\operatorname{Proj}(S')\to X$에 결부된 부분스킴이다.
\end{enumerate}
\end{proposition}

\begin{proof}
\begin{enumerate}
  \item[\rm(i)] $\varphi$가 전사라고 가정해도 된다
    (\hyperref[II.2.9.1-ko]{2.9.1}). 가정에 따라 $\varphi(S_+)$가
    $S'_+$를 생성하므로 $G(\varphi)=\operatorname{Proj}(S')$이다.
    한편 둘째 주장은 $X$ 위에서 국소적으로 확인할 수 있다. 따라서
    $f$를 $S_+$의 동차 원소라 하고 $f'=\varphi(f)$라 놓자.
    $\varphi$는 등급환의 전사 준동형이므로
    $\varphi_{(f')}:S_{(f)}\to S'_{(f')}$가 전사이고 그 핵이
    $\mathfrak{J}_{(f)}$임을 곧바로 확인할 수 있다. 이로써 (i)가
    증명된다 (\hyperref[I.4.2.3-ko]{I, 4.2.3}).
  \item[\rm(ii)] (i)에 따라, 표준 단사 준동형
    $j:\mathfrak{J}\to S$에서 유도되는 준동형
    $\widetilde{j}:\widetilde{\mathfrak{J}}\to\mathcal{O}_X$가
    $\widetilde{\mathfrak{J}}$에서 $\mathcal{J}$로 가는 동형임을
    확인하면 충분하다. 이는 (\hyperref[II.2.7.11-ko]{2.7.11})에서
    따른다.
\end{enumerate}
\end{proof}

$\mathfrak{J}$는
$\widetilde{j}(\widetilde{\mathfrak{J}'})=\mathcal{J}$를 만족하는
$S$의 등급 아이디얼 $\mathfrak{J}'$들 가운데 \emph{가장 큰} 것이다.
실제로 정의로 거슬러 올라가 곧바로 확인하면
(\hyperref[II.2.6.2-ko]{2.6.2}), 이 관계에서
$\alpha(\mathfrak{J}')\subset\Gamma_\bullet(\mathcal{J})$가 따른다.

\begin{corollary}[2.9.3]
\phantomsection
\label{II.2.9.3-ko}
(\hyperref[II.2.9.2-ko]{2.9.2}, (i))의 가정이 성립하고, 더 나아가
아이디얼 $S_+$가 $S_1$로 생성된다고 가정하자. 그러면 모든
$n\in\mathbf{Z}$에 대하여 $\Phi^*(\widetilde{S(n)})$은
$\widetilde{S'(n)}$과 표준적으로 동형이고, 따라서 모든
$\mathcal{O}_X$-가군 $\mathcal{F}$에 대하여
$\Phi^*(\mathcal{F}(n))$은 $\Phi^*(\mathcal{F})(n)$과 표준적으로
동형이다.
\end{corollary}

\begin{proof}
이는 (\hyperref[II.2.8.8-ko]{2.8.8})의 특수한 경우이며,
(\hyperref[II.2.5.10.2-ko]{2.5.10.2})를 고려한다.
\end{proof}

\begin{corollary}[2.9.4]
\phantomsection
\label{II.2.9.4-ko}
(\hyperref[II.2.9.2-ko]{2.9.2}, (ii))의 가정이 성립한다고 하자.
$X$의 닫힌 부분준스킴 $X'$이 정역 스킴이기 위한 필요충분조건은 등급
아이디얼 $\mathfrak{J}$가 $S$의 소아이디얼인 것이다.
\end{corollary}

\oldpage[II]{49}

\begin{proof}
$X'$은 $\operatorname{Proj}(S/\mathfrak{J})$와 동형이므로,
(\hyperref[II.2.4.4-ko]{2.4.4})에 따라 조건은 충분하다. 필요성을
보이기 위하여 완전열
$0\to\mathcal{J}\to\mathcal{O}_X\to\mathcal{O}_X/\mathcal{J}$를
생각하자. 이로부터 완전열
\[
  0\longrightarrow\Gamma_\bullet(\mathcal{J})
  \longrightarrow\Gamma_\bullet(\mathcal{O}_X)
  \longrightarrow\Gamma_\bullet(\mathcal{O}_X/\mathcal{J}).
\]
을 얻는다. $f\in S_m$, $g\in S_n$이고
$\Gamma_\bullet(\mathcal{O}_X/\mathcal{J})$에서
$\alpha_{n+m}(fg)$의 상이 영일 때, $\alpha_m(f)$와
$\alpha_n(g)$의 상 가운데 하나가 영임을 보이면 충분하다. 그런데
정의에 따라 이 상들은 정역 스킴 $X'$ 위의 가역
$(\mathcal{O}_X/\mathcal{J})$-가군
$\mathcal{L}=(\mathcal{O}_X/\mathcal{J})(m)$과
$\mathcal{L}'=(\mathcal{O}_X/\mathcal{J})(n)$의 단면들이다. 가정에
의하면 이 두 단면의 곱은 $\mathcal{L}\otimes\mathcal{L}'$에서 영이다
((\hyperref[II.2.9.3-ko]{2.9.3}) 및
(\hyperref[II.2.5.14.1-ko]{2.5.14.1})). 따라서
(\hyperref[I.7.4.4-ko]{I, 7.4.4})에 따라 둘 중 하나가 영이다.
\end{proof}

\begin{corollary}[2.9.5]
\phantomsection
\label{II.2.9.5-ko}
$A$를 환, $M$을 $A$-가군, $S$를 차수~$1$인 동차 원소들의 집합
$S_1$로 생성되는 등급 $A$-대수라 하자. $u:M\to S_1$을 $A$-가군의
전사 준동형이라 하고, $\overline{u}:\mathbf{S}(M)\to S$를 $u$를
확장하는 $M$의 대칭대수 $\mathbf{S}(M)$에서 $S$로 가는
($A$-대수의) 준동형이라 하자. 그러면 $\overline{u}$에 대응하는 사상은
$\operatorname{Proj}(S)$에서 $\operatorname{Proj}(\mathbf{S}(M))$로
가는 닫힌 몰입 사상이다.
\end{corollary}

\begin{proof}
가정에 따라 $\overline{u}$는 전사이므로
(\hyperref[II.2.9.2-ko]{2.9.2})를 적용하면 된다.
\end{proof}

\end{document}
